\documentclass[11pt,oneside,openany]{book}
\usepackage[T1]{fontenc}
\usepackage[utf8]{inputenc}
\usepackage{lmodern}
\usepackage{microtype}
\usepackage[a4paper,inner=30mm,outer=27mm,top=27mm,bottom=28mm]{geometry}
\usepackage{amsmath,amssymb,amsthm,mathtools}
\usepackage{mathrsfs}
\usepackage{booktabs,longtable,array,enumitem}
\usepackage[dvipsnames]{xcolor}
\usepackage[colorlinks=true,linkcolor=MidnightBlue,citecolor=MidnightBlue,urlcolor=MidnightBlue]{hyperref}
\usepackage{bookmark}
\numberwithin{equation}{chapter}
\newtheorem{theorem}{Theorem}[chapter]
\newtheorem{proposition}[theorem]{Proposition}
\newtheorem{lemma}[theorem]{Lemma}
\newtheorem{corollary}[theorem]{Corollary}
\theoremstyle{definition}
\newtheorem{definition}[theorem]{Definition}
\newtheorem{assumption}[theorem]{Assumption}
\theoremstyle{remark}
\newtheorem{remark}[theorem]{Remark}
\newtheorem{example}[theorem]{Example}
\newcommand{\RR}{\mathbb R}
\newcommand{\CC}{\mathbb C}
\newcommand{\ZZ}{\mathbb Z}
\newcommand{\NN}{\mathbb N}
\newcommand{\gLie}{\mathfrak g}
\DeclareMathOperator{\Ad}{Ad}
\DeclareMathOperator{\tr}{tr}
\newcommand{\Conf}{\operatorname{Conf}}
\setcounter{tocdepth}{1}
\setlist{itemsep=3pt,topsep=5pt}
\emergencystretch=2em
\hypersetup{pdftitle={Four-dimensional Yang--Mills theory: A constructive approach},pdfauthor={Working exposition from the supplied manuscripts}}

\title{Four-dimensional Yang--Mills theory\\[6pt]
 \large A constructive approach}
\author{Working exposition from the supplied manuscripts}
\date{Draft of 19 November 2026}
\begin{document}
\frontmatter
\hypersetup{pageanchor=false}
\maketitle
\hypersetup{pageanchor=true}
\chapter*{Main theorem}
\addcontentsline{toc}{chapter}{Main theorem}

The construction starts from Wilson lattice gauge theory on periodic
four-dimensional lattices. Its continuum observables are local
gauge-invariant fields. Their Euclidean correlations determine the
relativistic theory described below.

\noindent\textbf{Main theorem (continuum Yang--Mills fields).}
Let $G$ be a compact matrix group with simple nonabelian real Lie algebra
$\gLie$. Fix a faithful unitary representation and its normalized trace
inner product, a positive penalty parameter for the block transformation,
and a sufficiently large inverse coupling at the coarsest scale.
There is a choice of lattice spacings tending to zero, admissible physical
periods tending to infinity, bare Wilson couplings, and normalizations of
local observables with the following properties.

\begin{enumerate}[label=\textup{(\roman*)},leftmargin=*]
\item \textit{Continuum limit.} At each fixed admissible period the
correlations converge along the full sequence of lattice spacings. From
any increasing sequence of these periods one can extract a single
subsequence along which the correlations of all orders converge to smooth
functions at pairwise distinct points in $\RR^4$. The same subsequence
works for all derivatives. The resulting local fields form nested
finite-dimensional spaces, obtained by imposing successive bounds on the
scaling dimension of the polynomial labels (curvature counts as two, each
derivative as one, and dimensions add in products). They contain the identity and
admit commuting spacetime derivatives. At each bound the field space and
its gauge-invariant free Maxwell counterpart have the same dimension and
carry isomorphic representations of $O(4)$.

\item \textit{Relativistic theory.} The limiting correlations are averages,
with respect to a probability measure of full support on a compact metrizable
space of phases, of Euclidean covariant, reflection-positive correlations.
Every phase has the same field spaces and reconstructs a local,
Poincar\'e-covariant theory on four-dimensional Minkowski space. Its
Wightman distributions are tempered, its energy--momentum spectrum lies
in the closed forward light cone, and its invariant vacuum is unique up
to scalar multiples. The weak commutant of the fields consists of scalars.

\item \textit{Stress tensor and local products.} The phases share a real,
symmetric, conserved stress tensor whose local charges generate translations
and rotations through the Ward identities. Products of nearby fields admit
an operator product expansion to each fixed finite order, with coefficient
and remainder bounds uniform in the phase. A fixed Taylor subtraction
also extends the separated Euclidean correlations across coincident points
as distributions of locally finite order.
\end{enumerate}

Vacuum uniqueness here is a statement within each phase. Uniqueness of
the infinite-volume limit and a positive mass gap remain open in this construction.

\clearpage
\section*{Where the theorem is proved}
\addcontentsline{toc}{section}{Where the theorem is proved}

Theorem~\ref{ch19:conditional} (p.~\pageref{ch19:conditional}) assembles
the conclusions for one and the same family of Wilson measures. The main
steps of its proof can be read separately as follows.

\paragraph{The two limits in (i).}
Theorems~\ref{ch10:finite-history} and~\ref{ch10:convergence} tune the
bare coupling and prove convergence of the renormalization group trajectory
at each fixed scale. With the bounds on observable insertions from
Chapter~\ref{ch11:insertions}, Theorem~\ref{ch12:uv-theorem} proves the
full cutoff limit at fixed volume. Section~\ref{ch13:extraction} then
extracts one subsequence of volumes for all separated correlations and
their derivatives. Proposition~\ref{ch13:rank} identifies the dimensions
of the field spaces. Section~\ref{ch13:calculus} constructs their inclusions
and derivatives, and Theorem~\ref{ch14:rotations} identifies their $O(4)$
representations.

\paragraph{Phases and relativistic fields in (ii).}
Section~\ref{ch15:characters} gives the integral decomposition into phases.
Sections~\ref{ch15:phase-symmetry} and~\ref{ch15:fields} establish symmetry,
reflection positivity and equality of the field spaces in every phase.
Section~\ref{ch16:boundary-section} constructs the tempered boundary
values. Theorem~\ref{ch16:fields} reconstructs the local fields and proves
the spectrum condition. Theorem~\ref{ch16:vacuum} proves vacuum uniqueness
in every phase, and Proposition~\ref{ch16:irreducible} determines the
weak commutant.

\paragraph{The stress tensor and local products in (iii).}
Proposition~\ref{ch14:normalized-stress} fixes the normalization of the
conserved tensor by its Euclidean Ward identities. Section~\ref{ch16:charges-section}
continues them to the relativistic charge law~\eqref{ch16:charge-law}.
Theorem~\ref{ch18:ope-theorem} gives the finite-order expansion and its
remainder estimate. The distinct assertion about extension across the
collision diagonals is Theorem~\ref{ch18:extension-theorem}.

\paragraph{A further conclusion in one selected phase.}
Let $\Theta$ be the trace of the Euclidean stress tensor minus its
expectation in the phase selected in Section~\ref{ch15:selection}, and
let $G_\Theta(r)$ be its two-point function at separation $r>0$ in that
phase. The additional estimates of Chapter~\ref{ch17:trace} give
\[
 \lim_{r\downarrow0} r^8\log^2(1/r)\,G_\Theta(r)
     =\frac{3\dim_{\RR}\gLie}{\pi^4}.
\]
Theorem~\ref{ch17:all-radii} proves the limit for arbitrary radii tending
to zero. Section~\ref{ch17:exclusion} specifies the finite class of
mass-shell correlation functions excluded by this asymptotic.

\chapter*{About this draft}
\addcontentsline{toc}{chapter}{About this draft}

This book is being written from three supplied accounts of a proposed
construction of four-dimensional Yang--Mills theory, together with the
detailed outline accompanying Scott's account. The aim is a readable
mathematical exposition with definitions at their first use, consistent
notation, and precise separation of standard arguments from estimates
specific to the construction.

The present draft contains the nineteen planned chapters, including the
construction, reconstruction, quantitative trace law and finite-order
operator products. The accompanying blueprint gives bounded source
packets and a common notation registry; the audits record the detailed
hypothesis matching and corrections made while writing. This is a
working exposition, not an independent certification of the supplied
manuscripts or a substitute for further mathematical review.

The intended construction first removes the lattice spacing at each fixed
physical volume and then extracts an infinite-volume subsequence. Its
limiting correlations describe gauge-invariant local observables.
Reconstruction, the stress tensor and operator products are developed
after the continuum field spaces. No positive mass gap is asserted. Nor does the plan assert
uniqueness of the infinite-volume limit or independence of all choices
in the blocking procedure.

The reader is assumed to know basic probability, analysis, compact Lie
groups and the general aims of constructive quantum field theory.
Standard results are cited where used; their application is explained
without reproducing textbook proofs. Source-specific claims whose
hypotheses have not yet been matched are not treated as proved merely
because the source labels them theorems.

\chapter*{Reading guide}
\addcontentsline{toc}{chapter}{Reading guide}

The construction has two limiting operations. First the lattice spacing
is removed while the physical volume is held fixed. Then the volume
increases along a subsequence. The first operation requires a stability
argument, not just compactness. The second uses compactness only after
bounds uniform in volume have been established. Keeping these operations
separate also fixes the meaning of the local observables throughout the
book.

A lattice expression with the correct classical limit is not yet a
normalized quantum field. Integration can mix it with other expressions
of the same or lower degree, including the constant. We therefore begin
with finite-dimensional spaces of polynomial labels, construct their
local lattice representatives, and later determine their normalization
by transporting all coefficients together. Fields are finally defined
by equality of all separated correlation responses. This prevents a
formal equation of motion from being identified prematurely with a
vanishing interacting field.

The main analytic work is the block step. Its quadratic part determines
a Gaussian fluctuation field and a stationary background. The nonlinear
integral is separated into regions of small curvature, where Taylor and
Gaussian estimates apply, and complementary regions, where a positive
action cost is retained. Local terms are reorganized as polymer
activities: functions associated with connected supports, controlled by
weighted norms. The support weights must accommodate the full dependence
of an activity, including the cutoffs and local gauge choices used to
define it. This distinction matters when several insertions interact.

Once the block estimates are available, one scale schedule must satisfy
all fixed finite requirements. It is chosen slowly enough that each
stage's constants can be controlled before the next set of requirements
is introduced. A common trajectory then determines the couplings,
activities and observable normalizations for the entire Wilson family.
Observable insertions are tracked by commuting auxiliary variables whose
squares are zero. Coefficients involving several variables retain the
terms generated when insertions merge under blocking.

The continuum part uses this information in a different order from a
formal perturbative calculation. Finite families of free-field correlation functionals separate the
polynomial labels and control the dimension of the interacting field spaces. Coordinate
reflection positivity constructs Hilbert spaces and bounded slab
operators before continuous rotations are known. A conserved tensor,
local flux estimates and positivity then establish rotational Ward
identities and fix the stress normalization. Only then are phases
separated and relativistic fields reconstructed.

\begin{center}
\begin{tabular}{@{}p{.17\textwidth}p{.76\textwidth}@{}}
\toprule
Chapters & Role in the argument\\
\midrule
1--3 & Lattice conventions, local labels and exact microscopic equations;
free Maxwell correlations that separate polynomial labels.\\[5pt]
4--9 & Exact blocking, Gaussian kernels, large-field estimates,
local matching, polymer expansions and the nonlinear step.\\[5pt]
10--12 & One common trajectory, observable insertions and the full
ultraviolet limit at fixed physical volume.\\[5pt]
13--16 & Infinite-volume correlations, field quotients, Ward identities,
phases and relativistic reconstruction.\\[5pt]
17 & The quantitative trace coefficient and its all-radius consequence.\\[5pt]
18--19 & Finite operator products, specified contact extensions,
and assembly with the exact scope of the conclusions.\\
\bottomrule
\end{tabular}
\end{center}

The final two analytic refinements have different purposes. The trace
law needs a rate beyond convergence to the traceless free tensor; an
error that tends to zero may still be too large after multiplication by
the running coupling. An operator product expansion needs a signed
replacement estimate for a whole cluster, uniform under repeated copies.
Neither follows from the existence of separated correlations alone.
A conservative finite-order extension across collision diagonals,
by contrast, follows from polynomial separated bounds and a specified
Taylor subtraction.

The final assembly chapter collects the proved dependencies and
distinguishes the conclusions from further questions not addressed here. The
accompanying chapter audits provide source locations and review records;
they are working documents rather than prerequisites for reading the
mathematical exposition.

\tableofcontents
\mainmatter
\IfFileExists{chapters/01-lattice.tex}{% !TEX root = ../monograph.tex
\chapter{Lattice gauge theory and the continuum problem}
\label{ch01:lattice}

The starting point is a probability measure on finitely many group-valued
variables. Its definition requires neither gauge fixing nor a continuum
functional integral. The construction problem is to choose its coupling and
its local observables as the lattice is refined, and to control their
correlations uniformly enough to remove the cutoff.

Two limits have different roles. At a fixed physical volume, the lattice
spacing will tend to zero along a prescribed sequence. Only afterwards will
the volume tend to infinity, along a subsequence. Keeping these operations
separate is essential: compactness can produce a subsequence, but cannot by
itself give convergence along every ultraviolet cutoff.

\section{The gauge group and its normalization}
\label{ch01:gauge}

Fix a closed subgroup $G\subset U(n_G)$ in a specified faithful unitary
representation. Its real Lie algebra $\gLie$ consists of skew-Hermitian
matrices, and we assume that $\gLie$ is simple and nonabelian. The group itself
need not be connected. We use the invariant inner product
\begin{equation}
 \langle X,Y\rangle_{\gLie}=-\frac1{n_G}\operatorname{Re}\tr(XY),
 \qquad X,Y\in\gLie.
 \label{ch01:metric}
\end{equation}
The choice of representation is part of the normalization of the model,
not an inessential choice of matrix coordinates.

Write $d_{\gLie}=\dim_{\RR}\gLie$ and fix once and for all an arbitrary orthonormal basis $(e_\alpha)$
for \eqref{ch01:metric}. The adjoint Casimir is the positive number
\begin{equation}
 C_{\gLie}=-\frac1{d_{\gLie}}\sum_{\alpha=1}^{d_{\gLie}}
       \tr_{\gLie}\bigl(\operatorname{ad}(e_\alpha)^2\bigr).
 \label{ch01:casimir}
\end{equation}
Here $\operatorname{ad}(X)Y=[X,Y]$ and $\tr_{\gLie}$ is the real operator
trace. Invariance of the inner product makes $\operatorname{ad}(X)$
skew-adjoint. Simplicity implies the familiar identity
$-\tr_{\gLie}(\operatorname{ad}(X)\operatorname{ad}(Y))
=C_{\gLie}\langle X,Y\rangle_{\gLie}$; nonabelianity gives $C_{\gLie}>0$.
We shall use the constants
\begin{equation}
 b_0=\frac{11C_{\gLie}}{48\pi^2},\qquad
 \kappa_G=2b_0=\frac{11C_{\gLie}}{24\pi^2}.
 \label{ch01:slope}
\end{equation}
Their role in the construction will be to describe the leading logarithmic
growth of the inverse coupling. Defining these constants does not yet
establish that growth for the nonperturbative lattice model.

For $G=SU(n)$ in its defining representation, our metric is
$-\operatorname{Re}\tr(XY)/n$, and
\begin{equation}
 d_{\gLie}=n^2-1,\qquad C_{\gLie}=2n^2,\qquad
 \kappa_G=\frac{11n^2}{12\pi^2}.
 \label{ch01:su-n}
\end{equation}
The conventional Wilson parameter is $\beta=2n/g^2$. Thus
$\beta\sim\kappa_G\log(1/a)$ corresponds to
$g^{-2}\sim(11n/24\pi^2)\log(1/a)$.
This is the normalization of the one-loop coefficient of
Gross--Wilczek and Politzer~\cite{GrossWilczek1973,Politzer1973}.
The exact control of the error in that law is a later construction problem.

\section{Links, holonomies and the Wilson measure}
\label{ch01:wilson-section}

For an integer $s\ge1$, set
\[
 \Lambda_s=(\ZZ/s\ZZ)^4,\qquad
 E_s=\Lambda_s\times\{1,2,3,4\}.
\]
The element $(y,\mu)$ of $E_s$ represents the oriented link from $y$ to
$y+e_\mu$. A configuration is $U\in G^{E_s}$. The variable on the reversed
link is, by definition, $U(y,\mu)^{-1}$. A path holonomy is the ordered
product of its link variables, with this inverse convention on backward
steps. We use row transport: multiplying a row vector at the initial point
by that ordered product transports it to the final point.

The elementary positively oriented plaquette in the $(\mu,\nu)$-plane,
with $\mu<\nu$, has holonomy
\begin{equation}
 U_{y;\mu\nu}=U(y,\mu)U(y+e_\mu,\nu)
                 U(y+e_\nu,\mu)^{-1}U(y,\nu)^{-1}.
 \label{ch01:plaquette}
\end{equation}
Define the plaquette energy and the normalized Hilbert--Schmidt distance from the identity by
\begin{equation}
 w(u)=1-\frac1{n_G}\operatorname{Re}\tr u,
 \qquad
 \chi(u)=\frac{\|u-1\|_{\mathrm F}}{\sqrt{n_G}},
 \qquad \chi(u)^2=2w(u).
 \label{ch01:chord}
\end{equation}
Both functions are invariant under conjugation and inversion; they vanish
at the identity. The last identity follows directly from $u^*u=1$.

\begin{definition}[Wilson measure]
\label{ch01:wilson}
Let $dh$ be probability Haar measure on $G$. The Wilson energy and measure
on $\Lambda_s$ are
\begin{align}
 E_s^{\mathrm W}(U)&=\sum_{y\in\Lambda_s}\sum_{\mu<\nu}w(U_{y;\mu\nu}),
 \label{ch01:wilson-energy}\\
 d\mu_{s,\beta}(U)&=Z_{s,\beta}^{-1}e^{-\beta E_s^{\mathrm W}(U)}
                     \prod_{e\in E_s}dh(U_e),
 &Z_{s,\beta}&=\int e^{-\beta E_s^{\mathrm W}(U)}\prod_{e\in E_s}dh(U_e).
 \label{ch01:wilson-law}
\end{align}
\end{definition}
Every unoriented coordinate square is counted once, through $\mu<\nu$.
There is no additional factor of $1/2$ in \eqref{ch01:wilson-energy}.
The configuration space is compact and the integrand defining $Z_{s,\beta}$
is continuous and strictly positive, so $0<Z_{s,\beta}<\infty$ for every
real $\beta$. The reflection-positivity arguments below require
$\beta\ge0$.

A site gauge transformation $g:\Lambda_s\to G$ acts by
\begin{equation}
 U^g(y,\mu)=g(y)U(y,\mu)g(y+e_\mu)^{-1}.
 \label{ch01:gauge-action}
\end{equation}
If $\gamma$ is a path from $y$ to $z$, its holonomy transforms as
$U_\gamma^g=g(y)U_\gamma g(z)^{-1}$. Thus a closed holonomy transforms by
conjugation. Together with left and right invariance of Haar measure, this
proves gauge invariance of the energy and of the law. These are the basic
properties of Wilson's lattice formulation~\cite{Wilson1974}; no gauge
fixing is involved in defining the probability measure.

\begin{definition}[Locality]
\label{ch01:locality}
A function of $U$ is local on $X\subset\Lambda_s$ if it depends only on
links whose two endpoints belong to $X$. It is gauge invariant if its value
is unchanged by \eqref{ch01:gauge-action} for every site gauge transformation.
A finite stencil is a finite subset of $\ZZ^4$, placed in the torus by
translation followed by reduction modulo $s$.
\end{definition}
Locality and gauge invariance are different requirements. A single link
matrix entry is local but usually not gauge invariant. A traced holonomy
around a fixed lattice loop satisfies both.

Gauge transformations also explain why small fields must be defined
through plaquettes rather than individual links. The pure-gauge
configuration $U(y,\mu)=g(y)g(y+e_\mu)^{-1}$ can have links far from the
identity while every plaquette holonomy is the identity. Later,
small-field regions will be specified by bounds on the gauge-invariant
quantity $\chi(U_{y;\mu\nu})$ from \eqref{ch01:chord}.

The connection with classical curvature is most cleanly stated using exact
parallel transport. For a smooth $\gLie$-valued connection $A$, let the
transport on a segment from $x$ to $x+ae_\mu$ solve
\[
 u'(t)=u(t)aA_\mu(x+tae_\mu),\qquad u(0)=1,\qquad 0\le t\le1.
\]
The plaquette product then has the expansion
\begin{equation}
 U_{x;\mu\nu}=1+a^2F_{\mu\nu}(x)+O(a^3),
 \quad
 F_{\mu\nu}=\partial_\mu A_\nu-\partial_\nu A_\mu+[A_\mu,A_\nu].
 \label{ch01:classical-curvature}
\end{equation}
Consequently
$w(U_{x;\mu\nu})=\frac12a^4|F_{\mu\nu}(x)|_{\gLie}^2+O(a^5)$.
This familiar classical expansion motivates the action. It is not a
regularity assertion about typical lattice configurations, nor about an
eventual continuum random connection.

\section{Physical scales and lattice symmetries}
\label{ch01:scales-section}

We keep distances in lattice units distinct from physical distances.
For $y\in\Lambda_s$, let $r_s(y)$ be its unique representative in
$(-s/2,s/2]^4\cap\ZZ^4$. The periodic distance in lattice units is
\[
 \operatorname{dist}_s(y,z)=
       \min_{k\in s\ZZ^4}|\widetilde y-\widetilde z+k|_\infty,
\]
where any lifts $\widetilde y,\widetilde z$ may be used. This distance is
translation invariant and invariant under signed coordinate permutations.
At a centered-cell boundary the chosen representative itself need not
transform as an ordinary vector; periodic statements are unaffected.

Let $H_4$ denote the finite group of signed permutation matrices. An affine
lattice isometry $y\mapsto hy+z$, $h\in H_4$, acts on configurations by
pulling back oriented paths, using inverse links when orientation is
reversed. Plaquette traces and product Haar measure are preserved, so the
Wilson law is invariant under these isometries. Full Euclidean rotation
invariance is not a lattice symmetry. Its recovery in the continuum is a
separate part of the argument.

\begin{definition}[Scale geometry]
\label{ch01:scales}
Choose odd integers $L_m\ge3$ and set
\begin{equation}
 D_0=1,\qquad D_{m+1}=L_mD_m,
 \qquad a_m=D_m^{-1},\qquad \tau_m=\log D_m.
 \label{ch01:scale-sequence}
\end{equation}
Choose strictly increasing physical periods $R_p$, each an integer multiple
of four. At scale $m$ and volume index $p$, the lattice has
$s_{m,p}=R_pD_m$ sites per side and physical spacing $a_m$.
\end{definition}
Its physical side is exactly $a_ms_{m,p}=R_p$. Increasing $m$ means refining
the lattice; integrating out one scale goes from $m+1$ to $m$. The index
$N$ will denote the finest scale in a finite renormalization trajectory.
At cutoff $N$, a lattice site has physical placement
\begin{equation}
 \widehat y=a_N r_{R_pD_N}(y)\in(-R_p/2,R_p/2]^4.
 \label{ch01:placement}
\end{equation}
Further growth and divisibility requirements on these sequences will be
imposed when the estimates that require them have been established.

Odd block ratios give centered blocks without an asymmetric choice of
boundary sites. At a step of ratio $L$, every fine site can be written
uniquely as $Ly+z$, with a coarse site $y$ and
$z\in\{-(L-1)/2,\ldots,(L-1)/2\}^4$. The resulting parent map commutes
with signed coordinate permutations. The corresponding coarse point and
its fine representative have the same physical position.

If an observable uses a fixed stencil $\Sigma\subset\ZZ^4$, its physical
support has diameter at most $a_N\operatorname{diam}(\Sigma)$. Such an
observable becomes local at a point as $N\to\infty$. This is different
from a loop of fixed physical size, whose number of links grows with
$N$. The latter is not included merely by proving estimates for every
fixed stencil.

\section{Reflection positivity}
\label{ch01:reflection}

Reflection positivity is the link between Euclidean correlations and a
positive Hilbert-space inner product. On the lattice it also gives
chessboard estimates, which convert a bound on many reflected copies of a
local observable into a bound on a correlation of separated observables.
We use these classical results as inputs~\cite{OsterwalderSeiler1978,
FrohlichIsraelLiebSimon1978,Seiler1982}.

Here are the support and orientation conventions for the finite-volume
positivity statement. Fix an even side $s=2h$ and one coordinate, temporarily
called $t$. A site reflection acts by $t\mapsto-t$ modulo $2h$; a link
reflection acts by $t\mapsto-1-t$. Their positive sectors are respectively
$1\le t<h$ and $0\le t<h$. In the first case the two fixed time slices
are excluded. In either case observables in the positive algebra depend
only on links whose two endpoints are in the indicated sector.

If the reflected coordinate is $\mu$ and the reflection of sites is
$\vartheta$, the reflected configuration is
\begin{equation}
 (\Theta U)_\nu(y)=
 \begin{cases}
 U_\nu(\vartheta y),&\nu\ne\mu,\\
 U_\mu(\vartheta y-e_\mu)^{-1},&\nu=\mu.
 \end{cases}
 \label{ch01:reflection-links}
\end{equation}
The inverse in the second line is essential. For a complex observable set
$(\theta F)(U)=\overline{F(\Theta U)}$.

\begin{proposition}[Finite-volume reflection positivity]
\label{ch01:lattice-rp}
Let $\beta\ge0$. For either of the two reflection cuts above and every
bounded Borel gauge-invariant observable $F$ in the positive algebra,
\begin{equation}
 \int (\theta F)(U)F(U)\,d\mu_{2h,\beta}(U)\ge0.
 \label{ch01:rp}
\end{equation}
The same conclusion holds after translating the cut or choosing another
coordinate direction.
\end{proposition}
\begin{proof}
This is the finite-torus form of lattice gauge reflection positivity
\cite{OsterwalderSeiler1978,Seiler1982}. Polarization gives a positive
semidefinite pairing between finite linear combinations of positive-side
observables. The support restriction is part of the statement: it is not
an assertion about arbitrary dependence on links crossing the reflection
planes. Positivity of the Boltzmann density, by itself, would not imply
\eqref{ch01:rp}.
\end{proof}

For clarity, one version of the chessboard construction can be described
without observables on seams. Write $s=2q\ell$ and cut each coordinate into
$2q$ intervals of length $\ell$. Let the base cell consist of sites with
all coordinates strictly between $0$ and $\ell$. An observable local on
this cell reads no seam links. Reflecting the cell repeatedly through the
planes at multiples of $\ell$ gives its copies in all
$M=(2q)^4$ cells. Denote the resulting real pullback of a function $F$ to
cell $C$ by $\mathcal R_CF$; the orientation of links is always transformed
as in \eqref{ch01:reflection-links}. For a bounded gauge-invariant $F$ on
the base cell, put
\begin{equation}
 \|F\|_{\mathrm{cb},\mathscr T}
 =\left[\int\prod_{C\in\mathscr T}\mathcal R_C|F|(U)
                 \,d\mu_{s,\beta}(U)\right]^{1/M},
 \label{ch01:chessboard-root}
\end{equation}
where $\mathscr T$ is this specified tiling. The quantity in brackets is
nonnegative. This notation records the tiling as well as the observable;
it is not a volume-independent norm by definition.

The chessboard estimate bounds, for one base-cell observable $F_C$ per cell,
\begin{equation}
 \left|\int\prod_{C\in\mathscr T}\mathcal R_CF_C\,d\mu_{s,\beta}\right|
 \le\prod_{C\in\mathscr T}\|F_C\|_{\mathrm{cb},\mathscr T}.
 \label{ch01:chessboard}
\end{equation}
Apply the abstract chessboard theorem
\cite[Theorem~4.1]{FrohlichIsraelLiebSimon1978} to the nonnegative
observables $|F_C|$, one coordinate at a time. The slots are the $2q$
strips in that coordinate: translation by $\ell$ gives cyclicity, and
Proposition~\ref{ch01:lattice-rp} gives its reflection Cauchy--Schwarz
hypothesis across the strip seams. The involution reflects the local
strip coordinates, including link orientation. Repeating in the other
three coordinates gives the reflected copies in
\eqref{ch01:chessboard-root}. No power-of-two restriction on $2q$ is
needed. Empty cells are filled with the constant one.
More flexible placements used later require explicit control of support,
seams and divisibility. They are not consequences of changing the tiling
symbol in \eqref{ch01:chessboard-root} without checking those conditions.

\section{Correlations and the order of the limits}
\label{ch01:limits}

Recall that $a_N$ is the lattice spacing and $R_p$ the physical side of the
torus, as in Definition~\ref{ch01:scales}. For now let $\mathcal T$ be an
arbitrary set of local-observable labels. A family of insertions assigns to
$v\in\mathcal T$ and $y\in\Lambda_{R_pD_N}$ a continuous gauge-invariant
function $I_{N,p}(v,y;U)$. For each fixed $v$, require locality on a translate
of a finite stencil $\Sigma_v$ independent of $N$ and $p$.
The next chapter constructs such local representatives; their necessary
cutoff-dependent normalization is a later issue.

Choose bare couplings $\beta_{N,p}$ and abbreviate
$\mu_{N,p}=\mu_{R_pD_N,\beta_{N,p}}$, the Wilson probability from
\eqref{ch01:wilson-law}. For a list $v=(v_1,\ldots,v_n)$, define
\begin{align}
 M^{\mathrm{lat}}_{N,p}(v;y)
   &=\int\prod_{i=1}^n I_{N,p}(v_i,y_i;U)\,d\mu_{N,p}(U),\\
 C^{\mathrm{lat}}_{N,p}(v;f)
   &=a_N^{4n}\sum_{y\in\Lambda_{R_pD_N}^n}
                   f(\widehat y_1,\ldots,\widehat y_n)
                   M^{\mathrm{lat}}_{N,p}(v;y).
 \label{ch01:lattice-smear}
\end{align}
The factor $a_N^{4n}$ is the physical volume element for $n$ insertions.
No limit is built into these definitions.

\begin{definition}[Separated test functions]
\label{ch01:separated-test}
Set
\[
 \Conf_n=\{(x_1,\ldots,x_n)\in(\RR^4)^n:x_i\ne x_j\text{ for }i\ne j\}.
\]
A separated test function is an element of $C_c^\infty(\Conf_n)$, extended by zero
to $(\RR^4)^n$. It is admissible at volume $p$ if its support is contained
in $((-R_p/2,R_p/2)^4)^n$.
\end{definition}
Such test functions can depend jointly on all positions; they need not factor into
one-point test functions. Compactness of their support gives a positive distance
from the collision diagonals. We write
$\operatorname{sep}_\infty(x)=\min_{i\ne j}|x_i-x_j|_\infty$, with value
$+\infty$ for $n\le1$.

\begin{definition}[Ordered convergence]
\label{ch01:ordered-limit}
A lattice family converges to kernels $S_n(v;x)$ along an increasing
volume-index map $P$ if, for every $n$, every list $v$ and every separated
test function $f$, the following hold. At each admissible fixed volume $p$ there is
a limit
\begin{equation}
 \lim_{N\to\infty}C^{\mathrm{lat}}_{N,p}(v;f)=C_p^{\mathrm{UV}}(v;f),
 \label{ch01:uv-limit}
\end{equation}
and the retained volumes satisfy
\begin{equation}
 \lim_{k\to\infty} C_{P(k)}^{\mathrm{UV}}(v;f)
       =\int_{\Conf_n} f(x)S_n(v;x)\,dx.
 \label{ch01:volume-limit}
\end{equation}
The test function is eventually admissible along $P$.
\end{definition}
The first limit uses every cutoff $N$ in the chosen spacing sequence. The
second may require a subsequence of volumes. The desired construction has
the stronger quantifier order
\begin{enumerate}[label=\textup{(\roman*)}]
\item choose one lattice family;
\item prescribe any increasing sequence of volume indices;
\item extract one refinement that serves every label list and separated test function.
\end{enumerate}
In symbols, for every strictly increasing map $p_{\rm seq}:\NN\to\NN$
there should be a strictly increasing $q_{\rm seq}$ such that
$P=p_{\rm seq}\circ q_{\rm seq}$ works in
Definition~\ref{ch01:ordered-limit}. The lattice family, including its
normalizations, is chosen before $p_{\rm seq}$. Separate subsequences for
different observables would not suffice.

The bare inverse coupling is to satisfy
\[
 \frac{\beta_{N,p}}{\log(1/a_N)}\longrightarrow\kappa_G
 \quad(N\to\infty),\qquad\text{uniformly in }p,
\]
with $\kappa_G$ as in \eqref{ch01:slope}. This is one part of the target,
not a hypothesis that supplies the required continuum bounds. A fixed
parameter $b>0$ specifies the soft constraint in the block transformation;
it is distinct from $\beta$ and from the coefficient $b_0$ in
\eqref{ch01:slope}.

The eventual limit is a family of correlations of gauge-invariant local
fields. It is not presented as a probability measure on continuum
connections. Its decomposition into phases, reconstruction as relativistic
fields, and stress tensor require additional arguments. In particular,
reflection positivity alone does not establish the full reconstruction
theorem. Likewise, a unique vacuum in each reconstructed phase does not
imply a positive mass gap. Neither a mass gap, uniqueness of the
infinite-volume limit, nor independence of the block parameter $b$ is
included in the conclusions described here.

Finally, \eqref{ch01:volume-limit} concerns separated positions. Extending
correlations as distributions across coincident points is a later
renormalization problem. Such an extension is not automatically a limit
of the unsubtracted lattice correlations at collisions. This distinction
will remain important when products of local fields are considered.

\section{How the construction is organized}
\label{ch01:strategy}

The finite Wilson measure is easy to define; its behavior as the lattice
spacing vanishes is the problem. The number of integration variables grows,
and local observables must be rescaled and mixed before their correlations
can have nonzero finite limits. Weak-coupling perturbation theory predicts
some of these changes, but it does not bound the integral over all gauge
configurations. The construction therefore keeps the finite integral exact
while organizing its estimates by scale.

An exact block transformation integrates out information between two
successive lattice spacings and leaves a density on the coarser links.
The transformed density generally has many interactions, even when the
original action is a sum of plaquette terms. The useful question is whether
these new interactions admit a stable description: a prescribed local
action, a small remainder on configurations with small plaquette curvature,
and strongly suppressed contributions where that condition fails.
Here and throughout the book, small curvature means small plaquette
holonomy; it does not mean that every link is close to the identity in an
arbitrary gauge.

These two parts of configuration space require different estimates.
On the small-curvature part, a stationary background and its Gaussian
fluctuations organize the local integral. Taylor coefficients of low order
modify the prescribed action. The other terms must gain enough powers of
the block ratio or the inverse coupling to remain small after summing over
their positions. On the complement, positive action supplies an exponential
penalty. Connected sets of defects are summed with their full combinatorial
multiplicity. This is the role of the polymer expansion: it controls the
sum once the local integral and energy estimates have been proved.
It does not supply those estimates by itself.

The leading change in the coefficient of curvature squared is logarithmic
in the block ratio. Matching this change determines the constant
$\kappa_G$ in \eqref{ch01:slope}. One must then tune a whole finite-cutoff
trajectory so that its coarse endpoint has a prescribed coupling.
The same choice must work with all the estimates that will be required for
local observables. Choosing a different trajectory each time an observable
is introduced would not construct a single quantum field theory.

Local observables bring an additional finite-dimensional problem.
A monomial with $r$ curvature factors and $k$ derivatives has
\emph{scaling dimension} $2r+k$. Gauge-invariant polynomials in curvature and
its derivatives can mix under blocking. A transfer matrix records this mixing at each fixed scaling
dimension. Bounds for this matrix and for the remainder of each inserted
density are needed in addition to bounds for the density without
insertions. Several insertions also require estimates uniform away from
their collision diagonals. Bounds for lattice differences are essential:
bounded correlations alone would not give smooth continuum kernels.

The free Maxwell theory enters as a local comparison at short distances.
Its role is to distinguish field labels, identify possible equation
relations, and normalize certain symmetry identities. Such a comparison
does not assert that all correlations of the interacting theory are
Gaussian. It must be established for the same finite-cutoff family whose
continuum limit is being taken.

After these estimates are available, two separate arguments remove the
cutoffs. At a fixed physical volume, stability of the renormalization
trajectory gives a Cauchy property as the ultraviolet cutoff increases.
Uniform estimates in the physical volume then permit one common
infinite-volume subsequence, as specified in
Definition~\ref{ch01:ordered-limit}. Neither argument can replace the other.

Finally, lattice symmetries and positivity must be promoted to the
continuum statements needed for reconstruction. Coordinate reflection
positivity gives an initial Hilbert-space structure. Bounds on translation
operators, continuous Euclidean symmetries, and control of the correlations
must be established with compatible normalizations. Decomposition into
phases and relativistic reconstruction follow only after these steps.
The trace identity and the expansion of products at short separation
require further estimates of their own. This order explains why the book
keeps the exact finite model, the local analytic estimates, the choice of a
common family, and the continuum consequences distinct.
}{}
\IfFileExists{chapters/02-observables.tex}{% !TEX root = ../monograph.tex
\chapter{Local gauge-invariant observables}
\label{ch02:observables}

A continuum observable needs more than a lattice expression with the right
classical limit. Integration over fluctuations can mix it with other
observables, including the constant. We therefore separate two tasks. This
chapter constructs finite-dimensional spaces of classical polynomial labels
and local lattice representatives of those labels. Later, the
renormalization group will supply the cutoff-dependent linear maps that
normalize the representatives at a fixed physical scale.

The classical connections used below are test configurations. Their jets
provide a precise way to say which local polynomial a lattice observable
represents. We do not assume that a continuum random gauge field is a smooth
connection, or even that such a random connection has been constructed.
Throughout, $G\subset U(n_G)$ and the metric on its Lie algebra $\gLie$ are
those of Section~\ref{ch01:gauge}.

\section{Radial jets and polynomial labels}
\label{ch02:jets-section}

For a smooth connection near the origin, radial gauge is the condition
$x^\mu A_\mu(x)=0$. In this gauge the connection is recovered from curvature
by the usual radial identity
\begin{equation}
 A_\mu(x)=\int_0^1 t x^\nu F_{\nu\mu}(tx)\,dt.
 \label{ch02:radial-identity}
\end{equation}
Repeated coordinate indices are summed from $1$ to $4$. The identity
explains why a homogeneous potential of polynomial degree $k+1$ carries
the same scaling dimension as a $k$th derivative of curvature, namely
$k+2$. We use this observation to define labels algebraically.

\begin{definition}[Radial jets]
\label{ch02:radial-jets}
A formal radial jet is a sequence $J=(A_k)_{k\ge0}$, where
$A_k=(A_{k,\mu})_{\mu=1}^4$ is a $\gLie$-valued one-form with homogeneous
polynomial coefficients of degree $k+1$, satisfying
\[
 \sum_{\mu=1}^4x^\mu A_{k,\mu}(x)=0.
\]
Write $\mathcal J$ for the space of these sequences and
$\mathcal J_{\le d}$ for the finite-dimensional real space obtained by
retaining only $k+2\le d$. It is the zero space for $d<2$.
\end{definition}
Only finitely many jet coordinates enter any polynomial below. A coordinate
of $A_k$ has classical scaling dimension $k+2$ and polynomial degree one.
The scaling dimension of a monomial is the sum of those of its factors.
For a homogeneous polynomial in jet coordinates, we also call its scaling
dimension its \emph{weight}.
The ordinary polynomial degree counts factors of jet coordinates; to
distinguish it from scaling dimension we call it the \emph{factor count}.
Constant gauge transformations act on every coefficient by $\Ad(G)$.
Orthogonal changes of Euclidean coordinates act by the corresponding
one-form transformation. Both actions preserve scaling dimension and factor count.

\begin{definition}[Polynomial labels]
\label{ch02:labels}
For $d\in\NN$, let
\begin{equation}
 V_d=\CC[\mathcal J]^{\Ad(G)}_{\text{scaling dimension}\le d}
 \label{ch02:label-space}
\end{equation}
be the space of gauge-invariant complex polynomials of scaling dimension
at most $d$. Choose nested real bases $(b_{d,i})$ whose elements have definite scaling dimension
$w_i$ and factor count $f_i$, with the constant polynomial as one basis
vector. Conjugation means conjugating coefficients in these real bases.
Define the factor-count operator by
\begin{equation}
 N_{\mathrm{fac}}b_{d,i}=f_i b_{d,i},\qquad
 t^{N_{\mathrm{fac}}}b_{d,i}=t^{f_i}b_{d,i}\quad(t>0).
 \label{ch02:factor-scaling}
\end{equation}
\end{definition}
The spaces are finite-dimensional: only $A_k$ with $k+2\le d$ can occur,
and every nonconstant coordinate has scaling dimension at least two. The invariant
subspace decomposes according to scaling dimension and factor count, so the stated bases can be chosen once and for all. The
constant has bidegree $(0,0)$; the zero-factor subspace is precisely
$\CC\mathbf1$.

The curvature at the origin is
$F_{\mu\nu}=\partial_\mu A_{0,\nu}-\partial_\nu A_{0,\mu}$: the
commutator term vanishes there in radial gauge. Its coordinates have scaling dimension
two and factor count one. An important invariant polynomial is
\begin{equation}
 c_{F^2}=\frac14\sum_{\mu,\nu}
                 \langle F_{\mu\nu},F_{\mu\nu}\rangle_{\gLie},
 \label{ch02:curvature-square}
\end{equation}
with bidegree $(4,2)$. The inner product is extended complex bilinearly when
used in a complex polynomial, not sesquilinearly.

There are three distinct integers in this notation. A basis element has
scaling dimension $w_i$ and factor count $f_i$; it also belongs to a chosen space
$V_d$ with $w_i\le d$. We shall retain that last choice as part of the
observable label, writing $v=(d,c)$ with $c\in V_d$. The same polynomial
$c$ can occur at different tags $d$. Its eventual normalization depends on
the tag, so the resulting insertions must not be identified without a
comparison theorem. None of these integers is being asserted to be the
exact scaling dimension of the interacting quantum field.

\section{Test holonomies and Taylor normalization}
\label{ch02:taylor-section}

A label is a polynomial of a jet. A lattice observable is a function of
link matrices. To compare them, insert the parallel transports of a
truncated radial connection into the lattice observable.

For $J\in\mathcal J_{\le d}$, write
$A^{\le d}_\mu(x)=\sum_{k+2\le d}A_{k,\mu}(x)$. For a positive link
$(y,\mu)$ of $\ZZ^4$ and a real parameter $h$, define
\begin{equation}
 U(h,J)(y,\mu)=u(1),\qquad
 u'(t)=u(t)h A^{\le d}_\mu(h(y+te_\mu)),\quad u(0)=1.
 \label{ch02:test-holonomy}
\end{equation}
The solution remains in $G$. Standard dependence of solutions of analytic
ordinary differential equations on parameters gives analytic dependence on
the finitely many jet coordinates near zero; see
\cite{Teschl2012}. Reversed links have inverse transport, in the convention
of Section~\ref{ch01:wilson-section}.

Define the jet dilation
\begin{equation}
 \operatorname{sc}_hJ=(h^{k+2}A_k)_k.
 \label{ch02:jet-dilation}
\end{equation}
Homogeneity in \eqref{ch02:test-holonomy} gives
$U(h,J)=U(1,\operatorname{sc}_hJ)$. In particular, a homogeneous label of
scaling dimension $w$ satisfies $c(\operatorname{sc}_hJ)=h^wc(J)$.

\begin{definition}[Taylor map and local section]
\label{ch02:taylor-map}
For a finite-link function $O$ analytic near the identity configuration,
$\operatorname{Tay}_d O$ is the part of scaling dimension at most $d$ in
the Taylor expansion at $J=0$ of $O(U(1,J))$.
A local polynomial section is a complex-linear map
\[
 P_d^{\mathrm{loc}}:V_d\longrightarrow
       \{\text{polynomial functions of finitely many link entries and their
       conjugates}\}
\]
with $\operatorname{Tay}_dP_d^{\mathrm{loc}}[c]=c$.
\end{definition}
The term section simply means a right inverse of the Taylor map. The
polynomial condition is global in the link matrices; it is stronger than
analyticity near the identity. For a gauge-invariant observable the Taylor
polynomial is $\Ad(G)$-invariant, since constant gauge transformations
conjugate the test holonomies.

Equivalently, the section identity requires, for every $J$ and $0\le q\le d$,
\begin{equation}
 \left.\frac{d^q}{dh^q}P_d^{\mathrm{loc}}[c](U(h,J))\right|_{h=0}
 =\left.\frac{d^q}{dh^q}c(\operatorname{sc}_hJ)\right|_{h=0}.
 \label{ch02:taylor-identity}
\end{equation}
It fixes every homogeneous component of scaling dimension at most $d$. It also does not prescribe the value of the lattice observable
on a typical rough configuration.

\section{Finite-stencil representatives}
\label{ch02:sections}

The construction uses lattice curvature transported to a common point.
Extend the metric \eqref{ch01:metric} to the ambient real matrix space as
$\operatorname{Re}\tr(X^*Y)/n_G$, and let $\Pi_{\gLie}$ be orthogonal
projection onto $\gLie$. For a plaquette based at $z$, set
\[
 F_\square(U)=\Pi_{\gLie}\frac{U_\square-U_\square^{-1}}2.
\]
If $\gamma$ is a lattice path from the origin to $z$, its transported
curvature is
\begin{equation}
 \sigma_{\gamma,\square}(U)=\Ad(U_\gamma)F_\square(U).
 \label{ch02:transported-curvature}
\end{equation}
Under a site gauge transformation it is conjugated by the gauge element
at the origin. It is a real polynomial in the entries and conjugate
entries of the links it reads, because inverse unitary matrices are
conjugate transposes.

\begin{theorem}[Local polynomial representatives]
\label{ch02:local-sections}
For each $d$ there is a finite centered stencil $\Sigma_d\subset\ZZ^4$,
invariant under the signed permutation group $H_4$, and a local polynomial
section $P_d^{\mathrm{loc}}$ with the following properties.
Its values are gauge invariant and local on $\Sigma_d$. It preserves the
unit and complex conjugation, and satisfies
\begin{equation}
 P_d^{\mathrm{loc}}[c](h^*U)=P_d^{\mathrm{loc}}[h\cdot c](U),
 \qquad h\in H_4.
 \label{ch02:section-covariance}
\end{equation}
Here $h^*$ pulls back lattice links and
$(h\cdot c)(J)=c(h^{-1}\cdot J)$ is the induced action on labels.
Moreover, each basis column $P_d^{\mathrm{loc}}[b_{d,i}]$ is a polynomial
in finitely many transported curvatures, and every monomial in that
polynomial has at least $f_i$ factors.
\end{theorem}
The last assertion is what makes the section useful for estimates. An
arbitrary right inverse of the Taylor map need not preserve this lower
bound on factor count.

\begin{proof}
We give the finite-dimensional construction, invoking the usual inverse
function theorem rather than reproving it.

First consider the maps $J\mapsto\sigma_{\gamma,\square}(U(1,J))$.
Their linearizations at the zero jet separate $\mathcal J_{\le d}$.
Indeed, differentiating transport at the identity gives the line integral
of the connection along the link. The derivative of a plaquette curvature
is therefore its circulation. The derivative of the conjugating transport
in \eqref{ch02:transported-curvature} contributes nothing, since the
curvature itself vanishes at the identity.

To verify separation, suppose all these linearized plaquette measurements
vanish on $A=\sum_{k=0}^{d-2}A_k$. For each coordinate plane the circulation
around its unit square based at $z$ is the polynomial
\[
 N_{\mu\nu}(z)=\int_0^1\!\int_0^1
 (\partial_\mu A_\nu-\partial_\nu A_\mu)
                  (z+te_\mu+ue_\nu)\,dt\,du.
\]
It vanishes at every integer $z$, and hence is the zero polynomial. If $k$
is the largest degree index with $A_k\ne0$, its homogeneous part of degree
$k$ is $\partial_\mu A_{k,\nu}-\partial_\nu A_{k,\mu}$. Thus $dA_k=0$.
The radial identity and Euler's identity give
\[
 \sum_\mu x^\mu
  (\partial_\mu A_{k,\nu}-\partial_\nu A_{k,\mu})=(k+2)A_{k,\nu}.
\]
It follows that $A_k=0$, a contradiction. Descending through the degree
indices proves separation. Since $\mathcal J_{\le d}$ is finite-dimensional,
a finite collection of measurements already has injective joint derivative.

Fix such a collection and write
\[
 \mathscr S(J)=\bigl(\sigma_1(U(1,J)),\ldots,
                          \sigma_M(U(1,J))\bigr)\in\gLie^M,
 \qquad D=D\mathscr S(0).
\]
The map is analytic, $\mathscr S(0)=0$, and $D$ is injective. Both the jet
space and $\gLie^M$ have invariant real inner products. Hence the inverse
of $D$ on its range, composed with orthogonal projection onto that range,
is an equivariant linear map $C$ with $CD=1$. The inverse function theorem
applied to $C\mathscr S$ gives an analytic local inverse $g$; this is the
standard finite-dimensional theorem, for example
\cite[Theorem~9.24]{Rudin1976} together with its analytic form.
Set $f=gC$. Then $f(0)=0$ and $f\mathscr S(J)=J$ near zero.
Equivariance follows from uniqueness of the inverse as a germ.

Let $T_df$ be the ordinary Taylor polynomial of $f$ through degree $d$.
It has no constant term and is equivariant. For each real basis polynomial
$b_{d,i}$ define $Q_i(z)=b_{d,i}(T_df(z))$. Then $Q_i$ is invariant, is real
on real arguments, and has no monomial of degree below $f_i$. The unit
column is one. Consequently
\[
 \Psi_c(U)=\sum_i c_iQ_i(\sigma_1(U),\ldots,\sigma_M(U))
 \quad\text{for }c=\sum_i c_i b_{d,i}
\]
is a gauge-invariant polynomial cylinder function. Since
$T_df(\mathscr S(J))=J+O(|J|^{d+1})$, its Taylor map is $c$.
Each jet coordinate has scaling dimension at least two, so every monomial
of polynomial degree greater than $d$ also has scaling dimension greater
than $d$. No truncation of $Q_i$ after substitution is
needed.

Finally average over the finite lattice symmetry group:
\begin{equation}
 P_d^{\mathrm{loc}}[c](U)
   =\frac1{|H_4|}\sum_{h\in H_4}\Psi_{h^{-1}\cdot c}(h^*U).
 \label{ch02:section-average}
\end{equation}
Uniqueness of the parallel-transport equation gives
$h^*U(h_0,J)=U(h_0,h^{-1}\cdot J)$ for a scale parameter $h_0$.
Thus every summand still has Taylor map $c$. Reindexing the finite sum
proves \eqref{ch02:section-covariance}. Reality and the unit are preserved.
A signed image of a transported curvature is another such curvature,
possibly with a minus sign from orientation reversal. The group action
preserves factor count, so the lower bound on monomial degree survives
averaging. Enclosing all paths and plaquettes in a centered box produces
the finite stencil $\Sigma_d$.
For $d<2$ there are only constants and the construction reduces to the
constant section.
\end{proof}

All choices in this theorem are made on the covering lattice $\ZZ^4$,
before selecting a torus period. For a site $y\in\Lambda_s$, write
$P_d^{\mathrm{loc}}[c](y;U)$ for evaluation after translating the stencil
to $y$ and reducing its links modulo $s$. The construction remains
well-defined if different covering links become the same torus link.
Gauge invariance, covariance and the polynomial identities survive that
identification.

The abstract jet space admits an $O(4)$ action, but the covariance in
\eqref{ch02:section-covariance} concerns only $H_4$. Arbitrary rotations do
not act on the lattice links. This distinction prevents the classical
label symmetry from being mistaken for an already proved symmetry of the
continuum correlations.

\section{Small-field bounds for the representatives}
\label{ch02:bounds}

Recall that $V_d$ is the polynomial label space of scaling dimension at most $d$ from Definition~\ref{ch02:labels}; factor count counts curvature-jet factors instead of their scaling dimensions.

We now use the factor-count property of
Theorem~\ref{ch02:local-sections}. Recall from \eqref{ch01:chord} that
$\chi(u)=\|u-1\|_{\mathrm F}/\sqrt{n_G}$ measures the size of a plaquette.
On a fixed local stencil, a bound $\chi(U_\square)\le t$ gives a factor
$t$ for each transported curvature in \eqref{ch02:transported-curvature}.
The same factor is retained under suitably scaled derivatives.

For a collection $A=(A_e)_{e\in E_s}$ of Lie algebra elements, define the
left derivative
\[
 D_A^L F(U)=\left.\frac{d}{d\varepsilon}
       F\bigl((\exp(\varepsilon A_e)U_e)_{e\in E_s}\bigr)
          \right|_{\varepsilon=0}.
\]
An ordered list of these derivatives is used; no commutation of Lie
derivatives is assumed. Fix a nonnegative spatial exponent $p_0$ and a
site $y$. A permitted direction at threshold $t$ satisfies
\begin{equation}
 |A_e|_{\gLie}\le t\bigl(1+\operatorname{dist}_s(e^-,y)\bigr)^{p_0},
 \label{ch02:direction}
\end{equation}
where $e^-$ is the initial site of the positive link. This is a local
definition for the estimate; the later polymer norms will specify which
orders and spatial exponents they retain.

\begin{proposition}[Uniform section bounds]
\label{ch02:section-bounds}
For every $d,q,p_0$ there are finite constants $C_{d,q,p_0}$ and $C_d$,
independent of $s$, $y$ and $0<t\le1$, with the following properties.
If every plaquette used by the translated section has distance from the identity at most $t$,
then for every list of permitted directions of length $k\le q$,
\begin{equation}
 \left|D_{A_1}^L\cdots D_{A_k}^L
       P_d^{\mathrm{loc}}[b_{d,i}](y;U)\right|
       \le C_{d,q,p_0}t^{f_i}.
 \label{ch02:section-small-bound}
\end{equation}
For arbitrary configurations,
\begin{equation}
 \sup_{i,y,U}|P_d^{\mathrm{loc}}[b_{d,i}](y;U)|\le C_d.
 \label{ch02:section-global-bound}
\end{equation}
The same constants apply when the stencil wraps around the torus.
\end{proposition}
\begin{proof}
Projection is contractive and $U^{-1}=U^*$, so
\[
 |F_\square(U)|_{\gLie}
 \le\frac{\|U_\square-U_\square^*\|_{\mathrm F}}{2\sqrt{n_G}}
 \le\chi(U_\square).
\]
Transport by a unitary matrix preserves the metric. An undifferentiated
curvature measurement is therefore bounded by $t$ on the stated domain.

On the fixed stencil the right side of \eqref{ch02:direction} is bounded
by a constant times $t$, uniformly in the period: periodic distance is
no larger than distance to a chosen covering-space representative.
Differentiating a fixed path product $k$ times gives a finite sum of
products with exactly $k$ direction matrices and otherwise unitary
factors. Its derivatives are bounded by a constant times $t^k$.
The same conclusion holds for each transported curvature. For a positive
order $k$, this is in particular bounded by a constant times $t$, since
$t\le1$.

Every monomial representing the column $b_{d,i}$ has at least $f_i$
curvature factors. Apply the ordered product rule: each factor, whether
or not differentiated, contributes a factor bounded by a constant times
$t$. Summing the finitely many resulting terms proves
\eqref{ch02:section-small-bound}. Matrix entries of a unitary matrix have
modulus at most one, so the polynomial cylinder formula gives
\eqref{ch02:section-global-bound}.
If covering links coincide on the torus, one first differentiates the
covering expression and assigns the same direction to identified links.
The number of terms and their bounds do not increase with the period.
\end{proof}

The scaling operator $t^{N_{\mathrm{fac}}}$ from
\eqref{ch02:factor-scaling} is tailored to this estimate. By linearity,
permitted derivatives of $P_d^{\mathrm{loc}}[t^{-N_{\mathrm{fac}}}c]$
are bounded by $C_{d,q,p_0}\sum_i|c_i|$ on the small-field domain. The
powers $t^{f_i}$ cancel column by column. Outside that domain the global
bound instead gives $C_d\sum_i|c_i|t^{-f_i}$. The missing small-field gain
must later be supplied by the probability cost of a large-field region;
it cannot be obtained by extrapolating \eqref{ch02:section-small-bound}.

\section{Local representatives of the equations of motion}
\label{ch02:equation-section}

The ordinary section $P_d^{\rm loc}$ reproduces Taylor coefficients,
but need not reproduce equations of motion. We construct a second
section for this purpose. Everything in this section concerns the
finite-volume Wilson measure; no renormalization estimate is used.

A radial polynomial potential $A_j$, homogeneous of degree $j+1$,
is a \emph{Maxwell solution potential} if
\[
 \sum_\mu\partial_\mu(dA_j)_{\mu\nu}=0\quad\hbox{for every }\nu.
\]
Its curvature is automatically closed. Let $\mathcal I_d\subset V_d$
be the polynomials vanishing on all such solution jets through degree
$d$. This is the degree-$d$ part of the equation ideal; multiplication
is understood only when the resulting scaling dimension remains at most $d$.
The condition preserves both scaling dimension and factor count.
It is the kernel of restriction to solution jets, rather than a
statement about the interacting measure.

\begin{theorem}[An equation-adapted local section]
\label{ch02:equation-adapted}
There are a real, symmetry-equivariant linear section
$\mathcal V_d$ from $V_d$ to local continuous gauge-invariant functions of links, a real equivariant automorphism
$R_d$ of $V_d$, and finite constants $r_d,M_d$, with the following
properties:
\begin{enumerate}
\item $\mathcal V_d1=1$, $R_d1=1$, and
$\operatorname{Tay}_d\mathcal V_d p=p$.
The section has a fixed finite stencil, independent of the torus.
\item If $p$ has factor count $s$ and scaling dimension $w$,
then $R_dp-p$ is a sum of terms whose pair
$(\text{factor count},\text{scaling dimension})$ is strictly larger than $(s,w)$
in lexicographic order. Every term of $\mathcal V_dp$ is a product
of at least $s$ transported curvature coordinates.
\item For every $p\in\mathcal I_d$, every torus of side at least $M_d$,
every anchor $o$, every real $\beta\ne0$, and every bounded measurable
$S$ depending only on links with both endpoints farther than $r_d$
from $o$,
\begin{equation}
 \int \mathcal V_d(R_dp)(o,U)S(U)\,d\mu_\beta(U)=0.
 \label{ch02:microscopic-null}
\end{equation}
No gauge invariance or regularity of $S$ is assumed.
\end{enumerate}
The small-field derivative bounds for a fixed section proved above
also apply to $\mathcal V_d$, with its own finite constants.
\end{theorem}

The regrading $R_d$ is essential: the exact interacting relation is
$\mathcal V_d(R_d\mathcal I_d)$, not an assertion that
$P_d^{\rm loc}(\mathcal I_d)$ has zero Wilson correlations.
The proof below supplies the finite-stencil construction, whereas the
compact-group integration by parts underlying it is the usual lattice
equation-of-motion identity for Wilson's measure
\cite{Wilson1974,OsterwalderSeiler1978}.

\subsection{Curvature tests avoiding a prescribed link}

For $u\in\gLie$, pairing the transported curvature
\eqref{ch02:transported-curvature} with $u$ gives the real-valued function
$U\mapsto\langle u,\sigma_{\gamma,\square}(U)\rangle_{\gLie}$.
We call any such function a curvature coordinate.
Call it admissible for an oriented link $e$ if neither its plaquette
nor its transporting path uses the underlying unoriented link $e$.
Such a coordinate is independent of $U_e$.

\begin{lemma}\label{ch02:avoiding-tests}
For each $e$, the linearizations at the zero jet of admissible
curvature coordinates span the dual of the finite radial jet space
$\mathcal J_{\le d}$.
\end{lemma}
\begin{proof}
It suffices to show that their common annihilator is zero.
Removing one link does not disconnect $\ZZ^4$, so every plaquette
not containing $e$ admits a transporting path avoiding $e$.
Let $A$ be the polynomial one-form representing a jet in the annihilator.
At the zero connection the derivative of transport around an oriented
unit plaquette $\square$ in direction $A$ is
$\oint_{\partial\square}A=\int_\square dA$ by Stokes' theorem.
Transport to the anchor does not change this first variation.
Consequently $\int_\square dA=0$ for every plaquette that does not
contain $e$.

For each orientation $(\mu,\nu)$, this integral is a polynomial in
the base point $z$ of the plaquette. It vanishes at all but finitely
many integer points and hence identically, by the one-variable polynomial
identity applied successively in each coordinate. Let $A_j$ be the
highest nonzero homogeneous layer of $A$. Averaging a polynomial over
a translated unit square preserves its highest homogeneous part as a
polynomial in $z$. Thus the degree-$j$ part of the integral is
$(dA_j)_{\mu\nu}(z)$. Since the integral vanishes for every orientation,
$dA_j=0$. Radial gauge and Euler's identity give
\[
 0=\partial_i(x^\mu A_{j,\mu})
   =A_{j,i}+x^\mu\partial_i A_{j,\mu}
   =(j+2)A_{j,i},
\]
where closedness was used in the last equality.
This contradicts the choice of $j$. Descending through the layers
proves the assertion. Finite-dimensional duality now selects a
finite spanning collection of admissible coordinates.
\end{proof}

\subsection{Construction of the equation-adapted section}

\begin{proof}[Proof of Theorem~\ref{ch02:equation-adapted}]
We first represent the generators of the equation ideal by Wilson forces.
We then integrate by parts in a single link and assemble the local section.

\smallskip\noindent\emph{Step 1. Representation by Wilson forces.}

Write $S_W$ for the Wilson action without its factor $\beta$.
For a positive link $e$ and a direction $X\in\gLie$,
the force $F_e\in\gLie$ is characterized by
\[
 \partial_{e,X}S_W=\langle X,F_e\rangle,\qquad
 \partial_{e,X}f(U)=
 \left.\frac{d}{dt}f(U_e\leftarrow e^{tX}U_e)\right|_{t=0}.
\]
Differentiating the plaquettes incident on $e$ expresses $F_e$
as a signed sum of plaquette-curvature coordinates based at the
initial vertex of $e$. The precise signs follow the chosen
plaquette orientations. Transport the force to the anchor by a
path avoiding $e$, and pair it there with an arbitrary element of $\gLie$.
At the zero connection the force is the signed sum, to first order in
$A$, of the integrals $\int_\square dA$ over plaquettes incident on $e$.

Fix a layer $j$. These linear force functionals, restricted to jets
through layer $j$, span every functional which acts only on $A_j$
and annihilates the Maxwell solution potentials in that layer.
Indeed, if all force functionals vanish, their values at integer
base points form identically zero polynomials. Their top-degree
terms are
\[
 -\sum_\mu\partial_\mu(dA_j)_{\mu\nu}.
\]
For a link in direction $\nu$, the two incident $(\mu,\nu)$ plaquettes
enter with opposite signs. Their integrals therefore give a first
difference in the $\mu$ direction. Thus the top layer is a Maxwell
solution potential. Every functional of the indicated type
annihilates the common force kernel and is therefore a finite
linear combination of these forces. Notice that this argument
identifies the top layer only; it does not identify the entire
kernel of the discrete force system with the continuum solution
jet space.

Here is the corresponding description of the polynomial ideal.
Choose a linear projection of each jet layer onto its solution
subspace and let $\mathcal E$ be their direct sum. In the full,
not necessarily invariant, polynomial algebra the kernel of
$p\mapsto p(\mathcal EJ)$ is generated by linear functionals
vanishing on those solution subspaces. To check this without
an ideal-generation theorem, use
\[
 \prod_{\nu=1}^s u_\nu-\prod_{\nu=1}^s(\mathcal E^*u_\nu)
 =\sum_{i=1}^s
 (u_i-\mathcal E^*u_i)
 \prod_{\nu<i}\mathcal E^*u_\nu
 \prod_{\nu>i}u_\nu .
\]
Choose homogeneous linear coordinates, so both gradings are
preserved. Every homogeneous element of $\mathcal I_d$ is
therefore a sum of products $\lambda z_1\cdots z_{s-1}$,
where $\lambda$ is a one-layer equation functional.
Individual summands need not be gauge invariant.

Represent $\lambda$ by the finite force combination just constructed
on its truncated jet space. For each force summand involving $e$,
represent each $z_i$ by a linear combination of curvature coordinates
avoiding $e$, using Lemma~\ref{ch02:avoiding-tests} on the \emph{full}
jet space through degree $d$. Multiply these representatives by
that force summand. The resulting function $f$ has at least $s$
curvature factors. Its Taylor polynomial is
\begin{equation}
 \operatorname{Tay}_d f
   =\lambda z_1\cdots z_{s-1}
       +\{\hbox{terms strictly higher in }(s,w)\}.
 \label{ch02:equation-triangle}
\end{equation}
There are two distinct errors. The nonlinear part of a curvature
or force coordinate increases factor count. The linear force
representation was fixed only through the layer of $\lambda$;
its remaining linear terms increase the scaling dimension while preserving the factor
count. This is why the order in
\eqref{ch02:equation-triangle} puts factor count first.

\smallskip\noindent\emph{Step 2. Integration by parts and construction of the section.}

Every force term above has the form
$H\langle X,F_e\rangle$, where $X$ is the fixed direction at the
anchor transported back to the link, and $H$ is the product of the
remaining curvature coordinates and the exterior spectator.
Freeze the other links. Both $X$ and $H$ are then constant in $U_e$. Haar invariance
on the compact group gives
\[
 0=\int_G\partial_{e,X}\{e^{-\beta S_W}\}\,dU_e
   =-\beta\int_G\langle X,F_e\rangle
                          e^{-\beta S_W}\,dU_e.
\]
Multiplication by $H$ is legitimate because it is constant in
$U_e$. Fubini then proves the desired zero expectation, even for a
merely bounded measurable exterior spectator. This proof never
differentiates that spectator.

All lists used in constructing a basis are finite. Enlarge one
radius to contain their plaquettes, paths, and differentiated links.
Enlarge one period threshold so that distinct relevant cover links
cannot become identified with a differentiated link after
periodization. These two choices work for every anchor by translation.
They are independent of $\beta$ and of the exterior spectator.

Average the constructed functions over the gauge element at the
anchor. Every transported curvature transforms there by the adjoint
action, so this is a finite-dimensional average of tensors under the adjoint action.
It preserves the number of curvature factors and
\eqref{ch02:equation-triangle}. It also preserves the exact zero
identity: gauge invariance of the Wilson law transfers the gauge
transformation to the bounded exterior spectator without changing
its support. The result is fully gauge invariant.

Choose a basis of each homogeneous part of $\mathcal I_d$, perform
this construction there, and extend to a basis of $V_d$ by using
$P_d^{\rm loc}$ on a complement. The resulting linear map
$\mathcal F_d$ preserves the unit. Its Taylor map is
\[
 R_d=\operatorname{Tay}_d\mathcal F_d=I+T_d,
\]
where $T_d$ strictly increases the finite lexicographic grading.
Thus $T_d$ is nilpotent and
$R_d^{-1}=I-T_d+T_d^2-\cdots$ is a finite sum.
Average $\mathcal F_d$ simultaneously on domain and range over
signed coordinate permutations and over conjugation. These
operations preserve the ideal, both gradings, and the unit;
hence the same triangular conclusion remains valid.
Finally put
\[
 \mathcal V_d=\mathcal F_dR_d^{-1}.
\]
This gives the Taylor identity, all symmetries, and
$\mathcal V_dR_d p=\mathcal F_dp$ for $p\in\mathcal I_d$.
Both $R_d$ and its inverse preserve the lower bound on factor
count. The fixed finite products therefore obey exactly the
small-field derivative argument used for $P_d^{\rm loc}$,
with new constants for the new stencil.
\end{proof}

The identity \eqref{ch02:microscopic-null} is exact at each cutoff.
It gives no bound on a renormalized coefficient change.
To use it for continuum null fields one must compare this section
with the section used along the actual renormalization history,
with one invertible comparison map for all labels and all spectators.


\section{Examples and the remaining normalization problem}
\label{ch02:normalization}

The unit illustrates the simplest section: $P_d^{\mathrm{loc}}[\mathbf1]=1$
for every $d$. Its eventual physical normalization still depends on the
normalization index. There will be a reference scale for each $d$, and a choice of
coefficient making the renormalized insertion exactly one. This does not
follow by putting a universal factor $a_N^{-d}$ in front of the raw
constant section.

For the curvature-square polynomial $c_{F^2}$ in
\eqref{ch02:curvature-square}, a useful classical representative is the
clover average
\begin{equation}
 P^{\mathrm{cl}}(y;U)=\frac14\sum_{\square\ni y}w(U_\square),
 \label{ch02:clover}
\end{equation}
where the sum has four corner plaquettes in each of the six coordinate
planes, hence $24$ terms. Orientation does not matter because $w(u^{-1})=w(u)$.
For a smooth test connection the expansion
\eqref{ch01:classical-curvature} yields
\[
 P^{\mathrm{cl}}(y;U(h,J))
     =h^4 c_{F^2}(J)+O(h^5)
\]
at the origin, and thus $\operatorname{Tay}_4P^{\mathrm{cl}}=c_{F^2}$.
The factors are consistent: each plaquette contributes
$\frac12h^4|F_{\mu\nu}|_{\gLie}^2$, four copies occur per plane, and the
prefactor is $1/4$. This example fixes the normalization, but the general
construction is free to use a different section satisfying the same
Taylor identity.

The clover formula does not solve the quantum normalization problem.
Under blocking, the curvature square can mix with the constant. A general
renormalized version therefore involves both a multiplicative coefficient
and an additive constant, with coefficients depending on the cutoff and
volume. More generally, a label in $V_d$ can mix with other labels in the
same finite-dimensional space. The correct normalization is consequently
an invertible linear map on $V_d$, whose existence and estimates must be
proved along the chosen renormalization trajectory.

There is also a distinction between polynomial relations and vanishing
fields. Some polynomial combinations represent free Maxwell equations of
motion. The interacting lattice equations arise instead from Haar
integration by parts, and include nonlinear corrections. To identify the
null space of continuum observables requires both those identities and a
rank estimate. Neither the section theorem nor its classical Taylor
normalization identifies the interacting null space with the unmodified
free equation ideal.

The output of this chapter is therefore concrete but limited: a common
finite-stencil representative for every polynomial label, with symmetry
and threshold estimates uniform in the torus period. These are exactly
the local data needed to begin the renormalization of observables. The
uniform control of their correlations and their continuum normalization
remain the task of the subsequent construction.
}{}
\setcounter{chapter}{2}
\IfFileExists{chapters/03-maxwell.tex}{\chapter{The free Maxwell comparison theory}\label{ch03:maxwell}

The reference theory for short distances consists of independent Maxwell
fields, one for each direction in the Lie algebra. Its role is more precise
than a formal approximation to the action. It supplies finite-dimensional
spaces of local fields, correlation functionals that distinguish those
fields, and a normalization for the stress tensor. The interacting
construction must eventually reproduce these data in an appropriate
short-distance comparison. No interacting limit is assumed in this chapter.

We fix the positive inner product on the real Lie algebra $\gLie$ and write
$d_{\gLie}=\dim_{\RR}\gLie$. Use the orthonormal basis of Section~\ref{ch01:gauge},
$(e_\alpha)_{\alpha=1}^{d_{\gLie}}$. Greek letters $\mu,\nu,\rho,\sigma$
from the middle of the alphabet denote coordinate indices in
$\{1,2,3,4\}$; $\alpha,\beta$ denote indices in this basis. Time is the fourth
coordinate. The superscript $\mathrm{Max}$ always denotes the free
reference theory. We first define its correlations at distinct points;
this requires no multiplication of distributions at the same point.

\section{Curvature covariance}\label{ch03:covariance}

Let
\begin{equation}\label{ch03:green}
 \Gamma(x)=\frac{1}{4\pi^2|x|^2},\qquad x\ne0.
\end{equation}
Its locally integrable extension is the fundamental solution of
$-\Delta$ in four dimensions. A convenient way to specify the curvature
covariance is to start from independent Gaussian vector-potential
components with covariance
$\delta_{\alpha\beta}\delta_{\mu\nu}\Gamma(x-y)$, and take a curl in
each argument. The vector potential in this description is a device for
computing the covariance; it is not one of the gauge-invariant local
observables used in the construction.

Write $F^\alpha_{\mu\nu}=-F^\alpha_{\nu\mu}$ for the curvature
components in this basis. The covariance is the product of the inner
product on $\gLie$ and a spacetime kernel $C$: pairing the two fields
with $u,v\in\gLie$ gives $\langle u,v\rangle_{\gLie}C_{\mu\nu,\rho\sigma}(x-y)$.
In coordinates,
\begin{equation}\label{ch03:curvature}
 S^{\mathrm{Max}}_2(F^\alpha_{\mu\nu},F^\beta_{\rho\sigma};x,y)
 =\delta_{\alpha\beta}C_{\mu\nu,\rho\sigma}(x-y),
\end{equation}
where, with $h_{\mu\nu}=\partial_\mu\partial_\nu\Gamma$,
\begin{align}
 C_{\mu\nu,\rho\sigma}
 &=-\delta_{\nu\sigma}h_{\mu\rho}
   +\delta_{\nu\rho}h_{\mu\sigma}
   +\delta_{\mu\sigma}h_{\nu\rho}
   -\delta_{\mu\rho}h_{\nu\sigma},\label{ch03:kernel}\\
 h_{\mu\nu}(z)
 &=\frac{8z_\mu z_\nu-2\delta_{\mu\nu}|z|^2}
          {4\pi^2|z|^6}.\label{ch03:hessian}
\end{align}
The minus signs in the first and last terms arise because a derivative
in $y$ is minus a derivative in $z=x-y$. This convention matters both
for reflection positivity and for the stress normalization.

The covariance is translation invariant and transforms as a two-form in
each pair of indices under $O(4)$. At nonzero separation it satisfies
\begin{equation}\label{ch03:maxwell-equations}
 \partial_\lambda C_{\mu\nu,\rho\sigma}
 +\partial_\mu C_{\nu\lambda,\rho\sigma}
 +\partial_\nu C_{\lambda\mu,\rho\sigma}=0,
 \qquad
 \sum_\mu\partial_\mu C_{\mu\nu,\rho\sigma}=0.
\end{equation}
The first identity is the Bianchi identity; the second is the source-free
Maxwell equation. They hold off the origin. A distributional derivative
can have a contact term at the origin, so these identities must not be
used to infer a choice of contact extension.

For a multi-index $a=(a_1,\ldots,a_4)$, put
$|a|=a_1+\cdots+a_4$ and
$\partial^a=\partial_1^{a_1}\cdots\partial_4^{a_4}$.
Derivatives of the fields have covariance
\begin{equation}\label{ch03:jet-covariance}
 S^{\mathrm{Max}}_2(\partial^aF^\alpha_{\mu\nu},
                   \partial^cF^\beta_{\rho\sigma};x,y)
 =\delta_{\alpha\beta}(-1)^{|c|}
       \partial^{a+c}C_{\mu\nu,\rho\sigma}(x-y).
\end{equation}
In particular, for each fixed derivative order, this kernel is smooth
away from zero and homogeneous of degree $-4-|a|-|c|$. Thus its size on
$|x-y|\ge r$, $0<r\le1$, has a polynomial bound in $r^{-1}$.
This elementary estimate is the free counterpart of the much harder
uniform bounds required from the lattice construction.

\section{Local polynomials and Wick correlations}\label{ch03:wick}

A local polynomial is a polynomial in finitely many curvature derivatives
at one point. Its scaling dimension is computed by assigning degree $2$
to $F$ and degree $1$ to each derivative. For example,
$F^\alpha_{\mu\nu}F^\beta_{\rho\sigma}$ has degree $4$, whereas
$\partial_\lambda F^\alpha_{\mu\nu}$ has degree $3$.
Scaling dimension is an algebraic grading of the reference model; it is
not an assertion about anomalous dimensions of the interacting fields.

The polynomial attached to one insertion is Wick ordered. We use the
usual Gaussian pairing rule, whose operator formulation is due to
Wick~\cite{Wick1950}. Here it can be taken as the definition of the
separated correlation functions. Expand each polynomial into monomials,
regard each factor as a distinguishable occurrence, and sum over pairings
of all occurrences subject to the condition that paired occurrences
belong to different insertions. For each pair use
\eqref{ch03:jet-covariance}. The empty pairing has value one. The
resulting finite sum is denoted by
\[
 S_n^{\mathrm{Max}}(P_1,\ldots,P_n;x_1,\ldots,x_n),
 \qquad x_i\ne x_j\quad(i\ne j).
\]
There is no covariance at zero separation in this definition.

For instance, if $X_1,X_2,Y_1,Y_2$ are linear curvature derivatives,
\begin{align*}
 S_2^{\mathrm{Max}}({:}X_1X_2{:},{:}Y_1Y_2{:};x,y)
 &=C(X_1,Y_1;x-y)C(X_2,Y_2;x-y)\\
 &\quad+C(X_1,Y_2;x-y)C(X_2,Y_1;x-y).
\end{align*}
Here $C(X,Y;z)$ denotes the bilinear covariance obtained from
\eqref{ch03:jet-covariance}. By contrast, in a word containing several
insertions in the same time half-space, contractions between those
\emph{different} insertions are included. Wick ordering removes
self-contractions at each insertion, not all contractions inside a word.

The pairing rule gives complex multilinearity, symmetry under simultaneous
permutation of labels and positions, reality, and the unit identity
\[
 S_{n+1}^{\mathrm{Max}}(1,P_1,\ldots,P_n;x_0,x_1,\ldots,x_n)
 =S_n^{\mathrm{Max}}(P_1,\ldots,P_n;x_1,\ldots,x_n).
\]
Differentiating an insertion differentiates its incident covariances.
The product rule therefore agrees with ordinary differentiation of the
correlation function in that insertion position.

If the labels are homogeneous of degrees $d_1,\ldots,d_n$, then
\begin{equation}\label{ch03:homogeneity}
 S_n^{\mathrm{Max}}(P_1,\ldots,P_n;\lambda x_1,\ldots,\lambda x_n)
 =\lambda^{-\sum_i d_i}
 S_n^{\mathrm{Max}}(P_1,\ldots,P_n;x_1,\ldots,x_n)
 \quad(\lambda>0).
\end{equation}
For general labels of bounded degree, decompose into finitely many
homogeneous pieces. Equations \eqref{ch03:jet-covariance} and
\eqref{ch03:homogeneity} give bounds on all derivatives on separated
configuration spaces. In particular, smearing against a smooth compactly
supported function whose support avoids collisions is well defined.
This statement is weaker than specifying the correlations on test functions whose supports meet collision diagonals.

\section{The on-shell polynomial space}\label{ch03:shell}

The Maxwell identities impose linear relations on the jet variables.
They should be imposed before discussing strict positivity or the
number of independent fields. At each integer $j\ge0$, let
$\mathcal S_j$ be the real vector space of homogeneous degree-$j$
polynomial two-forms on $\RR^4$, with values in $\gLie$, satisfying
\[
 dF=0,\qquad \delta F=0.
\]
Here $\delta F=0$ means
$\sum_\mu\partial_\mu F_{\mu\nu}=0$ for every $\nu$; the harmless
choice of sign for the codifferential does not affect its kernel.
These are finite-dimensional spaces. They encode the possible homogeneous
Taylor coefficients of a source-free Maxwell curvature. A Taylor
coefficient of order $j$ is normalized as $D^jF(0)/j!$.

The phrase \emph{on shell} will mean restriction to these solution jets.
It does not refer to restricting a point of Euclidean spacetime.
A linear coordinate on $\mathcal S_j$ has scaling dimension $j+2$. As in Definition~\ref{ch02:labels}, we also call
`weight' the scaling dimension of a homogeneous jet polynomial.
Consequently polynomials of weight at most $d$ use only the spaces
$\mathcal S_j$ with $j\le d-2$. Constants are retained at every degree.
The action on $\gLie$ is the adjoint action of $G$; invariant polynomials
are the gauge-invariant local observables of the free reference theory.

To compare this definition with the radial-potential labels $V_d$,
associate to a closed homogeneous curvature $F^{(j)}$ the radial
potential
\begin{equation}\label{ch03:radial}
 A^{(j+1)}_\mu(x)=\int_0^1 t x^\nu F^{(j)}_{\nu\mu}(tx)\,dt
 =\frac{1}{j+2}\sum_\nu x^\nu F^{(j)}_{\nu\mu}(x).
\end{equation}
There is a sum over $\nu$ in the integral as well. This is the usual
radial homotopy formula: $x^\mu A_\mu=0$ and $dA=F$ when $dF=0$.
The connection here is the \emph{linear} Maxwell potential; no
nonabelian commutator is inserted into this formula. Restricting a label
in $V_d$ to these radial potentials gives a map
\[
 \omega_d:V_d\longrightarrow
 \{
 \text{complex polynomials on Maxwell solution jets of weight at most }d
 \}.
\]
It preserves the unit, conjugation and the geometric actions. Define
\begin{equation}\label{ch03:free-space}
 \mathcal M_d^{\mathrm{Max}}=\operatorname{im}\omega_d.
\end{equation}
Equivalently, this is $V_d/\ker\omega_d$ as a vector space. The image
formulation is useful: it does not assume a separate surjectivity theorem
onto a larger invariant-polynomial space.

The covariance identities \eqref{ch03:maxwell-equations} make the Wick
correlations well defined on these on-shell polynomial images. Products
of a Maxwell equation with another jet polynomial vanish in every
separated Wick correlation. Conversely, proving that \emph{all}
correlation-null polynomials vanish on shell requires a nondegeneracy
argument; it does not follow from writing down the equations.

The free spaces form a filtration. The inclusion
$\mathcal M_d^{\mathrm{Max}}\to\mathcal M_e^{\mathrm{Max}}$ for
$d\le e$ retains the same polynomial, while differentiation raises
weight by one. For a simple nonabelian Lie algebra there is no nonzero
invariant linear form on $\gLie$. Therefore every gauge-invariant field
of degree less than four in this polynomial image is a constant.
This low-degree statement follows from the weights: a nonconstant
monomial below degree four would have just one curvature factor or one
of its first derivatives.

\section{Reflection and the meaning of nondegeneracy}\label{ch03:reflection}

Let $\vartheta(\boldsymbol x,t)=(\boldsymbol x,-t)$.
Reflection acts on a two-form by pullback, so a component with one time
index changes sign and a purely spatial component does not. On a local
polynomial, $\theta$ combines this action with complex conjugation.
On a word it also reverses the order and reflects every insertion point.
This is the antilinear reflection used in the Osterwalder--Schrader
form, not merely conjugation of scalar coefficients.

For words $w_i$ supported at distinct points in the positive half-space,
the free field has reflection positivity:
\begin{equation}\label{ch03:rp}
 \sum_{i,j}\overline{c_i}c_j
 S^{\mathrm{Max}}(\theta w_i\,w_j)\ge0.
\end{equation}
The same inequality holds for smooth joint test functions supported on separated
positive-time configurations. This is the standard Gaussian
reflection-positive construction underlying
Osterwalder--Schrader reconstruction~\cite{OsterwalderSchrader1973,OsterwalderSchrader1975},
applied to the curvature covariance rather than to an unconstrained
vector-potential observable. Its hypotheses can be checked at the
one-particle level: spatial Fourier transformation of
\eqref{ch03:green} gives
\[
 \frac{e^{-|\boldsymbol p|(t+s)}}{2|\boldsymbol p|}
 \qquad(t,s>0),
\]
and taking the two curls with the reflected two-form indices gives the
positive transverse-polarization Gram form. Derivative jets act on the
same factors. The Gaussian Fock construction then gives
\eqref{ch03:rp}; contractions within each positive-time word are included
in its Fock vector. On compact separated supports the kernels are smooth,
so the corresponding form on joint test functions follows by integration. We use this
standard argument by citation rather than repeating the Fock-space
construction.

A simple sign check is informative. At separation $z=re_4$, the six
unordered curvature components have diagonal covariance. Three purely
spatial components have value $1/(\pi^2r^4)$; the three components
$F_{i4}$ have value $-1/(\pi^2r^4)$. Reflection changes the sign of
$F_{i4}$ in the first slot. The reflected one-jet form is therefore
positive in all six directions. Without the tensorial reflection, the
same displayed covariance would give the wrong positivity statement.

Positive semidefiniteness of \eqref{ch03:rp} alone does not identify its
null space. For later rank arguments one needs the stronger statement
that a nonzero on-shell polynomial has a strictly positive reflected
two-point pairing at two suitable points. Derivatives of different
orders may mix in that pairing. Checking one homogeneous order at a
time does not establish the statement for their sum.

\section{Strict positivity of finite jet families}\label{ch03:strict}

We now identify the null relations in the free polynomial theory.
Recall from Section~\ref{ch03:shell} that $\mathcal S_j$ is the space
of homogeneous degree-$j$ Maxwell curvature solutions. Write a member as
\[
 P_{ab}(v)=\sum_{i_1,\ldots,i_j}
 P_{ab;i_1\cdots i_j}v_{i_1}\cdots v_{i_j}.
\]
The derivative-index coefficients are symmetric. On this coefficient
space use the Frobenius inner product, summing both orders of the
antisymmetric field indices, all ordered derivative indices and components in the basis of Section~\ref{ch01:gauge}.
In particular, no factor $1/2$ is inserted for the field indices.
The linear jet probe represented by $P$ is
\begin{equation}\label{ch03:probe}
 \ell_P(F)=\frac1{j!}\sum_{a,b,i_1,\ldots,i_j,\alpha}
 P^\alpha_{ab;i_1\cdots i_j}
 \partial_{i_1}\cdots\partial_{i_j}F^\alpha_{ab}.
\end{equation}
The factor $1/j!$ matches our Taylor convention. Every linear functional
on the order-$j$ solution jets has this form.

Put $I(z)=z/|z|^2$ and
$R_z=1-2zz^{\mathsf T}/|z|^2$. For a two-form, $R_z^{(2)}$ means
applying $R_z$ to both form indices. The full reflection on a
polynomial solution is
$(R_z\cdot P)(v)=R_z^{(2)}P(R_zv)$: the derivative indices are
reflected as well. For real probes define
\[
 B_z(P,Q)=S_2^{\mathrm{Max}}(\ell_P,\ell_{R_z\cdot Q};z,0).
\]
Extend this bilinearly across different derivative orders. Covariance and
exchange of the two insertions show that $B_z$ is symmetric.

\begin{proposition}[Finite-jet nondegeneracy]\label{ch03:jet-positive}
For every $z\ne0$ and integer $m\ge1$, the form $B_z$ is positive
definite on $\bigoplus_{j<m}\mathcal S_j$.
\end{proposition}

The scalar differentiation identity needed in the proof is the harmonic
case of Hobson's differentiation formula~\cite{Hobson1892}. In our
normalization, if $P\in\mathcal S_j$, it says, componentwise,
\begin{equation}\label{ch03:hobson}
 \frac1{j!}P(\partial)\Gamma(y)
   =\frac{(-2)^j}{4\pi^2}|y|^{-2j-2}P(y),\qquad y\ne0.
\end{equation}
Here $P(\partial)$ means the ordered coefficient contraction without
the factor $1/j!$. Maxwell solutions are harmonic, as follows by
applying a divergence to the Bianchi identity. Thus the harmonic
hypothesis of the scalar formula holds in each component. We use that
classical differentiation result; the following finite-dimensional
argument explains its application to the mixed jet orders needed here.

\begin{proof}
Let $\mathcal R(P)(y)$ be the curvature response obtained by contracting
$\partial_y^j C(y)/j!$ with the coefficients of $P$. Taking the two
curls in \eqref{ch03:hobson} gives
\begin{equation}\label{ch03:kelvin}
 \mathcal R(P)(y)=c_j|y|^{-4}R_y^{(2)}P(I(y)),
 \qquad c_j=\frac{(j+1)(j+2)(-2)^j}{\pi^2}.
\end{equation}
For clarity, the identities determining the coefficient are as follows.
Set $V_i(y)=\sum_a y_aP_{ai}(y)$ and
$b_j=(-2)^j/(4\pi^2)$. The divergence of the differentiated Green
potential is
$D_i=-2(j+1)b_j|y|^{-2j-4}V_i$.
The Maxwell equations and homogeneity give
$\partial_iV_m-\partial_mV_i=(j+2)P_{im}$.
The response is $2(\partial_mD_i-\partial_iD_m)$, so its coefficient
is $4(j+1)(j+2)b_j=c_j$ and its two-form factor is
\[
 P_{im}(y)-2|y|^{-2}(y_iV_m(y)-y_mV_i(y))
   =(R_y^{(2)}P(y))_{im}.
\]
This verifies \eqref{ch03:kelvin}, including its sign, off the origin.

Fix the separation $z$. For $P\in\mathcal S_j$ define a finite tower
$\Sigma_zP$ by
\begin{equation}\label{ch03:triangular}
 (\Sigma_zP)_k
   =\frac{c_j}{c_k}\operatorname{Taylor}_k[P(\,\cdot-I(z))],
 \qquad 0\le k\le j.
\end{equation}
The operator $\operatorname{Taylor}_k$ extracts the homogeneous term of
degree $k$ at zero. Translation preserves the Maxwell equations, and
each homogeneous term separately satisfies them. Thus all components in
\eqref{ch03:triangular} lie in their declared solution spaces.
The degree-$j$ component is $P$ itself.

By the finite Taylor formula, the response of this tower is
\[
 \mathcal R(\Sigma_zP)(y)
   =c_j|y|^{-4}R_y^{(2)}P(I(y)-I(z)).
\]
Since $DI(z)=|z|^{-2}R_z$, its Taylor coefficients at $y=z$ vanish
through order $j-1$, while its order-$j$ coefficient is
$c_j|z|^{-2j-4}(R_z\cdot P)$.
For $Q\in\mathcal S_k$, differentiation in the second insertion gives
\[
 B_z(\Sigma_zP,Q)
 =(-1)^k\big\langle R_z\cdot Q,
       \operatorname{Taylor}_k[\mathcal R(\Sigma_zP)](z)
                 \big\rangle_F.
\]
It follows that this expression vanishes if $k<j$, and at $k=j$ equals
$(-1)^jc_j|z|^{-2j-4}\langle P,Q\rangle_F$.
Transforming the second argument introduces only lower orders in addition
to its unchanged top order. Symmetry therefore gives, on the full tower,
\begin{equation}\label{ch03:diagonal-gram}
 B_z(\Sigma_zP,\Sigma_zQ)
 =\sum_{j<m}\frac{(j+1)(j+2)2^j}{\pi^2|z|^{2j+4}}
                  \langle P_j,Q_j\rangle_F.
\end{equation}
Every coefficient is positive. Finally $\Sigma_z$ is invertible:
it is the identity plus a strictly degree-lowering map, hence its inverse
is a finite geometric sum. This proves positivity on the whole tower,
including mixtures of derivative orders.
\end{proof}

\begin{corollary}[Polynomial positivity and finite separating families]
\label{ch03:detectors}
Every nonzero element of $\mathcal M_d^{\mathrm{Max}}$ has strictly
positive reflected two-point pairing at a fixed separated pair.
For any basis $(P_1,\ldots,P_s)$ of this space and fixed distinct points
$x,y$, the linear map
\[
 P\longmapsto
 \big(S_2^{\mathrm{Max}}(P,P_j;x,y)\big)_{j=1}^s
\]
is injective. Consequently a label in $V_d$ vanishes in every separated
correlation if and only if its on-shell polynomial $\omega_d$ vanishes.
\end{corollary}

\begin{proof}
This is the standard symmetric-Fock-space consequence of the one-particle
Gram form and Wick's rule~\cite{Wick1950}. Explicitly, choose coordinates
$X_1,\ldots,X_M$ diagonalizing the positive form on a tower containing
all variables of the polynomial, with positive squared lengths $b_i$.
The reflected pairing of Wick monomials is
\[
 \langle X^a,X^c\rangle_{\rm ref}
   =\mathbf1_{a=c}\,a!\prod_i b_i^{a_i}.
\]
Thus the squared norm of a polynomial is the sum of the squared absolute
values of its coefficients with these positive weights; the constant
coefficient has weight one. Restricting to the gauge-invariant image
preserves positivity. If all basis-column pairings vanish, the pairing
with the reflected conjugate of the polynomial vanishes too, forcing
that polynomial to be zero. This also identifies the correlation-null
labels. On any fixed finite-dimensional degree space, the map of correlation pairings
has a positive minimum norm on the unit sphere. Small perturbations of
its matrix remain injective; applying this fact to interacting
correlations still requires a quantitative comparison estimate.
\end{proof}


\section{Quadratic fields and normalization checks}\label{ch03:quadratic}

We return to the free curvature covariance
\eqref{ch03:curvature}; in the basis of Section~\ref{ch01:gauge}, each component has the same
covariance and orthogonal components have zero cross-covariance. Define
\begin{align}
 Q_{\mu\nu}
 &=\sum_{\rho=1}^4
       {:}\langle F_{\mu\rho},F_{\nu\rho}\rangle_{\gLie}{:},
 & Q_{\mathrm{tot}}&=\sum_\mu Q_{\mu\mu},\label{ch03:quadratics}\\
 \mathcal E^{\mathrm{Max}}&=\tfrac14 Q_{\mathrm{tot}},
 &T^{\mathrm{Max}}_{\mu\nu}
   &=-Q_{\mu\nu}+\tfrac14\delta_{\mu\nu}Q_{\mathrm{tot}}.
 \label{ch03:stress}
\end{align}
The sign of $T^{\mathrm{Max}}$ is chosen to match the shell-flux
convention in the supplied construction. A complete stress-normalization
argument must check that convention against the translation Ward identity;
conservation alone does not determine the sign.

The tensor in \eqref{ch03:stress} is symmetric and traceless. Its
divergence vanishes as an on-shell Wick polynomial. These are the
ordinary algebraic Maxwell stress identities, obtained from the Bianchi
and source-free equations; the same identities hold in separated
correlations. The scalar $\mathcal E^{\mathrm{Max}}$ is nevertheless
nonzero. It is the curvature-square operator, not the trace of the free
stress tensor.

With the present normalization, Wick's rule gives
\begin{equation}\label{ch03:scalar-pair}
 S_2^{\mathrm{Max}}(\mathcal E^{\mathrm{Max}},
                    \mathcal E^{\mathrm{Max}};x,y)
 =\frac{3d_{\gLie}}{\pi^4|x-y|^8},\qquad x\ne y.
\end{equation}
For checking constants, at separation $re_4$ one has
$\mathcal E^{\mathrm{Max}}=\tfrac12\sum_{\mu<\nu}{:}\|F_{\mu\nu}\|_{\gLie}^2{:}$. The two cross-pairings of each square,
the six unordered components, and the prefactor $1/4$ give
$\tfrac14\cdot2\cdot6d_{\gLie}=3d_{\gLie}$.
Orthogonal covariance makes the answer independent of the direction of
$x-y$. The calculation uses only separated Wick contractions.

This normalization will be needed when comparing a renormalized scalar
field with the free curvature square. It does not itself identify the
interacting trace, its coefficient, or its logarithmic asymptotic.
Those conclusions require the interacting stress Ward identity and
quantitative estimates on moving comparison fields.

\section{The degree-four conserved tensor}\label{ch03:classification}

The finite-dimensional classification needed for the stress tensor is
small enough to state explicitly. In this section \emph{centered} means
that the constant Wick coefficient is zero. By the Wick prescription it
is also the one-point expectation. We consider real symmetric tensor
fields $J_{ab}$ of degree at most four, invariant under the adjoint action
and transforming as a rank-two tensor under the signed coordinate
permutations $H_4$. Conservation is an identity in the on-shell polynomial
space, or equivalently in every separated free correlation by
Corollary~\ref{ch03:detectors}.

\begin{proposition}[Conserved tensor line]\label{ch03:stress-line}
The space of such centered tensors has dimension three. Its conserved
subspace is the line $\RR T^{\mathrm{Max}}$, with
$T^{\mathrm{Max}}$ as in \eqref{ch03:stress}.
\end{proposition}

\begin{proof}
Simplicity of $\gLie$ excludes invariant linear forms on $\gLie$ and makes
an invariant symmetric bilinear form on $\gLie$ a multiple of the fixed
metric. Through degree four a centered invariant polynomial therefore
contains two undifferentiated curvature factors. Indeed, a derivative on
one of two factors would raise the degree to at least five, while a
single differentiated factor still has no invariant linear functional on $\gLie$.

Recall the quadratic fields $Q_{ab}$ and $Q_{\mathrm{tot}}$ from
\eqref{ch03:quadratics}. A basis of the tensor candidates is
\begin{align*}
 Z^{(1)}_{ab}&=\delta_{ab}Q_{aa},\\
 Z^{(2)}_{ab}&=\delta_{ab}(Q_{\mathrm{tot}}-2Q_{aa}),\\
 Z^{(3)}_{ab}&=(1-\delta_{ab})Q_{ab}.
\end{align*}
There is no sum over $a$ on the right sides of the first two formulas.
To see that these span, consider sign reversal of each coordinate.
A diagonal component can contain only squares of curvature components.
Permutations fixing its distinguished axis divide them into two classes:
those containing that axis, and those containing neither instance of it.
An off-diagonal component with indices $a,b$ must be odd in axes $a,b$
and even in the others. Its possible contractions are
${:}\langle F_{ac},F_{bc}\rangle_{\gLie}{:}$ with $c\notin\{a,b\}$;
permutations identify their coefficients. This gives exactly the three
candidates displayed above.

For an explicit rank check, let
$Q_{ij,ij}={:}\|F_{ij}\|_{\gLie}^2{:}$.
Evaluate the first component of the divergence of each $Z^{(l)}$ at
$e_1$, against $Q_{12,12}(0)$ and $Q_{23,23}(0)$. The two rows of
these evaluations, after division by $d_{\gLie}/(4\pi^2)^2$, are
\begin{equation}\label{ch03:divergence-matrix}
 \begin{pmatrix}-256&0&128\\0&-512&-128\end{pmatrix}.
\end{equation}
These numbers come directly from the curvature covariance
\eqref{ch03:kernel} and the two cross-pairings in a quadratic Wick
correlation. For example,
$S_2^{\mathrm{Max}}(Q_{11},Q_{12,12};te_1,0)
=2d_{\gLie}/(\pi^4t^8)$; its derivative at $t=1$ gives the first
entry after the stated division. The first two columns are independent.
Also $Z^{(1)}_{12}=Z^{(2)}_{12}=0$, whereas the correlation of
$Z^{(3)}_{12}(e_1)$ with
${:}\langle F_{13},F_{23}\rangle_{\gLie}{:}(0)$ equals
$-16d_{\gLie}/(4\pi^2)^2$. Thus the candidate tensors are themselves
independent, and the kernel of divergence has dimension at most one.

The Maxwell stress is conserved and
\[
 -T^{\mathrm{Max}}
 =\tfrac12 Z^{(1)}-\tfrac14 Z^{(2)}+Z^{(3)}.
\]
It is nonzero by the preceding off-diagonal correlation. Hence it spans that
kernel.
\end{proof}

This result fixes only a line of conserved tensors. Multiplying the
stress by a scalar preserves all the properties used in the proposition.
We next fix its normalization by a local Ward flux and then classify
linear responses on fields of arbitrary bounded degree.


\section{Local stress fluxes}\label{ch03:ward}

The sign in \eqref{ch03:stress} is fixed by a local identity, not by
conservation. Let $x_0,\ldots,x_n$ be distinct points and choose a real
$\chi\in C_c^\infty(\RR^4)$ equal to one near $x_0$, with support
avoiding $x_1,\ldots,x_n$. For a vector field $v$ define the flux
\begin{equation}\label{ch03:flux}
 \mathcal F_{\chi,v}(P_0;P_1,\ldots,P_n)
 =-\sum_{a,b}\int \partial_a\chi(y)v_b(y)
 S_{n+2}^{\mathrm{Max}}(T^{\mathrm{Max}}_{ab},P_0,\ldots,P_n;
                       y,x_0,\ldots,x_n)\,dy.
\end{equation}
The integrand is supported on $\operatorname{supp}\nabla\chi$, a compact set separated from all the insertions.
It is therefore an ordinary integral of a smooth function; no extension
of a product through $x_0$ is involved.

For a constant translation $v=e_b$, let $\delta_v=\partial_b$ act on
local polynomial labels. For $v(y)=K(y-x_0)$, $K^{\mathsf T}=-K$,
let $\delta_v=\mathcal L_K$ be the infinitesimal action induced on jets
by the pullback of two-forms under $y\mapsto x_0+e^{tK}(y-x_0)$.
Equivalently, on a smooth closed two-form $f$ its variation is
\begin{equation}\label{ch03:lie}
 (\mathcal L_v f)_{ij}
 =v_b\partial_bf_{ij}+ (\partial_i v_b)f_{bj}
                            +(\partial_jv_b)f_{ib}
 =(d\iota_v f)_{ij},
\end{equation}
with sums over $b$. On polynomial labels the action is the induced
Leibniz derivation, including the derivative indices of each jet.
For the dilation $v(y)=y-x_0$, this action at the center is the weight
operator
\begin{equation}\label{ch03:weight}
 \mathscr E=\sum_{j,\lambda}(j+2)X_{j,\lambda}
                                  \frac{\partial}{\partial X_{j,\lambda}},
\end{equation}
where $X_{j,\lambda}$ are any linear coordinates on order-$j$ solution
jets, including the $\gLie$ components. The sum is finite on each polynomial.

\begin{proposition}[Normalized free Ward identities]\label{ch03:free-ward}
For translations, centered rotations and centered dilation,
\begin{equation}\label{ch03:ward-identity}
 \mathcal F_{\chi,v}(P_0;P_1,\ldots,P_n)
 =S_{n+1}^{\mathrm{Max}}(\delta_vP_0,P_1,\ldots,P_n;x_0,\ldots,x_n).
\end{equation}
This holds for arbitrary finite Wick polynomials and arbitrary finite
spectator lists. The common coefficient is $+1$ for the stress
\eqref{ch03:stress} and the outward-flux convention \eqref{ch03:flux}.
\end{proposition}

The Ward mechanism is the standard free-field action of the stress
current. We include the normalization and the treatment of two
contractions at the enclosed insertion, since neither is supplied by a
formal integration by parts through a singularity. Wick's theorem is
used throughout in its usual form~\cite{Wick1950}.

\begin{proof}
Polarize our stress between closed two-forms $f,g$:
\[
 t_{ab}(f,g)=-\sum_c\bigl(\langle f_{ac},g_{bc}\rangle_{\gLie}+\langle g_{ac},f_{bc}\rangle_{\gLie}\bigr)
             +\tfrac12\delta_{ab}\sum_{c,d}\langle f_{cd},g_{cd}\rangle_{\gLie}.
\]
The Bianchi identity gives
\[
 \sum_a\partial_at_{ab}
 =-\sum_c\big(\langle (\operatorname{div}f)_c,g_{bc}\rangle_{\gLie}
                         +\langle f_{bc},(\operatorname{div}g)_c\rangle_{\gLie}\big).
\]
The tensor is symmetric and traceless. Thus contracting with a constant,
skew-linear or scalar-linear $v$ produces no additional divergence term
from $Dv$.

Suppose $f$ is a smooth Maxwell wave near $\operatorname{supp}\chi$,
and $g$ is the covariance pulse from $F_{ij}(x_0)$. Regularize its Green
function by convolution with a smooth compactly supported unit-mass
approximate identity $\rho_\epsilon$, so
$-\Delta\Gamma_\epsilon=\rho_\epsilon$.
The resulting closed pulse has divergence
\[
 (\operatorname{div}g_\epsilon)_c
   =\delta_{cj}\partial_i\rho_\epsilon(\,\cdot-x_0)
       -\delta_{ci}\partial_j\rho_\epsilon(\,\cdot-x_0).
\]
Integrate the displayed divergence identity against $\chi v$.
An integration by parts against $\rho_\epsilon$ shows that the flux
converges to
$\partial_i(v_bf_{bj})(x_0)-\partial_j(v_bf_{bi})(x_0)
=(\mathcal L_vf)_{ij}(x_0)$.
Here $\chi=1$ near the center. On $\operatorname{supp}\nabla\chi$,
the regularized covariance $g_\epsilon$ and its derivatives of each fixed
order converge uniformly to $g$ and the corresponding derivatives as
$\epsilon\to0$. Differentiating in the pulse center, while
holding $v,\chi$ fixed, gives the same identity for every jet probe.
For dilation, \eqref{ch03:lie} becomes
$(y-x_0)\cdot\nabla f+2f$; its order-$j$ jet at $x_0$ is
$(j+2)$ times the original jet. This gives \eqref{ch03:weight}.

In a Wick correlation, extract the two factors at the stress insertion.
The remaining sum has three types. If both partners are outside $x_0$,
the two waves are smooth throughout the cutoff support and conservation
gives zero flux. If exactly one partner is at $x_0$, the preceding
calculation differentiates that occurrence. Summing its possible
partners reconstructs Wick's rule, and summing occurrences gives the
Leibniz derivation $\delta_v$.

If both partners are at $x_0$, their jet orders $j,l$ give pulse
homogeneities $-4-j$ and $-4-l$. Conservation away from $x_0$ makes
the flux independent of the choice of $\chi$. Rescaling a cutoff by $t$
about $x_0$ multiplies this same flux by $t^{e-5-j-l}$, where
$e=0$ for translation and $e=1$ for rotation or dilation.
Taking $t=2$ forces the flux to vanish. The flux integral involves the two
pulses only away from their common singularity at $x_0$.
The polarization counts each ordered pair twice, canceling its factor
$1/2$. Combining the three types proves \eqref{ch03:ward-identity}.
\end{proof}

\section{Separating fields by insertions of bounded scaling dimension}\label{ch03:bounded-detector}

The polynomial labels and their scaling dimensions are those of Section~\ref{ch02:jets-section}; throughout this section they are evaluated in the free curvature field defined in Section~\ref{ch03:covariance}.

Corollary~\ref{ch03:detectors} detects a field of degree $d$ by another
field of degree $d$. A different statement is needed when a charge can
raise degree: its response must be tested by spectators whose degree
remains fixed. Their \emph{number} may be arbitrary.

We first make explicit the polynomial space used here. The free image
$\mathcal M_d^{\mathrm{Max}}$ in \eqref{ch03:free-space} is the full
space of invariant polynomials of weight at most $d$ on the solution
jets. Indeed, the radial map identifies each order of solution jets with
a linear subspace of the corresponding radial-potential space.
Choose an invariant inner product and its orthogonal projection onto
this subspace, separately in each weight. Composing a polynomial with
these projections extends it to the radial jets without changing its
weight or invariance. This is the usual compact-group averaging
argument; no equality with a larger space of unrestricted formal jets
is being assumed.

Let $\mathfrak P$ be the polynomial algebra on all solution jets and
$\mathfrak P^G$ its invariant subalgebra. For $D\ge0$, define
\begin{equation}\label{ch03:ND}
 \mathcal N_D=\{P\in\mathfrak P:
 S^{\mathrm{Max}}(P(x_0),P_1(x_1),\ldots,P_n(x_n))=0\},
\end{equation}
where the equality is required for every separated word with
$P_i\in\mathcal M_D^{\mathrm{Max}}$. There is no restriction on
the degree of $P$ or the length $n$.

\begin{proposition}[A separating family of insertions of bounded scaling dimension]
\label{ch03:uniform-detector}
There is a finite $D_*$ such that
$\mathcal N_{D_*}\cap\mathfrak P^G=\{0\}$.
\end{proposition}

\begin{proof}
First restrict all jet orders to at most $j$. Hilbert's finite-generation
argument~\cite{Hilbert1890}, using compact-group averaging, gives
finitely many invariant generators homogeneous in ordinary polynomial
degree. Let $b_j$ bound their weights. To apply that standard theorem,
the jet space is finite-dimensional, the polynomial ring is Noetherian,
and averaging preserves its invariant ideal and homogeneous pieces.
Thus an invariant polynomial of ordinary degree $k$ is a sum of products
of generators of total ordinary degree $k$, each of weight at most $b_j$.

If a polynomial in $\mathcal N_{b_j}$ is nonzero, take its lowest
nonzero ordinary-degree part $P_k$. Express its reflected conjugate as
such a sum of products, and put the factors at distinct spectator points
approaching a common point different from $x_0$. There are precisely $k$
spectator legs. Higher-degree parts of $P$ cannot pair with them, and
for $P_k$ every spectator leg pairs with a leg at $x_0$. No singular
spectator--spectator contraction survives. The limit is the strictly
positive two-point pairing of $P_k$ with its reflected conjugate,
contradicting \eqref{ch03:ND}. This proves detection at bounded jet order.

To obtain a uniform bound, we use a collision that lowers the highest
jet order. For a linear functional $p$ on scalar order-$j$ jets, let
$X_j(p)\in\gLie$ be its evaluation on the spacetime indices of the jet.
Write $X_{j,\alpha}(p)=\langle e_\alpha,X_j(p)\rangle_{\gLie}$ in
the basis of Section~\ref{ch01:gauge}. Fix a polynomial
$P$ with highest jet order $j$ and a quadratic spectator
${:}\langle X_k(p),X_l(q)\rangle_{\gLie}{:}$, where $l<k$.
Contracting just its order-$k$ leg with the top-order variables of $P$
leaves the polynomial
\begin{equation}\label{ch03:replacement}
 \mathcal R_{p,q,j}(z)P
 =\sum_{\alpha,\lambda}X_{l,\alpha}(q)
 C_{k,j}(p,e_\lambda;z)
           \frac{\partial P}{\partial X_{j,\alpha}(e_\lambda)}.
\end{equation}
The covariance $C_{k,j}$ is in a single unit direction in $\gLie$. We claim that
\begin{equation}\label{ch03:replacement-null}
 P\in\mathcal N_D,\quad k+l+4\le D
 \quad\Longrightarrow\quad
 \mathcal R_{p,q,j}(z)P\in\mathcal N_D\quad(z\ne0).
\end{equation}

Here is the necessary separation of singularities. A differentiated
curvature covariance of orders $k,j$ is harmonic, homogeneous of degree
$-r$, $r=4+k+j$, and becomes smooth after multiplication by
$|z|^{2(r-1)}$. A nonzero such harmonic function cannot become smooth
after multiplication by $|z|^{2(r-2)}$. Indeed, for a harmonic
homogeneous $H$ of degree $-r$,
\[
 \Delta(|z|^{2s}H)=4s(s+1-r)|z|^{2s-2}H.
\]
Starting with a smooth extension at $s=r-2$ and iterating would extend
$H$ smoothly to zero, contradicting its negative homogeneity unless it
vanishes. A product of two harmonic poles of orders $r_1,r_2$,
on the other hand, is cleared by $|z|^{2(r_1+r_2-2)}$.

Insert the quadratic spectator at $x_0+tz$ and keep all other spectators
fixed. Its two contractions with $P$ give finitely many homogeneous
product poles. Subtract them before multiplying by $t^r$.
The limit of the remaining expression is the correlation with
\eqref{ch03:replacement}: the lower single-contraction orders vanish,
and contractions with remote spectators are smooth. The original full
correlation is zero by \eqref{ch03:ND}. Successively removing the
highest powers of $t^{-1}$ shows that the surviving harmonic pole is
minus a product-pole coefficient of the same degree, or is zero.
The extra radial cancellation just proved forces it to be zero in
either case. This proves \eqref{ch03:replacement-null}; in particular
we have not assumed existence of the unmodified rescaled collision limit.

Put $g=d_{\gLie}$ and choose
$D_*\ge\max(b_g,2g+5)$. Use $k=g+1$ and $l=0,\ldots,g-1$.
For each $l$ choose a nonzero probe $q_l$ and a solution jet on which it
has value one. Such jets exist in every order: for example, the curl of
$e_3\operatorname{Re}(x_1+ix_2)^{l+1}$ is a nonzero homogeneous
Maxwell solution of degree $l$. If the highest order $j$ of $P$ exceeds
$g$, each polynomial \eqref{ch03:replacement} has smaller weight than
$P$. Strong induction on weight and \eqref{ch03:replacement-null}
therefore make it zero.

The matrix of polynomial entries $(X_{l,\alpha}(q_l))_{l,\alpha}$
has nonzero determinant: choose the jets so that $X_l(q_l)=e_{l+1}$
for each $l$. At these jets the matrix is the identity. Multiplication
by its adjugate, in the polynomial integral domain, shows that every
componentwise top-order contraction derivative in
\eqref{ch03:replacement} vanishes. As $p,z$ vary, the functionals
$C_{k,j}(p,\cdot;z)$ span the dual of the order-$j$ solution space.
To verify this last assertion, a vector annihilated by all of them has
a response wave whose order-$k$ jets vanish everywhere off zero.
It is then a polynomial of degree below $k$ there; its negative
homogeneity makes it zero. Formula~\eqref{ch03:kelvin} makes the
original vector zero. Hence all top-order partial derivatives of $P$
vanish. Descending removes every jet order above $g$, and the fixed-order
argument with $b_g$ concludes the induction.
\end{proof}

\section{Finite correlation criteria for infinitesimal symmetries}\label{ch03:finite-ward}

We now give the two finite-dimensional consequences used later to
normalize translations and remove an internal rotation defect.
Throughout this section the Lie algebra is simple and nonabelian.
Recall that $\mathcal M_d^{\mathrm{Max}}$ consists of invariant
polynomials of weight at most $d$ on the solution jets $\mathcal S_j$;
its elements are local Wick fields, not distributions defined on collision
sets. For maps $A$ on this space, a \emph{Ward sum} means
\[
 W_A(P_1,\ldots,P_n;x)=\sum_{i=1}^n
 S_n^{\mathrm{Max}}(P_1,\ldots,AP_i,\ldots,P_n;x).
\]
All points in this definition are distinct. The word length is arbitrary.

\begin{theorem}[Characterization of infinitesimal translations]\label{ch03:vector-tests}
Let $d\ge\max(4,D_*)$. Consider real-preserving complex-linear maps
$A_b:\mathcal M_d^{\mathrm{Max}}\to\mathcal M_{d+1}^{\mathrm{Max}}$
which kill the unit, send $\mathcal M_e^{\mathrm{Max}}$ into
$\mathcal M_{e+1}^{\mathrm{Max}}$ for $e\le d$, and transform as a vector
under signed coordinate permutations. Their Ward sums vanish identically
if and only if $A_b=c\partial_b$ for a common real scalar $c$.
There are finitely many evaluations of configuration derivatives of these
Ward sums, at collinear configurations with line axis different from
$b$, whose common kernel is precisely this translation line.
\end{theorem}

\begin{proof}
Translation invariance proves one direction. We prove the converse,
starting at weight four. Let $\mathcal S_j^{\mathrm{sc}}$ be the
scalar-valued solution jets, so $\mathcal S_j=\gLie\otimes
\mathcal S_j^{\mathrm{sc}}$, and write
$U_j=(\mathcal S_j^{\mathrm{sc}})^*$ for the scalar order-$j$ probes and
$Q(p,q)=\langle X_j(p),X_k(q)\rangle_{\gLie}$ for the corresponding
quadratic polynomial. The invariant metric on a simple Lie
algebra is unique up to scalar. Hence $\mathcal M_4^{\mathrm{Max}}$
is spanned by the unit and $Q(p,q)$ with $p,q\in U_0$;
weight five adds exactly the mixed quadratics with $p\in U_0,q\in U_1$.
Inversion of all coordinates acts by $(-1)^e$ at weight $e$.
The filtration and vector parity therefore imply that
$A_bQ(p,p)$ is a mixed quadratic. It has a unique representation by a
linear map $D_p:\mathcal S_0^{\mathrm{sc}}\to U_1$:
\[
 A_bQ(p,p)=\sum_\lambda Q(e^\lambda,D_p(e_\lambda)).
\]
Uniqueness follows by assigning independent zero- and first-order
solution jets in a single unit direction in $\gLie$. Here and until the coefficient comparison,
$b$ is fixed.

Denote the scalar covariance of probes by $C(p,q;z)$.
The three-quadratic Wick rule is
\begin{align*}
&S_3^{\mathrm{Max}}(Q(a,v),Q(p,p),Q(q,q);x,y,z)\\
&\quad=4d_{\gLie}\bigl[C(a,p;x-y)C(v,q;x-z)
             +C(v,p;x-y)C(a,q;x-z)\bigr]C(p,q;y-z).
\end{align*}
In the Ward sum for three quadratics at $tu,0,z$, multiply by $t^5$
and let $t\downarrow0$. The third response has no fifth-order pole.
The one-jet covariance is odd, so the other two terms have opposite
signs. The zero-order covariance is invertible, allowing the probe at
$z$ to be chosen so its contraction against any $q$ is $q(X)$,
for any prescribed $X\in\mathcal S_0^{\mathrm{sc}}$. Consequently
\begin{equation}\label{ch03:panel}
 C(D_p(X),q;u)q(X)=C(D_q(X),p;u)p(X).
\end{equation}
If $p(X)=0$, varying $q$ with $q(X)\ne0$ shows that $D_p(X)=0$.
The separation used here follows from \eqref{ch03:kelvin}: a probe
pairing to zero with all zero-order probes at every displacement has
zero response wave, hence is zero. Thus
$D_p(X)=p(X)f_b(p)$ for some $f_b(p)\in U_1$.
Cancellation in \eqref{ch03:panel} gives
\[
 C(f_b(p),q;u)=C(f_b(q),p;u).
\]
This identity and separation prove additivity and real homogeneity of
$f_b$; the original signed-coordinate covariance passes to $f_b$.

For completeness the remaining coefficient determination is very small.
Write $c_{ij}$ for the probe $F_{ij}$ and $\partial_a c_{ij}$ for its
first derivative. Parity allows only
$\partial_1c_{12},\partial_3c_{23},\partial_4c_{24}$ in $f_1(c_{12})$.
The swap of axes $3,4$ and the divergence equation reduce their span to
$\partial_1c_{12}$. Similarly the candidates for $f_1(c_{23})$ are
$\partial_1c_{23},\partial_2c_{13},\partial_3c_{12}$;
the swap of axes $2,3$ and the Bianchi equation reduce them to
$\partial_1c_{23}$. The two coefficients agree by reciprocity:
differentiation of the covariance gives
$C(\partial_1c_{12},c_{23};e_1+e_3)=-1/(2\pi^2)\ne0$.
Signed permutations then give $f_b=c_0\partial_b$ on every probe.
Since $\partial_bQ(p,p)=2Q(p,\partial_bp)$, we have
$A_b=c\partial_b$ on weight four, with $c=c_0/2$.

Set $B_b=A_b-c\partial_b$. It kills the stress and still has zero
Ward sums. Write its Ward sum for a word with an additional stress
insertion. Apply the flux formula \eqref{ch03:flux} to this insertion
with the dilation vector field $v(y)=y-x_0$. The term acting on the stress is zero. The free Ward identity
then gives the Ward sum with $\mathscr E B_bP$ in the distinguished slot
and $\mathscr EP$ in the other terms. Subtract the ordinary Ward sum
for the word starting with $\mathscr EP$. It follows that
$[\mathscr E,B_b]P\in\mathcal N_d$. Proposition
\ref{ch03:uniform-detector} makes this commutator zero, even though its
weight can be $d+1$. Thus $B_b$ preserves each weight eigenspace.
But vector covariance makes it anticommute with coordinate inversion,
which is the scalar $(-1)^e$ on that eigenspace. Therefore $B_b=0$.

Finally, if every configuration derivative of a Ward sum vanishes at
all strict collinear configurations along one axis other than $b$,
real analyticity implies that it vanishes everywhere. One can connect
any separated configuration to such a collinear configuration without
collisions: order its projections on that axis, break ties by tiny
ordered perturbations, and interpolate to the resulting line points.
The projected differences keep a strict sign until the separated endpoint.
Analytic continuation along the segment proves the assertion.
The space of permitted maps is finite-dimensional over $\RR$.
An intersection of kernels of linear functionals on it is already the
intersection of a finite subfamily (take a basis of their span in the
dual). Real and imaginary parts of the evaluations give the required
finite family.
\end{proof}

\begin{theorem}[A vanishing criterion for an internal rotation defect]\label{ch03:internal-tests}
For $d\ge4$, let
$\Delta:\mathfrak{so}(4)\times\mathfrak{so}(4)
\to\operatorname{End}_{\CC}(\mathcal M_d^{\mathrm{Max}})$
be real-bilinear and alternating. Assume its values preserve reality and
the weight filtration and kill the unit, and that
\[
 \Delta(sKs^{-1},sLs^{-1})=\rho(s)\Delta(K,L)\rho(s)^{-1}
\]
for orientation-preserving signed coordinate permutations $s$.
If every value of $\Delta$ has zero Ward sums, then $\Delta=0$.
It suffices to check that finitely many of these Ward sums vanish.
\end{theorem}

\begin{proof}
Fix $K,L$ and put $U=\Delta(K,L)$. At weight four write
$B_x(P,Q)=S_2^{\mathrm{Max}}(P,Q;x,0)$.
Recall $R_x=I-2xx^{\mathsf T}/|x|^2$, the reflection normal to $x$.
The zero-order covariance is
$\pi^{-2}|x|^{-4}\bigwedge^2R_x$ tensored with the identity on $\gLie$.
Wick's rule consequently gives, for $|x|=|y|$,
\[
 B_x(P,Q)=B_y(\rho(R_yR_x)P,Q)
 \qquad(P,Q\in\mathcal M_4^{\mathrm{Max}}).
\]
The two-point Ward sum says that $U$ is skew for every $B_x$.
These forms are nondegenerate by strict reflected positivity.
Comparing skewness for $x$ and $y$ therefore shows that $U$ commutes
with $\rho(R_xR_y)$ on this space. In particular it commutes with each
coordinate quarter-turn $q_{ab}$, by choosing $x$ parallel to $e_a$
and $y$ parallel to $e_a+e_b$.

An alternating bilinear map on $\mathfrak{so}(4)$ invariant under
simultaneous conjugation by these quarter-turns is zero, even when its
values lie in an arbitrary real vector space. To see this on the basis
$E_{ab}=e_ae_b^{\mathsf T}-e_be_a^{\mathsf T}$, a half-turn in the
plane spanned by the third and fourth axes fixes one member of a pair
sharing an axis and negates the other. For a disjoint pair, the even
permutation exchanging their two ordered axis pairs exchanges the
arguments, so alternation gives zero. Equal pairs vanish by alternation.
Bilinearity proves the claim. Covariance and the commutation just proved
apply this observation to $\Delta(K,L)P$ for every weight-four $P$.
In particular $U$ kills the stress.

For each $H\in\mathfrak{so}(4)$, repeat the preceding commutator
argument using the rotation vector field $v(y)=H(y-x_0)$ in the flux
formula \eqref{ch03:flux}. This gives $[\mathcal L_H,U]P\in\mathcal N_d$.
This time the commutator itself lies in $\mathcal M_d^{\mathrm{Max}}$,
so a single reflected two-point pairing shows that it vanishes. The uniform
bound $D_*$ is unnecessary. Exponentiation in the finite-dimensional
space shows that $U$ commutes with every quarter-turn on the whole
space. Covariance and the same alternating-map argument give
$\Delta(K,L)P=0$ for every $P$. The finiteness assertion again follows
by taking a basis of the span of the Ward functionals in the dual of
the finite-dimensional space of allowed $\Delta$.
\end{proof}

Each criterion involves finitely many evaluations on a specified space of possible responses. Their
application to the interacting construction requires proving that its
responses obey the stated filtration, covariance, and comparison
hypotheses. No such interacting input is part of the free-field results
proved here.
}{}
\setcounter{chapter}{3}
\IfFileExists{chapters/04-blocking.tex}{% !TEX root = ../monograph.tex
\chapter{The renormalization transformation and the effective density}
\label{ch04:blocking}

A renormalization transformation, called a step, is first an exact integration, and only then an
estimate. We define that integration on every fine configuration, including
large fields. Its output is generally not a Wilson density. We therefore
also introduce a representation of the effective density: a prescribed leading
action, a normalization constant, localized terms on small-field regions,
and a sum over polymers covering the large-field sites.

We use the gauge group and Haar measure of Chapter~\ref{ch01:lattice}
and the normalized matrix norm $\|\cdot\|_G$ of \eqref{ch01:matrix-norm}.
Left or right multiplication by a group element preserves this norm.
Recall that $a_m=D_m^{-1}$ is the lattice
spacing and $R_p$ the period, both in continuum units, and the odd ratio $L_m$ relates
the scales with indices $m+1$ and $m$. In this chapter we examine a single step
$m+1\to m$; unlike in~\cite{Balaban1984Propagators}, a step is not followed by a rescaling to the unit lattice. Write $L=L_m$, fix a positive integer \emph{smoothing order} $j$ (the number of convolution factors in the averaging weight $w^{\rm av}$ of Section~\ref{ch04:kernel-section}),
and set $R=(j+3)L$. Fine links $U$ live on $\Lambda_{R_pD_{m+1}}$ and coarse
links $V$ on $\Lambda_{R_pD_m}$. Patches are first described on the infinite
lattice $\ZZ^4$ and then placed periodically. The constants and path coefficients
introduced here are inputs; their simultaneous choice along a renormalization-group trajectory
will require the later estimates.

\section{Local gauge transformations and a symmetric choice}
\label{ch04:frames}

A choice of gauge on a patch assigns a frame to each site. Relative to
the original frames, this choice is specified by a family of matrices
$h_x\in G$: a contour variable $U_{xy}$ then has matrix $h_xU_{xy}h_y^{-1}$.
We choose $h_0=1$, so that $h_x$ expresses vectors at $x$ in the frame
at the patch centre. This lets us compare contour variables of the fine field before averaging. A deterministic axial gauge is convenient for bounds
but singles out coordinate directions. We combine a symmetric minimizing
choice near the flat field with a finite randomized choice elsewhere.
The use of gauge-covariant averages in a renormalization transformation belongs to
Ba\l{}aban's renormalization-group method~\cite{Balaban1987RG}; the precise
branch rules below are part of the present construction.

Let $\mathbb B_R=\{-R,\ldots,R\}^4$. For a permutation $o$ of the four
coordinate directions, denote by $T_U(x\to y;o)$ the contour variable along the
contour which changes coordinates in order $o$. For an initial point $z\in\mathbb B_R$,
put
\begin{equation}
 g_x^{z,o}=T_U(z\to x;o),\qquad
 h_x^{z,o}=(g_0^{z,o})^{-1}g_x^{z,o}.
 \label{ch04:axial-frame}
\end{equation}
These are \emph{local gauge transformations}, that is, $G$-valued functions on the sites of $\mathbb B_R$; they bring $U$ to an axial gauge with initial point $z$, and we call them \emph{axial}. The normalized one satisfies $h_0^{z,o}=1$. The unnormalized one gives
link coordinates
$U_{(x,\mu)}^{z,o}=g_x^{z,o}U(x,\mu)(g_{x+e_\mu}^{z,o})^{-1}$.
Their distances from the identity are unchanged if the normalized transformation is used instead,
since that change conjugates every transformed link by the same matrix.

A symmetric choice is possible if all relevant holonomies are close to the identity.
For $x,y\in\mathbb B_R$ and orders $o,o_0$, define the holonomy based at the
patch center by
\[
 \mathcal T_{xy;o}^{o_0}
   =g_x^{0,o_0}T_U(x\to y;o)(g_y^{0,o_0})^{-1}.
\]
Gauge transformations conjugate this holonomy, preserving its distance from the identity. Fix a small constant $\delta_{\rm fr}>0$.
If $\|\mathcal T_{xy;o}^{o_0}-1\|_G<\delta_{\rm fr}$ for every
$x,y\in\mathbb B_R$ and every $o,o_0\in S_4$, use the
normalized minimizer, with $h_0=1$, of
\begin{equation}
 \mathcal E_U(h)=\frac1{48|\mathbb B_R|^2}
   \sum_{x,y\in\mathbb B_R}\sum_{o\in S_4}
       \bigl\|h_xT_U(x\to y;o)-h_y\bigr\|_G^2.
 \label{ch04:frame-energy}
\end{equation}
The condition that all these distances be strictly below $\delta_{\rm fr}$, the \emph{loop condition}, includes every order $o_0$, so it does not privilege a preliminary
axial gauge.

If the loop condition fails, use axial gauge transformations. Fix $a_{\rm ax}\ge16$ and call a pair
$(z,o)$ eligible when
\begin{equation}
 \sum_{e\text{ internal to }\mathbb B_R}\|U_e^{z,o}-1\|_G
 \le a_{\rm ax}(2R+1)
      \sum_{\square\text{ internal to }\mathbb B_R}\|U_\square-1\|_G.
 \label{ch04:axial-eligible}
\end{equation}
Choose one pair $(z,o)\in\mathbb B_R\times S_4$ uniformly among all
eligible pairs. Use this same initial point and order throughout the patch:
the gauge at each $x\in\mathbb B_R$ is given by $h_x^{z,o}$ in
\eqref{ch04:axial-frame}. The coordinate paths from $z$ form a spanning
tree, and in this gauge every link of that tree is the identity. On the minimizing branch give all
pairs the same uniform probability weight, but let each pair specify the same minimizing
gauge transformation. The redundant values of the index $(z,o)$ allow both branches to be described by one
finite sum.

The following local assertion ensures that these rules are well-defined.
Its uniformity in the size of the patch is essential.

\begin{lemma}[Uniform choice of the local gauge transformation]
\label{ch04:frame-lemma}
There is a $\delta_{\rm fr}>0$, depending only on the matrix group, for
which the minimizer of \eqref{ch04:frame-energy} with $h_0=1$ is unique
whenever the loop condition holds. For every configuration and every fixed
coordinate order, at least one initial point satisfies \eqref{ch04:axial-eligible}
with $a_{\rm ax}=16$. These constants do not depend on $R$.
\end{lemma}
\begin{proof}
We first record the geometric input. The nearest-point projection $\pi$
onto the compact smooth submanifold $G$ is uniquely defined and smooth
on a uniform tubular neighborhood~\cite[p.~153, Lemma]{Foote1984}.
Its derivative at $1$ is the orthogonal projection $\Pi_{\gLie}$ onto the
Lie algebra. Hence on a sufficiently small convex ball about $1$,
\begin{equation}
 \|\pi(A)-\pi(B)-\Pi_{\gLie}(A-B)\|_G
       \le\tfrac14\|A-B\|_G.
 \label{ch04:projection-bound}
\end{equation}
This is the mean-value estimate for $D\pi-\Pi_{\gLie}$, not a new
projection theorem.

Fix $o_0\in S_4$ and use the axial gauge with initial point at the origin, with this
coordinate order. In this gauge the contour variable from $x$ to $y$ in order
$o$ is $\mathcal T_{xy;o}^{o_0}$. On the minimizing branch the loop condition
gives $\|\mathcal T_{xy;o}^{o_0}-1\|_G<\delta_{\rm fr}$.
Reversing a contour inverts its contour variable and permutes the coordinate orders.
A minimizer exists by compactness. Common left multiplication leaves the
energy invariant, so minimization with the normalization $h_0=1$ also gives
an unconstrained minimizer after allowing that common multiplier.
Comparison with the constant gauge transformation and the squared triangle inequality give
\[
 \frac{1}{|\mathbb B_R|^2}\sum_{x,y\in\mathbb B_R}\|h_x-h_y\|_G^2\le4\delta_{\rm fr}^2.
\]
Multiply all $h_x$ on the left by one group element
so that $1$ is nearest to their centroid $\bar h=|\mathbb B_R|^{-1}\sum_{x\in\mathbb B_R} h_x$.
Then $\|\bar h-1\|_G\le2\delta_{\rm fr}$. For each fixed $x$, minimization over the value at that point shows
that $h_x$ is nearest to
\[
 F_x(h)=\frac{1}{24|\mathbb B_R|}\sum_{y\in\mathbb B_R}\sum_{o\in S_4}h_y(\mathcal T_{xy;o}^{o_0})^{-1}.
\]
Indeed, when varying $h_x$ to $v$, retaining the old $h_x$ in the
diagonal term $y=x$ adds a nonnegative multiple of $\|v-h_x\|_G^2$,
which vanishes at $v=h_x$.
Since $\|F_x(h)-\bar h\|_G\le\delta_{\rm fr}$, all local means $F_x(h)$ lie
within $3\delta_{\rm fr}$ of $1$. For small enough $\delta_{\rm fr}$ they lie in the projection ball,
and $h_x=\pi(F_x(h))$.

Multiply two minimizers $h,k$ on the left in this way, so that $1$ is nearest to each centroid, and set
$u=\max_x\|h_x-k_x\|_G$. The differences of their local means satisfy
\[
 \|F_x(h)-F_x(k)\|_G\le u,\qquad
 \|F_x(h)-F_x(k)-(\bar h-\bar k)\|_G\le \sqrt{n_G}\,\delta_{\rm fr} u.
\]
For the second inequality, apply $\|AB\|_G\le\|A\|_G\|B\|_{\rm op}$
to each summand and use $\|B\|_{\rm op}\le\sqrt{n_G}\|B\|_G$.
The differences $\bar h-1$ and $\bar k-1$ are normal to $G$ at $1$, so
$\Pi_{\gLie}(\bar h-\bar k)=0$. At an index attaining the maximum,
\eqref{ch04:projection-bound} gives $u\le(\sqrt{n_G}\,\delta_{\rm fr}+1/4)u$.
Choosing $\sqrt{n_G}\,\delta_{\rm fr}\le1/4$
forces $u=0$. Thus minimizers differ only by a common left multiplier,
removed by $h_0=1$. Only averages occur in the contraction, which explains
why its constants do not grow with the patch.

For the axial gauge transformations, apply the usual plaquette telescoping
argument to a rectangular slice $X$ with $r$ varying coordinates. Write
$\Sigma_{\rm link}(X,U)$ and $\Sigma_{\rm plaq}(X,U)$ for the sums of
$\|U_e-1\|_G$ over its internal positive links and
$\|U_\square-1\|_G$ over its internal plaquettes, respectively. Choose a
hyperplane $X_0$ of the slice perpendicular to the last coordinate whose internal plaquette sum
is no greater than the average over such hyperplanes. Use the axial construction on
that hyperplane, then extend it along the last coordinate, setting those links
to one. Every remaining link differs from its translate, along the last coordinate, in the
chosen hyperplane by a product of intervening plaquettes. By the triangle
inequality and unitary invariance of $\|\cdot\|_G$, with $n=2R+1$ we obtain
\[
 \Sigma_{\rm link}(X,U^{\rm ax})
 \le n\bigl(\Sigma_{\rm link}(X_0,U^{\rm ax})+
                            \Sigma_{\rm plaq}(X,U)\bigr).
\]
Induction with the bound $4rn\Sigma_{\rm plaq}(X,U)$ closes because
$n\Sigma_{\rm plaq}(X_0,U)\le\Sigma_{\rm plaq}(X,U)$.
The construction follows the prescribed coordinate order; the chosen
hyperplanes of the successive slices, one in each dimension, specify an initial point. At $r=4$ the bound is \eqref{ch04:axial-eligible}
with constant $16$.
\end{proof}

Only smoothness of the nearest-point projection was used here. Analytic
estimates for the critical configuration will need more; they are not consequences of
this lemma alone.

\section{The gauge-covariant averaging operation and the exact kernel}
\label{ch04:kernel-section}

Translate the preceding construction of local gauge transformations to every patch
$Ly+\mathbb B_R$ centered at the fine site $Ly$ corresponding to the coarse site $y$, the \emph{patch of the averaging operation} at $y$.
For $z\in\mathbb B_R$, let $h_{y,z}$ be the gauge matrix at $Ly+z$,
normalized by $h_{y,0}=1$. For a coarse edge $e=(y,\mu)$ define the \emph{gauge-transformed parallel transporter}, for short the \emph{transporter},
\begin{equation}
 Y_{e,z}=h_{y,z}
       \left(\prod_{r=0}^{L-1}U(Ly+z+re_\mu,\mu)\right)
                     h_{y+e_\mu,z}^{-1}.
 \label{ch04:connector}
\end{equation}
Both ends are transformed by the local gauge transformations of the coarse endpoints, as shown
for two offsets in Figure~\ref{ch04:fig-block-paths}. In
particular, under a fine gauge transformation $g$, this transporter becomes
$g(Ly)Y_{e,z}g(Ly+Le_\mu)^{-1}$.

Let the \emph{averaging weight} $w^{\rm av}$ be the $j$-fold convolution of the uniform probability
on $\mathbb B_{(L-1)/2}$. It is nonnegative, has sum one, is invariant under
$H_4$, and is supported in $\mathbb B_{j(L-1)/2}\subset\mathbb B_R$.
Put $\bar Y_e=\sum_{z\in\mathbb B_R} w_z^{\rm av}Y_{e,z}$. On the first of its two branches, the value of the gauge-covariant \emph{averaging operation} is obtained by projecting the weighted
average of the transporters to the group; it then minimizes
\begin{equation}
 \mathcal P_e(u)=\tfrac12\sum_z w_z^{\rm av}\|Y_{e,z}-u\|_G^2,\qquad u\in G.
 \label{ch04:predictor-cost}
\end{equation}
Expanding about the weighted mean gives
\[
 2\mathcal P_e(u)
 =\sum_z w_z^{\rm av}\|Y_{e,z}-\bar Y_e\|_G^2
                         +\|\bar Y_e-u\|_G^2.
\]
Thus the minimizers are precisely the nearest group points to $\bar Y_e$.
If $\min_{u\in G}\mathcal P_e(u)<\varepsilon_{\rm pred}$, the identity
above gives $\min_{u\in G}\|\bar Y_e-u\|_G
<\sqrt{2\varepsilon_{\rm pred}}$. Choose $\varepsilon_{\rm pred}$ small
enough that $\bar Y_e$ then lies in the tubular neighborhood used in
Lemma~\ref{ch04:frame-lemma}. The value of the averaging operation is then unique, independently
of the number of transporters or the weights.

On this branch set $Q_e^{\rm pred}$ to that minimizer and assign uniform
probability weight to the redundant values $z\in\mathbb B_R$ of the index of the averaging operation. If the condition on the minimum fails,
let the index take the value $z$, and set $Q_e^{\rm pred}=Y_{e,z}$, with probability
$w_z^{\rm av}$. Equality in either the loop condition or the condition on
$\min_{u\in G}\mathcal P_e(u)$ selects the corresponding axial gauge transformation
or transporter. The averaging operation $U\mapsto Q_e^{\rm pred}$ is an analogue of Ba\l{}aban's averaging operation $\bar U$~\cite{Balaban1987RG}, not the same map: on the minimizing branch it is the nearest group point to a weighted mean, where Ba\l{}aban takes the exponential of a mean of logarithms.

\begin{figure}[htbp]
\centering
\begin{tikzpicture}[
    x=.85cm,y=.85cm,font=\small,
    transport/.style={draw=ymInk,line width=1.1pt,
      decoration={markings,mark=at position .64 with {\arrow{Stealth[length=2.6mm]}}},
      postaction={decorate}},
    frame/.style={draw=ymAccent,line width=.9pt,densely dashed,
      -{Stealth[length=2.2mm]},shorten <=2pt,shorten >=2pt}]
  \foreach \x in {0,...,8}
    \draw[ymGrid,thin] (\x,-1.3) -- (\x,2.3);
  \foreach \y in {-1,...,2}
    \draw[ymGrid,thin] (-.35,\y) -- (8.35,\y);

  % Two translated fine paths, each with seven bonds in the displayed slice.
  \draw[transport] (1,2) -- (8,2);
  \draw[transport] (0,-1) -- (7,-1);
  \foreach \r in {0,...,7}{
    \fill[ymInk] (1+\r,2) circle (1.5pt);
    \fill[ymInk] (\r,-1) circle (1.5pt);
  }
  \node[above=2pt] at (1,2) {$Ly+z$};
  \node[above=2pt] at (8,2) {$L(y+e_\mu)+z$};
  \node[above=2pt] at (4.5,2) {$L$ fine links};

  % Curves denote frame factors even when the frame is defined by minimization.
  \draw[frame] (0,0) to[bend left=16]
    node[pos=.50,left=3pt,fill=white,inner sep=1pt] {$h_{y,z}$} (1,2);
  \draw[frame] (8,2) to[bend left=16]
    node[pos=.50,right=3pt,fill=white,inner sep=1pt] {$h_{y+e_\mu,z}^{-1}$} (7,0);
  \draw[frame] (0,0) to[bend right=20] (0,-1);
  \draw[frame] (7,-1) to[bend right=20] (7,0);
  \foreach \x in {0,7}
    \fill (\x,0) ++(-2pt,-2pt) rectangle ++(4pt,4pt);
  \node[left=4pt,fill=white,inner sep=1pt] at (0,0) {$Ly$};
  \node[right=4pt,fill=white,inner sep=1pt] at (7,0) {$L(y+e_\mu)$};
\end{tikzpicture}
\caption[Fine paths and endpoint frames in the block average]{Two connectors in the matrix average $\bar Y_e$ for $e=(y,\mu)$,
shown in a coordinate plane. Each solid path traverses $L$ fine links.
The dashed arrows are frame factors, which can depend on the whole
patch. They give the framed transport $Y_{e,z}$ the same coarse
endpoints for every offset $z$.}
\label{ch04:fig-block-paths}
\end{figure}


Collect all choices made for the local gauge transformations and for the averaging operation in a finite \emph{auxiliary index} $\alpha$ and
denote their joint probability, the probability weight of $\alpha$, by $\lambda(U,\alpha)$. The choices for the local gauge transformations are
made first, followed by conditionally independent choices for the averaging operation. Thus
\begin{equation}
 \lambda(U,\alpha)\ge0,\qquad \sum_\alpha\lambda(U,\alpha)=1.
 \label{ch04:tag-law}
\end{equation}
Lemma~\ref{ch04:frame-lemma} supplies an eligible pair whenever it is needed.

Fix the \emph{penalty parameter} $b>0$ of the block kernel.\martin{Apparently $b$ will be fixed throughout and could just be
taken as $b=1$. It controls the soft conditioning used here.} The renormalization transformations for gauge fields in~\cite{Balaban1984Propagators,Balaban1985Averaging,Balaban1987RG} have a delta function in place of the densities $q_{b\beta}$ used below, so that there the coarse field equals the average exactly. For $\gamma\ge0$, let
\[
 q_\gamma(u)=Z_\gamma^{-1}e^{-\gamma\|u-1\|_G^2/2},\qquad
 Z_\gamma=\int_G e^{-\gamma\|v-1\|_G^2/2}\,dh(v).
\]
This is a probability density on $G$ with respect to Haar measure.

For a fixed fine field $U$ and choice $\alpha$, the value of the averaging operation
$Q_e=Q_e^{\rm pred}(U,\alpha)$ represents transport across the coarse
edge $e$. We impose agreement of $V_e$ with $Q_e$ by a soft constraint.
By unitary invariance its energy is $\tfrac12 b\beta\|V_e-Q_e\|_G^2$.
This penalizes the discrepancy in a gauge-invariant way, since $V_e$
and $Q_e$ transform by the same matrices at their endpoints. At large
$b\beta$ the conditional law thus concentrates near the value of the averaging operation.
Given $U$ and $\alpha$, we sample the coarse links independently with
these normalized densities. Averaging their product over $\alpha$ gives
the conditional density of $V$ given $U$; integrating against the fine
density gives its coarse marginal.

\begin{definition}[Exact block map]
\label{ch04:block-map}
The superscript $(L)$ denotes the block ratio: on a coarse torus
$\Lambda_s$ the fine torus is $\Lambda_{Ls}$. The period $s$, smoothing
order $j$ and penalty parameter $b$ are understood. For $\beta\ge0$, define the
block kernel of the step and its action on fine densities by
\begin{align}
 k^{(L)}_\beta(U,V)
   &=\sum_\alpha\lambda(U,\alpha)
        \prod_{e\text{ coarse}}
           q_{b\beta}\bigl(V_eQ_e^{\rm pred}(U,\alpha)^{-1}\bigr),
 \label{ch04:kernel}\\
 (\mathcal B^{(L)}_\beta f)(V)
   &=\int k^{(L)}_\beta(U,V)f(U)\,dU,
 \label{ch04:block-operator}
\end{align}
where $dU$ and $dV$ denote product probability Haar measure on their
respective link spaces.
\end{definition}

For fixed $\beta$ this is a bounded Borel kernel. Indeed, the loop
inequalities of Section~\ref{ch04:frames} and the axial eligibility
inequality \eqref{ch04:axial-eligible} compare continuous functions of
the fine links. Once the local gauge transformations are chosen, the condition for using the
minimizing branch of the averaging operation on edge $e$ is
$\min_{u\in G}\mathcal P_e(u)<\varepsilon_{\rm pred}$. This minimum is
continuous in the transporters $(Y_{e,z})_z$ by compactness of $G$.
On each minimizing branch the unique minimizer depends continuously on
its data. Thus the choices and their probabilities in the finite sum
\eqref{ch04:kernel} are Borel functions of $U$. Boundedness follows
from boundedness of $q_{b\beta}$ and \eqref{ch04:tag-law}.

\begin{proposition}[Normalization and covariance]
\label{ch04:kernel-properties}
The kernel is nonnegative and normalized in the coarse variable:
\begin{equation}
 \int k^{(L)}_\beta(U,V)\,dV=1.
 \label{ch04:kernel-normalization}
\end{equation}
It is covariant under fine gauge transformations and the induced coarse
gauge transformations $y\mapsto g(Ly)$. It is also covariant under signed
coordinate permutations and corresponding fine/coarse translations.
For every integrable fine density $f$,
\begin{equation}
 \int\mathcal B^{(L)}_\beta f\,dV=\int f\,dU.
 \label{ch04:exact-integral}
\end{equation}
\end{proposition}
\begin{proof}
For normalization, integrate each coarse-link factor using right invariance
of Haar measure, then sum \eqref{ch04:tag-law}. Fubini gives
\eqref{ch04:exact-integral}, with absolute integrability following from
the same computation applied to $|f|$.
For covariance, gauge transformations conjugate the holonomies of the loop conditions and
transform the endpoints of \eqref{ch04:connector}; signed lattice
isometries permute initial points, orders and offsets and invert reversed links.
The objective functions and eligibility conditions are unchanged. Uniqueness
on the minimizing branches, and reindexing of the finite alternatives,
then give the stated covariance of the kernel.
\end{proof}

In particular, if $U$ has normalized density $f/\int f$, the coarse
field has density $\mathcal B^{(L)}_\beta f/\int f$. The equality of
integrals \eqref{ch04:exact-integral} holds even when the coarse action
has terms beyond the Wilson plaquette action.

\section{A linear bound on coarse plaquette distances}
\label{ch04:predictor-geometry}

We bound the distance of a coarse plaquette holonomy from the identity
in terms of the fine plaquette holonomies in a surrounding box and the
discrepancies between the four coarse boundary links and the values of
the averaging operation on them. Fix a site $y$ of the coarse torus
$\Lambda_s$. Choose a lift $\tilde y\in\ZZ^4$ and form the fine lattice box
\[
 M_y=L\tilde y+\mathbb B_{(j+7)L}.
\]
For a fine plaquette $q$ on $\Lambda_{Ls}$, let $m_y(q)$ be the number
of its lifts with all four vertices in $M_y$. Changing the lift of $y$
translates the whole box by a fine period and leaves this count unchanged.
If $s>2(j+7)$, the box can contain at most one lift of each plaquette. Write
\[
 \mathcal P_y(U)=\sum_{q}m_y(q)\|U_q-1\|_G,
\]
where the sum ranges over all fine plaquettes on the torus, with one
orientation per square.

\begin{lemma}[Coarse plaquette estimate]
\label{ch04:predictor-comparison}
Fix the smoothing order $j$ and the constants in the choices of the local
gauge transformations and of the averaging operation in
Sections~\ref{ch04:frames} and~\ref{ch04:kernel-section}.
Suppose $s\ge100(j+8)$. Let $p$ be a coarse plaquette whose four
vertices lie within periodic sup distance $2$ of $y$ in coarse lattice
units. In particular, $y$ may be any vertex of $p$. For every fine field
$U$, coarse field $V$, and choice $\alpha$ with $\lambda(U,\alpha)>0$,
write $Q_e(U)=Q_e^{\rm pred}(U,\alpha)$. Then
\begin{equation}
 \|V_p-1\|_G\le\sum_{e\in\partial p}\|V_e-Q_e(U)\|_G
                         +C_jL^{-2}\mathcal P_y(U),
 \qquad \sup_q\sum_{y\in\Lambda_s}m_y(q)\le(2j+15)^4 .
 \label{ch04:macro-comparison}
\end{equation}
The constant $C_j$ is independent of $L$ and $s$.
\end{lemma}

\begin{lemma}[Mean distance from a constant gauge transformation]
\label{ch04:minimizing-frame-bound}
Let $U$ be a link field on a finite lattice box $X$, and let $h$ minimize
the objective \eqref{ch04:frame-energy} on $X$. Take uniform averages
over the indicated sites $x,z\in X$ and coordinate orders $o\in S_4$.
There exists $c\in G$ such that
\begin{equation}
 \operatorname{av}_x\|h_x-c\|_G
 \le C_G\operatorname{av}_{x,z,o}\|T_U(x\to z;o)-1\|_G ,
 \label{ch04:frame-first-power}
\end{equation}
where $C_G$ depends only on $G$.
\end{lemma}
\begin{proof}
Put $d_x=\operatorname{av}_{z,o}\|T_U(x\to z;o)-1\|_G$ and
$\bar d=\operatorname{av}_x d_x$. The constant gauge transformation has energy at most $\bar d$, since each path distance
is at most two. Minimality and the squared triangle inequality give
$\operatorname{av}_{x,z}\|h_x-h_z\|_G^2\le8\bar d$.
Let $m=\operatorname{av}_z h_z$ and
$m_x=\operatorname{av}_{z,o}h_zT_U(x\to z;o)^{-1}$, so
$\|m_x-m\|_G\le d_x$. As in the proof of Lemma~\ref{ch04:frame-lemma},
varying $h_x$ shows that it is a nearest group element to $m_x$.
Indeed, reversing the contour interchanges the two off-diagonal terms
involving $x$. Freezing the second endpoint of the diagonal term adds
a nonnegative squared distance that vanishes at $h_x$.

Choose $c$ nearest to $m$. Since each $h_x$ lies in $G$, Jensen's
inequality gives $\|m-c\|_G^2\le\operatorname{av}_x\|m-h_x\|_G^2
\le8\bar d$. Choose $\varepsilon>0$ so that the nearest-point
projection is Lipschitz on the ball of radius $2\varepsilon$ about $c$,
uniformly in $c\in G$, as in Lemma~\ref{ch04:frame-lemma}.
First suppose $8\bar d<\varepsilon^2$. At sites with $d_x<\varepsilon$
both $m$ and $m_x$ lie in this ball, so $\|h_x-c\|_G\le C_Gd_x$.
At the remaining sites use $\|h_x-c\|_G\le2\le2d_x/\varepsilon$.
Averaging gives \eqref{ch04:frame-first-power}. For
$8\bar d\ge\varepsilon^2$ the bound follows from $\|h_x-c\|_G\le2$.
The same proof applies on any finite
set for a uniformly averaged family of paths paired under reversal
with identity contour variables on the diagonal.
\end{proof}

\begin{proof}[Proof of Lemma~\ref{ch04:predictor-comparison}]
Let $X$ be a fine lattice box with $\ell$ sites in each direction and
$N=\ell^4$. For a link field $U$ on $X$, define the sum over its internal
positive links by
\[
 E_X(U)=\sum_{\mu=1}^4\ \sum_{\substack{x\in X\\x+e_\mu\in X}}
                       \|U(x,\mu)-1\|_G .
\]
For each coordinate order $o$ we have
\[
 \sum_{x,z\in X}\|T_U(x\to z;o)-1\|_G\le N\ell E_X(U).
\]
Indeed, choose $x,z$ independently and uniformly in $X$. On the
segment in direction $\mu$, the other three coordinates are uniform on
the $\ell^3$ transverse sites. Each is taken from $x$ or $z$ according
to the order $o$. The triangle inequality bounds the segment distance by
the sum of link distances along its entire coordinate line. Its expectation
is therefore at most $\ell^{-3}$ times the sum over all positive
$\mu$-links in $X$. Sum over $\mu$ and multiply by $N^2$.

Let $H$ be any gauge transformation on $X$. If $f$ is the minimizing local gauge transformation for $U$,
then $h_x=f_xH_x^{-1}$ minimizes the objective \eqref{ch04:frame-energy} for $H\cdot U$,
since substituting $h$ gives the same energy term by term. A common
left multiplication can restore the normalization at the centre.
Applying Lemma~\ref{ch04:minimizing-frame-bound} and the path bound gives
\[
 \sum_{x\in X}\|f_x-cH_x\|_G\le C_G\ell E_X(H\cdot U)
\]
for some $c\in G$. For an eligible axial gauge transformation $f$, put $g_x=f_xH_x^{-1}$.
The pure gauge $g\cdot1$ has contour variable $g_xg_z^{-1}$ from $x$ to $z$,
and the linkwise triangle inequality gives
$E_X(g\cdot1)\le E_X(f\cdot U)+E_X(H\cdot U)$. Applying the path
bound to $g\cdot1$ and averaging over $z$ yields some $c=g_z$ with
\[
 \sum_{x\in X}\|f_x-cH_x\|_G
 \le \ell\bigl(E_X(H\cdot U)+E_X(f\cdot U)\bigr).
\]
By the eligibility condition \eqref{ch04:axial-eligible},
$E_X(f\cdot U)\le a_{\rm ax}\ell\sum_{q\subset X}\|U_q-1\|_G$,
where the sum runs over plaquettes with all four vertices in $X$.

Now choose the axial gauge transformation $H$ of Lemma~\ref{ch04:frame-lemma} on $M_y$.
Since this box has side $2(j+7)L+1$, the axial estimate gives
$E_{M_y}(H\cdot U)\le C_j L\mathcal P_y(U)$. Lift the four vertices
$z$ of $p$ to $\tilde y+\mathbb B_2$. Their patches of the averaging operation
$Lz+\mathbb B_R$, with $R=(j+3)L$, and the straight paths between them
all lie in $M_y$. Each patch has side $2R+1=O_j(L)$. Applying the
preceding bounds on that patch, for either choice of local gauge transformation, gives $c_z\in G$
such that
\begin{equation}
 \sum_{i\in\mathbb B_R}\|h_{z,i}-c_zH_{Lz+i}\|_G
                                  \le C_j L^2\mathcal P_y(U).
 \label{ch04:patch-frame-sum}
\end{equation}
The averaging weights satisfy $w_i^{\rm av}\le L^{-4}$. Given the
last $j-1$ uniform points, at most one of the $L^4$ possible first
points yields a prescribed sum. Thus each endpoint local gauge transformation in the
transporter \eqref{ch04:connector} contributes at most
$C_jL^{-2}\mathcal P_y(U)$ to the weighted discrepancy. For its middle
contour, at each of the $L$ steps translation of the offset $i$ is injective.
The unweighted sum of the path distances in gauge $H$ is therefore at
most $L E_{M_y}(H\cdot U)$. Combining these bounds gives, for each link
of $\partial p$ written in its positive orientation $e=(z,\mu)$,
\[
 \sum_i w_i^{\rm av}\|Y_{e,i}-c_zc_{z+e_\mu}^{-1}\|_G
                                      \le C_jL^{-2}\mathcal P_y(U).
\]

For either branch of the averaging operation and any $u_*\in G$ we have
\[
 \|Q_e(U)-u_*\|_G\le2\max\{1,\varepsilon_{\rm pred}^{-1}\}
                      \sum_i w_i^{\rm av}\|Y_{e,i}-u_*\|_G .
\]
On the minimizing branch $Q_e(U)$ is a nearest group point to the
weighted mean of the $Y_{e,i}$, so the triangle inequality gives the
factor two. On the other branch the minimum of
\eqref{ch04:predictor-cost} is at least $\varepsilon_{\rm pred}$.
Every group distance is at most two, so
$\varepsilon_{\rm pred}\le\mathcal P_e(u_*)\le
\sum_i w_i^{\rm av}\|Y_{e,i}-u_*\|_G$ and $\|Q_e(U)-u_*\|_G\le2$.
Take $u_*=c_zc_{z+e_\mu}^{-1}$ on each link.
These four transports multiply to the identity around $p$, with
inverses on reversed links. Replacing one factor at a time, first by
the corresponding $Q_e(U)$ and then by $V_e$, proves the first
inequality in \eqref{ch04:macro-comparison}.

For the overlap bound, fix one lift of $q$ with base vertex $a\in\ZZ^4$.
Summing $m_y(q)$ over coarse sites is equivalent to fixing this lift
and counting all centres $Lz$, $z\in\ZZ^4$, whose boxes contain it.
Indeed, a lift translated by a fine period is in $M_y$ exactly when the
fixed lift is in the oppositely translated box. Each coordinate must
satisfy $|Lz_\mu-a_\mu|\le(j+7)L$, leaving at most $2(j+7)+1=2j+15$
integer choices for $z_\mu$. The four coordinates give the stated
bound $(2j+15)^4=O(j^4)$.
\end{proof}

\section{The regularity condition and the representation of the density}
\label{ch04:density-coordinates}

The threshold used to distinguish small from large fields must remain fixed
while the inverse coupling $\beta$ is varied for tuning. To express that distinction,
let $\bar\beta>1$ and use the reference values
\begin{align}
 \widehat\beta_m&=\bar\beta+\kappa_G\tau_m,&
 M_m^{\rm cut}&=\exp\bigl((\log\log\widehat\beta_m)^2\bigr),\\
 t_m&=\widehat\beta_m^{-1/2}M_m^{\rm cut},&
 B_m&=\lceil(\log\widehat\beta_m)^2\rceil.
 \label{ch04:thresholds}
\end{align}
Here $\tau_m=\log D_m$ is the logarithmic scale from
Definition~\ref{ch01:scales}. We use these formulas only where $0<t_m\le1$.
The inverse coupling $\beta$ ranges in
$\mathfrak B_m=[\widehat\beta_m/2,2\widehat\beta_m]$.
The reference value and threshold are not recomputed when that coupling
changes. Their eventual compatibility with the schedule of Chapter~\ref{ch10:trajectory} is a later
choice.

At the present scale write $\Lambda=\Lambda_{R_pD_m}$, and let
$\operatorname{Bad}_m(U)$, the \emph{large-field set} of $U$, be the sites incident to at least one plaquette
with $\|U_\square-1\|_G>t_m$. The configuration $U$ satisfies the \emph{regularity condition} on a set $X$, which is then called a small-field region, if
$X\cap\operatorname{Bad}_m(U)=\varnothing$; write the indicator of this condition as $h_X(U)$.
Thus every plaquette meeting $X$ must have $\|U_\square-1\|_G\le t_m$. The term regularity condition is that of~\cite{Balaban1985Spaces}, where the spaces of regular configurations carry bounds on the plaquette variables, depending on the scale, on a nested sequence of domains, and a second condition; here the condition is a single inequality at one scale, and the large-field set is neither enlarged nor summed over. This definition is
gauge invariant and involves a closed condition on the sites outside the large-field set.

We keep two kinds of local corrections: the localized terms, whose sum is the \emph{regular part of the action}, and the large-field activities.

\begin{definition}[Localized terms and large-field activities]
\label{ch04:activities}
A \emph{localized term} of the regular part of the action consists of a \emph{localization domain} $X_i$, a base point $o_i\in X_i$,
and a real Borel gauge-invariant function $E_i^{\rm reg}$ local on $X_i$.
It vanishes at the flat configuration and has the required number of
ordered left derivatives on configurations satisfying the regularity condition on $X_i$.

A \emph{large-field activity} has an \emph{exclusion support} $Y_i$, a nonempty \emph{prescribed large-field
set} $I_i^{\rm bad}\subset Y_i$, and an additional Borel gauge-invariant
condition, an event $\mathcal E_i^{\rm extra}$. Its real Borel value
$K_i^{\rm lf}$ depends on links internal to $Y_i$ and on the indicators of
large-field sites in $Y_i$, and vanishes outside
\begin{equation}
 \{U:\operatorname{Bad}_m(U)\cap Y_i=I_i^{\rm bad}\}
                         \cap\mathcal E_i^{\rm extra}.
 \label{ch04:eligibility}
\end{equation}
The additional condition reads the same links and large-field indicators.
\end{definition}
The names localized term, localization domain and regular part of the action are Ba\l{}aban's~\cite{Balaban1987RG,Balaban1988Expansions}. His localized terms are analytic, one for each connected localization domain, and bounded term by term; those used here are real Borel functions with finitely many derivatives, have a base point and an arbitrary localization domain, and are summed in a weighted $\ell^1$ norm. The word regular has two technical uses, as in his papers: the regular part of the action, and the regularity condition. Ba\l{}aban uses the word activity for the activities of a polymer expansion~\cite{Balaban1988Cluster}; its use for the large-field terms is particular to this book. Neither a localized term nor a large-field activity is assumed nonnegative. Both may be indexed by a
countable set; their summability will be part of the norm conditions.

There are two different meanings of support in this definition. The set
$Y_i$ determines exclusion and coverage: the $Y_i$ of each subfamily in the sum below will be pairwise disjoint.
But deciding whether a site belongs to $\operatorname{Bad}_m(U)$ reads
plaquettes meeting that site, whose links can extend beyond $Y_i$. The
\emph{dependence set} includes those links as well as the internal
links and any links read by other conditions. Two disjoint exclusion supports
can have overlapping dependence sets. A connected expansion based
on dependence must use the latter, whereas the sum over disjoint covering families below uses
the former. Disjoint exclusion supports do not imply probabilistic
independence.

\section{The prescribed leading action}
\label{ch04:leading-action}

The running coupling is a component of the representation relative to a prescribed leading
action. Its construction will come from Gaussian integration and local
matching; here we specify the form of that component without claiming
the estimates needed to select it.

At scale $m$, fix constants $c_{\rm dir},c_{\rm lin},R^{\rm rel}>0$ with
$16c_{\rm lin}<c_{\rm dir}$. Define
\begin{align}
 \operatorname{clip}_t(r)&=
 \begin{cases}
 c_{\rm dir}r^2/2,&r\le t,\\
 c_{\rm lin}tr,&r>t,
 \end{cases}
 \label{ch04:clip}\\
 \chi^{\rm rel}_{m,x}(U)&=
 \mathbf1\left\{\|U_{y;\mu\nu}-1\|_G
       \le t_m e^{\operatorname{dist}(x,y)/R^{\rm rel}}
              \text{ for every }y,\mu<\nu\right\}.
 \label{ch04:relaxed-test}
\end{align}
The distance is periodic and measured in lattice units. We call the condition in \eqref{ch04:relaxed-test} the \emph{relaxed regularity condition} at $x$: the regularity condition on the whole lattice, with $t_m$ replaced by the threshold $t_m e^{\operatorname{dist}(x,y)/R^{\rm rel}}$ growing exponentially with the distance from $x$. Although written
as a condition on every plaquette, \eqref{ch04:relaxed-test} reads only a finite
neighborhood: $\|U_{y;\mu\nu}-1\|_G\le2$ makes all conditions with
$t_m e^{\operatorname{dist}(x,y)/R^{\rm rel}}\ge2$ automatic. Any claim of
locality must include the plaquettes of the remaining conditions in its dependence
set.

For each orientation $\mathfrak p=(\mu,\nu)$, fix a summable linear
combination of transported curvatures,
\[
 \mathcal L^{\rm path}_{\mathfrak p}(U,x)
      =\sum_i r_{\mathfrak p i}\mathcal F_{\mathfrak p i}(U,x).
\]
Each $\mathcal F_{\mathfrak p i}$ is a plaquette curvature transported to
$x$, as in \eqref{ch02:transported-curvature}. Also specify local real
invariant multilinear expressions $A_i^{(q)}(U,x)$, $q=2,3,4$, in such
curvatures, each multiplied by the indicator of the regularity condition on its declared
support. The support of the expression contains the contours, plaquettes and base point; this
indicator may read additional adjacent plaquettes. Assume the coefficient sums
converge absolutely at the fixed finite volume. Set
\begin{align}
 P^{\rm raw}_{N,m}(\beta,U)
 &=\beta\sum_{x,\mathfrak p}
    \left[\operatorname{clip}_{t_m}(\|U_{x;\mathfrak p}-1\|_G)
       +\tfrac12\chi^{\rm rel}_{m,x}(U)
                    |\mathcal L^{\rm path}_{\mathfrak p}(U,x)|^2\right]
 \notag\\
 &\quad+\beta\sum_{x,i}(A_i^{(3)}+A_i^{(4)})(U,x)
                     +\sum_{x,i}A_i^{(2)}(U,x),\\
 P^{\rm path}_{N,m}(\beta,U)
 &=\frac1{|H_4|}\sum_{h\in H_4}P^{\rm raw}_{N,m}(\beta,h^*U).
 \label{ch04:path-action}
\end{align}
The coefficients are held fixed as $\beta$ varies. Thus this action is
affine in $\beta$. Its ultimate normalization requires the quadratic
coefficient in the $\beta$-dependent part to be the curvature square,
and no such coefficient in the $\beta$-independent quadratic term $\sum_{x,i}A_i^{(2)}(U,x)$. Those are
matching conditions, not consequences of the displayed form.

The function $\operatorname{clip}_t$ of \eqref{ch04:clip}, quadratic below $t$ and linear above, is not a smooth cutoff.
At $r=t$ it drops from $c_{\rm dir}t^2/2$ to $c_{\rm lin}t^2$ immediately
to the right. The terms of the representation of the effective density are therefore Borel functions,
not globally differentiable functions of the links. Derivative estimates
refer to the local expressions on the domains where the regularity condition holds, with the event indicators
held fixed. In particular, a derivative with respect to the inverse coupling $\beta$
does not differentiate the threshold or its indicator.

\section{Representations by disjoint covering families, and their norms}
\label{ch04:density-norms}

Recall that $Y_i$ is an exclusion support and $I_i^{\rm bad}$ its prescribed large-field
set, as in Definition~\ref{ch04:activities}. Let
$\mathfrak F_m(U)$ be the finite subsets $\mathfrak a$ of large-field
activity indices such that the $Y_i$, $i\in\mathfrak a$, are pairwise
disjoint and
$\bigcup_{i\in\mathfrak a}I_i^{\rm bad}=\operatorname{Bad}_m(U)$.
We call them the \emph{disjoint covering families} of the large-field sites (in combinatorics, exact covers). They do not select a
preferred partition into connected components.

\emph{Representation data} $\rho$ consist of $\beta$, a real constant
$v_{\rm sc}$ per site, the localized terms and the large-field
activities, with the fixed action \eqref{ch04:path-action}. Define
\begin{align}
 V_\rho(U)&=P^{\rm path}_{N,m}(\beta,U)
                +\sum_i h_{X_i}(U)E_i^{\rm reg}(U)-v_{\rm sc}|\Lambda|,
 \notag\\
 \operatorname{dens}(\rho)(U)
   &=e^{-V_\rho(U)}
          \sum_{\mathfrak a\in\mathfrak F_m(U)}
                         \prod_{i\in\mathfrak a}K_i^{\rm lf}(U).
 \label{ch04:effective-density}
\end{align}
The empty product is one. Because each $I_i^{\rm bad}$ is nonempty, if
there are no large-field sites the only disjoint covering family is empty and the sum
is one. Otherwise a single support may cover two regions, or two disjoint
supports may cover them separately; both alternatives belong in the sum.
Representation data satisfying the definitions alone need not give a positive density.
Positivity follows when a later construction proves that its density equals
an exact block integral of a positive input.

Fix an integer derivative order $q\ge2$, a nonnegative spatial exponent
$\varpi$, a support exponent $\sigma>0$ (the exponent of the support weight), and the threshold $t_m$.
For a localized term with base point $o_i$, use the seminorm
\begin{equation}
 \|E_i^{\rm reg}\|_m
 =\sup_{0\le k\le q}\ \sup_{U,A_1,\ldots,A_k}
        |D_{A_1}^L\cdots D_{A_k}^L E_i^{\rm reg}(U)|,
 \label{ch04:local-seminorm}
\end{equation}
where $U$ ranges over configurations that satisfy the regularity condition on $X_i$.
Each $A_r$ is a Lie algebra valued link field satisfying, on each link $e$ with initial site $e^-$,
$|(A_r)_e|_{\gLie}\le t_m(1+\operatorname{dist}(e^-,o_i))^\varpi$.
The left derivatives are those of Section~\ref{ch02:bounds}. The required
differentiability includes existence and uniform bounds at the boundary of the
domain defined by the regularity condition, interpreted through the local expression $E_i^{\rm reg}$ itself;
no derivative of $h_{X_i}$ is included.

For a set $X\subset\Lambda$, let the \emph{connected covering number} $\operatorname{cov}_{B_m}(X)$ be the
minimum number of radius-$B_m$ periodic balls whose union is connected
in the nearest-neighbor graph and covers $X$. Set $\operatorname{cov}_{B_m}(\varnothing)=0$. The covering union may
contain additional sites. Its exponential weight penalizes extended
supports without requiring a support itself to be connected; the covering number plays the role of the linear size of a localization domain in~\cite{Balaban1987RG}.
Choose majorants $e_i\ge\|E_i^{\rm reg}\|_m$ and
$k_i\ge\sup_U|K_i^{\rm lf}(U)|$, and define
\begin{equation}
 \begin{split}
 \mathcal E_m(\rho)&=\sup_{x\in\Lambda}
     \sum_{i:o_i=x}e^{\sigma\operatorname{cov}_{B_m}(X_i)}e_i,\\
 \mathcal K_m(\rho)&=\sup_{x\in\Lambda}
     \sum_{i:x\in Y_i}e^{\sigma\operatorname{cov}_{B_m}(Y_i)}k_i.
 \end{split}
 \label{ch04:weighted-sums}
\end{equation}
The first is a sum over the terms with base point $x$, an \emph{anchored} sum: each localized term is counted at one
site. The second counts a large-field activity at every site of its
exclusion support. Replacing an anchored sum by a sum over the terms whose localization domain contains $x$, a \emph{containing-site} sum like the second,
requires a bound for the number of possible base points; it is not a harmless
notational change. Likewise, replacing $Y_i$ by its dependence set
changes the weight and requires a geometric estimate.

Finite weighted sums ensure that \eqref{ch04:effective-density} is
well-defined at fixed volume. Indeed
\[
 \sum_i\|E_i^{\rm reg}\|_m\le|\Lambda|\mathcal E_m,\qquad
 \sum_i\sup|K_i^{\rm lf}|\le|\Lambda|\mathcal K_m,
\]
and the absolute sum over all finite subfamilies is at most
$\prod_i(1+k_i)\le\exp(|\Lambda|\mathcal K_m)$.
Restricting to disjoint covering families only lowers this bound. The estimate
is a finite-volume convergence statement, not a uniform bound on the
entire partition function.

\section{Conditions for a one-step induction}
\label{ch04:admissibility}

The preceding representation data describe possible densities; their smallness
conditions will describe the inputs of a one-step theorem. There are
additional normalization conditions. On a nonwrapping localization domain\mgubi{Issue 10: ``nonwrapping'' is used in several chapters and never defined. In the sources: $10\,\mathrm{dist}(x,o)<n$ for all $x$ in the support, $n$ the period; Chapter 17 uses the factor $10$.}, choose
centered coordinate lifts $\xi_{o_i}(y)$. For every pair of affine link
fields of the form
\[
 A_{(y,\mu)}=u_\mu+\sum_\nu\xi_{o_i}(y)^\nu v_{\mu\nu},
 \qquad u_\mu,v_{\mu\nu}\in\gLie,
\]
require the Hessian at zero of $A\mapsto E_i^{\rm reg}(\exp A)$ to vanish
on that pair. Together with $E_i^{\rm reg}(1)=0$ and the vanishing linear
term, this excludes the low-order local contributions assigned to the
leading action. The first derivative vanishes for a gauge-invariant localized
term by invariance under constant gauge transformations: each
linear coefficient would otherwise be an invariant linear functional on
the simple Lie algebra. The meaning of the affine Hessian condition is
fixed in the chosen local coordinates; no coordinate lift is imposed on
a wrapping localization domain by this definition.

One may prescribe the intended size bounds
\begin{equation}
 \mathcal E_m\le\widehat\beta_m^{-11/10},\qquad
 \mathcal K_m\le\exp\bigl(-(M_m^{\rm cut})^{1/5}\bigr),
 \label{ch04:activity-bounds}
\end{equation}
with the reference quantities of \eqref{ch04:thresholds}. In an induction,
$q,\varpi,\sigma$ may depend on the \emph{stage} $J_m$ in force at scale $m$ (a stage specifies finitely many derivative orders, support exponents, path widths and coefficient bounds; see Section~\ref{ch10:choices}); the stage index
is different from the scale index $m$. For now they remain specified input
parameters. Neither \eqref{ch04:activity-bounds} nor the density identity
proves that the block map preserves the class.

Parameter continuity is also stated for the representation data. A curve
$u\mapsto\rho_u$ has fixed localization domains, exclusion supports, base points, prescribed large-field sets,
additional conditions and event indicators. Its coupling and its constant $v_{\rm sc}$ are continuous. At each
$u_0$, differences of the derivatives of the localized terms and of the values of the large-field activities have
majorants $e_i(u),k_i(u)$ satisfying
\[
 \sum_i e^{\sigma\operatorname{cov}_{B_m}(X_i)}e_i(u)
 +\sum_i e^{\sigma\operatorname{cov}_{B_m}(Y_i)}k_i(u)
                  \longrightarrow0\quad(u\to u_0).
\]
Both sums here are over all indices, at a fixed finite volume. This
continuity condition is strong enough to control the density and its
integral; it differentiates no moving large-field boundary.

The analytic task remaining is now precise. Starting with representation data at the
fine scale, one must construct coarse representation data whose density is exactly
$\mathcal B^{(L_m)}_\beta\operatorname{dens}(\rho)$, whose coupling has the
claimed adjustment, and whose localized terms and large-field activities obey the coarse conditions.
A class defined as the outputs reachable from a Wilson starting density
can only be introduced after that construction. It must not be confused
with the much larger set of representation data satisfying the displayed
norm bounds.
}{}
\setcounter{chapter}{4}
\IfFileExists{chapters/05-gaussian.tex}{\chapter{Gaussian kernels and stationary backgrounds}
\label{ch05:gaussian}

The quadratic approximation serves two purposes. It identifies the
covariance of the fluctuations eliminated by a block transformation, and
it supplies the inverse needed to construct a stationary background for
the nonlinear action. These uses require different estimates. Exact matrix
identities hold on each finite lattice; uniform locality requires bounds
which survive changes of the lattice size and of the blocking history.

We first fix the normalization of the exact identities. We then isolate the
additional hypotheses for uniform analytic continuation and for a
stationary background. The general block transformation algebra is
standard; see Ba{\l}aban--Feldman--Kn\"orrer--Trubowitz
\cite{BalabanFeldmanKnoerrerTrubowitz2016}. The construction of local
propagators and backgrounds for lattice gauge theory originates in
Ba{\l}aban's work \cite{Balaban1984Propagators,Balaban1985Variational}.
Those references explain the method. Their results cannot be applied to a
different averaging map without checking its hypotheses.

\section{Lattice differential forms and zero modes}
\label{ch05:forms}

In this chapter distances and Fourier variables are in lattice units.
Let $\Lambda_s=(\ZZ/s\ZZ)^4$. A lattice $r$-form assigns an element of
$\gLie$ to each site and each increasing $r$-tuple of coordinate
directions. The scalar product is the sum over sites and increasing
tuples, using the Lie algebra metric of \eqref{ch01:metric}.
In particular there is no volume normalization in this scalar product.
Write $\Omega^r(\Lambda_s;\gLie)$ for this finite-dimensional space.
The forward difference is
$\nabla_\mu u(x)=u(x+e_\mu)-u(x)$, and the exterior derivative is
\[
 (d_r u)_{\mu_0\cdots\mu_r}
 =\sum_{i=0}^r(-1)^i\nabla_{\mu_i}
       u_{\mu_0\cdots\widehat{\mu_i}\cdots\mu_r}.
\]
Commutation of the differences gives $d_{r+1}d_r=0$.
An adjoint is always taken for the stated counting scalar products.

A linearized gauge transformation of a one-form $a$ is
$a\mapsto a+d_0\phi$. If $H$ is a positive self-adjoint operator on
two-forms, the quadratic action and its kinetic operator are
\begin{equation}
 S_0(a)=\tfrac12\langle d_1a,Hd_1a\rangle,\qquad K=d_1^*Hd_1.
 \label{ch05:kinetic}
\end{equation}
Thus $K$ vanishes on gradients. It also vanishes on constant one-forms.
On a periodic lattice the latter are harmonic modes and are not gradients
of periodic scalar functions. Removing gauge directions alone therefore
does not make $K$ invertible.

For $p\in(2\pi/s)\ZZ^4/(2\pi\ZZ)^4$ our Fourier convention is
\begin{equation}
 \widehat u(p)=\sum_{x\in\Lambda_s}e^{-ip\cdot x}u(x),
 \quad
 u(x)=s^{-4}\sum_p e^{ip\cdot x}\widehat u(p),
 \quad
 \|u\|_2^2=s^{-4}\sum_p|\widehat u(p)|^2.
 \label{ch05:fourier}
\end{equation}
Set $q_\mu(p)=e^{ip_\mu}-1$ and
$\lambda(p)=\sum_\mu q_\mu(-p)q_\mu(p)$.
On real momenta this is $\sum_\mu|q_\mu(p)|^2$.
The symbol of $d_r$ is exterior multiplication by $q$.
In particular, on one-forms,
\[
 |q\wedge a|^2+|\overline q\mathbin{\lrcorner}a|^2
      =\lambda(p)|a|^2.
\]
Here contraction denotes the adjoint of exterior multiplication.
For $p\ne0$ this gives the orthogonal decomposition into a gradient
and its orthogonal complement. At $p=0$ the whole one-form space is
harmonic.

Whenever a Gaussian probability measure is used below, its precision is
strictly positive on the specified coordinate space. One may work on a
fixed slice with the harmonic modes removed, or add a positive regulator.
A massless covariance at zero momentum is not defined by writing
$1/\lambda(0)$. Conversely, an identity between quadratic forms can
remain valid on a space where no normalized massless Gaussian exists.

\section{The exact Gaussian block transformation}
\label{ch05:schur}

Let $E_f,E_c$ be real finite-dimensional inner-product spaces, $Q:E_f\to
E_c$ a linear block map, and $b>0$ the block penalty. Suppose initially
that $K_f>0$ on $E_f$. Define
\begin{equation}
 A=K_f+bQ^*Q,\qquad
 W=A^{-1}bQ^*,\qquad
 K_c=bI-b^2QA^{-1}Q^* .
 \label{ch05:schur-data}
\end{equation}
These formulas describe the normalized Gaussian problem. When the
physical quadratic exponent is
$\beta(\langle A_f,K_fA_f\rangle+b\|QA_f-V\|^2)/2$,
the variables here are $a=\sqrt{\beta}A_f$ and
$v=\sqrt{\beta}V$. In unscaled link logarithms the fluctuation
covariance is $\beta^{-1}(K_f+bQ^*Q)^{-1}$.
Thus the block parameter remains $b$ after normalization, although
the original block kernel has penalty $b\beta$.

The map $W$ sends a coarse field to the stationary fine field. The
fluctuation covariance is $A^{-1}$. The Schur-complement formulas give
\begin{align}
 K_c&=(b^{-1}I+QK_f^{-1}Q^*)^{-1}, \label{ch05:coarse-cov}\\
 \langle a,K_fa\rangle+b\|Qa-v\|^2
 &=\langle v,K_cv\rangle+\langle a-Wv,A(a-Wv)\rangle .
 \label{ch05:completion}
\end{align}
These are the usual finite-dimensional Gaussian identities, or the
real quadratic specialization of the block transformation identities in
\cite{BalabanFeldmanKnoerrerTrubowitz2016}; no separate proof is needed.
The accompanying normalization is
\[
 \int_{E_f}e^{-\langle a,K_fa\rangle/2-b\|Qa-v\|^2/2}\,da
 =(2\pi)^{\dim E_f/2}(\det A)^{-1/2}
            e^{-\langle v,K_cv\rangle/2}.
\]
The Lebesgue measure here comes from the specified inner product.
In a normalized block kernel there is also the factor
$(b/2\pi)^{\dim E_c/2}$.

There is an equivalent probabilistic statement. If $a$ and $\xi$ are
independent centered Gaussians with covariances $K_f^{-1}$ and $I$,
then $v=Qa+b^{-1/2}\xi$ has covariance
$b^{-1}I+QK_f^{-1}Q^*$. Conditional on $v$, the fine field has mean
$Wv$ and covariance $A^{-1}$.

The identities \eqref{ch05:schur-data} and
\eqref{ch05:completion} also hold for $K_f\ge0$, provided $A$ is
invertible. Their derivation uses no inverse of $K_f$.
In this case \eqref{ch05:coarse-cov} and the joint probability
interpretation are unavailable unless the remaining null directions have
been treated. In the gauge problem, every use of this extension must
specify the slice or quotient on which $A$ is invertible.

The useful consequences, with the same domains and adjoints, are
\[
 K_fW=Q^*K_c,\qquad QW=I-b^{-1}K_c,\qquad
 W^*K_fW+b(QW-I)^*(QW-I)=K_c.
\]
The last identity says that the stationary fine action together with
its block penalty is exactly the coarse quadratic action.

\section{Averaging and the Fourier alias formula}
\label{ch05:alias}

To check the scaling factor in the covariance recursion, take
$\Lambda_{Ls}$ as fine lattice and $\Lambda_s$ as coarse lattice, with
odd $L\ge3$. Let $w(z)$ be a finitely supported real averaging kernel
of total mass one, periodized on the fine lattice, and use
\[
 (Ma)_\mu(y)=\sum_z w(z)\sum_{r=0}^{L-1}
                    a_\mu(Ly+z+re_\mu).
\]
The inner sum is a sum along a link of length $L$, not its average.
Define the symbols
\[
 \widehat w(\theta)=\sum_z w(z)e^{i\theta\cdot z},
 \qquad r_L(t)=L^{-1}\sum_{r=0}^{L-1}e^{irt}.
\]
The sign in $\widehat w$ is chosen for this formula for $M$.
The aliases of a coarse momentum $p$ are
$\theta_z=(p+2\pi z)/L$, with one representative for each
$z\in(\ZZ/L\ZZ)^4$.

Suppose an auxiliary, translation-invariant Gaussian one-form has
diagonal covariance symbol $\operatorname{diag}(c_\mu(\theta))$ in
the direction index. Color indices are suppressed and are unchanged
by this calculation. The symbol of the covariance of
$Ma+b^{-1/2}\xi$ is
\begin{equation}
 c'_\mu(p)=b^{-1}
 +L^{-2}\sum_{z\ {\rm mod}\ L}
 \widehat w(\theta_z)\widehat w(-\theta_z)
 r_L(\theta_{z,\mu})r_L(-\theta_{z,\mu})c_\mu(\theta_z).
 \label{ch05:alias-recursion}
\end{equation}
This formula is exact for a positive regulated covariance. It can
also be used for nonzero modes of a massless symbol with the zero
mode handled separately.

For clarity, the normalization check is short. Substitution of
\eqref{ch05:fourier} into $Ma$ gives
$\widehat{Ma}_\mu(p)=L^{-4}\sum_z
L\widehat w(\theta_z)r_L(\theta_{z,\mu})
\widehat a_\mu(\theta_z)$.
Fine Fourier covariances carry the volume $(Ls)^4$, whereas coarse
ones carry $s^4$. Consequently the covariance factor is
$L^{-8}L^2L^4=L^{-2}$. This is the factor used in
\eqref{ch05:alias-recursion}.

The particular smoothing in the supplied construction uses
\[
 \widehat w_L(\theta)=
 \prod_{\nu=1}^4\left(L^{-1}
       \sum_{|r|\le(L-1)/2}e^{ir\theta_\nu}\right),
 \qquad \widehat w(\theta)=\widehat w_L(\theta)^j,
\]
where $j$ is the integer smoothing order. Centering makes
$\widehat w_L$ even on real and complex momenta.
The recurrence then contains $\widehat w_L(\theta_z)^{2j}$.
This choice fixes an auxiliary covariance. Identification with the
quadratic form of a particular gauge-equivariant block map additionally
requires checking that its linearization differs from $M$ only by an
appropriate gradient, and that the chosen treatment of zero modes agrees.

\section{Gauge reduction and a regular curvature symbol}
\label{ch05:curvature}

At a real nonzero momentum suppose $c_\mu>0$ and put
$w_\mu=c_\mu^{-1}$. Minimizing over a scalar gauge parameter gives
\begin{equation}
 \min_\phi\sum_\mu w_\mu|a_\mu+q_\mu\phi|^2
 =\frac{\displaystyle\sum_{\mu<\nu}w_\mu w_\nu
                  |q_\mu a_\nu-q_\nu a_\mu|^2}
        {\displaystyle\sum_\mu w_\mu|q_\mu|^2}.
 \label{ch05:quotient}
\end{equation}
It is the weighted orthogonal projection formula. Thus a diagonal
two-form metric representing the quotient form is
\begin{equation}
 (H_0)_{\mu\nu}(p)=
 \left[c_\mu(p)c_\nu(p)
               \sum_\alpha\frac{q_\alpha(-p)q_\alpha(p)}
                                  {c_\alpha(p)}\right]^{-1}.
 \label{ch05:H0}
\end{equation}
Formula \eqref{ch05:H0} is initially an identity away from zero.
Its regular extension at zero is a separate issue.

Here is the cancellation used for that extension. Let $M$ now denote
the total dilation over several steps; it is a number, distinct from
the averaging operator of the previous section.
Put $\lambda_M(p)=M^2\lambda(p/M)$.
Suppose the accumulated covariance has the representation
\[
 c_\mu(p)=\frac{A_\mu(p)}{\lambda_M(p)}+G_\mu(p),
 \quad
 A_\mu(p)=F(p)r_M(p_\mu/M)r_M(-p_\mu/M),
 \quad B_\mu=A_\mu+\lambda_MG_\mu.
\]
The remainder $G_\mu$ collects the nonprincipal aliases and the
accumulated block noise. Assume that $F,G_\mu,A_\mu,B_\mu$ are
holomorphic on a common neighborhood and that $F,A_\mu,B_\mu$ do
not vanish there. The elementary identity
$q_\mu(p)r_M(p_\mu/M)^{-1}=M q_\mu(p/M)$, interpreted after
multiplication to remove possible zeros, gives
\[
 q_\mu(p)q_\mu(-p)F
       =M^2q_\mu(p/M)q_\mu(-p/M)A_\mu .
\]
Consequently, if
\begin{equation}
 h=F^{-1}-\sum_\mu
       \frac{q_\mu(p)q_\mu(-p)G_\mu}{A_\mu B_\mu},
 \qquad
 \lambda_M h=\sum_\mu\frac{q_\mu(p)q_\mu(-p)}{B_\mu},
 \label{ch05:pole-cancel}
\end{equation}
then \eqref{ch05:H0} becomes
\[
                 (H_0)_{\mu\nu}=(h B_\mu B_\nu)^{-1}.
\]
The second equality in \eqref{ch05:pole-cancel} follows by subtracting
$1/B_\mu$ from $1/A_\mu$; it is not an additional assumption.
If $h$ is also nonvanishing, the last formula proves holomorphic
extension through the apparent massless pole.
When $F(0)=A_\mu(0)=1$, it gives $H_0(0)=I$.

This calculation proves pole cancellation, not a uniform zero-free
strip. For that stronger conclusion the nonprincipal aliases,
accumulated noise, and reciprocal factors must be bounded uniformly
over every admitted blocking history. These are precisely the
estimates which cannot be replaced by finite-dimensional Gaussian
algebra.

One can modify the metric away from exact two-forms without changing
the action. For $\zeta\ge0$ define, at real momentum,
\begin{equation}
 H^{-1}=H_0^{-1}
       +\zeta(d_2H_0^{-1})^*(d_2H_0^{-1}).
 \label{ch05:metric-modification}
\end{equation}
If $H_0>0$, this defines a positive $H$, and
$H^{-1}H_0d_1=d_1$, because $d_2d_1=0$. Therefore
$Hd_1=H_0d_1$. Positivity on all two-forms and agreement on exact
two-forms are different assertions; this modification is useful
because later estimates need both.

\section{Uniform estimates over finite blocking histories}
\label{ch05:histories}

We now verify the uniform reciprocal bounds required in
\eqref{ch05:pole-cancel} for the actual averaging recurrence.
A finite history consists of odd ratios $L_1,\ldots,L_n\ge3$ and
smoothing orders $j_0,\ldots,j_n\ge4$, with
$j_{h-1}-j_h\in\{0,1\}$. Set $M_h=L_1\cdots L_h$, $M_0=1$.
The order used by the $h$th block map is its incoming order
$j_{h-1}$. Denote the centered box symbol $\widehat w_L$ of
Section~\ref{ch05:alias} by $w_L$ in the following formulas.

\begin{proposition}[Uniform curvature metric]
\label{ch05:history-prop}
Fix the final smoothing order $j=j_n\ge4$ and the penalty $b>0$.
Start with the nonzero-mode covariance $c_\mu(p)=1/\lambda(p)$
and iterate \eqref{ch05:alias-recursion} along any such finite
history. The resulting curvature metric $H_0$ has a holomorphic,
periodic extension to a common strip about real momenta, with
uniformly bounded entries and inverse entries. On real momenta
$H_0$ has uniform positive upper and lower bounds and $H_0(0)=I$.
The constants depend on $j,b$, but not on the history length,
ratios or torus period.
\end{proposition}

\begin{proof}
The box composition identity
$w_A(B\theta)w_B(\theta)=w_{AB}(\theta)$ holds for odd
$A,B$, by decomposing a centered interval of length $AB$
into centered intervals of length $B$.
Thus the accumulated filter is
\begin{equation}
 f(\theta)=w_{M_n}(\theta)^j
       \prod_{\{h:j_{h-1}=j_h+1\}}w_{M_h}(\theta).
 \label{ch05:history-filter}
\end{equation}
The directional filter is
$f(\theta)r_{M_n}(\theta_\mu)$: the link factors compose by
$r_A(Bt)r_B(t)=r_{AB}(t)$.

Here and below momentum estimates are on an enlarged cube
$|\operatorname{Re}p_\mu|<\pi+\eta$,
$|\operatorname{Im}p_\mu|<\eta$, for a sufficiently small fixed
$\eta>0$. Write $M=M_n$ and choose centered alias representatives
$z_\mu\in[-(M-1)/2,(M-1)/2]$. The finite geometric sum gives
\[
 |r_M((p_\mu+2\pi z_\mu)/M)|,\quad
 \left|M^{-1}\sum_{|r|\le(M-1)/2}
     e^{ir(p_\mu+2\pi z_\mu)/M}\right|
       \le \frac{C}{1+|z_\mu|}.
\]
Indeed the numerator of the geometric quotient is bounded, and
$M|e^{i(p_\mu+2\pi z_\mu)/M}-1|$ is bounded below by
$c(1+|z_\mu|)$ for $z_\mu\ne0$; the principal case is bounded
directly by the finite sum. At the principal real alias the
centered factor has a positive lower bound.

The additional factors in \eqref{ch05:history-filter} have a
uniform product bound. A factor of side $A\le M$ has modulus
at most $\exp(C\eta A/M)$ on the cube evaluated at $p/M$
or any other alias. The sides occurring in the product increase
at least geometrically, so their sum divided by $M$ is bounded.
On the real principal alias the same factors have product bounded
away from zero. To see this without multiplying a fixed lower
bound infinitely often, use
$|\log w_A(p/M)|\le C A/M$ on that real cube, with the value
at a removable zero interpreted continuously. The bound follows
from the geometric quotient near zero and compactness for
$A/M$ bounded below. Again $\sum A/M$ is bounded.
Consequently, with
$\Gamma(z)=\prod_\mu(1+|z_\mu|)$,
\[
 |f((p+2\pi z)/M)r_M((p_\mu+2\pi z_\mu)/M)|
       \le C_j\Gamma(z)^{-j},
 \qquad |f(p/M)|\ge c_j>0
\]
in the first estimate on a complex cube and in the second on its
real slice.

There is also an exact description of the accumulated noise.
The independent noise introduced at step $h$ is propagated only
through the suffix $h+1,\ldots,n$. If that suffix has total ratio
$D=M_n/M_h$, its contribution, before multiplication by $b^{-1}$,
is
\[
 D^{-2}\sum_{z\ {\rm mod}\ D}
   f_{\rm suffix}(\theta_z)f_{\rm suffix}(-\theta_z)
   r_D(\theta_{z,\mu})r_D(-\theta_{z,\mu}).
\]
This includes $1$ when $h=n$ and $D=1$.
To verify the formula, apply \eqref{ch05:alias-recursion} to one
constant covariance and split each residue modulo a product of
ratios into its successive residues. The filter composition
identities above identify the summands and the factors multiply
to $D^{-2}$. Summing over the noises gives precisely the noise
part of the iterated recurrence.
Each suffix contribution is bounded by $C_jD^{-2}$ because
\[
 \sum_z\Gamma(z)^{-2j}\le
 \left(\sum_{k\in\ZZ}(1+|k|)^{-2}\right)^4<81.
\]
As $D\ge3^{n-h}$, the total noise is at most
$C_j\sum_{r\ge0}3^{-2r}$, uniformly in $n$.

Separate the principal alias of the propagated initial
covariance from all other aliases. The former gives
$A_\mu/\lambda_M$ with the notation of
\eqref{ch05:pole-cancel}; the latter are
\[
 \sum_{z\ne0}
 \frac{f(\theta_z)f(-\theta_z)
                 r_M(\theta_{z,\mu})r_M(-\theta_{z,\mu})}
                      {M^2\lambda(\theta_z)}.
\]
For a nonprincipal real alias the denominator is bounded below
by $c(1+\max_\mu|z_\mu|)^2$. Its change on a complex strip of
width $\eta$ is at most
$C\eta(1+\max_\mu|z_\mu|)+C\eta^2$.
Decreasing $\eta$ gives the same lower bound in modulus.
The preceding summability therefore bounds this sum holomorphically.
Together with the suffix noise estimate it proves a uniform
holomorphic bound for $G_\mu$, which is nonnegative on the real
slice.

Recall that $F=f(p/M)f(-p/M)$,
$A_\mu=F r_M(p_\mu/M)r_M(-p_\mu/M)$ and
$B_\mu=A_\mu+\lambda_MG_\mu$.
The principal filter bounds give positive real lower bounds for
$F,A_\mu$, and hence for $B_\mu$. All three have uniform complex
upper bounds. Cauchy estimates on a slightly larger cube therefore
give uniformly bounded reciprocals on a smaller one.
The identity in \eqref{ch05:pole-cancel} now supplies the real
lower bound for $h$: since $B_\mu\le C_B$,
\[
 h(p)=\frac{\sum_\mu |q_\mu(p)|^2/B_\mu(p)}{\lambda_M(p)}
          \ge c/C_B\quad(p\ne0).
\]
We used uniform comparison of both $\lambda(p)$ and
$\lambda_M(p)$ with $|p|^2$ on the real cube. At zero $h=1$.
Its defining formula gives a holomorphic upper bound using the
reciprocals already obtained, and another Cauchy estimate gives
a uniformly bounded reciprocal for $h$.
Thus $(H_0)_{\mu\nu}=(hB_\mu B_\nu)^{-1}$ has all the asserted
bounds. The order of this argument is important: the lower bound
for $h$ uses $B_\mu$, not an assumed inverse bound for $h$.
Periodicity holds first on real overlaps by permutation of the
aliases in the original covariance recurrence. Analytic continuation
then joins the cubes to a strip. The same Fourier symbol works on
every finite-period grid.
\end{proof}

This proves uniform bounds for arbitrary finite histories.
Convergence for a prescribed sequence of growing histories is
stronger and is not a consequence of this proposition alone.

\section{Exact transports and the curvature defect}
\label{ch05:transports}

Recall that $H_0$ is the two-form metric in \eqref{ch05:H0}; $H$ denotes its positive completion from \eqref{ch05:metric-modification}, agreeing with $H_0$ on exact two-forms. Fine and coarse subscripts below specify the scale of each metric.

We describe the algebra at one real nonzero coarse momentum.
This also specifies which transport must be estimated. The use
of aliases and block-periodic Fourier coordinates is standard;
see \cite{BalabanFeldmanKnoerrerTrubowitzBloch2016}.
Write $\theta_l=(p+2\pi l)/L$,
$\delta_r=Ld_r(\theta_l)$ for the rescaled fine exterior derivative,
and $d_r=d_r(p)$ for the coarse one.
Let
\[
 D_l=L^{-2}\operatorname{diag}(c_{f,\mu}(\theta_l)),
 \quad D=\operatorname{diag}(c_{c,\mu}(p)),\quad
 x=d_0,\quad x_l=\delta_0 .
\]
For this paragraph $H,H_l$ denote the unmodified curvature metrics.
Set
\[
 R=d_1^*Hd_1,\quad R_l=\delta_1^*H_l\delta_1,\quad
 \rho=x^*D^{-1}x,\quad \pi=\rho^{-1}x^*D^{-1}.
\]
The gauge quotient identity gives
\[
 R=D^{-1}-\rho^{-1}D^{-1}xx^*D^{-1},\qquad
 DR=I-x\pi .
\]
The fine formulas have subscript $l$. Define
\[
 m_l=w_L(\theta_l)^j,\qquad
 F_{r,l}=m_l\bigwedge\nolimits^r
            \operatorname{diag}(r_L(\theta_{l,\mu})).
\]
The link identity and the covariance recurrence say, respectively,
\begin{equation}
 F_{r+1,l}\delta_r=d_rF_{r,l},\qquad
             \sum_lF_{1,l}D_lF_{1,l}^*=D-b^{-1}I .
 \label{ch05:transport-input}
\end{equation}
Introduce the positive scalar
$\alpha=\sum_l|m_l|^2\tr D_l$, and put
$e_l=\overline m_l\tr D_l/\alpha$,
$y_l=|m_l|^2\rho/\rho_l$, and $C=Hd_1Dd_1^*$.
The rescaled one-form and two-form transports are
\begin{align}
 \mathcal A_l&=D_lF_{1,l}^*R+e_lx_l\pi,\\
 \mathcal T_l&=H_l\delta_1D_lF_{1,l}^*d_1^*
                         +y_lF_{2,l}^*(I-C).
 \label{ch05:transport-formula}
\end{align}
Their defining identities are
\begin{align}
 R_l\mathcal A_l&=F_{1,l}^*R,&
 \delta_1^*\mathcal T_l&=F_{1,l}^*d_1^*,&
 \mathcal T_lHd_1&=H_l\delta_1\mathcal A_l,\\
 \sum_lF_{1,l}\mathcal A_l&=I-b^{-1}R.
 \label{ch05:transport-identities}
\end{align}
These follow by substituting the quotient formula for $R$ and
using \eqref{ch05:transport-input}; in the last equality the
gradient contribution is $x\pi\sum_lm_le_l=x\pi$.
In particular the term $I-C$ cannot simply be omitted: it is
annihilated by $d_1^*$ and annihilates $Hd_1$, but controls the
transport on the complementary two-forms.

The physical kernels from a coarse site to a fine site $Ly+r$
are obtained from the alias sums
\[
 L^{-1}\sum_l e^{i\theta_l\cdot r}\mathcal A_l(p),
 \qquad
 L^{-2}\sum_l e^{i\theta_l\cdot r}\mathcal T_l(p)
\]
by coarse Fourier inversion. Denote the resulting operators by
$W$ and $T$. Restoring counting adjoints gives
\begin{equation}
 T^*d_{1,f}=d_{1,c}Q,\qquad
 TH_cd_{1,c}=H_fd_{1,f}W,\qquad
 QW=I-b^{-1}d_{1,c}^*H_cd_{1,c}.
 \label{ch05:physical-transports}
\end{equation}
Here $Q$ is the straight averaging operator in this construction.
The same statements for a block average containing a gradient
correction require its corresponding choice of the gradient
part of $W$.

The apparent poles in \eqref{ch05:transport-formula} can be
removed using the principal covariance split. We record the
crucial cancellation to avoid concealing a uniformity assumption.
Write $s=\lambda_{M_c}(p)$, $D=s^{-1}\operatorname{diag}(B_\mu)$,
$D_0=s^{-1}\operatorname{diag}(B_{0,\mu})$, and
\[
 D=F_{1,0}D_0F_{1,0}^*+U,\qquad U=\operatorname{diag}(u_\mu).
\]
The tail $U$ consists of nonprincipal aliases and the current
noise; it is uniformly holomorphic by the estimates just proved.
Let $z_\mu=u_\mu/B_\mu$. Then
\[
 \frac{|m_0r_L(\theta_{0,\mu})|^2B_{0,\mu}}{B_\mu}
                      =1-sz_\mu .
\]
With $J=F_{2,0}^{-1}$, define the bounded analytic matrix
\[
 \mathcal B=H_0\left[
       \operatorname{diag}_{\mu<\nu}(z_\mu+z_\nu)
                    -s\bigwedge\nolimits^2\operatorname{diag}(z_\mu)
                         \right].
\]
Here $H_0$ means the fine metric at alias zero, not a new
modification of the metric. Entrywise multiplication gives
$s\mathcal B=H_0-y_0F_{2,0}^*HF_{2,0}$, and hence
\[
 \mathcal T_0=y_0F_{2,0}^*
   +\bigl(\mathcal B Jd_1\operatorname{diag}(B_\mu)
                                -H_0Jd_1U\bigr)d_1^*.
\]
All factors are regular at zero. For $\mathcal A_0$, the
corresponding principal denominator is
$s\alpha=|m_0|^2\sum_\mu B_{0,\mu}+s\alpha_{\ne0}$,
which is bounded below on the real cube. Its reciprocal is
therefore uniformly holomorphic. Subtracting the leading
diagonal term in $\mathcal A_0$ leaves the regular factor
\[
 m_0^{-1}\left(u_\mu/B_\mu-\alpha_{\ne0}/(s\alpha)\right).
\]
For nonprincipal aliases the denominators have the lower bounds
used in Proposition~\ref{ch05:history-prop}.
These formulas give bounded principal columns and bounds
$C\Gamma(l)^{-j}$ for both nonprincipal columns on a smaller cube.
Adjoints in these analytic formulas mean transpose at opposite
momentum, with $l$ replaced by $-l$.

Each fine difference in the alias sum contributes
$e^{i\theta_{l,\mu}}-1$, bounded by
$CL^{-1}(1+\max_\nu|l_\nu|)$.
The summability estimate for $\Gamma$ therefore gives, for
$r\le j-2$, exponentially weighted coarse Fourier coefficient
bounds of order $L^{-1-r}$ for $\nabla_f^rW$ and $L^{-2-r}$
for $\nabla_f^rT$. Fourier contour displacement is the usual
conversion of strip analyticity to exponential coefficient
decay. It is performed after summing the aliases, whose
permutation under $p\mapsto p+2\pi n$ gives periodicity.
When a fine difference crosses a block boundary, its coarse
argument is shifted by one block; this is included by the phase
$e^{i\theta_l\cdot r}$ and must not be dropped.
At zero momentum the principal columns are $I$ and the other
columns vanish, so the identities extend there without choosing
a massless Gaussian normalization.

\subsection{A positive factorization of the defect}

The next argument uses exterior-complex exactness together with
the bounds just obtained. Exactness for a regular sequence is
standard Koszul algebra \cite[Lemma 15.31.2]{StacksKoszul}.
The quantitative version needed here preserves exponential
coefficient bounds: if an exponentially summable matrix row
$A$ on positive-degree forms satisfies $Ad_r=0$, it can be
written $A=Vd_{r+1}$ on a smaller strip, with a fixed norm cost.
One obtains this version by successively dividing
$A(p)-A(p)|_{p_a=0}$ by $e^{ip_a}-1$ and contracting the
$a$th exterior index. The elementary quotient
\[
 \frac{e^{inp_a}-1}{e^{ip_a}-1}
   =\sum_{h=0}^{n-1}e^{ihp_a}\quad(n>0)
\]
and the corresponding negative sum for $n<0$ show that a
division costs at most $C/\eta$ when width $\eta$ is reduced
to $\eta/2$. There are only four coordinates. This is the
usual coordinate proof of Koszul exactness, with the coefficient
norm retained. Applying it on both sides shows that a symmetric
two-form row $E$ with $Ed_1=d_1^*E=0$ has
$E=d_2^*C_0d_2$, with $C_0$ symmetric and exponentially summable.
For completeness, first choose a left lift $d_2^*W_0=E$.
Since $d_2^*W_0d_1=0$, exactness gives
$W_0d_1=d_3^*V_0$. The top differential $d_3^*$ has trivial
kernel on analytic four-form coefficients, so $V_0d_0=0$.
A right lift gives $V_0=X_0d_1$. Thus
$W_1=W_0-d_3^*X_0$ satisfies both $W_1d_1=0$ and
$d_2^*W_1=E$. Lift once more to write $W_1=C_0d_2$,
and replace $C_0$ by $(C_0+C_0^*)/2$.
All these lifts act on the infinite lattice coefficients before
periodization; finite-period harmonic modes do not supply the
analytic exactness used here.

Put $B_c=H_c^{-1}$, $B_f=H_f^{-1}$ for the unmodified metrics,
and
\[
 Z_0=B_c-T^*B_fT-b^{-1}d_{1,c}d_{1,c}^*.
\]
Equation \eqref{ch05:physical-transports} gives
$Z_0H_cd_{1,c}=0$.
Applying the two-sided lift to $H_cZ_0H_c$ and the one-sided
lift to $d_{2,f}B_fTH_c$ produces
\begin{equation}
 Z_0=(d_{2,c}B_c)^*C_0(d_{2,c}B_c),\qquad
 Sd_{2,c}B_c=d_{2,f}B_fT.
 \label{ch05:defect-lifts}
\end{equation}
The estimates are
$\|C_0\|\le C_0'$ and $\|S^*S\|\le C_1L^{-2}$.
Indeed a fixed fine-residue entry of $T$ or $B_fT$ has
weighted coefficient bound $CL^{-2}$ and one fine difference
has bound $CL^{-3}$. A Gram product sums over $L^4$ residues.
Its bounds are thus $L^4(CL^{-2})^2=O(1)$ and
$L^4(CL^{-3})^2=O(L^{-2})$, respectively.
The lift acts on coarse exterior indices separately for every
fine residue, so it introduces no further residue sum.

Modify the inverse metrics as in
\eqref{ch05:metric-modification}, with coefficients
$\zeta_c,\zeta_f$. Write the modified inverse metrics as
$\widetilde B_c,\widetilde B_f$, and
$\widetilde H_c=\widetilde B_c^{-1}$.
Set
\[
 D_c=(I+\zeta_c d_{2,c}B_cd_{2,c}^*)^{-1},\qquad
 R_0=C_0+\zeta_c I-\zeta_f S^*S.
\]
Expansion of \eqref{ch05:defect-lifts} gives the exact identity
\begin{equation}
 \widetilde H_c
   (\widetilde B_c-T^*\widetilde B_fT
                         -b^{-1}d_{1,c}d_{1,c}^*)
       \widetilde H_c
       =d_{2,c}^*D_c^*R_0D_cd_{2,c}.
 \label{ch05:positive-defect}
\end{equation}
If $\zeta_c\ge C_0'+2$ and $L^2\ge\zeta_f C_1$, then
$R_0\ge I$. Since $D_c^{-1}$ has a uniform upper bound for
fixed $\zeta_c$, the three-form metric $D_c^*R_0D_c$ is
uniformly strictly positive. Its inverse and positive square
root have exponential coefficient bounds on a smaller strip
by the usual analytic functional calculus.
This proves positivity on the whole three-form space.

For each adjacent pair of stages choose the coarse modification
large enough for its finitely many incoming possibilities.
Fix these coefficients before choosing the block ratios.
Then the condition $L^2\ge\zeta_f C_1$ is a finite lower
bound on each permitted ratio. No constant uniform over all
smoothing orders is asserted: the estimates are uniform over
histories ending at a fixed order.

\section{Convergence at a fixed scale}
\label{ch05:horizon-convergence}

The curvature metrics and transport operators in this section are those of Sections~\ref{ch05:curvature} and~\ref{ch05:transports}; their histories now vary while the observation scale stays fixed.

The preceding uniform bounds also give convergence for the particular
nested histories used in the construction. The nesting is essential:
fix the ratios, smoothing orders and modification coefficients at every
scale, and increase only the terminal cutoff. At a fixed scale $m$, the
history with terminal scale $N+1$ extends the history with terminal scale
$N$ at its fine end. An arbitrary sequence of bounded histories need
not converge.

We first record how to identify limits on the infinite lattice. Write a
translation-invariant row as $K_{ab}(z)$, where $z\in\ZZ^4$ and $a,b$
range over fixed finite coordinate sets. For an adjacent-scale row,
include the finitely many fine residues in these coordinate sets. Suppose
\[
 \sum_z e^{\sigma|z|_1}|K_{N,ab}(z)|\le C
 \qquad(\sigma>0)
\]
uniformly in $N$. Every subsequence has a further subsequence converging
in each coefficient, by diagonal extraction. Its limit has the same
bound, by finite partial sums. For every $0\le\tau<\sigma$, coefficient
convergence implies
\begin{equation}
 \sum_z e^{\tau|z|_1}|K_{N,ab}(z)-K_{ab}(z)|\longrightarrow0.
 \label{ch05:horizon-weighted}
\end{equation}
Indeed each summand is bounded by
$2C e^{-(\sigma-\tau)|z|_1}$, a summable sequence. If the symbols have
a specified limit on a dense set of real momenta, any two cluster
symbols agree there, hence everywhere by continuity. Fourier inversion
then identifies every coefficient. Alternatively, agreement of cluster
matrices on an unbounded sequence of torus sides suffices: for fixed $z$,
the periodized entry is $\sum_h K(z+nh)$, and its terms with $h\ne0$
tend to zero as $n\to\infty$. Either criterion proves uniqueness of
the cluster limit and therefore convergence of the full sequence.
This argument identifies the lattice-cover coefficients, not merely a
matrix on one fixed torus.

\begin{proposition}[Convergence of the selected Gaussian rows]
\label{ch05:horizon-limit}
Fix a scale and the nested schedule just described. The unmodified
curvature metrics, their inverses, the adjacent one- and two-form
transports, and the modified metrics of
\eqref{ch05:metric-modification} converge in the sense of
\eqref{ch05:horizon-weighted}, at every width strictly smaller than a
common positive width available at that scale. The same conclusion
holds for the specified coordinate exterior lifts, their amended
three-form metric in \eqref{ch05:positive-defect}, and their positive
square roots. It also holds for a positive square root of
$\widetilde H-cI$ whenever the fixed direct coefficient $c$ is chosen
strictly below the uniform lower bound of $\widetilde H$.
All limits are for the actual chosen rows on the cover and their
periodizations. No choice of a subsequence, matrix eigenbasis or
square-root signs is made.
\end{proposition}

\begin{proof}
Let $d=N-m$ be the history depth and $M_{d,m}$ its total dilation.
Let $u_{d,m,\mu}(p)$ be the squared modulus of its accumulated directional
filter, evaluated at $p/M_{d,m}$. Thus $u_{0,m,\mu}=1$; the filter is the
product of the centered-box and link factors used in
\eqref{ch05:history-filter}. Write $g_{m,\mu}(\theta)$ for the
single-step directional filter at scale $m$, including the link average
$r_{L_m}(\theta_\mu)$. Composition gives
\[
 u_{d+1,m,\mu}(p)=|g_{m,\mu}(p/L_m)|^2
                         u_{d,m+1,\mu}(p/L_m).
\]
Every factor is an average of unit complex numbers on real momenta,
possibly multiplied by a unit phase. Its modulus is at most one.
Induction, with the schedule fixed at each scale, gives
$0\le u_{d+1,m,\mu}(p)\le u_{d,m,\mu}(p)\le1$.
Thus the directional filter squares converge at every real $p$, including
any fixed reciprocal translate.

The accumulated noise is also convergent. Extending the history adds one
independent noise contribution propagated through all the previously
fixed coarser steps. The old contributions are unchanged, and every
new contribution is nonnegative at real momenta. The accumulated noise
therefore increases with depth. The suffix estimate in
Proposition~\ref{ch05:history-prop} bounds it uniformly, and the noise
at the last blocking step contributes one. Its limit is finite and at
least one. This argument concerns the actual nested noise list; it is
not a monotonicity assertion for arbitrary histories of the same length.

For the initial covariance contribution, extend the finite nonprincipal
alias sum to all $z\in\ZZ^4\setminus\{0\}$ by zero outside the centered
representatives modulo $M=M_{d,m}$. Its summand at an eventually included
fixed $z$ is
\[
 \frac{u_{d,m,\mu}(p+2\pi z)}{M^2\lambda((p+2\pi z)/M)}.
\]
The denominator tends to $|p+2\pi z|^2$, which is positive for nonzero
$z$ and $p$ in the base cube. The alias estimates above dominate these
summands by $C\prod_\nu(1+|z_\nu|)^{-2}$, independently of depth.
Dominated convergence proves convergence of the sum. At $p\ne0$ the
principal summand converges by the same argument without an alias sum.
Consequently each directional covariance $c_{N,m,\mu}(p)$ has a finite
limit at real $p\ne0$, bounded below by $b^{-1}$.

The curvature metric is the rational expression
\eqref{ch05:H0} in these covariances and the exterior derivatives. Its
scalar quotient denominator is
$\sum_\mu|q_\mu(p)|^2/c_{N,m,\mu}(p)$; its limit is positive away from
zero. Thus its symbol converges there. The uniform analytic bounds
already proved and the cluster-limit criterion identify one limit on
the cover, including at zero. Inverses have the same conclusion.
This avoids assigning a finite massless covariance to the zero mode.

For adjacent transports the ratio $L=L_m$ is fixed as $N$ increases.
At $p\ne0$, every one of its finitely many aliases
$\theta_l=(p+2\pi l)/L$ is a nonzero real momentum. All covariance
factors in \eqref{ch05:transport-formula} therefore converge.
The denominators $\rho,\rho_l$ have positive limits, and so does
\[
 \alpha=\sum_l |m_l|^2\tr D_l:
\]
the principal centered-box factor is nonzero on the real base cube and
$D_l$ has a positive diagonal limit. Hence the finite alias formulas
and their Fourier reconstruction converge. The previously proved
uniform exponential bounds identify unique cover limits for $W,T$.
At exceptional zero momenta uniqueness follows from continuity of
cluster symbols, rather than division by a vanishing expression.

Products preserve this convergence by the weighted convolution
inequality, with the fixed finite coordinate multiplicity included.
For uniformly bounded inverse rows, cluster limits satisfy the inverse
identity on each torus; uniqueness of the finite-dimensional inverse
and recovery from unbounded periods identify their cover limit.
The coordinate exterior lifts used in \eqref{ch05:defect-lifts} are
also continuous after a width decrease. Each elementary operation is
coordinate evaluation or a divided difference
\[
 \frac{\widehat K_N(p)-\widehat K_N(p|_{p_a=0})}{e^{ip_a}-1}.
\]
Its numerator converges, its denominator is nonzero on a dense set,
and its smaller-width bound was proved in the exterior construction.
The cluster criterion consequently applies through the finitely many
coordinate operations. This requires the fixed coordinate lift of that
construction, not a fresh arbitrary solution of the exterior identity.

It follows that $C_0,S,D_c$ and $D_c^*R_0D_c$ in
\eqref{ch05:positive-defect} converge, as do the modified metrics and
inverses. Finally consider one of the stated positive square roots.
Its uniform analytic bound gives cluster rows. On each finite torus
their matrices are positive semidefinite and have the same square,
namely the limit of the underlying positive matrix. The positive
semidefinite square root is unique by the finite-dimensional spectral
theorem. Thus the cluster matrices agree on every torus, and period
recovery identifies the cover row. This proves convergence for the
chosen roots themselves, not merely for their squares.
\end{proof}

In particular all canonical path-row coefficients obtained from these
rows converge, with a strictly positive gap between the available width and
the width used to sum paths. Finite products, fixed exterior operations
and the fixed positive factorizations preserve this conclusion. The
actual block average correction introduced below is a fixed finite-range
row at the selected scale, so its application to the convergent
straight transport also preserves convergence. These are statements
about the quadratic data. They do not imply cutoff continuity of the
nonlinear activity arrays or of marked integral maps; those require
their own stability arguments.

\section{Uniform inverse bounds: the hypotheses that matter}
\label{ch05:inverse}

A spectral gap in Euclidean norm is not, by itself, a bound for an
inverse in supremum norm independent of volume. One sufficient
framework is an exponentially local matrix with a uniform positive
gap. To state its quantitative content, assign each coordinate $i$
a lattice site $a_i$, with uniformly bounded coordinate multiplicity
per site, and define
\[
 \|T\|_\eta=\max\left\{
 \sup_i\sum_j e^{\eta d(a_i,a_j)}|T_{ij}|,
 \sup_j\sum_i e^{\eta d(a_i,a_j)}|T_{ij}|\right\}.
\]
Here $d$ is periodic lattice distance. Suppose $H=H^*\ge mI$ and
$\|H\|_{\eta_0}\le C$ uniformly in the volume.
For some $0<\eta<\eta_0$, depending only on $m,C,\eta_0$ and the
coordinate multiplicity, $H^{-1}$ has uniformly bounded
exponentially weighted row and column sums, after reducing the
exponent once more if necessary. This is the usual weighted
resolvent estimate behind exponentially local lattice propagators;
compare \cite{Balaban1984Propagators}.
It follows by conjugating with
$e^{\eta d(a_i,a_{i_0})}$ and applying the resolvent identity to the
resulting small perturbation. In particular
$\|H^{-1}\|_{\infty\to\infty}$ is bounded independently of volume.

For analytic momentum-dependent matrices, an analogous perturbation
argument must use holomorphic coordinates. Suppose $K(p)$ is
holomorphic in a strip of width $\sigma$ and, in fixed coordinates,
$\|K(p)\|\le A$ there, while at real $x$ its least singular value
is at least $a>0$. Cauchy estimates and the inverse perturbation
lemma give a smaller strip, depending only on $a,A,\sigma$,
on which $K^{-1}$ is holomorphic and uniformly bounded. This
statement is dimension independent.

In weighted Fourier spaces the weights often involve
$|q(\operatorname{Re}p)|$ and are not holomorphic. The preceding
argument is justified if, at every real anchor $x$, the matrices
defining the weighted norms have holomorphic replacements in a
fixed neighborhood of $x$, uniformly equivalent to the original
norms, and coinciding with them on the vertical slice
$\operatorname{Re}p=x$. One applies Cauchy's estimate in these
replacement coordinates. A bound in nonholomorphic coordinates
alone does not justify differentiating the weighted matrix.
This distinction is essential when the weight degenerates with a
scale parameter.

\section{The actual block average and the forest coordinates}
\label{ch05:actual-coordinates}

We now identify the nonlinear coordinates behind the Gaussian calculation.
Fix $j,L$, a coarse period $s\ge100(j+8)$, and in each centered block
$Ly+\mathbb B_{(L-1)/2}$ a spanning tree rooted at $Ly$. Trees contain
only internal block edges. Denote by $\iota_F$ extension from free-edge
coordinates to all fine edges, with value zero on tree edges. For coarse
links $V$, their replica is
\[
 (\operatorname{rep}V)_e=
 \begin{cases}1,&\operatorname{parent}(e_-)=\operatorname{parent}(e_+),\\
 V_{y,\mu},&\operatorname{parent}(e_-)=y,
                    \operatorname{parent}(e_+)=y+e_\mu.
 \end{cases}
\]
The centered-block convention and odd $L$ make these the only possibilities
for a positive fine edge in direction $\mu$. In particular the replica
is one on every tree edge. The actual forest chart is
\begin{equation}
 U(v,\psi)_e=\exp((\iota_F\psi)_e)
                         (\operatorname{rep}\exp v)_e.
 \label{ch05:forest-chart}
\end{equation}
It is a chart on the tree slice, not a choice of independent variables on
tree edges. Its derivative at the identity is exactly
\begin{equation}
 DU(0,0)[h,z]=\operatorname{rep}h+\iota_Fz,
 \qquad DU(0,0)[0,z]=\iota_Fz.
 \label{ch05:forest-derivative}
\end{equation}
All color coordinates are orthonormal for the trace metric. Thus
$\|\iota_Fz\|_2=\|z\|_2$ in counting norms. This identifies the vertical
coordinates used by the Hessian with the arrays in the tree inequality
below; no volume-dependent coordinate conversion occurs.

\subsection{A common analytic branch for all local labels}

The minimizing frame and block average of Chapter~\ref{ch04:blocking} are
real analytic near the identity. For the frame this follows directly from
its chart equation: its quadratic objective at zero is
$\tfrac12\operatorname{av}_{r,r'}|h_r-h_{r'}|^2$, positive definite after
$h_0=0$ is imposed. For the block average, its Hessian in the group variable
is the identity. Complexify these finite-dimensional real analytic chart
equations and apply the analytic implicit-function theorem
\cite[Theorem 1.3.8]{LeblSCV}. The resulting critical points are the unique
minimizers already identified in Lemma~\ref{ch04:frame-lemma}, after
reducing the neighborhood if necessary.

At fixed $j,L$ every such equation reads a translate of one fixed finite
patch. Choose a radius in the supremum norm of all its input logarithms
on which these analytic extensions and their derivatives are bounded.
There are finitely many directions and local tree shapes at fixed $L$;
take their common radius and bounds. Periodic pullback to a patch has
supremum norm at most one, so these bounds are independent of $s$.
Decrease this radius until every loop test and minimum-cost test is
strictly on its minimizing branch. Every local label then selects the
same frame and block average as a function of the full fine field on the
whole common ball. Redundant uniform branch masses are constant there.
This is equality of the local action functions before stationarity is
imposed, not merely equality along a stationary graph. It concerns the
uncut action; shrinking smooth cutoff plateaux are treated separately.

Let $a$ be an infinitesimal fine field. Write $I_a(r\to r';o)$ for its
signed sum along a coordinate path in the frame patch at $Ly$. The
linearized minimizing-frame equation gives
\begin{align}
 \tau_y(r)&=-\operatorname{av}_{r',o}I_a(Ly+r\to Ly+r';o),\notag\\
 H_{y,r}(a)&=\tau_y(r)-\tau_y(0),&
 R a(y)&=-\sum_rw_r^{\rm av}H_{y,r}(a).
 \label{ch05:frame-linearization}
\end{align}
Indeed its equation is $H_r-\operatorname{av}_{r'}H_{r'}
=-\operatorname{av}_{r',o}I_a(r\to r';o)$, and $H_0=0$ fixes the additive
constant. Differentiating the actual connector and then its block average
therefore yields
\begin{equation}
 \mathcal Q a=M a+d_{0,c}R a,
 \qquad \mathcal Q d_{0,f}\phi=d_{0,c}(\phi|_{L\Lambda_s}).
 \label{ch05:actual-Q}
\end{equation}
Here $M$ is exactly the straight length-$L$ average with the convolution
weights used in \eqref{ch05:alias-recursion}. The second identity follows
also by differentiating gauge covariance at block roots. Each route has
length $O_j(L)$; hence $R$ has finite coarse range and absolute row sum
at most $C_jL$. These statements include the root normalization. Dropping
it would change $\mathcal Q$.

\subsection{The tree comparison and the actual Hessian}

For a fine edge $e$, join its endpoints to their block roots along the
chosen trees. In the tree slice the resulting walk has transport $U_e$.
It lies in one or two adjacent blocks and has length at most $2L^4-1$.
In the macro axial gauge of Section~\ref{ch04:predictor-geometry}, its
distance from the identity is bounded by this length times the sum of bond distances from the identity in the macro-box. The
first-power frame estimate controls the two root-to-frame discrepancies;
the block average estimate controls the coarse root-to-root transport. Thus
\begin{equation}
 \chi\bigl(U_e(\operatorname{rep}V)_e^{-1}\bigr)
 \le C_{j,L}\mathcal P_y(U)
       +\mathbf1_{\{e\text{ crosses blocks}\}}
                            \chi(V_{y,\mu}Q_{y,\mu}(U)^{-1}).
 \label{ch05:nonlinear-tree}
\end{equation}
For an internal edge the walk closes at its root and only the first term
is needed. For a crossing edge, insert the block average between the two
root frames and use the triangle inequality for the normalized Hilbert--Schmidt distance. The endpoint errors
are bounded by the sum in \eqref{ch04:frame-first-power}, so this step
requires no pointwise frame estimate not already proved. The constant is
allowed to depend on $L$; the stronger pre-$L$ factor in
\eqref{ch04:macro-comparison} is not required for this use.

Apply \eqref{ch05:nonlinear-tree} to $U=\exp(\epsilon a)$ with $a$ zero
on tree edges, and $V=1$. Divide by $|\epsilon|$ and let $\epsilon\to0$.
The first-order variation of the distance from the identity is the color norm, so
\[
 |a_e|\le C_{j,L}\sum_{p\text{ in the macro-patch at }y}|(d_{1,f}a)_p|
       +\mathbf1_{\{e\text{ crosses blocks}\}}|(\mathcal Qa)_{y,\mu}|.
\]
Square, use the finite macro-patch cardinality, and sum over edges. A fine
plaquette belongs to a bounded number of macro-patches, and each coarse
edge occurs for at most a fixed number of fine edges at fixed $L$.
Consequently
\begin{equation}
 \|a\|_2^2\le B_{j,L}\|d_{1,f}a\|_2^2
                       +C_{j,L}\|\mathcal Qa\|_2^2,
 \quad a=\iota_Fz.
 \label{ch05:tree-comparison}
\end{equation}
Both constants are uniform in the period and in the allowed trees.
The use of periodic lifts preserves the bounded multiplicity.

Let $\mathcal A_f$ be the uncut fine action, without the common inverse
coupling or inherited loop action, and set
\begin{equation}
 f(v,\psi)=\mathcal A_f(U(v,\psi))
       +b\sum_{e\text{ coarse}}w\bigl((\exp v)_e
                                   Q_e(U(v,\psi))^{-1}\bigr).
 \label{ch05:actual-action}
\end{equation}
Assume its actual quadratic Taylor coefficient is
$K_f=d_{1,f}^*H_fd_{1,f}$ with $H_f\ge\varepsilon I>0$. This coefficient
is read from the same path rows in $\mathcal A_f$: each quadratic path
square contributes the square of its linearized curvature row, and atoms
of curvature degree at least three contribute nothing at second order.
Identifying their sum with a chosen metric $H_f$ is an exact coefficient
condition on the incoming path action.

The first derivatives of both terms of \eqref{ch05:actual-action} vanish
at the identity. Therefore the second derivative of the forest chart does
not enter its Hessian. By \eqref{ch05:forest-derivative}, for
$a=\operatorname{rep}h+\iota_Fz$ and
$a'=\operatorname{rep}h'+\iota_Fz'$, it is precisely
\begin{equation}
 D^2 f(0)[(h,z),(h',z')]
   =\langle d_{1,f}a,H_fd_{1,f}a'\rangle
           +b\langle h-\mathcal Qa,h'-\mathcal Qa'\rangle.
 \label{ch05:actual-hessian}
\end{equation}
No derivative of the nonlinear block average beyond its actual first
derivative is missing here: its residual vanishes at zero. In particular,
\begin{equation}
 D_\psi^2f(0,0)\ge
       \frac{\min(\varepsilon,b)}{1+B_{j,L}+C_{j,L}}I.
 \label{ch05:actual-gap}
\end{equation}
This proves the vertical gap for the actual forest chart, rather than
assuming it from the flat Fourier calculation.

\subsection{A smooth representative of the stationary tangent}
\label{ch05:actual-tangent}

Let $W_s$ denote the straight-average transport already constructed in
Section~\ref{ch05:transports}. Its identities in physical coordinates are
\[
 K_fW_s=M^*K_c,\qquad MW_s=I-b^{-1}K_c.
\]
These hold for the same incoming metric and penalty as
\eqref{ch05:actual-hessian}. Since $K_cd_{0,c}=0$, replacing $M$ by
$\mathcal Q=M+d_{0,c}R$ does not change the right side of the first
identity. The second identity does require a gradient correction.

Choose a smooth compactly supported cardinal partition of unity
$\varphi$ on $\RR^4$: $\varphi(n)=\mathbf1_{n=0}$ for $n\in\ZZ^4$ and
$\sum_{n\in\ZZ^4}\varphi(x-n)=1$. For example, in one dimension divide a
smooth bump supported in $(-1,1)$ and positive there by the sum of its
integer translates, and take the fourfold product. Define a periodic
extension of coarse site fields to fine sites by
\[
 (E h)(x)=\sum_{y\in\Lambda_s}\sum_{k\in\ZZ^4}
                \varphi(x/L-y-sk)h(y).
\]
It equals $h(y)$ at $x=Ly$, is supported within one coarse block of its
input, and has fine differences bounded by
$\|\nabla_f^rE\|_{\infty\to\infty}\le C_rL^{-r}$. Now put
\begin{equation}
 W_a=W_s-d_{0,f}ERW_s.
 \label{ch05:corrected-transport}
\end{equation}
By \eqref{ch05:actual-Q},
\begin{equation}
 \mathcal QW_a=MW_s=I-b^{-1}K_c,
 \qquad K_fW_a=\mathcal Q^*K_c,
 \qquad d_{1,f}W_a=d_{1,f}W_s.
 \label{ch05:actual-transport-identities}
\end{equation}
Thus $W_a$ minimizes the actual quadratic energy and has the same reduced
quadratic form $K_c$. This holds also on constant modes: $E$ extends
constants to constants, gradients have zero curl, and no inverse of a
zero-momentum Laplacian has been introduced.

The correction preserves the needed smooth interpolation bounds. The
finite-range row $R$ has size $C_jL$ and the already proved kernel of
$W_s$ has size $C_jL^{-1}$ with exponential coarse decay. Their composition
$RW_s$ therefore has an $O_j(1)$ exponentially decaying coarse kernel.
Applying $r+1$ fine differences to $E$ gives
\begin{equation}
 |\nabla_f^rW_a(Ly+x,y')|
       \le C_{j,r}L^{-1-r}e^{-c_jd(y,y')},
 \quad r\le j-2,
 \label{ch05:actual-bridge}
\end{equation}
with the same permitted orders as for $W_s$. The extension adds only
bounded coarse range, so it preserves exponential summability. This
argument uses finite-difference bounds of a cardinal extension, not an
assumption that a discontinuous block-constant extension is smooth.

Finally adjust $W_av$ by the unique gradient $d_{0,f}\phi$ whose site
potential is zero at every block root and makes all tree edges vanish.
It is obtained successively along each tree. Such a gradient leaves
$\mathcal QW_av$ unchanged by \eqref{ch05:actual-Q}, and leaves its curl
unchanged. Write the resulting representative as $W_Fv$. Coercivity
\eqref{ch05:actual-gap} makes it the unique minimizer in the forest slice.
Once a stationary graph $\psi=\rho(v)$ has been constructed, differentiating
its equation and using \eqref{ch05:actual-hessian} proves
\begin{equation}
 \operatorname{rep}v+\iota_FD\rho(0)v=W_Fv.
 \label{ch05:stationary-tangent}
\end{equation}
The smooth representative $W_a$ and the forest representative $W_F$ are
distinct. Estimate~\eqref{ch05:actual-bridge} is asserted for $W_a$, not
for tree coordinates. Gauge-invariant quadratic insertions see the same
curl in either representative.

\section{A stationary background on a uniform neighborhood}
\label{ch05:stationary}

We first state the quantitative implicit-function estimate needed
for the nonlinear action. Let $x$ be the coarse coordinates and
$y$ the free fine coordinates after gauge fixing. Both carry the
supremum norm. Write $f(x,y)$ for the action including the block
penalty, and $F(x,y)=D_yf(x,y)$.
All norms below are operator norms induced by these supremum norms.

\begin{proposition}[Uniform stationary graph]
\label{ch05:stationary-prop}
Suppose $F$ is holomorphic on a product neighborhood
$\|x\|_\infty<r_0$, $\|y\|_\infty<r_0$, with $F(0,0)=0$.
Let $H_*=D_yF(0,0)$ be invertible and assume, uniformly in the
number of coordinates,
\[
 \|H_*^{-1}\|\le J,\qquad
 \|D_yF(x,y)-H_*\|\le C\max(\|x\|_\infty,\|y\|_\infty),
 \qquad \|D_xF(x,y)\|\le B .
\]
Choose $0<r_y<r_0$ and $0<r_x\le r_y$ so that
$JC r_y\le1/4$ and $JB r_x\le r_y/4$.
Then for $\|x\|_\infty<r_x$ there is a unique solution
$y=\rho(x)$ of $F(x,y)=0$ in $\|y\|_\infty\le r_y$.
It is holomorphic, satisfies $\rho(0)=0$ and
$\|\rho(x)\|_\infty\le r_y/2$, and obeys
\[
 D\rho(x)=-[D_yF(x,\rho(x))]^{-1}D_xF(x,\rho(x)),
 \qquad \|D\rho(x)\|\le \tfrac43 JB .
\]
\end{proposition}

This is the quantitative contraction form of the implicit-function
theorem: apply it to $y-H_*^{-1}F(x,y)$ on the closed $r_y$-ball.
The two displayed smallness conditions give contraction constant
$1/4$ and displacement of the origin at most $r_y/4$.
We use the standard theorem rather than repeat its proof; holomorphic regularity is the analytic implicit-function theorem \cite[Theorem 1.3.8]{LeblSCV}. The real
smooth version is also a direct application of the implicit-function
theorem \cite[Theorem 9.28]{Rudin1976}.
For a merely $C^N$ real action the conclusion is $C^{N-1}$ for
$\rho$, not holomorphy.

In applications the operator bounds in this proposition should be
deduced from local derivative estimates before any dimension
dependent conversion of norms. For example, suppose each coarse or fine coordinate
has an anchor $a_i$, the total number of coordinates anchored at any
site is bounded uniformly in the volume, and, for $2\le q\le N$, every rooted derivative
satisfies
\[
 |\partial_{i_1}\cdots\partial_{i_q}f(w)|
 \le C\exp\left(-\eta\sum_{h=1}^q d(a_{i_t},a_{i_h})\right)
 \quad\text{for every }t.
\]
The lattice shell sum
$S_\eta=\sup_v\sum_i e^{-\eta d(v,a_i)}$ is bounded uniformly in
the torus size. Summing the second derivative gives $B\le CS_\eta$;
summing the third gives the Hessian Lipschitz bound
$CS_\eta^2$. A real Hessian gap together with exponential locality
supplies the required $J$ by the previous section.
These implications yield a uniform real neighborhood. A uniform
complex neighborhood requires the corresponding complex estimates,
which do not follow from a fixed finite number of real derivatives.

At fixed derivative order, higher derivatives of the background and
of its Hessian can then be controlled by repeated implicit
differentiation. If only polynomial dependence on an auxiliary
volume is needed, this conversion can be postponed until after the
uniform neighborhood and the first derivative have been obtained.
It must not be used to choose the neighborhood itself.
For a positive real Hessian $H(x)$, the determinant contribution is
differentiated by the standard identity
$D\log\det H[h]=\tr(H^{-1}DH[h])$.

\subsection{Supplying the local derivative estimates}

We retain the distinction from Section~\ref{ch05:actual-tangent}: $W_a$ is the smooth representative of the actual stationary tangent, whereas $W_F$ is its representative in the forest coordinates. They have the same curvature, but the fine-difference bounds apply to $W_a$.

The derivative hypotheses have a direct origin in the
exponentially weighted path expansion. We spell it out because
a bound for the action on each finite torus would not supply a
uniform neighborhood.

Suppose the uncut classical action is a sum of local analytic
terms indexed by rooted finite path collections. Let $\ell$
be the total path length plus one. Require that each collection
have connected support containing its root, and that
\begin{equation}
 \sup_z\sum_{\gamma:\ {\rm root}(\gamma)=z}
                   |c_\gamma|e^{\sigma\ell(\gamma)}\le P .
 \label{ch05:path-budget}
\end{equation}
A bounded number of paths per term is sufficient; products of
their matrix entries, normalized traces and fixed invariant
contractions are allowed. The constants $\sigma,P$ and this
number are independent of the torus.
Choose local coordinate maps with uniformly bounded analytic
extensions and fixed support radius at the chosen block ratio.

A complex coordinate of modulus at most $r$ changes each matrix
factor's norm by at most $e^{Cr}$. Therefore a term of total
length $\ell$ has complex modulus at most $Ce^{Cr\ell}$.
Choose $r$ so that $Cr<\sigma/4$. Differentiating $q$ times
selects at most $C_q(1+\ell)^q$ positions; the same estimates
hold on the smaller complex ball. If the differentiated
coordinates have prescribed anchors $a_{i_1},\ldots,a_{i_q}$,
connectedness implies
\[
 \sum_{h=1}^q d(a_{i_t},a_{i_h})\le C_R q(1+\ell)
             \quad\text{for every }t .
\]
Here $C_R$ allows for the finite support radius of the coordinate
maps, fixed at this block ratio. A fraction of the exponential
weight in \eqref{ch05:path-budget} controls this distance, and a
further fraction absorbs the polynomial $(1+\ell)^q$.
The remaining exponential weight bounds the sum over possible
roots of a term visiting a fixed anchor: the number of roots
within distance $C_R(1+\ell)$ grows only polynomially in $\ell$.
It follows, for every fixed $N$, that
\[
 |\partial_{i_1}\cdots\partial_{i_q}f(w)|
 \le C_N\exp\left(-\eta_N
                 \sum_h d(a_{i_t},a_{i_h})\right),
 \qquad 2\le q\le N ,
\]
on a common complex ball. Block penalties and
local coordinate terms add finite-range contributions satisfying
the same estimate. Smooth cutoffs are not being analytically
continued: the action here is the uncut action whose germ at
zero defines the background.

For the actual forest chart \eqref{ch05:forest-chart}, the deterministic
tree comparison and the Hessian identification have already been proved
in \eqref{ch05:tree-comparison}--\eqref{ch05:actual-gap}. The chosen
block average is analytic on a common fixed-patch ball, and all its redundant
local labels agree there as functions of the full field.
The derivative bounds above give exponential locality of that
Hessian. Section~\ref{ch05:inverse} now supplies a supremum-norm
bound for its inverse independent of volume; a sufficiently
small complex perturbation preserves this bound.
Proposition~\ref{ch05:stationary-prop} therefore applies on a
uniform neighborhood.

This argument proves the uniform stationary construction for
uncut path actions satisfying \eqref{ch05:path-budget}, the
actual incoming quadratic coefficient condition in
\eqref{ch05:actual-hessian}. The actual block average and forest coordinates
supply the local analytic-coordinate hypotheses and the deterministic
comparison \eqref{ch05:tree-comparison}. It does not require
a later matching theorem. If two local descriptions agree
as functions on the same coordinate neighborhood, uniqueness
makes their stationary graphs agree there.
Equation~\eqref{ch05:stationary-tangent} identifies the tangent with
the forest representative of the actual corrected transport $W_a$.
The stationary chain rule and \eqref{ch05:actual-transport-identities}
give the coarse quadratic form as the Hessian of the reduced action. Equality of actions only on a
previously selected graph would not suffice for this argument.

The constants are uniform in the period after the block ratio,
smoothing order and path bound have been fixed. Uniformity in
an arbitrary block ratio is a separate weighted-norm estimate;
the Euclidean tree gap above is allowed to depend on $L$.

\begin{proposition}[Uniform actual stationary chart]
\label{ch05:actual-stationary}
Fix the group, $j,L,b$, an exponential path bound
\eqref{ch05:path-budget}, and the actual incoming quadratic metric with
$H_f\ge\varepsilon I>0$. For the uncut action
\eqref{ch05:actual-action} there are radii $r_V,r_\psi>0$ and constants
$c,J,C_\rho>0$, independent of the allowed period and block trees, such
that its unique stationary graph $\rho$ satisfies
\[
 \|\rho(v)\|_\infty<r_\psi/2,
 \qquad \|D\rho(v)\|_{\infty\to\infty}\le C_\rho
 \quad(\|v\|_\infty<r_V).
\]
On the corresponding smaller real product chart the vertical Hessian is
at least $cI$ and its inverse has supremum operator norm at most $J$.
All local tag descriptions define the same action and stationary graph
there. For every prescribed finite order $N$, ordinary derivatives of
$f$ through order $N$, of $\rho$ through order $N-1$, and of the centered
log determinant of its vertical Hessian through order $N-2$ are bounded
by fixed polynomials in the total coordinate number $n$. The radii and
first-derivative bound do not acquire a factor depending on $n$.
\end{proposition}

To check the last assertions, the third rooted derivative bound gives
$\|H(w)-H(0)\|_{2\to2}\le C\|w\|_\infty$ by row and column sums.
Reduce the common ball until this is at most half the gap in
\eqref{ch05:actual-gap}. The weighted inverse argument then applies
uniformly throughout that real ball. Proposition~\ref{ch05:stationary-prop}
gives the graph and its volume-independent Lipschitz constant. Repeated
implicit differentiation isolates $H D^k\rho$ at order $k$; every other
term involves lower derivatives of $\rho$ and derivatives of $f$ of order
at most $k+1$. Converting bounded coordinate entries to ordinary
multilinear norms costs only a fixed power of $n$ at fixed order. Induction
therefore proves the asserted polynomial bounds. The identities
$D H^{-1}=-H^{-1}(DH)H^{-1}$ and
$D\log\det H=\tr(H^{-1}DH)$ give the determinant assertion, with one
additional factor $n$ for the trace. These are standard implicit and
Gaussian determinant differentiation rules; no polynomial-in-volume
bound is used to select the radii. For an application requiring $k+4$
derivatives, take the path/coordinate derivative order at least $k+6$
so that the stated determinant margin is retained.

\section{Scope of the completed Gaussian construction}
\label{ch05:pending}

For each fixed smoothing order, the finite-history construction
now provides uniformly analytic positive curvature metrics,
the specified one- and two-form transports, and a positive
three-form factor for the amended curvature defect. Its
parameter choices precede the choice of the scale sequence.
The nonlinear stationary graph is supplied for each incoming
path action with the explicit local bound and actual quadratic
coefficient condition just stated; its tree comparison and block average
branch compatibility are now proved for the chosen block kernel.

For the nested selected schedule, Proposition~\ref{ch05:horizon-limit}
now also proves fixed-scale convergence of the actual quadratic rows
at every smaller exponential width. For the actual nonlinear
step, the chosen path representation and block average must be
checked against \eqref{ch05:path-budget} and the actual quadratic
coefficient condition \eqref{ch05:actual-hessian}. The common minimizing
branch and equality of its local descriptions have been checked above. These are precise
interface conditions, not consequences of Gaussian algebra.
These conditions will be imposed on the admissible path action and
retained by the nonlinear step; a Gaussian identity alone would not
justify dropping them.
}{}
\setcounter{chapter}{5}
\IfFileExists{chapters/06-large-fields.tex}{% !TEX root = ../monograph.tex
\chapter{Large fields and local Gaussian integration}
\label{ch06:large-fields}

A step integrates every configuration. The usefulness of
small-field coordinates therefore depends on a bound for their complement.
This chapter separates two mechanisms. Large plaquette fields and
large deviations of coarse links from their averages, both indexed by a large-field index triple (a set of fine plaquettes, a set of coarse edges and a set of coarse sites, defined below), retain a deterministic positive energy.
Large fluctuation-field events, the failures of smooth local cutoffs in the remaining coordinates, are controlled
by simultaneous Gaussian tail estimates. These mechanisms use different
estimates and must be combined under the same integral. Ba{\l}aban makes the same split, into large plaquette fields and large fluctuation fields~\cite{Balaban1988Expansions}; he adds the location of a large fluctuation field to the large-field region, whereas here the large-field set is not enlarged and these events are bounded by a joint Gaussian tail estimate.

We work at two consecutive scales, the fine scale $m+1$ and the coarse scale $m$. The algebraic decompositions below
are exact. The energy and differentiated integral estimates use
uniform bounds for the chosen path components and for the non-flat matrices of a sector (a term of the decomposition below with fixed discrete data).
We construct these bounds for the same local integral below; the flat
Gaussian identities alone would not supply them.

\section{Large-field index triples and the local energy}

Fix a stage, a block ratio $L$, the penalty parameter $b>0$, and the fine and coarse
path data of \eqref{ch04:path-action}. Throughout this chapter distances
are coarse lattice distances, except that path-length truncations use true fine
length. Write
\[
 t_f=t_{m+1},\qquad t_c=t_m,\qquad B=B_m,\qquad
 R_*=\lfloor B/100000\rfloor,\qquad
 t_{\rm hi}=\widehat\beta_{m+1}^{-49/100}.
\]
The thresholds $t_m=\widehat\beta_m^{-1/2}M_m^{\rm cut}$ and
$B_m=\lceil(\log\widehat\beta_m)^2\rceil$ were defined in
\eqref{ch04:thresholds}. At fixed $L$, as $\widehat\beta_m\to\infty$,
\begin{equation}
 \frac{t_f}{t_c}\longrightarrow1,\qquad
 B^n\frac{t_f}{t_{\rm hi}}\longrightarrow0,\qquad
 B^n\frac{t_{\rm hi}^2}{t_f}\longrightarrow0
 \quad(n\text{ fixed}).
 \label{ch06:scales}
\end{equation}
The coarse torus period is at least $100(B+1)$. All constants below are
independent of that period.

For a configuration pair $(U,V)$ and the value $Q(U)$ of the averaging operation of
Chapter~\ref{ch04:blocking}, put
\[
 r_p^f=\chi(U_p),\qquad r_p^c=\chi(V_p),\qquad
 n_e=\chi(V_eQ_e(U)^{-1}).
\]
Here $\chi$ is the normalized Hilbert--Schmidt distance from the identity, and $n_e$ is called the \emph{deviation} of the coarse link $V_e$ from its average $Q_e(U)$. A \emph{large-field index triple}, for short an \emph{index triple},
$A=(A_f,A_n,A_c)$ consists of fine plaquettes, coarse edges, and
coarse sites. Its size is the sum of the three cardinalities. Its location
contains the projected corners of its plaquettes, the endpoints of its
edges, and its sites. Write $A^{[r]}$ for the coarse $r$-neighborhood of
that location. Fix $W>0$. An index triple is consistent if, at all plaquettes and edges meeting
$A^{[5B]}$,
\begin{align*}
 &r_p^f\ge t_f/4\quad(p\in A_f),&
 &r_p^f\le t_f/2\quad(p\notin A_f),\\
 &n_e\ge Wt_f\quad(e\in A_n),&
 &n_e\le2Wt_f\quad(e\notin A_n),
\end{align*}
and $A_c$ is exactly the set of coarse sites incident to plaquettes with
$r_p^c>t_c$ in that neighborhood. Every chosen value of the auxiliary index of Chapter~\ref{ch04:blocking} for a local gauge transformation must have positive
probability weight. The overlaps in the first two threshold pairs permit a partition
with smooth cutoffs; consistency does not prescribe a unique index triple.

For the exact integral use a global partition of unity before forming components.
Choose a smooth function $\vartheta_f$ equal to zero below $t_f/4$
and to one above $t_f/2$, with values in $[0,1]$, and an analogous
$\vartheta_n$ for the thresholds $Wt_f,2Wt_f$. Give a set $A_f$ of fine plaquettes
the weight
\[
 \prod_{p\in A_f}\vartheta_f(r_p^f)
 \prod_{p\notin A_f}(1-\vartheta_f(r_p^f)),
\]
and use the analogous product for the sets $A_n$ of coarse edges. Summing over all
pairs $(A_f,A_n)$ gives one at every configuration. Take $A_c$ to be the
coarse large-field set of the configuration. Each index triple of nonzero weight obeys
the consistency inequalities at every plaquette and edge, not merely within $5B$.
The weaker local definition above suffices for the bound on the local energy $\mathcal E_A$ defined below;
later component factors retain the restrictions of the global partition of unity
on their larger dependence sets. This refinement adds local
eligibility conditions without discarding any part of the original integral.

A \emph{sector} is a term of this decomposition with fixed discrete data: it fixes the index triple, its assignment to connected
components of its neighborhoods, a rooted spanning tree in every block,
and the value of the auxiliary index. Tree links are set to the identity.
These choices fix the coordinates and supports before differentiation.
Let
\[
 D=\{p\in A_f:r_p^f>t_{\rm hi}\}\cup\{e\in A_n:n_e>t_{\rm hi}\}.
\]
Let $H$ be the coarse plaquettes based at distance at most $4B/5$ from $D$. Local
counting gives
\begin{equation}
 |H|\le K B^4|D|,\qquad |A^{[2B]}|\le K B^4|A|.
 \label{ch06:counts}
\end{equation}
The constants may depend on the fixed stage and ratio.

\subsection{An exact completion of squares}

Let $F(g)=P_{\gLie}(g-g^{-1})/2$, with the orthogonal projection taken in
the fixed trace metric. The plaquette and path components of the curvature are grouped as
\[
 X_f(p)=\bigl(\sqrt{c_{\rm dir,f}}F(U_p),L_{f,R_*}(U,p)\bigr),
 \quad
 X_c(p)=\bigl(\sqrt{c_{\rm dir,c}}F(V_p),L_{c,R_*}(V,p)\bigr).
\]
The $L_{*,R_*}$ are the path sums of \eqref{ch04:path-action}, truncated
at true fine length $R_*$. Let $\chi_p^{\rm rel}$ denote the indicator of the relaxed
regularity condition \eqref{ch04:relaxed-test}. Define coarse components, whose plaquette part carries the factor $\alpha_p$,
and local action densities by
\begin{align*}
 v_c(p)&=\bigl(\alpha_p\sqrt{c_{\rm dir,c}}F(V_p),
                   \chi_p^{\rm rel}(V)L_{c,R_*}(V,p)\bigr),
 &\alpha_p&=\min(1,t_c/r_p^c),\\
 K_f(p)&=\operatorname{clip}_{t_f}(r_p^f)
                  +\tfrac12\chi_p^{\rm rel}(U)|L_{f,R_*}(U,p)|^2,\\
 K_c(p)&=\operatorname{clip}_{t_c}(r_p^c)
                  +\tfrac12\chi_p^{\rm rel}(V)|L_{c,R_*}(V,p)|^2.
\end{align*}
Set $\alpha_p=1$ when $r_p^c=0$. Here $\operatorname{clip}_t$ is the Borel quadratic/linear
function in \eqref{ch04:clip}, with its respective scale coefficients.
In particular $|v_c(p)|\le Ct_f$.

For clarity the square-completion input can be stated without choosing a
particular factorization of the metrics. Suppose flat maps for the fine curvature, the penalty term, and the
complement, $M_f,M_n,N$, satisfy
\begin{equation}
 M_f^*M_f+M_n^*M_n+N^*N=I,
 \qquad M_f^*X_f^{\rm lin}+M_n^*Z^{\rm lin}=X_c^{\rm lin}.
 \label{ch06:flat-identities}
\end{equation}
The superscript $\mathrm{lin}$ denotes linearization at the flat field.
Here $a$ and $v$ are the fine and coarse link variations, respectively,
and $Z^{\rm lin}=\sqrt b(v-\mathcal Q a)$ is their discrepancy under the
linearized averaging operator $\mathcal Q$ of \eqref{ch05:actual-Q}. For the fixed quadratic data these identities follow from the
identities \eqref{ch05:physical-transports} and \eqref{ch05:positive-defect} of Chapter~\ref{ch05:gaussian}.
Here is their realization through a factorization of the metrics, including the exact factorization hypothesis.
Let $C_\nu$ map curvature to the plaquette and path components used at
scale $\nu\in\{f,c\}$, and require
$C_\nu^*C_\nu=\mathsf H_\nu$, where $\mathsf H_\nu$ is the chosen
modified curvature metric. This is the quadratic coefficient condition
of \eqref{ch05:actual-hessian}; it is not a conclusion for an arbitrary
path action. Put $\mathsf B_\nu=\mathsf H_\nu^{-1}$ and
$P_c=C_c\mathsf B_cC_c^*$. The positive plaquette component makes $C_c$
injective. Thus $P_c=C_c(C_c^*C_c)^{-1}C_c^*$ is the orthogonal
projection onto the range of $C_c$. Equation
\eqref{ch05:positive-defect} supplies an exponentially local operator $K$
with
\[
 K^*K=\mathsf B_c-T^*\mathsf B_fT-b^{-1}d_{1,c}d_{1,c}^*.
\]
Choose
\[
 M_f=C_f\mathsf B_fTC_c^*,\qquad
 M_n=b^{-1/2}d_{1,c}^*C_c^*,\qquad
 N=\begin{pmatrix}I-P_c\\KC_c^*\end{pmatrix}.
\]
Direct multiplication gives the first identity in
\eqref{ch06:flat-identities}. For the second, use
$X_f^{\rm lin}=C_fd_{1,f}a$,
$X_c^{\rm lin}=C_cd_{1,c}v$, and
$T^*d_{1,f}=d_{1,c}\mathcal Q$. The last identity is
\eqref{ch05:physical-transports} with the linearization $\mathcal Q$ of the averaging operation in
\eqref{ch05:actual-Q}: its extra coarse gradient has zero curl.
The metric modification does not change these identities on curls.
Thus both identities \eqref{ch06:flat-identities} hold for the same action components satisfying the factorization hypothesis and
the same linearization $\mathcal Q$. The path construction used here imposes this factorization:
in Section~\ref{ch07:classical} it chooses
$C_c=(\sqrt{c_{\rm dir}}I,\rho)^T$ with
$\rho^*\rho=\mathsf H_c-c_{\rm dir}I$, and realizes these very components
by parallel-transported plaquettes. Cubic and quartic matching terms do not
change their Hessian, while the $\beta$-independent quadratic term of the path action is not part
of the classical metric. The step in Theorem~\ref{ch09:step-theorem}
uses precisely this construction. At the Wilson initialization (the Wilson density at the ultraviolet cutoff) the plaquette
and additional components have coefficients $\sqrt{c_{\rm dir}}$ and
$\sqrt{1-c_{\rm dir}}$, so their Gram matrix is $I$; the
fine metric at the ultraviolet cutoff is unmodified ($\zeta_f=0$). The first coarse modification
and every later one obey \eqref{ch05:positive-defect}. Thus the factorization
is an invariant of this construction, although it remains a
necessary premise when applying the estimate to other path actions.
Parallel-transport the coefficients
along their fixed coordinate contours using the pull-back $\operatorname{rep}V$ of the coarse field,
and discard contours of fine length greater than $R_*$. Denote the resulting
maps by $\Phi,\Phi_n,\mathcal N$ and put
$Z(e)=\sqrt b F(V_eQ_e^{-1})$. Their adjoints below are the adjoints of the
parallel-transported maps.

The vectors $\Phi v_c$ and $\Phi_n v_c$ are the parallel-transported linear
comparison terms. Throughout this chapter, a comparison term means one
of these specified vectors, not an independently chosen approximation.

Define six local terms:
\begin{align}
 E_f&=K_f+\tfrac12|\Phi v_c|^2-\langle\Phi v_c,X_f\rangle,\notag\\
 E_n&=\tfrac b2 n_e^2+\tfrac12|\Phi_n v_c|^2
                                      -\langle\Phi_n v_c,Z\rangle,\notag\\
 E_c&=\langle v_c,X_c\rangle-\tfrac12|v_c|^2-K_c,\notag\\
 E_N&=\tfrac12|\mathcal N v_c|^2,\notag\\
 E_L&=\langle v_c,\Phi^*X_f+\Phi_n^*Z-X_c\rangle,\notag\\
 E_D&=\tfrac12\langle v_c,
       (I-\Phi^*\Phi-\Phi_n^*\Phi_n-\mathcal N^*\mathcal N)v_c\rangle.
 \label{ch06:six-terms}
\end{align}
Every expression is evaluated at its corresponding plaquette, edge, or
complement-component base point. On summing over the whole torus, the cross terms
and squares cancel exactly:
\begin{equation}
 \sum(E_f+E_n+E_c+E_N+E_L+E_D)
       =\sum K_f+\frac b2\sum_e n_e^2-\sum K_c.
 \label{ch06:global-energy}
\end{equation}
This equality is algebraic and does not assume flatness. The local energy
$\mathcal E_A$ is instead the sum of the six terms at base points in
$A^{[2B]}$. Cancellation across its boundary is unavailable. The two terms $E_L,E_D$, which are not of definite sign,
must therefore be estimated locally.

\subsection{Direct bounds and the signed comparison}

Assume the parallel-transported maps have absolute row and column sums bounded by
$CL^{-2}$ and $CL^2$ for $\Phi$, respectively, and by $C$ for
$\Phi_n,\mathcal N$. Bounds involving $L^{-2}$ below require $C$ fixed
before choosing $L$. Put
\[
 R_f=\sum_{p\in A_f:r_p^f>t_f}r_p^f,
 \qquad R_n=\sum_{e\in A_n}n_e.
\]
The fine plaquette term, when $r_p^f\le t_f$, is at least
$\tfrac12(\sqrt{c_{\rm dir,f}}r_p^f-|\Phi v_c(p,0)|)^2$.
Since $|\Phi v_c|\le CL^{-2}t_f$, a plaquette of $A_f$ in this range
contributes at least $c_{\rm dir,f}t_f^2/128$ for large $L$.
Above $t_f$ the linear clip retains a multiple of $t_fr_p^f$.

There is one additional loss: a path component whose relaxed regularity condition fails is no
longer part of a square. Such a component based at $p$ has a witness $q$ with
\[
 r_q^f>t_f\exp(\operatorname{dist}(p,q)/R^{\rm rel}).
\]
For a fixed witness the number of possible bases is at most
$C(1+R^{\rm rel}\log(r_q^f/t_f))^4\le C r_q^f/t_f$.
All witnesses lie within $B/50$ for sufficiently large reference coupling.
They belong to $A_f$ by consistency. Split path endpoints
into those with distance from the identity at most $t_f$ and the rest. The row bound handles
the former, and the column bound handles the latter, giving
\[
 \sum_{\text{omitted components }p}|L_{f,R_*}(U,p)|\le C R_f.
\]
Multiplication by $|\Phi v_c|\le CL^{-2}t_f$ absorbs this loss into the
linear clip after increasing $L$. Consequently
\begin{equation}
 \sum E_f\ge\frac{c_{\rm lin,f}}2 t_fR_f
  +\frac{c_{\rm dir,f}}{128}t_f^2
                         |\{p\in A_f:r_p^f\le t_f\}|.
 \label{ch06:fine-bound}
\end{equation}
The square of the deviation similarly gives, for sufficiently large $W$ after $L$,
\begin{equation}
 \sum E_n\ge\frac b8\sum_{e\in A_n}n_e^2.
 \label{ch06:mismatch-bound}
\end{equation}

The coarse path contribution cancels exactly in $E_c$. The plaquette part is
$c_{\rm dir,c}(\alpha_p-\alpha_p^2/2)|F(V_p)|^2
-\operatorname{clip}_{t_c}(r_p^c)$.
The squared normal displacement from the tangent space is $O(\chi(g)^4)$,
so $|F(g)|^2\ge\chi(g)^2-C\chi(g)^4$ near the identity. Choose fixed
$\delta>0$ so that also $|F(g)|^2\ge3\chi(g)^2/4$ below $\delta$.
The clip coefficients satisfy $c_{\rm dir,c}>16c_{\rm lin,c}$, and hence
\begin{equation}
 E_c(p)\ge
 \begin{cases}
 -Ct_c^4,&r_p^c\le t_c,\\
 c_{\rm dir,c}t_cr_p^c/4,&t_c<r_p^c\le\delta,\\
 -c_{\rm lin,c}t_cr_p^c,&r_p^c>\delta.
 \end{cases}
 \label{ch06:coarse-bound}
\end{equation}

The remaining ingredient is genuinely about the chosen nonabelian paths
and averaging operation. Its quantitative form is recorded explicitly.
\begin{proposition}[Local signed comparison from compatible components]
\label{ch06:signed-input}
Assume \eqref{ch06:flat-identities}, the stated row and column
bounds for the parallel-transported maps, and a common positive exponential bound for the uncut path
kernels, with their coefficients satisfying the factorization hypothesis for the curvature components used
in those identities. Then, on every consistent sector for the averaging operation
of Chapter~\ref{ch04:blocking}, there are $C,K$ fixed before $L,W$ and
$Z=Z_{j,L,W}$ such that
\begin{align}
 \sum_H(E_L+E_D)&\ge
 c_{\rm lin,c}t_c\sum_{r_p^c>\delta}r_p^c
       -Ct_f(L^{-2}R_f+R_n)-Zt_f^2B^4|D|,
 \label{ch06:affected}\\
 \sum_{A^{[2B]}\setminus H}(E_L+E_D)&\ge
                                  -\frac a4t_f^2|A|.
 \label{ch06:unaffected}
\end{align}
Every coarse plaquette with $r_p^c>\delta$ lies in $H$ and their number
is at most $K|D|$. Here
\[
 a=\min\{c_{\rm lin,f}/8,c_{\rm dir,f}/128,b/32,c_{\rm dir,c}/64\}.
\]
\end{proposition}

\begin{proof}
In the lifted macro-box
of radius $(j+7)L$, let $\mathcal P_y(U)=\sum_qm_y(q)r_q^f$ be the sum of
fine plaquette distances from the identity with the contour multiplicities $m_y(q)$. The required comparison
is
\begin{equation}
 r_p^c\le\sum_{e\in\partial p}n_e+C_jL^{-2}\mathcal P_y(U),
 \qquad \sup_q\sum_y m_y(q)\le K_j.
 \label{ch06:macro-comparison}
\end{equation}
This estimate does hold for the precise branches of
Chapter~\ref{ch04:blocking}, as the following argument shows. It is a
bound linear in the plaquette sum $\mathcal P_y(U)$; comparison of squared energies of the local gauge transformations alone would give
an inadequate square-root bound.

The estimate for the averaging operation, linear in the plaquette sum, is Lemma~\ref{ch04:predictor-comparison}; its comparison uses the same axial gauge on the macro-box and the same rooted contours as here.

The same bound localizes coarse plaquettes with $r_p^c>\delta$.
If every variable $r_q^f$ or $n_e$ in their fixed macro-box is at most
$t_{\rm hi}$, \eqref{ch06:macro-comparison} bounds $r_p^c$ by
$C_{j,L}t_{\rm hi}<\delta$ after increasing the lower bound on the reference coupling.
Thus each such plaquette has a plaquette or edge of $D$ in that macro-box. The fixed
macro-box fits inside the $B$-neighborhood used to define $H$, and bounded
multiplicity of the macro-boxes gives the asserted count $K|D|$.
Separating the small-distance contribution from the large-field contribution $\Delta_p$ gives
$r_p^c\le C_{j,L,W}t_f+\Delta_p$ with
$\sum_p\Delta_p\le C(L^{-2}R_f+R_n)$.
The constant $C$ in the last bound must remain independent of $L,W$.
Applying the row and column bounds for the parallel-transported maps then gives
\eqref{ch06:affected}; the positive displayed term is obtained by absorbing
its upper bound into the enlarged loss constants. It is not a new
positivity assertion for $E_L$.

Outside $H$, all fine plaquette distances from the identity and deviations $n_e$ in the
surrounding fixed macro-boxes are at most $t_{\rm hi}$. The variables at
plaquettes and edges outside $A_f$, $A_n$ obey this too, since $2Wt_f\le t_{\rm hi}$ eventually. The
comparison \eqref{ch06:macro-comparison} bounds coarse plaquette distances from the identity there by
$C_{j,L,W}t_{\rm hi}$. Choose a coarse axial gauge on a box of radius
$B/2$. Its sum of bond distances from the identity is bounded by $C_{j,L,W}B^5t_{\rm hi}$.
To compare a fine link with the corresponding link of $\operatorname{rep}V$, join its two
endpoints to their block roots along the tree. All contour variables along the tree
are the identity. The resulting contour lies in two neighboring blocks.
In an axial gauge on the macro-box its distance from the identity is bounded by the length of the contour times
the sum of bond distances from the identity in the macro-box. Comparing a crossing contour with the value of the averaging operation
adds the errors of the local gauge transformations at the two endpoints, already bounded in
Lemma~\ref{ch04:predictor-comparison}, and then the deviation $n_e$.
Consequently, with a constant depending on the fixed tree family and $L$,
\[
 \chi\bigl(U_e(\operatorname{rep}V)_e^{-1}\bigr)
       \le C_{j,L}\mathcal P_y(U)+n_{e_c}
\]
for crossing edges; omit the deviation for internal edges. Any tree in a
fixed block has uniformly bounded path length (at most $L^4-1$), which
suffices here because this constant may depend on $L$. This is not a
bound on $\Delta_p$ with a constant independent of $L$. The parent lift of the coarse gauge preserves
these distances between corresponding links. It therefore places both fine and coarse
logarithms in a common chart, throughout radius $2B/5$, with bound
$\delta_{\rm loc}=C_{j,L,W}B^{10}t_{\rm hi}$.

The derivative of the rooted averaging operation also matches the flat
identity. Near the identity all branches of the local gauge transformations and of the averaging operation chosen by the auxiliary index
are minimizing; the redundant values of the auxiliary index do not change their values. The
finite-patch minimizing local gauge transformation is smooth there by the same nondegenerate
implicit equation used in Lemma~\ref{ch04:frame-lemma}. Write
$h_{y,z}(\exp(\epsilon a))=1+\epsilon H_{y,z}(a)+O(\epsilon^2)$ and
$\overline H_y=\sum_z w_z^{\rm av}H_{y,z}$. Differentiating
the transporter \eqref{ch04:connector}, and then the averaging operation applied to it, gives
\[
 (\mathcal Qa)_{y,\mu}
    =(\mathcal Aa)_{y,\mu}+\overline H_y(a)
                                -\overline H_{y+e_\mu}(a),
\]
where $\mathcal A$ averages the straight length-$L$ link sums with the
same weights. Thus $\mathcal Q-\mathcal A$ is a coarse gradient. Its curl
vanishes exactly. This verifies the identification of the linearized averaging operation needed in
\eqref{ch06:flat-identities}. Their realization through the factorization of the metrics was derived
above from \eqref{ch05:physical-transports} and
\eqref{ch05:positive-defect}; the incoming path action must still satisfy
the exact quadratic coefficient condition in
\eqref{ch05:actual-hessian}. Second-order
Taylor bounds are uniform in the period because each value of the averaging operation reads a
fixed finite patch at fixed $j,L$.

Set
$\varepsilon_{\rm loc}=(1+R_*)\delta_{\rm loc}+e^{-\tau R_*}$,
where $\tau>0$ is a common path decay rate. Subtracting the flat identities
\eqref{ch06:flat-identities} before estimating terms gives the required
bounds
\begin{align*}
 |\Phi^*X_f+\Phi_n^*Z-X_c|&\le C\delta_{\rm loc}\varepsilon_{\rm loc},\\
 |(I-\Phi^*\Phi-\Phi_n^*\Phi_n-\mathcal N^*\mathcal N)v_c|
                                    &\le Ct_f\varepsilon_{\rm loc}.
\end{align*}
Indeed the contour variable of a contour of length $s$ differs from its flat value by at most
$Cs\delta_{\rm loc}$; the omitted tails have size $Ce^{-\tau R_*}$.
The count \eqref{ch06:counts} and the limits \eqref{ch06:scales} yield
\eqref{ch06:unaffected}. This completes the verification of Proposition~\ref{ch06:signed-input}.
It uses the flat identities, with their factorization, for the chosen path family, not merely
a decay estimate for a flat Fourier kernel.
\end{proof}

\begin{proposition}[Energy bound under the local comparisons]
\label{ch06:energy}
Under the hypotheses of Proposition~\ref{ch06:signed-input}, choose
$L$ sufficiently large, then $W$ sufficiently large, then the reference
coupling sufficiently large. Uniformly in the period and consistent sector,
\[
 \mathcal E_A(U,V)\ge\kappa t_f^2|A|,\qquad \kappa=a/2>0.
\]
The lower bound on the reference coupling may depend on $j,L,W$; $a$ is fixed by the coefficients $c_{\rm dir}$, $c_{\rm lin}$ of the
clip functions and by $b$.
\end{proposition}
\begin{proof}
The bounds \eqref{ch06:fine-bound} and \eqref{ch06:mismatch-bound} leave
$a t_f^2$ per plaquette of $A_f$ or edge of $A_n$, $a t_ft_{\rm hi}$ per element
of $D$, and enough linear surplus to absorb
$Ct_f(L^{-2}R_f+R_n)$. The choices are in that order: increase $L$ for
the term $L^{-2}R_f$, and $W$ for the term $R_n$.

Assign each coarse large-field site to one incident plaquette with distance from the identity above
$t_c$. A plaquette has at most four such sites. Since $t_f\le2t_c$
eventually, the middle range of \eqref{ch06:coarse-bound} contributes
$4at_f^2$ per plaquette and therefore $at_f^2$ per assigned site. The
exceptional plaquettes with distance from the identity above $\delta$ number at most $K|D|$;
their negative energies cancel the positive term in
\eqref{ch06:affected}. Their missing site contributions amount to at most
$4aKt_f^2|D|$. Also $E_N\ge0$. The total remaining losses are bounded by
\[
 (4aK+ZB^4)t_f^2|D|+Ct_c^4KB^4|A|+\frac a4t_f^2|A|.
\]
Increase the lower bound on the reference coupling until
$(4aK+ZB^4)t_f^2\le at_ft_{\rm hi}$ and
$Ct_c^4KB^4\le at_f^2/4$. Equation~\eqref{ch06:scales} justifies both
choices. The first loss is absorbed by the positive energy contribution retained for $D$, leaving
at least $at_f^2|A|/2$.
\end{proof}

\section{The Gaussian of a sector at a non-flat background}
\label{ch06:sector}

Fix the sector and the conditioning variables $Y$: the non-tree links that are held fixed in the Gaussian integral and integrated afterwards, called \emph{frozen} below. (The word conditioning is Ba{\l}aban's~\cite{Balaban1988Cluster}; in~\cite{Balaban1988Expansions} the integrations over such variables are left to later steps, whereas here they are carried out in the same step.) Tree links are not integration
variables. \emph{Active} coordinates, the integration variables, are logarithms (Lie-algebra coordinates) of the remaining non-tree links
outside $A^{[10B]}$. Linearize the retained square
terms at the pull-back $\operatorname{rep}V$ of the coarse field, using the chosen averaging operation
and keeping the cutoffs fixed. Collect the resulting affine expressions
in a single vector $\mathsf D u+b_Y$, where $u$ contains the active
coordinates and $b_Y$ is the value at $u=0$. Their quadratic energy is
\begin{equation}
 Q_Y(u)=\tfrac12|\mathsf D u+b_Y|^2,\qquad
 \mathsf A=\mathsf D^*\mathsf D,\qquad f_Y=\mathsf D^*b_Y.
 \label{ch06:actual-matrix}
\end{equation}
If hypercubic orientations are averaged, this affine map uses the square roots
of their weights. A component is retained when its \emph{link support}, the fixed set of links from which its expression is built, meets an active
edge. This rule differs from retaining components whose base points are a specified
distance from the location of $A$. A comparison between the corresponding
matrices requires a proof. Identity coordinates may be adjoined for
inactive variables, with their trivial integrals normalized to one.

The link support contains the links on which the averaging operation depends and every endpoint of a retained
path; here $r_{\rm pred}$ is the largest coarse distance at which the averaging operation reads a link. Its coarse radius is at most $R_*+5+r_{\rm pred}\le B/50$ for large
coupling. Thus the link support of a retained component lies in the small-field region when its
active edge is outside $A^{[10B]}$. This geometric
observation permits local estimates at the background field $\operatorname{rep}V$; it does not by
itself prove a lower bound on $\mathsf A$.

\subsection{Transferring a local estimate in the block axial gauge to the non-flat components}

The gap needed below can be obtained without identifying non-flat and flat
matrices. The geometric comparison just proved implies a local inequality in the
specified block axial gauge:
\begin{equation}
 |a_e|\le C_L\sum_{p\in\mathcal P_e}|(d_fa)_p|
                    +|\mathcal Q a(e_c)|.
 \label{ch06:tree-control}
\end{equation}
Here $\mathcal P_e$ is the set of fine plaquettes of a fixed macro-box,
$e_c$ is the coarse edge in the direction of $e$, and $\mathcal Q$ is
the derivative of the averaging operation at the identity. The inequality
is for arrays vanishing on the chosen tree, extended by zero on frozen
edges. For internal edges the last term may be omitted. Its nonlinear
origin is worth specifying: join the endpoints of $e$ to their block
roots along the same tree, compare the resulting contour variable with the value of the
averaging operation in a common axial gauge on the macro-box, and bound the difference by the
plaquette sum of the macro-box. Differentiating this comparison at the identity
gives \eqref{ch06:tree-control}. The contour constants and endpoint bounds are those proved above for the
same trees and the same values of the auxiliary index. No coordinate change to an unrelated tree is made.

The Poincar\'e-type inequality \eqref{ch05:tree-comparison} applies to these
coordinates: the pull-back $\operatorname{rep}V$ is varied by left multiplication on non-tree bonds,
tree coordinates are zero, and the averaging operation is the rooted one, with the linearization $\mathcal Q$
of \eqref{ch05:actual-Q}. Thus the retained plaquette components and the retained components containing the averaging operation, for short \emph{averaging components},
are exactly its derivatives, multiplied by their fixed positive
square weights. Each such derivative differs, in one common gauge on its
macro-box, from its flat derivative by at most $\varepsilon$ times the
sum of input norms on its fixed link support. This follows from continuity of
the finite-patch derivative on the minimizing branch once clearance and
small coarse plaquettes put that entire patch in a small common gauge.
Insert the comparison in \eqref{ch06:tree-control}, square, and sum over
active edges. Fixed multiplicities of the link supports give
\[
 \|a\|_2^2\le C_L\sum_p|D_p a|^2+C_L\sum_e|D_e^{\rm pred}a|^2
                         +C_L\varepsilon^2\|a\|_2^2.
\]
All norms here are unnormalized $\ell^2$ norms, without division by volume. The choice of retained
components introduces no error: an omitted component has a link support disjoint
from the active coordinates, hence annihilates their zero extension.
For $C_L\varepsilon^2\le1/2$ and positive lower plaquette and penalty
coefficients, this proves
$\mathsf D^*\mathsf D\ge c_LI$. The other components add nonnegative squares.
The gauge may depend on the edge at which the local inequality is applied.
It is fixed during differentiation, and invariance of the norm on $\gLie$
permits the subsequent sum. The resulting lower bound $c_L$ may depend
on $L$. The constant in the block comparison
\eqref{ch06:macro-comparison} must already be fixed when $L$ is chosen.

\subsection{A localized approximate inverse}

For an edge coordinate $i$, let $E_i$ be its two projected endpoints.
Two useful symmetric windows are
\begin{align*}
 \Omega_i^{\rm sec}&=\{i':\max_{x\in E_i,y\in E_{i'}}
                                      \operatorname{dist}(x,y)\le B/20\},\\
 \Omega_i^{\rm loc}&=\{i':\min_{x\in E_i,y\in E_{i'}}
                                      \operatorname{dist}(x,y)\le B\}.
\end{align*}
Restrict them to the active coordinate set. Both contain a ball of radius
$B/60$ about $i$ for large $B$. Use one prescription throughout a given
integral and define
\begin{equation}
 P_{i'i}=\mathbf1_{i'\in\Omega_i}
                    (\mathsf A[\Omega_i,\Omega_i]^{-1})_{i'i},
 \qquad \Gamma=\tfrac12(P+P^*),\qquad m_Y=-P^*f_Y.
 \label{ch06:windows-def}
\end{equation}
Here $\mathsf A[\Omega_i,\Omega_i]$ denotes the principal submatrix on
$\Omega_i$: to obtain column $i$ of $P$, solve the restricted system with
right-hand side the $i$th coordinate vector and extend the solution by
zero. The comparison below shows that the finite-range symmetric matrix
$\Gamma$ is positive definite and hence defines a Gaussian covariance.

\begin{proposition}[Window estimates from bounds on the non-flat matrix]
\label{ch06:windows}
Assume $cI\le\mathsf A\le CI$ and an exponentially weighted Schur bound
$\sup_i\sum_{i'}e^{\tau d(i,i')}|\mathsf A_{ii'}|\le C$, with the analogous
column bound. Assume bounded coordinate multiplicity per coarse site.
Then, for some $\eta,c'>0$, uniformly in the active set, period, and
principal compression,
\begin{equation}
 \|P-\mathsf A^{-1}\|_\eta+
 \|\Gamma-\mathsf A^{-1}\|_\eta\le C'e^{-c'B}.
 \label{ch06:window-error}
\end{equation}
Here $\|\cdot\|_\eta$ is the maximum of the exponentially weighted row and
column sums. For sufficiently large $B$, $c_1I\le\Gamma\le C_1I$.
Fixed-order derivatives obey corresponding bounds if the matrix $\mathsf A$
has them, with all windows and retained coordinate sets held fixed.
\end{proposition}
\begin{proof}
The weighted conjugation argument of Section~\ref{ch05:inverse} gives
exponential inverse bounds for $\mathsf A$ and all its principal
compressions. The constants are uniform because a principal compression
preserves the gap and the weighted bound. The column $Pe_i$ solves
$\mathsf APe_i=e_i$ inside $\Omega_i$. Outside that window its residual
crosses a distance at least $B/60$ from $i$. The convolution of two
exponentially decaying kernels therefore gives
$\|\mathsf AP-I\|_\eta\le C'e^{-c'B}$ after decreasing the decay rate.
Multiplication by $\mathsf A^{-1}$ proves the first estimate in
\eqref{ch06:window-error}; adjunction proves the second. Operator norm
is bounded by Schur norm, so comparison with $\mathsf A^{-1}\ge C^{-1}I$
proves positivity.

For derivatives, use $D\mathsf A^{-1}=-\mathsf A^{-1}(D\mathsf A)
\mathsf A^{-1}$ and its iterates, also on the compressions. Every term in
a differentiated residual still crosses the window boundary. Polynomial
spatial weights are moved along a kernel using
$1+d(o,i')\le(1+d(o,i))(1+d(i,i'))$ and absorbed by a smaller exponential
rate. These operations leave the factor $e^{-c'B}$ at any fixed order.
\end{proof}

\subsection{The inhomogeneous term and derivatives on a fixed sector}

The inhomogeneous term $b_Y$ in \eqref{ch06:actual-matrix} uses a bounded extension of the
frozen logarithms. This detail is essential. Let $v_Y$ be the patched,
smoothly clipped frozen displacement, extended by zero off the frozen
coordinates, with $|v_Y|_\infty\le2d$ and $d\le r t_f$, for fixed $r$.
If $b(V)$ is the value of the components at $\operatorname{rep}V$ and
$\mathsf D_{\rm all}(V)$ is its derivative on all non-tree coordinates,
including the frozen ones, the inhomogeneous term is
\begin{equation}
 b_Y=b(V)+\mathsf D_{\rm all}(V)v_Y,\qquad
 f_Y=\mathsf D^*b_Y .
 \label{ch06:clipped-source}
\end{equation}
The weights of the components are included in both terms. The original integrand is
preserved by putting the difference from this affine square into the
local correction factors. No estimate below replaces an arbitrary
frozen Haar logarithm by a small logarithm.

Here is why the inhomogeneous term is uniformly small. A retained component has an active
coordinate outside $A^{[10B]}$. Its whole
link support has radius at most $B/50$; hence its base point, endpoints, and fixed
macro-boxes have more than $9B$ clearance for large $B$. All coarse
variables $r_p^c$, $n_e$ used by the component are therefore small. In the common local gauge,
the four-link curvature formula, the absolute path-coefficient bounds,
and the smooth branch of the averaging operation give
\[
 |b(V)|_\infty\le C_b t_f,\qquad
 \|\mathsf D_{\rm all}(V)\|_{\infty\to\infty}\le C_a.
\]
The comparison term is extended without its large-field indicator when the field is
varied. Thus \eqref{ch06:clipped-source} implies
$|b_Y|_\infty\le(C_b+2rC_a)t_f$, uniformly over the frozen Haar variables.
The column bound for $\mathsf D$ now gives
\begin{equation}
 |f_Y|_\infty+|m_Y|_\infty\le C t_f,\qquad
 |\mathsf A m_Y+f_Y|_\infty\le C t_f e^{-cB}.
 \label{ch06:actual-forcing}
\end{equation}
The second inequality uses
$\mathsf A m_Y+f_Y=(I-\mathsf A P^*)f_Y$ and the window estimate.
In particular $g^{-1}(\mathsf A m_Y+f_Y)$ remains exponentially small
for $g=\beta^{-1/2}$ on the coupling band: $g^{-1}t_f=M_f^{\rm cut}$
up to fixed constants grows slower than $e^{\epsilon B}$.

For clarity, the corresponding derivative assertion has a different
scale. For any fixed allowed order $k$, use dimensionless variation
parameters whose coarse link increments and clipped frozen displacements
have derivatives bounded by $C_k t_f$. Keep the choice of retained components, the spanning trees, and
windows fixed. Differentiating a contour variable contributes a fixed power of its
length. The exponential path bound absorbs that power after decreasing
its decay rate. A derivative with base point $o$ can occur only on a contour variable
whose dependence set contains $o$; its distance to the component is bounded
by the length of the contour plus the fixed radius of the averaging operation. Consequently the
same argument bounds the derivatives of $\mathsf D$ in exponentially
weighted row and column norms, uniformly in volume. Products give the
bounds for $\mathsf A=\mathsf D^*\mathsf D$. All fixed positive-order
derivatives of $b_Y,f_Y,m_Y$ are $O(t_f)$ in these directions of size $O(t_f)$,
and the derivatives of the residual $\mathsf Am_Y+f_Y$ retain the factor $e^{-c_kB}$.
Unscaled derivatives instead carry the corresponding powers of $t_f^{-1}$.

These derivatives exist on a common neighborhood, not just on the
set satisfying the large-field indicator inequalities. Indeed the formulas for the components with the event indicator omitted
extend smoothly there. If their first derivative norm is at most $C$,
a ball of radius $c_L/[2(1+C)]$ preserves the gap $c_L/2$ of $\mathsf A$.
Every principal compression remains invertible on that ball. Inverse
differentiation and the proof of Proposition~\ref{ch06:windows} therefore
apply with the retained components fixed even if nearby fields change their large-field indicators.
The radius need not be uniform in arbitrary unnormalized coordinates.

\subsection{The exact integral and its normalization}

Let $n$ be the real dimension of the active coordinate space. Write $g=\beta^{-1/2}$, and let $\gamma_\Gamma$ denote the centered
Gaussian measure on $\RR^n$ with positive definite covariance $\Gamma$.
For any $m\in\RR^n$ and any $F$ integrable against the weight on the left,
the change of variables $u=m+gz$ gives
\begin{align}
 \int_{\RR^n}e^{-\beta Q_Y(u)}F(u)\,du
 &=\mathcal C_Y\int F(m+gz)
       \exp\bigl(\eta\cdot z+\tfrac12 z^*(\Gamma^{-1}-\mathsf A)z\bigr)
                                      \,d\gamma_\Gamma(z),\notag\\
 \mathcal C_Y&=(2\pi g^2)^{n/2}(\det\Gamma)^{1/2}e^{-\beta Q_Y(m)},\notag\\
 \eta&=-g^{-1}(\mathsf Am+f_Y).
 \label{ch06:gaussian-identity}
\end{align}
This is the usual finite-dimensional completion of squares with an
arbitrary reference covariance. We use it as an identity, not an
approximation. In particular $m=m_Y$ and \eqref{ch06:windows-def} do not
remove the remaining linear or quadratic term in the exponent.

In a sector integral $F$ includes the chart indicators, the Haar density
of every active link, the probability weights of the auxiliary index, incoming activities, and all
remaining action terms. It is extended by zero off its original chart.
The frozen variables are subsequently integrated with probability Haar
measure $dY$. When $\Gamma$ depends on $Y$, the resulting reference law is
\[
 dY\,d\gamma_{\Gamma(Y)}(z).
\]
It is a conditional Gaussian law, generally not a product of an independent
Haar field and a fixed Gaussian. The positive scalar $\mathcal C_Y$ must
be retained with the Haar Jacobians and relative determinants. A change of
tree coordinates preserves this combined integral, not its determinant
factor in isolation. Zero-dimensional active sectors have determinant one
and no Gaussian variables.

\section{Differentiation on a common dependence set}

The cutoffs and large-field indicators in the sector definition are fixed Borel conditions.
Their indicators are not differentiated. Even on a fixed sector an
incoming large-field activity can be merely measurable. Differentiating
it as a function of its fine links would be unjustified. A local change
of variables instead keeps those physical fine links fixed.

Write an active physical link as $U_e=\exp(u_e)(\operatorname{rep}V)_e$.
At a base field $V_0$, set $u=m_Y(V_0)+gz$. Within the logarithmic chart of the active links, where these logarithms are defined (its radius $d_*$ is fixed below), define
\begin{align}
 u_V(e)&=\log\!\left(\exp(u_e)(\operatorname{rep}V_0)_e
                                  (\operatorname{rep}V)_e^{-1}\right),\notag\\
 z_V&=g^{-1}(u_V-m_Y(V)).
 \label{ch06:protected-change}
\end{align}
Then $\exp(u_V(e))(\operatorname{rep}V)_e=U_e$ exactly. The following
extension makes this a legitimate change on the whole Gaussian space.

Choose a fixed logarithm radius $d_*>0$ and a smooth radial function
$\chi$ equal to one on $|u|\le d_*$ and zero on $|u|\ge2d_*$. For a
group element $M$ sufficiently near the identity put
\[
 F_M(u)=u+\psi(M,u),\qquad
 \psi(M,u)=\chi(u)\bigl(\log(\exp(u)M)-u\bigr).
\]
The logarithm is used only on the support of $\chi$. At $M=1$,
$\psi=0$ and $D_u\psi=0$. Compactness of that support permits one common
neighborhood of $1$ on which
$|\psi|<d_*/2$ and $\|D_u\psi\|<1/2$.
For every target $v$, the map $u\mapsto v-\psi(M,u)$ is a contraction
of the complete Euclidean space, so $F_M$ is bijective.
Its derivative is invertible, and the inverse function theorem gives a
smooth inverse. Its determinant is positive by continuity from $M=1$.
All fixed parameter derivatives of $\psi$ and of this determinant are
bounded uniformly in $u$, since they vanish outside a fixed compact set.
These are the usual finite-dimensional contraction and inverse-function
arguments.

Apply this map separately on each active bond with
$M_e(V)=(\operatorname{rep}V_0)_e(\operatorname{rep}V)_e^{-1}$, and set
\[
 T_Vz=g^{-1}\bigl(F_{M(V)}(m_Y(V_0)+gz)-m_Y(V)\bigr).
\]
The extension of $m_Y$ described above, without the event indicator, is used off the valid
sector. The map is globally smooth and invertible, and its Jacobian is
the product of the positive bond Jacobians; the factors $g$ cancel.
Call a factor \emph{protected} when it is cut off inside this logarithmic chart. If a protected factor is supported where $|u_V(e)|\le8a_*t_f$
and $8a_*t_f\le d_*/2$, then $|u_e|<d_*$ by the displacement bound.
There $F_M$ equals the logarithmic map and holds the physical link fixed.
Nested cutoffs, with the smaller support inside $|u_V|\le4a_*t_f$,
leave a strict margin. Thus a protected measurable factor is constant
on a neighborhood of the base parameter whenever it contributes; it
does not have to be differentiated. Cutoff remainders (defined in the next section) must likewise
vanish on a neighborhood, rather than only at the base point.

Use one coordinate change on the union $I$ of all coordinates read by the
product under consideration. If $\Gamma_I$ is its marginal covariance,
the density ratio contributed by that change is
\begin{equation}
 \mathcal R_I(V,z)=|\det D_z z_V|\,
 \frac{(\det\Gamma_I(V))^{-1/2}
        e^{-z_V^*\Gamma_I(V)^{-1}z_V/2}}
      {(\det\Gamma_I(V_0))^{-1/2}
        e^{-z^*\Gamma_I(V_0)^{-1}z/2}}.
 \label{ch06:density-ratio}
\end{equation}
The marginal is the principal covariance block $\Gamma_I$, whose inverse
is generally not a principal block of $\Gamma^{-1}$. Applying separate
changes to overlapping factors would not reproduce this common ratio.

The familiar derivative formulas
\[
 D\log\det\Gamma_I=\tr(\Gamma_I^{-1}D\Gamma_I),\qquad
 D\Gamma_I^{-1}=-\Gamma_I^{-1}(D\Gamma_I)\Gamma_I^{-1}
\]
show that the logarithmic derivative of the Gaussian part is
\[
 -\tfrac12\tr(\Gamma_I^{-1}\dot\Gamma_I)
 +\tfrac12z^*\Gamma_I^{-1}\dot\Gamma_I\Gamma_I^{-1}z
 -z^*\Gamma_I^{-1}\dot z.
\]
Higher fixed derivatives are polynomials in the coordinates, with
coefficients controlled by the window estimates and the coordinate
change. Ordinary Gaussian moment formulas, in particular
Isserlis' formula~\cite{Isserlis1918}, then give polynomial bounds in $|I|$
for fixed derivative order. The trace contributes a factor proportional to
$|I|$; it is not uniformly bounded independently of the dependence set.
For local factors with $s$ base points one typically has $|I|\le CB^4s$.

A factor's dependence set includes every cutoff condition, every contour of the averaging operation,
window matrix entry and frozen coordinate needed to evaluate it. This set
is different from the exclusion support $Y_i$ in
Definition~\ref{ch04:activities}. Disjoint exclusion supports imply no
independence. Even disjoint coordinate sets give conditional Gaussian
independence only when their cross covariance vanishes. Integration over
shared frozen variables can destroy that independence after conditioning
is removed.

\subsection{Differentiation under the common Gaussian law}

There is a useful distinction between the bounds used to justify
differentiation and those used to estimate the resulting derivatives.
The first may depend on the fixed finite volume. Since $\psi$ has
compact support and the frozen logarithms are clipped,
$T_Vz-z$ and its fixed parameter derivatives are bounded uniformly in
$z$ and in the frozen variables on one neighborhood. If
$\Gamma_I(V_0)^{-1}\ge\lambda I$, shrink that neighborhood so
$\Gamma_I(V)^{-1}\ge\lambda I/2$. The Gaussian density ratio and its
first $k$ derivatives are then bounded by
\[
 C\,p_{\Gamma_I(V_0)}(z)^{-1}
       e^{-\lambda|z|^2/8}(1+|z|)^{2k},
\]
with a finite constant $C$. Multiplication by the reference Gaussian
density makes this integrable. The Jacobian has bounded derivatives.
For the local correction factors, the endpoint cutoffs ensure bounded
arguments whenever the factor is nonzero. Hence their derivatives are
also bounded on that neighborhood. Dominated differentiation follows.
For a quadratic correction this last assertion requires both endpoints
of every nonzero matrix term to carry the protecting cutoff; it is not
a consequence of a cutoff at only one endpoint.

At the base point $T_{V_0}z=z$ and the ratio is one. Its differentiated
logarithm is a polynomial of degree at most two in $z$ and the
derivatives of $T_Vz$, and repeated exponential differentiation gives
a fixed-degree polynomial bound. Matrix products are estimated in row
and column norms; a trace contributes a factor $|I|$. Therefore, if $s$ counts the
base points of the factors and $|I|\le C B^4s$, there are fixed integers $r_k,m_k$
such that
\begin{equation}
 |D^j\mathcal R_I(V_0,z)|
 \le A_k(\widehat\beta_f)(1+s)^{m_k}
             (1+|I|+\sum_{i\in I}|z_i|)^{r_k},\quad j\le k,
 \label{ch06:base-scores}
\end{equation}
where $A_k$ is subpower, meaning
$A_k(x)\le C_\epsilon x^\epsilon$ eventually for each $\epsilon>0$.
In this statement the normalized variation derivatives of $g^{-1}m_Y$
and of $g^{-1}\psi$ must have subpower bounds; these follow from
\eqref{ch06:actual-forcing} and $t_f/g\asymp M_f^{\rm cut}$.
The relevant marginal inverses have uniform operator bounds; when
row bounds are used, a polynomial factor in $|I|$ suffices and is
already allowed in \eqref{ch06:base-scores}. No assertion about
independence of the coordinates is involved.

\section{Simultaneous estimates for large fluctuation-field events}

A cutoff remainder, a cutoff factor minus one, is bounded but need not be pointwise small. Its
smallness comes from a large fluctuation-field event. An event label specifies one
particular cutoff and its base point; repeated appearances of that event label
must be counted once. Suppose each remainder with event label $w$ vanishes
unless
\begin{equation}
 \mathcal E_w=\{\exists i:\ d(i,w)\le 2B,
                                   \ |z_i|\ge T\},\qquad
 T=\xi t_f/(2g),
 \label{ch06:tail-event}
\end{equation}
for fixed $\xi>0$. Verifying this implication uses the mean $m_Y$, the
derivative of the averaging operation, and the cutoff factors. A covariance estimate by itself cannot
prove it.

\begin{lemma}[Joint tails with bounded local multiplicity]
\label{ch06:joint-tails}
Let $z$ be zero-mean Gaussian with covariance $C\le M I$. Suppose an
event label has at most $N$ possible witness coordinates, and each coordinate
can witness at most $Q$ distinct event labels. For $r$ distinct event labels,
if $T\ge M/2+HQ$ with $H>0$, then
\[
 \mathbb P\left(\bigcap_{w=1}^r\mathcal E_w\right)
                                      \le(2N e^{-H})^r.
\]
No independence between the events is assumed.
\end{lemma}
\begin{proof}
This is the exponential Markov bound for a Gaussian, with the local
multiplicities kept explicit. For any distinct coordinate set $J$ and
signs $\sigma_i$, the variance of $\sum_{i\in J}\sigma_i z_i$ is at most
$M|J|$. Therefore
\[
 \mathbb P\{|z_i|\ge T\ (i\in J)\}
            \le2^{|J|}\exp\bigl((M/2-T)|J|\bigr).
\]
Assign one witness to each event label. There are at most $N^r$ assignments.
For the distinct coordinates in an assignment, $|J|\ge r/Q$ and
$|J|\le r$. The displayed estimate is at most $2^re^{-Hr}$.
Summing assignments proves the assertion.
\end{proof}

For the four-dimensional lattice, bounded multiplicity of coordinates
and event labels per coarse site gives $N,Q\le C_jL^d(1+B)^4$ for a fixed
finite power $d$. For $\beta$ in the coupling band,
$t_f/g\asymp M_f^{\rm cut}=\exp((\log\log\widehat\beta_f)^2)$.
It grows faster than every power of $\log\widehat\beta_f$. Thus, for any
fixed $\tau>0$, one can arrange
$H\ge2\tau\log\widehat\beta_f+\log(2N)$ and obtain
\begin{equation}
 \mathbb P\left(\bigcap_{w=1}^r\mathcal E_w\right)
                                      \le\widehat\beta_f^{-2\tau r}.
 \label{ch06:power-tail}
\end{equation}
The same statement holds uniformly under conditioning on the frozen variables if the
covariance bound and local vanishing assertion are uniform in that
conditioning.

\subsection{The deterministic estimate behind a cutoff remainder}

To distinguish the analytic issue from the probability estimate, write
$B_p(U,V)$ for a fine curvature component or averaging component minus its comparison term, and
$R_p(v)$ for its parallel-transported affine approximation on the same local
patch. On the complement of the large fluctuation-field event, $v=m_Y+gz$ has size $O(t_f)$.
Taylor expansion on the common coarse gauge, with its polynomial-in-$B$
gauge bound, gives
\[
 |B_p(U,V)-R_p(v)|
 \le C B^a t_f^2+C B^a t_f e^{-\sigma R_*}=o(t_f).
\]
The first term contains the product-logarithm and parallel-transport errors;
the second is the omitted path tail. Here the fixed path-length truncation
$R_*$ is a positive fixed multiple of $B$, so its exponential dominates
all powers of the coupling. The row bound also gives
$|R_p(v)-R_p(m_Y)|\le C\xi t_f$ if all components $gz_i$ of the fluctuation field in its
dependence set obey $|gz_i|\le\xi t_f$.

It remains to bound the affine component at the mean, not just the mean:
\begin{equation}
 |R_p(m_Y)|\le C B^a t_f^2+C B^a t_f e^{-\sigma R_*},
 \label{ch06:mean-target-needed}
\end{equation}
or the corresponding bound after subtracting the prescribed nonzero
linear target. If that target is within the strict plateau margin, choose
$\xi$ so $C\xi$ is less than one third of the margin and increase the
lower bound on the reference coupling to absorb the two errors. The local factors of the two partition-of-unity weights (the weights given to the sets $A_f$ and $A_n$ by the global partition of unity of the first section, with local factors $\vartheta_f(r_p^f)$ or $1-\vartheta_f(r_p^f)$ and $\vartheta_n(n_e)$ or $1-\vartheta_n(n_e)$)
and the normalized weights of the auxiliary index are then equal to one. The same strict
inequalities persist on a parameter neighborhood and supply the
zero-germ assertion used in \eqref{ch06:pointwise-jets}.

There is a concrete comparison identity for the mean behind this estimate. Let $v$
be one candidate field on a common clean ball of radius $8B$, zero on
the tree, with $|v|_\infty\le C B^a t_f$. Suppose every affine component in
that ball satisfies $|\mathsf Dv+b_Y|\le\epsilon$, where
$\epsilon= C B^a t_f^2+C B^a t_f e^{-\sigma R_*}$.
Put $E=\mathsf AP-I$. Then, with inactive identity coordinates included,
\[
 m_Y-v=E^*v-P^*(\mathsf Av+f_Y),\qquad
 \mathsf Av+f_Y=\mathsf D^*(\mathsf Dv+b_Y)+\Pi_{\rm inactive}v.
\]
These are exact identities. For a target at distance greater than $30B$ from the location of $A$,
all the windows in this identity lie in the $8B$ ball about the target,
which contains no plaquette of $A_f$, edge of $A_n$ or site of $A_c$. The inactive term vanishes
there, since all non-tree coordinates are active and $v$ is zero on
the tree. Hence
\[
 |m_Y(i)-v(i)|\le C\epsilon+C e^{-cB}|v|_\infty
 \quad\text{within distance }B\text{ of the target}.
\]
The row bound transfers the same estimate to the affine component at the target.

\begin{proposition}[Cancellation for the affine mean]
\label{ch06:mean-cancellation}
Use the quadratic forms of the two consecutive scales as in
\eqref{ch06:flat-identities}, the representative $W_F$, in the block axial gauge, of the linearization of their critical configuration, and the comparison term
defined by parallel-transporting and truncating the untruncated flat maps $M_f,M_n,N$. Then
\eqref{ch06:mean-target-needed} holds uniformly at targets farther than
$30B$ from the location of $A$. The constants may depend on the fixed step.
\end{proposition}
\begin{proof}
Let $a=W_F h$ be the fine field that the linearization of the critical configuration, in its representative $W_F$ of \eqref{ch05:stationary-tangent}, assigns to a coarse
logarithmic field $h$, and set $w=X_c^{\rm lin}(h)$ and
$p=(X_f^{\rm lin}(a),Z^{\rm lin}(a,h))$. Write
$T w=(M_fw,M_nw)$. Quadratic Gaussian matching gives
$|p|_2^2=|w|_2^2$, with the unnormalized $\ell^2$ norms. The two identities
\eqref{ch06:flat-identities} give $T^*p=w$ and $T^*T+N^*N=I$.
Consequently
\[
 \|p-Tw\|_2^2
 =\|w\|_2^2-2\|w\|_2^2+\|Tw\|_2^2=-\|Nw\|_2^2.
\]
Both sides have opposite signs, so $p=Tw$ and $Nw=0$.
In particular every linear component minus its comparison term vanishes. This argument
also applies to each signed orientation with the image of the spanning trees under the symmetry;
it never assumes that the operator $W_F$ for rooted trees commutes with the symmetry.

Choose one coarse gauge on the clean $8B$ ball and let $h$ be its
logarithmic field, extended by zero outside. The local axial gauge
construction bounds $\|h\|_\infty\le C B^a t_f$.
The operator $W_F$ has a bounded row norm at fixed $L$:
the exponentially decaying kernel of the linearization is summed absolutely,
and the correction bringing it to the block axial gauge is a sum along finitely many tree
edges per block. Thus the relative candidate
$v=W_Fh-\operatorname{rep}h$ is zero on the tree and has
$\|v\|_\infty\le C B^a t_f$.

Compare a nonlinear component in the inner ball with its untruncated linear component.
Taylor's theorem for the four-link curvature and the fixed
branch of the averaging operation bounds the nonlinear remainder of the component by $C\|h\|_\infty^2$.
For a path term, differentiating its products adds powers of the
length, summable against the exponential path bound.
Parallel transports change by at most their length times
$C\|h\|_\infty$; multiplying by an input of the component contributes a further
$O(\|h\|_\infty)$. The same estimates apply to the comparison term, since
the comparison term uses exactly the parallel-transported untruncated maps appearing in $T$.
Its clipping is inactive on this clean link support.
Omitted paths contribute at most
$C B^a t_f e^{-\sigma R_*}$ after decreasing $\sigma$.

Finally the derivative in the fluctuation field is evaluated at $\operatorname{rep}V$,
whereas the parallel-transported flat derivative is evaluated at zero before
parallel transport. Their difference has row norm at most
$C B^a t_f$, by the same second-derivative and parallel-transport estimates.
Applied to $v$, it contributes $C B^{2a}t_f^2$. No equality of these two
matrices is used. The frozen part of the inhomogeneous term is absent on the clean inner ball:
all its non-tree coordinates are active. The exact linear cancellation
therefore gives
$|\mathsf Dv+b_Y|\le C B^{a'}t_f^2+
 C B^{a'}t_f e^{-\sigma R_*}$ simultaneously on every required component,
using the one field $h$. The comparison identity for the mean above proves
\eqref{ch06:mean-target-needed}.
\end{proof}

This conclusion uses the path realization, with its factorization, of the untruncated flat maps in
\eqref{ch06:flat-identities}. A different comparison term with the same numerical
norm bounds would not satisfy it. For event indicators far from the location of $A$ it establishes the
implication as follows. The plaquette part of the fine comparison term is at most
$C L^{-2}t_f$, and the comparison term of an averaging component at most $C t_f$.
Choose $L$ and the plateau width $W$ for the deviations $n_e$ so these bounds lie
strictly inside their respective plateau thresholds, as in the energy
estimate. The preceding proposition, the nonlinear Taylor error, and
the bound on the fluctuation field then leave the fine plaquette distance from the identity below $t_f/4$
and the deviation $n_e$ below $Wt_f$. The corresponding factors of the partition-of-unity weights and the
normalized weights of the auxiliary index there are equal to one. The bond increment is bounded
by $Ct_f$ on this event, so a sufficiently broad fixed chart cutoff is
also one. By convention, a cutoff factor other than a chart cutoff whose location (the projected corners of its plaquette, the endpoints of its edge, or its site) meets $A^{[300B]}$, and a chart cutoff on a link whose projected endpoints meet $A^{[304B]}$, is not written as one plus a cutoff remainder: it is kept whole as a factor of the sector, and its cutoff remainder is set to zero. Section~\ref{ch06:primitive-bounds} below imposes the corresponding conditions through distance $300B$ and retains the chart cutoffs through $304B$.\mgubi{Convention taken from Lean (YM.\allowbreak ConnectedActivityE6.\allowbreak Coat.\allowbreak mandatory, mayor, aura; definitions read by two agents) and existence/\allowbreak gaussian-warnings.\allowbreak tex; v3 did not state it. Not checked: that Lean's factor types are exactly v3's cutoff factors.} Every other cutoff remainder lies outside these neighborhoods and has these small-field neighborhoods.
Replacing $\xi$ by $\xi/2$ gives strict margins on a neighborhood and
hence the zero-germ support condition. This proves the implication
\eqref{ch06:tail-event} for these event indicators, with the fixed comparison term
and plateau choices just stated.

\subsection{Derivatives of products of local factors}
\label{ch06:bounded-jet-route}

The local expansion offers a direct way to combine differentiation
and rare events. It is useful even when no aggregate quadratic majorant
has been established. Let $F(V,z)$ be a product of local factors after the
single change of variables \eqref{ch06:protected-change} above. Suppose, at the base point, its derivatives
through order $k$ satisfy
\begin{equation}
 |D^jF(V_0,z)|\le A_k(x)(1+s)^{m_k}
       \mathbf1_{\mathcal E}(z)\prod_i a_i,\qquad j\le k,
 \quad x=\widehat\beta_f .
 \label{ch06:pointwise-jets}
\end{equation}
Here $\mathcal E$ is the intersection of the events for $r$ distinct
cutoff remainders. This support assertion concerns all derivatives:
each remainder must vanish identically on a parameter neighborhood
outside its event. Assume the local endpoint protections needed for
dominated differentiation in the preceding subsection.

Then the common Gaussian integral, including its density ratio, obeys
\begin{equation}
 \left|D^j\mathbb E_{\Gamma_I(V_0)}
                  [\mathcal R_I(V,z)F(V,z)]_{V=V_0}\right|
 \le \widetilde A_k(x)(1+s)^{\widetilde m_k}
                   x^{-\tau r}\prod_i a_i,\qquad j\le k,
 \label{ch06:actual-jet-integral}
\end{equation}
provided \eqref{ch06:power-tail} holds with exponent $2\tau r$.
The constants are uniform in the frozen variables, which can subsequently
be integrated against their Haar probability measure.

Here is the moment argument in full. If every Gaussian coordinate has
variance at most $M$, the elementary exponential estimate gives
$\mathbb E(1+|z_i|)^q\le2e^{Mq^2/2}$ for each positive integer $q$.
Writing $N=|I|$ and $F_I=1+N+\sum_I|z_i|$, one has
$F_I\le(1+2N)(1+\max_I|z_i|)$ and hence
\[
 \mathbb E F_I^q\le(1+2N)^q(1+2N e^{Mq^2/2}).
\]
This uses no independence. Apply Leibniz' rule,
\eqref{ch06:base-scores}, and \eqref{ch06:pointwise-jets}.
Cauchy--Schwarz gives
$\mathbb E[\mathbf1_{\mathcal E}F_I^q]
 \le\mathbb P(\mathcal E)^{1/2}(\mathbb E F_I^{2q})^{1/2}$.
Since $N\le CB^4s$, all remaining factors are polynomial in $s$ and
subpower in $x$, proving \eqref{ch06:actual-jet-integral}.

For its application, \eqref{ch06:pointwise-jets} must be verified for the
factors of the expansion, with their full dependence sets. Ordinary
exponentiated correction terms satisfy it once their logarithmic
derivatives have small amplitude: the partition formula for derivatives
of $e^h-1$ gives a fixed multiple of that amplitude when $|h|\le1$.
Large-field components use their smooth local extension only when the
entire required coarse neighborhood is small-field. Their physical
indicator is not differentiated. The residual Gaussian density, the exponential factor in \eqref{ch06:gaussian-identity}, is first
decomposed into the protected linear and quadratic correction
factors; it is not copied once for every activity. These are specific
conditions on the exact expansion, not automatic consequences of the
Gaussian gap.

\subsection{One majorant for the product}

Here $z$ has the conditional law $d\gamma_{\Gamma(Y)}$, with frozen links $Y$ integrated in probability Haar measure. The normalization $\mathcal C_Y$ and the residual exponential are both those of \eqref{ch06:gaussian-identity}.

The residual Gaussian density in \eqref{ch06:gaussian-identity}, the
Jacobian ratios, and differentiated factors produce unbounded functions
of $z$. Bounds on their separate integrals cannot simply be multiplied.
The correct estimate is under a single common Gaussian law.

For a decomposition that retains an unbounded exponential instead, the following alternative criterion is available. Suppose a finite product of
local factors has, after up to $k$ allowed derivatives, a majorant of the
form
\begin{equation}
 C_k(\widehat\beta_f)(1+s)^d\,P_k(z)
       \exp\bigl(\eta\cdot z+\tfrac12z^*Kz\bigr)
       \prod_i a_i\prod_w M_w\mathbf1_{\mathcal E_w},
 \label{ch06:majorant-form}
\end{equation}
where $s$ counts base points, $P_k$ is a fixed-degree polynomial in absolute
coordinates with controlled moments, $K\ge0$, and the product over $w$
uses distinct event labels. The numbers $a_i$ bound ordinary local factors;
$M_w$ bound the cutoff factors or the probability weights of the branches. Their number must
include all allowed local values of the auxiliary index. Assume
\begin{equation}
 \|\Gamma^{1/2}K\Gamma^{1/2}\|\le1/8,
 \qquad \tr(\Gamma^{1/2}K\Gamma^{1/2})
              +|\Gamma^{1/2}\eta|^2\le C\sum_i\lambda_i.
 \label{ch06:quadratic-control}
\end{equation}
These are bounds on the sum of all local corrections. Smallness of each
summand separately does not suffice when the number of factors grows.
Exponential kernel bounds and bounded overlap of their full dependence sets
are one way to obtain them. For the linear term, a useful estimate
uses weights $w_{ia}>0$ with $\sup_a\sum_iw_{ia}\le C$; then
$|\sum_i\eta_i(a)|^2\le C\sum_i|\eta_i(a)|^2/w_{ia}$.
It controls the cross terms in the common square.

\begin{proposition}[Common local integral bound]
\label{ch06:common-majorant}
Assume \eqref{ch06:majorant-form}--\eqref{ch06:quadratic-control}, with
uniform polynomial moment bounds for $P_k$ under the corresponding
tilted Gaussians. Assume the local vanishing implication
\eqref{ch06:tail-event} for the same coordinates and event indicators. Then the
absolute differentiated integral is at most
\[
 C'_k(\widehat\beta_f)(1+s)^{d'}
       e^{C'\sum_i\lambda_i}\prod_i a_i
                  \prod_{w\text{ distinct}}(M_w\widehat\beta_f^{-\tau}).
\]
The statement holds after integration of the frozen Haar variables when
all inputs are uniform in those variables. It requires no independence
between the local factors or the large fluctuation-field events.
\end{proposition}
\begin{proof}
The standard Gaussian quadratic-exponential formula, with
$\widetilde K=\Gamma^{1/2}K\Gamma^{1/2}$, is
\[
 \mathbb E_\Gamma e^{\eta\cdot z+z^*Kz/2}
 =\det(I-\widetilde K)^{-1/2}
   \exp\left(\tfrac12\eta^*\Gamma^{1/2}
           (I-\widetilde K)^{-1}\Gamma^{1/2}\eta\right).
\]
Use it also with doubled exponent to control the square of the majorant.
The strict norm inequality in \eqref{ch06:quadratic-control} bounds the tilted
covariance and mean; Gaussian polynomial moments remain controlled.
Cauchy--Schwarz and \eqref{ch06:power-tail} give the asserted bound.
\end{proof}

Finally, Proposition~\ref{ch06:energy} contributes
$e^{-\kappa\beta t_f^2|A|}$. Throughout the coupling band,
$\beta t_f^2\asymp(M_f^{\rm cut})^2$. Any verified
logarithmic losses per element of $A$ of size
$O(B^d(1+\log\widehat\beta_f)^d)+\beta O(B^dt_f^3)$ are
$o(\beta t_f^2)$. They can be included while retaining
$e^{-\kappa'\beta t_f^2|A|}$ for any fixed $0<\kappa'<\kappa$ once the
lower bound on the reference coupling is increased. This last comparison concerns losses
actually assigned per element of $A$. It cannot absorb an unexplained
volume-dependent normalization.

\section{Weighted bounds for the elementary local factors}
\label{ch06:primitive-bounds}

We now retain the individual correction terms, including repeated paths,
instead of estimating only their sum as an operator. Fix a derivative
order $k$, a support exponent $\alpha>0$, and a multiplier $h$ such that the $hB$-neighborhoods are
large enough to contain every window and protected variation used above.
For a support $S$, let $c_B(S)$ be the least number of radius-$B$
balls whose union is connected and contains $S$. For each factor
$\lambda$, let $S_\lambda$ contain all sites on which it depends.
We bound its nonnegative amplitude $a_\lambda$ by requiring
\[
 \sup_x\sum_{\lambda:x\in S_\lambda}
               e^{\alpha c_B(S_\lambda)}a_\lambda\le\epsilon.
\]
The statement concerns the positive derivative orders available in the
incoming path bounds. Additional fixed polynomial factors in $B$ or
$M_f^{\rm cut}$ are subpower in $x=\widehat\beta_f$.

\subsection{Gaussian correction series}

Put $D=\mathsf A\Gamma-I$. The window estimates give exponentially small
weighted entry and fixed derivative envelopes for $D$. For the
interpolating precision $\mathsf A_s=I+s(\mathsf A-I)$, $0\le s\le1$,
the same argument applies uniformly: its gap is at least
$\min(1,c_L)$ and its weighted bounds are uniform.
Let $P_s$ be its localized approximate inverse and $E_s=\mathsf A_sP_s-I$.
The exact inverse and trace expansions are
\[
 \mathsf A_s^{-1}=P_s\sum_{n\ge0}(-E_s)^n,\qquad
 \Gamma^{-1}-\mathsf A=\sum_{n\ge0}(-D)^{n+1}\mathsf A .
\]
All series converge in weighted matrix norm. Integrating
$\tr((\mathsf A-I)\mathsf A_s^{-1})$ gives $\log\det\mathsf A$.
The $n=0$ term is its local diagonal normalization; the remaining terms of the series
are $(\mathsf A-I)P_s(-E_s)^{n+1}$.
The second determinant series comes from
$\log\det(I+D)=\log\det\mathsf A+\log\det\Gamma$.
This determinant identity is legitimate even if $D$ is not symmetric:
$I+D=\mathsf A\Gamma$ is similar to a positive matrix, and the
convergent real logarithm series agrees with its positive determinant.

Each term is indexed by its ordered sequence of bonds, with the intermediate Lie algebra indices contracted in each matrix product. For each fixed derivative order $k$,
the corresponding derivative of the amplitude is bounded by
\[
 C_k(B)e^{-cB}\bigl(C_k(B)e^{-cB}\bigr)^n,
\]
with possibly a smaller positive $c$. The endpoint cutoffs in a quadratic term
bound each Gaussian endpoint and its normalized derivatives by
a subpower function. A remaining linear term has one such endpoint.
Inactive endpoints give a zero germ because both precision and
covariance have identity inactive blocks. Thus all these estimates hold
for the protected factors on the entire sample space.

For completeness, the sum over the supports of the terms is also small.
Consecutive bonds have distance at most $3B+2$.
There are at most $K=C_L(1+B)^4$ possible next bonds.
A term of length $n+u$, for fixed $u$, has support cover at most
$(n+u+1)s_h$, where $s_h$ depends only on the multiplier $h$.
If that support contains a prescribed site, its starting bond has
at most $C_L(1+B)^4(n+u+1)^4$ possible locations.
Fixed derivatives contribute an additional polynomial factor in $n+u$.
Absorb these polynomials by $C_k2^{(k+4)n}$. The containing-site sum
is consequently bounded by
\[
 C_k(B)e^{-cB}\sum_{n\ge0}
 \bigl(C_k(B)e^{-cB}2^{k+4}
                e^{\alpha s_h}K\bigr)^{n}.
\]
Its ratio is at most $1/2$ eventually. The sum is smaller than
$x^{-A}$ for every prescribed $A>0$. This proves smallness of the
individual determinant and inverse-correction terms with their full
supports, including repeated bonds, rather than merely a small
operator norm.

\subsection{Short local and inherited terms}

The remaining estimates for the elementary local factors can be displayed together. Each bound
below is a containing-site bound with full fixed $hB$-neighborhoods and fixed-order
normalized derivatives, on the coarse small-field neighborhoods needed
for those derivatives:
\[
\begin{array}{l|l}
 \text{family}&\text{bound}\\ \hline
 \text{nonlinear action remainders far from }A
       & C_k B^{a_k}\beta t_f^3\\
 \text{extracted mean squares far from }A
       & C_k B^{a_k}\beta t_f^4\\
 \text{local determinant differences}
       & C_k B^{a_k}t_f\\
 \text{local Haar logarithms}
       & C_k B^{a_k}t_f^2\\
 \text{linear Gaussian residual and correction terms}
       & x^{-A}\quad\text{for every fixed }A\\
 \text{inherited localized terms}
       & C_k B^{a_k}\delta_f\\
 \text{inherited $\beta$-independent quadratic terms}
       & C_k B^{a_k}t_f^2\Lambda_f\\
 \text{inherited cubic and quartic terms}
       & C_k B^{a_k}\beta(t_f^3+t_f^4)\\
 \text{omitted path tails}
       & C_k B^{a_k}\beta t_f^2 e^{-cR_*}.
\end{array}
\]
Here $\delta_f$ is the incoming bound on the localized terms and $\Lambda_f$
the bound on the coefficients of the $\beta$-independent quadratic term. For example
$\delta_f\le x^{-1/2}$ and $\Lambda_f\le x^{1/8}$ make every displayed
family small by a fixed negative power of $x$.
The factors from differentiation and counting can be absorbed by
reducing the decay exponent if necessary.

The first two lines use the affine square with the exact
derivative at the background field. Taylor's remainder is quadratic in the displacement;
multiplying by the base affine component, of size $O(t_f)$ on the protected
chart, gives $O(t_f^3)$. The square of the remainder is smaller.
The extracted mean square is $O(B^a t_f^4)$ by
Proposition~\ref{ch06:mean-cancellation}. For derivatives the directions
are the directions of size $O(t_f)$ used above; applying the
same finite Taylor formula with parameter derivatives gives the same
powers, with the fixed factors from the path bounds. No derivative of a large-field indicator is taken.

For the third line, compare non-flat and flat matrices in the common
small gauge on the local window. Their difference is $O(B^a t_f)$.
The resolvent identity for each principal compression transfers this
to the diagonal integral defining the local logarithmic normalization.
The window size and trace introduce only powers of $B$.
The fourth line follows from $\log J(u)=O(|u|^2)$ for the Haar density
in exponential coordinates and from the protected chart bound.
Every site count in an $hB$-neighborhood is polynomial in $B$.

For inherited terms there is a support issue in addition to these
pointwise powers. If an incoming support $X$ has fine connected covering number
$d=c_{B_f}(X)$, with $B_f\le2B$, then its projected $hB$-neighborhood satisfies
\[
 c_B(S)\le P_h+\frac{10P_h}{L}d,\qquad
 x_0\in S\Longrightarrow
 d(x_0,\pi o_X)\le(h+20)B(1+d),
\]
for a fixed $P_h$. To prove the first bound, traverse a spanning tree
of the fine covering centers, sample that walk every $LB$ steps, and
cover each projected $hB$-neighborhood by at most $P_h$ coarse balls.
The walk has length at most $10Bd$, giving the displayed coefficient.
The same walk proves the bound on the distance to the base point. Thus an incoming
base-point weight $e^{\sigma d}$ absorbs the coarse support weight and all
polynomial multiplicities provided
\[
 \alpha\,10P_h/L+4\epsilon<\sigma
\]
for some fixed $\epsilon>0$. Counting possible base points contributes
$C_L B^4(1+d)^4$; derivative factors contribute another fixed power of
$1+d$, absorbed by $e^{\epsilon d}$. Summation of the remaining
geometric series proves the inherited site bounds.
For path terms the analogous cost is their true length, including
winding and repetitions; the unused exponential path factor absorbs its
support and its omitted tail. This explains why a positive incoming
additional support decay, not merely an unweighted coefficient sum, is required.

\subsection{Components and assembly under one law}

All constituent factors keep their original full dependence sets, including cutoff and path conditions. Their exclusion supports, defined in Definition~\ref{ch04:activities}, determine the disjoint covering families, not independence; the parameter $\sigma$ in the weighted support sums is the exponential locality weight of \eqref{ch04:weighted-sums}.

The relevant energy is the remainder after Gaussian extraction,
not the local energy $\mathcal E_A$ alone. We specify the assignment to prevent a
positive affine square from being lost in this estimate.
Let $O$ be a component of $A$ with $n$ elements. Its assigned action
terms have base points within $200B$ of $O$; their \emph{core} consists of the
terms within $2B$. The remaining terms form an \emph{annulus} $Q$.
Form a graph on the subsets of $A$ by joining two sets whenever their
closed $500B$ neighborhoods intersect. Group the sets by the connected
components of this graph. Distinct components of this graph then have
disjoint $500B$ neighborhoods (separation greater than $1000B$ suffices),
so their $200B$ action and $304B$ Haar assignments cannot overlap.
The enlarged finite-range adjacency changes only the fixed polynomial
counting constants. This geometric disjointness does not assert
probabilistic independence.

For this component estimate impose the weak conditions of the partition of unity
through distance $300B$: every fine plaquette not in $A_f$ has distance from the identity at
most $t_f/2$, every deviation $n_e$ with $e\notin A_n$ is at most $2Wt_f$,
the prescribed coarse large-field indicators equal those of the configuration, and no coarse
large-field site occurs away from the sites of $A_c$. The retained values of the auxiliary index for the local gauge transformations
have positive probability weight. These are conditions on the factor itself;
they are stronger than the initial $5B$ consistency definition.
For an annular base point at distance greater than $2B$ from the location of $A$,
its required $B$ neighborhood and the smaller comparison patches
are consequently free of elements of $A$. All such reads fit inside
the retained $300B$ region. Smooth chart cutoff factors are retained
throughout $304B$.
Local determinant terms are assigned through $200B$, whereas Haar
terms are assigned through $304B$; these two sets are not identified.

Let $H$ be the fine curvature and averaging components extracted from the
annulus by the link-support rule of Section~\ref{ch06:sector}. The link support of a retained component has radius
at most $B/50$. A retained component meets an active coordinate outside
$A^{[10B]}$, so its entire link support lies outside
$A^{[(10-2/50)B]}$ and hence outside the $2B$ core.
We first work in one signed orientation, before its probability weight
is applied. Denote the derivative and the inhomogeneous term of its components by
$\mathsf D_\sigma,b_{Y,\sigma}$, including the plaquette and penalty
coefficients but excluding the square root of the orientation weight.
The global matrix in \eqref{ch06:actual-matrix} stacks
$\sqrt{\omega_\sigma}\mathsf D_\sigma$.
Suppress $\sigma$ temporarily and set
\[
 X_H(\widetilde u)=\tfrac12
       \sum_{j\in H}|(\mathsf D_\sigma\widetilde u_{\rm active}+b_{Y,\sigma})_j|^2.
\]
The orientation probability is not included in this half-square; it
will appear once in the average below. The coefficients of the components and
the frozen part of the displacement field $\widetilde u$ are included once: equivalently apply the all-coordinate derivative to
the full displacement field and add the value of the components at $\operatorname{rep}V$. The \emph{displacement field} $\widetilde u$
is the active increment, the clipped frozen logarithm on frozen
non-tree bonds, and zero on tree bonds. Where the retained cutoff factors are nonzero
it reconstructs the original fine field on every retained link support.
Indeed the clean annular fine plaquettes and deviations $n_e$,
together with the common-tree comparison, put each frozen mismatch
within its logarithm cutoff's inner plateau. On active bonds the
cutoff support gives $|\widetilde u|\le2w t_f$; take the fixed width $w$
larger than the frozen comparison constant. Both the tree field and $\operatorname{rep}V$
are identity on tree bonds. Thus this same displacement field, uniformly bounded
by $2w t_f$, may be used in the affine approximation; an unrelated
logarithm or an unrestricted Gaussian sample cannot be substituted.

Write $W(Q)$ for the sum of the six local energy terms
\eqref{ch06:six-terms} on the annulus, with the same fine clip threshold
as the core. Then
\begin{equation}
 W(Q)-X_H(\widetilde u)\ge-\epsilon(x)n,\qquad
 \epsilon(x)\le C B^a t_f^3+C B^a t_f^2e^{-cR_*}
                         =o(t_f^2).
 \label{ch06:annulus-extraction}
\end{equation}
To verify the inequality, work in one clean local gauge about each
annular action term. Its nonlinear component minus its comparison term has norm
$M\le C B^a t_f$ and differs from its affine approximation
by $E\le C B^a t_f^2+C B^a t_f e^{-cR_*}$.
Taylor expansion gives this for the derivative at $\operatorname{rep}V$
directly; comparing first to parallel-transported flat components adds an error of
the same order. Hence the difference of the two half-squares is at
most $E(2M+E)/2$. The energies $E_f$ and $E_n$ equal the squares of their
nonlinear components minus comparison terms, plus the nonnegative difference between the squared distance from the identity and
the squared curvature representative. If a component is not retained, no affine square is subtracted,
so only its nonnegative square remains.
The coarse counterterm contributes $O(t_f^4)$.
The linear term of the complement component vanishes by the cancellation in
Proposition~\ref{ch06:mean-cancellation}; its nonlinear square is
$O(B^a t_f^4)$. The two signed mismatch terms are a comparison term of size
$O(t_f)$ paired with the quadratic Taylor/parallel-transport errors, hence
$O(B^a t_f^3)+O(B^a t_f^2e^{-cR_*})$.
There are at most $C B^4n$ such terms, proving
\eqref{ch06:annulus-extraction}.
The clean annular fields are below the fine threshold, so the fine
clip is its quadratic branch. The retained indicators of the relaxed regularity condition also equal
one: within their nontrivial range the patch is clean, and beyond
that range their permitted bound exceeds the uniform bound on the distance from the identity in the group. This checks the formula for the components used in the argument.

For precision, let $\sigma$ range over the signed lattice orientations
with their probability weights $\omega_\sigma$, transforming the field,
spanning trees, auxiliary index, index triple, and components together. Every term inside the
following average is the unweighted single-orientation quantity just
defined. Thus each orientation probability is counted exactly once. The energy assigned to this component
after Gaussian extraction is exactly
\begin{equation}
 \mathcal E_O^{\rm rem}
 =\operatorname{Av}_{\sigma}\bigl(
       \mathcal E_{\sigma O}+W(Q_\sigma)
       -X_{H_\sigma}(\widetilde u_\sigma)+C_\sigma(m_Y)\bigr),
 \qquad C_\sigma(m_Y)\ge0 .
 \label{ch06:component-extracted-energy}
\end{equation}
This is a finite identity between action terms: every near action term is assigned once;
a retained fine curvature or averaging component subtracts its affine square;
its extracted mean square is retained with positive sign.
Components that are not retained subtract zero. The mean-square sum is $C_\sigma$.
The near action terms partition into the core and annulus, so
\eqref{ch06:component-extracted-energy} and
\eqref{ch06:annulus-extraction}, followed by Proposition~\ref{ch06:energy},
give
\[
 \mathcal E_O^{\rm rem}\ge(\kappa t_f^2-\epsilon(x))n
                         \ge\tfrac12\kappa t_f^2 n
\]
at large coupling. This is the needed bound on the extracted energy.
Dropping $X_H$ before proving the annular estimate would be invalid.

The remaining local normalization, Haar, auxiliary-index, and frozen-coordinate
contributions are bounded by $C B^a(1+\log x)^a n=o(\beta t_f^2)n$.
The component amplitude consequently satisfies
\[
 a_{\rm comp}\le\exp\bigl(-\kappa'\beta t_f^2n\bigr)
\]
for a fixed positive $\kappa'$ smaller than the retained energy
coefficient. Fixed derivatives of its smooth extension on a fully
small-field coarse neighborhood add polynomial factors per component
and are absorbed by a further decrease of $\kappa'$. The physical Borel
indicator is held fixed. This assertion is where the retained cutoff factors are nonzero,
under the precise core/annulus assignment just specified; where a
cutoff factor kept inside the component vanishes, the full component factor is zero.
Connected index triples with prescribed root have at most
$\bigl(C(1+B)^a\bigr)^{n}$ possibilities: list them by a rooted spanning tree,
using their finite-range adjacency. This entropy and
$e^{\alpha c_B(S)}\le e^{C\alpha n}$ are absorbed by the same
remaining positive energy bound. A component's physical Borel indicator is held fixed;
the derivative assertion concerns its smooth extension on the stated
small-field coarse neighborhood.

Finally replace each small logarithm $f_\lambda$ by
$e^{f_\lambda}-1$. For $|f_\lambda|\le1$, the exponential partition
formula preserves its amplitude up to a constant depending on $k$.
The previous bounds furnish the pointwise product jets
\eqref{ch06:pointwise-jets}; derivative factors polynomial in the total
support size are absorbed by a small further decrease of the support exponent $\alpha$.
Cutoff remainders carry the joint-event factors proved above.
Equation~\eqref{ch06:actual-jet-integral} then gives a product majorant
under the one conditional Gaussian law, followed by the frozen
Haar probability integral. This is the local integration estimate
required before the connected expansion. Incoming large-field
activities retain their given amplitudes and participate in this same
product bound; they are not claimed to become pointwise small by
Gaussian integration.

\subsection{Comparison with the linear local integral}
\label{ch06:linear-comparison}

The estimates also give the small parameter needed to compare a local
nonlinear integral with its linear Gaussian counterpart. Fix the same
spanning trees, coordinate set, windows, and clean coarse field throughout the
comparison. Write $\Gamma_{\rm lin}$ for the window covariance formed
from the corresponding flat components, and $v$ for the
linear relative field $W_Fh-\operatorname{rep}h$ of the proof of Proposition~\ref{ch06:mean-cancellation} on the common clean neighborhood.
Taylor expansion of the nonlinear components and the principal resolvent identity
give, on those windows,
\[
 \|\Gamma-\Gamma_{\rm lin}\|\le C B^a t_f,\qquad
 |m_Y-v|_\infty\le C B^a t_f^2+C B^a t_f e^{-cB}.
\]
The second estimate is the comparison for the mean proved above; the first
is valid in weighted row norms, or with a harmless local polynomial
factor in an unweighted matrix norm. Both have fixed-order versions
under the same available path derivative bounds. In particular the
mean shift in Gaussian units satisfies
\[
 g^{-1}|m_Y-v|_\infty
       \le C B^a g(M_f^{\rm cut})^2+
                         C B^a M_f^{\rm cut}e^{-cB}.
\]
Thus it is of size $x^{-1/2+o(1)}$, rather than merely bounded.

The physical coordinate comparison uses the same block axial chart.
Its nonlinear map and its first-order Taylor map differ by
\[
 C B^a(t_f+g|z|)^2
       \le C'B^a(t_f^2+g^2|z|^2)
\]
on their protected common chart. Dividing by the fine small-field
radius gives $C B^a(t_f+(g^2/t_f)|z|^2)$.
All its fixed Gaussian moments are $x^{-1/2+o(1)}$ on a region
with polynomial-in-$B$ dimension. This is a Taylor estimate with
polynomial moment control, not a uniform bound on unbounded samples;
the chart-complement events are controlled separately.

For an interpolation of the positive covariances and of the centered
means, differentiation of its common Gaussian density ratio therefore
has local polynomial moments bounded by $x^{-1/2+o(1)}$, with the
appropriate polynomial factor in the size of the dependence set.
Indeed the logarithmic derivative is the trace, quadratic, and linear
score displayed after \eqref{ch06:density-ratio}; each coefficient
contains one of the covariance or normalized mean differences just
estimated. Integrating this derivative compares the two laws.
For an inherited activity with only $k$ available derivatives, this
comparison must be implemented by holding its physical fine links fixed
with the change of variables \eqref{ch06:protected-change}. Differentiate the common
density, Jacobian, smooth cutoff factors, and at most the available
$k$ activity derivatives. Taylor expansion of
$D^k E(U_{\rm nonlinear})-D^k E(U_{\rm linear})$ would require an
unavailable $(k+1)$st derivative and is not justified here.
The estimate above for the change of variables $T_V$ supplies bounds on
$T_V$ and its density scores; it is not permission to lose an extra
derivative of the incoming activity.
This comparison applies on a local region where all the windows lie in
the small-field region. It does not compare the Gaussian law of a sector whose index triple has nonempty $A_f$ or $A_n$ with the law of the sector with $A_f=A_n=\emptyset$,
does not replace $W_F$ by the representative in a different gauge,
and does not identify separately chosen event indicators. Such replacements need
their own tail or support decomposition.

\section{What must be checked for a concrete step}

The local integral estimate now follows from the gap of the non-flat matrix $\mathsf A$, the clipped
inhomogeneous term, cancellation of the comparison terms under the factorization hypothesis, strict margins in the event inequalities, and the
bounds for the elementary local factors above. Its hypotheses are the compatible untruncated
quadratic identities of the two consecutive scales, the rooted averaging operation, exponential path bounds
through the prescribed derivative order, and the incoming activity bounds
with a strictly larger support exponent than the one used in the conclusion. These are incoming data, not estimates
that can be created by choosing the schedule of Chapter~\ref{ch10:trajectory}. All matrices,
comparison terms, spanning trees, values of the auxiliary index, and frozen inhomogeneous terms in the argument belong to
that same pair of scales.

The chapter controls local amplitudes and their integrals. Reconstruction
of the exact block density additionally requires assigning every
nonlinear component, determinant contribution, Haar factor, and event indicator exactly
once. In particular the two near sets for local determinant terms and
Haar terms need not coincide, and the flat determinant scalar is
extracted only once for the whole family of spanning trees. Those identities, and the
subsequent connected regrouping, are distinct from the analytic bounds
proved here.
}{}
\setcounter{chapter}{6}
\IfFileExists{chapters/07-matching.tex}{% !TEX root = ../monograph.tex
\chapter[Local matching and the one-loop coefficient]{Local matching and the one-loop coefficient}\label{ch07:matching}

A renormalization step has two distinct normalization tasks. Its
classical action at the critical configuration must retain the prescribed quadratic form
and agree with a gauge-invariant local action through degree four.
Its Gaussian determinant changes the coefficient of that quadratic
form. The first task is a local Taylor construction; the second requires
a comparison of determinants uniform in the block ratio. This chapter
separates the two, so that a perturbative coefficient does not stand in
for an estimate on the block integral.

Fix one step $m+1\to m$, an odd block ratio $L=L_m$, and the Lie
algebra metric of Chapter~1. Momentum is measured in lattice units at
the indicated scale. The adjoint Casimir $C_G$ is defined by
$\tr_{\gLie}(\operatorname{ad}X\operatorname{ad}Y)
=-C_G\langle X,Y\rangle$; set
\[
 \kappa_G=\frac{11C_G}{24\pi^2}.
\]
The group is compact with simple nonabelian Lie algebra. Constants may
depend on fixed stage bounds, but constants described as uniform in $L$
are chosen before $L$. All local Taylor expansions below concern the
analytic action near the flat field, where the indicators of the regularity condition are
constant. They do not differentiate at the boundary of the small-field region.

\section{Infinite-lattice Taylor kernels and Ward identities}\label{ch07:kernels}

Write $C_n(E)=D^nE(0)/n!$ for the normalized symmetric Taylor
coefficient of an action $E$. A translation-invariant coefficient is
specified by its values on plane waves. After division by the number
of lattice sites, its symbol is a tensor
$f(p_1,\ldots,p_{n-1})$, with $p_0=-\sum_{i>0}p_i$.
The tensor has one spacetime index and one argument in $\gLie$ for each field.
The symbol is periodic with period $2\pi$ in each independent momentum.
An \emph{infinite-lattice symbol} means one function chosen on the infinite lattice whose
restriction to each finite momentum grid describes the corresponding
torus coefficient. This requirement is stronger than a separate smooth
extension of each finite grid.

For $\sigma>0$, write $|f|_\sigma\le A$ when every component is
holomorphic and bounded by $A$ on the strip
$|\operatorname{Im}p_i^a|<\sigma$. There are only finitely many direction
and Lie algebra components in the basis of Section~\ref{ch01:gauge} at fixed $n$. Reducing the strip gives an
exponentially weighted sum of Fourier coefficients, with a constant
depending on $n,\sigma,G$. Thus strip bounds for infinite-lattice symbols imply locality
uniformly in the period. We say a symbol has a zero of order at least $r$
if all its Taylor coefficients of total degree below $r$ vanish at zero
external momenta.

Let $q_a(p)=e^{ip_a}-1$ and $d_1(p)v=q(p)\wedge v$ be the lattice
curl. Suppose an analytic gauge-invariant difference $E$ has vanishing
Taylor coefficients below degree $n$. Its infinitesimal gauge identity
at that degree contains only the linear gauge variation $d_0\omega$:
the nonlinear part multiplies lower coefficients, which are zero.
Consequently its first nonzero symbol satisfies
\begin{equation}\label{ch07:ward}
 \delta_i f=0\quad(0\le i<n),\qquad
 \delta_i=\iota_{q(p_i)}.
\end{equation}
Contraction is algebraic, without complex conjugation; it acts only on
the direction index of the indicated \emph{leg} (field argument of the symbol). Reality means
$\overline{f(-p)}=f(p)$ on real momenta. Constant gauge transformations
make the tensor invariant under the diagonal adjoint action. A Ward identity verified on
all finite momentum grids passes to the continuous infinite-lattice symbol by
density. This last step requires an infinite-lattice kernel.

\section{Division by the lattice curl}\label{ch07:division}

The algebraic Poincar\'e lemma explains why gauge-invariant leading
coefficients are built from curvature; see
\cite[Section~4]{BarnichBrandtHenneaux2000}. Its local continuum form does
not give the periodic strip bounds needed here. We give the quantitative
adaptation, including the constraint on the sum of momenta.

\begin{lemma}[Quantitative curvature representation]\label{ch07:curl}
Let $f$ have $n$ one-form legs, satisfy \eqref{ch07:ward}, and obey
$|f|_\sigma\le A$. For $n\ge3$ there are $\eta>0,K<\infty$,
depending only on $n,\sigma$, and a periodic holomorphic tensor $g$
with two-form legs such that
\[
 f=\delta_0\cdots\delta_{n-1}g,\qquad |g|_\eta\le KA.
\]
For $n=2$ the same assertion holds if $f$ has order $r+2$, $r\ge0$;
then $g$ can be chosen with order $r$. Reality and adjoint-invariant
coefficients can be imposed with a fixed additional norm factor.
\end{lemma}

\begin{proof}
The elementary division is
\[
 D_af(p)=\frac{f(p)-f(p|_{p_a=0})}{e^{ip_a}-1}.
\]
The expression $p|_{p_a=0}$ means that only the $a$th coordinate is set
to zero. Periodicity removes the apparent singularities at $p_a\in2\pi\ZZ$.
For the Fourier monomial $e^{iz\cdot p}$ the quotient is a finite
geometric sum with $|z_a|$ terms and no frequency larger
than $|z_a|$. A smaller strip absorbs this polynomial factor into the
exponential Fourier decay. Hence $D_a$ is bounded from a fixed strip to
a smaller one and loses at most one Taylor order.

For a cycle $\delta f=0$ in one exterior factor, write
$f=e_a\wedge u+v$, with $u,v$ containing no $e_a$. Write $\delta'$ for contraction with $q$ with its $a$th component omitted.
Then $\delta'v=-q_au$, and direct substitution gives
\[
 f=\delta(e_a\wedge D_av)+v|_{p_a=0}.
\]
Repeat over the four axes. At positive exterior degree the last
remainder has no possible component and is zero. At degree zero the
last remainder is $f(0)$, so that case is solvable precisely when
$f(0)=0$. The construction preserves auxiliary constraints independent
of the momentum being divided. It raises the exterior degree by one,
loses at most one Taylor order, and has a bounded smaller-strip norm.

Only the dependent momentum $p_0$ needs extra care. The following
induction gives it. With $t$ independent momenta and exterior degrees
$(d,1,\ldots,1)$, a cycle in every factor has a zeroth-factor primitive
preserving the other cycles whenever $d+2t>4$.
For $t=0$ its degree exceeds four and it is zero. For $t>0$, replace
$p_1$ by $p_0$ as independent variable and apply the preceding division.
This invertible integral shear only reduces the strip by a fixed factor.
It gives $\delta_0X=f$, preserving the cycles in factors $2,\ldots,t$.
Apply another zeroth-factor division to $\delta_1X$, obtaining
$\delta_0Y=\delta_1X$. At $p_1=0$ the tensor $Y_0$ is a cycle with
zeroth degree $d+2$ and $t-1$ independent momenta. Induction gives
$\delta_0Z=Y_0$. Extend $Z$ independently of $p_1$ and set
$Y'=Y-\delta_0Z$. It vanishes at $p_1=0$, so scalar division in factor
one gives $\delta_1D=Y'$. The corrected primitive
$X-\delta_0D$ has contraction $f$ in factor zero and zero in factor one.
All other cycles were retained. The invariant quantity $d+2t$ makes
this a terminating induction.

Initially $d=1,t=n-1$, so the induction applies for $n\ge3$.
After the zeroth-factor primitive, divide the independent legs successively.
For $n=2$, start instead with an order-$r+2$ tensor. The first primitive has
order $r+1$; contraction in factor one restores one order, and the next
primitive $Y$ has order $r+1$. Thus $Y(0)=0$ without induction. Its scalar
primitive has order $r$, the corrected primitive has order $r+1$, and the
last division leaves order $r$. All operations are bounded on a common
smaller strip and their number is fixed.

Finally replace $g$ by $(g(p)+\overline{g(-p)})/2$ and project each
coefficient tensor on $\gLie$ onto the invariant real tensors. Contraction acts
on direction indices and commutes with this finite-dimensional
projection. These operations preserve the reconstruction and all
retained orders, with a fixed norm factor.
\end{proof}

Choose coordinate contours from the origin to the Fourier displacements
$z_1,\ldots,z_{n-1}$. The transported curvature along such a
contour is
\begin{equation}\label{ch07:slot}
 F_{\gamma,ab}(U,x)=\Ad(U_{x+\gamma})
 \operatorname{proj}_{\gLie}\frac{U_{x+z;ab}-U_{x+z;ab}^{-1}}2,
 \qquad \gamma:0\longrightarrow z.
\end{equation}
All factors transform at the same base point $x$, so an invariant
multilinear form on $\gLie$ is gauge invariant. Its length is the
sum of the contour lengths plus four per plaquette. The linearization of
each factor is the corresponding curl. The Fourier coefficients of
$g$ therefore produce an exact leading Taylor coefficient $f$.
Symmetrizing the legs implements the convention $C_n=D^n/n!$;
it introduces no extra factorial.

If $a$ denotes such a multilinear term and $|a|$ the sum of absolute
components in the basis of Section~\ref{ch01:gauge}, the result is
\begin{equation}\label{ch07:path-norm}
 \sum_a |a|e^{w\operatorname{length}(a)}\le K_w A
 \qquad(0<w<\eta).
\end{equation}
Averaging over signed coordinate symmetries preserves matching and this
bound. Absolute convergence permits periodization on every torus.
For the two-leg case with $r=2$, the curvature coefficients $c(z)$ obey
\begin{equation}\label{ch07:moments}
 \sum_zc(z)=0,\qquad \sum_z z_ac(z)=0\quad(1\le a\le4).
\end{equation}
These are the value and first derivatives of its symbol at zero.
They hold separately for each orientation and each pair of Lie algebra arguments.

\section{The marginal quadratic coefficient}\label{ch07:quadratic}

A quadratic Taylor symbol includes the factor $1/2!$. For its transverse
entry, define
\begin{equation}\label{ch07:alpha}
 \alpha(f)=2\operatorname{Re}[t^2]f_{22}^{uu}(te_1),
\end{equation}
where $[t^2]$ extracts the coefficient of $t^2$ at $t=0$, and the
superscript $uu$ means that the same unit vector $u\in\gLie$ is inserted in
both arguments. The quadratic homogeneous term of the kinetic symbol $b_H$
of $J$, evaluated on $v,w\in\gLie$, is
\[
 (b_H)_{ab}^{vw,[2]}(p)=\tfrac12
 (|p|^2\delta_{ab}-p_ap_b)\langle v,w\rangle_{\gLie},
\]
so $\alpha(b_H)=1$. The factor two in \eqref{ch07:alpha} is therefore
fixed by the Taylor convention, not by a choice of coupling scheme.

\begin{lemma}[Removal of the kinetic coefficient]\label{ch07:unique}
An analytic quadratic symbol from an invariant real action, with
exchange symmetry $f_{ab}^{uv}(p)=f_{ba}^{vu}(-p)$ and signed-coordinate symmetry, has
$f-\alpha(f)b_H=O(|p|^4)$. Its coefficient $\alpha$ is independent
of the transverse axis pair and the unit vector $u$.
\end{lemma}

\begin{proof}
The Ward identity first annihilates the constant term. Its linear term is
also zero: inversion of all coordinates acts positively on the two
one-form indices and negatively on a linear momentum polynomial.
Edge translation phases have constant term one; their higher terms
multiply the already vanishing constant coefficient.

At quadratic order the symmetric diagonal forms on $\gLie$ are scalar
multiples of the invariant metric. A self-adjoint operator commuting
with the adjoint action has invariant eigenspaces, which are ideals;
simplicity leaves only one eigenvalue. Coordinate reflections allow
only $p_c^2$ in a diagonal direction entry and $p_ap_b$ in an
$a,b$ off-diagonal entry. Permutation symmetry identifies their
coefficients. The transverse Ward equations then force the tensor to
be a scalar multiple of $|p|^2I-pp^{\mathsf T}$.
One transverse evaluation determines that scalar. Subtracting the
kinetic symbol $b_H$ itself preserves all symmetries, not merely their leading
parts. Its residual cubic coefficient vanishes by inversion, since all
lower coefficients now vanish. Thus the residual has order four.
\end{proof}

Apply the two-leg curvature representation to this residual. The
resulting quadratic curvature terms have the zero moments
\eqref{ch07:moments}. They cannot change the coupling: their curvature
symbol is $O(|p|^2)$ and their link symbol is $O(|p|^4)$.
The coefficient of curvature square is carried entirely by $\beta$.
This statement refers to Taylor coefficients at the flat field; it does
not identify a nonlinear effective action with its quadratic truncation.

\section{Haar measure and the gauge-fixing term}\label{ch07:completion}

The Gaussian contribution at a critical configuration $U_v$ of the fine field is
\begin{equation}\label{ch07:thermal}
 T(v)=\mathcal L_{m+1}(U_v)+\tfrac12\log\det Q(v)-\log J_{\rm Haar}(v).
\end{equation}
Here $Q$ is the fluctuation Hessian without its factor $\beta$,
$J_{\rm Haar}$ is the density of Haar measure in the chosen chart, and
$\mathcal L_{m+1}$ is the $\beta$-independent quadratic term of the path action, called the \emph{inherited loop action}.
Field-independent constants may be subtracted throughout. This does
not alter any Taylor coefficient of positive degree.

The following finite-dimensional identity accounts for the ghost
factor. Let $\mathsf J$ be the joint Hessian in left tangent coordinates,
$\mathcal R$ the infinitesimal gauge action, and $\mathcal I$ a
slice embedding. Assume
\[
 \mathsf J\mathcal R=0,\quad |\det(\mathcal R,\mathcal I)|=1,
 \quad Q=(Ds)^*\mathcal I^*\mathsf J\mathcal I\,Ds,
 \quad J_{\rm Haar}=|\det Ds|.
\]
These are identities at the critical configuration. Given a gauge
condition $F$ with invertible $G=F\mathcal R$, put
$C=\mathsf J+F^*F$. Block triangular factorization in the basis
$(\mathcal R,\mathcal I)$ gives
\[
 \det C=(\det G)^2\det(\mathcal I^*\mathsf J\mathcal I),
 \qquad
 \tfrac12\log\det Q-\log J_{\rm Haar}
 =\tfrac12\log\det C-\log|\det G|.
\]
The fluctuation form is positive on the slice. With a positive regulated
ghost operator the absolute value can be dropped. This is the
finite-dimensional Faddeev--Popov determinant identity; the moving
Haar Jacobian is part of it. Omitting that Jacobian would change the
quantity being compared. Background-field perturbation theory uses
exactly this vector-minus-ghost combination~\cite{Abbott1982}.

For the lattice comparison choose $\mu=L^{-1}$ and
\[
 P_\mu=\mu^2(\Delta_U+\mu^2I)^{-1},\qquad
 m_\mu=\mu^2P_\mu^9,\qquad \Delta_U=d_U^*d_U.
\]
If $X\mathcal R=I$ is the derivative of the local gauge transformation, take
$F=d_U^*P_f+m_\mu X$; then $G=\Delta_U+m_\mu$.
The regulated fine operator is
$A_\mu=A+d_Ud_U^*+M_\mu$, where $A$ is the fine kinetic Hessian
and $M_\mu$ applies $m_\mu$ to each one-form component.
Its determinant response is
$\Pi_\mu=\frac12\log\det A_\mu-\log\det G$.
The same ghost appears in both expressions and cancels in their
difference.

\section{The regulated lattice determinant}\label{ch07:one-loop}

Perturbative matching for the Wilson lattice action is treated in
L\"uscher--Weisz~\cite{LuscherWeisz1995}; our comparison additionally
controls the selected non-Wilson path action and its averaging operation.
The continuum coefficient is the classical calculation of
Gross--Wilczek and Politzer~\cite{GrossWilczek1973,Politzer1973}.
We use the background-field formulation~\cite{Abbott1982}; the following
normalization check and integrable error criterion are the parts needed
for the lattice application. Recall that $\alpha$ includes the factor
two converting a normalized quadratic Taylor coefficient into the
coefficient of $\frac12\langle d_1v,d_1v\rangle$.

Let $k$ denote the loop momentum and $E_{ab}$ the matrix unit on
spacetime indices. For external momentum $te_1$, polarization $e_2$,
and unit vector $u\in\gLie$, the leading ghost vertex is
$-2ik_2\operatorname{ad}u$. The vector
vertex has four copies of this orbital term and the spin term
$2it(E_{21}-E_{12})\operatorname{ad}u$. The orbital--spin trace is zero.
With vector and ghost bubble coefficients $-1/4$ and $+1/2$, respectively,
and $\tr(\operatorname{ad}u)^2=-C_G$, the coefficient extracted by
\eqref{ch07:alpha} has homogeneous integrand
\begin{equation}\label{ch07:pole}
 h(k)=4C_G\left(\frac{k_2^2}{|k|^6}
       -\frac{4k_1^2k_2^2}{|k|^8}-\frac1{|k|^4}\right).
\end{equation}
This is the standard background-field bubble calculation, with its
Lie algebra normalization expressed in our metric. It is not a claim that
the full lattice integrand equals $h$.

\begin{proposition}[The lattice integrand remainder]\label{ch07:integrand-input}
Let the incoming analytic path action have a fixed positive exponential
coefficient bound, a uniformly positive curvature metric, and kinetic
normalization $H(0)=I$. Form its left-coordinate Hessian (the Hessian in left increments of the proof below), add
the gauge-fixing term $d_Ud_U^*$ to it, and use the covariant regulator $m_\mu$ defined above.
The regulated quadratic coefficient has an integrand $I_\mu$ satisfying
\begin{align*}
 \alpha(\Pi_\mu)&=(2\pi)^{-4}\int_{[-\pi,\pi]^4}
                                    \operatorname{Re}I_\mu(k)\,dk,\\
 |I_\mu(k)|&\le A(|k|+\mu)^{-4},\\
 |I_\mu(k)-h(k)|&\le D(|k|^{-3}+\mu^2|k|^{-6})
                  \quad(\mu\le|k|\le1).
\end{align*}
The constants depend on the fixed path bounds, on the lower bound for the metric and on the group,
but not on $\mu\in(0,1]$ or the torus period. The finite-torus Hessian
is the corresponding momentum mesh average, with an exponentially small
per-site error relative to the integral at each fixed $\mu$.
\end{proposition}

\begin{proof}
We specify the derivatives used in the comparison. In one vector or
ghost channel let $a(k)$ be the zero-background inverse propagator,
$b(p;k)$ its first derivative in the background field, for short its first \emph{background derivative}, and $c(p,p';k)$ its second
background derivative. These are exact derivatives of the same operator; the
internal momentum of $b(p;k)$ changes from $k$ to $k+p$.
Their contribution to the logarithmic Hessian is
\begin{equation}\label{ch07:response-integrand}
 \tr\bigl(a(k)^{-1}c(-p,p;k)
       -a(k)^{-1}b(-p;k+p)a(k+p)^{-1}b(p;k)\bigr).
\end{equation}
Take half the vector expression minus the ghost expression and apply
the normalized external coefficient functional. This defines $I_\mu$.
Finite Fourier inversion gives the identical formula with the internal
integral replaced by the mesh mean. Thus the integrand is associated
with the fine action, rather than with a separately chosen
continuum model.

We first check its local Taylor data. A transported plaquette curvature
has, under the formal scaling $A(x)=\epsilon a(\epsilon x)$,
leading term $\epsilon^2(da+[a,a])$ in the corresponding orientation.
Its remaining terms have scaling dimension at least three; a Taylor derivative or
a factor of $A$ counts as one unit. This follows by expanding the four
edge exponentials and their parallel transports. At scaling dimension four the row squares
and direct term therefore have coefficient on $\gLie$-valued one-forms
$c_{\rm dir}I+\widehat\rho(0)^*\widehat\rho(0)=I$.
Cubic and quartic curvature terms have scaling dimensions at least six and eight.
The leading local density is consequently the Wilson curvature-square
density. All differences from it have scaling dimension at least five.
Exponential path summability justifies these expansions term by term
and bounds their Taylor remainders: each derivative contributes a polynomial
in path length, absorbed by a smaller \emph{width}, that is, a smaller exponential decay rate of the path coefficients.

The Hessian in left increments is obtained from
$S(\log(e^{x_+}e^{x_-}))$ by two $x_+$ derivatives. The
Baker--Campbell--Hausdorff correction starts with two fields and hence
raises the scaling dimension relative to its linear term. Accordingly, after two
fluctuation and $n$ background derivatives, the difference from the Wilson
response has momentum order at least $3-n$, for $0\le n\le2$.
This is a statement about the combined momenta of all legs. It gives
uniform compact derivative bounds as well as the indicated orders.

In particular the gauge-fixed vector inverse propagator is
\[
 a_{\rm loc}(k)=d_1(k)^*H(k)d_1(k)+d_0(k)d_0(k)^*
              =|k|^2I+O(|k|^3).
\]
The ghost inverse propagator is $\omega(k)I$, with
$\omega(k)=\sum_a|q_a(k)|^2=|k|^2+O(|k|^4)$.
For a mark of direction $a$ and Lie algebra direction $u$, expanding the Wilson cubic
vertex and adding the gauge-fixing vertex gives the linear part of the first vector
background derivative
\[
 b_{\rm lin}(p;k)_{ij}
 =\bigl[-i(2k_a+p_a)\delta_{ij}
              +2ip_j\delta_{ia}-2ip_i\delta_{ja}\bigr]
                         \operatorname{ad}u.
\]
The linear part of the first ghost background derivative is $-i(2k_a+p_a)\operatorname{ad}u$.
These expressions can also be read from the continuum quadratic
operator in background gauge. The lattice remainder in each is
$O((|k|+|p|)^2)$, by the comparison of scaling dimensions just proved. The second
background derivative and its first two external momentum derivatives are uniformly
bounded on the real momentum cube. No cancellation of a seagull term
has been assumed.

For completeness we give the bounds that propagate this information
through the regulator and the inverses. At a fixed internal momentum
$k$ put $T=|k|+\mu$, and for a matrix path $M(t)$ write
\[
 \|M\|_{T,2}=\max_{0\le j\le2}T^j\|M^{(j)}(0)\|.
\]
The paths have displaced momenta $k+te_1$, with velocities of bounded
size. The product rule bounds products by a fixed multiple of the
product of these norms. The two inverse derivative identities give
$\|M^{-1}\|_{T,2}\le C/T^2$ whenever
$\|M\|_{T,2}\le CT^2$ and $\|M(0)^{-1}\|\le C/T^2$.

The lower bound on the inverse propagators at real momenta holds globally on the centered cube.
Indeed $\omega(k)\ge c|k|^2$, the exterior identity controls the sum of
transverse and longitudinal parts, and
\[
 x+\mu^2\left(\frac{\mu^2}{x+\mu^2}\right)^9
                    \ge2^{-10}(x+\mu^2)\quad(x\ge0).
\]
Thus the inverse propagators of both channels are bounded below by $cT^2I$.
The massive ghost operator $\Delta_U+\mu^2I$ has bounds for its inverse propagator and its
first and second background derivatives, of respective sizes
$T^2,T,1$ in the scaled derivative norms. Inverting gives sizes
$T^{-2},T^{-3},T^{-4}$: a first background derivative is $-R bR$, and a second
is the sum of $-RcR$ and the two ordered $RbRbR$ terms.
Multiplication by $\mu^2$, taking nine products, and multiplying once
more by $\mu^2$ gives regulator sizes
\begin{equation}\label{ch07:regulator-budgets}
 \mu^2,\qquad \mu^2/T,\qquad \mu^2/T^2.
\end{equation}
This uses $\mu^2/T^2\le1$. Product differentiation through order two
preserves the same sizes in $\|\cdot\|_{T,2}$; the number nine is fixed.
The four constant maps $J_a$ in the vector regulator $M_\mu=\sum_aJ_a^*m_\mu J_a$, where $J_a$ takes a one-form to its component in direction $a$, do not
change these powers.

The first term of \eqref{ch07:response-integrand}, differentiated twice
in $t$, is bounded by $CT^{-2}+C\mu^2T^{-6}$. Its unshifted inverse
is independent of $t$; the local second background derivative has a unit-scale
bound, while its regulator part uses \eqref{ch07:regulator-budgets}.
The second term has two inverses of size $T^{-2}$ and two first background derivatives
of size $T$, with two derivatives contributing $T^{-2}$. Its bound is
$CT^{-4}$. Since $T$ stays bounded above on the cube, these estimates
give the first asserted radial bound.

For the comparison use $T=|k|$ and restrict to $\mu\le T\le1$.
Let $\delta=T+\mu^2/T^2$. The regulated lattice inverse propagator minus the free
quadratic one has scaled norm at most $C\delta T^2$.
The resolvent identity therefore bounds the inverse difference by
$C\delta T^{-2}$. Likewise the first background derivative minus its displayed
linear part is bounded by $C\delta T$: the local remainder is
$O(T^2)$ and the regulator contribution is $O(\mu^2/T)$.
Subtract the free bubble by telescoping its four factors. Each term
contains one difference, so its second external derivative is bounded
by $C\delta T^{-4}=C(T^{-3}+\mu^2T^{-6})$.
The seagull estimate above is bounded by the same expression on this
annulus. The free vector-minus-ghost bubble is exactly
\eqref{ch07:pole}, proving the second radial inequality.

At fixed $\mu>0$ the inverse propagators are invertible on a smaller complex
strip by their positive real lower bound and analyticity. All three response
layers come from the same infinite-lattice path series. The integrand
thus has a common holomorphic extension for that fixed $\mu$.
For a periodic analytic function, the difference between its mesh mean
and its integral is the sum of nonzero Fourier coefficients whose
indices are multiples of the mesh side. Exponential coefficient decay
bounds this error by $C_\mu e^{-c_\mu n}$.
The constants in this last mesh estimate may depend on $\mu$;
the radial constants proved above do not. This proves the final assertion
and justifies the coefficient integral.
\end{proof}

\section{Uniform comparison of determinants}\label{ch07:comparison}

Two estimates must be distinguished. An absolute determinant Taylor
bound may grow like the number $L^4$ of fine sites per block. The
\emph{difference} between the blocked and regulated determinant must
have a bound independent of $L$ in order to identify the logarithmic
coefficient. The relevant mechanism is finite rank together with a
trace-class regulator.

Here is the structure behind the trace-norm requirement. Writing $H$
for the derivative of the averaging operation, $B$ for the Hessian of its penalty term, $Z$ for the
coarse-frame projection, and $K=m_\mu X$, cancellation of the fine
kinetic Hessians leaves
\begin{align}\label{ch07:theta}
 \Theta={}&H^*BH-Z^*Z+P_f^*d_UK+K^*d_U^*P_f+K^*K\\
 &\hspace{12mm}-P_f^*M_\mu P_f.\notag
\end{align}
The first two terms factor through spaces of dimensions at most
$4\dim\gLie$ and $\dim\gLie$ per coarse momentum. If $W$ maps into such an intermediate space of dimension $r$ and $V$
maps out of it, $VW$ has at most $r$ nonzero singular values. Thus its
trace norm satisfies $\|VW\|_1\le r\|V\|\|W\|$, uniformly in $L$.
Derivatives preserve the intermediate spaces.

For the regulator, at zero background the alias eigenvalues of $P_\mu$
are bounded by $C(1+|z|)^{-2}$ in a smaller complex strip. Consequently
\[
 \|P_\mu^3\|_1\le C\sum_{z\in\ZZ^4}(1+|z|)^{-6}<\infty.
\]
Given bounded first two background derivatives of $P_\mu$, distribute at most two
derivatives among the three factors in $(P_\mu^3)^3$. Every term retains
one undifferentiated factor $P_\mu^3$, which is trace class; the other factors have operator
bounds. Hence the first two background derivatives of $m_\mu$ have trace norm
$O(\mu^2)$. This explains the ninth power and avoids an erroneous
factor $L^4$. To apply this trace estimate, we establish the comparisons of weighted and Euclidean
norms, the background derivatives at the critical configuration and the scaled inverse bounds below.

\subsection{The critical configuration and the contact term of the averaging operation}
\label{ch07:contact}

We verify the assertion about the contact correction to the determinant, defined below, for the averaging operation used in the
block integral. Let $\zeta(x)$ be the representative of its critical configuration in centered
coordinates (in which the local gauge transformation is normalized in the mean, Section~\ref{ch07:weighted}), with $\zeta(0)=0$ and tangent $(Wx,r_x)$, where $r_x$ is its
coarse scalar component after recentering.
\mgubi{Issue 10: which $W$? Both $W_a$ (then $r_x=-RW_ax$) and $W_F$ are possible. Independence of the choice needs invariance of $f$ under fine gauge transformations; only constant ones are stated.}%
For each scalar determinant function $f(x,z)$, constant gauge transformations
act orthogonally on its matrices and leave $f$ invariant. Every invariant
linear functional on the simple Lie algebra vanishes, since it annihilates
brackets and $[\gLie,\gLie]=\gLie$. Applying this to each coordinate gives
$Df(0)=0$. Thus
\begin{equation}\label{ch07:tangent-chain}
 D^2[f(x,\zeta(x))](0)[u,v]
   =D^2f(0)[(u,Wu,r_u),(v,Wv,r_v)].
\end{equation}
Second derivatives of the critical-configuration map have disappeared for this
specific reason. Matrix Hessians themselves need not have vanishing first
derivative; the statement is about their invariant scalar determinants.

Use rooted coordinates (in which the local gauge transformation is normalized at the root of each block) for the following exact calculation. Let $w_x(v)$
be the logarithmic change in the averaging operation under a fluctuation left increment
$v$, with $w_x(0)=0$. Let $r(x)$ be the derivative of the penalty term
with respect to that increment of the averaging operation, and set
\[
 S_x(a)_{ij}=\langle a,D_v^2w_x(0)[e_i,e_j]\rangle,
 \qquad Z(x)=S_x(r(x)),\qquad A(x)=M_{\rm core}(x).
\]
The Hessian chain rule gives $M_{\rm act}=A+Z$, with $r(0)=Z(0)=0$; we call $Z$ the \emph{contact term} of the averaging operation, and the difference in \eqref{ch07:contact-exact} below the \emph{contact correction}.\mgubi{Issue 10: $M_{\rm act}$, $M_{\rm core}$ are not defined. From the Lean definitions: $M_{\rm act}$ is $C=J+F^*F$ at the critical configuration; in $M_{\rm core}$ the penalty Hessian is replaced by $H^*BH$ (probably $C_0$ up to the change of coordinates; not verified).}
The change of fluctuation coordinates from rooted to centered is triangular with determinant one,
so it does not change the determinant. Define $Y(x)$ by
$\langle a,Y(x)\rangle=\tr(A(x)^{-1}S_x(a))$.
The averaging operation is covariant, and the matrices on the fluctuation coordinates transform by the
same congruence. Cyclicity of trace therefore makes $Y$ a covariant
coarse vector. Constant gauge transformations imply $Y(0)=0$.

The linear block identity $QW=I-b^{-1}K_c$ fixes the derivative of the
penalty gradient:
$Dr(0)=-K_c=:J$. Expanding $\log\det(A+Z)-\log\det A$ through degree two
now gives
\begin{align}\label{ch07:contact-exact}
 D^2(\log\det M_{\rm act}-\log\det M_{\rm core})(0)[u,v]
 ={}&\langle Ju,DY(0)v\rangle+\langle Jv,DY(0)u\rangle\\
 &-\tr(A(0)^{-1}S_0(Ju)A(0)^{-1}S_0(Jv)).\notag
\end{align}
In particular the derivative of $A^{-1}$ is included in $DY$;
$Y(0)=0$ removes the term with $D^2r$.

There is an exact identity $DY(0)d_{0,c}=0$. Gauge the trivial
configuration by a coarse gauge transformation, then restore the chosen
block axial gauge. Gauge invariance preserves stationarity; uniqueness of the
local critical configuration identifies the restored configuration with the
critical configuration on this curve. Covariance and $Y(0)=0$ show that $Y$
vanishes there. Differentiation proves the identity. This uses the nonlinear
averaging operation, rather than an independently chosen linear averaging operator.

At fixed $L$, the finite-stencil averaging operation and the analytic inverse of
$A(0)$ give holomorphic infinite-lattice kernels for every term in
\eqref{ch07:contact-exact}. Normalized finite-torus traces are momentum
mesh averages. Their limit is the integral of these same kernels;
exponential Fourier decay gives uniform convergence also along varying
external mesh momenta. Write $M(p)$ for the mixed kernel of $DY$.
Its Ward identity along $p=te_\nu$ reads
$q_\nu(te_\nu)M_{j\nu}(te_\nu)=0$ as an identity on $\gLie$.
Hence every column of $M(0)$ is zero and $M(p)=O(|p|)$.
Since $J(p)=-K_c(p)=O(|p|^2)$, the first two terms of
\eqref{ch07:contact-exact} have order at least three, and the last has
order at least four. Therefore the contact correction has zero $\alpha$.
Only qualitative analyticity at fixed $L$ was used here. Its constants
need not be uniform in $L$, because this is an exact vanishing statement.

\subsection{Bounds for the centered matrices}\label{ch07:weighted}

The smooth representative of the linearization $W$ of the critical configuration is $W_a$ from \eqref{ch05:corrected-transport}; its representative in the block axial gauge is $W_F$ from \eqref{ch05:stationary-tangent}. As above, $\beta$ is the inverse coupling, and $b$ is the penalty parameter of the block kernel.

We establish the operator estimates for the selected path action and
averaging operation. The inputs are their positive normalized metric, fixed
exponential path bounds, and the identity \eqref{ch05:stationary-tangent} and smooth
representative of Section~\ref{ch05:transports}; all refer to the same
two consecutive scales. The coarse scalar part of the smooth representative
may be nonzero after recentering. Its derivatives below are merely
required to be bounded, not to vanish.

Recall that $\mu=L^{-1}$ is the inverse block ratio. Use the unitary
block Fourier transform $P(p)$ on positions within the block, and set
\[
 h_l(p)=(|q(\operatorname{Re}(p+2\pi l)/L)|^2+\mu^2)^{1/2},
 \qquad
 \|u\|_{p}^2=\sum_l h_l(p)^2|(P(p)^{-1}u)_l|^2.
\]
The joint space also has a coarse scalar component with its ordinary
Euclidean norm. Opposite momenta are paired in the dual norm.
On a fixed small coarse strip, $P$ and $P^{-1}$ have bounded operator
norms, $h_l\ge\mu$, and $|q(k_l)|\le5h_l$.
Thus a unit fluctuation direction has normalized $L^2$ norm over positions within the block
$O(\mu)$ and its first fine difference has normalized $L^2$ norm
$O(\mu^2)$. Background smooth bridge legs have the same powers in
supremum norm, by \eqref{ch05:actual-bridge}. Extending to the fixed
padded box of radius $(j+5)L$ around the patches of the averaging operation only changes constants: its number of residue
blocks is fixed at the given smoothing order, and its Bloch phases
are uniformly bounded.

\paragraph{Background derivatives of the fine kinetic form.}
The zero-background kinetic form is a curl on each fluctuation leg.
Its first background derivative is a sum of products with one finite
difference among three legs; its second background derivative needs none among
four legs. These representations have exponentially summable
coefficients with uniform bounds. For the first background derivative, subtract
its constant momentum coefficient, which is zero by the
Taylor calculation at scaling dimension four in Proposition~\ref{ch07:integrand-input}, and use
coordinate division. The zero-background representation is exactly
$d_1^*Hd_1$; the second background derivative already has the required form.
The number of small factors is
\[
 (n+2)+\text{number of differences}=4\qquad(0\le n\le2).
\]
H\"older with two fluctuation $L^2$ factors and the remaining supremum
factors gives $O(\mu^4)$ in normalized counting measure. Summing over the
fine base point contributes $L^4$, canceling it. A translation by $z$ introduces a
Bloch factor of modulus at most $C e^{\sigma|z|/L}$. The exponential
decay of the coefficients makes the resulting sum over $z$ uniformly bounded. Thus the fine kinetic Hessian and its first two
background derivatives have bounded operators between the stated
weighted spaces, uniformly in $L$. This is an estimate for the
left-coordinate response, whose definition by derivatives in left increments was used above.

\paragraph{The centered local gauge transformation and the penalty term.}
A small supremum bound alone is not enough for two fluctuation directions.
Use the uniform probability measure on the patch $Ly+\mathbb B_R$, $R=(j+3)L$, of the averaging operation and the three
norms $L^\infty,L^2,L^1$. Background directions use the first, one fluctuation
direction the second, and two fluctuation directions the third.
Products are controlled by $L^2L^2\subset L^1$, and averages contract
these norms. A contour of $O(L)$ links, each of whose positive derivatives
is $O(\mu)$ in its assigned norm, has bounded derivatives of every fixed
order. Indeed adjoining one more link changes the next derivative by
$O(\mu)$ times a polynomial in already bounded lower derivatives;
summing $O(L)$ such changes gives $O(1)$. For the coordinate contours of the transporters,
the link read at a fixed position sits at a translate of a point having some coordinates of each endpoint.
On the endpoint pair space this is a measure-preserving coordinate
permutation, followed by a bounded-measure restriction and translation.
The argument therefore applies to the averaged transporters.

Here is why centering gives uniform derivatives of the local gauge transformation. Let $F_x$ be the
minimizing local gauge transformation after recentering, $E_{xy}$ its averaged contour variable,
$\nu$ normalized counting measure on its block, and $\omega$ the
smoothing probability. If $\mathcal P$ is the nearest-group map and
$\Pi=D\mathcal P(1)$, define
\[
 J_0=\sum_y\nu(y)F_y,\qquad
 B_x=\sum_y\nu(y)F_y(E_{xy}-1),
\]
\[
 R(T)=\mathcal P(T)-\Pi T,\qquad
 S_x=\Pi B_x+R(J_0+B_x).
\]
The minimizing equation $F_x=\mathcal P(J_0+B_x)$ and centering
$\sum_x\omega(x)\Pi F_x=\Pi1$ imply exactly
\[
 F_x=\Pi1-\sum_y\omega(y)\Pi S_y+S_x.
\]
Now $E_{xy}(0)-1=0$ and $DR(1)=0$. At derivative order $n+1$ the
right side depends on derivatives of the local gauge transformation only through order $n$.
Induction using the preceding product bounds gives uniform derivatives
through order four, also with two fluctuation directions. The smoothing
probability satisfies $\omega\le C_j\nu$: its reduction modulo $L$
is uniform, so every atom has mass at most $L^{-4}$, and its support
has $O_j(L^4)$ sites. This controls the change of averaging measure.
The argument estimates the selected minimizer; its existence and local
uniqueness come from Chapter~\ref{ch04:blocking}.

The averaging of transporters and the fixed smooth penalty term are compositions of
these operations. Their mixed derivatives with up to two background and
two fluctuation directions are uniformly bounded. The
logarithmic derivative $X$ of the centered local gauge transformation has normalized $L^2$ output bounded by a
constant, so its unnormalized output is at most $CL^2$. Rooted and
centered descriptions differ by a triangular coordinate change of
determinant one; the same averaged-contour bound controls its background
coarse coordinates. Consequently $F=d_U^*P_f+m_\mu X$ has a uniform
weighted operator bound, since $m_\mu=O(\mu^2)$ cancels $L^2$.
Covariance of the local gauge transformation gives $X\mathcal R=I$ in the same coordinates, and
hence $F\mathcal R=\Delta_U+m_\mu$. Gauge signs are fixed here by
taking the fine component of $\mathcal R$ to be $d_U$.

\paragraph{Uniform inverse bounds.}
At zero background the regulated fine form has the weighted lower bound
$2^{-10}\min(1,\varepsilon)$, where $H_f\ge\varepsilon I$.
This follows from the exterior identity and the lower bound for the scalar regulator
proved in Proposition~\ref{ch07:integrand-input}. The kinetic form with the penalty
term added, the form $\mathcal G$ below, needs a quotient estimate. In centered coordinates it is
\[
 \mathcal G(u,r)=\langle u,K_fu\rangle
                         +b|M u-d_{0,c}r|^2.
\]
Here $M$ is the linear averaging operator. Indeed the rooted
linearization of the averaging operation is $Q=M+d_{0,c}R$ by \eqref{ch05:actual-Q}; setting
$r=r_{\rm root}-Ru$ converts $Qu-d_{0,c}r_{\rm root}$ to the
displayed expression. Under a fine gauge variation, the new scalar
coordinate changes by the smoothed scalar mean, since
$R d_{0,f}$ is root evaluation minus that mean. This is the same
centering used in the calculation for the local gauge transformation, and the triangular coordinate
change has determinant one.
Subtract a gauge vector to make every nonzero alias of $u$
transverse and eliminate $r$ using the scalar zero alias. The
latter scalar mean is $\mu^2$ times a product of centered box sums,
bounded away from zero after division by $\mu^2$ on the principal
real cube. Write the resulting fine aliases as $u_l$ and set
$A_h=\sum_{l\ne0}h_l^2|u_l|^2$. The lower bound on the fine metric gives control
of $A_h$, since $|q(k_l)|\ge c\mu$ for $l\ne0$.
The zero-alias coefficient of $M$ has least singular value $c_j\mu$:
its sum along the transporter is of order $L$ and the unitary block Fourier
coefficient is $L^{-2}$. For the upper bound on $M$, fix one of the $L$
steps of the transporter. Averaging the link at that step over the offsets gives coefficients
$O_j(L^{-4})$ on $O_j(L^4)$ sites, hence Euclidean row norm
$O_j(L^{-2})$. Summing over its $L$ positions along the transporter and converting
a weighted unit vector with $h_l\ge L^{-1}$ gives a fixed bound.
Thus the full $M$ has bounded weighted norm.
Consequently
\[
 \mu^2|u_0|^2\le C_j(|Mu|^2+A_h),
 \qquad
 \inf_a\|(u,r)-\mathcal R a\|_p^2\le C_j\mathcal G(u,r).
\]
These estimates include zero coarse momentum; the scalar zero
alias removes the coarse scalar coordinate there as well.

Moreover $\|\mathcal R a\|_p\le C|\Gamma_\mu a|$, where
$\Gamma_\mu=\Delta+m_\mu$. For the fine component this is
$h_l|q_l|\le h_l^2\le C\Gamma_\mu(k_l)$; the coarse averaging
component has norm $O(\mu^2)$ and obeys the same bound. Since
$F\mathcal R=\Gamma_\mu$ and $F$ is bounded, the quotient estimate
implies a positive weighted lower bound for $\mathcal G+F^*F$.
All constants are fixed before $L$.

The weights are not holomorphic. To continue this lower bound, at each real
base point hold $h_l$ constant while retaining the holomorphic block Fourier
matrix. On a fixed-radius coarse ball these constant weights are uniformly
comparable to the moving weights, because $h_l\ge\mu$ and
$|q(k_l(p))-q(k_l(x))|\le C\mu|p-x|$.
Cauchy's inequality in the coordinates with constant weights bounds the change in
operator norm, and a fixed smaller strip preserves half the real lower bound.
This proves uniform holomorphic inverse bounds for both matrices.
Differentiating $M^{-1}M=I$ then bounds the first two background derivatives of the inverses.

\paragraph{Trace norms and finite-volume coefficients.}
The kinetic response cancels identically, including its two derivatives,
in \eqref{ch07:theta}. The finite-rank terms retain their fixed-dimensional
coarse intermediate spaces. For the regulator, the argument given before this
subsection for the trace-class factor $P_\mu^3$ applies. In the terms containing
$m_\mu X$, its factor $\mu^2$ compensates for the factor $L^2$ of the derivative $X$ of the local gauge transformation.
In the separate term $P_f^*M_\mu P_f$, it compensates for two comparisons of weighted and Euclidean
norms, each of size $O(\mu^{-1})$. These large factors occur in
different terms; the same small factor is not used to compensate for both. Positive background derivatives of $d_U$ have a fine-link factor
$O(\mu)$; the single norm comparison contributes at most $O(\mu^{-1})$.
The norms of these first two background derivatives are therefore bounded in both
placements needed for $d_U^*d_U$ and $d_Ud_U^*$. The remaining
products use exactly these compatible input/output norms. Thus
$\Theta$ and its first two background derivatives have uniform trace norm. This proves all operator and trace bounds used in
Proposition~\ref{ch07:comparison-prop} for this same action and local gauge transformation.

Finally the infinite-lattice Fourier representations follow from the infinite
path coefficients and finite stencils of the averaging operation. All derivatives are
taken before restricting to a torus grid. Products and inverse derivatives
preserve this identification, with intermediate momenta shifted by the
inserted external momentum. The Fourier mesh estimate at the end of
Proposition~\ref{ch07:integrand-input} gives the finite-volume exponential
error. The preceding contact proof identifies the coefficient of
the Gaussian contribution \eqref{ch07:thermal}. Hence the comparison concerns the same density,
averaging operation and critical configuration throughout.


\begin{proposition}[Common comparison estimates]\label{ch07:comparison-input}
On a fixed positive strip of coarse momenta, the background derivatives at the critical
configuration and all intermediate shifted momenta obey the following
bounds, uniformly in $L$ and in the period. After removing the
contact term of the averaging operation, the blocked completed matrix $C_0$ and
$\widetilde A_\mu=\operatorname{diag}(A_\mu,I)$ have bounded inverses
and bounded first and second normalized background derivatives in
compatible weighted Hilbert norms. Their difference $\Theta$ and its
first two background derivatives have bounded trace norm in these same norms.
The infinite-lattice quadratic symbols represent the finite-torus coefficients
with an error tending exponentially to zero with the coarse period.
The linearization of the critical configuration is the one built with the averaging operation and obeys the
exact coarse--fine matching identities.
\end{proposition}

\begin{proof}
The preceding weighted inverse, background-derivative and trace estimates prove this
proposition for the normalized path action and the averaging operation.
\end{proof}

\begin{proposition}[Block comparison]\label{ch07:comparison-prop}
For the same normalized path action and averaging operation, the curvature-square
coefficients of the blocked response and the regulated fine determinant
satisfy $|\alpha_{N,m}-\alpha(\Pi_\mu)|\le C$, uniformly in $L,N,p$.
\end{proposition}

\begin{proof}
For a positive matrix continued holomorphically near zero background,
\[
 D_vD_w\log\det A
 =\tr(A^{-1}A_{vw}-A^{-1}A_vA^{-1}A_w).
\]
Subtract this formula for $C_0$ and $\widetilde A_\mu$, and use
$C_0^{-1}-\widetilde A_\mu^{-1}
=-C_0^{-1}\Theta\widetilde A_\mu^{-1}$.
Telescoping produces terms containing a trace-class difference or one
of its derivatives, with all other factors bounded. The trace ideal
inequality gives a uniform bound for the difference symbol on a common
strip. The ghosts cancel. The contact correction has zero $\alpha$ by
the preceding order argument.

Folding the fine symbol into coarse momentum also preserves $\alpha$.
The regulated determinant is gauge invariant, so its quadratic symbol
annihilates longitudinal legs. For a nonzero alias the equation
$K_f(k_l)N_l=A_lK_c(p)$ and the positive transverse fine inverse give\mgubi{Issue 10: $N_l$ is introduced only in the alias paragraph near the end of the chapter, and $A_l$ nowhere ($\mathcal A_l$ of Chapter~\ref{ch05:gaussian} is another object). Say that the alias amplitudes are those of $W_s$.}
$N_l^\perp=O(|p|^2)$, since $K_c(p)=O(|p|^2)$. Thus nonzero
aliases contribute order four even when their longitudinal Fourier coefficients
do not vanish. The zero alias has factors $L^4$ from
cell volume, $L^{-2}$ from the two interpolated link fields, and
$L^{-2}$ from extraction at momentum $p/L$. Their product is one.
The inherited loop has order four and contributes no $\alpha$.
Finally Cauchy's inequality on a circle inside the common strip bounds
the second Taylor coefficient of the difference by a fixed multiple of
its supremum. This proves the assertion.
\end{proof}

\subsection{The one-loop shift of the coupling}

\begin{proposition}[One-loop shift of the coupling]\label{ch07:slope}
For two consecutive scales with the normalized path action, the lower bound on the metric and the
identity \eqref{ch05:stationary-tangent} used in Section~\ref{ch07:weighted},
\[
 |\alpha_{N,m}+\kappa_G\log L|\le C
\]
with a constant depending only on the fixed stage inputs.
\end{proposition}

\begin{proof}
On the unit three-sphere the angular means are
$\mathbb E\theta_2^2=1/4$ and
$\mathbb E\theta_1^2\theta_2^2=1/24$. Therefore
\[
 (2\pi)^{-4}\int_{\mu<|k|<1}h(k)\,dk
 =\frac{8\pi^2C_G}{(2\pi)^4}
   \left(\frac14-\frac16-1\right)\log\mu^{-1}
 =-\kappa_G\log\mu^{-1}.
\]
The two error integrals on this annulus are uniformly bounded:
$\int_\mu^1 r^3r^{-3}dr\le1$ and
$\mu^2\int_\mu^1r^3r^{-6}dr\le1/2$.
The ball $|k|\le\mu$ has volume $O(\mu^4)$ and integrand
$O(\mu^{-4})$; the region outside the unit ball is also uniformly
integrable. Thus $\alpha(\Pi_\mu)=-\kappa_G\log\mu^{-1}+O(1)$.
Proposition~\ref{ch07:comparison-prop} and $\mu=L^{-1}$ finish the proof.
\end{proof}

The sign is for the step from scale $m+1$ to scale $m$: the coefficient $\beta$ of the
kinetic action decreases by $\kappa_G\log L$ at one loop.
This proposition does not assert the exact nonlinear recursion for
$\beta$. Connected interactions contribute to that recursion and are
handled by the one-step integral identity.

\section{Classical matching through fourth order}\label{ch07:classical}

\subsection{Uniform classical coefficients at the critical configuration}
\label{ch07:classical-uniform}

To construct the cubic and quartic matching terms, we first need
classical coefficient bounds fixed before $L$. The fixed-$L$ chart for the critical configuration of
Chapter~\ref{ch05:gaussian} provides existence and differentiation, but
its inverse in the block axial gauge alone does not provide these stronger bounds.
We obtain them from the weighted quotient just constructed.

Let $\mathscr S(v,z)$ be the joint classical action, with coarse field
$v$ and fluctuation (internal) coordinates $z$, the fine field together with the local gauge transformation. It contains the fine
kinetic action and the penalty term built with the averaging operation, without $\beta$ or the
inherited loop. At zero its first derivative vanishes. Write
$\mathscr H=D_z^2\mathscr S(0)$ and let $\mathcal R$ denote its
infinitesimal gauge action on the fluctuation coordinates. Its gauge-fixed inverse is
\[
 \mathscr C^{-1}=(\mathscr H+F^*F)^{-1},
\]
with the uniform weighted bound just proved. For a coarse direction $u$
let $T u=(u,W_a u,r_u)$ be the full linearization of the critical configuration in the smooth
centered representative. In particular
\begin{equation}\label{ch07:tangent-orthogonal}
 D^2\mathscr S(0)[Tu,(0,z)]=0
 \quad\text{for every internal }z.
\end{equation}
This is the differentiated stationarity equation, with the
identification for the averaging operation in Chapter~\ref{ch05:gaussian}. Stationarity
in all joint fluctuation directions follows from stationarity on the
slice of the block axial gauge and gauge invariance: that slice together with the joint
gauge directions spans the tangent space of the fluctuation coordinates. Thus coarse-frame
directions are included in \eqref{ch07:tangent-orthogonal}.
Define the covector on fluctuation directions
\[
 B(u,v)[z]=D^3\mathscr S(0)[Tu,Tv,(0,z)].
\]

Both the cubic and quartic vertices with background legs and $B$ have uniformly
bounded infinite-lattice Fourier kernels. Here is the power count in their
weighted norms. A fine-action cubic vertex has at least one
finite difference among its three legs; a quartic vertex needs none.
Their exponentially summable coefficient bounds follow from the same
expansion of the normalized path action and coordinate division used for the background
derivatives of the kinetic form above. An undifferentiated background leg contributes $O(\mu)$ in
supremum norm, and its differentiated leg contributes $O(\mu^2)$. A fluctuation
leg of unit weighted norm contributes these same powers in normalized $L^2$ norm.
Thus three background legs with one difference, four background legs,
and two background legs with one fluctuation leg and one difference each
contribute $O(\mu^4)$ in normalized summation over positions within the block. In the last case
$L^2\subset L^1$ on probability measure is sufficient. The sum over
$L^4$ fine base points cancels this factor. Finitely many coordinate-change
terms have the same or higher scaling dimension. Penalty vertices are bounded by
the induction on mixed derivatives of the centered local gauge transformation above, which allows up to
four derivatives and includes the bounded coarse scalar legs $r_u$.
The bounds for $B$ are dual weighted norms, not Euclidean coordinate
bounds with a dimension-dependent norm comparison.

The covector $B$ annihilates gauge directions. To check this, write the
gauge identity in the fluctuation coordinates as
$D_z\mathscr S(v,z)[\mathcal R(v,z)a]=0$ and differentiate twice along
$Tu,Tv$. The term containing $D\mathscr S(0)$ is zero. The two terms
with a derivative of $\mathcal R$ vanish by
\eqref{ch07:tangent-orthogonal}, since those derivatives are fluctuation
vectors. The remaining term is
\[
                       B(u,v)[\mathcal R(0)a]=0.
\]
Therefore $w=\mathscr C^{-1}B(u,v)$ is a representative of the quotient
solution $\mathscr H w=B(u,v)$. Indeed, pairing
$\mathscr Cw=B(u,v)$ with $\mathcal R a$ gives
$\langle F\mathcal R a,Fw\rangle=0$. The regulated gauge matrix
$F\mathcal R$ is invertible, so $Fw=0$ and $\mathscr H w=B(u,v)$.
Any other solution differs by a gauge vector, against which $B$ vanishes.
This justifies using the bounded gauge-fixed inverse in the exchange term;
it does not replace a degenerate Hessian by an arbitrary inverse.

If $E(v)$ is the value at the critical configuration, implicit differentiation now gives
\begin{align}\label{ch07:uniform-stationary-jets}
 D^3E(0)[u_1,u_2,u_3]
   &=D^3\mathscr S(0)[Tu_1,Tu_2,Tu_3],\\
 D^4E(0)[u_1,u_2,u_3,u_4]
   &=D^4\mathscr S(0)[Tu_1,Tu_2,Tu_3,Tu_4]\notag\\
   &\quad-\sum_{\{ij\}\{kl\}}
          B(u_i,u_j)\mathscr C^{-1}B(u_k,u_l),\notag
\end{align}
where the sum is over the three partitions of $\{1,2,3,4\}$ into pairs.
Each summand evaluates the covector $B(u_i,u_j)$ on the vector
$\mathscr C^{-1}B(u_k,u_l)$, using the counting inner product to
identify the latter covector with a vector before applying the inverse.
These formulas follow by differentiating stationarity; the higher
derivatives of the critical-configuration map multiply its first derivative or
\eqref{ch07:tangent-orthogonal}. On the diagonal the last sum is
$3B(u,u)\mathscr C^{-1}B(u,u)$, agreeing with the earlier formula.
The gauge-annihilation argument makes the formulas independent of the
representative chosen to differentiate the critical-configuration map.

For Fourier waves, each exchange inverse has intermediate coarse
momentum equal to a sum of two external momenta. Restrict each external
strip by a fixed factor so that all these sums remain in the inverse
strip. The vertex and inverse bounds then give a common strip and a
uniform bound for both symbols in
\eqref{ch07:uniform-stationary-jets}, before choosing $L$.
Their finite-volume identities are exact: every vertex and inverse is
the restriction of the same infinite-lattice Fourier matrix, and momentum
conservation leaves the coarse volume factor exactly once. Dividing by
$3!$ and $4!$ gives normalized infinite-lattice Taylor symbols with the required
bounds. No neighborhood for the critical configuration of
size uniform in $L$ has been inferred from this argument; only its
fixed-order coefficients need that uniformity here.

\subsection{Construction of the matching curvature terms}

Let $E(v)$ be the classical action at the critical configuration of Chapter~5, with the
factor $\beta$ removed. Its Hessian is prescribed by the exact free
block relation. Write the corresponding plaquette metric as
$H=c_{\rm dir}I+\rho^*\rho$, with $H(0)=I$ and $c_{\rm dir}>0$.
The direct Wilson term together with the squares of the transported
rows of $\rho$ gives an analytic gauge-invariant action $J(v)$ with
this Hessian. Indeed the transported rows linearize to $\rho d_1v$,
so
\[
 D^2J(0)[v,v]=\langle d_1v,Hd_1v\rangle=D^2E(0)[v,v].
\]
The construction uses the coarse metric matched to the incoming
Gaussian form, not an arbitrary positive coarse metric.

The preceding coefficient estimates at the critical configuration and the fixed row bounds
give the hypotheses of the following construction.

\begin{proposition}[Classical matching]
\label{ch07:classical-match}
Suppose $E$ and $J$ are invariant analytic germs on every torus, have
zero value and first derivative and identical Hessians. Suppose their
normalized coefficients of orders three and four have infinite-lattice symbols with strip
bounds, and the row coefficients of $J$ have an exponential bound.
There are invariant cubic and quartic curvature terms $S_3,S_4$, chosen
once for all periods, such that
\[
 E-J-S_3-S_4=O(v^5).
\]
Their weighted sums are bounded by a fixed multiple of the input Taylor
bounds on a smaller positive width.
\end{proposition}

\begin{proof}
The leading coefficient of $E-J$ is cubic and satisfies the homogeneous
Ward identities. Lemma~\ref{ch07:curl} and the transported construction
produce $S_3$ with that coefficient. The fourth coefficient of
$E-J-S_3$ is again transverse, since its lower coefficients vanish.
Apply the lemma again to obtain $S_4$. It is essential to subtract the
fourth coefficient of $S_3$: a nonlinear transported cubic curvature
term has nonzero higher Taylor coefficients. The exponential path
bound controls each fixed Taylor derivative on a smaller complex link
ball; factors of path length are absorbed by a width loss. Thus the
fourth coefficient of $S_3$ has the required common bound. This proves
both the matching and the bounds. Polarization converts the equality on
rays to equality of all derivatives through order four.
\end{proof}

\subsection{Preserving the classical coefficient bound}

These estimates can also be arranged to preserve a fixed classical
coefficient bound, rather than producing an unspecified function of the
incoming bound. Keep the Gaussian metric bounds fixed, and denote the
bound on the remaining incoming path coefficients by $A$. Compare the
fine action with the Wilson action, evaluated on the \emph{same}
linearizations $T$ and with the \emph{same} gauge-fixed inverse $\mathscr C^{-1}$.
The normalization at scaling dimension four proved above implies that their cubic
vertex difference has at least two finite differences, and their
quartic vertex difference has at least one. Coordinate division of their
infinite-lattice symbols gives these factors with exponentially summable
coefficients bounded by a polynomial $P(A)$. A slightly stronger incoming
exponential width absorbs the polynomial path-length factors.

When all legs are background legs, the extra difference gives $L^4\mu^5=\mu$.
For the difference of the covectors $B$, move differences on the single
fluctuation leg to the two background legs by momentum conservation. In a
coordinate this uses
\[
 q(-p_1-p_2)=-e^{-i(p_1+p_2)}
                \bigl(q(p_1)+e^{ip_1}q(p_2)\bigr).
\]
It is a discrete integration by parts, with bounded translation factors;
it leaves at most two differences on any background leg. The
bounds on the smooth representative $W_a$ through order two therefore give the same gain.
The fluctuation leg is estimated only in normalized $L^2$, without asking
for a second weighted derivative. Thus the vertex and dual-vertex
remainders are bounded by $P(A)/L$.

The Wilson vertices, penalty term, and their exchange pairing have
bounds depending only on the fixed Gaussian metric and the data of the averaging operation.
The inverse in the exchange term depends only on the quadratic action,
so higher incoming coefficients do not change it. Expanding the two
factors $B=B_{\rm W}+\delta B$ in
\eqref{ch07:uniform-stationary-jets} shows that the normalized classical
symbols, and then their curvature families, satisfy a bound of the form
\begin{equation}\label{ch07:classical-budget}
             A_{\rm out}\le a_{\rm stage}+\frac{P_{\rm stage}(A)}{L}.
\end{equation}
The squared remainder contributes $P(A)^2/L^2$ and is absorbed into the displayed
polynomial. Subtracting the prescribed coarse quadratic row and the
quartic contribution of the constructed cubic family preserves this
form: those operations have fixed coarse-scale bounds, and the latter
operation is linear in that cubic family.
The output strip is chosen from the Gaussian coarse strip. It need not
shrink with each step: an incoming coefficient displaced by $z$ contributes a factor
at most $e^{c\sigma|z|/L}$ on that output strip, which is absorbed by the
fixed incoming width after choosing $L$ large enough. Finally choose an
outgoing classical bound larger than $a_{\rm stage}$ and then $L$
so that the second term in \eqref{ch07:classical-budget} is smaller than the difference between these two bounds.
For a finite collection of stage data these choices are simultaneous.

Consequently the cubic and quartic curvature families are constructed
for this same action at the critical configuration with bounds fixed at the stage before
$L$. Their input Hessian is the matched coarse metric.
This completes the classical matching construction without invoking the schedule of Chapter~\ref{ch10:trajectory}
or a nonlinear local-integral estimate.

\section{Inherited loops: four differences and alias suppression}

Taylor degree here counts field factors. It is distinct from scaling dimension, which also counts derivatives, as in Section~\ref{ch02:jets-section}; a momentum difference therefore improves scaling dimension without adding a field factor.
\label{ch07:inherited-proof}

Let $\Lambda_{m+1}^{\rm loop}$ denote the incoming weighted loop sum.
The estimates of Chapter~\ref{ch05:gaussian} for the representatives of the linearization of the critical configuration give, with ``pulled back'' meaning composed with the critical-configuration map,
\begin{equation}\label{ch07:inherited}
 \|\mathcal L_{m+1}\hbox{ pulled back}\|_{w'}
 \le C L^{-2}\Lambda_{m+1}^{\rm loop}
\end{equation}
To prove this bound, group incoming quadratic curvature terms by their relative displacement
$z$. Their coefficient matrix $r(z)$ has exponential sum bounded by
$\Lambda_{m+1}^{\rm loop}$ and zero zeroth and first moments. Thus the
fine link symbol
\[
 S(p)=d_1(-p)^{\mathsf T}\widehat r(p)d_1(p)
\]
has order four, satisfies $S d_0=d_0^*S=0$, and has strip bound
$C\Lambda_{m+1}^{\rm loop}$. This bound holds componentwise in the
Lie algebra and hence also for the symmetrized bilinear form, with a
constant depending on $\dim\gLie$.

An analytic periodic function with an order-four zero is a sum of
products of four coordinate differences $q_a(p)$ times analytic periodic
functions on a smaller strip, with linearly bounded coefficient norms.
To see this quantitatively, write
$f(p)-f(0)=\sum_a q_a(p)D_a f(p_1,\ldots,p_a,0,\ldots,0)$
using the divisions of Lemma~\ref{ch07:curl}, and repeat on the quotients,
which retain one fewer vanishing order. There are finitely many terms.
For a quadratic symbol, move two differences to the opposite leg by
$q_a(p)=-e^{ip_a}q_a(-p)$. The phases merely translate a coefficient
by finitely many lattice sites. Each term now has two differences on
each of its two field arguments.

The representative $W_s$, built with the linear averaging operator, has exponentially local infinite-lattice
kernels with $\nabla_f^rW_s=O(L^{-1-r})$ for $r\le2$, by
Section~\ref{ch05:transports}. By \eqref{ch05:corrected-transport}, the linearization built with the averaging operation differs
from this representative by a fine gradient; the gauge-invariant
quadratic symbol $S$ annihilates that difference on both sides.
Consequently it is legitimate to use $W_s$ in this calculation, even
if the representative $W_F$ in the block axial gauge itself has poor difference bounds.
Each differentiated leg contributes $CL^{-3}$. Summing over the fine base point in
one coarse block contributes $L^4$, so the pulled-back quadratic symbol has
bound
\[
 C L^4 L^{-3}L^{-3}\Lambda_{m+1}^{\rm loop}
       =CL^{-2}\Lambda_{m+1}^{\rm loop}.
\]
Displacing an argument by a fine displacement $z$ contributes a factor at most
$Ce^{\sigma|z|/L}$ on a coarse strip of width $\sigma$.
Choose $L$ after the fine coefficient width $w$ so that $\sigma/L<w$;
absolute summability controls these displacements. Holomorphy and the
infinite-lattice representation of the finite-torus coefficients follow from normal convergence.
A final fixed strip reduction converts this to the weighted coefficient
bound required in \eqref{ch07:inherited}.

The order-four zero on the coarse lattice requires another argument.
For aliases $k_l(p)=(p+2\pi l)/L$, let $N_l(p)$ be the Fourier
coefficient at the alias $l$ of the linearization of the critical configuration. The coarse symbol is
\[
 H(p)=L^4\sum_l N_{-l}(-p)^{\mathsf T}S(k_l(p))N_l(p).
\]
For $l\ne0$, use $K_f(k_l)N_l=A_lK_c(p)$, the exact equation for
that linearization. Project $N_l$ orthogonally off the gradient
$q(k_l)$. The positive fine curvature metric implies that
$K_f(k_l)+d_0(k_l)d_0(k_l)^*$ is invertible near $p=0$ for this fixed
$l,L$. As $K_c(p)=O(|p|^2)$, the transverse coefficient is $O(|p|^2)$.
The Ward identities of $S$ remove the longitudinal coefficients, so each
nonzero alias has order four. The zero alias has order four
because $S(p/L)$ does. This argument requires no lower inverse bound
uniform in $L$ for an individual nonzero alias.

Finally the pullback by the nonlinear critical-configuration map has the same Hessian as this
linear pullback: the inherited action has zero first derivative, so its
chain rule has no term involving the second derivative of the critical-configuration map.
This proves the weighted gain and retained moment cancellations for the
inherited action, once the linearization built with the averaging operation and the representative $W_s$
have been identified up to gauge as above.

\section{Output coefficients and the remaining analytic width}\label{ch07:output}

Suppose the determinant symbol at the critical configuration has bound $\Phi$ at a common
width, the inherited loop symbol has bound $R$ and order four, and the
kinetic symbol has bound $B_H$. Subtract the extracted coefficient
$\alpha b_H$. Lemma~\ref{ch07:unique} and the two-leg division produce
a quadratic curvature family with bound
\[
 K(\Phi+|\alpha|B_H+R)
\]
and the zero moments \eqref{ch07:moments}. Together with the prescribed
rows and the cubic/quartic families of
Proposition~\ref{ch07:classical-match}, this defines one common output
path action. If the infinite-lattice determinant symbol approximates its finite
period values with error $|\Lambda|\epsilon$, the same coefficient
error is retained by linear subtraction. For real plane waves with
component amplitudes at most one, the Hessian error is at most
$2(4\dim\gLie)^2|\Lambda|\epsilon$: expansion in component waves
contributes the square of their number, and $D^2=2C_2$.

All sums should be bounded first at a width $w'>w>0$.
Each fixed differentiation of a transported term contributes a polynomial factor in
its path length; the factor $e^{-(w'-w)\operatorname{length}}$ absorbs
it. This gives the fixed derivative orders needed for local expansion.
Parameter derivatives require the corresponding derivative bounds on
the infinite-lattice input kernels. Linear curvature division preserves such
bounds because its operations and widths are chosen independently of
the particular kernel. No parameter regularity follows merely from
matching the kernels pointwise.

For a quantitative descent, use the inherited estimate
\eqref{ch07:inherited}. If the new loop
bound is $A L^4+B\log L+C_0$, choose $L$ so that $CL^{-2}\le1/4$,
then the lower bound on the reference coupling so that
\[
 A L^4+B\log L+C_0\le\tfrac12\Lambda_m^{\rm loop},\qquad
 \Lambda_{m+1}^{\rm loop}\le2\Lambda_m^{\rm loop}.
\]
Their sum is at most the outgoing bound. This is the order of choices:
widths and fixed amplitudes, block ratio, then the lower bound on the reference coupling.

The comparison bounds and lattice integrand estimates have now been
derived from the same normalized analytic path action, positive metric and
averaging operation at the critical configuration. The inherited estimate follows from the
identified linearization of the critical configuration and the bounds on the infinite-lattice kernels of $W_s$. These results are
one-step statements: their application to an entire selected history
requires that its incoming path actions obey the fixed stage bounds and its
charts for the critical configuration satisfy the identifications of Chapter~\ref{ch05:gaussian}.
They do not replace the nonlinear local-integral bounds or prove the
existence of an infinite sequence of scales.
}{}
\setcounter{chapter}{7}
% !TEX root = ../monograph.tex
\chapter{Polymer expansions and connected correlations}
\label{ch08:chapter}

An effective density is useful only if it remains local after integration.
Its logarithm should therefore be expressed as a sum of terms supported on
connected sets, with a bound that does not grow with the surrounding volume.
The cluster expansion provides this passage from local activities to connected
terms. Its convergence is a standard result of statistical mechanics. The
work specific to the gauge theory is to put the activities produced by the
block integration within its hypotheses.

This chapter states the abstract results in the form needed for that
application. It does not prove another cluster-expansion theorem. We use
Koteck\'y--Preiss convergence\cite{KoteckyPreiss1986}, in the explicit rooted
form of Ueltschi\cite[Theorem~1]{Ueltschi2004}, and explain how support weights,
activity derivatives and large-field constraints enter. The distinction
between a general cluster estimate and a bound on the actual gauge-theory
activities will remain explicit throughout.

\section{Activities, supports and incompatibility}
\label{ch08:activities}

Let $\Lambda$ be a finite set of lattice sites. A \emph{polymer label}
$\gamma$ belongs to a finite or countable set $\mathcal P$, has a nonempty
support $X_\gamma\subset\Lambda$, and carries a complex number
$z_\gamma$, called its \emph{activity}. Different labels may have the same
support. The labels can retain information about an integration sector,
derivative or local Taylor coefficient; they are not identified solely by
their supports.

A symmetric, reflexive relation $\gamma\not\sim\eta$ specifies
\emph{incompatibility}. Reflexivity means $\gamma\not\sim\gamma$: two occurrences of the same
label are incompatible. We assume
\begin{equation}\label{ch08:overlap}
 \gamma\not\sim\eta\quad\Longrightarrow\quad
 X_\gamma\cap X_\eta\ne\varnothing.
\end{equation}
When incompatibility is defined by overlap of supports, the implication
is an equivalence. We use this particular rule for the constrained gas
in Section~\ref{ch08:forced}.

The polymer partition function is
\begin{equation}\label{ch08:gas}
 Z_{\rm pol}(z)=\sum_{\substack{A\subset\mathcal P\text{ finite}\\
                         A\text{ pairwise compatible}}}
                       \prod_{\gamma\in A}z_\gamma.
\end{equation}
The empty collection contributes one. This sum is a polynomial for finitely
many labels. For countably many labels it is well defined whenever
$\sum_\gamma|z_\gamma|<\infty$, since its absolute sum is at most
$\exp(\sum_\gamma|z_\gamma|)$. The cluster criterion below implies this
summability in finite volume. Complex or signed activities are allowed;
they need not be probability weights.

We keep track of spatial size with a nonnegative function $c(X)$ on
nonempty subsets of $\Lambda$. The needed properties are
\begin{equation}\label{ch08:cover-properties}
 |X|\le h\,c(X),\qquad
 c(X\cup Y)\le c(X)+c(Y)\quad\text{if }X\cap Y\ne\varnothing,
\end{equation}
where $h>0$ is fixed. One choice is the number of sites, with $h=1$.
A useful multiscale choice is the minimum number of radius-$B$ lattice
balls whose union is connected and contains $X$. In that case
$h=(2B+1)^4$ is an upper bound on the size of a ball, including on a torus.
The symbol $c$ will always denote this support-size function in this chapter,
not a field coefficient.

For $\alpha\ge0$ define the site norm
\begin{equation}\label{ch08:site-norm}
 \|z\|_\alpha=
 \sup_{x\in\Lambda}\sum_{\gamma:x\in X_\gamma}
                  |z_\gamma|e^{\alpha c(X_\gamma)}.
\end{equation}
It counts every polymer at each site in its support. A norm which counts
only the chosen anchor of a polymer is different. Passing from an anchored
bound to \eqref{ch08:site-norm} requires a support-size estimate or a loss
in the exponential weight. This distinction matters when checking a
cluster criterion for an effective density.

\section{The logarithm and its convergence criterion}
\label{ch08:criterion}

For a tuple $\boldsymbol\gamma=(\gamma_1,\ldots,\gamma_n)$,
form the graph on the \emph{occurrences} $1,\ldots,n$, joining distinct
vertices whenever their labels are incompatible. Define
\begin{equation}\label{ch08:ursell}
 \varphi^T(\boldsymbol\gamma)=
 \sum_{\substack{H\text{ connected spanning graph}\\
                   H\text{ uses incompatibility edges}}}
                 (-1)^{|E(H)|}.
\end{equation}
For $n=1$ its value is one; for a disconnected incompatibility graph it
is zero. Repeated labels are allowed in this tuple, even though the gas
\eqref{ch08:gas} contains each label at most once. They are necessary in
its logarithm: for a gas with one label, $Z_{\rm pol}=1+z$ and
$\log Z_{\rm pol}=z-z^2/2+\cdots$. The $n!$ below compensates for
ordering the occurrences; it is not part of the definition
\eqref{ch08:ursell}.

\begin{theorem}[Standard polymer convergence criterion]
\label{ch08:kp}
Suppose $A_\gamma\ge0$ and
\begin{equation}\label{ch08:kp-condition}
 \sum_{\eta:\eta\not\sim\gamma}
       |z_\eta|e^{A_\eta}\le A_\gamma,
 \qquad \sum_\gamma |z_\gamma|e^{A_\gamma}<\infty.
\end{equation}
Here the first inequality is required for each $\gamma\in\mathcal P$.
Then the cluster series
\begin{equation}\label{ch08:log}
 \mathcal L(z)=\sum_{n\ge1}\frac1{n!}
       \sum_{\gamma_1,\ldots,\gamma_n}
        \varphi^T(\boldsymbol\gamma)\prod_{i=1}^n z_{\gamma_i}
\end{equation}
converges absolutely and $Z_{\rm pol}(z)=e^{\mathcal L(z)}$.
For a fixed label $\gamma_1$, its absolute rooted series satisfies
\begin{equation}\label{ch08:root}
 \sum_{n\ge1}\frac1{(n-1)!}
 \sum_{\gamma_2,\ldots,\gamma_n}
 |\varphi^T(\boldsymbol\gamma)|\prod_{i=2}^n|z_{\gamma_i}|
 \le e^{A_{\gamma_1}}.
\end{equation}
\end{theorem}

\begin{proof}
This is the hard-core specialization of the cited convergence theorem:
the Mayer function is $-1$ on incompatible pairs and zero otherwise.
There is no need to assume $Z_{\rm pol}>0$ before taking the logarithm.
The series defines the branch continued from zero activity. For real
activities it is real, and the exponential identity then gives
$Z_{\rm pol}>0$. For complex activities the conclusion is nonvanishing.
\end{proof}

The connected graph sum must be taken before its absolute value. Bounding
the sum of the absolute values of all connected graphs would discard the
cancellation on which convergence depends. The standard tree-graph
inequality is one way to establish \eqref{ch08:root}; we use the result
rather than repeat its combinatorial proof.

\section{A spatially local bound with a geometric tail}
\label{ch08:tail}

The site norm from \eqref{ch08:site-norm} gives a particularly transparent
criterion. Suppose that, for some $\alpha\ge0$, $\delta>0$ and
$\gamma_0>0$,
\begin{equation}\label{ch08:reserve}
 \|z\|_{\alpha+\delta}\le\gamma_0,
 \qquad 0<q<1,\qquad h\gamma_0\le q\delta.
\end{equation}
The additional decay $e^{-\delta c(X)}$ controls the sum over polymers
whose supports meet $X$. The parameter $q$ measures
smallness and controls the number of occurrences in a cluster.

Write $X(\boldsymbol\gamma)=\bigcup_i X_{\gamma_i}$. A direct
weighted application of Theorem~\ref{ch08:kp} yields, for every site $x$
and integer $k\ge1$,
\begin{equation}\label{ch08:weighted-tail}
 \sum_{n\ge k}\frac1{n!}
 \sum_{\boldsymbol\gamma:x\in X(\boldsymbol\gamma)}
 |\varphi^T(\boldsymbol\gamma)|
 \prod_{i=1}^n|z_{\gamma_i}|
 e^{\alpha c(X(\boldsymbol\gamma))}
 \le\gamma_0 q^{k-1}.
\end{equation}
This form is slightly stronger than the bound with an additional
$(1-q)^{-1}$ obtained by separately summing estimates at each occurrence
number. Either form is sufficient for the applications.

To prove \eqref{ch08:weighted-tail}, apply the rooted bound
\eqref{ch08:root} to the rescaled activities
$\widetilde z_\gamma=q^{-1}|z_\gamma|e^{\alpha c(X_\gamma)}$
and take $A_\gamma=\delta c(X_\gamma)$. By
\eqref{ch08:overlap} and \eqref{ch08:cover-properties}, the left side of
\eqref{ch08:kp-condition} is bounded by
$q^{-1}|X_\gamma|\gamma_0\le\delta c(X_\gamma)$.
For a nonzero connected coefficient the union weight is at most the
product of the support weights. Designating one occurrence whose support
contains $x$ changes $1/n!$ to $1/(n-1)!$ and gives an upper bound;
the rooted sum is at most $\gamma_0/q$. Restoring the factor $q^n$
on terms with $n\ge k$ gives \eqref{ch08:weighted-tail}. Nonempty supports
and finite $\Lambda$ verify the total summability required by the theorem.

More generally, let $\mathcal Q\subset\mathcal P$ be a distinguished
set of labels with
\[
 \sum_{\gamma\in\mathcal Q}|z_\gamma|
                  e^{(\alpha+\delta)c(X_\gamma)}\le b_{\mathcal Q}.
\]
Replacing the condition $x\in X(\boldsymbol\gamma)$ by
``at least one occurrence belongs to $\mathcal Q$'' in
\eqref{ch08:weighted-tail} gives $b_{\mathcal Q}q^{k-1}$ on the right.
This is useful when a cluster has to meet a prescribed insertion.
An insertion selected in this way is a distinguished occurrence in the
combinatorial expansion, not yet a renormalized quantum field.

For a family satisfying \eqref{ch08:reserve} with the same constants in
every volume, the right-hand side of \eqref{ch08:weighted-tail} is uniform
in $|\Lambda|$. Summing the $k=1$ estimate over $x$ bounds $|\mathcal L(z)|$ by
$|\Lambda|\gamma_0$. Thus extensive free-energy bounds coexist with
volume-independent estimates on connected terms meeting a fixed site.
It is the latter estimates that preserve locality in a renormalization
step.

\section{Deleting polymers and controlling local changes}
\label{ch08:deletion}

Recall that $Z_{\rm pol}$ is the gas partition function
\eqref{ch08:gas}, and $X_\gamma$ is the full support of label $\gamma$.
For a set $Y\subset\Lambda$, delete all labels whose supports meet $Y$:
put $z^Y_\gamma=z_\gamma$ if $X_\gamma\cap Y=\varnothing$ and zero
otherwise. Under \eqref{ch08:reserve}, both cluster series converge and
\begin{equation}\label{ch08:deletion-identity}
 \frac{Z_{\rm pol}(z^Y)}{Z_{\rm pol}(z)}
  =e^{-\mathcal L_Y(z)},\qquad
 \mathcal L_Y(z)=\sum_{n\ge1}\frac1{n!}
 \sum_{\substack{\boldsymbol\gamma\in\mathcal P^n\\
 X(\boldsymbol\gamma)\cap Y\ne\varnothing}}
 \varphi^T(\boldsymbol\gamma)\prod_{i=1}^n z_{\gamma_i}.
\end{equation}
This is simply subtraction of two absolutely convergent logarithms.
Its local bound is
\begin{equation}\label{ch08:deletion-bound}
 |\mathcal L_Y(z)|\le |Y|\gamma_0,
\end{equation}
and clusters with at least $k$ occurrences contribute at most
$|Y|\gamma_0q^{k-1}$. Counting a cluster more than once when it meets
several sites of $Y$ only increases the upper bound.

The same mechanism controls changes of activities in a bounded region.
If two families of activities agree away from $Y$, their logarithmic
cluster terms cancel unless a cluster contains a changed activity.
A quantitative Lipschitz estimate additionally needs a bound on the
difference of those activities, or a uniformly admissible interpolation
between them. Agreement of supports alone gives locality of the
difference, not its smallness.

\section{Parameters, derivatives and connected insertions}
\label{ch08:parameters}

Suppose $z_\gamma(u)$ is holomorphic in a complex parameter $u$ on a
disc and that \eqref{ch08:reserve} holds uniformly on each of its compact
subdiscs. The cluster series then converges normally there. Consequently
its sum is holomorphic and can be differentiated term by term. If the
bound for a local cluster sum is $M$ on $|u-u_0|<\rho$, Cauchy's estimate
gives
\begin{equation}\label{ch08:cauchy}
 |\partial_u^r\mathcal L_Y(u_0)|\le r!\rho^{-r}M.
\end{equation}
There is no separate combinatorial proof for each derivative order.
Several parameters are treated on a polydisc, with the product of the
corresponding Cauchy factors. Only derivatives of fixed finite order are
asserted unless additional uniform estimates are supplied.

One can introduce formal sources for bounded quantities $F_j$ by the
finite-volume generating function
\[
 Z(\boldsymbol t)=Z(0)\,
       \mathbb E\exp\!\left(\sum_{j=1}^r t_jF_j\right).
\]
The mixed derivative of $\log Z$ at zero is the connected correlation
(or cumulant) of the $F_j$. For example,
\[
 \left.\partial_{t_1}\partial_{t_2}\log Z(\boldsymbol t)
 \right|_{\boldsymbol t=0}
 =\mathbb E(F_1F_2)-\mathbb E(F_1)\mathbb E(F_2).
\]
The cluster representation expresses such derivatives through clusters
carrying the distinguished source occurrences; see
\cite[\S3]{Ueltschi2004}. To use this representation in the gauge theory,
the source-deformed activities must satisfy a uniform criterion on an
actual complex neighbourhood of zero. Boundedness of the undeformed gas
does not imply that estimate.

There is a second way to insert observables: differentiate the unnormalized
density and expand the result directly. This is preferable if some activity
can vanish, because division by that activity is then meaningless.
Bounds are assigned to the inserted activity itself. The usual mnemonic
``differentiate the logarithm of each activity'' is not a valid definition
in that situation.

\section{Large-field constraints and exact reassembly}
\label{ch08:forced}

An effective gauge density includes regions where a small-field expansion
is not available. Fix a configuration of the coarse variables for the
moment, and let $D\subset\Lambda$ be the set of sites declared to be in
the large-field region. In addition to the regular activities
$z_\gamma$, introduce large-field activities $K_\eta$ with supports
$Y_\eta$ and nonempty sets $I_\eta\subset D\cap Y_\eta$.
The set $I_\eta$ records the large-field sites assigned to that activity.
Let $\mathfrak C(D)$ consist of the allowed finite families $B$ of such
labels with disjoint supports and
\[
 \bigcup_{\eta\in B} I_\eta=D.
\]
The exact rules for eligibility are part of the definition of an effective
density; they may exclude some otherwise possible covers. When $D$ is
empty, the empty family is permitted. Figure~\ref{ch08:fig-exact-cover}
shows two covers of the same set \(D\) and how the exclusion rule can fail
even at a site outside \(D\).

% !TEX root = ../monograph.tex
\begin{figure}[htbp]
\centering
\begin{tikzpicture}[x=1cm,y=.70cm,font=\small,
  support/.style={draw=ymInk,line width=.85pt,rounded corners=2pt},
  lattice/.style={draw=ymGrid,thin}]
  \foreach \offset/\case/\titletext in
    {0/1/One support,4.4/2/Two disjoint supports,8.8/3/Excluded pair}{
    \begin{scope}[xshift=\offset cm]
      \ifcase\case\or
        \draw[support,fill=ymInk!7] (.67,.65) rectangle (2.33,2.35);
      \or
        \draw[support,fill=ymInk!7] (.67,.65) rectangle (1.33,2.35);
        \draw[support,fill=ymInk!7] (1.67,.65) rectangle (2.33,2.35);
      \or
        \draw[support,fill=ymInk!7] (.67,.65) -- (1.33,.65) --
          (1.33,2.68) -- (2.36,2.68) -- (2.36,3.36) --
          (.67,3.36) -- cycle;
        \draw[support] (1.73,.65) rectangle (2.27,3.25);
      \fi
      \foreach \x in {0,...,3}
        \draw[lattice] (\x,.44) -- (\x,3.48);
      \foreach \y in {1,2,3}
        \draw[lattice] (-.15,\y) -- (3.15,\y);
      \draw[ymAccent,line width=1.35pt] (1,1) rectangle (2,2);
      \foreach \x in {1,2}
        \foreach \y in {1,2}
          \fill (\x,\y) circle[radius=2.2pt];
      \ifnum\case=3
        \draw[ymInk,line width=.8pt,fill=white] (2,3) circle[radius=2.5pt];
      \fi
      \node at (1.5,-.04) {\titletext};
    \end{scope}
  }
\end{tikzpicture}
\caption[Exact covers of a large-field set]{Exact covers in a coordinate plane. One plaquette exceeds the
large-field threshold. Its four vertices (dots) form \(D\).
The outlined supports in either of the first two panels cover \(D\)
exactly, assuming the additional eligibility tests hold. In the middle
panel the indicators in both supports read the four links of the same
highlighted plaquette. In the last panel the prescribed large-field sets
still partition \(D\), but the supports share the circled site outside \(D\).}
\label{ch08:fig-exact-cover}
\end{figure}


For the simple overlap incompatibility rule, the constrained partition
sum has the form
\begin{equation}\label{ch08:constrained}
 Z_D=\sum_{B\in\mathfrak C(D)}
       \left(\prod_{\eta\in B}K_\eta\right)
       Z_{\rm pol}(z^{Y(B)}),\qquad
 Y(B)=\bigcup_{\eta\in B}Y_\eta.
\end{equation}
Applying \eqref{ch08:deletion-identity} gives the exact factorization
\begin{equation}\label{ch08:reassembly}
 Z_D=e^{\mathcal L(z)}
       \sum_{B\in\mathfrak C(D)}
         \left(\prod_{\eta\in B}K_\eta\right)
         e^{-\mathcal L_{Y(B)}(z)}.
\end{equation}
For finite sums this is algebra; the same identity holds for absolutely
convergent sums. It is the basic reason that removing regular polymers
near a large-field region can be absorbed into modified activities.

The expression $\mathcal L_{Y(B)}$ need not split as a sum over individual
members of $B$. A connected regular cluster can meet two large-field
supports. Such a cluster couples them, and the modified activity belongs
on their union together with the cluster support. Grouping by connected
components is the appropriate reorganization. An arbitrary product of
independently corrected large-field activities would omit these terms.

Formula \eqref{ch08:reassembly} does not by itself bound the sum over
$B$. That estimate requires suppression of large-field activities by the
action, bounds on the number and geometry of eligible supports, and enough
additional exponential support decay to absorb the regular clusters attached to them.
These are analytic inputs to the nonlinear renormalization step. The
abstract theorem bounds the regular cluster sum once those inputs have
been established; it does not establish them.

Two support notions must be distinguished in this application. The set
used in an exclusion or covering rule need not contain every link read
by an activity. For example, a large-field condition at a boundary site can
inspect a plaquette extending outside its assigned covering set. Disjoint
covering sets therefore need not imply disjoint dependence on links.
The abstract gas above needs the overlap implication
\eqref{ch08:overlap} for its chosen exclusion supports; it does not assume
statistical independence. When a local integration or factorization
argument also requires full dependence supports, those enlarged sets
and their geometric estimates must be specified separately. One must not
silently replace the covering sets by enlarged sets in an exact covering
identity.

\section{What an application must supply}
\label{ch08:application}

The reusable conclusions of this chapter are the connected expansion
\eqref{ch08:log}, the local tail \eqref{ch08:weighted-tail}, parameter
differentiation under a uniform complex-domain bound, and exact removal
and reassembly formulas. They are statements about a specified family
of activities. To apply them after a block integration one must give:
\begin{enumerate}
 \item an exact expansion of the integrated density, including the
 scalar factor and every large-field sector;
 \item supports containing every link and cutoff condition on which an
 activity depends, and a definite incompatibility relation;
 \item a site-norm bound with a strictly positive support-weight reserve,
 satisfying \eqref{ch08:reserve} uniformly in the period;
 \item corresponding estimates for the source-deformed or inserted
 activities whenever a connected observable is required;
 \item an absolutely summable bound for the constrained large-field
 families and the clusters attached to them.
\end{enumerate}
These requirements explain why the cluster expansion is not the main
obstruction in the construction. The general summation mechanism is
already available in the literature. What must be proved for the Wilson
model is that exact blocking, stationary integration, the large-field
decomposition and the chosen observable normalization produce activities
obeying those requirements on the same trajectory.

\setcounter{chapter}{8}
\IfFileExists{chapters/09-renormalization.tex}{\chapter{The nonlinear renormalization step}
\label{ch09:renormalization}

A block transformation acts on the density, not on its chosen list
of local terms. Our task is to express its exact output in the
coordinates of Chapter~\ref{ch04:blocking}, while improving the
bounds enough to repeat the operation. The organization by local
interactions and large-field regions follows the constructive
renormalization method of Ba{\l}aban
\cite{Balaban1988Expansions}. We use the cluster theorem of
Chapter~\ref{ch08:chapter}; the work here is the identification,
normalization and estimation of its concrete activities.

Fix one adjacent pair of scales, with fine scale $m+1$ and coarse
scale $m$. Write $L=L_m$, $B=B_m$, $t=t_m$, and
$\mathfrak b=\widehat\beta_{m+1}$. The actual incoming inverse coupling is
$\beta'\in[\mathfrak b/2,2\mathfrak b]$.
All event indicators use the fixed reference thresholds of
\eqref{ch04:thresholds}. They do not move when $\beta'$ varies.
Write the prescribed coarse path action as
\[
 P^{\rm path}_{N,m}(\beta,U)=\beta\mathcal K(U)+\mathcal L(U),
\]
where $\mathcal K$ includes its cubic and quartic terms and
$\mathcal L$ is the quadratic loop correction.
This identity defines $\mathcal K,\mathcal L$ from
\eqref{ch04:path-action}. Neither is the unknown effective action.

For this induction the derivative orders and spatial exponents
are stronger than the minimum allowed in the density definition:
take $q_{m+1}\ge q_m\ge6$ and
$\varpi_{m+1}\ge\varpi_m\ge2$, with the prescribed derivative
reserves for the local integral estimates. In particular the
normalized affine directions used below belong to the activity
seminorm. At the finest scale choose $0<c_{\rm dir}\le1$.

The estimates below distinguish proved algebraic implications
from the analytic inputs of Chapters~\ref{ch06:large-fields}
and~\ref{ch07:matching}. In particular a product estimate under
one Gaussian law, a local low-order matching estimate, and a
one-loop drift estimate are different inputs.

\section{The exact local integral and its first expansion}
\label{ch09:integral}

Let
$\mathscr Z_{\rm blk}=\mathcal B_m^{\beta'}\operatorname{dens}(\rho')$,
with the exact kernel \eqref{ch04:block-operator}.
In each sector one first performs the local coordinate change
on its active links and leaves the other links in probability
Haar measure. Chapter~\ref{ch06:large-fields} gives the exact
Gaussian identity, including the difference between the chosen
finite-range covariance and the inverse precision.

Here is the scalar bookkeeping in a form useful for checking
reassembly. Suppose there are $N_{\rm free}$ free tree-gauge
bonds, of which $n_{\rm act}$ use Lie algebra coordinates, and
let $D$ be the dimension after dummy Gaussian coordinates have
been added. Set $g=(\beta')^{-1/2}$, $q=\dim\gLie$, and
$c_{\rm H}$ equal to the density of probability Haar measure
in the chosen exponential coordinates at zero.
With $a=\mu+gz$, the measure multiplier before the Haar
logarithm is
\[
 (2\pi)^{-(D-n_{\rm act}q)/2}g^{n_{\rm act}q}c_{\rm H}^{n_{\rm act}}
 =(2\pi)^{-D/2}\mathfrak m^{N_{\rm free}}
                  e^{-(N_{\rm free}-n_{\rm act})\log\mathfrak m},
 \quad
 \mathfrak m=c_{\rm H}(2\pi g^2)^{q/2}.
\]
The dummy Gaussian integral equals one after the inactive Haar
variables have been frozen; Fubini then permits their integration.
The flat reference determinant uses the same forest as the
sector. It is extracted once, not separately for every eventual
connected component. The local precision correction, Haar
logarithm, missing-coordinate factor and stationary residual
remain among the local exponential factors.
Keeping them is what recovers the original Haar integral.

Expand exponentials as $e^V=1+(e^V-1)$ and retain all sector
allocation indicators. A complete dependence region contains
every Haar coordinate, tag variable and Gaussian covariance
halo read by its factor. If two such supports are disjoint,
the Haar and tag coordinates are independent; the Gaussian
cross covariance is zero because the chosen covariance has
finite range. Their integrals therefore factor.
Connected components of intersecting dependence regions define
the first polymer activities.

This factorization concerns dependence regions. It does not
assert that large-field activities with disjoint exclusion
supports $Y_i$ have disjoint full dependence regions.
The large-field indicators in their eligibility criteria can read outside $Y_i$.
Those conditions remain in every exact-cover expression.

The required simultaneous estimate for a connected product is
\begin{equation}
 \int G_M\,d\nu\le
 P_j(\mathfrak b)(1+s_M)^{d_j}
       e^{-r_j\beta't_{m+1}^2|A_M|}
       \prod_{a\in M}u_a .
 \label{ch09:joint-input}
\end{equation}
Here $\nu$ is the one sector Gaussian/Haar law,
$G_M$ majorizes the product and its prescribed ordered
derivatives, $s_M$ counts the full dependence region in
radius-$B$ blocks, and $A_M$ is its selected large-field condition set.
The function $P_j$ has subpower growth:
$\log P_j(\mathfrak b)/\log\mathfrak b\to0$ at fixed stage and
ratio. Distinct cutoff-failure labels each contribute their
own factor to the product. The bounded-jet estimate
\eqref{ch06:actual-jet-integral}, applied to the primitive factors
constructed in Chapter~\ref{ch06:large-fields}, proves this joint
bound on their common law. The alternative quadratic-tilt criterion
in Proposition~\ref{ch06:common-majorant} is not required here.
Separate bounds on $\int|F_a|\,d\nu$ cannot
replace this estimate.

The elementary power counting is as follows. The entries include
the prescribed finite number of derivatives.
\[
\begin{array}{l|l}
\text{factor after deterministic subtraction}&\text{majorant}\\ \hline
\text{stationary cubic remainder}&P_j\beta'g^3\\
\text{centered determinant variation}&P_j(t_{m+1}+g)\\
\text{Haar remainder}&P_j(t_{m+1}+g)^2\\
\text{inherited quadratic-loop variation}
   &P_j\Lambda_{m+1}^{\rm loop}(t_{m+1}+g)g\\
\text{covariance and stationary residual}
   &P_j\mathfrak b^{c_j}e^{-c'_jB}\\
\text{incoming regular activity}&P_j\epsilon_{m+1}^{\rm reg}\\
\text{cutoff failure at a distinct label}&P_j\mathfrak b^{-1/2}
\end{array}
\]
For example $\beta'g^3=g$, and the stationary mean
is $O(gM_{m+1}^{\rm cut})$, so a cubic Taylor term is bounded
by $g$ times a polynomial in $M_{m+1}^{\rm cut}$ and Gaussian
coordinates. The latter is estimated inside the common integral.
The other rows follow from the determinant inverse identity,
the second-order vanishing of the Haar correction, and
the window errors of Chapter~\ref{ch06:large-fields}.
The loop coefficient bound must grow subpolynomially at fixed ratio.

After increasing the inverse coupling threshold these rows are bounded
by $\mathfrak b^{-7/16}$, since a subpower factor can be bounded
by $\mathfrak b^{1/16}$. Read-count polynomials are absorbed by
additional exponential support decay. The first component expansion
and then Theorem~\ref{ch08:kp} give an absolutely convergent
logarithm of the ordinary gas without large-field sites. Components attached to
a prescribed large-field site set remain in the exact-cover sum.
The deletion identity \eqref{ch08:deletion-identity} accounts
for ordinary polymers excluded by a selected large-field region.
Repeated polymer labels in the logarithm are counted at every
occurrence, with the factor $1/n!$ of \eqref{ch08:log}.

All rearrangements can first be made for finite selections.
Under \eqref{ch09:joint-input} and the strict cluster criterion,
the absolute envelopes allow passage to the countable sums.
Thus this procedure gives an identity for $\mathscr Z_{\rm blk}$,
rather than an asymptotic expansion for its logarithm.

\section{Improving the regular remainder}
\label{ch09:improvement}

The first bound $\mathfrak b^{-7/16}$ is too large for the
target $\epsilon_m^{\rm reg}=\widehat\beta_m^{-11/10}$.
There are three different sources of improvement.

\subsection{Terms containing several elementary factors}

Give every elementary ordinary factor a degree one, and give
a first-expansion polymer the sum of its factors' degrees.
A cluster degree is the sum over its polymer occurrences.
Multiply each elementary majorant by
$h=\mathfrak b^{3/8}$ before performing either expansion.
The resulting local smallness parameter is
$h\mathfrak b^{-7/16}=\mathfrak b^{-1/16}$.
The weighted rooted bound of
Section~\ref{ch08:tail}, applied with these inflated activities,
bounds the whole sum by $P_jh\mathfrak b^{-7/16}$.
Restoring the original weights for terms of total degree at
least three gives
\begin{equation}
 P_jh^{-3}(h\mathfrak b^{-7/16})
       =P_j\mathfrak b^{-19/16}
       \le\mathfrak b^{-23/20}.
 \label{ch09:degree-gain}
\end{equation}
The gap $19/16-23/20=3/80$ absorbs the remaining subpower factor.
This counts elementary factors, not just vertices in the
second cluster expansion; a single polymer may already have
degree three. The same calculation applies after placing a
fixed number of derivatives.

An incoming regular activity accompanied by a further factor
is bounded by the rooted estimate with that incoming activity
distinguished:
$P_j\epsilon_{m+1}^{\rm reg}\mathfrak b^{-1/2}$.
Long path terms have an extra factor $e^{-c_jB}$.
Both fit the right side of \eqref{ch09:degree-gain} after
increasing the threshold.

\subsection{The linear inherited regular term}

The component linear in an incoming regular activity must use
its affine-Hessian cancellation. In local radial gauge, the
lifted potential splits into parts with inverse-ratio sizes
$L^{-1},L^{-2},L^{-3}$: a constant part, its affine variation,
and its remainder. The constant part is removed by a local
gauge transformation. Curvature therefore begins at weight
two. Invariance under constant gauge transformations remove linear curvature terms,
and the incoming affine-Hessian condition removes the
quadratic curvature term of weight four. Every surviving
selection in the ordered Taylor expansion has weight at least
five.

Here is the support estimate needed to sum that local
calculation. Let $d$ be the incoming covering number.
For $d^3<L$ the inspected support admits a nonwrapping
radial gauge and the per-label gain is $L^{-5}$.
For the other supports use the original bound. Both cases
are bounded by
\[
 \delta_L(d)\le L^{-5}(1+d)^{15},\qquad
 \delta_L(d)=
 \begin{cases}L^{-5},&d^3<L,\\1,&d^3\ge L.\end{cases}
\]
The full chain rule adds a fixed polynomial in $d$.
Choose the scale ratio and incoming support exponent so that, after
support compression, the outgoing weight uses at most half of the
incoming exponential decay. The remaining decay absorbs the polynomial
in $d$.
There are $L^4$ fine anchors above each coarse anchor.
Consequently the inherited principal norm is at most
$V_jL^{-1}\epsilon_{m+1}^{\rm reg}$, with $V_j$ fixed before
$L$.

For the comparison with the actual integral use the same forest
and physical local coordinates on both sides. The estimates of
Section~\ref{ch06:linear-comparison} give a covariance difference
$O(P(B)t_f)$ and a mean difference, in Gaussian units, of order
$P(B)g(M_f^{\rm cut})^2+P(B)M_f^{\rm cut}e^{-cB}$,
where $t_f=t_{m+1}$ and $g=(\beta')^{-1/2}$. Each is
$\mathfrak b^{-1/2+o(1)}$. To see where the gain occurs, for
a covariance interpolation $C_s$ with fixed positive lower bound,
the logarithmic density derivative is
\[
 \frac12\left\{z^*C_s^{-1}\dot C_s C_s^{-1}z
                      -\tr(C_s^{-1}\dot C_s)\right\}.
\]
A varying mean adds its linear Gaussian score. Every coefficient
contains a covariance or normalized-mean difference. Its fixed
moments therefore contain the small factor just stated, times a
polynomial in the local dimension.

Hold the physical links in the incoming activity fixed by the
protected coordinate change of Chapter~\ref{ch06:large-fields}.
The nonlinear coordinate correction is then differentiated through
its Jacobian and density, where the same gain is available.
This avoids requiring one extra derivative of the incoming activity
when adjacent levels have equal derivative orders. Telescope all
remaining factor products, marking one changed factor. The
nonquadratic local weights contribute a gain of order $g$ with
subpower losses; a second incoming activity contributes
$O((\epsilon_{m+1}^{\rm reg})^2)$ instead. Common Gaussian moments
and the unused support exponent sum the resulting polynomial costs.
Together with the principal estimate this proves
\[
 \mathcal E_{\sigma+3}(\text{inherited term})
 \le \left(V_jL^{-1}+
       \omega_{j,L}(\widehat\beta_m)\right)
                    \epsilon_{m+1}^{\rm reg},
 \qquad
 \omega_{j,L}(\widehat\beta_m)
       =\widehat\beta_m^{-1/2+o(1)}.
\]
For this comparison the local windows and their dependence halos
must be clean. A window meeting a selected large-field component is
assigned to that component; a failed remote cutoff has the joint
Gaussian-tail bound of Chapter~\ref{ch06:large-fields}. Large
remaining supports gain $e^{-cB}$ from the unused incoming support
weight. Thus no comparison between unrelated event indicators or a seeded and
an unseeded law has been made. The actual interpolation estimates
are essential; the scaling-dimension count alone would not prove
the comparison remainder.

\subsection{Low-order terms and localization}

One- and two-factor terms must be combined before estimating them.
The degree-four classical matching and quadratic determinant matching
must apply to the same integral and stationary branch. Rather than
use a formal cumulant remainder, we retain a finite stationary
region and estimate its complement. The resulting centered
auxiliary bound, with $h=\mathfrak b^{-7/16}$, is
\begin{equation}
 P_j(\mathfrak b)\left\{
 \mathfrak b^{-1}h+\mathfrak b h^5+h^3\right\}
                         +\mathfrak b^{-K},
 \label{ch09:aux-input}
\end{equation}
where $K$ can be any fixed positive number. Its largest power is
$\mathfrak b^{-19/16+o(1)}$, which is smaller than
$\mathfrak b^{-23/20}$. The argument below specifies the actual
sector and stationary estimates needed to obtain this bound.

% !TEX root = ../monograph.tex
\subsection{The stationary integral with its cutoff retained}
\label{ch09:stationary-remainder}

We now give the local analytic implication behind the auxiliary bound.
It is an application of Laplace's method, for which we refer to
\cite[Chapter II]{Wong2001}, with three features that must be checked
here: the integration dimension grows with the auxiliary period, the
Hessian at the critical configuration depends on the coarse field, and the integration
cutoff cannot be discarded before estimating its complement.
All actions and matrices in the following argument belong to the same
\emph{pure} small-field block integral, the block integral of the path action alone: incoming activities are omitted, the
prescribed path action and block kernel are unchanged, and the fine-field
indicator $\mathbf1_{\{\chi(U_p)\le t_f\ \text{for every fine }p\}}$
is retained, where $t_f=t_{m+1}$. Thus its integrand is zero on the fine large-field
complement. The other sectors at the fine scale have already been attributed to the
large-field expansion.

Put $z=\widehat\beta_{m+1}$, $g=\beta'^{-1/2}$ and
$h=z^{-7/16}$. The letter $z$ is a scalar inverse coupling parameter in this
subsection; the Gaussian integration variable is $\xi$. Auxiliary
periods and coordinate counts are bounded by fixed powers of
$1+\log z$. Write $M=M_{m+1}^{\rm cut}$ for the cutoff amplitude
in \eqref{ch04:thresholds}, and take $D=d_0M$ with a fixed small
$d_0>0$. We shall use the following two distinct facts:
\[
 M=z^{o(1)},\qquad (1+\log z)^k=o(M^\epsilon)
 \quad\text{for every fixed }k,\epsilon>0.
\]
Both follow directly by taking logarithms of
$M=\exp((\log\log z)^2)$.

Here and below $z^{o(1)}$ in an upper bound means that for every
$\epsilon>0$ the quantity is $O_\epsilon(z^\epsilon)$ as $z\to\infty$
at fixed stage and block ratio.

Let $x$ be the coarse logarithmic coordinates, with their supremum
norm, and $y\in\RR^n$ the free fine coordinates after fixing the block axial
gauge. The classical action including the penalty term of the block kernel is
$\phi(x,y)$; write $\rho(x)$ for its critical configuration and
$A(x)=D_y^2\phi(x,\rho(x))$ for its positive Hessian.
The extra exponent $\psi(x,y)$ includes the fine loop term and minus
the logarithm of the product Haar-chart Jacobian. Suppose the critical-configuration map $\rho$,
action, extra exponent and relative logarithmic determinant have the
fixed derivative orders required below on a common neighborhood,
with bounds $z^{o(1)}$. The Hessian gap is uniform in the auxiliary
volume at the fixed block ratio. These are the stationary and path
estimates of Chapters~\ref{ch05:gaussian}--\ref{ch07:matching}; no
different flat covariance is substituted here.

For a fine-coordinate displacement $g\xi$, with $\xi\in\RR^n$, set
$s(x)=\tfrac12\log\det A(x)$ and define the second-order Taylor remainder
divided by $g^2$ through its integral formula
\[
 V(g,x,\xi)=\int_0^1(1-t)
 D_y^2\phi(x,\rho(x)+tg\xi)[\xi,\xi]\,dt.
\]
Stationarity gives
$\phi(x,\rho(x)+g\xi)-\phi(x,\rho(x))=g^2V(g,x,\xi)$,
including the limiting value
$V(0,x,\xi)=\tfrac12\xi^*A(x)\xi$.
With $\gamma_0=N(0,A(0)^{-1})$, put
\begin{align}
 a(g,x,\xi)&=s(x)-s(0)+V(0,0,\xi)-V(g,x,\xi)\notag\\
 &\quad+\psi(x,\rho(x))-\psi(x,\rho(x)+g\xi),\notag\\
 R(g,x)&=\int_{|\xi|\le D}e^{a(g,x,\xi)}\,d\gamma_0(\xi),
 \qquad L(g,x)=\log R(g,x).
 \label{ch09:truncated-integral}
\end{align}
The determinant factor is essential: at $g=0$ the exact identity is
\[
 R(0,x)=N(0,A(x)^{-1})\{\,|\xi|\le D\,\},
\]
which is the Gaussian measure of a ball, not identically one. Changing $\xi$ to
$-\xi$ also gives the exact parity $R(-g,x)=R(g,x)$.

Here are the uniform bounds needed for this elementary application
of Laplace's method. Differentiating the divided difference on
$|\xi|\le D$ costs a fixed polynomial in $D$ and the stationary
derivative bounds. It is therefore $z^{o(1)}$ at every fixed order.
The exponent vanishes at $(g,x)=(0,0)$, so on
$|g|,\|x\|_\infty\le h$ it is at most one in absolute value for
large $z$. If the reference Gaussian measure of the ball is at least $1/2$, then
$e^{-1}/2\le R\le e$, and the same derivative bounds $z^{o(1)}$ hold
for $L$. To obtain derivatives through order $r$ and the two
derivatives in $g$ used for parity, it suffices to retain derivatives
of $\phi,\rho$ through order $r+4$ and of the other terms through
order $r+2$. These are finite-stage requirements.

If $A(x)^{-1}\le cI$ and $n$ is at most a fixed polynomial in
$\log z$, a coordinate Gaussian tail bound gives
\[
 1-R(0,x)\le 2n\exp\bigl(-D^2/(2cn)\bigr).
\]
This is smaller than every fixed inverse power of $z$.
The same conclusion for its fixed derivatives follows either by
the Gaussian score formula or by finite-order interpolation on a
ball of radius $z^{-1}$ inside the parameter domain. We recall the
needed interpolation fact: a superpolynomially small value and
polynomially bounded derivatives through order $r+1$ imply
superpolynomially small derivatives through order $r$ on a smaller
ball. It follows successively from Taylor's estimate
$\|DF\|\le2\sup|F|/u+u\sup\|D^2F\|$ on a radius-$u$ ball,
applied to successive derivatives with shrinking radii.
No coordinate-dimension factor is needed when operator norms are
used. Parity and Taylor's theorem in $g$ now yield, for every fixed
$K>0$,
\begin{equation}
 \|D_x^r L(g,x)\|\le C_r(z)\beta'^{-1}+z^{-K},
 \qquad C_r(z)=z^{o(1)}.
 \label{ch09:even-gain}
\end{equation}
This includes $r=0$; the cutoff error at $g=0$ has been retained
throughout the argument.

\subsection{Identifying the retained region in the pure block integral}
\label{ch09:stationary-exterior}

Write $\mathcal D_x(y)$ for the configuration map of the block axial gauge.
Consider the region
$\mathcal S_x=\mathcal D_x(\rho(x)+g\overline B(0,D))$.
Let $\mathscr Z(x)$ be the pure block integral, and set
\[
 \Phi(x)=\beta'\phi(x,\rho(x))+\psi(x,\rho(x))
                         +\tfrac12\log\det A(x).
\]
If $N$ is the number of free fine links and $c_H$ the Haar density
constant, its normalization constant is
$\Omega=\pi_{\beta'}(\sqrt{2\pi}g)^n c_H^N$;
$\pi_{\beta'}$ is the block-kernel normalizing prefactor.
When $\mathcal S_x$ lies in the chart, all its fine plaquettes satisfy
the regularity condition $\chi(U_p)\le t_f$, and the actions for all values of the auxiliary index with positive probability weight coincide with $\phi$, the
change of variables gives exactly
\begin{equation}
 \int_{\mathcal S_x}\text{(pure integrand)}
                  =\Omega e^{-\Phi(x)}R(g,x).
 \label{ch09:retained-identity}
\end{equation}
The probability weights of the auxiliary index sum to one, so there is no additional factor counting its values.
The hypotheses about this region must be established before
using the identity.

Here is a verification using the geometry and the Hessian gap.
By the estimate \eqref{ch05:nonlinear-tree} in the block axial gauge, a fine
configuration lies in a fixed closed coordinate ball of the chart\mgubi{Issue 10: in which norm? existence/ takes the supremum ball and uses its convexity.}, the \emph{safe chart}, whenever its fine
plaquettes and the deviations of the coarse links from the values of the averaging operation are below a fixed $\delta>0$.
On that chart assume $D_y^2\phi\ge\gamma I/2$. Stationarity gives
\[
 \phi(x,y)-\phi(x,\rho(x))
                      \ge\gamma|y-\rho(x)|^2/4.
\]
Thus a chart point outside $\mathcal S_x$ adds at least
$\Delta=\gamma D^2/4$ to the exponent. A point outside the chart
with all fine plaquettes small must have such a deviation exceeding
$\delta$, and contributes a fixed positive multiple of $\beta'\delta^2$.
The possibly negative cubic, quartic and loop terms are bounded by
the path coefficient bounds. If $W$ bounds the auxiliary coordinate and site
counts, their total contribution to the exponent is at most
\[
 W B_{\rm loop},\qquad
 B_{\rm loop}=\beta'C(t_f^3+t_f^4)+\Lambda t_f^2,
 \qquad 0\le\Lambda\le2M.
\]
All such terms and the logarithmic Haar/determinant normalization
term are $o(M^2)$. In particular, the normalization term
$-N\log[c_H(\sqrt{2\pi}g)^{\dim\gLie}]$ is only a polynomial
in $\log z$. The bound on $\phi(x,\rho(x))$ on $\|x\|<h$
also leaves the fixed penalty for a deviation exceeding $\delta$ available, because
that value tends to zero and $\Delta/\beta'\to0$.
Consequently the exterior integrand is bounded by
$\Omega e^{-\Phi(x)}e^{-\Delta+o(M^2)}$.

It remains to verify the regularity condition on the fine plaquettes on the retained ball. For the
auxiliary configurations used here, which are reflected periodic extensions (pullbacks of the original field to an auxiliary torus under the coordinatewise reflection map $\pi_v$ of Section~\ref{ch09:local-density}), this can be done directly
from the induced map on curvatures, without a lower bound from a sector
expansion. A full coordinate cycle of a reflected periodic extension traverses an
original path and then its inverse. Its holonomy is therefore exactly
one at every transverse base point. Equivalently, choose a coordinate
axial gauge on the original cube and pull that gauge back under the
reflection map. Plaquette filling bounds give logarithmic coordinates
with
\[
 \|x\|_\infty\le P(B)t_c,\qquad
 \|d_{1,c}x\|_\infty\le C t_c+P(B)t_c^2,
 \qquad t_c=t_m.
\]
The second estimate follows by expanding each plaquette product in
these small link coordinates. Crucially, its leading constant is
independent of $B$; the polynomial loss occurs only in the quadratic
error. This argument would fail for an arbitrary small-curvature
periodic field with nontrivial winding holonomies.

Let $\mathcal U(x)=\mathcal D_x(\rho(x))$ be the critical configuration as a physical
fine field. Its linearization is the representative $W_F$ in the block axial gauge of
Chapter~\ref{ch05:gaussian}. Its curl agrees with that of the smooth
representative $W_a$. The curvature identity and the
fine-row bound of the induced map $T$ on curvatures give
\[
 d_{1,f}W_a=H_f^{-1}TH_c d_{1,c},\qquad
 \|H_f^{-1}TH_c\|_{\infty\to\infty}\le C_jL^{-2}.
\]
The operator norm here sums each fine-output row over coarse
inputs; the opposite column sum has a different scaling.
The quadratic remainder of the critical-configuration map and the plaquette
product is bounded using their common second derivatives. Thus
\begin{equation}
 \sup_p\chi(\mathcal U(x)_p)
       \le C_jL^{-2}t_c+P_{j,L}(B)t_c^2.
 \label{ch09:stationary-cutoff-margin}
\end{equation}
Choose $L$ so that $C_jL^{-2}t_c\le t_f/8$, using the uniform
bound $t_c/t_f\le2$. Then increase the lower bound on $z$ so that
the second term is at most $t_f/8$. These choices are in the
required order: the constant in the first term is fixed before
$L$, while the polynomial in the second term may depend on $L$.

Plaquette functions have a fixed Lipschitz constant $L_0$ in the
physical fine-link supremum coordinates. Local chart maps have
fixed-stage derivative bounds. Consequently variation of $x$ by
$z^{-1}$ and of the fine coordinates at the critical configuration by $g\xi$,
$|\xi|\le D$, changes every plaquette distance from the identity by at most
$C_{j,L}z^{-1}+L_0gD$, after incorporating the fixed chart norm
conversion into $L_0$. Choose $d_0$ sufficiently small, so that
$L_0gD\le t_f/8$, and then raise the lower bound on $z$ for the term
$C_{j,L}z^{-1}$. Every point of each retained region now satisfies the regularity condition on the fine plaquettes,
with a strict margin. The common chart and the ball on which the actions for all values of the auxiliary index with positive probability weight coincide
contain these regions because $h+gD\to0$.
We have therefore obtained, on a full microscopic neighborhood of
each tested coarse configuration satisfying the regularity condition,
\begin{equation}
 0\le\mathscr Z(y)-\Omega e^{-\Phi(y)}R(g,y)
       \le\Omega e^{-\Phi(y)}e^{-c_0M^2},\qquad c_0>0.
 \label{ch09:exterior-bound}
\end{equation}
This deduction uses the linearization of the critical configuration (through its smooth representative $W_a$), its second
derivatives, the compatible path bounds and the geometric
estimate \eqref{ch05:nonlinear-tree}. The direct proof of the regularity condition applies to the
auxiliary fields that are reflected periodic extensions and to their microscopic neighborhoods, which
are precisely the configurations used for localization.

\subsection{Centering the full logarithm}
\label{ch09:centered-log}

Put $E=\log\mathscr Z-\log\Omega+\Phi-L(g,\cdot)$.
Since $R\ge e^{-1}/2$, \eqref{ch09:exterior-bound} gives
$0\le E\le2e\,e^{-c_0M^2}$. To pass to derivatives, write
the positive pure integrand at a fixed fine field and a fixed value of the auxiliary index as
$a(\omega)e^{p(x,\omega)}$, with $a\ge0$ independent of $x$.
Only the coarse-link kernel is differentiated. Its exponent has
fixed-order derivatives bounded by $C_r\beta' n_b$, where
$n_b$ is the number of coarse links. The exponential chain rule
and integration give relative bounds
\[
 \|D^r\mathscr Z(x)\|\le P_r(\beta'n_b)\mathscr Z(x).
\]
The logarithm chain rule thus gives polynomial bounds without
dividing an absolute error by a possibly tiny partition function.
The other terms in $E$ have polynomial derivative bounds from the
stationary estimates. Interpolation on the microscopic neighborhoods
in \eqref{ch09:exterior-bound} proves that every fixed jet of $E$
is smaller than every inverse power of $z$.

Let $a_c=\mathcal K$ be the coarse classical action and
$d_c=\mathcal L+\alpha\mathcal Q$ the prescribed coarse quantum contribution:
the inherited loop action of the coarse path action and its quadratic part
multiplied by the coefficient $\alpha$ fixed by the matching, as in
\eqref{ch09:early}.\mgubi{Identified from the Lean definition YM.\allowbreak Precision.\allowbreak draft ($\alpha$ times the quadratic term plus the loop term), existence/\allowbreak stationary-output-construction.\allowbreak tex and v2 D-section; no $\Delta$ and no constant per site. Open: Lean uses $\mathcal K$, $\mathcal L$, $\mathcal Q$ without event indicators, while v3 defines them with indicators; they coincide on the small-field ball.} Define
\[
 C=a_c-\phi(\cdot,\rho),\qquad
 W=d_c-\psi(\cdot,\rho)-\tfrac12\log\det A.
\]
The corrected logarithm has the exact decomposition
\begin{equation}
 \log\mathscr Z+\beta'a_c+d_c
                  =\log\Omega+L+E+\beta'C+W.
 \label{ch09:pure-decomposition}
\end{equation}
Classical matching gives $D^rC(0)=0$ for $0\le r\le4$.
The quadratic quantum matching bounds $D^2W(0)$ by an
exponentially small finite-period error, with polynomial coordinate
norm costs. Its higher fixed derivatives are $z^{o(1)}$.
The full corrected logarithm has zero first derivative at zero
by constant gauge invariance and simplicity of $\gLie$.
Thus $DW(0)=-DL(0)-DE(0)$; no separate parity claim about each
determinant contribution is needed.

For a function on the radius-$h$ ball write its scaled derivative
bound through order $r$ as
$\max_{0\le k\le r}h^k\sup\|D^kf\|$.
Taylor's theorem, \eqref{ch09:even-gain}, and the two matching
cancellations show that the left side of
\eqref{ch09:pure-decomposition} minus its value at zero has this bound at most
\begin{equation}
 C_r(z)\bigl(\beta'^{-1}h+\beta'h^5+h^3\bigr)+z^{-K},
 \qquad C_r(z)=z^{o(1)},
 \label{ch09:pure-small-bound}
\end{equation}
where $K$ is arbitrary and the exponentially small quadratic
finite-period error has been included in the last term.
For $h=z^{-7/16}$ the three powers are respectively
$z^{-23/16}$, $z^{-19/16}$ and $z^{-21/16}$.
Their largest exponent is $-19/16<-23/20$.
Original small-field logarithmic coordinates and the directions admitted by the activity
seminorm \eqref{ch04:local-seminorm}, restricted to the cube $N_{\bar w}$ about the localization site and carried to the reflected periodic extension ($\bar w$ and the reflected extension are those of Section~\ref{ch09:local-density} below; the cube $N_r(x)$ is used there without a definition), are bounded by a polynomial in $B$ times
$t_m$, which is $o(h)$.\mgubi{All directions of the seminorm are meant, not only the affine ones (v2, D-section). v2 adds that the reflected torus must admit the charts and derivative bounds used here, and says that this and the finite-period comparison on a common analytic strip ``remain under verification''; v3 does not prove it either.} They therefore obey the same estimate.
This proves the auxiliary remainder needed for localization from
the explicit stationary, matching and actual-sector inputs above.
It does not require a formal expansion of an untruncated
exponential of a cubic Gaussian polynomial.


\subsection{Recovering a local density}

The passage from auxiliary sums to local activities has an
exact finite-dimensional explanation. Truncate the selected
low-order formula, retaining a weighted tail
$\varepsilon_{\rm tr}$. Give an ordinary selected term $f_a^0$
the shares $p_{xa}=\mathbf1_{X_a}(x)/|X_a|$, and a path term
the share $\mathbf1_{\{o_a\}}(x)$. These shares sum to one.
Set
\[
 F_x=\sum_{a\ {\rm selected}}p_{xa}f_a^0,\qquad
 M_x=\sum_{a\ {\rm selected}}p_{xa}h_{X_a}f_a^0.
\]
Choose a common dependence radius $R$ about every site receiving
a nonzero share. Set $w=4R+100(1+B)$, $\bar w=w+5$, and use
auxiliary tori $\Lambda_i$ of sides $2(w+i)$, $0\le i\le4$.
The relevant reflection cubes must embed in the original torus.
For a cyclic coordinate $q$ of period $2v$, define
$\pi_v(q)=\min(\bar q,2v-\bar q)$, where $0\le\bar q<2v$.
Apply this map coordinatewise. Write $U_{i,z}$ for the pullback of
the original field under $q\mapsto z+\pi_{w+i}(q)$ on $\Lambda_i$.
Reversed links are inverted, so small plaquette distances from the identity stay small.
Let $F_y^{(i)}$ denote the same local formula on $\Lambda_i$
with its flat value subtracted, and set
\begin{align}
 \lambda_i&=\frac{(-1)^i\binom4i}{2^4\,4!},&
 G_x&=\sum_{i=0}^4\lambda_i\sum_{y\in\Lambda_i}
                              F_y^{(i)}(U_{i,x}),\notag\\
 L_x&=M_x-F_x(1)-
   \sum_{i=0}^4\lambda_i\sum_{y\in\Lambda_i}
       h_{N_{\bar w}(x-\pi_{w+i}(y))}
                      F_y^{(i)}(U_{i,x-\pi_{w+i}(y)}).
 \label{ch09:localization}
\end{align}
In the second line the original cube is based at
$x-\pi_{w+i}(y)$, so its site
$y$ corresponds to the original site $x$.
This dependence on $y$ is essential.

When $N_{2\bar w}(x)$ is small-field, all these event indicators are one.
Classify sites of $\Lambda_i$ by their distances, through $R$,
from the reflection faces. With the recentering just specified,
the local expression in each class is independent of $i$.
The multiplicity of each class is polynomial in $i$ of
degree at most four; only the all-interior class has degree
four, with leading coefficient $2^4$.
The fourth-difference weights therefore give
$\sum_i\lambda_i\sum_y F_y^{(i)}(U_{i,x-\pi_{w+i}(y)})
=F_x(U)-F_x(1)$, and $L_x=0$.
This is not a pointwise assertion that $G_x=F_x-F_x(1)$.

Sum \eqref{ch09:localization} over $x$ and change variables
$x-\pi_{w+i}(y)\mapsto x$ in each finite auxiliary sum.
Translation covariance makes $F_x(1)=F_0(1)$, giving exactly
\[
 \sum_x M_x=|\Lambda|F_0(1)+
              \sum_x h_{N_{\bar w}(x)}G_x+\sum_xL_x .
\]
Thus $G_x$ is the small-field local activity and $L_x$ is a
logarithmic correction supported where its enlarged
neighborhood contains a large-field site.

Auxiliary volumes and reflected derivative norms grow only
polynomially in $B$. Thus \eqref{ch09:aux-input} implies
\[
 \|G_x\|_m\le P_j(\mathfrak b)
           (\mathfrak b^{-23/20}+\varepsilon_{\rm tr}),
 \qquad \sup|L_x|\le P_j(\mathfrak b)\mathfrak b^{-7/16}.
\]
The first fits an arbitrarily small fraction of
$\epsilon_m^{\rm reg}$, since $23/20>11/10$.
The second tends to zero and is incorporated into
large-field activities by the exact operation below.

The selected low-order terms contain no incoming regular or
large-field activity. Their coefficients depend only on the
prescribed path, actual inverse coupling and fixed labels. Their
weighted continuity on the compact inverse coupling band therefore
permits one truncation for all inputs with those data.
Terms involving incoming activities have already been
estimated separately. More generally compactness along a
single activity curve would give a truncation only along
that curve. All choices here are retained under restriction
of the parameter interval.

\section{A precise preparatory output}
\label{ch09:preparatory}

The preceding estimates, when their concrete local inputs
hold, give local functions $f_i$, preliminary large-field
activities $k_i$, and a scalar $\tau$ satisfying
\begin{equation}
 \mathscr Z_{\rm blk}(U)=
 e^{|\Lambda|\tau-\beta'\mathcal K(U)-\mathcal L(U)
                       -\alpha\mathcal Q(U)-\sum_i h_{X_i}f_i(U)}
                         \mathcal F^0(U).
 \label{ch09:early}
\end{equation}
Here $\mathcal F^0$ is the exact-cover sum of the $k_i$,
and $\mathcal Q$ is the quadratic part, restricted by its event indicators, of the
prescribed path action. The coefficient $\alpha$ is
the matched determinant coefficient, not an adjustable
trial inverse coupling. The bounds needed below are
\begin{equation}
 \sup_x\sum_{o_i=x}e^{(\sigma+3)c(X_i)}a_i\le H_{\rm pre},
 \quad a_i\ge\|f_i\|_m,\qquad
 \sup_x\sum_{i:x\in Y_i}e^{(\sigma+1)c(Y_i)}k_i^{\rm maj}
       \le e^{-c_j(M_{m+1}^{\rm cut})^2},
 \label{ch09:early-bounds}
\end{equation}
where $c(X)=\operatorname{cov}_{B_m}(X)$ and
$H_{\rm pre}=\theta c_j^{\rm pre}\epsilon_m^{\rm reg}$,
$0<\theta\le1$. The positive stage constant $c_j^{\rm pre}$
is fixed from the subsequent subtraction and path-replacement
operator norms before choosing the ratio.

To obtain these bounds choose $L$ first so that
$V_jL^{-1}\epsilon_{m+1}^{\rm reg}$ fits the desired fraction of
$H_{\rm pre}$. This is a condition on the fixed stage constants:
the ratio $\epsilon_{m+1}^{\rm reg}/\epsilon_m^{\rm reg}$ is bounded
on the permitted neighboring reference-inverse-coupling bands.
Then choose the inverse coupling threshold for the inherited
comparison, \eqref{ch09:degree-gain},
\eqref{ch09:aux-input}, and the small logarithmic corrections.
Choose the truncation accuracy after its polynomial
localization costs have been specified.
Allocate distinct positive fractions of the exponent in
\eqref{ch09:joint-input} to sector multiplicities, local
normalization losses, and the final large-field bound.
Polynomial support costs use support-weight reserves.
An unexplained extensive normalization cannot be absorbed into
a single selected large-field condition.

Equation \eqref{ch09:early} and its bounds are the concrete
interface for the rest of the chapter. The remaining
operations and estimates will now be proved directly from it.

\section{Local quadratic subtraction}
\label{ch09:subtraction}

For $a<b$ define the clover curvature at $o$ by
\[
 F^{\rm cl}_{ab}(U,o)=\tfrac14\sum_{\varepsilon,\eta=\pm1}
 \varepsilon\eta\,\operatorname{proj}_{\gLie}
       \frac{U_{o;\varepsilon a,\eta b}
                         -U_{o;\varepsilon a,\eta b}^{-1}}2 .
\]
The loop follows the four signed edges and starts at $o$.
This field transforms by $\Ad$ at $o$ and reads a radius-two
neighborhood. For an affine potential
$A_\mu(y)=u_\mu+\sum_\nu y^\nu v_{\mu\nu}$ its derivative
at the identity is $v_{ba}-v_{ab}$.

For a preliminary $f$ on $X$ anchored at $o$, put
$Z=X\cup N_2(o)$. On a nonwrapping $Z$, let
$A^{q,T}_c(y)=\mathbf1_{\{c=b\}}\xi_o(y)^aT$ for $q=(a,b)$,
and fix a unit $T_*\in\gLie$. Define
\[
 q_f(q,r)=D^2[f(\exp\cdot)](0)[A^{q,T_*},A^{r,T_*}],
 \quad
 Q_f(U)=\tfrac12\sum_{q,r}q_f(q,r)
       \langle F^{\rm cl}_q(U,o),F^{\rm cl}_r(U,o)\rangle.
\]
Set $q_f=0$ if $Z$ wraps. Gauge invariance makes the first
derivative at the identity vanish and the Hessian annihilate
local gradients. The zero-curl part of an affine potential
is a discrete gradient of a quadratic site function.
It can be extended periodically using a cutoff outside
the dependence region. Therefore the Hessian on affine fields
depends only on their constant curls. Invariance under the adjoint action and simplicity of $\gLie$ make its
bilinear form a multiple of the invariant metric.
Thus $Q_f$ has the same affine Hessian as $f$.
This is the local Ward-identity argument of
Section~\ref{ch07:kernels} in finite coordinates.

Set
\[
 R_f=f-f(1)-Q_f,\qquad
 \zeta_f=(1-h_Z)\{f(1)+h_{\{o\}}Q_f-h_Xf\}.
\]
An exact two-case identity is
\begin{equation}
 h_Xf=f(1)+h_{\{o\}}Q_f-\zeta_f+h_ZR_f .
 \label{ch09:subtract-identity}
\end{equation}
On a small-field $Z$ all event indicators are one; otherwise $h_Z=0$
and the definition of $\zeta_f$ proves the equality.
The remainder vanishes at the flat field and has zero
affine Hessian on nonwrapping supports.
The function $\zeta_f$ vanishes on small-field $Z$.
Its large-field indicator dependence must still be recorded in its full
dependence region; no derivative of an event indicator is taken.

The evaluations in normalized affine directions and polarization give
$t^2|q_f(q,r)|\le3a_f$. Every allowed derivative of
$F^{\rm cl}$ is $O_j(t)$ on its small-field domain, so
\begin{equation}
 \|R_f\|_m+\sup|\zeta_f|\le C_j a_f .
 \label{ch09:subtraction-bound}
\end{equation}
In the supremum term, $f$ and $Q_f$ occur only under their
own small-field conditions. Enlarging $X$ to $Z$ costs a bounded
number of covering balls.

The conversion of anchored bounds to containing-site bounds
is explicit. If $x\in Z$ and $c(Z)=n$, its anchor lies
within distance $CBn$ of $x$. There are at most
$C(1+B)^4(1+n)^4$ choices. An additional exponent in the support weight
absorbs this polynomial. Consequently the discrepancy
series has containing-site norm at exponent $\sigma+1$
bounded by $C_j(1+B)^{d_j}H_{\rm pre}$.

By translation and hypercubic covariance,
\[
 \sum_{o_i=o}q_{f_i}(q,r)=\Delta\delta_{qr}.
\]
An axis reflection kills off-diagonal orientation pairs;
permutations identify the diagonal pairs.
Absolute convergence follows from the derivative bound,
and
\begin{equation}
 t^2|\Delta|\le C_jH_{\rm pre}.
 \label{ch09:delta-bound}
\end{equation}
Thus, with
$\mathcal C(U)=\frac12\sum_o h_{\{o\}}\sum_q|F_q^{\rm cl}(U,o)|^2$,
the sum of the extracted quadratic functions is
$\Delta\mathcal C(U)$.

\section{Replacing the local quadratic term by the path action}
\label{ch09:path-reset}

The path action is the tested and clipped density of \eqref{ch04:path-action}. Its path and cutoff data remain fixed during this replacement, as do the frame choices defining the actual probability kernel \eqref{ch04:block-operator}.

Expand the quadratic part into plaquette channels and ordered
pairs of entries in each row square. For a channel $d$ let
$a_dP_d$ be its quadratic expression without its event indicator, $a_dQ_d$
its constant-curvature clover substitute, and $L_d$ its
actual clipped or tested expression. Its support $Y_d$
contains both paths and every nonautomatic relaxed small-field condition.
On the small-field set of $Y_d$, $L_d=a_dP_d$.
Define
\[
 F_d=\Delta a_d(Q_d-P_d),\qquad
 J_d=(1-h_{Y_d})\{(\alpha+\Delta)L_d
                              -\Delta a_dh_{\{o_d\}}Q_d\}.
\]
Then
\begin{equation}
 (\alpha+\Delta)L_d+h_{Y_d}F_d
 =\alpha h_{Y_d}a_dP_d+
       \Delta a_dh_{\{o_d\}}Q_d+J_d .
 \label{ch09:channel}
\end{equation}
The equality follows on each of the two health domains.
Every $F_d$ has zero affine Hessian and every $J_d$
vanishes on its small-field support.

The selected free metric satisfies
$c_{\rm dir}I+S(0)^*S(0)=I$, where $S(0)$ is the row
coefficient matrix at zero momentum.
Therefore summing the clover substitutes gives exactly
$\mathcal C$, while the raw channels restricted by their event indicators sum to
$\mathcal Q$. Add the regular terms
$-(\alpha+\Delta)A_d$ for the cubic and quartic atoms.
After the common signed-symmetry average, let $F$ be the
resulting sum of regular terms restricted by their event indicators and $J$ the sum of the $J_d$.
Equation \eqref{ch09:channel} becomes
\begin{equation}
       (\alpha+\Delta)\mathcal K+F
                   =\alpha\mathcal Q+\Delta\mathcal C+J.
 \label{ch09:path-identity}
\end{equation}
Combining \eqref{ch09:early},
\eqref{ch09:subtract-identity} and
\eqref{ch09:path-identity} gives, at every field,
\begin{equation}
 \mathscr Z_{\rm blk}=
 e^{|\Lambda|\tau-\sum_i f_i(1)
       -(\beta'+\alpha+\Delta)\mathcal K-\mathcal L-E}
       \mathcal F^0\,e^{\sum_i\zeta_{f_i}+J},
 \qquad E=\sum_i h_{Z_i}R_{f_i}+F .
 \label{ch09:before-fusion}
\end{equation}
No division by the exact-cover gas was used, so this equality
also holds where that gas vanishes or is signed.

Weighted path sums and additional exponential decay in path length give constants
$C_{\rm reset},L_0,L_1,L_h$ such that
\begin{align}
 \mathcal E_m(E)&\le C_{\rm reset}H_{\rm pre}
     +C_j\{H_{\rm pre}L_0+
                  (|\alpha|t^3+H_{\rm pre})L_1\},\\
 \|\zeta+J\|_{\sigma+1,{\rm site}}
 &\le C_j(1+B)^{d_j}H_{\rm pre}
   +C_j\{t|\alpha|+H_{\rm pre}/t+H_{\rm pre}\}L_h .
 \label{ch09:reset-cost}
\end{align}
Here the site norm sums supremum majorants of logarithmic
terms over supports containing a site.
The low regular channels contain two curvatures, the high
channels at least three; each derivative costs a polynomial
in path length, controlled by the exponential path norm.
The actual clipped channel can have size $O(t)$ on a
large-field configuration, which explains the term
$t|\alpha|$ in the second line.

Choose $c_j^{\rm pre}$ in \eqref{ch09:early-bounds}
small enough to absorb the terms linear in $H_{\rm pre}$
in the first line. At fixed ratio $|\alpha|$ is bounded,
and
\[
 t^3/\epsilon_m^{\rm reg}
       =\widehat\beta_m^{-2/5}(M_m^{\rm cut})^3\longrightarrow0.
\]
Thus the first line is at most $\epsilon_m^{\rm reg}$.
The second line tends to zero even after multiplication
by $(2B+1)^4$: in particular
$H_{\rm pre}/t=O(\widehat\beta_m^{-3/5}/M_m^{\rm cut})$.
This is the smallness needed for the exact regrouping.

\section{Absorbing logarithmic corrections into large-field activities}
\label{ch09:fusion}

Write the logarithmic corrections in
\eqref{ch09:before-fusion} as $\ell_a$ on nonempty exclusion
supports $Z_a$, with $\ell_a=0$ when $Z_a$ is small-field.
Their extra large-field indicator dependence regions are retained.
For each finite connected selection $P$ of preliminary
large-field labels and logarithmic labels, require a
nonempty large-field part and pairwise disjoint preliminary
exclusion supports. Connectivity refers to intersections
of these exclusion supports and the $Z_a$.
Set $Y_P$ equal to their union and $I_P$ equal to the union
of the preliminary prescribed large-field site sets. Define
\[
 K_P(U)=\mathbf1_{\{Y_P\cap\operatorname{Bad}_m(U)=I_P\}}
     \prod_{i\in P_{\rm lf}}k_i(U)
     \prod_{a\in P_{\log}}(e^{\ell_a(U)}-1),
\]
with the conjunction of all inherited additional conditions.

\begin{proposition}[Exact regrouping]
\label{ch09:fusion-prop}
Suppose $\sum_a\sup|\ell_a|<\infty$. The exact-cover sum of
the $K_P$ equals $\mathcal F^0 e^{\sum_a\ell_a}$.
If nonnegative majorants satisfy
\[
 |k_i|\le e^{-\kappa}v_i,\qquad
 |e^{\ell_a}-1|\le v_a,\qquad
 \sup_x\sum_{a:x\in Z_a}e^{(\sigma+1)c(Z_a)}v_a\le\gamma,
 \quad (2B+1)^4\gamma\le1,
\]
where the unified index also includes the preliminary
large-field supports, then
\[
               \mathcal K_m(K)\le e^{-\kappa}(e^\gamma-1).
\]
\end{proposition}

\begin{proof}
Expand $e^{\sum_a\ell_a}=\prod_a[1+(e^{\ell_a}-1)]$.
In a nonzero term each selected logarithmic label contains
a large-field site. Exact coverage assigns that site to a preliminary
prescribed large-field site set, so its connected component contains
a large-field label. Taking components gives disjoint union
supports and the displayed new prescribed-site equality.
Conversely flattening a compatible family of new labels
recovers the old exact cover and the logarithmic selection.
No constituent label is repeated because its support is
nonempty. All additional conditions are retained. This is a bijection
of finite terms and does not require independence of their
large-field indicator reads.

For the bound retain $e^{-\kappa}$ from one large-field
label per component, and bound the others by $v_i$.
Subadditivity of $c$ on connected unions and the rooted
tree bound underlying Theorem~\ref{ch08:kp} sum the
connected selections through a fixed site. One additional unit of exponential support decay bounds the number of attachments, since a support has at most
$(2B+1)^4c(Z)$ sites. Their sum is bounded by
$\gamma$, hence by $e^\gamma-1$.
The resulting absolute envelope justifies the countable
regrouping.
\end{proof}

Take $\kappa$ to be half of the exponent retained in the
second bound of \eqref{ch09:early-bounds}.
The remaining half makes the preliminary large-field
majorants small in the attachment norm.
For the logarithmic terms use
$|e^\ell-1|\le e^{\sup|\ell|}\sup|\ell|$ and
\eqref{ch09:reset-cost}.
Then $\gamma\to0$ and a strict margin
$(2B+1)^4\gamma\le1/2$ is available.
Finally
\[
 e^{-\kappa}(e^\gamma-1)
       \le \exp(-(M_m^{\rm cut})^{1/5})
             =\epsilon_m^{\rm lf}
\]
for a sufficiently large inverse coupling threshold, provided the
neighboring reference scales have the prescribed comparison.
The square $(M_{m+1}^{\rm cut})^2$ in $\kappa$ dominates
the exponent $1/5$ and all polynomial logarithmic costs.

\section{The normalization of the field-independent term and continuity}
\label{ch09:continuity}

The output inverse coupling and scalar are now fixed:
\begin{equation}
 \beta=\beta'+\alpha+\Delta,\qquad
 v_{\rm sc}=\tau-|\Lambda|^{-1}\sum_i f_i(1).
 \label{ch09:reset}
\end{equation}
Equation \eqref{ch09:before-fusion} and
Proposition~\ref{ch09:fusion-prop} prove
$\operatorname{dens}(\rho)=\mathscr Z_{\rm blk}$ pointwise.
Flat normalization also gives
$v_{\rm sc}=|\Lambda|^{-1}\log\mathscr Z_{\rm blk}(1)$.
The exact block integral of a positive input is positive;
signed individual activities cause no contradiction.

Projection already provides the useful estimate
\[
 |\Delta|\le C_jH_{\rm pre}/t_m^2
       \le C_j\theta c_j^{\rm pre}
                   \widehat\beta_m^{-1/10}(M_m^{\rm cut})^{-2}
       \le1 .
\]
If the matched coefficient has
$\alpha=-\kappa_G\log L+O_j(1)$ as in
Proposition~\ref{ch07:slope}, then
\begin{equation}
       \beta=\beta'-\kappa_G\log L+e,\qquad |e|\le d_j+1.
 \label{ch09:drift}
\end{equation}
No estimate of $\Delta$ was used to define its local
projection or prove the exact algebra.

There is a second, intrinsic way to identify the reset.
On an even-sided torus choose direction-one links alternately
equal to $u$ and $1$ on even and odd second-coordinate
slices, with $\chi(u)=rt$ for sufficiently small fixed $r>0$.
Every nontrivial plaquette has this distance from the identity. The direct
positive plaquette term dominates the cubic and quartic
terms when $r$ is small, so this small-field configuration
$U_*$ satisfies $\mathcal K(U_*)\ge c_*|\Lambda|t^2$.
Hold the constructed activities fixed and vary only a trial
inverse coupling correction $d$. The scalar discrepancy is
\[
 f(d)=(\beta'+\alpha+d)\mathcal K(U_*)+\mathcal L(U_*)+
 E(U_*)-\log\frac{\mathscr Z_{\rm blk}(1)}{\mathscr Z_{\rm blk}(U_*)}
       =(d-\Delta)\mathcal K(U_*).
\]
It has the unique zero $d=\Delta$.
Two evaluations identify the two scalars only because the
full-field identity has already been established.
Matching two fields alone would not prove that identity.

For a compact continuous input curve choose fixed local
supports, cutoff conditions, truncations and reflected tori.
Projection, subtraction and channel replacement are bounded
linear operations at these fixed labels. Products are
continuous by telescoping with one difference factor.
For changes of inverse coupling hold physical field coordinates
fixed; for a centered Gaussian of precision $\beta A$,
\[
 \frac{d\nu_{\beta_2}}{d\nu_{\beta_1}}(a)
  =(\beta_2/\beta_1)^{D/2}
             e^{-(\beta_2-\beta_1)\langle a,Aa\rangle/2}.
\]
Its derivative is the same density ratio times
$D/(2\beta_2)-\langle a,Aa\rangle/2$.
On a small relative interval it is dominated by a
polynomial times a Gaussian with a smaller positive
precision. Mean and matrix changes give analogous
Gaussian scores. Apply the common-law estimate to these
scores; do not differentiate a merely measurable incoming
large-field activity.

Infinite regrouping requires more than coefficientwise
continuity. Given a sequence of parameters approaching
$u_0$, pass to a subsequence whose weighted input
discrepancies have sum at most $\varepsilon$.
Their sum added to the majorant at $u_0$ is a common
summable envelope. The strict attachment margin allows
this addition for small $\varepsilon$.
The rooted estimate then gives one summable output envelope;
dominated convergence proves weighted continuity on this
subsequence. Every potentially nonconvergent sequence
would have such a convergent further subsequence, a
contradiction.

At fixed volume, the exact-cover sums obey the elementary
telescoping estimate
\[
 |\mathcal F_u(U)-\mathcal F_{u_0}(U)|
 \le e^{\sum_i M_i}\sum_i\sup|K_{i,u}-K_{i,u_0}|,
 \quad M_i=\max(\sup|K_{i,u}|,\sup|K_{i,u_0}|).
\]
Thus the densities are continuous uniformly in $U$.
The positive block-kernel normalizers are continuous and
bounded away from zero on compact inverse coupling intervals;
dominated convergence in probability Haar measure proves
continuity of $\mathscr Z_{{\rm blk},u}$.
The scalar formula and projection show continuity of
$v_{\rm sc}$ and $\beta$. Restriction to a closed
subinterval retains the construction and all its labels.

\section{Wilson initialization}
\label{ch09:wilson}

The finest-scale representation can be constructed directly
from the Wilson action \cite{Wilson1974}, without a limiting
argument. Write $W(u)=1-\operatorname{Re}\tr(u)/n_G$ for the Wilson
plaquette energy and $F(u)=\operatorname{proj}_{\gLie}(u-u^{-1})/2$.
The map $X\mapsto F(e^X)$ has derivative $I$ at zero and a
local inverse $\psi$. The map and its inverse are odd, and
$W(u^{-1})=W(u)$, so $W(e^{\psi(Y)})$ is even.
Let $T_4$ be the degree-four Taylor
term of the Wilson plaquette energy evaluated at
$e^{\psi(Y)}$. Then
\[
 R(u)=W(u)-\tfrac12|F(u)|^2-T_4(F(u))
\]
vanishes through degree four at the identity.
Choose the finest row
$\sqrt{1-c_{\rm dir}}F$, no cubic or loop term, and quartic
atom $(1-c_{\rm dir})T_4$ with its prescribed event indicator.

For a plaquette $p$ write $h_p$ for the health event indicator on its
vertices, $\chi_p^{\rm rel}$ for its row cutoff condition, and
$k_p=(1-c_{\rm dir})|F(U_p)|^2/2$.
Define the residual energy
\[
 \mathcal W_p=
 \begin{cases}
 W(U_p)-c_{\rm lin}t\chi(U_p),&\chi(U_p)>t,\\
 h_pk_p+(1-h_p)(1-c_{\rm dir})W(U_p)
                   -\chi_p^{\rm rel}k_p,&\chi(U_p)\le t.
 \end{cases}
\]
Since $|F(u)|\le\chi(u)$ and $W(u)=\chi(u)^2/2$,
this is nonnegative; when $\chi(U_p)>t$ it is at least
$t^2/4$. The relaxed small-field condition implies that the plaquette
itself is small, and the condition can also fail because of a nearby bad
plaquette. Direct substitution
gives
\[
 W(U_p)=L_p+(1-c_{\rm dir})h_pT_4(F(U_p))
           +(1-c_{\rm dir})h_pR(U_p)+\mathcal W_p ,
\]
where $L_p$ is the actual clipped quadratic channel.
If the relaxed small-field condition inspects a larger neighborhood, its
failure is included in $\mathcal W_p$ by the same formula;
the identity remains valid because on a small-field plaquette
the residual is $k_p-\chi_p^{\rm rel}k_p$.

Use $\beta(1-c_{\rm dir})R(U_p)$ as a regular activity.
Its normalized derivatives are bounded by $C_j\beta t^5$.
For orders at most five, the missing powers after
differentiation are supplied by Taylor's formula; for
higher fixed orders, the direction factors supply at
least $t^5$. Hence
\[
 \beta t^5/\epsilon_N^{\rm reg}
              =O(\widehat\beta_N^{-2/5}(M_N^{\rm cut})^5)\to0 .
\]
Assign the nonnegative residual energies to connected
components of large-field sites whose radius-$2B$ neighborhoods
overlap. A component site set $I$ has exclusion support
$Y_I=N_{2B}(I)$ and activity
\[
 K_I=\mathbf1_{\{Y_I\cap\operatorname{Bad}_N=I\}}
                         e^{-\beta d_I},
\]
where $d_I$ is the sum of assigned residual energies with
the signed average. Each large-field site meets a large-field plaquette,
and a plaquette has at most four vertices, so
$d_I\ge c\,t^2|I|$ for a fixed $c>0$.
Every residual term is assigned exactly once; its relaxed small-field condition has a witnessing large-field site within distance $B$.
The only nonzero exact-cover selection is the actual set
of these components, so the density identity is exact.

A connected prescribed set of $r$ large-field sites is bounded in number
by a rooted spanning-walk count. The containing-site
weighted sum is at most
\[
 C_j(1+B)^4\sum_{r\ge1}
       \left[C_j(1+B)^8e^{5^4\sigma-c\beta t^2}\right]^r .
\]
Since $\beta t^2\ge(M_N^{\rm cut})^2/2$, this fits
$\epsilon_N^{\rm lf}$ for a sufficiently large threshold.
The activity labels and eligibility criteria are independent of $\beta$.
The resulting full Wilson curve is continuous, and its
scalar normalization is zero.

\section{The one-step conclusion and its input data}
\label{ch09:step}

\begin{theorem}[One renormalization step]
\label{ch09:step-theorem}
Assume the derivative orders and spatial exponents stated at
the start of the chapter. Fix an adjacent quadratic pair with the
positive metrics, exact transport identities and actual gauge-covariant block average
of Chapter~\ref{ch05:gaussian}. Let the incoming path have the
normalized quadratic coefficient and exponential coefficient and
derivative bounds of Chapters~\ref{ch04:blocking} and
\ref{ch07:matching}. The incoming activities obey the small-field
cancellations, large-field eligibility rules and weighted bounds of
Chapter~\ref{ch04:blocking}, with the support reserve used above.
Assume also that the represented input density is strictly positive;
this holds at the Wilson start and is preserved by the exact block
kernel. Smallness of signed activities alone is not a substitute
for this condition.
Use the same forest, path realization, clipped frozen source and
cutoff choices throughout the local integral construction of
Chapter~\ref{ch06:large-fields}. Construct the outgoing path by the
classical and quantum matching in Chapter~\ref{ch07:matching}.
Choose the stage constants, ratio floors, minimum
auxiliary periods and inverse coupling thresholds in the order
specified above.
Require also that the actual coarse period exceeds the largest
local read diameter and twice the largest auxiliary reflection
half-period used in the construction. This includes the minimum
period $100(B_m+1)$ in Chapter~\ref{ch06:large-fields}; enlarging
the stage-dependent period floor ensures that every localization
cube and its dependence neighborhood embeds in the actual torus.

Then every admissible continuous input curve has a
continuous output representation with fixed labels,
satisfying the regular and large-field bounds
\eqref{ch04:activity-bounds} and the pointwise identity
\[
 \operatorname{dens}(\rho_u)
       =\mathcal B_m^{\beta'_u}\operatorname{dens}(\rho'_u).
\]
Its inverse coupling is \eqref{ch09:reset} and obeys \eqref{ch09:drift}, uniformly in
volume, cutoff and curve parameter. The construction
retains restriction compatibility and all exact-cover
cutoff conditions. Output inverse coupling-band membership additionally
requires the numerical band margin in the choice of
scales.
\end{theorem}

\begin{proof}
The proof consists of the convergent local expansion,
the three remainder estimates, the exact identities
\eqref{ch09:subtract-identity} and
\eqref{ch09:path-identity}, and
Proposition~\ref{ch09:fusion-prop}. Positivity and
normalization refer to the actual block integral.
\end{proof}

This is a theorem for one compatible adjacent pair. Its premises
include the actual quadratic coefficient and fixed path coefficient bounds;
they cannot be replaced by unrelated bounds of the same numerical
size. The common Gaussian estimate is supplied by the complete
factor construction in Chapter~\ref{ch06:large-fields}, and
Sections~\ref{ch09:stationary-remainder}--\ref{ch09:centered-log}
give the auxiliary estimate for the same pure integral. The
classical bound contraction and inherited-loop contraction of
Chapter~\ref{ch07:matching} allow the path bounds to be retained
with the stated order of choices.

What remains for a whole trajectory is to make those choices
simultaneously, tune the Wilson inverse coupling to the prescribed coarse
endpoint, and prove stability as the ultraviolet cutoff changes.
These are additional assertions, not consequences of applying an
arbitrary finite step repeatedly. In particular, no higher parameter
analyticity is inferred from the continuity proved here.
}{}
\setcounter{chapter}{9}
\IfFileExists{chapters/10-common-trajectory.tex}{% !TEX root = ../monograph.tex
\chapter[Common trajectories and stability]{A common trajectory and fixed-scale stability}
\label{ch10:trajectory}

The step must be used with a single schedule (Section~\ref{ch10:schedule})
and a consistent choice of activity arrays. There are two
separate problems. First one must tune each finite Wilson theory to the
same coarse inverse coupling. Then one must compare theories with different
ultraviolet cutoffs. Continuity of a finite step solves the first
problem, but does not by itself solve the second.

We first construct the numerical schedule and the tuned finite
histories from Theorem~\ref{ch09:step-theorem}. We then describe the
additional quantitative comparison needed for convergence of these
histories. Its abstract consequences are proved here. The distinction
between that comparison and ordinary continuity remains explicit;
in particular we do not assert analytic dependence on the tuned
endpoint. The elementary contraction arguments are applications of
Banach's principle~\cite{Banach1922}; we give the estimates specific to
the oppositely directed inverse coupling and activity recursions.

\section{What is fixed before choosing the scales}
\label{ch10:choices}

A stage $j\ge1$ specifies finitely many derivative orders, support
exponents, path widths and coefficient bounds. These are numerical
bounds for the representations of Chapter~\ref{ch04:blocking}, not a
new collection of physical theories. Stages are nested: the conditions
of stage $j+1$ include those already required at stage $j$. At a fixed
stage there are only finitely many permitted \emph{stage pairs} (the stages of the two scales of one step). Choose
constants for all these pairs simultaneously.

The order of choices is important. Fix the quadratic metric bounds,
path widths and classical bounds first. Choose the odd block ratio
large enough for the bound
\eqref{ch07:classical-budget}, the bound
\eqref{ch07:inherited}, and the margins in the support exponent for the local integral.
Only afterwards choose the lower bounds on the reference coupling and minimum auxiliary
periods. Enlarging such a lower bound does not change the block
kernel or the coefficient construction. The Wilson initialization can
be made smaller than any prescribed fixed fraction of its activity
bounds by this enlargement, as follows from the estimates in
Section~\ref{ch09:wilson}.

Let $C_j^{\rm drift}$ be a stage bound for the error in the downward inverse coupling step:
\begin{equation}\label{ch10:step-drift}
 \left|\beta_{m+1}-\beta_m-\kappa_G\log L_m\right|
       \le C^{\rm drift}_{J(m+1)}.
\end{equation}
Here $J(m+1)$ is the stage used for this step and
$\kappa_G=11C_{\gLie}/(24\pi^2)>0$. The normalization of the field-independent term is included in $C_j^{\rm drift}$.
Chapter~\ref{ch07:matching} fixes this coefficient with the same
normalization of the Lie algebra metric as the block integral.
Its error constant is fixed before the block ratio. Choose
nondecreasing odd integers $\mathscr L_j\ge3$ above all lower bounds on the block ratio at stage $j$,
and also require
\begin{equation}\label{ch10:drift-floor}
 \frac{C_j^{\rm drift}}{\log\mathscr L_j}
 \le \min\{\kappa_G/4,j^{-2}\}.
\end{equation}
Any further fixed logarithmic size bound $K_j$ can similarly be
required to satisfy $K_j/\log\mathscr L_j\le j^{-2}$.
Factors depending on the chosen $\mathscr L_j$ are covered by later lower bounds on the starting scales, not by redefining the constants used in
\eqref{ch10:drift-floor}.

Future applications may prescribe countably many finite numerical
requirements on the schedule, provided their analytic estimates have already been
established for a fixed stage. At stage $j$ take the maximum of the
first $j$ such requirements. This diagonal choice absorbs finitely many constants;
it does not establish an estimate merely because it has been listed.
In particular the comparison hypotheses below must be proved for the
same coefficient and activity construction before they can be used.


One particular later requirement can already be made explicit.
Fix the equation-adapted sections of
Theorem~\ref{ch02:equation-adapted} once and for all.
For $d\le j$, let $a_{dj}$ be their derivative constants at the
orders and spatial exponents of stage $j$, and let $q_d$ bound their
union stencil with $P_d^{\rm loc}$. These constants depend only on
fixed cover paths, finite polynomials and the stage orders;
Proposition~\ref{ch02:section-bounds} makes them independent of
the block ratio, period and cutoff. Fix the stage column bound
$C_j$, support exponent $\sigma_j$, and finite symmetry norm
$\mathcal O_d$ before choosing the ratio. With
\[
 Q_{dj}=2+e^{81\sigma_j}(a_{dj}+C_j+2q_d),
 \qquad \delta_d=(2\mathcal O_d^2)^{-1},
\]
include the additional lower bound on the block ratio
\begin{equation}\label{ch10:equation-floor}
 \mathscr L_j^{1/8}\ge
       \max_{d\le j}2(1+C_j)Q_{dj}\delta_d^{-1}.
\end{equation}
For the bounded coarse polynomial trials also include
$\mathscr L_j^{1/8}\ge8(1+C_j)Q^{\rm coarse}_{dj}$ for $d\le j$,
where $Q^{\rm coarse}_{dj}$ is their small-field chart bound, fixed before
the ratio. On that chart the indicator of the relaxed regularity condition is one. Its
complementary global value is controlled by the same later
square-root large-field bound condition described below.
Here $C_j$ is the prescribed bound, not a norm recomputed after
choosing $\mathscr L_j$. No radius at the ultraviolet cutoff enters this lower bound:
the raw stencil is fixed; the requirement at the ultraviolet cutoff that its cover
be at most $81$ is imposed later by increasing the lower bound on the
reference coupling and hence the available localization radius.
The remaining value norms at the ultraviolet cutoff contain inverse powers of the
small-field threshold. The lower bounds on the starting scales require their product with the
square root of the large-field bound at the cutoff to be at most one,
as well as the fixed-stencil support conditions and the conditions on the number of insertion points per coarse site.
They are absorbed by choosing sufficiently large lower bounds on the reference coupling,
not included in $Q_{dj}$. Those lower bounds on the starting scales also impose
$\widehat\beta_N\ge a_{d,J(N)}$ and all reserve inequalities at the cutoff
for $d\le J(N)$. Since only finitely many degrees are checked at a
stage, this is a noncircular addition to the schedule. Its use is
the near-identity comparison in
Section~\ref{ch11:equation-calibration}.

\section{A slowly increasing stage schedule}
\label{ch10:schedule}

We construct explicitly the \emph{schedule}: the stage and the block ratio at every scale, including the initial scales.
Let $b_j,c_j$ be nonnegative integers specifying respectively the
minimum starting scale and a requested length for a \emph{stage plateau}, a maximal interval of scales on which the stage is constant.
If $\jmath_j$ is the largest smoothing order used in stage $j$ or a
permitted stage pair, include
$b_j\log3\ge2\log(\jmath_j+7)$ for $j>1$. This ensures that
the full dependence radius of a probability-kernel pullback shrinks in
continuum units, including the variables that choose its local gauge transformations.
Starting with $A_1=0$, define
\[
 k_j=c_j+2A_j+1,\qquad
 A_{j+1}=A_j+b_{j+1}+2k_j+1,
 \qquad J(m)=j\quad(A_j\le m<A_{j+1}).
\]
Then $J$ is nondecreasing, tends to infinity, and increases by at
most one in a single step. If $J(m+1)=j>1$, then $m\ge b_j$.
Moreover the interval
$[k_j-\lfloor k_j/2\rfloor,k_j+\lfloor k_j/2\rfloor]$
lies in the $j$th plateau: its lower endpoint is at least $A_j$,
and its upper endpoint is less than $A_{j+1}$.

Set
\begin{equation}\label{ch10:scales}
 L_m=\mathscr L_{J(m+1)},\qquad
 \tau_0=0,\quad \tau_{m+1}=\tau_m+\log L_m,
 \qquad a_m=e^{-\tau_m}.
\end{equation}
Thus $a_m$ is the lattice spacing, and $a_m^{-1}$ is an integer.
Monotonicity gives
\[
 m\log3\le\tau_m\le m\log\mathscr L_{J(m)}.
\]
Choose a prescribed value $\bar\beta>0$ of the coupling at unit scale, the \emph{coarse endpoint}, and reference couplings
\begin{equation}\label{ch10:reference}
             \widehat\beta_m=\bar\beta+\kappa_G\tau_m.
\end{equation}
The lower bound of the first stage on the reference coupling is satisfied by $\bar\beta$. Every subsequent
such lower bound $d_j$ is satisfied by increasing its starting scale until
$\widehat\beta_m\ge d_j$. We may simultaneously require
\begin{equation}\label{ch10:late-entry}
 \log\mathscr L_j\le\tau_m/j,\qquad
 V_j(\tau_m)\le e^{\tau_m},
 \quad J(m+1)=j>1,
\end{equation}
Also require, for the original support radius and its per-site
normalization exponent, $B_{m+1}^2\le2B_m^2$ and
$(L_m B_m^2)^4q_{m+1}^{\rm site}/\tau_{m+1}\le1/j$ at stage $j$.
Here $q_{m+1}^{\rm site}$ is bounded by
$C_j(1+\log\widehat\beta_{m+1})$ using the Haar-ball normalization
argument. At a fixed stage the numerator is a fixed polynomial in
$\log(2+\tau_m)$, so this is a finite lower bound on the starting scale. It will keep
replication costs subpower at a nearby sparse stopping scale. For the
later grouping of finitely many local insertions, also require
$(\sigma_{m+1}+1)h_m/\tau_m\le1/j$, where
$h_m=16[1+L_m(B_m+B_{m+1}+1)^2]^4$ and $\sigma_{m+1}$ is the
prescribed support exponent. This is again a fixed polynomial in
$\log(2+\tau_m)$ at each stage. It ensures that moving a finite
number of insertion points contributes only a subpower factor.
In \eqref{ch10:late-entry}, $V_j$ is any prescribed fixed polynomial bound for the required
lattice side. The same argument works for a subexponential $V_j$.
Choose $b_j$ so these inequalities hold whenever
$\tau_m\ge b_j\log3$. They then hold at every step of stage $j$,
not just at its first occurrence.

For comparison estimates requiring an interval of scales $w_j(s)=o(s)$ around
one scale of stage $j$, increase $c_j$ so that
$w_j(s)\le\lfloor k/2\rfloor$ whenever
$k\ge c_j$ and $k\log3\le s\le k\log\mathscr L_j$.
At $k=k_j$ the whole prescribed interval stays on the same plateau.
This is useful when a fixed finite calculation must be isolated from
changes in derivative order or support exponent.

Only finitely many steps belong to the first plateau. Choose an integer
$M$ large enough that the lattice side $R_p/a_m$ meets the required
lower bound at every step of this plateau for all $p\ge0$, where
\begin{equation}\label{ch10:periods}
                 R_p=4(M+p+1),\qquad p\ge0.
\end{equation}
On later plateaux \eqref{ch10:late-entry} ensures the required lattice side. These periods also contain
arbitrarily late multiples of every positive integer: for a given
integer $c$, take $p=c(M+z+1)-(M+1)$ with $z$ arbitrarily large.
Finally
\[
 \frac{\log L_m}{\tau_m}\longrightarrow0.
\]
This follows from \eqref{ch10:late-entry} outside the first plateau.
No limit involving $\tau_0^{-1}$ is intended.

\section{Tuning a finite history to the coarse endpoint}
\label{ch10:shooting}

The coupling band at scale $m$ is
$[\widehat\beta_m/2,2\widehat\beta_m]$. Define inner barriers
\[
 \ell_m=\bar\beta+\tfrac12\kappa_G\tau_m,
 \qquad u_m=\bar\beta+2\kappa_G\tau_m.
\]
Both lie in the band. If a step starts at $\ell_{m+1}$, its output
is at most $\ell_m-\kappa_G\log L_m/4$ by
\eqref{ch10:step-drift}--\eqref{ch10:drift-floor}. Starting at
$u_{m+1}$ gives an output at least
$u_m+3\kappa_G\log L_m/4$.

We use only the intermediate value theorem. If a continuous function
$f$ on $[x,y]$ satisfies $f(x)\le d<u\le f(y)$, take its first
point $w$ with value $u$ and its last point $z\le w$ with value $d$.
Then $f([z,w])\subset[d,u]$. An excursion below $d$ would force a
later crossing of $d$; an excursion above $u$ an earlier crossing of
$u$. When $d=u$, choose a single point of that level set. In
particular no monotonicity of the output inverse coupling is required.

\begin{theorem}[Tuned finite histories]\label{ch10:finite-history}
Assume that every block step used below satisfies Assumption~\ref{ch09:forest-assumption}, which is not proved in this monograph. Fix the common numerical choices above and the compatible finite
quadratic and path constructions of Chapters~\ref{ch05:gaussian}
and \ref{ch07:matching}. For every cutoff $N$ and period $R_p$ there
is a finite history starting at a Wilson density and ending with
$\beta_{N,p,0}=\bar\beta$. It has the admissible activity bounds and
exact block identities at every scale $0\le m<N$. Its inverse couplings
stay between $\ell_m$ and $u_m$.
\end{theorem}

\begin{proof}
At scale $N$ use the Wilson curve with bare inverse coupling in
$[\ell_N,u_N]$. Apply the continuous step of
Theorem~\ref{ch09:step-theorem}. The barrier inequalities and the
interval argument give a compact subinterval on which the output
stays between $\ell_{N-1}$ and $u_{N-1}$ and attains these endpoints.
Restrict the entire previously constructed history to that interval.
Restriction compatibility preserves the indices, activity bounds and
exact identities. Repeat downwards. At scale zero the barriers
coincide at $\bar\beta$. For $N=0$ use the Wilson family directly.
Every input to an applied step lies in its band; no continuation
through a forbidden inverse coupling interval has been used.
\end{proof}

Choose one such history for each $N,p$, retaining its path arrays,
local integral choices (made for every forest and averaged over the forests as in Section~\ref{ch09:forest-average}) and output representations. The
choice need not be unique. The paths fixed before the period are
kept identical when the period changes. The exact identity is
\[
 \operatorname{dens}(\rho_{N,p,m})
 =\mathcal B_m^{\beta_{N,p,m+1}}
                  \operatorname{dens}(\rho_{N,p,m+1}).
\]
It is an identity of the densities, not merely of their representations, including the additive
normalization constant. That constant is not asserted to converge.

The same construction reaches any coarse endpoint
$|b-\bar\beta|\le h$, where
$h=\min\{\bar\beta/2,\kappa_G\log L_0/4\}$.
Descend to scale one and use the last barrier crossing to reach $b$.
This proves existence in a common endpoint interval, but neither
uniqueness of the bare coupling nor continuity of its choice.

\section{Uniform logarithmic running}
\label{ch10:running}

Telescoping from the fixed endpoint gives
\begin{equation}\label{ch10:running-error}
 |\beta_{N,p,m}-\widehat\beta_m|
 \le\sum_{i<m}C^{\rm drift}_{J(i+1)},\qquad N\ge m.
\end{equation}
For each $\epsilon>0$, choose $r$ so that
$J(i+1)^{-2}\le\epsilon/2$ for $i\ge r$. The right side is at most
\[
 \sum_{i<r}C^{\rm drift}_{J(i+1)}+\tfrac\epsilon2(\tau_m-\tau_r).
\]
The finite initial sum is eventually at most $\epsilon\tau_m/2$.
Consequently
\begin{equation}\label{ch10:running-limit}
 \beta_{N,p,m}=\bar\beta+\kappa_G\tau_m+o(\tau_m)
\end{equation}
uniformly in $N\ge m$ and $p$. In particular the bare Wilson
inverse coupling obeys $\beta_{N,p,N}/\log(a_N^{-1})\to\kappa_G$ uniformly
in $p$. This conclusion does not use convergence of histories at a
fixed scale. A bounded stage error is not summable by itself; the
small ratio in \eqref{ch10:drift-floor} is essential.

By using a smaller fixed constant than $\kappa_G/4$ in that lower bound,
the ratios $\beta_{N,p,m}/\widehat\beta_m$ can all be kept in an
arbitrarily prescribed fixed neighborhood of one. We henceforth
choose it small enough that
\[
 17/20\le\sqrt{\beta_{N,p,m}/\widehat\beta_m}\le23/20.
\]
For example an error at most $\kappa_G\log L_m/10$ at each step
implies a relative error at most $1/10$ after summation. This also
covers the initial indices, including $m=0$.

\section{The coefficient comparison before the activity comparison}
\label{ch10:path-stability}

A representation by transported curvature terms is not unique.
For comparisons we fix the linear curl division and the choice of coordinate contours
of Chapter~\ref{ch07:matching} once on the common
index sets of the coefficients. An alternative is the orthogonal
representative modulo the kernel of all periodic Taylor coefficients;
then the orthogonal complement must be taken inside the admissible
subspace, so that its invariance and moment constraints are preserved.
Equality of the action on one torus is insufficient to identify its
coefficient array.

At a fixed stage enlarge each incoming classical or loop coefficient bound
by a factor $64$, one family at a time, and choose the lower bound on the block ratio
for that enlarged bound. The bounds in
\eqref{ch07:classical-budget} and \eqref{ch07:inherited} then retain
unenlarged output bounds. This enlargement does not concern the
positive quadratic metric, which is kept fixed. The fixed-order
calculation at the critical configuration is valid after reducing its local analytic
neighborhood if necessary.

With quadratic data fixed, the outgoing cubic coefficient is affine
in the incoming cubic coefficient. With the cubic coefficient also
fixed, the same holds for the quartic coefficient. Finally the
quadratic loop coefficient is affine in the incoming loop. The
implicit-differentiation formula \eqref{ch07:uniform-stationary-jets} proves the
first two assertions: its exchange term is fixed in the quartic
comparison. The inherited Hessian and determinant formula prove the
third. Fixed linear coefficient functionals preserve these identities.

Normalize each coefficient norm by its stage bound. If an affine
map $T$ sends the admissible ball of radius $64$ into the output ball
of radius one, then for $\|H\|,\|K\|\le1$, writing
$\delta=\|H-K\|>0$, the admissible extrapolation
$K_*=K+32(H-K)/\delta$ has norm at most $33$ and gives
\begin{equation}\label{ch10:affine-contraction}
 \|T(H)-T(K)\|\le\delta/16.
\end{equation}
The extrapolation preserves all linear symmetry and moment conditions.
This is an affine estimate, not an assertion that the full nonlinear
step is affine.

Proposition~\ref{ch05:horizon-limit} gives convergence of the
quadratic kernels and free transports at every fixed scale, with a
strictly stronger exponential width than the required width. Choose
the direct coefficient strictly below its uniform positive metric
gap, as required there.
Their fixed-window inverses converge by the resolvent identity and the
common inverse bound. In each successive coefficient comparison the
change in these preceding data contributes an error
$e_{N,N',i}\to0$ at fixed $i$. The current coefficient discrepancy
therefore satisfies
\[
 \delta_i\le\tfrac14\delta_{i+1}+e_{N,N',i}.
\]
For $k>m$ and $N,N'>k$, iteration gives
\begin{equation}\label{ch10:path-truncation}
 \delta_m\le2(1/4)^{k-m}
             +\sum_{i=m}^{k-1}(1/4)^{i-m}e_{N,N',i}.
\end{equation}
First send $N,N'\to\infty$ at fixed $k$, then $k\to\infty$.
This proves the Cauchy property in succession for cubic, quartic and
loop coefficients, for the quadratic kernels of the nested blocking sequence and the fixed common coefficient functionals.

There is a useful precise version of the strict width margin. If coefficients
$c_N(\gamma)$ on a common index set converge coordinatewise and
$\sum_\gamma |c_N(\gamma)|e^{w'\ell(\gamma)}\le C$ with $w'>w$,
then their tails at width $w$ and length greater than $r$ are bounded
by $Ce^{-(w'-w)r}$. If only finitely many indices have
bounded length, finite-coordinate convergence and this tail bound
prove convergence at width $w$. For a Hilbert encoding instead, one
must verify summability of the encoding weights on that index set;
a bound on individual coefficients alone does not bound the number of
paths of a given length.

\section{The stronger local comparison required for cutoff limits}
\label{ch10:comparison-input}

Fix a period $R_p$ in continuum units. At scale $m$ let $\mathcal X_{m,p}$ be
the completion of the activity arrays in their normalized
weighted derivative and value norm. Zero seminorm arrays are
identified. Equality of their evaluated partition functions is not
this equivalence relation. Index sets, event indicators, supports and prescribed-site
rules are fixed at this scale independently of $N$. Families are extended
by zero to absent indices.

To state the comparison, set
\[
 \chi_m=\widehat\beta_m^{-1/32},\qquad
 \mathfrak b_m=\widehat\beta_m^{31/32},\qquad
 \lambda_m=\chi_m^{-1}\sqrt{\beta_m/\widehat\beta_m},\qquad
 b_m(\lambda)=\widehat\beta_m(\chi_m\lambda)^2.
\]
The symbol $b_m(\lambda)$ in this section is the inverse coupling coordinate
map; it is unrelated to the penalty parameter $b$. On the narrow band
$17/20\le\chi_m\lambda\le23/20$ we have
\begin{equation}\label{ch10:inverse-coordinate}
 \mathfrak b_m|\lambda-\lambda'|\le
               |b_m(\lambda)-b_m(\lambda')|.
\end{equation}
The extra factor in the right side is
$\chi_m(\lambda+\lambda')\ge17/10$.

Write the step maps at cutoff $N$ as
\[
 x_m=F_{m,N}(\lambda_{m+1},x_{m+1}),\qquad
 q_m=Q_{m,N}(\lambda_{m+1},x_{m+1}),
 \quad
 \beta_m-\beta_{m+1}-\alpha_{N,m}=q_m.
\]
Here $x_m$ is the activity array, the family of localized terms and large-field activities, and $q_m$ is the coefficient $\Delta$ of the coupling constant
renormalization of the step. It is not the vacuum energy renormalization. The coefficient
$\alpha_{N,m}$ fixed by the matching is the one already used in the exact step.

The following strengthens the continuous finite-step statement of
Theorem~\ref{ch09:step-theorem}. Its proof uses the additional numerical reserves already allowed
in the scale choices.
\begin{proposition}[Quantitative comparison on the common data]
\label{ch10:stability-assumption}
The step maps are defined for every prescribed narrow-band inverse coupling
sequence and every allowed activity array of norm at most $1/8$ (here an allowed array is in particular equivariant in the sense of Definition~\ref{ch09:equivariant}, a linear condition with fixed bijections of labels; the maps are those of Theorem~\ref{ch09:step-theorem} and are therefore subject to Assumption~\ref{ch09:forest-assumption}),
including zero terminal arrays. They preserve that ball, satisfy
$|Q_{m,N}|\le s_m/2$, and, on the same cutoff, obey
\begin{align}\label{ch10:lipschitz}
 \|F_{m,N}(z)-F_{m,N}(z')\|&\le\tfrac1{16}\|z-z'\|,\\
 |Q_{m,N}(z)-Q_{m,N}(z')|&\le4s_m\|z-z'\|,\notag
\end{align}
where $\|z-z'\|=\max\{|\lambda-\lambda'|,\|x-x'\|\}$.
The strengthened Wilson terminal arrays have norm at most $1/8$.
The numerical choices satisfy $\sum_{i<m}s_i\le \mathfrak b_m/128$.
The paths and $\alpha_{N,m}$ converge at each fixed scale with the
strict width margin above. Finally, for two sequences of admissible
inputs at different cutoffs, common limits of their scalar coordinates
and Cauchy activity arrays with difference tending to zero imply
output differences tending to zero, in both the activity norm and
the coupling constant renormalization. This last assertion applies also to every
zero-terminal-array history with prescribed inverse couplings.
\end{proposition}

The last condition compares changing maps $F_{m,N}$, not just changing
inputs to one fixed map. Uniform contractions by themselves do not
imply it. For example $F_N(x)=(-1)^N/10$ has Lipschitz constant zero
and no cutoff limit. Nor does smallness of the output imply the
small Lipschitz constant in \eqref{ch10:lipschitz}.

\begin{proof}[Proof of Proposition~\ref{ch10:stability-assumption}]
We derive the contraction from finite parameter evaluations (finite linear combinations of values along a real line), bound those
evaluations using the local expansion, sum the errors over scales, and
finally prove continuity as the cutoff is removed.

\begin{enumerate}[label=\emph{Step \arabic*.},ref=\arabic*,leftmargin=*,align=left]
\item \emph{Contraction from finite parameter evaluations.}
One sufficient quantitative estimate uses only real parameter lines.
For a line $z+tv$ of admissible inputs, $|t|\le1$, suppose every finite
choice of weights $a_i$ and points $t_i$ satisfies
\begin{align}\label{ch10:finite-tests}
 \left\|\sum_i a_i F(z+t_i v)\right\|
 &\le d\sup_{n\ge0}\left|\sum_i a_i(t_i/r)^n\right|,\\
 \left|\sum_i a_i Q(z+t_i v)\right|
 &\le f_m\sup_{n\ge0}\left|\sum_i a_i(t_i/(2r))^n\right|,
 \qquad r\ge2.\notag
\end{align}
These are estimates of the arrays in their derivative norm,
not just of the value of a partition function. They imply an
absolutely convergent one-variable expansion. Indeed the map
$((t/r)^n)_{n\ge0}\mapsto F(z+tv)$ extends linearly and boundedly
from its span in $c_0$ to its closure. That closure contains every
coordinate vector: subtract the first $n$ coordinates, divide by
$(t/r)^n$, and let $t\downarrow0$. Applying the bounded map to the
geometric expansion gives coefficients of norm at most $d$.
Consequently
\[
 \|F(z+tv)-F(z)\|\le d|t|/(r-|t|).
\]
The identical argument applies to $Q$. This short bounded-operator
argument records the precise information in the parameter evaluations; it is not an
inference of analyticity from continuity.

For two inputs at distance $0<h\le1/4$, use the line with direction
$(z'-z)/(4h)$ and evaluate at $t=4h$. Activity norm along this line
is at most $1/8+1/4<1$. An inverse coupling reserve from the narrow band to
$33/40\le\chi_m\lambda\le47/40$ makes the whole line admissible
when $\chi_m\le1/128$. The resulting bounds are at most $4dh$ and
$4f_mh$. If $h>1/4$, use the separate output bounds; if the distance
is zero, use the same estimate with a vanishing seminorm direction
and let the parameter tend to zero. Thus $d\le1/256$ and
$s_m=2f_m$ imply \eqref{ch10:lipschitz} with room to spare.
The finite-evaluation inequalities require the analytic estimates on the
local expansion established in the next step.

\item \emph{Parameter estimates from the local expansion.}
\label{ch10:parameter-supplier}

Here is the parameter estimate for the local factors. It uses
the protected change of variables and real Gaussian majorants of
Chapter~\ref{ch06:large-fields}, retaining their fixed derivative
order and margins in the support exponent. The incoming array is affine along the
parameter line. The path, event indicators, held fine reads and coefficient
projection are fixed. In particular the finite low-order collection in
Chapter~\ref{ch09:renormalization} is chosen on the compact inverse coupling
band independently of incoming activity arrays. The estimates below
therefore compare the same output construction on different lines.
For the estimates one uses the explicit local expansion on the larger
affine activity domain. Its polynomial products and convergent logarithmic
series define the array and local scalar outputs there even when an
extrapolated array has not been identified with a positive density.
Positivity is required for the histories of densities, not for this algebraic
extension of their local comparison formulas.

Put $y=\widehat\beta_{m+1}$, $R=y^{1/256}$ and
$\delta=y^{-1/128}$. The inverse coupling is a square of an affine function
of the line parameter. The relative change on $|t|\le6R$ is
$O(Ry^{-1/32})=O(y^{-7/256})=o(\delta)$. Choose a slightly smaller
reference inverse coupling $\beta_0$ so that
\[
 \beta(t)=\beta_0\ell(t)^2,\quad \ell(t)=a+ct,
 \qquad 1\le a-6R|c|\le a+6R|c|\le1+\delta.
\]
All these real inverse couplings remain in the reserved band. A compensating
rescaling keeps each fine field read fixed: the Gaussian coordinate is rescaled
by $\ell(t)$ while its fluctuation coefficient is divided by
$\ell(t)$. The indicators of large fluctuation-field events and measurable incoming values
are consequently not differentiated in this parameter.

For a parameter series use the norm
$|f|_R=\sum_{n\ge0}|f_n|R^n$. Its product estimate is the ordinary
Cauchy-product inequality. A real local energy factor, including its
coordinate prefactor, has the form
\[
 J(t,U)=A(U)\ell(t)^n\exp\{(1-\ell(t)^2)\Gamma(U)\}.
\]
The energy $\Gamma$ can have either sign. It cannot be replaced by
its absolute value before estimating the exponential.
For $q=\Gamma(U)$ the exponential series norm is at most $e^{H(q)}$,
where
\[
 H(q)=(1-a^2)q+2aR|cq|+R^2c^2|q|.
\]
Choose the sign of $\xi$ so that $q c\xi\le0$ and take
$3\le|\xi|\le4$. Since $a\ge1+6R|c|$, direct expansion gives
$H(q)\le(1-\ell(R\xi)^2)q$. After $j$ coarse-field derivatives,
factoring out $\ell(t)^n e^{(1-\ell(t)^2)q}$ leaves a polynomial
of degree at most $2j$. Interpolate that polynomial at $2j+1$
distinct such real nodes. Their scaled Lagrange polynomials have a
fixed coefficient norm. Thus real jet bounds $C d^i$ at those nodes
imply
\begin{equation}\label{ch10:real-node-bound}
 |D^iJ(\cdot,U)|_R\le K_j(1+\delta)^n C d^i,
                   \qquad i\le j.
\end{equation}
The node sign is chosen only for a pointwise bound and is never
differentiated. The same argument applied to the Gaussian density
ratio in $D$ local real coordinates yields
\[
 \left|\ell(t)^D e^{(1-\ell(t)^2)Q(z)}\right|_R
 \le(1+\delta)^D,\qquad Q(z)\ge0.
\]
Only coordinates read by the local term enter $D$.

The large-field corrections, which are small, require a different estimate. Their
parameter dependence is polynomial of degree at most two; real
values and derivatives of their exponential minus one have bound
$v d^i$ at $-1,0,1$. For $v\le1/2$, taking the real logarithm at
these three nodes preserves a bound $C_jv d^i$. Polynomial
interpolation contributes a factor $C_jR^2$, and exponentiating in the coefficient
algebra gives a bound $C_jR^2v d^i$ when $C_jR^2v\le1$.
Affine activities with a prescribed large-field set have the same polynomial interpolation
bound without taking a logarithm. These steps require no
analyticity of the incoming activity as a function of the gauge
field beyond the already imposed finite derivative order.

For the primitive bound $v=y^{-7/16+o(1)}$, this radius factor
changes its exponent to
$-7/16+2/256=-55/128<-2/5$.
The local factor estimate \eqref{ch10:real-node-bound} retains its negative
large-field energy: its additional $(1+\delta)^n$ is absorbed into
the fixed constant per local factor. A local read set in support $S$ has
at most $C_L(1+B)^4 C_B(S)$ coordinates. Since $B$ is polynomial
in $\log y$, its Gaussian prefactor satisfies, for any fixed
$\eta>0$ and sufficiently large $y$,
\[
 (1+\delta)^{C_L(1+B)^4C_B(S)}\le e^{\eta C_B(S)}.
\]
This is absorbed by a margin in the support exponent, independently of the
volume. The joint probability of large fluctuation-field events and Gaussian score moments
are those already used for the common majorant; Cauchy--Schwarz
retains their gain from these events without assuming that they are independent.

We can therefore run the connected expansions of
Chapter~\ref{ch08:chapter} in this coefficient norm. A derivative
of order at most $j$ of a product with $k$ factors contributes a factor at most
$k^j\le2^{jk}$, absorbed by enlarging each vertex bound by $2^j$.
Polynomial support moments are absorbed by another fixed margin in the
support exponent. The first connected expansion has bound $y^{-2/5+o(1)}$;
the logarithmic expansion and spatial derivative weights consequently
have bound $y^{-3/8}$ at radius $R/2$, after increasing the lower bound on the starting scale. The low-order projection uses finitely many auxiliary
volumes of side $O(B)$ and supports coverable by a fixed number of
$B$-balls. Its polynomial factors have exponent zero. Hence the
raw projected regular family has finite-evaluation bound
\begin{equation}\label{ch10:broad-test}
 \mathcal N\left(\sum_i a_i H(t_i)\right)
 \le y^{-1/4}\sup_{n\ge0}\left|\sum_i a_i(2t_i/R)^n\right|.
\end{equation}
Here $\mathcal N$ is the full unnormalized regular derivative bound.
The margin between $-3/8$ and $-1/4$ absorbs the projection factors.
A connected covering number merely growing subpolynomially would not suffice
inside an exponential weight; its uniform bound is used here.

For completeness we explain the sharpening which makes this useful.
At each real point, the stronger one-step remainder estimate gives
\[
 \mathcal N(H(t))\le 3\varepsilon C_{\rm st}x^{-11/10}
                         +O(x^{-23/20+o(1)}),\qquad
 x=\widehat\beta_m,
\]
where the fixed fraction $\varepsilon$ can be made as
small as desired before choosing the lower bound on the reference coupling. This is the
small remainder estimate of Chapter~\ref{ch09:renormalization}
applied to this same collection.
Let a finite parameter evaluation have moment norm $M_q$ at a fixed radius $q$.
Interpolation at $k+1$ fixed real nodes separates its moments through
degree $k$ from the remainder. The interpolating evaluation functional has total
variation at most $U_{k,q}M_q$; the remaining evaluation functional has moment norm
at radius $v\ge q$ at most
$(1+U_{k,q})(q/v)^{k+1}M_q$. Applying the pointwise bound to the
first part and \eqref{ch10:broad-test} to the second yields
\[
 M_q\{U_{k,q}\sup_t\mathcal N(H(t))
             +(1+U_{k,q})(2q/R)^{k+1}y^{-1/4}\}.
\]
Take $k=300$. The last power is $y^{-365/256}$, strictly smaller
than $y^{-23/20}$. First choose $\varepsilon$ to absorb the fixed
interpolation constant, then increase the lower bound on the reference coupling. The
result is an arbitrarily small multiple of $x^{-11/10}M_q$.
The scalar projection contributes the squared inverse field threshold as a factor;
thus its bound has the form
\[
 f_m=C\varepsilon x^{-11/10}/t(x)^2,
 \qquad t(x)=x^{-1/2}e^{(\log\log x)^2}.
\]
The coupling constant renormalization is linear on finite parameter evaluations, and its fixed path contribution
is $O(|\alpha|t(x)^3)=o(x^{-11/10})$.

Arrays with a prescribed large-field set retain one required large-field factor throughout
the exact regrouping. Work with their coefficientwise majorants, not separately
chosen value bounds. A raw hard bound $a_h\le C e^{-cP(y)^2}$,
where $P(y)=e^{(\log\log y)^2}$, is negligible compared with any
fixed fraction $\phi$ of $e^{-P(x)^{1/5}}$. The soft correction has
series bound $a_s\le y^{-1/8}$. Put $p=y^{1/16}$ and
$\zeta=-\log(\phi p)\ge0$. Then
$I=e^\zeta a_h+2a_s\le(8p)^{-1}+2a_s$ eventually, the support
ball count times $I$ tends to zero, and
\[
 e^{-\zeta}(e^I-1)\le\phi(1/4+4pa_s)\le\phi.
\]
This is precisely the estimate of the exact regrouping, with the hard factor
retained once. Radius reduction by a fixed factor converts the
finite-evaluation coefficient bounds into one summable majorant at every
support. Indicators of the eligibility conditions are parameter independent and
preserve it. The result proves both inequalities
\eqref{ch10:finite-tests}, with arbitrarily small fixed $d$, for the
protected local construction. The inverse coupling bound above gives
$s_m=2f_m\le C_jx^{-1/10}$ needed for cumulative control.

This argument uses the real-node bounds for the same protected local
sectors on the reserved inverse coupling band. Those bounds follow by applying
the real estimates of Chapter~\ref{ch06:large-fields} with their
constants fixed on that band; the rescaling keeps all fine reads and
event indicators fixed. It adds a parameter estimate to that construction,
not a differentiation of measurable cutoff indicators.

\item \emph{Summation over scales.}

The cumulative scalar bound is compatible with the schedule.
Suppose at a fixed stage its bound has the proved form
$s(x)\le C_j x^{-1/10}$, where $x$ is the reference coupling.
For the stage increment $a=\kappa_G\log\mathscr L_j>0$ and $x\ge a$,
\[
 (x+a)^{31/32}-x^{31/32}
 \ge\tfrac{31}{32}2^{-1/32}a x^{-1/32}.
\]
Since $x^{-1/10}=o(x^{-1/32})$, a lower bound on the starting scale ensures
$s(x)\le((x+a)^{31/32}-x^{31/32})/128$.
Summation gives the required bound by telescoping, including all
initial scales once $\bar\beta$ exceeds the lower bound of the first stage.
The decay estimate for $s(x)$ is required before this numerical
argument; an arbitrary stage constant would not suffice.

\item \emph{Continuity for a common input.}
\label{ch10:cutoff-supplier}

Use now the quadratic and path coefficient limits with the strict
width margin established in
Section~\ref{ch10:path-stability}. We verify the cutoff-comparison
part of Proposition~\ref{ch10:stability-assumption} on the
allowed inputs. All assertions in this paragraph are at one fixed
scale and period; constants need not be uniform in the scale.

For a finite collection of elementary indices and bounded path lengths,
the fine-coordinate dimension is finite. Write each integral in the
same physical coordinates at the two cutoffs. Convergence of the path
coefficients with a width margin gives convergence of its finitely
many required derivatives, uniformly on the smaller local chart.
The inverse identity
$H_N^{-1}-H^{-1}=H_N^{-1}(H-H_N)H^{-1}$, the common positive Hessian
bound and uniqueness of the critical configuration give convergence of means, precisions
and jets of the critical configuration. Differentiated Gaussian densities are their
ordinary polynomial scores. On a sufficiently late tail their
precisions have a common smaller positive lower bound, so a fixed
Gaussian exponential times a polynomial dominates them. The protected
change of variables leaves measurable fine reads fixed. Thus the
same input activity discrepancy is bounded in its prescribed norm;
no derivative of a measurable indicator is introduced.

Use the physical fine coordinates also for the event indicators. At a fixed
scale their threshold is determined by the fixed reference coupling,
and the list of large fluctuation-field events, the index set of supports and the clipping rule
are independent of the ultraviolet cutoff. Thus the event indicator is the
same measurable function in both integrals. Although its expression
in a centered Gaussian variable would move with the mean or precision,
no such parameter-dependent expression is compared here. The local change of
density includes its Jacobian and leaves that physical event indicator fixed.
The coarse small-field event indicators are likewise identical. This avoids any
uniform assertion about the measure of a parameter-dependent threshold surface.

To obtain the suprema in the output derivative bound, choose the
protected Gaussian and chart estimates through one order beyond the
outgoing order. This is one additional derivative order for the
integration kernels, and does not require one additional derivative
of an incoming activity. Indeed in each differentiated integral the
protected coordinate change holds its entire fine activity read
fixed. An additional coarse increment differentiates the smooth
density, chart and Jacobian scores, while the old measurable factors
remain unchanged. The bounds of Chapter~\ref{ch06:large-fields} apply
at every fixed prescribed score order; increase that order before
selecting the stage constants. Direct path terms and the finite
local projection are smooth and satisfy the same extra order by
their exponential coefficient bounds. Thus this additional order is available
also on a plateau with equal incoming and outgoing activity derivative
orders. On each common small-field chart it gives equicontinuity of the
finite collection of output derivatives. A finite net on the compact
coarse configuration domain and dominated convergence then give
uniform convergence there. The estimates extend to its closure from
the protected neighborhood; multiplication by the same small-field event indicator
preserves the resulting supremum bound.

Next remove the finite-index restriction. Long paths have uniformly
small tails by the strict exponential width margin. Large connected
supports and high cluster multiplicities have uniformly small tails
by the margins in the support exponent and the attachment reserves of the polymer estimates.
Remaining multiplicities inherited from the input arrays are treated
by the Cauchy tail argument at the end of this chapter: a fixed late
array has a summable value tail, and the activity-norm discrepancy
bounds the tail of every later array. In a protected derivative the
input factors are kept at fixed fine reads, so these are value tails;
the differentiations fall on the controlled Gaussian and chart
factors. If an input derivative is explicitly present, its difference
is already controlled by the corresponding component of the input
norm. First truncate with these uniform errors, then compare the
finite integrals. This proves convergence in the output norm.

The coupling constant renormalization is a finite linear functional of the same
converging local derivatives and therefore also converges. Local
subtraction preserves this conclusion by linearity; the exact regrouping preserves
it by its convergent product expansion with a strict margin in the
site-norm bound. The finite auxiliary projection uses the same fields and
periods at both cutoffs. Thus this comparison concerns the
returned arrays and $Q_{m,N}$, not some other representation of their
densities.

The same proof applies when the terminal activity is zero and the
inverse couplings below it are prescribed in advance. At the initial auxiliary
scale the density is the positive bare path exponential, and all
activity bounds are zero. After each step its activity array has the
small bound proved above. Replacing the returned inverse coupling by the next
prescribed inverse coupling multiplies the represented density by a strictly
positive path exponential; its activity gas is unchanged. Therefore
positivity is retained, and the next step is an allowed input to
Theorem~\ref{ch09:step-theorem}. All conditions on the critical configuration, supports and cutoffs
depend on the common path and reference band, which have
not changed. In particular this induction supplies precisely the
histories with prescribed inverse couplings used in the convergence proof; it does not assume
that every abstract element of a completed norm ball represents a
positive density.

The quadratic limit of Proposition~\ref{ch05:horizon-limit}, the affine
coefficient comparison and the preceding local-integral argument give
the common-input continuity required here, including all zero-array
histories with prescribed inverse couplings. Together with the finite-evaluation
estimates and cumulative bound, this proves
Proposition~\ref{ch10:stability-assumption}. \qedhere
\end{enumerate}
\end{proof}

\section{Convergence of the coupled histories}
\label{ch10:fixed-scale}

Recall the schedule from Section~\ref{ch10:schedule}: $J(m)$ is the stage at scale $m$, $L_m=\mathscr L_{J(m+1)}$ is the block ratio, and $\tau_m=\log a_m^{-1}$. The reference coupling $\widehat\beta_m=\bar\beta+\kappa_G\tau_m$ is distinct from the tuned coupling $\beta_m$.

\begin{theorem}[Fixed-scale convergence]
\label{ch10:convergence}
Use the common trajectory above with its strengthened comparison
reserves and fixed coefficient functionals. If the coarse endpoints
converge to $b$, then for every fixed
$m,p$ their inverse couplings and activity arrays converge along the entire
cutoff sequence. Their coefficients $q_m$ of the coupling constant renormalization converge as well.
The limits depend on $b$ and the fixed common trajectory, and not
on the choices made in tuning the bare coupling. No convergence of additive vacuum
normalizations is asserted.
\end{theorem}

\begin{proof}
First prescribe two inverse coupling sequences having the same limit at each
fixed scale. At a scale $k>m$ replace the activity array by zero and
run the same maps downwards using the prescribed inverse coupling sequence, without
imposing the scalar recursion on this auxiliary history. The
hypothesis on its domain is needed here. Repeated use of
\eqref{ch10:lipschitz} gives
\begin{equation}\label{ch10:activity-truncation}
 \|x_m-x_m^{[k]}\|\le(1/16)^{k-m}.
\end{equation}
Figure~\ref{ch10:fig-cutoff-histories} shows the common interval
of scales on which the two truncated arrays are compared.

\begin{figure}[!hbp]
\centering
\begin{tikzpicture}[font=\small, x=1cm, y=1cm, >=Stealth,
    history/.style={ymInk, line width=.9pt},
    truncated/.style={ymAccent, dashed, line width=.9pt}]
  \node[anchor=east] at (-.3,.66) {level};
  \foreach \x/\level in {0/0,2.8/m,6.2/k,8.8/N,10.4/{N'}}
    \node at (\x,.66) {$\level$};
  \foreach \x in {2.8,6.2}
    \draw[ymGrid] (\x,.42) -- (\x,-1.65);

  \foreach \y/\cutoff/\name in {0/8.8/N,-1.17/10.4/{N'}} {
    \begin{scope}[shift={(0,\y)}]
      \node[anchor=east] at (-.3,0) {cutoff $\name$};
      \draw[history,->] (\cutoff,0) -- (6.29,0);
      \draw[history,->] (6.2,0) -- (2.89,0);
      \draw[history,->] (2.8,0) -- (.09,0);
      \foreach \x in {0,2.8,6.2,\cutoff}
        \fill[ymInk] (\x,0) circle[radius=1.7pt];
      \draw[truncated,->] (6.2,-.35) -- (2.8,-.35);
      \draw[ymAccent,fill=white,line width=.8pt]
        (6.14,-.41) rectangle (6.26,-.29);
      \node[anchor=west,inner xsep=3pt] at (6.27,-.35) {$x_k^{[k]}=0$};
    \end{scope}
  }
\end{tikzpicture}
\caption[Comparing two cutoff histories at a fixed level]{Activity comparison at a fixed period $R_p$. The two solid
rows reach the same tuned inverse coupling $\bar\beta$ at level zero.
In each dashed run the activity is reset to zero at $k$, retaining
that row's inverse couplings and block maps down to $m$.
At $m$ the error from this replacement is at most $16^{-(k-m)}$ in
each row, uniformly in its cutoff.}
\label{ch10:fig-cutoff-histories}
\end{figure}


For fixed $k$, the truncated arrays converge by finite downward
induction and the common-input cutoff comparison. This statement
is Cauchy convergence: apply that comparison to arbitrary pairs of
subsequences. Then choose $k$ so the two errors in
\eqref{ch10:activity-truncation} are small. Completeness gives a
common activity limit for both prescribed inverse coupling sequences, and common limits of the coupling constant
renormalizations follow from the same local comparison.

At each fixed scale the inverse coupling coordinate lies in a compact real
interval. A diagonal subsequence therefore makes all inverse coupling
coordinates converge; the preceding paragraph supplies activity
limits without compactness of their norm balls. Compare two such
limiting histories. The finite scalar recursion, summed from their
common coarse endpoint, gives
\begin{equation}\label{ch10:prefix}
 b_m(\lambda_m)-b_m(\lambda'_m)
                 =-\sum_{i<m}(q_i-q'_i).
\end{equation}
The limits of the common coefficients $\alpha_{N,i}$ fixed by the matching cancel.
There is no infinite sum or ultraviolet boundary term.

The estimates \eqref{ch10:lipschitz} pass to the two limits. To
justify this without extending a map arbitrarily to the completed
space, at each finite cutoff run the second limiting inverse coupling sequence
with a zero terminal array. The first paragraph identifies its limit
with the second history. Compare it with the first history at the
same cutoff and pass to the limit.
Put
$d_m=\max\{|\lambda_m-\lambda'_m|,\|x_m-x'_m\|\}$.
The activity discrepancy is at most $1/4$. By
\eqref{ch10:inverse-coordinate}, \eqref{ch10:prefix}, and the
individual scalar bounds, its inverse coupling discrepancy is at most
$\mathfrak b_m^{-1}\sum_{i<m}s_i\le1/128$. Hence
$D=\sup_m d_m\le1/4$. The sharper scalar comparison now gives
\[
 |\lambda_m-\lambda'_m|
 \le4\mathfrak b_m^{-1}\sum_{i<m}s_iD\le D/32,
 \qquad \|x_m-x'_m\|\le D/16.
\]
Thus $D\le D/16$ and $D=0$. Every subsequence has a further
subsequence with these same limits. A subsequence staying a positive
distance away at one fixed scale would contradict this fact, so the
whole sequence converges. Finally $b_m$ is continuous, giving the
asserted convergence of the physical inverse couplings.
\end{proof}

\section{Identifying the analytic comparison on the selected family}
\label{ch10:identification}

For clarity we record what the additional input requires of the local
integral, so it can be checked independently of the preceding abstract
argument. The finite path comparisons must use the same fixed
coefficient functionals and rooted estimates as the history. Replacing the
Wilson normalization at the ultraviolet cutoff must leave every fixed lower-scale
construction unchanged eventually, or have an explicitly vanishing
comparison error. A new choice of equally bounded coefficients at
each occurrence would not suffice.

For a finite set of activity indices, write both local integrals in
common physical coordinates. Converging means and positive precisions,
with their uniform inverse and derivative bounds, give convergence
against a common Gaussian majorant. The protected change of variables keeps the physical fine event indicators fixed,
including their reference thresholds, and retains its common Jacobian. The common majorant of
Chapter~\ref{ch06:large-fields} concerns the inhomogeneous term and the event indicator;
it cannot be replaced by a bound for an unrelated flat Gaussian.
The coupling constant renormalization also needs convergence in its local
derivative norm. Convergence of the gas alone does not identify it.

The infinite index sets require a separate tail argument. For an
input array $f$, let $b_f(i)$ be the supremum of the absolute value of the
activity, including its event indicator. On a fixed finite torus the activity bounds
give $\sum_i b_f(i)\le C_m\|f\|$. A Cauchy sequence has uniformly
small value tails: compare with one late fixed member, truncate that
member's summable tail, and use
$b_f(i)\le b_{f-g}(i)+b_g(i)$. This does not assert density of finite
arrays in every derivative norm. Output derivative tails require the
larger margin in the support exponent in the connected expansion. The local
projection is finite and the sum of the exact regrouping has a strict convergence
margin, so these estimates, when proved on the common inputs, survive
both operations.

The auxiliary zero-array histories used in the proof require the same
checks. Their incoming activity bounds vanish, but their path,
critical configuration, cutoff geometry and differentiated Gaussian factors
do not vanish. One must use the same linearization of the critical configuration, clipping data and
critical-configuration map. Uniform implicit-function uniqueness identifies an
auxiliary map of this kind built from the same repeated steps with the selected one; merely invoking the
existence of some admissible critical-configuration map is insufficient.

The schedule,
finite-history tuning and logarithmic running follow from the finite-step theorem. The fixed-scale conclusion is
Theorem~\ref{ch10:convergence}. Steps~\ref{ch10:parameter-supplier} and
\ref{ch10:cutoff-supplier} in the proof of Proposition~\ref{ch10:stability-assumption}
establish the stronger local comparisons
from the protected expansion and converging background coefficients.
Proposition~\ref{ch05:horizon-limit} supplies the fixed-scale quadratic
limit for the nested blocking sequence. This completes the fixed-scale
construction, which remains distinct from subsequent convergence of
all continuum correlations.
}{}
\setcounter{chapter}{10}
\IfFileExists{chapters/11-marked-insertions.tex}{% !TEX root = ../monograph.tex
\chapter{Insertions and finite-scale comparison}
\label{ch11:insertions}

A step acts linearly on an inserted density even though
its effective-action coordinates evolve nonlinearly. Keeping the density
itself in view resolves two difficulties: several insertions generate
joint terms, and an auxiliary sequence of block kernels can preserve
all original correlation integrals without preserving their intermediate
densities. This chapter establishes that finite-step calculus and the
local comparison estimates it needs. Statements about a limiting quantum
field require further uniform estimates; they are not consequences of
the source algebra alone.

Throughout, the polynomial label space and its local section are
$V_d$ and $P_d^{\rm loc}$ from Chapter~\ref{ch02:observables}.
All block kernels are the normalized kernels of
Chapter~\ref{ch04:blocking}. A \emph{history} means one fixed choice of these
kernels, inverse couplings, adjacent path pairs, and local representations; the local representations are those of Chapter~\ref{ch09:renormalization}, whose activities are averaged over the forests (Section~\ref{ch09:forest-average}). Throughout this chapter we assume Assumption~\ref{ch09:forest-assumption}, which is not proved in this monograph, together with its analogue for the restarted steps of Section~\ref{ch11:restart-construction}, and, for the steps with insertions, its extension stated as Hypothesis~(H) in Section~\ref{ch11:marked-step}, which is not proved either.
The estimates of Chapters~\ref{ch06:large-fields} and
\ref{ch09:renormalization} are used only on that same history.

The argument has four parts. First the squarefree source algebra
and support estimates retain every joint insertion term. Next the
single-insertion estimates and the joint Gaussian comparison construct
the restarted history and the normalized product-moment bounds.
Backward comparison then relates local trials to the
renormalized fields and microscopic equations. Finally a common
localization map onto the polynomial label space and a larger source weight give the bound from a finite separating family, from which uniform finite-difference bounds follow.

\section{An algebra for finitely many insertions}

Let $\mathcal I$ be a finite set of insertion indices and define
\[
 \mathcal A_{\mathcal I}
 =\CC[z_i:i\in\mathcal I]/(z_i^2:i\in\mathcal I),\qquad
 z_S=\prod_{i\in S}z_i.
\]
The variables commute; a monomial with a repeated index is zero.
Every element has the unique form $\sum_{S\subseteq\mathcal I}a_Sz_S$.
If $u$ has zero constant coefficient, then $u^{|\mathcal I|+1}=0$,
so the series for $\exp u$ and $\log(1+u)$ terminate.
For positive weights $\omega_i$, put
\[
 \|a\|_{\omega}=\sum_{S\subseteq\mathcal I}
                         |a_S|\prod_{i\in S}\omega_i.
\]
This norm is submultiplicative by the disjoint-subset convolution
formula. The weights rescale source variables; they are not probabilities.

Write $[z_S]M=M_S$ for the coefficient of $z_S$ in $M$. If
$M_\varnothing=1$ and $S\ne\varnothing$,
\begin{equation}
 [z_S]\log M
 =\sum_{\pi\in\mathcal P(S)}
    (-1)^{|\pi|-1}(|\pi|-1)!\prod_{I\in\pi}M_I .
 \label{ch11:cumulants}
\end{equation}
Here $\mathcal P(S)$ denotes the partitions of $S$ into nonempty blocks,
and $|\pi|$ is the number of blocks of $\pi$.
This is the standard moment--cumulant identity; we use the classical
cumulant calculus of Leonov and Shiryaev~\cite{LeonovShiryaev1959},
rather than develop a second connected-expansion formalism.
Nilpotence makes the formula an algebraic identity even when there
is no positive probability law behind the coefficients.
On bounded coefficient sets, $\log$ is Lipschitz with a constant
depending only on $|\mathcal I|$ and the coefficient bound:
expand the finite sum and telescope each product.

\subsection{The exact density and its integral}

At cutoff $N$, let $\rho_N\ge0$ be the density without insertions, with
$Z_N=\int\rho_N\,d\lambda_N>0$, where $\lambda_N$ is product Haar
probability measure. For dimensionful insertions $F_i$, assume every
$\rho_N\prod_{i\in S}F_i$ is integrable and set
\begin{equation}
 q_N(z,U)=\rho_N(U)\prod_{i\in\mathcal I}(1+z_iF_i(U)),
 \qquad q_m=\mathcal B^{(L_m)}_\beta q_{m+1}.
 \label{ch11:source-density}
\end{equation}
Here $\beta$ is the inverse coupling used for the step $m+1\to m$ on the chosen history. The block map of \eqref{ch04:block-operator} acts coefficientwise. Consequently
$Z_N^{-1}\int[z_S]q_m\,d\lambda_m$ is the original cutoff moment
of the insertions indexed by $S$.

If one instead uses the dimensionless one-insertion density of degree
$d$, its recursion is $H_m=L_m^d\mathcal B^{(L_m)}_\beta H_{m+1}$.
The scalar rescaling belongs to this convention; it is absent from
\eqref{ch11:source-density}, where the dimensionful normalization is
already in $F_i$. Mixing these two conventions changes the identity-field
insertion. In particular the reference-scale coefficient of the dimensionful
identity field in a degree-$d$ convention is $a_{r(d)}^d e_0$, not the coefficient
obtained from the individual scaling dimension of $e_0$.

\begin{proposition}[Changing the block history]
\label{ch11:mass}
Replace any kernels below a fixed scale index $k$ by other normalized positive
block kernels. Transfer the full polynomial $q_k$ through them.
Every coefficient integral is unchanged. A field-independent polynomial
$J\in\mathcal A_{\mathcal I}$ commutes with every such transfer; if
$J_\varnothing=1$, the density without insertions and its partition integral
are unchanged.
\end{proposition}
\begin{proof}
For any integrable coefficient $h$, normalization and Fubini give
$\int\mathcal B h\,d\lambda=\int h\,d\lambda$ and
$\|\mathcal B h\|_1\le\|h\|_1$.
Iteration proves the first assertion. Multiplication by $J$ is a fixed
linear map of the finite coefficient space, so it commutes with the
integral. Its constant coefficient proves the last assertion.
\end{proof}

There is also a conditional-expectation interpretation. Draw the terminal
field with law $\rho_N\,d\lambda_N/Z_N$ and then follow the selected
block kernels. The conditional expectation of $\prod_{i\in S}F_i$
given the scale-$m$ field is
$[z_S]q_m/[z_\varnothing]q_m$ almost everywhere where the denominator
is positive. This follows by testing against a bounded function of the
scale-$m$ field. Its value on the zero-density set is immaterial.
Two histories generally have different conditional laws; the proposition
asserts equality of their integrated moments.

\section{Joint local terms and large-field components}

At an intermediate scale a local representation has the form
\begin{equation}
 e^{\Phi(U)+\sum_{S\ne\varnothing}z_S h_S(U)}
 \sum_{\mathcal F\ {\rm eligible}}
       \prod_{Y\in\mathcal F}
          \left(K_Y(U)+\sum_{S\ne\varnothing}z_S k_{Y,S}(U)\right).
 \label{ch11:local-source}
\end{equation}
Each regular term carries its required regularity condition;
each large-field term retains its original eligibility and large-field conditions.
An index $S$ may contain several insertion indices. Such terms are
generated by integration even if the regular source at the ultraviolet cutoff has only
terms linear in the source variables, that is, of \emph{source degree} one. The dependence set includes the link supports of the
insertions, event indicators, paths, and Gaussian windows. It need not coincide
with the exclusion support of a large-field activity.

One must never divide a coefficient with insertions by $K_Y$: that activity may
vanish. For example, the large-field parts of two insertion factors on a selected
large-field support contribute
\[
 K_Y(1+z_1b_1)(1+z_2b_2)
 =K_Y+z_1K_Yb_1+z_2K_Yb_2+z_1z_2K_Yb_1b_2.
\]
The joint coefficient contains the original activity once.

More generally expand the terminal product into its small-field
parts and large-field parts. Merge the link support of each large-field part with
every selected large-field support it meets, and merge intersecting
link supports of large-field parts. The exact cover ensures that the link support of a large-field part
containing a failure of the regularity condition meets at least one selected support.
Merge each resulting component, retaining the constituent indices,
original eligibility criteria, and union of their prescribed large-field sets.
Its value is the product of those constituents, with each selected
activity included once. Expanding the constituent indices recovers
the original term uniquely; hence this is an exact coefficient identity.
This algebraic construction gives all joint terminal large-field terms,
including at zeros of the original activities.

For a bound uniform in the number of inserted points, assume at most one
insertion point per site and a fixed finite list of link-support radii and
degrees. Then a site belongs to only a bounded number of link supports of
insertions. To apply the component tree estimate of Chapter~\ref{ch08:chapter}
to the union of the link supports, use a larger support exponent in the
bounds on the individual factors.
At the ultraviolet cutoff, with the Wilson action, every new component contains an original
large-field support, whose decay factor is
$\exp(-c\beta_Nt_N^2 n)$. Keep half this exponent aside before summing
attachments. A polynomial in $B_N,\widehat\beta_N$ times the remaining
decay factor tends to zero faster than every inverse power.
This proves small joint terminal bounds when the weighted large-field
parts of the insertions themselves have a summable site bound. That latter bound
is a quantitative hypothesis on the local section; it is not implied by the
finite polynomial identity.

\section{Support bounds for several components with insertions}

Fix $n$ insertion indices in each reflected replica cell. Suppose every
elementary principal term has at least one insertion point, contains its
insertion points in its dependence set, and that set is within
distance $\rho$ of each such point. Suppose insertion points in different
cells are separated by more than $2\rho$.
A component formed by intersecting these dependence sets then lies in one
cell. Indeed insertion points on adjacent terms have distance at most $2\rho$,
so they belong to the same cell. Connect the distinct insertion indices used by
the component. Erasing cycles gives a path of at most $n-1$ edges.
All its insertion points are therefore within $2n\rho$ of one another,
and every point of its dependence set is within $(2n+1)\rho$ of each
insertion point. This bound is independent of the number of elementary
terms in the component.

A second estimate controls increasing the source weights. Let each
component carry a nonempty index set $S$, dependence set $X$, and let
$t(S)$ be the number of replica cells represented in $S$.
Suppose at most $n$ indices occur per cell. Increase the common source
weight from $\omega_-$ to $\omega=H^U\omega_-$.
For components within one cell the factor is at most $H^{Un}$.
For components meeting several cells, retain an additional margin in the support exponent
$\sigma>0$. If
\begin{equation}
 \sigma c_B(X)\ge Un\,t(S)\log H,
 \label{ch11:spare-support}
\end{equation}
then the increased source weight and decreased support weight give
\[
 H^{U|S|}e^{-\sigma c_B(X)}\le1.
\]
Consequently an incidence bound $p_1$ at insertion points for single-cell terms
at exponent $e$, and $p_2$ for other terms at exponent $e+\sigma$,
give the bound $H^{Un}p_1+p_2$ at exponent $e$.

The geometric lower bound \eqref{ch11:spare-support} and the stronger
incoming norm are both required. An estimate proved only at exponent
$e$ cannot be reused at $e+\sigma$. Large-field components require
their own joint attachment estimate because one such component may
carry insertion points from many cells. The half of the decay exponent kept aside above, together
with the stronger support norm, absorbs this additional weight.
The support exponent is therefore chosen before the expansion with insertions
and before the history thresholds are chosen.

\section{One step with insertions}
\label{ch11:marked-step}

We give quantitative bounds before applying the Gaussian comparison.
This also specifies the margin in the support exponent needed when the number of
inserted points grows with the volume. Fix an output exponent $e\ge0$,
put $v=e+5$, and let $B_f\le2B_c$ be the adjacent support radii.
For the prescribed enlargement $X^\uparrow$ of a fine support, the
covering argument of Section~\ref{ch06:primitive-bounds} gives
\[
 Lc_{B_c}(X^\uparrow)\le15^4\bigl(L+10c_{B_f}(X)\bigr).
\]
Consequently, if $10\cdot15^4(e+7)\le L\sigma_f$,
\begin{equation}
 e^{(v+2)c_{B_c}(X^\uparrow)}
 \le T_e e^{\sigma_f c_{B_f}(X)},\qquad
 T_e=\exp(15^4(e+7)).
 \label{ch11:cover-transfer}
\end{equation}
The numerical constant describes this choice of support enlargement;
the essential feature is its fixed additive cost and $L^{-1}$ fine
cover cost.

In one sector, which includes one forest (one rooted spanning tree in every block), and for one value of the auxiliary index, distinguish three kinds of factors:
local corrections without insertions, incoming large-field factors, and insertion factors
with nonempty index sets. Their dependence sets determine
the intersection graph. The exact-cover condition requires the local corrections without insertions in each
component to cover every site prescribed by its incoming large-field factors. In particular a component
with no insertion-free correction has no incoming large-field factor.
This is a condition on the exact cover decomposition, retained when
factors with insertions are attached, rather than a consequence of spatial
intersection alone.

Let $\delta$ be the containing-site bound for insertion-free correction
amplitudes. Suppose the incoming large-field amplitudes, already
weighted by $e^{\sigma_f c_{B_f}}$, have site bound $\eta$.
Let $a$ bound the insertion incidence sum at a fixed insertion point after
\eqref{ch11:cover-transfer}, at exponent $v+2$.
If at most $g$ insertion points project to a coarse site, set
\[
 {\cal W}(B)=(33B)^4\sum_{j\ge0}(j+1)^4e^{-j},\qquad
 \Sigma=\delta+{\cal W}(B_c)L^4T_e\eta+{\cal W}(B_c)g a.
\]
The insertion and background incidence bound at a specified insertion point is
$T_e\eta+a$; no $L^4$ occurs there. The site bound is $\Sigma$ at
exponent $v+1$. To see the latter conversion, root a support at one of
its fine sites or insertion points. Its covering number is at least
$\lfloor d(o,x)/(8B_c)\rfloor$. The number of possible roots in shell
$j$ is at most $(33B_c)^4(j+1)^4$, while losing one support exponent
provides $e^{-j}$. Background roots have $L^4$ possible positions within the block;
roots at insertion points have at most $g$ choices.

Choose inflation factors $m,l\ge1$. The common local integral bound of
Chapter~\ref{ch06:large-fields} may be applied with the insertion-free
amplitudes inflated by $m$ for factors with a nonempty prescribed
large-field set, and by $l$ for other corrections. Smallness remains available
by increasing the lower bound on the reference coupling at fixed $m,l$.
A nonprincipal insertion factor is either pointwise at most its bounded
amplitude divided by $l$, or supported on an event of conditional
probability at most $l^{-1}$ and bounded by that amplitude.
For a compatible connected collection $C$, therefore,
\begin{align}
 t_C\left|\int\prod_{j\in C}F_j\,d\nu\right|
 &\le\prod_{j\in C}b_j,\notag\\
 t_C&=\begin{cases}
 m,&\text{some factor prescribes large-field sites},\\
 l,&\text{no sites are prescribed and some factor is nonprincipal},\\
 1,&\text{purely principal}.
 \end{cases}
 \label{ch11:integrated-component}
\end{align}
Here $\nu$ is the single Gaussian/Haar law of the sector, formed with its forest; by Hypothesis~(H) below the bounds $b_j$ and the factors $t_C$ do not depend on the forest.
If an insertion-free correction occurs, its inflated common product bound
absorbs $t_C$, while the other factors are bounded pointwise.
If none occurs, the condition on prescribed large-field sets removes incoming large-field
factors; use either the pointwise or event bound for one nonprincipal
insertion. No independence is used in either case.

Assume $(2B_c+1)^4\Sigma\le1$. The rooted-tree bound gives connected
component site sum at most $\Sigma$ at exponent $v$, and, at insertion points,
incidence at most $T_e\eta+a$. Indeed let
$c_i=c_{B_c}(X_i)$, $a_i=b_i e^{vc_i}$ and
$\widehat b_i=b_i e^{(v+1)c_i}$. The sum of $\widehat b_j$ over
regions meeting $X_i$ is at most $|X_i|\Sigma\le c_i$.
Iterating $R_i^{(0)}=a_i$ and
$R_i^{(q+1)}=a_i\exp(\sum_{j:X_j\cap X_i\ne\varnothing}R_j^{(q)})$
gives $R_i^{(q)}\le\widehat b_i$. Every finite connected collection
has a spanning tree and is included in this upper bound.
This is the rooted-tree majorization underlying the polymer criterion
of Chapter~\ref{ch08:chapter}; it is applied here with the full
dependence sets of terms with insertions, rather than proved anew for each subset of insertions.

The split of an incoming principal factor $Wz_S$ by a local cutoff is
exact:
\[
 (1+\gamma Wz_S)(1+(1-\gamma)Wz_S)=1+Wz_S .
\]
The second term is nonprincipal and the joint tail bound for cutoff remainders gives
its supporting event probability at most $x^{-1}$.
Thus $l\le x$ absorbs that term. If principal and arbitrary remainder
incidence bounds before enlargement are $h$ and $z$, respectively,
one may take
\begin{equation}
 a=2T_eh+lT_ez.
 \label{ch11:marked-incidence}
\end{equation}
Both pieces retain the same incoming event indicator and dependence set.
Purely principal connected components have the one-cell support
bound proved earlier. Their integral is the joint moment
in Proposition~\ref{ch11:joint-gaussian}, since it contains neither
an insertion-free correction nor an incoming large-field factor.

Finally normalize the insertion-free local gas by the same exact
connected expansion as in Chapter~\ref{ch09:renormalization}, that is, with its activities averaged over the forests as in Proposition~\ref{ch09:forest-average-prop}; the normalization constant does not depend on the forest (part~(i) of that proposition). Assumption~\ref{ch09:forest-assumption} concerns only the construction without insertions, so the combination of the fixed-forest computation above with this average requires the following hypothesis, which is not proved in this monograph. \emph{Hypothesis~(H) (forests with insertions).} For every forest $F$, the exact identity for the block integral of the incoming representation with insertions \eqref{ch11:local-source} has the form \eqref{ch09:fixed-forest-gas}, with the same constant $\mathscr N_F$ and function $\mathscr S$, with labels, supports $X_M$ (containing the dependence sets of the insertions carried by $M$) and prescribed large-field sets $I_M$ independent of $F$, and with activities $a_M^F(z,U)$ in $\mathcal A_{\mathcal I}$, Borel and gauge-invariant in $U$ and local on $X_M$, satisfying $\sum_M\sup_U\|a_M^F(\cdot,U)\|_\omega<\infty$; the labels comprise those of \eqref{ch09:fixed-forest-gas}, with the same supports and prescribed large-field sets and with coefficient of $z_\varnothing$ equal to the activity there, and further labels, which carry an insertion factor and whose coefficient of $z_\varnothing$ is zero; every coefficient $[z_S]a_M^F$ satisfies part~(b) of Assumption~\ref{ch09:forest-assumption}; the bounds of this section, in particular \eqref{ch11:integrated-component}, \eqref{ch11:marked-incidence} and the bounds $\Sigma$, $T_e\eta+a$, $k$, $\psi$, and every other estimate of this chapter which is made for one forest, in particular those of Section~\ref{ch11:actual-reader}, hold for every $F$ with right sides independent of $F$; and, for every $\varsigma\in\mathfrak G$, there is a bijection $M\mapsto\varsigma\cdot M$ from the labels for the insertions $F_\bullet$ onto those for the insertions $F_\bullet\circ\varsigma^*$, transporting supports and prescribed large-field sets as in part~(d), with $a^{\varsigma F}_{\varsigma\cdot M}[F_\bullet\circ\varsigma^*](z,U)=a^F_M[F_\bullet](z,\varsigma^*U)$, where $[F_\bullet]$ indicates the insertions from which the incoming representation with insertions is built. Under~(H), the proof of Proposition~\ref{ch09:forest-average-prop}(i)--(ii) applies coefficient by coefficient, because $\mathcal A_{\mathcal I}$ is commutative and the coefficients of activities with disjoint supports depend on the trees of disjoint sets of blocks: the identity holds with the averaged activities $\mathbb E_Fa_M^F$, which obey the same bounds, and its coefficient of $z_\varnothing$ is \eqref{ch09:averaged-gas}. Thus the density with insertions and the density without insertions are averaged by the same average over the forests. The estimates of this section are made for one forest and are used, through their uniformity in the forest, for the averaged activities.
If its local normalization corrections have total site bound $k$
at exponent $v$ after the $l$ inflation, choose an output bound
$\psi$ with $e^\Sigma-1+k\le\psi$.
The regular nonprincipal amplitudes are then bounded by $l$, the selected
components by $m$, and the incidence at insertion points by $T_e\eta+a$.
The corrections with empty index set contribute zero to that incidence.
In particular output incidence at most one follows from
$T_e\eta+a\le1$; it is not an extra independence or finite-arity
assumption. The estimates hold for every finite index set with
the stated bound $g$ on the number of insertion points per coarse site, uniformly in the number of replica cells.

\subsection{The one-insertion graph estimates}
\label{ch11:one-graph}

Fix a degree $d$ with coarse derivative order at least
$\max(2,d+1)$, fine spatial exponent at least $d$, fine smoothing order
at least $d+2$, and coarse smoothing order at least $d+3$ for the
simultaneous calculation for normalized local insertions and the regular and large-field remainders.
Assume the period is at least $100(1+B_c)$.
The fixed normalized local insertion is tested with one additional
derivative beyond each requested derivative used in a coefficient evaluation; this requirement
does not impose an additional derivative on an arbitrary incoming
activity.

Write $\kappa_i=(\epsilon_i^{\rm for})^{1/2}$ for the square root
of the large-field bound at scale index $i$.
The remainder norm used here is the sum of the anchored,
support-weighted regular majorants and $\kappa_i^{-1}$ times the
anchored, support-weighted large-field majorants, with their original
indices of support and prescribed large-field set. Thus the fine and coarse remainder norms
use $\kappa_f$ and $\kappa_c$, respectively, not a common value.

For this degree write $w_i$ and $f_i$ for the scaling
dimension and curvature-factor count of a homogeneous basis vector.
The leading rescaling is
$(S_Lx)_i=L^{d-w_i}x_i$, so $\|S_L^{-1}\|\le1$.
Taylor localization, the extraction of the local polynomial part by Taylor expansion of the fine local section on the free interpolated
coarse jet, gives the one-step localization map $T$, a triangular change of basis with diagonal
$(t_c/t_f)^{f_i}$. Its entries and inverse are bounded at fixed
degree and stage: the differentiated interpolation kernel has bound
$C_{d,r}L^{-1-r}e^{-c|y-y'|}$ for $r\le d$, summable against the
polynomial growth of a radial jet. Thus the leading map is $TS_L$.

Let $G$ be an incoming linear graph of regular vanishing-jet and
large-field remainders, with norm at most $s_f$ per coefficient. As in the graph transform for invariant manifolds, the step with insertions maps coefficient graphs at the fine scale to graphs at the coarse scale and halves the distance between two of them.
The step with insertions gives a linear coefficient map $M_G$ and a
raw remainder map $R_G$. For a common error bound $\epsilon$ above the bound $C_d/L+o_x(1)$, which is valid before $L$ and the lower bound on the reference coupling are chosen ($\epsilon$ may be taken to be $L^{-2/5}$ after those choices), it satisfies
\begin{align}
 |M_Gx-TS_Lx|&\le\epsilon(|x^\circ|+\|Gx\|_{\rm rem}),\notag\\
 \|R_Gx\|_{\rm rem}&\le C_d|S_Lx^\circ|
                                  +\epsilon\|Gx\|_{\rm rem}.
 \label{ch11:graph-estimates}
\end{align}
The circle deletes the identity-field coordinate. Finite symmetry averaging
only enlarges $C_d,\epsilon$ by fixed degree-dependent constants.

Here is the smallness argument for the remainder, including large
supports. Root the connected expansion at its unique factor with insertions;
the remaining factors have the common bound already established.
For a regular remainder on support $X$, put $r=c_{B_f}(X)$.
If its coarse cover exceeds $B_c$, one retained unit of exponential support decay gives
$E_L(\widehat\beta_f)e^{-B_c/2}$, where $E_L(x)\le C_Lx^{a_L}$ includes the Gaussian-scale costs and fixed derivative polynomials. This still tends to zero because $B_c\asymp(\log x)^2$.
Otherwise distinguish $r^3<L$ and $r^3\ge L$.
In the first case the support fits in the fixed centered nonwrapping
lift. Vanishing radial holonomy derivatives imply cancellation of screened
polynomial link moments: if direction $v_j$ is polynomial of degree
at most $a_j-1$ on the tested links and $\sum_j a_j\le d$, then
$D^k(R\circ\exp)(0)[v_1,\ldots,v_k]=0$.
To see this, write the local analytic coordinates as
$\log(e^{q(x)}H(J)_e e^{-q(x+e_\alpha)})$, with a site gauge
polynomial $q$ and radial coefficients $J$.
Their differential maps parameters of scaling dimension $a$ onto polynomial link
fields of degree at most $a-1$: decompose a homogeneous polynomial
one-form into its radial part and the gradient of
$(x\cdot A(x))/(a)$, then correct the lower-degree link-average error.
Higher derivatives preserve the summed scaling dimension.
Gauge invariance removes $q$; polarization of the vanishing holonomy derivatives and induction in the chain rule remove all partitions except
the desired derivative. Locality permits screening to the tested
links, always in the same centered lift.
The discrete Taylor decomposition of the actual interpolating field
therefore cancels its polynomial part through degree $d$; its
remainder contributes $L^{-(d+1)}$.
In the second case use
$1\le L^{-(d+1)}(1+r)^{3(d+1)}$.
The additional incoming exponential support decay absorbs this polynomial and
all derivative costs. Multiplication by the insertion rescaling $L^d$
therefore leaves $C_d/L$ in both cases.
For a large-field input of physical support-weighted mass $m$,
the incoming normalized mass is $m/\kappa_f$. The large-field factor carrying this input lies in no purely
principal component. The nonprincipal and selected-component estimates,
followed by local subtraction, give
\[
 C_{\rm orb}|x_{\rm out}|\le\eta m/\kappa_f,\qquad
 \|R_{\rm out}\|_{{\rm rem},\kappa_c}\le\eta m/\kappa_f,
 \qquad \eta=L^{-2/5}.
\]
The underlying real/imaginary input decomposition contributes $2m$, not
$2m/\kappa_f$. The factor $\kappa_f^{-1}$ enters only the input norm;
the outgoing selected terms are weighted by $\kappa_c^{-1}$ in the
local subtraction. The common large-field energy bound absorbs that
outgoing normalization before the displayed estimate is formed.
Both adjacent bounds are retained in this comparison.
A normalized local insertion contributes its extracted polynomial and a bounded
Taylor remainder; its coefficient error has the same small factor.
These are the three terms in \eqref{ch11:graph-estimates}.
A wrapping small-field remainder is never declared large-field.

Choose $L$ and then the lower bound on the reference coupling so the coefficient
perturbation is smaller than half the inverse bound of $T$.
Absorption gives $|x|\le s_f|M_Gx|$ after fixing a suitable stage
column constant $s_f$. Thus $M_G$ is invertible.
The outgoing graph $G'=R_GM_G^{-1}$ lies in the prescribed ball
$\|G'\|\le s_c$ if $s_c$ is chosen above the finite constant from
\eqref{ch11:graph-estimates}. This ball invariance is part of the
choice, not a consequence of invertibility alone.
The identity field has zero remainders and is sent exactly to $L^d$ times itself.

All arrays and maps taking local polynomial coefficients are fixed independently of $G$.
For graphs $G,H$ with $a=\|G-H\|/s_f$, the distinguished factor in
their difference is $(G-H)x$; the normalized local insertion cancels.
The same rooted sums therefore prove
\begin{equation}
 |(M_G-M_H)x|+\|(R_G-R_H)x\|_{\rm rem}
                         \le\delta a|x|,
 \qquad \delta=C_d\epsilon s_f .
 \label{ch11:graph-difference}
\end{equation}
The harmless $C_d$ covers the finite symmetrization.
Increase the block ratio and the lower bound on the reference coupling until $4\delta s_f\le1$,
together with the preceding inverse and invariance inequalities.
These are raw-map difference estimates; continuity alone
would not supply them. They are the input for the contraction estimate
of the graph transform on the coefficient graphs.

For fixed incoming data the maps are continuous in the inverse coupling
and the canonical path coefficients in weighted $\ell^1$ with a strict margin in the decay width, and in the incoming activity arrays in the same strong norm. All local integrands use the same
protected physical-coordinate change and fixed common index sets.
The common integrable bounds in Chapter~\ref{ch06:large-fields}
justify dominated convergence on each fixed finite truncation. The strict margins in the path width and in the support exponent make the omitted countable tails uniformly small. Taking the finite truncation first and then its size to infinity proves continuity in these weighted norms. Finite Taylor localization and symmetry
averaging preserve continuity. Inversion preserves it on the common
inverse-norm ball. No extra derivative of an incoming activity is
taken when comparing equal-regularity adjacent stages.

\subsection{Propagating the already formed terms of source degree at least two}
\label{ch11:higher-modes}

Suppose every pair of distinct ancestors with insertions has periodic
maximum-distance at least $B^2$. If a dependence set contains $r\ge2$
such ancestors, then
\begin{equation}
 Br\le64c_B(X).
 \label{ch11:separated-cover}
\end{equation}
A doubled spanning traversal of a minimal connected cover has length
at most $6B(c_B(X)-1)$. Assign each ancestor to a covering center.
The assigned centers are distinct and separated by at least $B^2-2B$.
Their cyclic traversal costs at least $r(B^2-2B)$. For $B\ge8$
the stated constant follows. This argument uses periodic distances.

Retain two cover units beyond those needed by the connected expansion.
If the source weights grow by $\widehat\beta_f$ per insertion, then eventually
\[
 \widehat\beta_f^{|I|}e^{-c_B(X)}
       \le e^{-c_B(X)/2},\qquad |I|\ge2,
\]
because $64\log\widehat\beta_f/B\le1/2$.
The second unit absorbs site conversion and attachments. Existing terms of source degree at least two thus gain at least $e^{-B/64}$ while retaining their final support
weights. This uses the full incoming majorant for terms of source degree at least two before
changing weights; the single-insertion bound alone would not give it.

Let $H_f$ be that stronger incoming incidence bound and $\zeta$ the
full weighted incidence sum for terms of source degree one. With a strict margin in the
site criterion for insertion-free factors of Section~\ref{ch11:marked-step}, root each
component at an existing factor of source degree at least two or at a factor of source degree one
in a newly merged component. The common integral and rooted-tree estimates give
\begin{equation}
 H_c\le C_e\bigl(e^{-B/64}H_f+\zeta\bigr).
 \label{ch11:higher-induction}
\end{equation}
The constant includes the fixed change of cover radius and the finite list
of degrees, but not the period or number of insertions.
A new component of source degree at least two has at least two individual insertion indices; using the bound for
one gives the deliberately weaker bound by $\zeta$.
Existing large-field coefficients of source degree at least two are treated as whole incoming factors
with their prescribed large-field sets retained. Newly selected components use the
margin in the decay exponent of Chapter~\ref{ch06:large-fields} to absorb the fixed
inflations in \eqref{ch11:integrated-component}.
Choose the stage threshold so $C_e e^{-B/64}\le1/2$ and
$C_e\zeta\le1/2$. Then $H_f\le1$ implies $H_c\le1$.
The joint terminal construction starts the induction.

If the input is specified at exponent $\sigma_f-4$, its change
of cover radius must still leave the two additional coarse units. For example,
\[
 10\cdot15^4(e+9)\le L(\sigma_f-4)
\]
suffices for the enlargement used here.
The four removed fine units and the two retained coarse units are
different costs. A declaration that exponents lie in $[98,100]$
would not prove this inequality.

\section{A joint Gaussian comparison}

The comparison of several local polynomials must preserve their cross
covariances. Separate Gaussian approximations for individual factors
are insufficient. The following finite-degree statement is the useful
form, permitting a different cutoff for every subset of insertions.

\begin{proposition}[Joint polynomial comparison with cutoffs]
\label{ch11:joint-gaussian}
Fix $n,D$ and $A,B\ge1$. For each nonempty $I\subseteq\{1,\ldots,n\}$,
let $p_I$ be a gauge-invariant polynomial of degree at most $D$ and
coefficient $\ell^1$ norm at most $A$.
For every member $t$ of a finite probability mixture, let one probability
space carry a centered joint Gaussian array $\eta_t$, common to all
subsets, with covariance $C_t$. Suppose
\[
 |a_j|\le B,\quad |(C_t)_{ij}|,|(C_*)_{ij}|\le B,\quad
 |(C_t-C_*)_{ij}|\le\varepsilon,\qquad 0\le\varepsilon\le1.
\]
For each $I$ let $0\le\gamma_{t,I}\le1$ satisfy
$\mathbb E_t(1-\gamma_{t,I})\le\varepsilon^2$.
On its positive set suppose the actual normalized local curvature
coordinates, after an allowed gauge rotation, obey
$|\xi_{t,I,j}-a_j-\eta_{t,j}|\le\varepsilon$ on all used coordinates.
The polynomial's invariance removes the rotation.

With the same probability weights $\lambda_t$ for all subsets and the empty product,
define
\[
 M_S=\sum_{\pi\in\mathcal P(S)}\sum_t\lambda_t
       \mathbb E_t\prod_{I\in\pi}\gamma_{t,I}p_I(\xi_{t,I}),\qquad
 M_S^*=\sum_{\pi\in\mathcal P(S)}
       \mathbb E_{C_*}\prod_{I\in\pi}p_I(a+\eta),
\]
and set $M_\varnothing=M_\varnothing^*=1$.
The symbol $\mathbb E_{C_*}$ may denote the Gaussian pairing
functional for a symmetric array, without assuming positivity.
There is a polynomial $P_{n,D}$ such that
\[
 \sum_{S\ne\varnothing}|M_S-M_S^*|
       +\|\log M-\log M^*\|_1
       \le\varepsilon P_{n,D}(A,B).
\]
Its constants are independent of the number of unused coordinates
and of the number of mixture members.
\end{proposition}
\begin{proof}
Fix a partition with $\ell$ blocks and put
$\gamma=\prod_{I\in\pi}\gamma_{t,I}$.
Then $1-\gamma\le\sum_I(1-\gamma_{t,I})$, so
$\mathbb E_t(1-\gamma)\le\ell\varepsilon^2$ without independence.
Telescope each product monomial to replace the physical coordinates
by $a+\eta_t$. Every difference contains one factor of size
$\varepsilon$ and at most $\ell D-1$ Gaussian factors.
Fixed marginal moments and H{\"o}lder's inequality bound the result by
$\varepsilon A^\ell C_{n,D}(1+B)^{nD}$.
Only coordinates in that monomial occur; no ambient-dimension count
is introduced.

Cauchy--Schwarz removes $\gamma$ at cost
$\sqrt\ell\,\varepsilon$ times the second moment of the polynomial.
Finally the Gaussian pairing formula~\cite{Isserlis1918} expresses each
reference moment as a finite sum of products of means and covariances.
Telescope these products to replace $C_t$ by $C_*$.
This contributes another $\varepsilon$ times a fixed polynomial in $A,B$.
Average with the common probability weights and sum over finitely many
partitions. The finite polynomial formula \eqref{ch11:cumulants} gives
the logarithm bound. The empty coefficient is exactly one because
the empty product has no cutoff.
\end{proof}

The proposition applies to the purely principal components of an
expansion with insertions, after the normalization constant has been extracted once.
It does not identify the whole interacting density with a Gaussian.
Deleting an additive remainder for an estimate may leave the precision,
affine mean and prescribed path data unchanged, but it changes the
interacting density. Exact reconstruction always uses the full
density. This distinction is indispensable when matching the Gaussian
comparison to the step.

\subsection{The exact cross-covariance in the block axial gauge}
\label{ch11:tree-covariance}

We identify the flat reference used in the preceding comparison.
Use coordinates with the unnormalized $\ell^2$ scalar products on the two adjacent tori.
Let $D,D'$ be their bond-to-plaquette curls and let
\[
 K=D^*M_fD,\quad K'=(D')^*M_cD',\quad
 A=K+bQ^*Q ,
\]
where $M_f,M_c$ are the modified positive curvature metrics and $Q$
is the rooted linearization of the averaging operation. Let $I_F$ insert the non-tree
coordinates and $A_F=I_F^*AI_F$. The tree gap gives $A_F>0$.
Set $X=I_FA_F^{-1}I_F^*$.
The rooted gauge projection $\mathcal P$ onto the slice of the same block axial gauge
satisfies $A\mathcal P=A$ and $D\mathcal P=D$, since its gauge
potential vanishes at each block root. Thus $XA=\mathcal P$
and $DXA=D$; no inverse of the gauge-degenerate full $K$ is used.

Choose matrices $Z,Z'$ solving $KZ=D^*$ and $K'Z'=(D')^*$.
They exist because $\ker K=\ker D$ and the symmetric matrix $K$
has range $(\ker D)^\perp$; similarly for $K'$.
The matrices $P_f=DZ$ and $P_c=D'Z'$ are independent of the choices,
symmetric and positive semidefinite.

For the linearization $W$ of the critical configuration, taken in its representative $W_F$ in the block axial gauge, the adjacent
Gaussian identities of Chapter~\ref{ch05:gaussian} read
\[
 KW=Q^*K',\quad AW=bQ^*,\quad DW=JD',
\]
where $J=H_{0,f}^{-1}T H_{0,c}$ is the curvature response.
Here $T$ is the induced map on curvatures and $H_{0,f},H_{0,c}$
are the unmodified diagonal operators; the modified metrics
$M_f,M_c$ still define $K,K'$ and agree with them on curls.
Thus $J$ is not the induced map $T$.
Transpose the first identity and multiply by $Z$:
$K'QZ=(D')^*J^*$. Multiplying by $(Z')^*$ gives
$D'QZ=P_cJ^*$. On the other hand $AZ=D^*+bQ^*QZ$.
Multiply by $DX$ and use $DXA=D$ and $DXbQ^*=DW$ to obtain
\begin{equation}
 (DI_F)A_F^{-1}(DI_F)^*
       =DXD^*=P_f-JP_cJ^*.
 \label{ch11:conditional-curvature}
\end{equation}
This proves every cross-entry simultaneously, including different
plaquettes and coordinates in the basis of Section~\ref{ch01:gauge}. It retains both modified metrics,
even when an adjacent pair straddles a change of stage.
The right side is positive because the left is a covariance.
Harmonic and gauge modes were handled through the range equations,
not discarded.

At zero background the Gram matrix of the components is precisely $A_F$:
the absolutely summed path coefficients reproduce their curvature
factorization, the linearization of the averaging operation is $Q$, and each signed
orientation has the same quadratic form. Their probability average
therefore adds no extra factor. The path-length truncation and the non-flat
background are handled by the inverse-window comparison, not by
changing this exact identity.

\subsection{Nonlinear local coordinate and cutoff estimates}
\label{ch11:actual-reader}

Fix a clean neighborhood and one choice of spanning trees of the blocks (a forest; the identity \eqref{ch11:conditional-curvature} holds for every forest and its right side does not depend on the forest, and the estimates below are used, by Hypothesis~(H) of Section~\ref{ch11:marked-step}, with bounds independent of the forest, so that they pass to the average over the forests), one value of the auxiliary index, one coarse gauge and one Gaussian
sample for all coordinates read by the polynomial family.
Section~\ref{ch06:linear-comparison} gives, against the flat model, mean error
$O(B^a t_f^2)$ and covariance error $O(B^a t_f)$, plus exponentially
small window tails. After multiplication by the finitely many curl
rows and summing the finite-range window entries, these estimates
remain $x^{-1/2+o(1)}$.
Combining with \eqref{ch11:conditional-curvature} gives the full
conditional curvature array, rather than only its diagonal.
All comparisons use the same local gauge transformation based at the root; the gauge-invariance of each
polynomial then permits its physical evaluation with a different local
gauge transformation without changing its value.

Here is an explicit scale choice that makes the nonlinear error small.
Put $R=x^{-3/8}$ and $g=\beta^{-1/2}$ on the coupling band.
The coarse logarithm on the required patch has size
$x^{-1/2+o(1)}=o(R)$.
Restrict the finitely many Gaussian read coordinates so the physical
fluctuation is at most a fixed fraction of $R$.
Since $R/g=x^{1/8+o(1)}$, the Gaussian tail bound makes failure of
this restriction smaller than every inverse power of $x$, even after
the polynomial-in-$B$ count of the local reads.
On this event, the product-logarithm and curvature Taylor remainders
are $O(B^a R^2)$. Dividing by $g$ gives
$O(B^a x^{-1/4})$.
The mean error against the flat model, in Gaussian units, is
$O(B^a t_f^2/g)=x^{-1/2+o(1)}$.
The omitted path tails are smaller than every power of $x$.

For a smooth cutoff also requiring the fine small-field cutoff, use the
conclusion of Proposition~\ref{ch06:mean-cancellation} for event indicators far from the large-field set.
The linear fine comparison component is $O(L^{-2}t_f)$; choose the fixed
block ratio before the lower bound on the reference coupling so it lies strictly within
the fine plateau. The comparison for the averaging operation has its independently
chosen strict plateau width. Restricting the local fluctuation variables to the smaller
fixed fraction of $t_f/g$ required by that proposition has failure
smaller than every power of $x$, because $t_f/g=M_f^{\rm cut}$ grows
faster than every power of $\log x$. Intersect these cutoff events
with the radius-$R$ event. Their failure bound is still summable over
the finitely many local reads.

Thus for every fixed family of polynomial reads with the prescribed
clearance and support, the hypotheses of
Proposition~\ref{ch11:joint-gaussian} hold with, for example,
$\varepsilon=x^{-1/32}$ after increasing the lower bound on the reference coupling.
The covariance error is smaller than $\varepsilon$ and cutoff loss
smaller than $\varepsilon^2$. One joint Gaussian sample and one linear prediction of the fine
curvature from the coarse field are used for all subsets. The normalization constant, which does not depend on the forest (Proposition~\ref{ch09:forest-average-prop}(i)), is the one
extracted from the integral without insertions before the insertions are chosen; all
remaining principal products are integrated against its normalized
Gaussian/Haar probability law. Incoming large-field and correction
factors belong to the other components already bounded in
\eqref{ch11:integrated-component}, and are not absorbed into a
subset-dependent probability measure.

\section{Restarting below a comparison scale}

A useful auxiliary history, the \emph{restarted history}, keeps the inverse coupling constant below a
comparison scale while retaining the paths, adjacent quadratic identities,
and block geometry. Proposition~\ref{ch11:mass} preserves all original
normalized coefficient integrals along this history. The construction
below provides its local representations and estimates on that same
restarted history.

Two changes of representation are involved. For a large reference coupling parameter
$s$, define
\[
 H_s=e^{(\log\log s)^2},\qquad h_s=e^{(\log\log s)^{3/2}},\qquad
 t=s^{-1/2}H_s,\quad \vartheta=s^{-1/2}h_s .
\]
The restarted large-field threshold is $\vartheta$, not $t$; its
support radius is the square of the original logarithmic-square radius.
Eventually $2\vartheta\le t$. For the piecewise quadratic--linear function
$d_t$, with linear coefficient $c_{\rm lin}<c_{\rm dir}/16$,
\[
 d_t(r)-d_\vartheta(r)\ge0,\qquad
 d_t(r)-d_\vartheta(r)\ge c_{\rm lin}\vartheta^2
                      \quad(r>\vartheta).
\]
For $\vartheta<r\le t$ this follows from
$c_{\rm dir}r^2/2-c_{\rm lin}\vartheta r$; for $r>t$ it is
$c_{\rm lin}(t-\vartheta)r$.
The positive difference controls newly selected large-field sites.
The next construction assigns their event indicators, eligibility criteria, joint
insertion terms and path tails consistently.

% !TEX root = ../monograph.tex
\subsection{Constructing the restarted representations}
\label{ch11:restart-construction}

We now construct the auxiliary history just described. The physical
insertions retain their original calibration; only the integration
below the comparison level is restarted. All support norms in this
subsection use connected covers and include full dependence regions.
Their exponents are chosen with the reserves required by the marked
expansion before choosing an inverse coupling threshold.

Write $x_i=\widehat\beta_i$ and, for $x$ large, put
\[
 H(x)=e^{(\log\log x)^2},\quad h(x)=e^{(\log\log x)^{3/2}},\quad
 b(x)=\lceil(\log x)^2\rceil,\quad
 \vartheta(x)=x^{-1/2}h(x).
\]
The new representation at level $i$ has radius $B_i=b(x_i)^2$
and threshold $\vartheta(x_i)$. Its clip has the original two
coefficients, evaluated at this new threshold. Its path coefficients
and block-kernel branches are the original ones. In particular,
``retaining the clip'' means retaining its coefficients and formula,
not asserting equality of two different threshold functions.

\paragraph{Choice of reference parameter for the cutoff.}
There is a unique sufficiently large $r(x)$ satisfying
\[
 r(x)^{-1/2}H(r(x))=\vartheta(x).
\]
Indeed the function on the left is eventually strictly decreasing;
its values at $x$ and $x^2$ bracket the right side. Thus
$x\le r(x)\le x^2$, and taking logarithms gives
$r(x)=x^{1+o(1)}$. With this choice the original threshold formula evaluated at $r(x)$
gives the new threshold at $x$. The inverse coupling in the block
integral still lies in the band determined by $x$. In particular
\[
 b(x)\le b(r(x))\le4b(x),\qquad
 25b(r(x))+30\le B_i\le b(r(x))^2
\]
eventually. On each fixed adjacent stage pair,
$x_{i+1}=x_i+\kappa_G\log L_i$ implies
$B_{i+1}\le2B_i$ and
$\vartheta(x_{i+1})/\vartheta(x_i)\to1$.
The local chart, clipped source, Gaussian precision and tail events
are those of Chapter~\ref{ch06:large-fields} with reference parameter $r(x_i)$
and threshold $\vartheta(x_i)$. Its elementary scale inequalities
continue to hold: the radii are logarithmic powers and the thresholds
are $x^{-1/2+o(1)}$. The actual inverse coupling band remains expressed in
$x_i$, not in $r(x_i)$. In particular one must not invoke a lower
bound $\beta\ge r(x_i)/2$. For every fixed $\eta>0$, changing the reference parameter and
the actual band instead give, eventually,
\[
 y^{1-\eta}\le\beta\le2y,\qquad y=r(x_i).
\]
For a prescribed finite local estimate, fix the number of powers
of $\beta^{-1}$ and $g=\beta^{-1/2}$ that occur before choosing
$\eta$. Replacing a comparison power of $y$ by the corresponding
actual inverse coupling power then loses at most $y^{C\eta}$, where $C$
is fixed by that estimate. Choose $\eta$ small enough to retain
its negative exponent. The dimensionless quadratic gap and path
bounds are unchanged; physical covariance factors retain their
actual $\beta^{-1}$ normalization. Likewise, the cutoff-to-noise
ratio is computed exactly as
$\vartheta(x_i)/g\asymp h(x_i)$, and the selected-component energy
is $\beta\vartheta(x_{i+1})^2\ge c h(x_{i+1})^2$.
These comparisons, not a fictitious inverse coupling band in $y$, are
used below.

\begin{proposition}[Re-expression of the endpoint]
\label{ch11:restart-endpoint}
At a sufficiently late comparison level $k$, the original density has an
exact presentation with radius $b(x_k)^2$, threshold
$\vartheta(x_k)$, the original path coefficients, clip coefficients,
inverse coupling and scalar, regular bound no larger than its original
bound, and large-field site bound
\[
 G(x_k)=\exp\bigl(-a_*h(x_k)^2\bigr).
\]
Here $a_*>0$ is chosen sufficiently small at each fixed stage before
its threshold for the starting scale. The construction retains the old large-field
labels and their eligibility criteria inside the new activities.
\end{proposition}
\begin{proof}
Abbreviate $s=x_k$, $t=s^{-1/2}H(s)$, $\vartheta=\vartheta(s)$,
$M=b(s)$ and $B=M^2$. Write $A_t,A_\vartheta$ for the path actions
with their respective event indicators and clips, and $E_t,E_\vartheta$ for
the original regular coefficients with the corresponding event indicators.
The old regular bound is at most $s^{-11/10}$; the old
large-field bound is at most $w_s=e^{-H(s)^{1/5}}$.
These are the selected trajectory's endpoint bounds.

Eventually $2\vartheta\le t$. The clip comparison proved above gives
$d_t(r)-d_\vartheta(r)\ge c_{\rm lin}\vartheta^2$ whenever
$r>\vartheta$, and a nonnegative difference otherwise.
Distribute each plaquette's difference among its incident sites.
Allowing five sites is a harmless upper bound. For the set
$F_\vartheta$ of new large-field sites this gives local energy contributions $d_z$ with
\[
 \sum_zd_z=\sum_p(d_t(r_p)-d_\vartheta(r_p)),\qquad
 d_z\ge c_{\rm lin}\vartheta^2/5\quad(z\in F_\vartheta).
\]
The omitted quadratic path rows give another nonnegative energy contribution.
Truncate each row at true path length $B$ before squaring, and keep
its half-square only where the old retention condition holds and the new one
fails. The relaxed small-field condition reads a radius-$B$ neighborhood: outside it
$2\le\vartheta e^{B/c}$ makes the condition automatic. The truncated
half-square therefore reads radius at most $2B$ and vanishes if its
tested neighborhood is new-small-field. Distribute its value equally
among the new large-field sites in that neighborhood, making no division
when the set is empty. The resulting local energy contributions $b'_z\ge0$ depend only on a neighborhood of radius
at most $4B$, and their sum is exactly the sum of truncated
half-squares. Perform this operation separately for each signed
orientation, retaining its averaging weight.

Choose $a_*\le c_{\rm lin}/20$, put $u=a_*h(s)^2$, and for large-field sites
$z$ put $X_z=\exp[-\beta(d_z+b'_z)]$. The actual band gives
$\beta\vartheta^2\ge h(s)^2/2$, hence $|X_z|\le e^{-2u}$.
The remaining exponent difference is
\[
 D=-(A_t-A_\vartheta)+\beta\sum_z(d_z+b'_z)
                               +(E_\vartheta-E_t).
\]
It consists of the omitted quadratic row pairs, the cubic, quartic
and loop event indicator differences, and the regular event indicator differences.
Each term vanishes on a new-small-field full dependence region.
At a fixed larger support exponent $a+2$ its site bound is
\begin{equation}
 \epsilon_s\le C B^C\bigl(
 \beta C_{\rm path}^2e^{-\sigma B/4}
 +\Gamma_s t^2+\beta C_{\rm path}(t^3+t^4)+s^{-11/10}\bigr)
                  =s^{-1/2+o(1)}.
 \label{ch11:restart-endpoint-error}
\end{equation}
Here $\sigma>0$ is a retained true-length path rate and
$\Gamma_s\le2H(s)$ is the loop coefficient bound.
For the first term expand the difference of squares into ordered
path pairs, one of length exceeding $B$, and retain part of their
exponential length weight. For the atoms restricted by event indicators use their old
bounds $t^j$ and the cover bound $1+2j+\operatorname{length}/B$.
For an old support $Y$, its fixed halo obeys
$c_B(Y^{[8]})\le C+C(M/B)c_M(Y)$. Eventually its old exponent
absorbs this new exponent and the extra unit used to turn an anchored
bound into a site bound. These are the support conversions of
Section~\ref{ch06:primitive-bounds}; their losses are polynomial in
$B$. The same argument bounds enlarged old large-field labels by
$CB^Cw_s$. No path multiplicity or winding path is removed.

Coalesce the absolutely summable terms by full support into $D_Y$.
For large $s$, put $W_Y=e^{D_Y}-1$, so $|W_Y|\le2|D_Y|$ and
$W_Y$ vanishes when its support is new-small-field. The identity
\begin{equation}
 e^{-A_t-E_t}=e^{-A_\vartheta-E_\vartheta}
                \prod_Y(1+W_Y)\prod_{z\in F_\vartheta}X_z
 \label{ch11:restart-endpoint-identity}
\end{equation}
is exact. Multiply by the old large-field gas and expand the optional
$W_Y$ factors. Keep the old compatibility and eligibility criteria.
Group using intersections of full dependence supports, including the fixed
halo for old large-field conditions and the radius-$4B$ supports of $X_z$.
Every nonzero component contains a selected $X_z$: an old activity
contains an old large-field site, hence a new large-field site, and every nonzero $W_Y$
reads a new large-field indicator. The factor at that site is mandatory and intersects
that component. Assign the component its exact new set of large-field sites and also
retain the union of its old sets of large-field sites and all their conditions.
Conversely, those constituents recover the original compatible
collection and optional subset. Thus the regrouping is a bijection
of terms, not a replacement of old eligibility by a new condition.
Disjoint new exclusion supports are the unions of these full read
supports. Their eligibility criteria require exact coverage of the new large-field sites there.

Reserve $e^{-u}$ per selected site factor. The remaining site sum
at exponent $a+1$ is at most
\[
 \rho_s=2\epsilon_s+CB^Cw_s+
       (8B+1)^4e^{(a+1)9^4}e^{-u}=s^{-1/2+o(1)}.
\]
A radius-$4B$ box has a connected $B$-cover of at most $9^4$
boxes; at most $(8B+1)^4$ such boxes contain a given site.
The connected-set bound of Chapter~\ref{ch08:chapter}, losing one
support exponent, applies since $(2B+1)^4\rho_s\le1$.
Every component retains at least one reserved factor, so its total
site bound is at most $e^{-u}(e^{\rho_s}-1)\le G(s)$.
There is no extra sum over prescribed sets of large-field sites: selected site factors already
specify them, and exact-coverage conditions only remove terms.
Absolute convergence justifies countable labels and coalescing.

The regular coefficients themselves are unchanged. Their domains
and allowed directions shrink with the threshold, and increasing
the cover radius decreases their support weights. Their old bounds
therefore persist. The exact identity and the gas bijection prove
the assertion, with unchanged scalar. All operations preserve gauge
invariance, reality and signed symmetry, by reindexing the original
orientation average.
\end{proof}

\paragraph{The genuine density and its neutral reference.}
At an intermediate level write the genuine presentation as
$R=e^{-E}R^0$, where $R^0$ is obtained by deleting only the regular
additive family $E=\sum_\alpha E_\alpha$. Its path, clip, large-field
activities and all their conditions remain unchanged. If the anchored
weighted value bound of that family is $q$, then
$\sum_\alpha|E_\alpha|\le q|\Lambda|$; this proves integrability
at finite volume, but is not used as a volume-uniform estimate.
Instead insert local factors
\[
 h_\alpha=e^{-E_\alpha}-1,\qquad
 \prod_\alpha(1+h_\alpha)=e^{-E}.
\]
Since $|h_\alpha|\le e^{|E_\alpha|}|E_\alpha|$, their local site
bound is the old anchored bound times a fixed polynomial in the
radius. They are inserted into the same expansion as the neutral
Gaussian integral. Thus the object integrated is always $R$, also
for every coefficient of a marked density. Deleting $E$ for the
neutral estimates never deletes it from the density identity.

\begin{proposition}[One restarted history on the comparison window]
\label{ch11:restart-family}
Fix the finite degree and derivative requests and the comparison
window with top level $k$, reference inverse coupling $s=x_k$, and at most
$(\log\log s)^4$ steps. Take its thresholds for the starting scales large enough
that $3s/4\le x_i\le s$ and
$|\beta_{N,p,k}-s|\le s/8$. On the original adjacent path data there
are presentations $P_i$, $i\le k$ in that window, with the bounds
$(B_i,\vartheta(x_i))$, common actual inverse coupling
$\beta=\beta_{N,p,k}$, large-field bound $G(x_i)$ and additive
value cost at most
\[
 q_{k-i}=s^{-1/200}A_s^{4(k-i)},\qquad
                    A_s=e^{(\log\log s)^3}.
\]
Their exact densities satisfy
\begin{equation}
 D_k=\rho_{N,p,k},\qquad
 D_i=\mathcal B_{i,\beta}D_{i+1}=\operatorname{density}(P_i),
 \qquad \int D_i=Z_{N,p}.
 \label{ch11:restart-density}
\end{equation}
Here the block operator has the original geometric branches, evaluated
at the common inverse coupling. The same local factors, normalization and
all deletion-of-support identities supply both weighted marked-step
uses. The constants are uniform in cutoff, period and the compatible
local fields, for the fixed requests.
\end{proposition}
\begin{proof}
Start with Proposition~\ref{ch11:restart-endpoint}.
The reference-inverse-coupling band gives $x_i/2\le\beta\le2x_i$ on every step.
Only the actual inverse coupling is frozen; the reference inverse couplings, original path pairs,
metric floors, gauge-covariant block average derivatives and forest charts keep their
original adjacent values. The recording-scale inequalities above
verify the radii and threshold comparisons in the local integral.
The fine path and clip in that integral are exactly the ones in the
incoming presentation. A change of stage uses its own adjacent pair,
not the previous stage's quadratic form.

Suppose the genuine incoming value cost is $q_n$, $n=k-i-1$.
Its neutral version has zero such cost and the same large-field
bound and loop coefficient bound. Thus it meets the small incoming
regular bound in Chapter~\ref{ch06:large-fields}; we do not claim
that the larger $q_n$ meets that derivative bound. The genuine
additive factors are handled by the local multipliers just defined.

We spell out the quantitative margins on this actual expansion.
The primitive exponent decomposition in
Section~\ref{ch06:primitive-bounds}, with reference parameter
$y=r(x_i)$, gives an unselected local exponent of site size
$t_y$ with $2t_y\le y^{-1/48}$ after enlargement of its threshold.
Indeed the corresponding exterior row, determinant and Haar
bounds are $y^{-1/4+o(1)}$, including radius polynomials and the
actual inverse coupling factors just discussed; the linear residual is
bounded by $y^{-4}y^2$ before these subpower losses. The two
determinant-word series and inverse correction are exponentially
small in $b(y)$, including the $y^4$ factor for two Gaussian
coordinates and any fixed inverse powers of the actual inverse coupling. The inherited path
terms have costs $O(\Gamma\vartheta^2)$ and
$O(\beta(\vartheta^3+\vartheta^4))$, with true-length reserves.
The inherited bounds before changing the reference parameter can be weakened to
$y^{-1/12}$, with fixed finite powers of the actual inverse coupling;
choose the preceding $\eta$ so their loss is less than $1/48$.
Thus all these are negative powers with a strict margin beyond
$1/24$. One may weaken their sum successively to $y^{-1/24}$
and then $y^{-1/48}$
to include exponentiation, support conversion and fixed constants.
The neutral additive part is zero. The local exponents are bounded
extensions on the same chart event indicators; outside them the failed cutoff conditions
are retained as tail-event factors, not evaluated as small exponents.

Selected components have, after the extracted Gaussian and local
normalization factors have been bounded, bound
$\exp[-K h(x_{i+1})^2 r]$ for $r$ selected cutoff conditions, by the actual
component energy in Chapter~\ref{ch06:large-fields}. Choose
$a_*>0$ also so small that $4a_*\le K/2$.
For both marked-step uses at once take
\[
 m=\max\{G(x_i)^{-2},G(x_{i+1})^{-1}\},\qquad
                 1\le l\le\max(1,H(s)^{q-3}),
\]
where $q$ is the fixed requested finite power.
The ratio $h(x_i)^2/h(x_{i+1})^2$ tends to one uniformly in the
window: its logarithm is $O(\sqrt{\log\log s}/\log s)$.
Thus $\log m\le4a_*h(x_{i+1})^2$ eventually, retaining at least
half the component energy. Apply the factor $m$ only to components with an
actual selected prescribed set of large-field sites, and to a selected component within each
such component. The original prescribed set of large-field sites coverage ensures that an
incoming large-field factor cannot occur without such a component.
Factors $l$, finite decorations and local counting costs have
logarithms $O((\log\log s)^2+\log B_i)$ per selected cutoff condition and fit
inside the remaining margin.

The simultaneous tail estimate of Chapter~\ref{ch06:large-fields}
uses these very chart coordinates and event indicators under the same Gaussian
law, including any bounded genuine multipliers. It therefore bounds
all compatible factors from failed cutoff conditions jointly. Its cutoff exponent
and the retained selected-component exponent dominate every multiple
of $\log s$. The unselected components after inflation cost at most
$s^{-1/48+o(1)}$. The projected old large-field site sum is bounded
by ${\cal W}(B_i)L_i^4T_eG(x_{i+1})$, also smaller than every inverse
power of $s$. It need not include $m$ a second time: the selected
component estimate already includes that factor.

The genuine additive multipliers have uninflated site bound
$C(42B_i+1)^4q_n$. Every fixed product of radius factors, support
conversion constants and $l$ is eventually at most $A_s^{1/4}$.
Hence their inflated site bound is at most $q_nA_s^{1/4}$.
For $n\le(\log\log s)^4$ and every fixed $r$,
\[
 q_nA_s^r\le s^{-1/400}.
\]
This follows by taking logarithms: the positive contribution is
$O((\log\log s)^7)$, whereas the negative term is
$-(\log s)/200$. The connected expansion can therefore be rooted
at an additive multiplier when one occurs; its site bound fits
$\omega_n=q_nA_s$. The other components fit
$\delta_n-\omega_n$, where $\delta_n=q_nA_s^3$, since
$s^{-1/48+o(1)}=o(q_n)$. This proves the inflated product bound and
its site summability on the genuine, not a substituted, density.
All finite stage factors in this paragraph are absorbed simultaneously
by its choice of a sufficiently large starting scale; no uniform supremum over infinitely many stages
is used.

It remains to reconstruct the coarse presentation. The regular
connected logarithm has anchored cost at most $16\delta_n$.
Keep the original coarse path action $A_i$. Its difference from
the base action extracted in the Gaussian integral consists of the
cubic, quartic and quadratic-loop terms restricted by their event indicators, with cost at most
\[
 \beta e^{7e}C_i\vartheta(x_i)^3+
 \beta e^{9e}C_i\vartheta(x_i)^4+
 e^{5e}\Gamma_i\vartheta(x_i)^2\le s^{-1/200}.
\]
The cover bound for a degree-$j$ path atom is
$1+2j+\operatorname{length}/B_i$; the path decay absorbs the last term.
Add this correction to the regular logarithm. Its total value cost
is at most $s^{-1/200}+16\delta_n\le q_{n+1}$.
The sign is fixed by
$e^{-\beta B_i^{\rm ref}}e^{-\ell_i}
=e^{-A_i}e^{-(\ell_i+\beta B_i^{\rm ref}-A_i)}$;
$B_i^{\rm ref}$ here is the base action, not the radius.

At the sliced-integrand level the event indicators retain the finite tag-count
factor; it cancels the denominator of the normalized tag average
when the slices are summed. Signed orientations retain their original
probability weights. Neither averaging introduces a second determinant
or a second extracted scalar.
The positive Gaussian prefactor determines a real scalar logarithm.
No branch of a complex logarithm is chosen. The remaining
large-field factors, including the regular logarithms with holes,
are grouped at their full support and exact old prescribed sets of large-field sites as in
\eqref{ch08:reassembly}. The raised hole correction is bounded by
$16l\delta_n e^{16\delta_n}$, using
$|e^{-c}-1|\le e^{|c|}|c|$.
This and the other component costs are $s^{-1/400+o(1)}$, within
the available $H(s)^{-1/8}$ reserve. The inflated connected bound
gives outgoing large-field bound at most $m^{-1}\le G(x_i)^2$,
hence at most $G(x_i)$. No residual path correction is left outside
the regular family. Averaging the regular pieces over the finite
signed motion group retains their anchored bound and invariant sum.
The exact factorization and reassembly thus give
$D_i=\mathcal B_{i,\beta}D_{i+1}$ on every coarse configuration,
including configurations with failed cutoff conditions. There is no division
by a large-field activity, which can vanish.

At each step choose this expansion before specifying the marked
family and retain it for subsequent use. Smaller support exponents
and smaller inflations are then estimates on the same coefficients.
Only the newly constructed lower entry is installed; the endpoint
and all already installed higher entries are unchanged. This proves
finite descending induction, including adjacent stage changes.
Finally positivity of $D_i$ follows from the positive endpoint and
probability kernels, not from positivity of individual activities.
Their mass-one identity and finite-volume integrability give
$\int D_i=Z_{N,p}$ by Fubini. The same equalities hold coefficientwise
for finite marked polynomials. This proves the asserted normalization
and the common-history statement.
\end{proof}


Finally let $C^{\rm ref}$ be a field-independent polynomial of prescribed
connected reference correlations with zero constant coefficient.
The corrected source is
\[
 Jq_N,\qquad J=\exp(-C^{\rm ref}),
\]
and its exact image under the block maps is $Jq_k$. Its coefficient of $z_\varnothing$ remains the
original density without insertions. One must start this comparison from the
full $q_k$, including all regular and large-field coefficients of source degree at least two.
Multiplying a freshly formed product of single insertions would
discard those coefficients and would not compare the original correlations.

\section{The normalized integral with several insertions}
\label{ch11:replicated}
Suppose the representation at the stopping scale, on $V$ sites, has $n$ copies,
weighted single-insertion bounds $x,y$, and higher regular and large-field
incidence at most one each. Assign each term with insertions to one of its
insertion points and each insertion-free term to one site in its support. The
coefficient norm is bounded by
\[
 e^{v_0+V(b_{\rm act}+\kappa)+n(x+y+2)}.
\]
Here $v_0$ is the constant term, $b_{\rm act}$ the absolute action cost per
site and $\kappa$ the insertion-free activity site sum. The estimate uses
$1+t\le e^t$ for the eligible activity product and retains every
coefficient of source degree at least two. The two unit bounds on the terms of source degree at least two come from
\eqref{ch11:higher-induction}.

Choose $0<\rho\le1$ below the small-field threshold. A maximal
packing in the ambient real matrix space of dimension $D_G$, with
Haar invariance, gives $\mu_G\{\chi(U)\le r\}\ge(r/4)^{D_G}$.
Restrict every forward link to normalized Hilbert--Schmidt distance from the identity at most $\rho/4$; the event has
probability at least $(\rho/16)^{4D_GV}$ and all plaquettes have
normalized Hilbert--Schmidt distance from the identity at most $\rho$. If the exact action there is bounded by
$V[b_{\rm act}+\beta b_{\rm loc}(\rho)]$, positivity gives
\[
 Z\ge e^{v_0-V[b_{\rm act}+\beta b_{\rm loc}(\rho)
                         +4D_G\log(16/\rho)]}.
\]
Signed regular costs are included in $b_{\rm act}$. Write $q$ for the full source density at the stopping level, as in
\eqref{ch11:source-density}, and $q^{\rm site}$ for the resulting
per-site cost. Dividing by the lower bound on $Z$ gives
\begin{align}
 \left|Z^{-1}\int[z_{\mathcal I}]q\,d\lambda\right|
 &\le\frac{e^{Vq^{\rm site}+n(x+y+2)}}{\prod_i\omega_i},\notag\\
 q^{\rm site}&=2b_{\rm act}+\kappa+\beta b_{\rm loc}(\rho)
                  +4D_G\log(16/\rho).
 \label{ch11:replicated-bound}
\end{align}
The constant term cancels. Proposition~\ref{ch11:mass} identifies the
coefficient integrals with the original ones. A reflected tiling at
scale $a_m$, stopped at $a_k$, has $(a_m/a_k)^4$ sites per cell.
Taking the number-of-cells root, with common weight $\omega_k$, gives
\begin{equation}
 {\cal R}_m(F)\le\omega_k^{-1}e^{(a_m/a_k)^4q_k^{\rm site}+x+y+2}.
 \label{ch11:root-bound}
\end{equation}
Here ${\cal R}_m$ denotes the chessboard seminorm of $F$ and
$q_k^{\rm site}$ the per-site cost at level $k$.
Multiplying an admissible residual weight by $c$ divides the bound
by $c$, uniformly in the period when the local bounds are uniform.


\subsection{An independent modulus of continuity in the insertion points}
\label{ch11:positional}

We use the usual reflection-positive transfer representation and
chessboard inequality for the Wilson measure
\cite{OsterwalderSeiler1978,FrohlichIsraelLiebSimon1978}.
These apply to the original positive Wilson law; no reflection
positivity of an individual effective activity is asserted.
The product inequality bounds insertions in distinct reflection cells
by their chessboard seminorms. Expanding a copied product of
$F+G$ over subsets of cells proves the triangle inequality for these
chessboard seminorms, including when one of them vanishes.

For correlation evaluation positions separated in every coordinate, the required
modulus follows independently of difference insertions. Represent a
cyclic correlation as a cyclic product of transfer operators
$\tr\prod_{i=1}^q B_iT^{H_i}/Z$, where $T$ is positive,
$B_i$ is a local slab kernel, the slab width is $w$, and the total
length is $M=qw+\sum_iH_i=2sw$.
The reflected-copy bound for a slab and its adjoint gives
$\|B_i\|_{2s}\le A Z^{1/(2s)}$.
The standard Schatten--H\"older inequality
\cite[Theorem~6.32]{TroppMatrixAnalysis2022} now bounds the product with
complex separations $h_i$ by
\[
 \left|\tr\prod_iB_iT^{h_i}\right|
 \le\prod_i\|B_i\|_{2s}\,
                 Z^{\sum_i\operatorname{Re}h_i/M}
 \le A^q Z ,
\]
provided their real parts are nonnegative and their sum equals the
original total free length. This uses the spectral identity
$\|T^h\|_{M/\operatorname{Re}h}=Z^{\operatorname{Re}h/M}$.

Moving one insertion by $l$ lattice units changes two separations to
$H_0-z,H_{q-1}+z$. If both original separations exceed $2R$,
the normalized product is analytic and bounded by $A^q$ for $|z|\le2R$.
Cauchy's estimate, integrated from $0$ to $l\le R$, gives a
difference at most $A^q l/R$. With $R=\rho/a_N$, this is
$A^q la_N/\rho$, a Lipschitz bound in continuum units, independent of the mesh.
The chessboard bound is chosen at one fixed interior cell scale below
the coordinate separation $\rho$.

The transfer kernels act on a compact continuous slice space.
The finite-matrix estimate passes to them by common tagged partitions:
replace each continuous kernel $K(x,y)$ by the matrix
$\sqrt{\mu(C_u)\mu(C_v)}K(\xi_u,\xi_v)$.
Positivity of $T$ is preserved. Every fixed cyclic trace converges
to its cyclic kernel integral by uniform continuity, including the
partition trace and the finitely many norms obtained from reflected copies of the slab.
Since $Z>0$, the integrated inequalities for the normalized products pass to the limit.
Thus the argument does not assume a finite-dimensional transfer space.
It gives the modulus at the separated correlation configurations needed
in Lemma~\ref{ch11:absorption}, without using the difference estimate
that lemma is intended to prove.

Recall that $\mathsf Z_{N,p,d}$ transfers a cutoff coefficient to
its fixed reference scale. A comparison coefficient is admissible
when its inverse coordinates under $\mathsf Z_{N,p,d}$, its \emph{cutoff-scale coordinates}, have a
uniform subpower bound. The graph norm still combines the regular
majorant with the large-field majorant multiplied by
$\kappa_m^{-1}=(\epsilon_m^{\rm for})^{-1/2}$.

\section{Backward comparison of a local trial}
\label{ch11:trial-descent}
The inhomogeneous graph estimate controls local trials as well.
Put $s_i=L_i^{1/8}$, $e_i=s_i^{-3}$ and $A_i=32(1+C_i)^2$.
Choose $100000K_i^6(1+C_{i+1})^2\le s_i$, so the raw error is
at most $e_i$. Subtract one normalized local insertion density from the trial at the ultraviolet cutoff,
choosing its reference coefficient so the residual coefficient at the coarsest scale
is zero. Invertibility permits this choice; exact linearity propagates
the same subtraction through every step.
Let $x_i$ be the residual coefficient and $c_i$ its full remainder norm bound, including the subtracted graph of normalized local insertions. With
$y_i=|S_{L_i}x_{i+1}|$, the raw inequalities give
\[
 y_i\le K_i\bigl(|x_i|+e_i(y_i+c_{i+1})\bigr),\qquad
 c_i\le K_iy_i+e_ic_{i+1},\qquad |x_{i+1}|\le y_i.
\]
If $A_i|x_i|\le c_i$, substitution gives
$y_i\le4K_ie_ic_{i+1}$. The ratio choice ensures
$4K_iA_{i+1}e_i\le1$ and $(4K_i^2+1)s_ie_i\le1$.
Hence $A_{i+1}|x_{i+1}|\le c_{i+1}$ and $s_ic_i\le c_{i+1}$.
The induction starts at the coarsest scale with zero coefficient.
If the trial at the ultraviolet cutoff has coefficient plus remainder cost at most $w$,
the bound on the graph of normalized local insertions gives $c_N\le w+C_N(w+|x_N|)$.
Absorb $|x_N|\le c_N/A_N$ to obtain $c_N\le2(1+C_N)w$.
Thus
\begin{equation}
 c_i+|x_i|\le4(1+C_N)w e^{-(\tau_N-\tau_i)/8}.
 \label{ch11:descent-error}
\end{equation}

For a trial $f=a_N^eF$ with fixed link support, extract its degree-$e$
coefficient $z$ and subtract the corresponding dimensionful section.
Split the difference with zero polynomial coefficient into its small-field part on the
union $D$ of the link supports and its large-field part. Attach the latter to
the original selected large-field factors; the terminal density is
unchanged. If $c_{N,p}$ bounds local derivatives and $b_{N,p}$
the global value, sufficient bounds are
\[
 q=K_ec_{N,p},\quad h=c_{N,p}+C_Nq,\quad
 r=b_{N,p}+2A_e\mathcal N_Nq,
\]
\[
 Q_N=1+q+e^{81\sigma_N}(h+2|D|r),\qquad
 E_{N,p}=4(1+C_N)Q_N+
       e^{81\sigma_N+162}r\zeta_{e,N}u_N^{-1}.
\]
Here $A_e$ bounds the global section, $\mathcal N_N$ is its finite
sum of inverse threshold powers, $\sigma_N$ the support exponent,
$u_N>0$ the strict margin for attachments at the ultraviolet cutoff, and $\zeta_{e,N}$ its
source weight. Eventually the fixed link support has connected covering number at most 81;
two extra support units contribute the factor $e^{162}$.
The first term absorbs the coefficient and small-field remainder; the
second absorbs the attachment of the large-field part with its original prescribed large-field set.
If $c_{N,p},b_{N,p},C_N,e^{81\sigma_N},\mathcal N_N,u_N^{-1}$
are subpower in $a_N^{-1}$, then so is $E_{N,p}$.

The backward comparison and the product-moment estimate now give, at every fixed
pair $(m,k)$ of cell scale and stopping scale, a reference label $v_{N,p}$ with cutoff-scale
coordinates bounded by $2E_{N,p}$ and
\begin{equation}
 \mathcal R_m\bigl(F-O_{N,p}(e,v_{N,p})\bigr)
       \le C_{e,m,k}E_{N,p}e^{-\tau_N/8}.
 \label{ch11:trial-error}
\end{equation}
Indeed the residual source weight is inversely proportional to its
bound in \eqref{ch11:descent-error}. All constants are independent
of period when the terminal norm bounds are. Thus the discrepancy for local trials tends to zero in the chessboard
seminorm. The same estimate applies to finite differences formed from
translates, using the derivative bounds on the translated link supports.
It bounds cutoff-scale coordinates, not the raw reference labels.

\section{Finite-dimensional detection before taking limits}
\label{ch11:detectors}

A useful part of the later difference argument is entirely finite
dimensional. We record it now, keeping its analytic hypotheses visible.
Let $V$ be a fixed finite-dimensional label space, $E_i\subset V$
the equation subspace at a finite cutoff, and $\Pi_i$ a bounded
projection with kernel $E_i$. Such projections need not converge.
Choose a fixed invariant Euclidean inner product and average an
orthogonal projection with this kernel over the finite symmetry group,
assuming $E_i$ is invariant. Each conjugate annihilates $E_i$ and differs
from the identity by a map into $E_i$. The average has these two
properties, which imply idempotence: if $v-\Pi_i v\in E_i$, then
$\Pi_i(v-\Pi_i v)=0$. Its norm remains bounded uniformly in $i$.

Let $L_i:V\to\CC^r$ be the finite family of correlation functionals used
to separate the quotient. Suppose
\begin{equation}
 |\Pi_i v|\le Z|L_i\Pi_i v|
 \label{ch11:detector-bound}
\end{equation}
with one constant $Z$, independent of cutoff and of the volume in continuum units.
This follows from a sufficiently accurate finite-cutoff comparison
with an injective free matrix of correlation pairings: if $|c|\le K|Gc|$ and
$\|Q-G\|\le(2K)^{-1}$, then
$|c|\le K|Qc|+\tfrac12|c|$, hence $|c|\le2K|Qc|$.
The separating family of free correlations comes from the finite Maxwell quotient; no
continuum rank assertion is needed for this argument.

\begin{lemma}[Absorption of a possibly large equation component]
\label{ch11:absorption}
Let $x_i\in V$ be encoded coefficients of a normalized lattice
difference and put $y_i=\Pi_i x_i$. Suppose $|L_i x_i|\le C$ and,
for
\[
 z_i=\frac{x_i-y_i}{1+|y_i|},
\]
the microscopic equation comparison gives
$|L_i z_i|\le(2Z)^{-1}$ eventually, uniformly in the volume in continuum units.
Then $|y_i|\le2ZC+1$.
\end{lemma}
\begin{proof}
Since $L_i y_i=L_i x_i-(1+|y_i|)L_i z_i$,
\eqref{ch11:detector-bound} gives
$|y_i|\le ZC+\tfrac12(1+|y_i|)$.
Move the last half-norm to the left.
\end{proof}

The normalized equation coefficient $z_i$ lies in $\ker\Pi_i$.
Applying an equation estimate to it still requires the estimate's
growth hypothesis. For instance, if $x_i$ has a subpower coefficient
bound, then $z_i$ has one: its first term is a bounded scalar multiple
of $x_i$, while $y_i/(1+|y_i|)$ is bounded by one.
Thus the argument does not presuppose a uniform bound for the raw
difference coefficient $x_i$.

There is also an exact way to test finite differences using bounded
separated observations. For a lattice spacing $a$, a coordinate unit vector $e$, and an integer
$h\ge1$,
\[
 \frac1h\sum_{j=0}^{h-1}
 \frac{F(x+(j+1)ae)-F(x+jae)}a
 =\frac{F(x+hae)-F(x)}{ha}.
\]
For a fixed physical distance $b>0$, taking $h=\lceil b/a\rceil$ bounds the right side by
$2\sup|F|/b$. To replace averaged correlation evaluation positions by the desired
positions, however, one needs a modulus of continuity in the insertion points established
independently of the normalized difference estimate being proved.
That modulus is an essential analytic input, not another consequence
of this telescoping identity. Likewise the microscopic equation
insertion must vanish against the distant spectator product
at finite cutoff, uniformly over the volumes in continuum units under
consideration. An identity obtained only after an ordered limit
would not justify the absorption lemma's uniform hypothesis.

\section{Collected estimates and local calibration}

The chapter now gives the quantitative single-insertion raw-map estimates
\eqref{ch11:graph-estimates}--\eqref{ch11:graph-difference},
including ball invariance and fixed-input continuity; the full
recurrence for terms of source degree at least two \eqref{ch11:higher-induction}; and the normalized
product-moment estimate \eqref{ch11:replicated-bound}.
The exact tree covariance \eqref{ch11:conditional-curvature} and the
local estimates in Section~\ref{ch11:actual-reader} identify the
joint Gaussian law used for the principal components.
The source weights and positive margins in the support exponent appearing in these
statements must be included in the selected stage requirements.

For a finite cutoff, composing the invertible single-insertion coefficient
maps from $N$ to the reference scale defines
$\mathsf Z_{N,p,d}$. It is real and symmetry-equivariant (the equivariance uses that of the insertion-free representations, which rests on Assumption~\ref{ch09:forest-assumption}, and the transport part of Hypothesis~(H) of Section~\ref{ch11:marked-step}; it is not proved here without them), and the exact
rule for the identity field gives its coefficient $a_{r(d)}^de_0$ at the reference scale.
The normalized local insertion is therefore
\[
 O_{N,p}(d,c;x)=a_N^{-d}
       P_d^{\rm loc}[S_{1/t_N}\mathsf Z_{N,p,d}^{-1}c](x),
\]
with the factor-degree scaling $S_t$ of the local section.
Uniform column estimates follow by multiplying the displayed
one-step inverse bounds; their subpower growth requires the
corresponding inequalities on the schedule.

There remains a distinct local calibration issue before claiming the
all-period finite-difference theorem. A translated finite difference
of degree $d$ has a protected local representative at degree $d+1$:
\[
 a_N^{d+1}\Delta_\alpha O(x)
 =a_N^dO(x+a_Ne_\alpha)-a_N^dO(x).
\]
This cancels the explicit $a_N^{-1}$, but moving the derivative-weight
origin contributes the factor $2^{\varpi_N q_N}$.
The lower bound on the starting scale of the stage must include
$2^{\varpi_Nq_N}\le\widehat\beta_N$ to keep that factor subpower.
The same requirement applies to repeated fixed finite differences.

Section~\ref{ch11:trial-descent} now gives the reference label $v_{N,p}$ of the trial, with
subpower cutoff-scale coordinates and a vanishing chessboard-seminorm
error, under its explicit terminal and stage bounds.
The microscopic equation section and its comparison with the dimensionful
section are given next. The argument with the separating family proves the finite-cutoff comparisons
needed for uniform differences.

% !TEX root = ../monograph.tex
\section{Calibrating the microscopic equations on the same history}
\label{ch11:equation-calibration}

Fix a degree $d$. The section $\mathcal V_d$ and automorphism $R_d$
are those of Theorem~\ref{ch02:equation-adapted}; neither is chosen
again when the cutoff or period changes. The physical section remains
$P_d^{\rm loc}$. We compare the two on the history already used to
define $\mathsf Z_{N,p,d}$.

We first specify an additional finite stage requirement.
At stage $j$, let $a_{dj}$ bound the derivatives of $\mathcal V_d$
requested by that stage, for every $d\le j$. These are finite constants
by the powered-curvature-factor estimate in Chapter~\ref{ch02:observables}.
Let $q_d$ bound the fixed union stencil and put
\[
 Q_{dj}=2+e^{81\sigma_j}(a_{dj}+C_j+2q_d).
\]
Here $\sigma_j$ is the support exponent and $C_j$ the column bound.
These quantities are fixed before the block length.
Let $\mathcal O_d\ge1$ bound both the finite symmetry actions on
coefficient $\ell^1$ norm and their averaging coefficients, and put
$\delta_d=(2\mathcal O_d^2)^{-1}$. Include the finite maximum
\begin{equation}
 L^{1/8}\ge
 \max_{d\le j}2(1+C_j)Q_{dj}\delta_d^{-1}
 \label{ch11:calibration-floor}
\end{equation}
in the stage-$j$ block floor. Include
$\widehat\beta_N\ge a_{d,J(N)}$ in the adjacent stage conditions on the starting scale. These requests are compatible with the selection order:
the section and its derivative constants are fixed first, then the
length, then the inverse coupling threshold. Raising a later threshold does
not change the section. Thus they can be imposed for all stages
before selecting the common history, with no dependence on $p$.

Define the normalized microscopic section at the origin by
\[
 O^{\mathcal V}_{N,p}(v)
   =a_N^{-d}\mathcal V_d[S_{1/t_N}v](0).
\]
The factor scaling $S_{1/t_N}$ commutes with $R_d$ only when
$R_d$ preserves factor count, which is not asserted here.
Accordingly its exact equation labels at this normalization are
\begin{equation}
 I_{N,d}^{\mathcal V}
    =S_{t_N}R_d\mathcal I_d.
 \label{ch11:scaled-equation-labels}
\end{equation}
Since $\mathcal I_d$ is homogeneous in factor count, one can equally
write $I_{N,d}^{\mathcal V}=
(S_{t_N}R_dS_{1/t_N})\mathcal I_d$.
This conjugation must be retained when using a cutoff-independent
unscaled equation section.

\begin{proposition}[One comparison map for all equation labels]
\label{ch11:equation-comparison}
Under \eqref{ch11:calibration-floor} and the terminal attachment
bounds of Section~\ref{ch11:trial-descent}, there is, for every
sufficiently large $N$ and every $p$, one real equivariant map
$P_{N,p,d}:V_d\to V_d$ such that
\[
 \|P_{N,p,d}-I\|_{1\to1}\le\tfrac12,\qquad
 \|P_{N,p,d}^{-1}\|_{1\to1}\le2.
\]
For each fixed stopping pair $(m,k)$,
\begin{equation}
 \mathcal R_m\!\left\{
 O^{\mathcal V}_{N,p}(v)
   -O_{N,p}(d,\mathsf Z_{N,p,d}P_{N,p,d}v;0)\right\}
 \le C_{d,m,k}\|v\|_1 E_N e^{-\tau_N/8},
 \label{ch11:equation-root-error}
\end{equation}
where $E_N$ is subpower in $a_N^{-1}$, uniformly in $p$.
The same map works for every spectator; it is constructed before
any spectator is specified.
\end{proposition}
\begin{proof}
For each basis vector, the exact terminal decomposition used in
Section~\ref{ch11:trial-descent} applies to the fixed section
$\mathcal V_d$. Its part with zero polynomial coefficient is localized on the union
stencil; the complementary part is attached to the original
large-field selection, with no change of density.
It gives coefficient $e_j$, ordinary norm bound
$Q_N=Q_{d,J(N)}$, and complete source norm bound
\[
 E_N=4(1+C_N)Q_N+
 e^{81\sigma_N+162}Z_N^{\rm val}\zeta_{d,N}u_N^{-1}.
\]
Here
$Z_N^{\rm val}=(A_{\mathcal V}+2A_{\rm loc})
 \sum_jt_N^{-s_j}$, with $A_{\mathcal V},A_{\rm loc}$
the global value norms of the two fixed sections.
The ordinary terminal norm bound uses a sharper attachment estimate
than the general trial bound. If $w_N^{\rm for}$ is the terminal
large-field bound and $D$ the union stencil, its contribution
to the remainder norm normalized by $\sqrt{w_N^{\rm for}}$ is
\[
 2|D|\sqrt{w_N^{\rm for}}\,
       e^{\sigma_N c_{B_N}(D)}Z_N^{\rm val}.
\]
Indeed attach the complementary local function to the original
nonempty prescribed large-field set and sum its possible distinguished
sites in $D$; the remaining large-field suppression bounds this
normalized mass. The same exact attachment retains reserved mass
$e^{(\sigma_N+2)c_{B_N}(D)}Z_N^{\rm val}$.
Impose the additional conditions on the starting scale
\[
 |D|w_N^{\rm for}\le1,\qquad c_{B_N}(D)\le81,\qquad
 \sqrt{w_N^{\rm for}}Z_N^{\rm val}\le1.
\]
Then $|D|\le q_d$ bounds the ordinary contribution by
$2q_de^{81\sigma_N}$, explaining exactly the definition of $Q_N$.
These are post-length threshold conditions. At fixed adjacent
stages the value norm is a fixed polynomial in the inverse
small-field threshold, whereas the large-field bound decays
faster than every such power. Their finite simultaneous threshold
therefore exists. The unweighted reserved mass is not discarded.
This last term is required: the ordinary derivative bound alone
does not bound the complementary large-field contribution.

Backward comparison gives a reference coefficient $v_j$ and the
chessboard-seminorm error in \eqref{ch11:trial-error}. Put
$u_j=\mathsf Z_{N,p,d}^{-1}v_j$.
There is a slightly stronger endpoint estimate than the displayed
bulk descent bound. At the last step its cone inequality gives
\[
 |e_j-u_j|\le
 4K_{N-1}L_{N-1}^{-3/8}c_N
 \le L_{N-1}^{-1/8}c_N
 \le 2(1+C_N)Q_N L_{N-1}^{-1/8}.
\]
The middle inequality follows from the previously imposed block
floor. Requirement \eqref{ch11:calibration-floor} bounds this by
$\delta_d$.

Let $Ue_j=u_j$. Average $U$ by conjugation with the finite geometric
symmetry group, and then with the real structure, obtaining $P$.
For the conjugation term use $\overline{U\overline v}$, so the result
is still complex linear. The norm of $P-I$ is at most
$\mathcal O_d^2\delta_d=1/2$, and the Neumann series gives the stated
inverse bound. Apply the same averaging to the field errors.
Linearity and the triangle inequality for the chessboard seminorm give \eqref{ch11:equation-root-error}, enlarging its fixed
constant to cover the finite orbit. This preserves one common
history, all source monomials, and one common map.

Finally $C_N$, $e^{81\sigma_N}$, $t_N^{-1}$,
$\zeta_{d,N}$, and $u_N^{-1}$ have the subpower bounds already
required for trial descent. The extra condition on the starting scale bounds
$a_{d,J(N)}$ by $\widehat\beta_N$, which is subpower.
Finite sums and products preserve this property, proving the claim
for $E_N$.
\end{proof}

The varying equation subspace in reference coordinates is therefore
\begin{equation}
 \mathsf E_{N,p,d}
   =\mathsf Z_{N,p,d}P_{N,p,d}
               S_{t_N}R_d\mathcal I_d.
 \label{ch11:actual-equation-subspace}
\end{equation}
For $v_N\in I_{N,d}^{\mathcal V}$ of subpower norm, the first term
in \eqref{ch11:equation-root-error} has exactly zero correlation
against every bounded spectator outside its microscopic stencil,
by \eqref{ch02:microscopic-null}. For fixed positive physical
separation, this exclusion holds at every sufficiently fine mesh,
uniformly in period. Holder's chessboard bound then shows that the
physical insertion with coefficient
$\mathsf Z_{N,p,d}P_{N,p,d}v_N$ has vanishing separated correlations
against any finite family of spectators with bounded chessboard seminorms. This proves the equation inclusion needed for the
argument using a finite separating family; equality with the full null space still
uses the free nondegeneracy and detection estimates.

The scaling in \eqref{ch11:actual-equation-subspace} also has a
useful check. Since $R_d-I$ never lowers factor count and $t_N\to0$,
the conjugate $S_{t_N}R_dS_{1/t_N}$ remains bounded.
Its strictly factor-raising terms tend to zero. Terms at equal
factor count can still raise scaling dimension and need not tend
to zero under this scaling alone. Thus the raw triangular
construction and the near-identity comparison $P$ do not, by
themselves, prove that this equation subspace approaches the
unmodified Maxwell ideal. No such identification is required here: use bounded projections
with kernel exactly \eqref{ch11:actual-equation-subspace} and establish
the bound from the separating family \eqref{ch11:detector-bound} on their complements.
The terminal map $\mathsf Z$ already includes the scaling-dimension
evolution; its physical convention has not been changed.


\section{The cutoff-dependent equation space and uniform finite differences}
\label{ch11:uniform-differences}

All eventual assertions in this section mean: there is $N_0$ such
that the assertion holds for every $N\ge N_0$ and every period $p$.
A bound after first fixing $p$ is insufficient.
Write $i=(N,p)$, $X_i=\mathsf Z_{N,p,e}$, and
\[
 K_i=X_iP_{N,p,e}S_{t_N}R_e\mathcal I_e.
\]
This is precisely the equation space
\eqref{ch11:actual-equation-subspace}. Every factor is invertible,
so $\dim K_i=\dim\mathcal I_e$ at each finite cutoff.
Choose the bounded equivariant projections $\Pi_i$ with kernel $K_i$
as in the construction of the separating family above, and fix
$A\ge\sup_i\|\Pi_i\|$. Their rank is
$r=\dim(V_e/\mathcal I_e)$, without any limiting rank assertion.

Call a coefficient family $z_i$ admissible here if
$\|X_i^{-1}z_i\|_1$ is bounded by a sequence subpower in $a_N^{-1}$,
independent of $p$. The inverse-column bound makes every bounded
coefficient family admissible. This class is closed under finite sums
and multiplication by scalar families of modulus at most one.

\begin{lemma}[Vanishing chessboard seminorms for the equation space]
\label{ch11:uniform-sweep}
If $z_i$ is admissible and $\Pi_i z_i=0$, its normalized local insertion
has chessboard seminorm tending to zero, uniformly in $p$, at every
sufficiently fine fixed cell scale. The assertion includes all its
finitely many signed and conjugated orientations.
\end{lemma}
\begin{proof}
Put $v_i=P_{N,p,e}^{-1}X_i^{-1}z_i$.
Its norm is subpower, and $v_i\in S_{t_N}R_e\mathcal I_e$.
The error in \eqref{ch11:equation-root-error} is therefore at most
$e^{-\tau_N/16}$ eventually, uniformly in $p$, after increasing
the fixed lower bound on $N$.

The microscopic comparison insertion itself has chessboard seminorm exactly
zero. Indeed, every other reflected copy lies outside its fixed
microscopic equation neighborhood once the mesh is sufficiently
fine relative to the fixed cell in continuum units. Apply
\eqref{ch02:microscopic-null} to one copy and the product of all
other copies. That product is bounded and measurable at each
finite cutoff, regardless of the size of its coefficients.
There is no limiting argument in this identity.
The smallest allowed period in continuum units is positive, so one lower bound
on $N$ ensures the required nonwrapping and separation for
every period. The triangle inequality for the chessboard seminorm proves the claim.
\end{proof}

In the following local comparison, $S_t$ multiplies each basis
coordinate by $t$ to its curvature-factor count, whereas $W_\rho$
uses scaling dimension. Neither is the coefficient map $\mathsf Z$.
The equation correction $R_d$ and its ideal $\mathcal I_d$ are
those of Theorem~\ref{ch02:equation-adapted}; the separate map
$P_{N,p,d}$ compares the two local sections.

\subsection{A common localization map for the separating family}
\label{ch11:common-extraction}

The map used to normalize free comparison columns cannot be chosen
independently of the polynomial that is integrated.
Here is the construction from the single-insertion representation.
Fix a polynomial degree $D\ge d$, to be increased below, before the
stage derivative order; require $D+1$ derivatives.
At the comparison scale $k$ put
\[
 s=\widehat\beta_k,\quad H=e^{(\log\log s)^2},\quad
 M=\lceil(\log s)^2\rceil,\quad B=M^2,\quad b=8B,
 \quad \gamma=\frac1{16(D+2)}.
\]
The original insertion norm uses radius $M$, whereas the restarted
representation uses radius $B$. These are different radii.

Let $A u$ be the actual local polynomial coefficient in dimensionful
normalization. The calibrated single-insertion bound controls the full
regular series by $C_k\|Au\|_1$.
Split it at original connected covering number $M$.
The part with covering number at most $M$ is contained in the
radius-$b$ patch: a connected chain of $M$ original-radius cubes
has diameter at most $4M^2+4M\le8M^2$.
The remaining part has norm at most
$\eta C_k\|Au\|_1$, with $\eta=e^{-eM}$, by retaining one
positive support exponent $e$.
Keep its original event indicators and all original large-field terms.
Only the near part is replaced below.

Use logarithmic link coordinates $U_e=\exp(\theta x_e)$,
$\theta=s^{-1/2}H$, in this patch. The Taylor polynomial through
ordinary degree $D$ of the near regular sum is well defined.
Indeed the series of derivatives is absolutely bounded by the
regular norm. Polarizing the directional derivatives bounds its
coefficient sum by $C_D(1+n_b)^D C_k\|Au\|_1$, where
$n_b=O(b^4)$ is the number of real link coordinates.
The Taylor remainder is bounded by
$C_k\|Au\|_1\|x\|^{D+1}$, up to a fixed constant.
The finite degree-$d$ localization by radial Taylor expansion of the near sum is its
prescribed local polynomial coordinate: all regular remainders have
zero extracted polynomial. Applying the extraction to a Taylor polynomial instead of the
holonomy curve itself changes this by at most
$C(1+b)^a\theta$ times its coefficient norm.
This follows from the local holonomy equation: the normalized
logarithm of a threshold holonomy and its linearized edge
integral differ by $\theta$ times a fixed-degree polynomial
in the link location. Only finitely many terms contribute
through degree $d$.

All constants just used are fixed before $s$.
A lower bound on the starting scale of the stage can impose
\[
 C_k\le H^\gamma,\qquad
 C_D(1+n_b)^D\le H^\gamma,\qquad
 C(1+b)^a\le H^\gamma.
\]
The second and third inequalities follow from
$\log b=O(\log\log s)=o(\log H)$; the first is a lower bound on the starting scale for
the already selected stage column bound.
The discarded-support factor $\eta$ and the nonlinear ray
error $\theta$ are smaller than every fixed inverse power of $H$.
Consequently the ordinary linear radial Taylor localization of the resulting
polynomial agrees with $Au$ up to relative error $H^{-1}$,
uniformly in $N,p,u$.

Convert links to linear sums of based plaquette curvatures in the
axial gauge on this same patch. Let $W_u$ be the resulting raw
polynomial in curvature divided by $g=\beta_{N,p,k}^{-1/2}$.
Each axial row has coefficient sum at most $4bg/\theta$.
The coupling band gives
$g/\theta\le\sqrt{8/7}\,H^{-1}$, so this row sum is eventually
at most one. Thus $\deg W_u\le D$ and
$\|W_u\|_{\rm coeff}\le H\|Au\|_1$, after the finite increases of the starting scale.
The construction is linear in $u$ because the Taylor polynomial
of the function itself is unique; it does not require a linear
choice of its separate activity representatives.

The raw polynomial also approximates the same normalized local insertion
on the restarted small-field region. Its threshold is
$\vartheta=s^{-1/2}e^{(\log\log s)^{3/2}}$.
Put $a=4b\vartheta$. The nonlinear axial link logarithm divided
by $\theta$ and the linear sum of plaquette curvatures divided
by $\theta$ differ by at most $C_Ga^2/\theta$, and their size is
at most $2a/\theta=H^{-1+o(1)}$.
Taylor's remainder and replacement of one monomial factor at
a time therefore give
\begin{equation}\label{ch11:raw-polynomial-error}
 |F_u^{\rm near}(U)-W_u(g^{-1}F(U))|
       \le H^{-D}\|Au\|_1.
\end{equation}
For example the Taylor term is at most
$4H^\gamma H^{(-1+\gamma)(D+1)}\|Au\|_1
 =4H^{-D-15/16}\|Au\|_1$.
The nonlinear coordinate error and discarded-support error
are smaller than any required inverse power.
For the solutions $u=u_{j,N,p}$ of the lift equations \eqref{ch11:simultaneous-detector-lifts} below, which satisfy
$\|Au\|\le H^{bd+c_d}$, choose $D>bd+c_d+v$ before the stage
orders to obtain error $H^{-v}$.

Let $K_k$ be the original same-insertion covariance of the
Gaussian comparison. Define the polynomial used in Wick calculations by
\[
 \mathscr P u=\exp\!\left(\tfrac12\Delta_{K_k}\right)W_u.
\]
This sign is deliberate: its evaluation with covariance whose
diagonal block is decreased by $K_k$ equals the Gaussian
evaluation of $W_u$ with the undecreased covariance.
The identity follows by multiplying the commuting finite
operators $\exp(\tfrac12\Delta_{K_k})$ and
$\exp(-\tfrac12\Delta_{K_k})$.
It remains valid for indefinite symmetric covariances.

Set $z=\theta/g\ge\sqrt{7/8}\,H$.
For $p_u(y)=W_u(zy)$,
\[
 (\exp(\tfrac12\Delta_{K_k})W_u)(zy)
   =\exp(\tfrac12\Delta_{z^{-2}K_k})p_u(y).
\]
The rescaled polynomial $p_u$ is the Taylor polynomial with
the factor $g/\theta$ in its axial substitution cancelled.
Its remaining axial rows have norm at most $4b$.
Since $(4b)^D\le H^\gamma$ eventually, its coefficient bound is
$4H^{3\gamma}\|Au\|_1\le H^\alpha\|Au\|_1$,
where $\alpha=4\gamma=1/(4(D+2))$.
The finitely many radial Taylor coefficient functionals likewise have coefficient
norm at most $H^\alpha$ after increasing the starting scale.
These bounds are needed before applying the Wick estimate.
The covariance itself has entry bound at most $H^{1/4}$ after
an increase of the starting scale, so the rescaled covariance entries are at most
$H^{-1}$ eventually.
For polynomials of degree at most $D$, the coefficient norm of
the difference between this Wick operator and the identity
is at most
\[
 H^{-1}w_D,\qquad
 w_D=\tfrac12D^2e^{D^2/2}.
\]
Indeed a contraction has coefficient norm at most $D^2/2$;
sum the finite exponential series, bounding its terms in absolute
value. This bound is independent of the number of ambient
coordinates because the polynomial coefficient norm already
sums their coefficients.

Let $\mathcal E$ send $v\in V_d$ to its polynomial, of weighted degree at most $d$, in the fixed radial jet variables, expanded in the fixed radial polynomial basis.
It is injective between finite-dimensional spaces, so one may choose $c_d\ge1$ with
$\|v\|_1\le c_d\|\mathcal Ev\|_{\rm coeff}$.
The extraction from $\mathscr Pu$ of the coefficients of low scaling dimension belongs to this basis:
the original covariance is a real spatial matrix tensored with
the identity on $\gLie$, so its contractions commute with the
adjoint action. Axial conversion and weighted truncation
also commute with that action. Define $Tu$ to be precisely
these basis coefficients, retaining the inverse coupling scaling above.
Using $\alpha=1/(4(D+2))$, the preceding coefficient bounds,
linear substitution, and the Wick estimate give
\[
 \|Au-Tu\|_1
 \le c_d\bigl(H^{-1}+D2^DH^{-1+\alpha(D+2)}
                  +w_DH^{-1+\alpha(D+1)}\bigr)\|Au\|_1
 \le\tfrac12\|Au\|_1
\]
eventually. Constants and conjugation are preserved at every
operation, hence $T1=1$ and $T\bar u=\overline{Tu}$.
This constructs the same real linear map for every column;
it is not an arbitrary replacement for $A$.
The inverse bound and simultaneous lift equations stated below
therefore apply to this common localization map onto the polynomial label space.

Finally let $\chi$ be the restarted small-field indicator on the
whole patch. Define $f_u^{\rm near,mask}$ to be only the near sum
with its original event indicators restored, including its constant term.
It equals $F_u^{\rm near}$ without its event indicator wherever $\chi=1$,
because the restarted threshold is smaller than the original one.
Then
\[
 f_u^{\rm near,mask}=\chi W_u+\chi(F_u^{\rm near}-W_u)
             +(f_u^{\rm near,mask}-\chi F_u^{\rm near})
\]
is exact. The second term has the bound
\eqref{ch11:raw-polynomial-error}; the third vanishes wherever
$\chi=1$ and is assigned to the original prescribed large-field set.
The distant regular sum and original large-field terms were
retained throughout. Substituting this identity into the finite
squarefree exponential preserves every higher source variable coefficient.
Thus the density, polynomial $\mathscr P$, common
localization map $T$, and normalized lifts all refer to the same
terminal insertion.


The common localization map may also be chosen equivariant under the
finite signed coordinate group. Average the raw polynomial
construction on domain and range before defining $T$.
The insertion, covariance and radial Taylor localization have the
corresponding covariance identities, so this preserves both
the exact density comparison and its extraction of the coefficients of low scaling dimension.
All errors are enlarged by at most a fixed finite symmetry
norm. Since their relative bounds tend to zero, choose the
threshold to make the averaged error at most one half.
Average with conjugation at the same time. This uses one
polynomial construction for all transformed columns, rather
than symmetrizing independently chosen lifted vectors.

At a fixed comparison scale $k$ the one-step coefficient bounds in the direction of the transfer, applied successively from the reference scale to $k$, give
$\|u\|_1\le D_k\|Au\|_1$, uniformly in $N,p$.
The relative half-error therefore implies
$\|u\|_1\le2D_k\|Tu\|_1$, so $T$ is injective and hence
invertible. Let $v_1,\ldots,v_r$ be real representatives of the
free separating family. Write $W_\rho$ for scaling-dimension
dilation and $S_t$ for factor-count dilation. With
$\rho=H^b$ solve all the equations
\begin{equation}\label{ch11:simultaneous-detector-lifts}
 S_{g/\theta}T u_{j,N,p}=W_\rho v_j.
\end{equation}
They have unique real solutions using the same $T$.
The coupling band gives $\theta/g\le2H$, and therefore
\[
 \|u_{j,N,p}\|_1
 \le2D_k\|S_{2H}W_\rho v_j\|_1.
\]
Their sum is a finite bound $B_k$ independent of cutoff and
period. No bound uniform in the comparison scale is needed.
Moreover $\|Au_j\|\le2\|Tu_j\|$ is bounded by a fixed
power $H^{bd+c_d}$, where the additional exponent $c_d$
depends on factor counts and the finite target list, not on
$b$. This is the growth factor used when choosing the
Taylor order and the geometric magnification.

Here $c_B(X)$ is the connected covering number at radius $B$.
The support exponent multiplies $c_B(X)$ in an exponential norm,
whereas a source weight is raised to the number of distinct
insertion indices. Changing one of these weights does not change
the coefficient or its original prescribed large-field set.

\subsection{Absorbing the larger source weight}
\label{ch11:weight-restoration}

A bound by one at the original source weight is not enough for the
comparison of separating correlations. The following argument absorbs the necessary change
of weight uniformly in the number of reflected cells. Write $n\le2$
for the number of insertions in each cell, fix $U\ge1/16$, and put
\[
 \omega_-=H^{1/16-U},\qquad \omega=H^{1/16}.
\]
For an activity with insertions, with support $X$ and source variable set $I$, its
weighted absolute bound contains $\omega_-^{|I|}e^{e_*c_B(X)}$.
Choose $e_*=e+2$ before constructing the expansion with insertions.
The same expansion, estimated at $e_*$, is used below; there is
no change of coefficients when its support exponent is lowered.

Each prescribed support is contained in the fixed enlargement of a
radius-$AB$ ball around its cell's finitely many insertion points, with $A$
chosen before the inverse coupling. A connected grid of radius-$B$ cubes
covers it with at most $C(1+A)^4$ cubes. Thus its support bound
at exponent $f$ is bounded by
$R_f=\exp\bigl(fC(1+A)^4\bigr)$, independent of $B$ and the period.
At successive stages these are finite pre-threshold constants.
The stage choices of sufficiently large starting scales impose a common bound $H^{\epsilon_j}$
on $R_{e+2}$, $C_*$, $L^4T_{e+2}$, and the fixed incidence
constants, with $\epsilon_j\downarrow0$.
Consequently every finite product of these constants is
$H^{o(1)}$ along the selected comparison scales.
This uniform lower bound on the starting scale, rather than treating a changing stage
constant as an absolute constant, justifies the exponent
comparisons below.
This bounded-cover assertion is used at the middle scales.
It is not an assertion that the separated opening sites
at distance $H^b$ fit in one radius-$AB$ ball.
The supports of their terms of source degree at least two are instead estimated by
\eqref{ch11:opening-higher-weight}.
At an earlier step of the \emph{upper interval}, from the comparison scale to the middle scale, a nonzero connected principal component
has radius at most $(6n+3)B$ around each of its own insertion points;
a term of source degree two with more widely separated insertion points is absent.
Only these nonzero components enter the principal support bound.

Let $Y_c$ contain all prescribed supports in cell $c$. Suppose
these sets have mutual distance at least $D'$, and every activity
support contains the sites of all its insertion points. Let $t(I)$ count
the cells represented in $I$, and let $t(X)$ count the sets $Y_c$
met by $X$. Then $|I|\le nt(I)\le nt(X)$.
For $t(X)\ge2$, a connected $B$-cover gives
\begin{equation}\label{ch11:dispersal}
 c_B(X)\ge\lambda t(X),\qquad
 \lambda=\frac{D'}{3(2B+1)},
\end{equation}
once $D'/3>2B+1$.
To see this, choose one point of $X\cap Y_c$ per represented cell,
and follow a path in its connected cover until distance
$\lfloor D'/3\rfloor$ from that point. A cube covers at most
$2B+1$ of the successive distance values. The resulting initial
path segments belong to disjoint distance balls, and a cube
cannot meet two of them. Summing proves the inequality.
Here connected cover means the connected cover used in the
activity norm; one cannot replace it by an arbitrary disconnected
cover of the insertion points.

\begin{lemma}[Regular remainders at the larger source weight]\label{ch11:restore-regular}
Suppose the regular remainder has incidence bound at most
$H^{\kappa-q}$ at $(\omega_-,e+2)$, with support containing its
insertion points, and $\lambda/\log H\to\infty$.
If $q\ge\kappa+Un+2$, its incidence bound at $(\omega,e)$ is at most
$H^{-1}$ for all sufficiently large $H$.
\end{lemma}
\begin{proof}
Separate activities meeting one cell from those meeting at least
two. The first class gains at most $H^{Un}$ on changing weight.
For the second class retain one additional unit of the support exponent and use
\[
 H^{U|I|}e^{-c_B(X)}
 \le H^{Unt(I)}e^{-\lambda t(I)}\le1
\]
eventually. Thus the new bound is at most
$(H^{Un}+1)H^{\kappa-q}\le2H^{-2}\le H^{-1}$.
The estimate is independent of the total number of cells.
\end{proof}

The support containment is preserved by the algebra.
Initially an insertion contains its site. A coarse activity
contains the projected union of the fine supports, and its source variable
set is their union. Prescribed support enlargement, event indicator splitting
and the connected regrouping all retain those projected points.
This induction applies to both regular and large-field activities.
For a window whose cell side at its lower comparison scale is at least
$(\log s)^{20}$ times the middle lattice spacing, whereas
$B\le4(\log s)^4$, the cell neighborhoods have
$D'\ge(\log s)^{20}/2$ after their fixed-radius enlargements.
With $H=\exp((\log\log s)^2)$ this proves
$\lambda/\log H\to\infty$, uniformly in period.
The condition concerns the distance between replica cells, not
between the two insertions within one cell.

\begin{lemma}[Large-field activities at the larger source weight]
\label{ch11:restore-forced}
Suppose the old large-field site bound is $\eta_-$ at
$(\omega_-,e+2)$, where $\eta_-$ decays faster than every inverse
power of $H$. Suppose each cell has at most $2^n$ additional
\emph{failed-cutoff factors} (the terms $-(1-\chi)Rz_I$ of the factors $1-(1-\chi)Rz_I$ in the proof below, which vanish where the cutoff $\chi$ equals one), supported in $Y_c$, whose coefficient
bounds before source variable weighting are at most $C_*=H^{o(1)}$ and whose
support bounds at exponent $e+1$ are at most $R_{e+1}=H^{o(1)}$.
Every nonzero such factor must meet an old large-field activity.
Then their exact connected regrouping has site bound at most
\[
 4H^{Un+2^n}\eta_-
\]
at $(\omega,e)$, for all sufficiently large $H$.
\end{lemma}
\begin{proof}
A failed-cutoff factor of nonconstant source monomial $I$ occurs in the exact
identity
\[
 \exp_*(\chi W-R)
   =\exp_*(\chi(W-R))\,\exp_*(-(1-\chi)R).
\]
Its last exponent is square zero for each source monomial.
Expand it as a product of factors $1-(1-\chi)Rz_I$.
Whenever such a factor is nonzero, its support contains a large-field indicator
covered by an original prescribed large-field set. Thus every surviving
connected component contains an old large-field activity.
This is a pointwise identity of the full density with insertions.

Put $F=\log H$. A new failed-cutoff factor has bound
$p_f\le H^{1/8}C_*R_{e+1}$ and assign it auxiliary bound
$w_f=e^{-F}p_f$. Assign an old activity $o$ the auxiliary bound
\[
 w_o=e^g e^{2^nFt(X_o)}p_o,
\]
where $p_o$ is its bound at $(\omega,e+1)$.
The original bound $p_o^-$ at $(\omega_-,e+2)$ gives
\[
 w_o\le e^gH^{(Un+2^n)t(X_o)}
                     e^{-c_B(X_o)}p_o^-.
\]
For activities meeting at least two cell neighborhoods,
\eqref{ch11:dispersal} absorbs the power of $H$.
For the remaining activities its factor is at most $H^{Un+2^n}$.
The old auxiliary site sum is therefore at most
$e^gH^{Un+2^n}\eta_-$.

Set $\delta=H^{-1/2}$ and choose
$e^g=\delta/(2H^{Un+2^n}\eta_-)\ge1$.
The new factors have auxiliary site sum at most
$2^nC_*R_{e+1}H^{-7/8}$, so the total auxiliary sum
$\Gamma$ is at most $\delta$ eventually.
Assign each failed-cutoff factor to an old activity it meets.
At most $2^nt(X_o)$ factors are assigned to $o$.
In each nonzero connected component, their discarded
$e^F$ factors are thus bounded by the $e^{2^nFt(X_o)}$
factors, leaving at least one factor $e^{-g}$.
The connected-set estimate of Chapter~\ref{ch08:chapter}
gives the site sum at exponent $e$ at most
\[
 e^{-g}(e^\Gamma-1)
 \le2\delta e^{-g}=4H^{Un+2^n}\eta_-.
\]
Its counting condition $(2B+1)^4\delta\le1$ follows from
$B=O((\log s)^4)$ and the definition of $H$.
No independence of large-field indicators or activities is used.
\end{proof}

These lemmas explain the required stronger support estimate.
For the same expansion on the upper interval let its regular remainder satisfy
\[
 z_0\le H^{-q},\qquad z_{j+1}\le(\alpha+2)z_j,\qquad
 (\alpha+2)^{k-\ell}\le H^\kappa
\]
at exponent $e+2$ and weight $\omega_-$.
This implies the premise of Lemma~\ref{ch11:restore-regular}.
The constants at exponent $e+2$ must be fixed before $\alpha$
and the block ratio. Applying an estimate proved only at exponent
$e$ a second time does not establish this premise.

If the corrected principal coefficients at the middle scale are at
most $H^{-31/32}$, their contribution to the incidence bound at the larger source
weight is at most $2^nR_eH^{-27/32}$.
Together with the two preceding lemmas this gives regular
bound
\[
 H^{-1}+2^nR_eH^{-27/32}\le H^{-3/4}
\]
and large-field bound at most $H^{-q}$.
The reference connected polynomial is subtracted once, on the
same sample and with the same subset moments.
Its large-field part is precisely the part handled above.

On the lower interval suppose the step with insertions has the
quantitative bound
\begin{equation}\label{ch11:lower-weighted-step}
 z_{j+1}\le\alpha z_j+H^{1/2}z_j^2
\end{equation}
at the larger source weight, and its large-field part stays in the
allowed interval. If
$(\alpha+1)^{\ell-w}\le H^{1/4}$, induction gives
\[
 z_j\le(\alpha+1)^jH^{-3/4}\le H^{-1/2}.
\]
Thus every step remains in the domain
$H^{1/2}z_j\le1$ used in the induction.
At the output the source variable incidence and large-field site bounds
are at most one. On a torus with $q_0$ sites and $nq_0$ insertions,
the scalar majorant of the exponential with insertions gives
\begin{equation}\label{ch11:weighted-output-envelope}
 \omega^{q_0|I|}|[z_I^{\rm all}]\rho_{\rm marked}(U)|
 \le\exp\bigl(\text{unmarked exponent}(U)+q_0(1+n)\bigr).
\end{equation}
Every cell carries the source monomial $I$ in the coefficient on the left.
The integral without insertions is still the original partition function
by \eqref{ch11:source-density}.
This gives the numerator needed in the finite separating family
comparison, once the stronger upper recurrence and
\eqref{ch11:lower-weighted-step} have been proved for the combined estimate with insertions. The bound by one of \eqref{ch11:higher-induction} for terms of source degree at least two has neither
the source weight nor the small initial remainder needed here.

\paragraph{The two weighted step estimates.}
The inflations required here depend on the inverse coupling, so the
fixed-inflation statement alone is insufficient.
Split the decay exponent $Kh(x)^2r$ of a selected component with $r$ selected cutoff conditions, given by
Chapter~\ref{ch06:large-fields} (Section~\ref{ch11:restart-construction}), into fractions kept aside and choose
the reported bound $G(x)=\exp\{-a_*h(x)^2\}$ with $4a_*\le K/2$ so that multiplication by $G(x)^{-2}$
leaves a fixed positive fraction. This choice is made before the
thresholds. The amplitudes without prescribed large-field sites have a negative inverse coupling power
times subpower factors; multiplying by $H^c$ for fixed $c$ retains
that power, since $\log H=o(\log x)$. These margins establish the common
product bound \eqref{ch11:integrated-component} for the following
inflations, including their product over distinct prescribed large-field sets.

Use $e_*=e+2$, put $T=T_{e_*}$ and choose $\alpha\ge8T$.
For $H^{-q}\le z\le H^{-4}$ and principal incidence at most $H^{-4}$,
take
\[
 l=\max\left(1,\frac{H^{-3}}{\alpha z+H^{-q}}\right),
 \qquad m=G(x)^{-2}.
\]
Then $l\le\max(1,H^{q-3})$ and $l/x\le1$ eventually.
Thus the actual probability of a large fluctuation-field event absorbs the complementary
principal cutoff. The insertion incidence satisfies
\[
 a=lTz+2TH^{-4}\le\tfrac12H^{-3}.
\]
If $l>1$, use $lTz\le(T/\alpha)H^{-3}$; if $l=1$, use
$z\le H^{-4}$. Since $\mathcal W(B)=H^{o(1)}$ and $G(y)$
decreases faster than every inverse power of $H$, choosing
$\delta=H^{-1/8}/16$ gives $\Sigma\le H^{-1/8}/4$ and
$TG(y)+a\le H^{-3}$.
The connected bound and division by $l,m$ give
\[
 z_{\rm nonprincipal}\le H^{-3}/l
       \le\alpha z+H^{-q}\le(\alpha+1)z,\qquad
 z_{\rm principal}\le H^{-3},\qquad
 z_{\rm forced}\le G(x)^2.
\]
When $l=1$, its definition implies
$H^{-3}\le\alpha z+H^{-q}$, so the first estimate still holds.
Reassigning a principal term to its prescribed enlarged support
adds at most $2^nR_{e_*}y$ if its coefficient difference is $y$.
The joint Gaussian estimate at the requested power gives
$2^nR_{e_*}y\le H^{-q}$.
Its large-field part retains the original
prescribed large-field set. With $e^{-g_i}=G(x)$ the bound
\[
 e^{g_i}G(x)^2+nH^{-2}R_{e_*+1}\le H^{-1}
\]
gives outgoing large-field bound
$G(x)(e^{H^{-1}}-1)\le G(x)$.
This proves the upper recurrence
$z_{\rm new}\le(\alpha+2)z$ at the stronger exponent.

Below the middle scale there are no principal factors.
Use exponent $e$, $l=1$, $m=G(y)^{-1}$ and incoming large-field
bound at most $z\,a_{\rm low}(y)$, where
$0\le a_{\rm low}(y)\le1$. The outgoing incidence is at most
\[
 T_ez(1+a_{\rm low}(y))
 \le2T_ez\le\tfrac12\bigl(\alpha z+H^{1/2}z^2\bigr).
\]
The site conditions hold for $H^{-q}\le z\le H^{-1/2}$
by the same estimates, and the large-field output is at most $G(y)$.
This proves \eqref{ch11:lower-weighted-step}.
For the polynomial attenuation used here, the exponential
large-field decay kept aside gives $G(y)\le H^{-q}a_{\rm low}(y)$. Since
$z\ge H^{-q}$, the outgoing large-field bound is at most
$z\,a_{\rm low}(y)$, as required for the next step.

This calculation requires the decay kept aside and the inflated product bound
\eqref{ch11:integrated-component} for the factors without insertions of the restarted representation.
It does not transfer a bound from a different Gaussian law or
change an incoming bound after the history is fixed.
The restarted representation keeps the original path and clipping coefficients and Gaussian law (Section~\ref{ch11:restart-construction}).

The parameter $H=e^{(\log\log s)^2}$ measures the slowly growing
comparison scale, not a covariance or a Hilbert space.
The Gaussian fluctuation scale at the inverse coupling $\beta$ is $g=\beta^{-1/2}$.
The density $q_k$ with insertions remains the full transferred polynomial of
\eqref{ch11:source-density}, including its higher coefficients;
it is unrelated to the fixed exponent $q$ of the initial remainder bound.

\subsection{The covariance on a comparison window}
\label{ch11:window-covariance}

We record the covariance estimates used in the comparison of separating correlations,
including their finite-period errors. Fix enough smoothing orders for
the finite number of differences below. Let $H(p)$ denote the original
two-form metric symbol of Chapter~\ref{ch05:gaussian}, with
$H(0)=I$, a positive lower bound, and uniform exponential coefficient
bounds. Put $q_\mu(p)=e^{ip_\mu}-1$, let $X$ be exterior
multiplication by $q$ from one-forms, and let $U$ be exterior
multiplication from scalars. The curvature covariance is
\[
 P(p)=X(p)\bigl(X(p)^*H(p)X(p)+U(p)U(p)^*\bigr)^{-1}X(p)^*,
 \qquad P(0)=0.
\]
Indeed the completed inverse differs from the original potential
inverse only in gradient directions, annihilated by $X$.
The exterior identity gives coercivity at least
$c|q|^2$. Taylor expansion at zero, with differentiated bounds, yields
\[
 D^jP(p)=O(|p|^{-j}),\qquad
 D^j(P(p)-P_0(p))=O(|p|^{1-j}),
 \quad P_0=X_0|p|^{-2}X_0^*,\quad X_0=ip\wedge.
\]
All constants are uniform in the kernel history within the
already fixed quadratic bounds.

Use a smooth compact frequency partition whose $2\pi\ZZ^4$
translates sum to one, equal to one near zero. This leaves the
finite Fourier sum unchanged, including its zero-frequency value.
For difference orders $j,h$, subtract the Gaussian-regularized
continuum multiplier. The remainder has order $j+h+1$ at zero
and rapidly decreasing derivatives away from zero.
Splitting its Fourier integral at radius $|x|^{-1}$ and
integrating by parts $j+h+6$ times outside that ball gives
$C|x|^{-j-h-5}$. The inner integral has the same bound.
The finite Fourier sum differs from its integral by at most
$C(1+|x|)/N_{\rm per}$: on frequency cubes of side
$N_{\rm per}^{-1}$ use the derivative bound $C/|p|$,
which is integrable in four dimensions, and the phase difference.
This also handles the frequency cube containing $p=0$.
Gaussian regularization of the continuum potential has kernel
$(4\pi^2|x|^2)^{-1}(1-e^{-|x|^2/4})$ in radian convention.
Differentiating it therefore changes the separated Maxwell
covariance only by an exponentially decreasing term.
Consequently, for $|x|\ge1$,
\begin{equation}\label{ch11:finite-covariance}
 |\Delta^j\Delta^h\kappa_k(x)-\partial^j\partial^h C_{\rm Max}(x)|
 \le C|x|^{-j-h-5}+C(1+|x|)/N_{\rm per}.
\end{equation}
The lattice-side lower bound for the smallest period in continuum units
makes the second term negligible on every fixed-power-$H$ window,
uniformly in all periods.

For the installed curvature response $\Psi_i$, one also
needs to move differences through the curvature response, not merely bound
its differentiated row. The box-multiplier identities in the alias
formula give
\[
 \Delta_w\Psi_i=Q_{i,w}\Delta_w,\qquad
 \sup_x\sum_y|Q_{i,w}(x,y)|\le R L_i^{-|w|-2}.
\]
On the principal alias divide by the coarse difference using
$(e^{ip/L}-1)/(e^{ip}-1)=
[L\,\operatorname{avg}_L(p/L)]^{-1}$.
On other aliases consume one box factor for each difference;
their geometric-series identity has the same coarse factor.
Sufficient smoothing order leaves the remaining powers
nonnegative. The quotient is regular also on coordinate
hyperplanes. Its exponential coefficient bound follows at a
narrower decay width from
$(e^{inp}-1)/(e^{ip}-1)=\sum_{a=0}^{n-1}e^{iap}$
for $n>0$, with the analogous negative sum for $n<0$.
The extra factor $|n|$ is absorbed by that strict margin in the decay width.
This proves the row bound with $R$ fixed before $L$.
The curvature response preserves closed cochains, including harmonic ones,
by the exact exterior-transport identities of Chapter~\ref{ch05:gaussian}.

Let $\mathcal L=\Psi_{k-1}\cdots\Psi_\ell$,
$m=k-\ell$ and $S=a_\ell/a_k$.
A sign-valued test function on a specified row, without requiring that function to be closed, gives
\[
 S^{j+2}\sum_y|\Delta_w\mathcal L(x,y)|
       \le2^d R^m\quad(j=|w|\le d-1).
\]
Choose $m=\lceil(b+2)\log H/\log L\rceil$ with $b\ge6$,
and choose $L$ after $R$ so that
$4(b+2)\log R\le(\lambda/4)\log L$.
Then $S\ge H^{b+2}$ and the right side is at most
$P=H^{\lambda/4}$ after an increase of the starting scale.
The middle-scale field gives such a closed background field,
rather than an independently chosen one. In axial gauge on its
radius-$B$ patch take $g^{-1}$ times the Lie-algebra curvature
coordinate of each link, extend this one-cochain by zero, and
take its discrete exterior derivative $u$. Then $du=0$ exactly,
including across the artificial patch boundary, and
$\|u\|_\infty\le16B\vartheta/g=H^{o(1)}\le P$.
The truncation at the patch boundary is included in the
comparison error. Retain only response paths lying in the protected
patch and put the omitted paths in the regular remainder.
The exponential response bound makes their absolute contribution
at most $C B^a e^{-cB}g^{-1}$, with the same bound for every fixed
requested difference. The clipped field outside the patch is
bounded before multiplication by $g^{-1}$.
Since $B$ grows as $(\log s)^4$, this error is smaller than every
fixed inverse power of $H$, even after $g^{-1}=O(\sqrt s)$.
Thus zero extension is not asserted to preserve the full
conditional mean exactly.
On the interior stencil its difference from the nonlinear normalized
plaquette curvature is $O(B^2\vartheta^2/g)$, smaller than every
fixed inverse power of $H$. This follows by subtracting the
linear sum from the four-link product in that same gauge.
For this closed middle-scale background field $u$,
the mean $\mu=\mathcal Lu$ and the conjugated coarse covariance
$B_0=\mathcal L\kappa_\ell\mathcal L^T$ satisfy
\[
 S^{j+2}|\Delta^j\mu|\le P^2,\qquad
 S^{j+h+4}|\Delta^j\Delta^hB_0|\le P^3.
\]
The exact conditional covariance identity
\eqref{ch11:conditional-curvature} telescopes through the window.
After the same-insertion normal ordering used for
$\mathscr P$, its covariance is
\[
 \Gamma_{ij}(z,w)=\mathbf1_{i\ne j}
       \kappa_k(z_i+z,z_j+w)-B_0(z_i+z,z_j+w).
\]
In particular its diagonal is $-B_0$, not zero.
It need not be positive; only the polynomial Wick evaluation
uses this covariance.

Take the measurement magnification $\rho=H^b$.
The preceding bounds, \eqref{ch11:finite-covariance}, and
$\rho/S\le H^{-2}$ imply weighted mean and covariance errors
at the retained jet coordinates at most $H^{-3}$.
The highest difference and remainder coordinates have bound $P$.
Here the continuum approximation error contributes
$O(\rho^{-1})=O(H^{-b})$, the mean contributes
$H^{-4}P^2$, and the conjugated covariance contributes
$H^{-4}P^3$. The finite-period term is smaller than these
after the common lower bound on $N$.

To turn these into one finite coordinate system, choose dual
finite-difference functionals $\ell_\alpha$ on the polynomial
closed curvature cochains $t_\alpha$, with
$\ell_\alpha(t_\beta)=\delta_{\alpha\beta}$.
They exist by downward induction from the highest nonzero homogeneous
curvature, exactly as in the radial-jet construction.
Retain coordinates $\rho^{g_\alpha}\ell_\alpha f$ and,
at every stencil point $z$, the residual coordinate
$\rho^{d+1}(f-\sum_\alpha\ell_\alpha(f)t_\alpha)(z)$.
This gives an exact reconstruction of every stencil value.
For a closed cochain, repeated finite-difference summation along
coordinate paths bounds its residual by the top difference times
a fixed power of the stencil radius.
These factors are $H^{o(1)}$ and are absorbed by $H^{\lambda/4}$.
The common enlarged mean and covariance are therefore bounded by
$H^\lambda$, and their retained jet block differs from the free
one by at most $H^{-2}$. The free comparison has zero residual
coordinates. All subsets use this one enlarged system.

\subsection{From the common localization map to the corrected pair bound}
\label{ch11:corrected-pair-comparison}

Solve the simultaneous lift equations using the map $T$
constructed in Section~\ref{ch11:common-extraction}.
For the finitely many free polynomials, their coefficient sizes
are fixed. The solved lift has a bound $H^{bd+c_d}$ in normalized local insertion
coordinates. Choose $b$ larger than $c_d+3$, then choose $D$
large enough for \eqref{ch11:raw-polynomial-error}; all these
choices precede the derivative orders and final scale thresholds.
After the common finite-coordinate substitution, the retained
weighted polynomial is exactly the intended free polynomial.
Every remaining monomial has scaling dimension at least $d+1$, so its
coefficient sum is at most $H^{-2}$ by the factor
$\rho^{-(d+1)}$ and this choice of $b$.

For completeness, if a polynomial has coefficient norm $A$ and
degree at most $D_0$, Wick recursion bounds its expectation by
$A[2B(D_0+1)]^{D_0}$ when all means and covariance entries are
bounded by $B\ge1$. Changing these entries by at most $\epsilon$
changes it by at most the same expression times $\epsilon$.
This is proved by choosing the first variable of a monomial
and summing its mean term and at most $D_0-1$ pair terms;
subtract products one factor at a time.
It is independent of the number of coordinates.

Set $\lambda=1/[4(2D+D+1)]$ and $B=H^\lambda$.
The common coordinate estimate and the coefficient tail just
proved give, for every subset of the two insertions, moment
error at most $C H^{-2}H^{2D\lambda}\le H^{-3/2}$.
The empty moment is exactly one.
The nonlinear-to-polynomial error has the same bound:
use \eqref{ch11:raw-polynomial-error}, the
Gaussian estimate under the joint law with a smooth cutoff, and telescope the finitely many polynomial
transfer steps. If a step has coefficient norm at most
$A_{2D}$ and difference at most $H^{-v}$, the product error
over $m$ steps is at most $mH^{-v}A_{2D}^m$.
Choose the block length after $A_{2D}$ and then $v$ so this is
at most $H^{-3/2}$. The local-coordinate error is a negative
inverse coupling power times subpower factors and hence is below every
fixed $H^{-v}$ on the comparison band.
The comparison uses the same mean, covariance and common
cutoff factors at every subset; independently chosen Gaussian
samples would not justify the product estimate.

For two source variables the connected coefficients are simply the two
means and the covariance of the polynomial values. Subtraction
of these expressions gives a fixed polynomial bound in the
subset moments times their errors. Their free values remain
bounded on the chosen separated configurations.
Thus the corrected principal coefficients at the middle scale
are at most $H^{-31/32}$, with room to spare.
Here is the explicit estimate for the terms of source degree at least two already present
at the opening scale. Write $C=\lceil(\log s)^2\rceil$, $B=C^2$,
let $\theta_0>0$ be their original support exponent, and let
$r_0=s^{-a/2}$ be their original source variable radius. Their normalized
rooted bound is at most one by \eqref{ch11:higher-induction}.
The new radius is $r=H^f\le1$, where $f=1/16-U\le0$.
For a support containing $m\ge2$ distinct insertion points, the
chosen configurations have pairwise distance at least $C^2$.
The separated-site bound and cover enlargement give
\[
 Cm\le64c_C(X),\qquad
 (e_*+2)c_B(X^{+8})
 \le100000(e_*+2)+\tfrac12\theta_0c_C(X).
\]
For any prescribed $q$, set
$\Gamma=(q+10)\log H$ and $u_0=(a/2)\log s$.
The stage thresholds can impose
\[
 64(u_0+\Gamma+100000(e_*+2))\le C\theta_0/2,\qquad
 100000C(e_*+2)\le B\theta_0/2.
\]
This is possible because $C$ grows like $(\log s)^2$,
whereas $u_0=O(\log s)$ and $\log H=o(\log s)$.
Consequently, for each term of the original expansion with insertions,
\begin{equation}\label{ch11:opening-higher-weight}
 r^m e^{(e_*+2)c_B(X^{+8})}
 \le e^{-\Gamma}r_0^m e^{\theta_0c_C(X)}.
\end{equation}
This changes weights on the full family of terms of source degree at least two;
it does not compare $r$ directly with $r_0$.
Root-to-site conversion contributes only the factor $\mathcal W(B)=H^{o(1)}$.
At the opening all distinct sites, including those in one
measurement cell, are separated: their distance is a fixed
positive multiple of $H^b$, and $C^2=o(H^b)$.

Choose $U$ so the normalized single-insertion bound at radius $H^f$
is at most $H^{-8}$. Its near Taylor error is at most
$H^{f+bd+c_d-D}$, and its distant-support part is at most
$H^{f+bd+c_d}e^{-\theta_0 C}$. Increase $D$ to make both
at most $H^{-q-2}$. Together with
\eqref{ch11:opening-higher-weight} these give the initial
regular remainder $H^{-q}$ at exponent $e_*$.
For the large-field part perform the exact change of representation at the comparison scale
jointly with these factors with insertions, before summing components.
Every nonzero failed cutoff and every old large-field activity
meets a new large-field site. Its selected site factor has size at
most $e^{-2u}$, with $u=a_*h(s)^2$.
Keep $e^{-u}=G(s)$ aside once per component.
The remaining rooted sum is bounded by
\[
 H^{-q-8}
 +\mathcal W(B)\bigl(O(H^{-8})+O(H^{-q-8})\bigr)
 +H^{o(1)}\sqrt{e^{-H^{1/5}}}
 +(8B+1)^4e^{-u+O(e_*)},
\]
which is at most $H^{-1}$ and satisfies the connected-set
smallness condition after the increases of the starting scales.
Thus its regrouped site sum is at most
$G(s)(e^{H^{-1}}-1)\le G(s)$.
All original and new prescribed large-field sets and compatibility conditions are
retained, and component decomposition is reversible in every
squarefree coefficient. This proves the full opening bound;
no higher coefficient of the transferred density is discarded.

The upper interval has
$m=\lceil(b+2)\log H/\log L\rceil$ steps.
Before choosing $L\ge3$, fix
$\kappa=1+2(b+2)\log(\alpha+2)/\log3$.
Then $(\alpha+2)^m\le H^\kappa$ eventually for every allowed $L$.
Choose $q\ge\kappa+Un+2$ before the Taylor order $D$ and
the later lower bound on the block ratio involving its polynomial transfer norm.
This prevents a circular choice between $q,D$ and $L$.
The lower interval only increases the cell side by a factor
$C^{10}$, so it uses at most
$\lceil10\log C/\log L\rceil$ steps.
Since $\log C=O(\log\log s)=o(\log H)$, its growth factor
$(\alpha+1)^{\ell-w}$ is at most $H^{1/4}$ eventually.
The selected stage plateau must include these two finite windows;
their total length is $O((\log\log s)^2)$ at a fixed stage
and lies inside the short interval admitted by the restart
construction.

Sections~\ref{ch11:weight-restoration} and
\ref{ch11:common-extraction}, together with the restarted
representation, now give the weighted output majorant
\eqref{ch11:weighted-output-envelope}.
To normalize it one must retain the Gaussian power in the
partition function. Axial gauge fixing leaves $3V+1$ free links
on a four-dimensional torus of $V$ sites. Restrict them to a
ball for the normalized Hilbert--Schmidt distance of radius $g/4$ to obtain
\[
 Z\ge e^{\sigma-Vc_1}(c_-g^{\dim\gLie})^{3V+1}.
\]
For the numerator, successive axial plaquette elimination
integrates $3V-3$ distinct links, leaving the four wrapping
links. An explicit elimination order verifies that count.
Represent sites by $0,\ldots,N_{\rm per}-1$ and let $p(x)$
be the largest nonzero coordinate. The tree uses
$(x-e_{p(x)},p(x))$ for $x\ne0$.
Retain the four wrapping links $(-e_i,i)$.
For any other free link $(x,i)$ order it by
\[
 \bigl(t,\mathbf1_{\{i=t\}},\textstyle\sum_jx_j\bigr),
 \qquad t=\max(i,p(x)),
\]
lexicographically. If $i<t$, choose its plaquette in directions
$(i,t)$ based at $x-e_t$. If $i=t$, freeness implies
$x_i=N_{\rm per}-1$; put $w=x+e_i$, take $j=p(w)<i$,
and choose the $(j,i)$ plaquette based at $x-e_j$.
Every other free link in the chosen plaquette has smaller priority.
The chosen plaquettes are distinct, since equality would require
each of two links to have smaller priority than the other.
Integrate in decreasing order. The chosen link occurs once,
as $aUb$ or $aU^{-1}b$, so Haar invariance evaluates its factor.
All other plaquette factors are at most one and may be discarded.
This proves the stated $3V-3$ integrations, including the
wrapping links except for the four explicitly retained ones.
The Haar estimate
$\int_Ge^{-\chi(u)/(g/c)}du\le c_+(g/c)^{\dim\gLie}$
then gives the corresponding upper bound.
It follows that the normalized coefficient with weight $z$
obeys
\[
 z|\mathbb E F|
 \le \exp\bigl(VC_{\rm stage}-4\dim(\gLie)\log g\bigr).
\]
The constant term $\sigma$ cancels. A genuine additive value error
$\delta$ contributes $2V\delta$, one copy in numerator and
one in denominator, and is included in $C_{\rm stage}$.
These estimates use the same nonnegative output density.

At the output $V=q_0$. Its smallest-period site count grows
exponentially in the output logarithmic scale, so
$2\dim(\gLie)\log\beta/q_0\le1$ eventually.
Choose the later lower bound on the starting scale with
$C_{\rm stage}+1\le I_k=\log H/32$.
Applying the preceding estimate to
\eqref{ch11:weighted-output-envelope} gives
\[
 \|\hbox{corrected coefficient }I\|_{\rm root}
 \le e^{I_k}\omega^{-|I|}
 \le e^{-I_k},\qquad \varnothing\ne I\subseteq\{1,2\}.
\]
This is the desired corrected pair estimate on the original
partition denominator. At each fixed deep comparison scale its
simultaneously normalized columns are uniformly bounded in $N,p$.
Undoing the finite connected subtractions and then taking a
sufficiently deep comparison scale establishes
Proposition~\ref{ch11:finite-detector-input}, using the small value and joint product bounds of
Proposition~\ref{ch11:restart-family}. The estimates hold before any subsequence is chosen and their
thresholds are independent of the period.


\subsection{Finite-arity comparison on a common finite set of scales}
\label{ch11:finite-arity-comparison}

The pair argument extends to any fixed arity, with the same
partition normalization and one polynomial localization map for each degree.
We state the result in the form needed for smooth free-field
comparison. All constants may depend on the fixed finite arity,
degree list and compact configuration set.

\begin{proposition}[Common finite-arity Gaussian comparison]
\label{ch11:n-mark-comparison}
Fix $n\ge1$, finitely many real polynomial labels of bounded
degree, and a compact set of point configurations with positive
pairwise distance. Fix finitely many magnification exponents
$b_1,\ldots,b_t$ large enough for the preceding local estimates.
The selected stage requirements can be imposed for this finite
request so that one common localization map $T_{k,N,p,d}$ onto the polynomial label space
per degree and one restarted history serve all these scales.
At scale $\rho_j=H^{b_j}$ use the exact simultaneous lifts
\[
 S_{g/\theta}T_{k,N,p,d}u=W_{\rho_j}v.
\]
They are real and equivariant, and bounded in $N,p$ for each
fixed comparison scale. Set $G_\varnothing=1$ and let $G_I$
be the free moment of the chosen labels at the rescaled
configurations. Define the corrected coefficients, the Wick products of the $F_i$ with respect to the free field, by the recursion
\[
 R_\varnothing=1,\qquad
 R_I=\prod_{i\in I}F_i-
       \sum_{\varnothing\ne J\subseteq I}G_J R_{I\setminus J}.
\]
Then, for every nonempty $I\subseteq\{1,\ldots,n\}$, the
chessboard bound on the corresponding comparison cell is
\begin{equation}\label{ch11:n-mark-root}
 \|R_I\|_{\rm root}\le H^{-1/(16n)}
\end{equation}
at all sufficiently large cutoffs, uniformly in period and
in the stated compact configuration set.
The cell side can be chosen at most $a_kH^{b_j+5}$.
The estimate requires a positive distance between the points, but no
positive lower bound on individual coordinate differences.
\end{proposition}
\begin{proof}
Make all finite requests before choosing the stage orders,
ratios and starting scales. Construct $T$ at the comparison scale $k$ from
the insertion and its original same-insertion covariance.
Neither its near-support Taylor construction nor that covariance
depends on the choice of middle or output scale.
Choose the Taylor order and derivative margins for the
largest of the finitely many requests. Restart once on the
longest required window. Restriction of that history gives
the shorter windows, with unchanged localization map and density at the comparison scale.
There is no new normalization for individual spectator subsets.

Work throughout in the full $n$-source squarefree algebra.
The polynomial $\sum_I G_Iz_I$ has constant coefficient one,
so its inverse exists there. Multiplying the density with insertions
by this inverse gives exactly the displayed recursion.
Equivalently subtract its finite connected logarithm.
This operation commutes with every block integration and leaves
the density without insertions, hence its partition function, unchanged.

Use
\[
 \omega_-=H^{1/(8n)-U},\qquad
 \omega=H^{1/(8n)},\qquad I_k=\frac{\log H}{16n}.
\]
All support and regrouping arguments above already use $n$
only through $n$, $2^n$, and the fixed polynomial degree.
Choose $U\ge1/(8n)$ to control the finite list of single insertions,
then $\kappa,q,D$, and finally the ratios and starting scales in the
order already specified. The initial estimate for terms of source degree at least two
\eqref{ch11:opening-higher-weight} applies to all cardinalities;
the full incoming polynomial is retained.

For the Gaussian calculation under the joint law use degree bound $nD$
and, for example, $\lambda\le1/[32n^2(D+1)]$.
The coefficient tails and the covariance errors on the retained jet block are
$H^{-2}$, so the polynomial Wick estimate makes all subset
moment errors at most $H^{-3/2}$.
The finite logarithm is a polynomial of degree at most $n$
in these subset moments. Its derivative is bounded on a fixed
neighborhood of their free values, since every free subset
moment is bounded on the compact separated configuration set.
Hence the connected coefficient errors are also $O(H^{-3/2})$
and in particular at most $H^{-31/32}$ eventually.
This gives the same corrected principal bound as before.

Each source monomial gains at most
$\omega^{|I|}\le H^{1/8}$ at the larger source weight.
Lemma~\ref{ch11:restore-regular} therefore gives $H^{-1}$;
the corrected principal contribution is at most
$2^nR_eH^{-27/32}$; and the regrouping of activities constrained to cover their prescribed large-field sets has bound
$4H^{Un+2^n}\eta_-$.
The upper and lower recurrences then give incidence and site
bounds at most one at the output. Their scalar majorant contributes
at most $1+n$ per site. The partition
estimate absorbs this finite factor by imposing
$C_{\rm stage}+1\le I_k$.
Thus the corrected coefficient has chessboard seminorm at most
$e^{I_k}\omega^{-|I|}\le e^{-I_k}$, proving
\eqref{ch11:n-mark-root}.

For the geometric bound, the scale ratio of the upper interval is at most
$L H^{b_j+2}$ and the lower ratio at most $L C^{10}$.
The lower bounds on the starting scales give $L^2C^{10}\le H^3$, so the total
cell side in continuum units is at most $a_kH^{b_j+5}$.
All covariance estimates depend only on the full displacement
between distinct points. On the compact configuration set its minimum
length is positive, also after rounding at sufficiently large
$H$. The opening-support estimate uses the same positive
maximum-coordinate distance. Neither argument divides by a
single coordinate difference. This proves the asserted
uniformity, including configurations arbitrarily close to
coordinate coincidences.
\end{proof}

Different scales in this proposition use their own normalized
lifts, but the coefficient map is identical. Consequently, if
$b_0>b_j$, the lift normalized at $b_0$ is exactly the linear
combination of the degree-$d$ lifts at $b_j$ obtained by
applying $W_{H^{b_0-b_j}}$ to its free label.
On a fixed finite-dimensional degree space that contributes a factor at most
$C H^{(b_0-b_j)d}$ in coefficient norm.
The chessboard estimate for a single insertion and its free expectation therefore
give the same factor for its chessboard seminorm.
This is the comparison needed for the size and displacement
bounds at three scales; it does not assert that one fixed
unnormalized label has a small corrected error at all three.


We can now state the finite-cutoff comparison given by these estimates.

\begin{proposition}[Comparison of the separating correlations at finite cutoff]
\label{ch11:finite-detector-input}
On the common trajectory with the stage orders, lower bounds on the block ratio and
choices of sufficiently large starting scales specified above, for each $e$ there is an injective free measurement matrix
$G:\CC^r\to\CC^s$ and the following finite-cutoff approximations.
For arbitrarily small $\eta>0$ one can choose a fixed comparison
scale $k$, bounded real local polynomial coefficient families
$u_{1,i},\ldots,u_{r,i}$ with $\|u_{j,i}\|_1\le B_k$, and finitely
many bounded families of normalized local spectator insertions, all on the selected
history, such that their point-correlation matrix $Q_i^0$
satisfies
\[
 \|Q_i^0-G\|_{\max}\le\eta
\]
eventually for every period. The correlation configurations and small
neighborhoods of them have a positive common coordinate distance
and lie inside the smallest period in continuum units. All these bounded
normalized local insertion families have uniform chessboard bounds at every
sufficiently fine fixed cell depth.
\end{proposition}
\begin{proof}
Choose the free basis and separating configurations from
Chapter~\ref{ch03:maxwell}. Construct their common
localization map and simultaneous lifts by
Section~\ref{ch11:common-extraction}, and use the restarted family
of Proposition~\ref{ch11:restart-family}.
The opening estimate, stronger support bounds, larger source weight and output normalization in
Section~\ref{ch11:corrected-pair-comparison} give corrected
subset chessboard-seminorm error $e^{-I_k}$, uniformly in cutoff and period.
Undoing the connected subtractions gives point-matrix error at
most $C_ke^{-I_k}$, where $C_k$ is a polynomial in the finitely
many free subset moments. These moments remain bounded as the
rounded, rescaled configurations approach their fixed separated
limit configurations. Thus choose one sufficiently deep $k$ to obtain the
requested $\eta$. The exact lift equations give a fixed bound
$B_k$ for all cutoffs and periods at that scale.
The product-moment estimate gives the chessboard bounds for normalized local insertions.
No convergence or choice of a limiting equation subspace is used.
\end{proof}

This is a comparison before any cutoff limit. Its common localization
map, auxiliary density and normalized lifts are the ones
constructed above.

Recall the exact cutoff-dependent equation kernel
$K_i=\mathsf Z_{N,p,e}P_{N,p,e}S_{t_N}R_e\mathcal I_e$.
Its bounded projection $\Pi_i$ is distinct from the near-identity
section comparison $P_{N,p,e}$. The latter is invertible; the
former has kernel $K_i$.

\begin{proposition}[A separating correlation family on the entire cutoff-dependent complement]
\label{ch11:uniform-detector}
Using Proposition~\ref{ch11:finite-detector-input}, for every fixed
finite difference word $l$ there are averaged separating correlation functionals $L_i$ and
a constant $Z$, independent of $N,p$, such that
\[
 \|\Pi_i v\|_1\le Z\|L_i\Pi_i v\|_\infty
 \quad(v\in V_e)
\]
eventually for every period. Their values on a finite difference
$D_lF$ of a normalized local insertion with bounded chessboard seminorm are uniformly bounded.
Their values on any trial whose chessboard seminorm tends to zero tend to
zero uniformly in period.
\end{proposition}
\begin{proof}
Choose $K$ with $\|c\|_\infty\le K\|Gc\|_\infty$ and put
$\delta=(2Kr)^{-1}$. The case $r=0$ is immediate.
Use the assumed comparison with entrywise error at most $\delta/4$.
The modulus of continuity in the insertion points for normalized local insertions from
Section~\ref{ch11:positional} holds on the separated neighborhoods
with one constant and threshold for this finite list.

Choose $b>0$ so small that moving the distinguished point by at most
$2|l|b$ stays in those neighborhoods and changes every normalized local insertion
measurement by at most $\delta/4$. Put $h_N=\lceil b/a_N\rceil$.
Average over the shifts $0,\ldots,h_N-1$ once in each coordinate
listed in $l$, always in the distinguished slot. These normalized
averages define $L_i$. For the empty word take no averages.

The vectors $(I-\Pi_i)u_{j,i}$ are bounded, hence admissible.
Lemma~\ref{ch11:uniform-sweep} and the chessboard inequality with one distinguished insertion
show that the values of the averaged functionals on them tend to zero.
Thus the matrix
\[
 Q_i=(L_i\Pi_i u_{j,i})_{j=1}^r
\]
is within $\delta$ of $G$ eventually. For any $c\in\CC^r$,
\[
 \|c\|_\infty
 \le K\|Q_ic\|_\infty+Kr\delta\|c\|_\infty,
 \qquad
 \|c\|_\infty\le2K\|Q_ic\|_\infty.
\]
The $r$ vectors $\Pi_i u_{j,i}$ are therefore independent.
They span $\operatorname{ran}\Pi_i$, since its dimension is
exactly $r$. Writing $\Pi_i v=\sum_jc_j\Pi_i u_{j,i}$ gives
\[
 \|\Pi_i v\|_1\le ArB_k\|c\|_\infty
       \le2ArB_kK\|L_i\Pi_i v\|_\infty.
\]
Take $Z\ge\max(1,2ArB_kK)$.

For the assertions on the values of the averaged functionals, inserting a distinguished trial of chessboard seminorm $u_i$ with the bounded spectator fields gives a bound $Cu_i$
uniformly throughout the neighborhoods. Averaging preserves it.
For a difference word, differences commute with these averages,
and each coordinate telescopes exactly:
\[
 \frac1{h_N}\sum_{j=0}^{h_N-1}
 \frac{H(x+(j+1)e_\alpha)-H(x+je_\alpha)}{a_N}
 =\frac{H(x+h_Ne_\alpha)-H(x)}{h_Na_N}.
\]
Iteration bounds the averaged difference by
$(2/b)^{|l|}$ times the undifferentiated correlation bound.
All endpoints stay in the chosen neighborhoods.
This proves the bound on these values without assuming any pointwise
bound on the differentiated field.
\end{proof}

\begin{theorem}[Uniform bounds for finite differences]
\label{ch11:all-period-differences}
For the common trajectory with the trial estimates of
Section~\ref{ch11:trial-descent} and the stage requirements already
specified for their uniformity, every fixed finite list of
fields, every finite difference word applied separately to its
slots, and every $r_0>0$, their finite-volume correlation is bounded
by one constant for all sufficiently large $N$, every period $p$,
and every tuple with pairwise periodic maximum-coordinate distance
at least $r_0$ in continuum units.
\end{theorem}
\begin{proof}
Start with a bounded local polynomial coefficient and a difference word $l$.
The exact physical difference identity and the translated local
section bound make this a trial at degree $e=d+|l|$ satisfying the hypotheses of \eqref{ch11:trial-error}.
The factor $a_N^{|l|}$ in the normalization at its filtration index cancels all inverse mesh
powers. Its link support uses at most $2^{|l|}$ fixed translates;
its derivative bound grows by at most
$(1+2^{\varpi_Nq_N})^{|l|}$, a subpower quantity under the stage
requirement on the starting scale. The backward comparison gives an admissible coefficient
$x_i$ and an error $E_i=O(x_i)-D_lF$ with chessboard seminorm tending to zero.

The correlation pairings just proved give $\|L_ix_i\|\le C$.
Apply Lemma~\ref{ch11:absorption} with
$y_i=\Pi_i x_i$ and
$z_i=(x_i-y_i)/(1+\|y_i\|_1)$.
The latter family is admissible: its first term is a bounded
scalar multiple of $x_i$, and its second is norm bounded.
It lies in the exact kernel of $\Pi_i$, so the values of the correlation functionals on it tend to
zero by Lemma~\ref{ch11:uniform-sweep}.
Consequently $\|y_i\|_1\le2ZC+1$.
The chessboard bound for normalized local insertions applies to $O(y_i)$,
and the unnormalized admissible kernel family $x_i-y_i$
has vanishing chessboard seminorm. The identity
\[
 D_lF=O(y_i)+O(x_i-y_i)-E_i
\]
gives a uniform chessboard bound for the difference trial.
Repeat for the finite list of orientations and slots.

For clarity, two cell scales are needed to place a separated tuple in distinct cells.
Choose a fixed coarse cell side $a_\ell<r_0$.
Choose a finer fixed depth $k\ge\ell$, beyond all the finitely many
scale thresholds for the chessboard bounds, with $6na_k<a_\ell$.
At mesh $N$ the corresponding integer cell sides $t,T$ satisfy
$a_Nt=a_k$ and $a_NT=a_\ell$.
In each coordinate one point forbids at most $6t$ translations
of the coarse tiling if a margin $3t$ is required.
The $n$ points therefore leave a permissible translation, since
$6nt<T$. Distinct points occupy distinct coarse cells: sharing
a cell would make their periodic maximum-coordinate distance
at most $a_\ell<r_0$.
The chessboard-seminorm comparison between aligned fine and coarse
tilings transfers the depth-$k$ bounds to these interior cells;
it follows from the repeated chessboard inequality (maximizing
the position of a repeated copy in its containing cell).
Translation invariance alone is not this comparison.
Placing the finitely many trials in distinct cells now bounds the correlation by
the product of their fixed chessboard bounds.
All cell sides, margins and lower bounds on $N$ were chosen independently
of $p$. The periodic difference identity also holds at wrapped
sites, completing the proof.
\end{proof}

}{}
\setcounter{chapter}{11}
\IfFileExists{chapters/12-ultraviolet-limit.tex}{% !TEX root = ../monograph.tex
\chapter{The ultraviolet limit at fixed volume}
\label{ch12:ultraviolet}

We keep the physical period $R_p$ fixed and let the ultraviolet cutoff
$N$ tend to infinity. The level $m$ has spacing $a_m=D_m^{-1}$;
blocking proceeds from $m+1$ to $m$. The purpose of this chapter is to
compare different cutoffs on one fixed configuration space. Compactness
by itself would give only a subsequence. The additional ingredient is an
estimate for stopping the marked evolution at an intermediate level,
with a constant independent of that level.

The common trajectory of Chapter~\ref{ch10:trajectory} and the
marked estimates of Chapter~\ref{ch11:insertions} supply the stability
and finite-difference estimates used here. We state their precise
forms to distinguish a change of cutoff from comparison of two
insertions under one fixed transformation.

\section{Fixed-level convergence data}
\label{ch12:inputs}

Fix finitely many labels $v_i\in V_{d_i}$ and distinct physical origins.
At a fixed level $m$ above their reference levels, index the activities
by the same local supports and prescribed sets of large-field sites for
every cutoff $N>m$. The small-field indicators and eligibility criteria
are also fixed at level $m$. A large-field profile is the sum of activities
having one particular support and prescribed set of large-field sites; the original activity
list can be infinite even though the number of such profiles on a
fixed finite torus is finite.

The required background convergence consists of Cauchy convergence of
$\beta_{N,p,m}$, of the path coefficients in their exponentially
weighted sums, and of the regular and large-field activity arrays in
the norms of Section~\ref{ch04:activities}. The scalar multiple of the
density need not converge. For the marked transformation we need
fixed-input continuity across cutoffs: at each fixed level, applying
the cutoff-dependent marked map to any fixed admissible remainder
graph gives Cauchy output graphs and coefficient maps. This condition
is stronger than contraction for two graphs acted on by the same map.
For the common construction, Theorem~\ref{ch10:convergence} supplies
the background limits. The marked estimates and fixed-input comparison
in Section~\ref{ch11:one-graph} supply the corresponding single-insertion
maps. We keep the needed forms explicit to show how the two
comparisons enter this argument.

Here is a useful way to obtain the required path convergence. For
canonical path indices $a$ of length $\ell(a)$, suppose the coefficients
$c_{N,a}$ converge individually and obey a common bound at width
$w>v$. Writing $c_a=\lim_{N\to\infty}c_{N,a}$ and $C$ for the common bound, we have
\[
 |c_{N,a}-c_a|e^{v\ell(a)}\le 2C e^{-(w-v)\ell(a)}.
\]
The canonical indices enumerate displacement vectors and finitely many
plaquette directions, rather than all walks of a given length. For a
row their length is $|x|_1+4$, and
\[
 \sum_{x\in\ZZ^4}e^{-\delta(|x|_1+4)}
 =e^{-4\delta}\left(\frac{1+e^{-\delta}}{1-e^{-\delta}}\right)^4<\infty.
\]
For an atom of degree $q$ the analogous sum has $q-1$ displacement
vectors and exponent $4(q-1)$. Dominated convergence therefore gives
convergence in the smaller width $v$. A strict width reserve makes
coordinatewise coefficient convergence sufficient for the actual
path-action norm; convergence of an unrelated finite-volume matrix
would not do so.

\section{Contraction of the complete marked remainder}
\label{ch12:graphs}

For a fixed degree, let $\mathcal R_m$ denote the space of remainder
arrays at level $m$. Its seminorm is the infimum of the sum of the
regular derivative majorants and the large-field value majorants
multiplied by $\kappa_m^{-1}$, where
$\kappa_m=(\epsilon_m^{\rm for})^{1/2}$ is the square root of the
actual level-$m$ large-field bound, as in
Section~\ref{ch11:one-graph}. We keep the arrays themselves, without
identifying different arrays merely because their evaluations agree.
Evaluation of the regular sum and of each large-field profile is a
bounded linear map in this seminorm. This follows first for any
admissible majorant and then by taking its infimum.

An admissible graph $G:V_d\to\mathcal R_m$ obeys
$\|Gx\|_{\mathcal R_m}\le s_m\|x\|_1$, where $s_m\ge1$ is fixed
by the finite-degree construction. Set
\[
 d_m(G,H)=\sup_{x\ne0}
       \frac{\|(G-H)x\|_{\mathcal R_m}}{s_m\|x\|_1}.
\]
Thus any two admissible graphs have distance at most two. A marked step
from $m+1$ to $m$ sends the input $(x,Gx)$ to $(I_Gx,J_Gx)$, where
$I_Gx$ is the new polynomial coefficient and $J_Gx$ the new remainder.
Using the new coefficient as independent variable gives the output graph
$\Phi G=J_GI_G^{-1}$.

\begin{lemma}[Graph contraction]
\label{ch12:contraction}
Suppose the marked step maps admissible graphs into admissible graphs,
and the common-support marked difference estimates give, for
$a=d_{m+1}(G,H)$,
\[
 \|I_G-I_H\|\le\delta_m a,\qquad
 \|(J_G-J_H)x\|_{\mathcal R_m}\le\delta_m a\|x\|_1,
 \qquad \|I_G^{-1}\|\le s_{m+1}.
\]
If $4\delta_m s_{m+1}\le1$, then
$d_m(\Phi G,\Phi H)\le a/2$.
\end{lemma}
\begin{proof}
For $y=I_G^{-1}x$, use $J_H=(\Phi H)I_H$ to write
\[
 (\Phi G-\Phi H)x=(J_G-J_H)y+(\Phi H)(I_H-I_G)y.
\]
Its norm is at most $(1+s_m)\delta_m a s_{m+1}\|x\|_1$,
which is at most $as_m\|x\|_1/2$.
\end{proof}
The difference estimates are for the complete regular and large-field
remainder on common supports. Linearity of a single insertion
allows its local polynomial part to cancel. A contraction only for the regular
part would not prove this lemma for the actual inserted density.
For the actual transformation, \eqref{ch11:graph-estimates} gives
ball invariance and the inverse bound, while
\eqref{ch11:graph-difference} gives these two differences with
$\delta_m=C_d\epsilon_m s_{m+1}$.
The selected ratio and thresholds for the starting scales impose
$4\delta_m s_{m+1}\le1$. The fine and coarse remainder norms retain
their distinct factors $\kappa_{m+1}^{-1}$ and $\kappa_m^{-1}$.
Weighted path continuity at a strict width reserve and activity-norm
continuity in Section~\ref{ch11:one-graph}, together with
Theorem~\ref{ch10:convergence}, give the fixed-input cutoff
continuity as well.

\begin{proposition}[Independence of the terminal graph in the cutoff limit]
\label{ch12:forgetting}
For the selected trajectory, the graph contraction and fixed-input
cutoff continuity just verified imply that the selected
graphs and coefficients normalized at one reference level are Cauchy
at every fixed level.
\end{proposition}
\begin{proof}
Fix $m$ and a positive integer $h$. Replace the selected graph at
$m+h$ by one fixed admissible graph and propagate down $h$ steps.
The resulting graphs are Cauchy as $N\to\infty$: fixed-input
continuity and contraction give this by finite induction. Their
distance from the actual graphs is at most $2^{1-h}$. First let $h$
be large and then let both cutoffs tend to infinity. This proves the
Cauchy assertion without an infinite product of inverse maps.

The coefficient at its reference level is the prescribed label.
For each subsequent fixed level solve the finite-dimensional relation
$x_{m,N}=I_{m,N}x_{m+1,N}$. The maps are Cauchy and their inverses
uniformly bounded at this fixed level. The identity
\[
 \|I_{m,N}^{-1}x_{m,N}-I_{m,M}^{-1}x_{m,M}\|_1
 \le C_m\bigl(\|x_{m,N}-x_{m,M}\|_1
       +\|I_{m,N}-I_{m,M}\|\,\|I_{m,M}^{-1}x_{m,M}\|_1\bigr)
\]
proves the assertion by finite induction. Evaluation of $G_{m,N}x_{m,N}$
is Cauchy by the graph bound and its bounded evaluation maps.
\end{proof}

Consequently the complete functions multiplying a single source $f_{i,N,m}(U)$ and
large-field profiles linear in the sources $k_{i,N,m,Q}(U)$ converge uniformly in
$U$. Here $Q$ is a support together with its exact prescribed set of large-field sites.
Multiplication by a fixed Borel event indicator preserves uniform convergence;
no continuity of the event indicator as a function of $U$ is required.

\section{The scalar normalization cancels}
\label{ch12:normalization}

Write $\rho_{N,m}$ for the actual positive density at the fixed period,
$s_{N,m}$ for its extensive scalar term, and
\[
 \widetilde\rho_{N,m}=e^{-s_{N,m}}\rho_{N,m},\qquad
 Z^0_{N,m}=\int\widetilde\rho_{N,m}\,d\mu_{m,p},
\]
where $\mu_{m,p}$ is product probability Haar measure. Background
convergence makes $\widetilde\rho_{N,m}$ uniformly Cauchy. Indeed,
curvatures are bounded on the compact group, so weighted path
convergence implies uniform action convergence. Summing a rooted
activity norm over the finitely many sites gives uniform convergence
of the regular sum and every large-field profile. The exact-cover gas
is a finite polynomial in the latter profiles and fixed eligibility
indicators. Its product with the exponential of the regular action
therefore converges uniformly.

The limit of $Z^0_{N,m}$ is strictly positive. Choose a fixed link ball
$\mathcal U_m$ about the identity small enough that every inspected
plaquette satisfies the level-$m$ cutoff. On this set there are no
large-field sites, so the exact-cover sum is one. The action and the
regular sum are bounded above by a constant $A_{m,p}$ uniformly on
the cutoff tail. Since Haar measure assigns positive mass to every
nonempty open link ball,
\begin{equation}
 Z^0_{N,m}\ge e^{-A_{m,p}}\mu_{m,p}(\mathcal U_m)>0.
 \label{ch12:positive-partition}
\end{equation}
Nonnegativity of the full density justifies restricting the integral
to $\mathcal U_m$. Positivity of each partition function alone would
not imply a positive limit.

Let $z_1,\ldots,z_n$ be commuting variables with $z_i^2=0$, and write
$z_S=\prod_{i\in S}z_i$. The exact marked density at level $m$ has
regular terms linear in the sources $\sum_i z_i f_i$, higher regular entries
$H_X$, and large-field factors
\[
 K_Q+\sum_i z_i k_{i,Q}+K_Q^{\ge2}.
\]
Every term in $H_X$ and $K_Q^{\ge2}$ contains at least two distinct
source variables. All insertion origins belong to its support. Keep the complete
terms linear in the sources and set only these higher entries to zero. Call the
resulting density $\mathscr X^0_{N,m}(z,U)$. The expression
\begin{equation}
 C_{N,p,m}(v;x)=
 \frac{\int[z_{\{1,\ldots,n\}}]
            e^{-s_{N,m}}\mathscr X^0_{N,m}(z,U)\,d\mu_{m,p}}
      {Z^0_{N,m}}
 \label{ch12:stopped}
\end{equation}
is Cauchy in $N$ for fixed $m$. Its numerator is a finite polynomial
in uniformly convergent grouped profiles and functions multiplying a single source,
multiplied by a uniformly convergent exponential. Its denominator has
the positive limit in \eqref{ch12:positive-partition}. The same scalar
has been removed from both. In particular, deleting terms involving at least two insertions does
not mean deleting the large-field profiles linear in the sources.

\section{A stopping estimate uniform in the stopping level}
\label{ch12:stopping}

Recall that $B_m$ is the support radius from \eqref{ch04:thresholds}
and $\operatorname{cov}_{B_m}(X)$ counts a connected cover of $X$ by
radius-$B_m$ boxes. Suppose the physical origins belong to one common
lattice and are distinct modulo $R_p$. Their separation at level $m$
is at least $\eta/a_m$ for some $\eta>0$. Since $a_mB_m^2\to0$,
choose a fixed base level $r$ so that separation is at least $B_m^2$
for every $m\ge r$. Enlarge $r$ to include all reference levels.

We use the following marked norm conclusion. Terms involving at least two insertions have rooted
majorant sums at most one at support exponent $\sigma_m$. The transfer
accepts exponent $\sigma_{m+1}-4$ and restores exponent $\sigma_m$.
It retains all inventories and insertion origins. Moreover, after transfer
to the fixed base level, normalized coefficient extraction is bounded
by
\begin{equation}
 C_r=\frac{\exp\bigl(|\Lambda_{r,p}|q_r^{\rm site}+n(x_r+y_r+2)\bigr)}
                 {\prod_i w_i}.
 \label{ch12:root-bound}
\end{equation}
Here $q_r^{\rm site}$ is the per-site normalization cost in \eqref{ch11:replicated-bound}; $w_i>0$ are fixed source weights; $x_r$ bounds the weighted regular sizes of the terms linear in the sources; and $y_r$
bounds their weighted large-field profile sums. Such weights can be
chosen uniformly on the cutoff tail: the fixed-level coefficients linear in the sources are bounded, so scale each weight down to its positive
allowed size. Crucially all four quantities in \eqref{ch12:root-bound}
are evaluated at $r$, not at the later stopping level $m$.
These statements concern the complete actual marked transfer.
They follow from the estimate for terms involving at least two insertions
\eqref{ch11:higher-induction} and the normalized chessboard bound
\eqref{ch11:root-bound}.

\begin{lemma}[Removal of terms involving at least two insertions]
\label{ch12:omission}
Under these marked norm and transfer bounds, for each $m>r$ and all
sufficiently large $N$,
\[
 |M_{N,p}(v;x/a_N)-C_{N,p,m}(v;x)|
       \le 2C_r e^{-B_m/64}.
\]
The constant $C_r$ is independent of $m,N$; the cutoff threshold can
depend on $m$.
\end{lemma}
\begin{proof}
A support containing two origins separated by $B_m^2$ has cover
number at least $B_m/64$. Multiply every term of source degree at least two, and only those
entries, by a complex scalar $u$. For $|u|<S_m=e^{B_m/64}$ its
majorant still fits exponent $\sigma_m-2$, since
\[
 |u|e^{(\sigma_m-2)\operatorname{cov}_{B_m}(X)}
       \le e^{\sigma_m\operatorname{cov}_{B_m}(X)}.
\]
Transport the entire modified marked density to $r$. The support
reserve allows its first and subsequent steps, and the constant and
coefficients linear in the sources are unchanged. If a source monomial contains $q$ origins,
the stronger geometric bound $B_mq\le64\operatorname{cov}_{B_m}(X)$
also absorbs a change of each source weight by at most $e^{B_m/128}$.
This is the weight reserve needed when using the marked transfer.

Let $F_{N,m}(u)$ be the normalized integral of its full-source
coefficient. Exact probability-kernel integration preserves that
integral and its vacuum denominator. The fixed-base estimate gives
$|F_{N,m}(u)|\le C_r$ throughout the disk. The source algebra is
nilpotent; after grouping the finitely many support and prescribed-site profiles,
$F_{N,m}$ is a polynomial in $u$, hence holomorphic. For $S_m>2$,
Cauchy's derivative estimate on circles of radius $S_m/2$ about points
of $[0,1]$ gives
$|F_{N,m}(1)-F_{N,m}(0)|\le2C_r/S_m$. These two values are the
actual cutoff moment and \eqref{ch12:stopped}, respectively.
\end{proof}

\begin{proposition}[Cauchy convergence on common lattice points]
\label{ch12:dense-cauchy}
For the selected trajectory, the correlations converge along the whole cutoff sequence at every
configuration of distinct points belonging to one common lattice.
\end{proposition}
\begin{proof}
Given an error tolerance, first choose $m$ so that the two stopping
errors in Lemma~\ref{ch12:omission} are small. Keep $m$ fixed. The
stopped expression \eqref{ch12:stopped} is Cauchy as both cutoffs tend
to infinity. The triangle inequality proves the assertion. No new
subsequence is chosen.
\end{proof}

\section{Smoothness and all lattice differences}
\label{ch12:smooth}

The remaining passage is a standard discrete version of the
Arzel\`a--Ascoli theorem and convergence of derivatives; compare
\cite[Theorems~7.17 and~7.25]{Rudin1976}. We record the precise
application, since the rounded lattice functions are not continuous.
Let $U\subset\RR^D$ be open, $a_N\downarrow0$, and let $F_N$ be
functions on $a_N\ZZ^D$. Write
\[
 \Delta_{N,j}F_N(x)=a_N^{-1}(F_N(x+a_Ne_j)-F_N(x)).
\]
At a point $x\in U$, a rounded value means evaluation at a lattice point
within one mesh step in each coordinate, chosen by a fixed rule. Assume every fixed iterated difference is bounded on each compact
subset of $U$, eventually in $N$, and the rounded values are Cauchy
on a dense subset. Then there is $F\in C^\infty(U)$ such that all
rounded iterated differences converge locally uniformly to the
corresponding derivatives of $F$.

Only the discrete-to-continuous adaptation needs explanation. For
nearby lattice points, telescoping along coordinate edges gives
$|G_N(y)-G_N(x)|\le C|y-x|_1$ whenever the first differences of $G_N$
are bounded on a slightly larger compact neighborhood. Rounding adds
at most $2DCa_N$ to this modulus. A finite net from the dense Cauchy
set therefore proves locally uniform convergence of $F_N$.
Bounded second differences give the discrete Taylor estimate
\[
 |G_N(y)-G_N(x)-\sum_j(y_j-x_j)\Delta_{N,j}G_N(x)|
       \le C|y-x|_1^2.
\]
For fixed small $h>0$, compare $\Delta_{N,j}G_N(x_N)$ to the secant
with interval $h_N=a_N\lceil h/a_N\rceil$. The error is at most
$Ch_N$. The secants converge first as $N\to\infty$; then $h\to0$
proves that the first differences are Cauchy. Their own first-difference
bounds give locally uniform convergence. Passing to the displayed
Taylor inequality identifies their limits as the derivatives of the
limit of $G_N$. Iteration proves the assertion for every finite order.
All paths and shifted stencils lie in the enlarged compact neighborhood.

Apply this to periodic lifts on
\[
 U_{p,n}=\{(x_1,\ldots,x_n)\in(\RR^4)^n:
                         x_i-x_j\notin R_p\ZZ^4\text{ for }i\ne j\}.
\]
The union of the nested physical lattices is dense here. Its points
supply Proposition~\ref{ch12:dense-cauchy}; the all-order separated
finite-difference bounds supply the other hypothesis. Points sharing
one or several coordinate projections are allowed. Only coincidences
of the full four-dimensional points modulo the period are excluded.
The limiting function, denoted $M^{\rm UV}_{p,v}$, is smooth and
periodic on this domain.
The difference hypothesis is supplied by
Theorem~\ref{ch11:all-period-differences}: on a compact subset
of $U_{p,n}$ the periodic separation has a positive lower bound,
and any fixed shifted difference stencil stays in a slightly
larger separated compact set on the cutoff tail.

\section{The fixed-volume conclusion}
\label{ch12:conclusion}

\begin{theorem}[Full ultraviolet limit]
\label{ch12:uv-theorem}
For the common selected trajectory with the strengthened stage
requirements of Chapters~\ref{ch10:trajectory} and
\ref{ch11:insertions}, the correlations have the smooth
full-sequence limit
$M^{\rm UV}_{p,v}$ at every fixed physical period. For every
$f\in C_c^\infty(U_{p,n})$ supported in an admissible period cell
and every finite coordinate list $\alpha$,
\[
 \lim_{N\to\infty}a_N^{4n}\sum_{y}
    f(a_Ny)\Delta_{N,\alpha}M_{N,p}(v;y)
       =\int f(x)\partial^\alpha M^{\rm UV}_{p,v}(x)\,dx.
\]
Here differences in the summand are divided by their physical step
lengths. Their shifted sites must lie in the slightly enlarged
separated set, which holds on a cutoff tail.
\end{theorem}
\begin{proof}
The preceding section gives the full smooth limit on separated configurations.
The last assertion is a Riemann-sum consequence of locally uniform
convergence. Fixed forward, backward or centered stencils give the
same limit: telescope each stencil as a finite average of shifted
unit forward differences. The averaging weights sum to one and the
shifted points approach the same point.
\end{proof}

The convergence threshold in this theorem may depend on $p$. Bounds
on compact separated sets that were uniform in $p$ remain uniform
for the limiting kernels, by passage to the limit. These two
quantifier statements serve different purposes: the former removes
the ultraviolet cutoff; the latter permits the subsequent
infinite-volume extraction. Finite reflection inequalities also pass
to the limit whenever every correlation in their finite quadratic
form is supported away from collisions. Neither this argument nor
reflection positivity asserts an unsubtracted limit at coincident
points. Distributional extensions and reconstruction require their
own arguments.
}{}
\setcounter{chapter}{12}
\IfFileExists{chapters/13-continuum-fields.tex}{\chapter{Continuum correlations and local fields}
\label{ch13:continuum}

The ultraviolet limit leaves a theory on each fixed physical torus.
We now let its period grow. The construction may select a subsequence
of volumes, but it does not select a new subsequence of ultraviolet
cutoffs for the correlations. Local fields are then defined by their
responses in all separated correlations. Polynomial labels that give
the same responses represent the same field.

We use the full fixed-volume limit of
Theorem~\ref{ch12:uv-theorem}, the volume-uniform finite-difference
bounds of Theorem~\ref{ch11:all-period-differences}, and the equation comparisons and the comparisons using separating families in Chapter~\ref{ch11:insertions}.
Compactness gives the volume extraction; the separate equation
inclusion and rank bounds identify its local fields.

\section{One subsequence of volumes}
\label{ch13:extraction}

Recall that $\Conf_n$ is the open set of pairwise distinct points in
$(\RR^4)^n$, without any restriction on their coordinate projections.
Let $p(k)$ be any strictly increasing sequence of volume indices.
On every compact subset of $\Conf_n$, the periodic kernels
$M^{\rm UV}_{p(k),v}$ are eventually defined with separation bounded
below independently of $k$. Each of their derivatives is uniformly
bounded there. The cutoff threshold used to prove the bound may
depend on the fixed volume; after taking that ultraviolet limit it
disappears from the bound.

Arzel\`a--Ascoli and diagonal extraction
\cite[Theorem~7.25]{Rudin1976} give one refinement $P(k)=p(q(k))$
and a multilinear family $S_n$ such that
\begin{equation}
 M^{\rm UV}_{P(k),v}\longrightarrow S_n(v;\cdot)
                    \quad\hbox{in }C^\infty_{\rm loc}(\Conf_n).
 \label{ch13:volume-limit}
\end{equation}
The diagonal family is countable: arities, degree lists, basis-label
lists, derivative lists, and compact exhaustions of the configuration
spaces. Multilinearity extends its conclusions to all labels.
Integration by parts identifies the extracted derivative limits with
the derivatives of $S_n$. No diagonalization over all smooth test functions is
needed. For each fixed separated test function $f$ the order of limits is
\begin{equation}
 \lim_{k\to\infty}\lim_{N\to\infty}
 a_N^{4n}\sum_y f(a_Ny)M_{N,P(k)}(v;y)
                   =\int f(x)S_n(v;x)\,dx.
 \label{ch13:ordered-limit}
\end{equation}
The inner limit is the whole cutoff sequence. Every further volume
refinement preserves \eqref{ch13:volume-limit} and
\eqref{ch13:ordered-limit}.

Permutation symmetry, conjugation, the distinguished unit and
hypercubic covariance pass from the finite correlations by these
limits. Translation by an arbitrary real vector is obtained by
approximating it on the same lattice; the uniform first-difference
bound controls the rounding error. Full orthogonal covariance will
require a separate argument.

\section{Separated cells and errors in local observables}
\label{ch13:cells}

The chessboard seminorm of Section~\ref{ch01:reflection} bounds an
observable by its repeated reflected copies in a tiling. For local
observables $F_1,\ldots,F_n$ interior to distinct cells of the same
translated tiling $\mathcal T$, the chessboard inequality gives
\begin{equation}
 \left|\left\langle\prod_iF_i\right\rangle_{N,p}\right|
                   \le\prod_i\|F_i\|_{\mathcal T}.
 \label{ch13:chess}
\end{equation}
It is the lattice inequality already cited from
\cite{FrohlichIsraelLiebSimon1978}, with the constant observable in every unused cell.
The estimate does not grow with the distances between the cells.

Here is the required geometric placement. For a tiling of physical
side $a$, the translations in one coordinate bringing one of $n$
anchors within distance $w$ of a boundary have total length at most
$2nw$ modulo $a$. If $2nw<a$, a translation avoids all these strips;
a lattice translation does so for sufficiently fine cutoff. Choose
one in each coordinate. Anchors with pairwise maximum-coordinate
distance greater than $a$ cannot share a cell. Hence if the centered
local bounds apply at scale $k$ with margin $3a_k$, and
$6na_k<a_m$, anchors separated by more than $a_m$ admit distinct
side-$a_m$ cells with the required margins.

We shall use an error family $E_{N,p}$ with the following property:
for any fixed arity and sufficiently small fixed cell side, its
chessboard seminorm is at most $\eta_{N,p}\to0$ in the ordered
limits, uniformly over placements satisfying a specified interior
margin. The cell side can be chosen after the arity and a requested
lower scale, before the cutoff. The spectator insertions have
uniform centered seminorm bounds on sufficiently fine fixed scales.
Place the error and all spectators in distinct cells, choosing a
still finer scale for the spectators' margins. Equation
\eqref{ch13:chess} yields
\begin{equation}
 \left|\left\langle E_{N,p}(y_i)
                  \prod_{h\ne i}O_{N,p}(v_h,y_h)\right\rangle\right|
               \le C(v,\delta)\eta_{N,p},
 \label{ch13:vanishing-error}
\end{equation}
for configurations separated by $\delta>0$. The weighted number of
lattice points in the fixed compact support of a test function is bounded. The same
error therefore vanishes after separated smearing. This conclusion
uses the stated order of choices for the error norm; pointwise
smallness at one configuration would not suffice.

\section{Fields as equivalence classes of labels}
\label{ch13:quotient}

A label includes its degree space: the same formal polynomial in
$V_d$ and $V_e$ can have different cutoff normalizations. For
$c\in V_d$, a fixed spectator list $v$, and a separated smooth test function
$f$, define the one-slot response
\[
 L_{v,f}(c)=\int f(x_0,x)S_{n+1}(c(x_0),v(x))\,dx_0\,dx.
\]
Let
\begin{equation}
 N_d(S)=\bigcap_{v,f}\ker L_{v,f},\qquad F_d=V_d/N_d(S).
 \label{ch13:null}
\end{equation}
The intersection permits every finite spectator list, of arbitrary
degrees. Smoothness makes this equivalent to vanishing in every
pointwise separated correlation. Permutation symmetry ensures that
nullity has the same meaning in every slot. The quotient $F_d$ is
finite dimensional, but its dimension has not yet been identified.

We use the bounded equivariant idempotents $\Pi_{d,N,p}$ of
Section~\ref{ch11:uniform-differences}, with kernel
$\mathsf Z_{N,p,d}P_{N,p,d}S_{t_N}R_d\mathcal I_d$.
Choose a common ordered limit $\Pi_d$ after a coefficient refinement. More
precisely, diagonal extraction over the countably many degrees and
volumes gives one increasing sequence $N=n(j)$ along which the bounded
coefficient arrays converge at each fixed volume. A further volume
refinement makes those limits converge as well. All coefficient
limits below use these refined ordered tails. This auxiliary selection
does not change the already existing full cutoff limits of the
correlations. Write
$r_d=\dim\mathcal M_d^{\rm Max}$ for the dimension of the free Maxwell
quotient in Chapter~\ref{ch03:maxwell}. The construction gives:
\begin{enumerate}
\item the ranks of $\Pi_{d,N,p}$ are at most $r_d$;
\item every bounded family in $\ker\Pi_{d,N,p}$ has vanishing
separated responses in the ordered limits;
\item there are finite matrices of limiting separated responses,
evaluated on comparison labels, that approach an invertible
$r_d\times r_d$ free matrix of correlation pairings.
\end{enumerate}
The first statement is exact rank equality, since every map defining
the finite equation space is invertible. For the second,
Lemma~\ref{ch11:uniform-sweep} makes its chessboard seminorms tend
to zero; bounded rows have uniformly subpower inverse terminal
coordinates, and \eqref{ch13:vanishing-error} supplies every spectator
response. For the third, Proposition~\ref{ch11:finite-detector-input}
gives the actual finite-cutoff matrices with error tending to zero
as their comparison depth increases. Include their finitely many
bounded label columns at each of the countably many comparison depths
in the same coefficient refinement. At each fixed depth, their
ordered limits preserve the matrix estimate by multilinearity and
the uniform correlation bounds. An injective rectangular matrix of correlation pairings
has an invertible square minor; retain those rows.

\begin{proposition}[Identification by two rank bounds]
\label{ch13:rank}
For this selected continuum family,
$N_d(S)=\ker\Pi_d$ and $\dim F_d=r_d$.
\end{proposition}
\begin{proof}
For fixed $c$, the bounded family $(1-\Pi_{d,N,p})c$ has vanishing
responses. Taking limits gives $L_{v,f}((1-\Pi_d)c)=0$ for all
$v,f$, so $\ker\Pi_d\subset N_d(S)$. Converging idempotents have
eventually constant ranks: if their norm difference is less than
one, each is injective on the other's range. Thus
$\dim F_d\le\operatorname{rank}\Pi_d\le r_d$.
Each comparison matrix factors through $F_d$, since its rows are
separated responses. Matrices sufficiently close to the invertible
free matrix of correlation pairings have rank $r_d$. Hence $\dim F_d\ge r_d$.
Equality of dimensions and the kernel inclusion prove the result.
\end{proof}

\subsection{Representatives chosen before the spectators}

Approximate inclusion and derivative labels need not stay bounded.
The following finite-dimensional observation explains the normalized
null estimate required to replace them by fixed vectors. It also
avoids choosing a separate representative for every test function.

Let $V$ be finite dimensional, $\Pi_u\to\Pi$ idempotents, and
$L_{i,u}\to L_i$ linear functionals, along a proper filter of indices.
Assume $\ker\Pi=\bigcap_i\ker L_i$. If vectors $x_u$ satisfy
\begin{equation}
 L_{i,u}(x_u)\to t_i,\qquad
 L_{i,u}\left(\frac{x_u-\Pi_ux_u}{1+\|\Pi_ux_u\|}\right)\to0
                              \quad\hbox{for every }i,
 \label{ch13:normalized-null}
\end{equation}
there is one $z\in\operatorname{ran}\Pi$ with $L_i(z)=t_i$ for
all $i$. To see this, put $b_u=(1+\|\Pi_ux_u\|)^{-1}$ and
$z_u=b_u\Pi_ux_u$. A common cluster point $(b,z)$ satisfies
$b+\|z\|=1$, $\Pi z=z$ and $L_i(z)=bt_i$. If $b=0$, then
$z\in\ker\Pi\cap\operatorname{ran}\Pi$, a contradiction.
The vector $z/b$ is the desired representative. Compactness is used
only in one fixed finite-dimensional space; all response identities
hold at the same cluster point because they hold along the filter.

In the application, the filter is generated by ordered tails:
$p=P(k)$ with $k\ge k_0$, and $N=n(j)$ with $j\ge j_0(k)$.
The normalized null condition is available for every admissible
coefficient family in the sense of Section~\ref{ch11:uniform-differences}.
Indeed, if $x_{N,p}$ has subpower inverse terminal coordinates, then
$x/(1+\|\Pi x\|)$ has the same property, whereas
$\Pi x/(1+\|\Pi x\|)$ is bounded and therefore admissible.
Their difference lies in the exact finite equation kernel.
Lemma~\ref{ch11:uniform-sweep}, followed by
\eqref{ch13:vanishing-error}, proves
\eqref{ch13:normalized-null} for all spectators on the same tail.
This argument applies even when the original $x_{N,p}$ are unbounded.

\section{Inclusions and derivatives}
\label{ch13:calculus}

For each basis label $c\in V_d$, the protected trial construction
of Section~\ref{ch11:trial-descent} gives approximate representatives
in $V_e$ for $e\ge d$ and in $V_{d+1}$ for a physical lattice
difference. To check the inclusion trial, multiply the original
cutoff observable by $a_N^e$; its local representative is
$a_N^{e-d}$ times the degree-$d$ section, so its derivative and
value bounds remain subpower. For a difference, the extra factor
$a_N$ cancels the inverse step, and the fixed translated-stencil
cost is bounded by the coupling-threshold requirements used in
Theorem~\ref{ch11:all-period-differences}.
Trial descent yields admissible coefficients and vanishing chessboard seminorm errors, uniformly in period. The preceding paragraph gives
their normalized null estimates.
All representatives and comparison errors are chosen before
the spectators and test functions.
Applying the preceding
finite-dimensional observation gives fixed linear maps on labels,
$\widehat\iota_{d,e}$ and $\widehat\partial_{d,\mu}$, such that
\begin{align}
 S(\widehat\iota_{d,e}c(x_0),v(x))&=S(c(x_0),v(x)),\label{ch13:inclusion-response}\\
 S(\widehat\partial_{d,\mu}c(x_0),v(x))
                      &=\partial_{x_0^\mu}S(c(x_0),v(x)).
 \label{ch13:derivative-response}
\end{align}
For the derivative target, exact torus summation by parts is
\[
 a_N^{4n}\sum_y f(a_Ny)\Delta_{N,i,\mu}M_{N,p}(y)
 =a_N^{4n}\sum_y
 \frac{f(a_Ny-a_Ne_{i,\mu})-f(a_Ny)}{a_N}M_{N,p}(y).
\]
The quotient on the right tends uniformly to $-\partial_{i,\mu}f$;
the zero-order separated bound controls its error. This identifies
the distributional derivative target. Smoothness gives
\eqref{ch13:derivative-response} pointwise.

These maps descend to unique maps
$\iota_{d,e}:F_d\to F_e$ and $\partial_\mu:F_d\to F_{d+1}$
with the displayed responses. They preserve null classes by their
definition. The inclusions are injective and satisfy
\[
 \iota_{d,d}=1,\qquad
 \iota_{e,f}\iota_{d,e}=\iota_{d,f},\qquad
 \partial_\mu\partial_\nu=\partial_\nu\partial_\mu.
\]
Each identity follows by testing against every spectator and using
\eqref{ch13:null}; the derivative maps also commute with inclusions.
Conjugation and hypercubic covariance descend in the same way.
The union of the $F_d$ is therefore a filtered space of local fields,
with derivatives raising the filtration by one. No multiplication of
two fields at the same point has been defined by this construction.

\section{Reflection positivity and separated power bounds}
\label{ch13:positivity}

Let a word mean a finite list of labels and positions, with correlation
denoted $S(w)$. Write $w^*$ for conjugation of its labels and reversal
of its order. Time is coordinate four. Let $\theta w$ reflect that
coordinate and act on its labels by the corresponding signed
transformation, as well as applying $*$. For words with strictly
positive times, lattice bond reflection positivity passes to
\[
 \sum_{i,j}\overline{c_i}c_j S(\theta w_i\,w_j)\ge0.
\]
It also passes for joint smooth test functions supported on positive-time
separated configurations. Indeed their lattice supports stay away
from both reflection planes on the ordered tail. The lattice plane
$x^4=a_N/2$ tends to zero, and the uniform separated bound controls
the moving test functions. Taking the two limits preserves positivity of each
finite matrix. Pointwise positivity follows by concentrating smooth
bumps. This uses the finite lattice result of
Section~\ref{ch01:reflection}; it asserts no positivity for unspecified
extensions at coincident points.

For later use take the centered chessboard estimate
$\|O\|_k\le C a_k^{-(d+1)}$, uniform on a bounded set of labels
of degrees at most $d$ and every sufficiently fine fixed level $k$.
It follows by choosing $\epsilon=1$ in
Proposition~\ref{ch13:sharp-centered}, whose proof below uses only
the bound for reflected tilings of Chapter~\ref{ch11:insertions} and the
selected schedule. Slow growth of the block
ratios gives a constant $c>0$ with
\begin{equation}
                         a_{k+1}\ge c a_k^2.
 \label{ch13:scale-thickness}
\end{equation}
The first scale below a length $r\le1$ is consequently at least a
constant times $r^2$. Choose a cell side below $r/2$, then a centered
scale below that side divided by $6(n+1)$. The latter is at least
$c_n r^4$. The separated-cell estimate yields
\begin{equation}
 |S_n(v;x)|\le C_{n,v}\,r^{-4(d+1)n},
 \qquad \min_{i\ne j}|x_i-x_j|_\infty\ge r.
 \label{ch13:power}
\end{equation}
The constant is independent of the diameter of the configuration.
Applying the same argument to the fixed degree-raised labels in
\eqref{ch13:derivative-response} gives a polynomial separated bound
for every derivative.

\section{Repeated reflected words}
\label{ch13:repeated}

The preceding bound cannot be used with arity tending to infinity
without controlling its constants. We need a more specific estimate.
Let $w$ have $n\ge1$ insertions, degrees at most $d$ and coefficient
norms at most $R$,
pairwise separation at least $r$, and times in $[r,h-r]$, where
$0<r\le1$. We call $r$ its buffer in the slab of width $h$.
For a positive-time exterior word $v$, set
\[
 A_wv=w\,\tau_{se_4}v,\qquad w^\dagger=\tau_{se_4}\theta w,
 \qquad \|v\|_{\rm OS}^2=S(\theta v\,v).
\]
At this point $A_w$ acts only on words; boundedness on a Hilbert
space is not assumed.

\begin{proposition}[Uniform bound for repetitions]
\label{ch13:word-bound}
Under the centered chessboard bounds and their interior-cell
comparison used above, choose $0<s\le h$ so that $w$ still has
buffer $r/2$. There are $C=C(d,R)$ and $\kappa=4(d+1)$ such that,
for every exterior word $v$,
\[
 \|(A_{w^\dagger}A_w)^jv\|_{\rm OS}^2
       \le B_v\,[C^n(n/r)^{\kappa n}]^{4j},\qquad j\ge0.
\]
Here $s,C,\kappa$ are chosen before $v$, and $B_v<\infty$ is
independent of $j$.
\end{proposition}
\begin{proof}
Choose a cell scale $a_m<r/32$ as the first permitted such scale,
and $s=a_m\lfloor h/a_m\rfloor$. Equation
\eqref{ch13:scale-thickness} gives $a_m\ge c_1r^2$.
Choose a centered scale $a_k<\alpha(r/n)^2$, with $20\alpha<c_1$.
Again taking the first such scale gives
$20na_k<a_m$ and $a_k\ge c_2(r/n)^4$.

The expanded reflection form contains $4jn$ main insertions and two
exterior copies. Since $s$ is a multiple of $a_m$, the main points
have at most $2n$ residue patterns modulo the side-$a_m$ tiling,
independently of $j$. The strip count therefore places all main
points in distinct cells with margins $5a_k$, using constants
independent of the number of repetitions.

Choose finer aligned cell and centered scales for the finitely many
exterior points, with exterior cell side at most $a_k$.
Their two copies have times $\pm(2js+t)$, so a
choice depending on their minimum positive time and internal
separation works for every $j$. A further translation smaller than
the exterior cell side adjusts their margins while leaving the main
margins at least $3a_k$. The odd integer ratios preserve the fine
cell grid under coarse reflections. Apply the chessboard inequality
first to the coarse cells, regarding all exterior factors in one
coarse cell as a single observable, then to the fine cells for those
factors. If $X$ bounds the exterior centered seminorms, the result is
\[
 X^{2|v|}(C_0a_k^{-(d+1)})^{4jn}.
\]
Take $B_v=\max(1,X^{2|v|})$ and use the lower bound on $a_k$.
Every use of the lattice inequality is at fixed $j$ before taking
the ordered limits. A cutoff threshold uniform in $j$ is unnecessary.
\end{proof}

This estimate is the input for bounded slab insertions in the
reflection Hilbert space. Its derivation precedes that operator
construction and does not use orthogonal covariance or a phase
decomposition. The equation comparison, admissible trial encoding,
estimates for finite separating families and centered estimates all refer to the same selected
family. Full orthogonal covariance and products at coincident points
remain separate constructions.

% !TEX root = ../monograph.tex
% Fragment for Chapter13: insert after the coarse separated bound.
\section{Sharp short-distance bounds for a fixed number of insertions}
\label{ch13:sharp-section}

For a fixed number of insertions the loss in
\eqref{ch13:power} can be reduced to an arbitrarily small excess over
the sum of their degrees. This uses the normalized integral over the reflected tiling,
not the stronger Gaussian comparison used for the separating families. We keep
all incoming joint components with insertions; replacing them by a product of
terms of source degree one would invalidate the estimate.

Write $\tau_m=\log(1/a_m)$, so
$\widehat\beta_m=\bar\beta+\kappa_G\tau_m$. Let $C_j\ge1$ include
the inverse-column, finite-section and source constants at stage $j$.
The stage ratio and the lower bound on the reference coupling may be chosen to satisfy
\begin{equation}
 \frac{\log C_j}{\log \mathscr L_j}\longrightarrow0,\qquad
 B_{i+1}^2\le2B_i^2,\qquad
 \frac{(L_iB_i^2)^4 q_{i+1}^{\rm site}}{\tau_{i+1}}
                          \longrightarrow0.
 \label{ch13:sharp-schedule}
\end{equation}
Here $B_i$ is the support-block radius, and $q_i^{\rm site}\ge1$ is
the per-site normalization exponent of \eqref{ch11:replicated-bound};
it is not a density with insertions or the coefficient of the coupling constant renormalization.
For the path-action representation of the density one can use
$q_i^{\rm site}\le C_{J(i)}(1+\log\widehat\beta_i)$.
Indeed the exponent without its constant term is at most $C_{J(i)}V$. Set
$\rho=\widehat\beta_i^{-1/2}$ and restrict every link to
$\|U_e-1\|_G\le\rho/4$, using the matrix norm \eqref{ch01:matrix-norm}.
On this event the exponent is at least $-C_{J(i)}V$ and all large-field
indicators vanish. Exponential summability in the path length bounds the
leading positive path quadratic by $C_{J(i)}\widehat\beta_i\rho^2V$.
Signed regular costs and the absolute
activity site sum satisfy the bounds \eqref{ch04:activity-bounds}. The lower bound for the partition function
of Section~\ref{ch11:replicated} therefore adds only
$C_{J(i)}V+2D_GV\log\widehat\beta_i$ to the exponent, where $D_G$ is an ambient
real matrix dimension for the gauge group.

These requirements are compatible with sufficiently large lower bounds on the reference coupling. At each fixed
stage $B_i=O_j((\log(2+\tau_i))^2)$ and
$q_i^{\rm site}=O_j(1+\log(2+\tau_i))$, so the last quotient in
\eqref{ch13:sharp-schedule} tends to zero at its fixed ratio. Choose
the ratio after $C_j$, then increase the lower bound on the reference coupling until that quotient is at
most $1/j$ for both adjacent-stage choices. The same choice gives
the displayed bound on consecutive $B_i$,
$\log L_i=o(\tau_i)$, and $\log C_{J(i+1)}=o(\tau_i)$
at transitions. Choose each ratio against both adjacent-stage column
constants, which are fixed beforehand. These are conditions on the single selected
schedule; no new schedule is chosen after fixing a correlation.

\subsection{A normalized local insertion}

Fix a scaling-dimension bound $d$ and a ball of radius $R$ in the
coefficient space of $V_d$. Set
\[
 E_k=C_{d,R}^{\rm sec} C_{J(k)}\prod_{i<k}C_{J(i+1)},\qquad
 \omega_k=a_k^d\widehat\beta_k^{-w_d}E_k^{-1}.
\]
The fixed section constant $C_{d,R}^{\rm sec}$ includes the coefficient radius and all its signed and
conjugated images. Choose the fixed integer $w_d$ large enough to absorb
the finite powers of the inverse small-field threshold in the degree-$d$
section and its counterpart with a prescribed large-field set, as well as the source attachment bound. Multiplication of the single-insertion inverse bounds gives the
factor $E_k$; the scaling normalization is exactly $a_k^{-d}$.
Thus multiplication of the \emph{full} insertion by $\omega_k$
makes the exponential bound for its regular term of source degree one
less than three and the bound for its large-field term of source degree one
less than one. Its remaining normalized size can be placed below the
attachment bound by increasing $w_d$.

The bounds in Chapter~\ref{ch11:insertions} for terms of source degree at least two retain every
incoming source term and every original eligibility event. Decreasing all source
weights only decreases their sums of absolute coefficients. The source attachment
inequalities at the intervening scales propagate this chosen small
weight; they require the same inverse-column and threshold powers already
included in $E_k,w_d$. Consequently \eqref{ch11:root-bound}, with
$x=3,y=1$, gives
\begin{equation}
 \|O_{N,p}^{(d)}(c)\|_{\mathrm{cell}(m)}
 \le \omega_k^{-1}
           \exp\bigl((a_m/a_k)^4q_k^{\rm site}+6\bigr).
 \label{ch13:sharp-root}
\end{equation}
It holds for all sufficiently large $N$ uniformly in period and in the coefficient
ball, whenever the reflected insertion points at the stopping scale have spacing
at least $B_k^2$. Their spacing stays adequate at finer scales because
$L_i\ge3$ whereas $B_{i+1}^2/B_i^2\le2$.

The logarithm of $E_k$ is $o(\tau_k)$: split its sum into a finite
initial segment and a tail where
$\log C_{J(i+1)}\le\epsilon\log L_i$, then let $\epsilon\downarrow0$.
The separate terminal column is absorbed by the same ratio bounds and
$\log L_{k-1}=o(\tau_k)$. Also
$\log\widehat\beta_k=o(\tau_k)$. Therefore
\begin{equation}
                        6-\log\omega_k=d\tau_k+o(\tau_k).
 \label{ch13:weight-size}
\end{equation}
This is the reason the prefactor from the reflected tiling does not introduce a fixed
extra scaling degree.

\subsection{The stopping scale and centered estimate}

Choose $k=k(m)$ as the first scale $k\ge m$ with
$B_k^2\le a_m/a_k$. It exists because $B_k$ grows subexponentially
in $\tau_k$, whereas $a_m/a_k=e^{\tau_k-\tau_m}$. By minimality,
\[
 (a_m/a_k)^4q_k^{\rm site}
       \le (L_{k-1}B_{k-1}^2)^4q_k^{\rm site}=o(\tau_k).
\]
If the first scale already qualifies the same conclusion follows from
the fixed-stage bound. Since $q_k^{\rm site}\ge1$,
\[
 0\le\tau_k-\tau_m=\log(a_m/a_k)
       \le(a_m/a_k)^4q_k^{\rm site}=o(\tau_k).
\]
It follows that $\tau_k/\tau_m\to1$.
Combining this with \eqref{ch13:sharp-root} and
\eqref{ch13:weight-size} proves the following statement.

\begin{proposition}[Sharp centered bound]
\label{ch13:sharp-centered}
For every scaling-dimension bound $d$, coefficient radius $R$ and $\epsilon>0$, there
is $C_{d,R,\epsilon}$, fixed before the cell scale $m$, such that
\[
 \|O_{N,p}^{(d)}(c,y)\|_{\mathrm{cell}(m,y)}
        \le C_{d,R,\epsilon}a_m^{-d-\epsilon},\qquad\|c\|\le R.
\]
For each fixed $m$ this holds for all sufficiently large cutoffs,
uniformly in $p,y$ and in signed and conjugated images. The insertion
is interior in that cell on the same tail.
\end{proposition}
\begin{proof}
For all sufficiently large $m$, the logarithm of the right side of
\eqref{ch13:sharp-root} is at most $(d+\epsilon)\tau_m$ by the
preceding limits. For the finite remaining set of $m$, use its chosen
fixed stopping scale and take the maximum of the finite constants.
The fixed link support becomes interior at every fixed cell
scale in continuum units as $N\to\infty$. Translation covariance makes the lower bound on $N$
independent of $y$.
\end{proof}

\subsection{Arbitrary separated points}

The slow ratio condition gives a sharper estimate of the first scale below a length
than \eqref{ch13:scale-thickness}. For any fixed $c>0$, let $m(r)$
be the first scale with $a_m<cr$. Then
\begin{equation}
            \frac{\log(1/a_{m(r)})}{\log(1/r)}\longrightarrow1
                         \quad(r\downarrow0).
 \label{ch13:first-below}
\end{equation}
Indeed $a_{m-1}\ge cr>a_m=a_{m-1}/L_{m-1}$, and the logarithm of
the overshoot is at most $\log L_{m-1}=o(\tau_{m-1})$.
Any fixed finite iteration of such choices of a first scale has the same property.

For a fixed number $n$ of insertions and mutually separated insertion points, first choose a
cell side below a fixed small multiple of their minimum separation $r$.
Translate its boundaries away from all insertion points. Then choose a centered
scale smaller than this side divided by $6(n+1)$.
The interior-cell comparison and the chessboard inequality apply.
Both chosen scales have logarithm $(1+o(1))\log(1/r)$ by
\eqref{ch13:first-below}. Apply Proposition~\ref{ch13:sharp-centered}
with a sufficiently small excess to each degree separately. We obtain,
for every $\epsilon>0$,
\begin{equation}
 |S_n(A_1,\ldots,A_n;x)|
    \le C_{\boldsymbol A,n,\epsilon}
             r^{-\sum_i d_i-\epsilon},\qquad 0<r\le1,
 \quad A_i\in F_{d_i},\quad\min_{i\ne j}|x_i-x_j|_\infty\ge r.
 \label{ch13:sharp-separated}
\end{equation}
The same bound is uniform over bounded sets of fixed raw representatives.
Its constant is independent of the diameter of the configuration.
For fixed derivatives, replace the differentiated field by its degree-raised
label; their total order is added to $\sum_i d_i$.


% !TEX root = ../monograph.tex
\section{Bounded real approximations for separated words}
\label{ch13:bounded-approximations}

The later phase argument needs real approximants bounded uniformly in
the cutoff. Compactness of the link space at each cutoff is insufficient:
the physical normalization of a local field grows as the mesh shrinks.
We instead stop at one fixed coarse level and pull back a bounded coarse
function through the probability block kernels.

We record the precise marked estimate used in this construction. It is
the local-trial comparison for reflected tilings, with every term involving at least two insertions
retained, not only its one-column coefficient estimate. At a coarse
level $h$, let $J_h$ map the coefficients of a trial in the chosen basis
to the coefficients of its selected level-$h$ polynomial. Write $C_h$
for the prescribed coefficient-column bound,
$Q_h$ for the ordinary norm bound of a normalized local trial, and
$E_h$ for its complete bound, including the complementary large-field
part. For the selected trajectory these bounds and all fixed coefficient
products are subpower in $a_h^{-1}$. The last-step endpoint estimate
and chessboard seminorm estimate for the residual required here are
\begin{align}
 \|J_h-I\|_{1\to1}&\le
           2(1+C_h)Q_hL_{h-1}^{-1/8}\le\tfrac14,
 \label{ch13:coarse-matrix}\\
 \|\text{trial residual}\|_{\mathcal C_k}
       &\le C_{k,m,d}E_h e^{-(\tau_h-\tau_m)/8}.
 \label{ch13:coarse-residual}
\end{align}
Here $k,m$ are fixed before $h$; $\mathcal C_k$ is the centered
cell tiling at level $k$ from Section~\ref{ch13:cells}. Both estimates must concern the history restricted from
the actual cutoff $N$, uniformly in $N\ge h$ and the physical period.

The first inequality is the endpoint comparison used in
Proposition~\ref{ch11:equation-comparison}. Its final $1/4$ is an
additional finite stage requirement: require $L_{h-1}^{1/8}\ge8(1+C_h)Q_h$. For the normalized small-field polynomial the local smooth cutoff
is identically one on the ordinary chart, so its ordinary derivative
constant is fixed before this block choice. Its global value is controlled by
the square-root large-field bound, by a larger coupling threshold as in
Section~\ref{ch11:equation-calibration}. This does not replace $E_h$ by
$Q_h$ in the second estimate.

The second inequality requires more than a small coefficient linear in the sources.
We prove its complete marked version next. The source weight is increased
across successive scales; it is not increased by its entire final factor
on an already formed coarse component.

\begin{lemma}[The coarse-trial residual on reflected tilings]
\label{ch13:coarse-residual-supplier}
On the selected histories, after the finite coupling-threshold requirements
specified below, \eqref{ch13:coarse-residual} holds with all incoming
regular and large-field terms involving at least two insertions retained. Its constant is independent
of the outer cutoff, the period, and the number of replica cells.
\end{lemma}
\begin{proof}
Fix the degree ceiling before the schedule. At level $i$, choose a
positive number $\zeta_i$ so that a term linear in the sources with complete normalized
size at most $\zeta_i$ fits every local attachment and adjacent
comparison bound of Chapter~\ref{ch11:insertions}. Include the finite
signed symmetry orbit in these bounds. More explicitly, if $D_i$ is
the complete adjacent comparison constant and $u_i>0$ the terminal
attachment bound, require
\[
 C_i\zeta_i\le1,\qquad 2D_i\zeta_{i+1}\le1,\qquad
 e^{81\sigma_i+162}Z_i^{\rm val}\zeta_i/u_i\le1.
\]
Here $Z_i^{\rm val}$ bounds the global section and its complementary
cutoff part; this symbol denotes a bound, not a partition function.
Also decrease $\zeta_i$ to put the incidence sum for terms linear in the sources below the
small constant in \eqref{ch11:higher-induction}. These choices concern
normalized coefficients, before the physical factor $a_i^{-d}$.

These positive upper bounds can be chosen with
$|\log\zeta_i|=O_j(1+\log\widehat\beta_i)$ on every fixed stage.
Indeed the section derivatives, inverse small-field thresholds, and
adjacent comparison constants have finite powers of
$\widehat\beta_i$, with stage-dependent constants. The terminal attachment bound may be taken as $u_i=\tfrac14\widehat\beta_i^{-1/2}$; the original large-field suppression
is reserved separately in the attachment estimate and need not be
inverted to define $u_i$. The factor $e^{81\sigma_i+162}$ is constant
on each stage. Taking the minimum of this finite list of reciprocal
bounds gives such a choice. At transitions between adjacent stages use the minimum
of both lists. All requirements below therefore follow by increasing
the coupling thresholds, since $B_i$ grows as a positive constant times
$(\log\widehat\beta_i)^2$ on a fixed stage.

Choose $0<\phi_i\le1$ such that
\begin{equation}
 \phi_i L_i^{d+1}\max(1,C_{i+1})\zeta_i\le\zeta_{i+1},
 \qquad \phi_i^{-1}\le e^{B_i/128}.
 \label{ch13:residual-height}
\end{equation}
For example take the minimum of one and the ratio in the first
inequality. The logarithm of its inverse is bounded by a finite
stage constant times $1+\log\widehat\beta_i$, plus
$(d+1)\log L_i$, and hence is eventually below $B_i/256$.
Any finite local incidence bound $\delta_i>0$ can simultaneously
be included so that $\delta_i\phi_i e^{B_i/128}\ge1$:
its inverse has at most another fixed power of
$\widehat\beta_i$. This verifies the comparison and attachment bounds as well as the displayed, weaker upper bound on
$\phi_i^{-1}$. These are countably many finite degree requests to
the single schedule of coupling thresholds.

Put $G(h,i)=e^{(\tau_h-\tau_i)/8}$ and choose
\[
 w=\frac{\zeta_m a_m^dG(h,m)}{2E_h},\qquad
 \omega_i^{\rm res}=w\prod_{j=m}^{i-1}\phi_j.
\]
The residual coefficients from \eqref{ch11:descent-error}, with the
fixed harmless column factors included in $E_h$, satisfy
$c_i+|x_i|\le E_h/G(h,i)$. Since the ratio of
$a_i^{-d}/G(h,i)$ at successive levels is $L_i^{d+1/8}$,
\eqref{ch13:residual-height} gives
\begin{equation}
 \omega_i^{\rm res} a_i^{-d}\bigl(c_i+|x_i|\bigr)
       \le \omega_i^{\rm res} a_i^{-d}E_h/G(h,i)\le\zeta_i/2
                  \quad(m\le i\le h).
 \label{ch13:residual-height-bound}
\end{equation}
In particular the negative normalized local insertion used in the subtraction has
normalized cost at most $\zeta_h$ at $h$. Above $h$, its inverse
coefficient bounds grow by at most $C_{i+1}$ per step; the same bound
recurrence controls these factors and the physical factor $L_i^d$.
Thus its terminal source satisfies the joint attachment bound at $N$,
and its complete push to $h$ is available at the weights $\omega_i^{\rm res}$.

Here is the support calculation for its terms involving at least two insertions, including those
formed above $h$. Open the centered cells at $m$ with spacing at least
$B_m^2$. This spacing persists because $L_i\ge3$ and the adjacent
support radii satisfy the schedule of \eqref{ch13:sharp-schedule}.
Apply reweighting on the incoming fine support $X$. If it contains
$r\ge2$ distinct marked ancestors, then
$B_{i+1}r\le64c_{B_{i+1}}(X)$ by \eqref{ch11:separated-cover}.
Since $B_i\le B_{i+1}$ and $\omega_i^{\rm res}/\omega_{i+1}^{\rm res}=\phi_i^{-1}$, one factor from the incoming exponential support weight gives
\[
 (\omega_i^{\rm res}/\omega_{i+1}^{\rm res})^r e^{-c_{B_{i+1}}(X)}
 \le e^{B_i r/128-c_{B_{i+1}}(X)}
 \le e^{-c_{B_{i+1}}(X)/2}.
\]
To make this argument iterable, let $A_i>4$ be the full exponent
maintained in the norm for terms involving at least two insertions at level $i$, and impose
\[
 10\cdot15^4(A_i+9)\le L_i(A_{i+1}-4).
\]
Reserve four units of the incoming $A_{i+1}$ norm for this weight
increase, site conversion, and attachments. The support contraction
of Chapter~\ref{ch11:insertions} converts the remaining exponent
$A_{i+1}-4$ into the full outgoing exponent $A_i$, including its local
reserves. Thus the estimate does not lose two exponent units at each
successive step. Choosing the bounds for terms linear in the sources small enough makes
\eqref{ch11:higher-induction} preserve the bound one in these full norms
for both arrays of terms involving at least two insertions. This calculation applies to every cardinality and
period. A large-field array is treated as one incoming factor with
its complete original prescribed set of large-field sites, not as a product of its coefficients linear in the sources. The terminal joint construction reserves the original
large-field decay for the complementary factors. Hence there is no
uncontrolled reweighting by $w^r$ at the intermediate level $h$.

At that level insert the positive local trial in every cell.
If $\mathscr A_h$ is the complete pushed negative source, the new polynomial is
\[
 \mathscr A_h\exp_{\rm nil}\Bigl(\sum_v z_vF_v\Bigr).
\]
In its local representation this adds $F_v$ only to the regular
entry linear in the sources. All regular terms involving at least two insertions and all original large-field
terms remain unchanged. The scalar part is exactly the coefficient of the identity in the global exponential.
The nonconstant part multiplied by the smooth cutoff has zero complementary profile because
its small-field restriction is included in its definition.
Linearity makes the new term linear in the sources precisely the residual whose
coefficients satisfy \eqref{ch13:residual-height-bound}; its adjacent
regular and large-field comparisons have bounds at most one by the
choice of $\zeta_i$. Descend this whole polynomial to $m$. Its regular terms linear in the sources have bound at most two, and the
large-field terms linear in the sources have bound at most one, and its two incidence bounds for terms involving at least two insertions are at most one.

Let $\mathsf P_{h,N}$ push densities from level $N$ to level $h$ by
successive applications of \eqref{ch04:block-operator} along the chosen
history. Its adjoint $\mathcal P_{N,h}$ pulls observables back from $h$ to $N$.
Let $\mathcal I$ index the cells of the reflected tiling and let $-U_v$
denote the terminal insertion used for cell $v$. In the product above
the coefficient of $\prod_{v\in\mathcal I}z_v$ at level $h$ is
\[
 \sum_{A\subseteq\mathcal I}
 \mathsf P_{h,N}\left(\rho_N\prod_{v\notin A}(-U_v)\right)
                       \prod_{v\in A}F_v.
\]
Duality and the disjoint-cone identity
\eqref{ch13:cone-factorization} identify its integral with the original
partition function times the expectation of
$\prod_v(\mathcal P_{N,h}F_v-U_v)$.
No product of separately pushed densities for single insertions has replaced it.
Equation~\eqref{ch11:root-bound} now yields
\[
 \|\mathcal P_{N,h}F-U\|_{\mathcal C_k}
 \le w^{-1}\exp\bigl((a_k/a_m)^4q_m^{\rm site}+5\bigr)
 =\frac{2e^{(a_k/a_m)^4q_m^{\rm site}+5}}{\zeta_ma_m^d}
       E_h e^{-(\tau_h-\tau_m)/8}.
\]
The prefactor depends only on the fixed stopping data. This proves
\eqref{ch13:coarse-residual}, once the cones have the room verified
below. All identities and estimates use the original selected history.
\end{proof}

\begin{proposition}[Coarse bounded approximation]
\label{ch13:bounded-singleton}
Use the selected histories with the stage requirements above. Fix a scaling-dimension bound $d$, a real reference coefficient $c$, a
sufficiently fine centered tiling $k$, and $\delta>0$. There are real
observables $\mathcal B_{N,p}(c,y)$ with a finite bound independent of
$N,p,y$, such that
\[
 \|\mathcal B_{N,p}(c,y)-O_{N,p}(d,c;y)\|_{\mathcal C_k(y)}\le\delta
\]
on the ordered cutoff and volume tails. Both observables are interior
to the same cell. The construction includes their signed and conjugated
images with the same bound.
\end{proposition}
\begin{proof}
Fix $m\ge k$ with $B_m^2\le a_k/a_m$. At a finer level $h$, use the
coarse trial
\[
 V_h(v;U)=a_h^{-d}\left(v_0+
       \mathbf1_{\text{small-field stencil}}
                    P_d^{\rm loc}[S_{1/t_h}v^\circ](U)\right).
\]
The scalar part is retained exactly. The full inspected stencil includes
the links needed by the relaxed small-field condition. For fixed $h$ this is
a bounded measurable function, with bound uniform in the outer cutoff
and period: its coefficient range and allowed couplings are compact,
and its section is a fixed polynomial in compact group variables.

Let $\mathcal K_{N,h}(U,dV)$ be the composition of the one-step probability kernels
$k^{(L_i)}_\beta(U,V)\,dV$ of \eqref{ch04:kernel}, in order from
$i=N-1$ down to $h$ and with the coupling of each step, and set
$\mathcal P_{N,h}F(U)=\int F(V)\mathcal K_{N,h}(U,dV)$.
It preserves real values and contracts the supremum norm. Locality here
must include every random variable used to choose a gauge-covariant block average or gauge
frame. Conditional on the preceding link field, use independent primitive
random variables for those choices and for each coarse-link sample.
The full ancestor cone of a stencil contains its fine-link reads and
all these random variables at every intermediate level. For disjoint
cones, conditioning and integrating the independent variables gives
\begin{equation}
 \prod_C\mathcal P_{N,h}F_C
                   =\mathcal P_{N,h}\left(\prod_CF_C\right).
 \label{ch13:cone-factorization}
\end{equation}
This does not assert that coarse links remain conditionally independent
after a shared frame choice has been averaged out. The full cones,
including that choice, must be disjoint. Here is an explicit uniform
radius bound. Let $j_i$ be the smoothing order in the step $i+1\to i$.
In Chapter~\ref{ch04:blocking}, each endpoint frame reads a box of
fine-lattice radius $(j_i+3)L_i$; the other endpoint is at distance
$L_i$, and the connector paths stay in these boxes and the intervening
strip. Thus every input link and every shared frame seed for an output
link lies within physical distance $(j_i+7)a_i$ of its initial endpoint,
in the maximum norm. This bound also includes eligibility conditions and
one boundary-link margin. If the coarse inspected stencil has radius
$r_d$ in level-$h$ units, its full ancestor cone therefore has radius at
most
\[
 r_da_h+\sum_{i=h}^{N-1}(j_i+7)a_i.
\]
The growing smoothing orders require a schedule condition, not merely
finiteness at each step. Impose $(j_i+7)\le a_i^{-1/2}$ at all sufficiently
late levels. This is a permitted finite coupling-threshold requirement: if
$\jmath_q$ is the largest smoothing order used at stage $q$ or its
permitted adjacent pairs, require
$b_q\log3\ge2\log(\jmath_q+7)$ among the coupling-threshold conditions of
Section~\ref{ch10:schedule}. When $J(i+1)=q>1$, that schedule has
$i\ge b_q$ and $a_i\le3^{-i}$, giving the desired inequality.
The finitely many initial levels are irrelevant after increasing $h$.
Since $a_{i+1}\le a_i/3$, the radius is bounded, independently of $N$, by
\[
 r_da_h+\frac{a_h^{1/2}}{1-3^{-1/2}}\longrightarrow0.
\]
Choose $h$ so that this is smaller than the fixed interior margin of
each reflected level-$k$ cell. The cones, including their shared seeds,
then lie strictly inside disjoint cells; sufficiently large physical
periods prevent wrapping identifications.

The bound \eqref{ch13:coarse-matrix} gives
$\|J_h^{-1}\|_1\le2$. Thus the reference insertion with coefficient
$c$ is a linear combination of the trial columns, with sum of absolute
coefficients bounded by
$2\|c\|_1\exp(\sum_{i<h}\log C_{i+1})$. Replace these columns by
$\mathcal P_{N,h}V_h(e_j)$ and call the resulting observable
$\mathcal B'$. Its supremum is finite for fixed $h$ and independent
of $N,p,y$.

Replicate the discrepancies before applying the kernel. At level $h$
the coarse factors multiply the complete source density, including its
terms involving at least two insertions. Equation~\eqref{ch13:cone-factorization} identifies the
resulting residual with the product of the pullback discrepancies.
Lemma~\ref{ch13:coarse-residual-supplier} and recombination give
\[
 \|\mathcal B'-O(c)\|_{\mathcal C_k}
 \le C_{k,m,d,c}\exp\left(\sum_{i<h}\log C_{i+1}\right)
                       E_h e^{-(\tau_h-\tau_m)/8}.
\]
The prefactor is $e^{o(\tau_h)}$, so choose $h$ to make this at most
$\delta$, before choosing the cutoff. For real $c$ take
$\mathcal B=\operatorname{Re}\mathcal B'$; conjugation invariance
and the triangle inequality for the chessboard seminorm preserve the
bound. Translations and the finite symmetry orbit use the same stencils
and constants. Outside the ordered tails set the approximants to zero
if a bound for every index is desired.
\end{proof}

\begin{corollary}[Real word differences]
\label{ch13:bounded-word}
On these selected histories, the bounded-real approximation used for phase
reflection in Section~\ref{ch15:phase-symmetry} holds on a cofinal set of
generic grids. Its supremum bound is independent of copy number and copy
position; approximation thresholds may depend on each fixed pair.
\end{corollary}
\begin{proof}
Choose the grid level after the degree bound and any prescribed lower
level, then fix a generic real word $w$ of length $n$. Include all its
signed and conjugated labels in a finite family. Choose centered cells
with margins fitting between its distinct coordinate projections, and
let $C$ bound the single-insertion seminorms there. Apply
Proposition~\ref{ch13:bounded-singleton} with error $\delta\le1$.
The approximants have seminorm at most $C+\delta$, regardless of their
possibly much larger supremum norms.

For a product of $q$ such separated factors, telescoping and the
chessboard inequality give
\[
 \left|\left\langle\prod_{a=1}^q\mathcal B_a-\prod_{a=1}^qO_a
               \right\rangle\right|
                  \le q\delta(C+\delta)^{q-1}.
\]
Form the difference $g_j$ of the bounded products for a translated word
and its linear time reflection, using the same single-insertion approximation
at every translated copy. If their supremum bounds for single insertions are $M_i$,
then $|g_j|\le2\prod_iM_i=:R_\varepsilon$, independent of all copy
positions and indices. For two sufficiently widely separated copies,
expand their product of differences into four products of $2n$ factors.
Choose $\delta$ with $8n\delta(C+1)^{2n-1}\le\varepsilon$.
The preceding inequality gives the desired pair error. Only finitely
many configurations are involved at each fixed pair, so ordered tails
can be intersected. The empty word difference is zero. Increasing the
initial lower-level request gives the required cofinal grids.
\end{proof}


% !TEX root = ../monograph.tex
\section{Scaling dimension and the physical quotient}
\label{ch18:graded-quotient-section}

An estimate on a polynomial coefficient is useful for local product replacement only if
it retains the degree of the field that coefficient represents.
We derive the required quotient estimate from the triangular marked
transfer. This step uses the corrected physical inclusions and the
normalized-null statement of Section~\ref{ch13:quotient}; those statements
were established using the exact microscopic equation space.

Fix an output scaling dimension $t$. For $v\in V_t$, let
$x_l(v)=(x_{l,\alpha}(v))_\alpha$ be the vector of polynomial coefficients
after transport to level $l$, in the fixed basis of polynomials of
scaling dimension at most $t$. Write $|x|_1=\sum_\alpha|x_\alpha|$. The leading polynomial map at one step has the form
$M_l\mathsf d_l$, where
\[
 (\mathsf d_lx)_\alpha=L_l^{t-|\alpha|}x_\alpha.
\]
Here $|\alpha|$ is the scaling dimension of the basis polynomial, not
its factor count. The leading map fixes the constant and is triangular:
$(M_l)_{\alpha\beta}=0$ off the diagonal whenever
$|\alpha|\le|\beta|$. Both $M_l$ and its inverse have norm at most the
prescribed stage constant $K_l$. The full map differs from it by the
marked remainder bound of Section~\ref{ch11:one-graph}.
Enlarge the prescribed column bounds $C_l$ and block floors so that
\begin{equation}
 64K_l^2\le C_l,C_{l+1},\qquad C_{l+1}\le L_l,
 \qquad 2K_lL_l^{-3/8}(1+C_{l+1})\le1.
 \label{ch18:graded-room}
\end{equation}
These are finite-stage inequalities, with constants fixed before the
block ratio, and the schedule gives
$\sum_{l<m}\log C_{l+1}=o(\tau_m)$.

\begin{proposition}[Graded quotient estimate]
\label{ch18:graded-estimate}
For $0\le i\le t$ there are constants $C_{t,i}(m)$ with
$\log C_{t,i}(m)=o(\tau_m)$ such that the following holds.
Suppose $v_{N,p}\in V_t$ converges on an ordered cutoff and volume
refinement to $v$, and $|x_m(v_{N,p})|_1\le\epsilon$ on that tail.
Then
\begin{equation}
 \operatorname{dist}_{F_t}([v],\iota_{i-1,t}F_{i-1})
       \le C_{t,i}(m)a_m^{-(t-i)}\epsilon,
 \label{ch18:graded-distance}
\end{equation}
where $F_{-1}=0$. Constants are independent of the refinement and the
convergent label family. The norm on $F_t$ is fixed once.
\end{proposition}
\begin{proof}
For $i>0$ let $P$ project onto scaling dimensions below $i$.
Triangularity makes $\ker P$ invariant under $M_l$ and its inverse.
The degree scaling expands low coordinates by at least
$L_l^{t-i+1}$, and high coordinates by at most $L_l^{t-i}$.
The one-step marked bound and \eqref{ch18:graded-room} imply
\begin{align}
 |x_{l+1}|_1\le C_{l+1}|Px_{l+1}|_1
 &\ \Longrightarrow\quad
 |x_l|_1\le C_l|Px_l|_1,\quad
 L_l^{t-i+1}|Px_{l+1}|_1\le C_{l+1}|Px_l|_1,
 \label{ch18:low-cone}\\
 Px_{l+1}=0
 &\ \Longrightarrow\quad
 |x_l|_1\le C_{l+1}L_l^{t-i}|x_{l+1}|_1.
 \label{ch18:high-step}
\end{align}
To check the first implication, put $Y=|P\mathsf d_lx_{l+1}|_1$.
The premise and $C_{l+1}\le L_l$ give
$|\mathsf d_lx_{l+1}|_1\le2Y$ and $|x_{l+1}|_1\le Y$. Applying the inverse leading map
to the low part and absorbing the marked error gives
$Y\le2K_l|Px_l|_1$. The full output is at most $3K_lY$;
\eqref{ch18:graded-room} gives both assertions. When the low part
is zero, the high-degree scaling and $K_l+1\le C_{l+1}$ give
\eqref{ch18:high-step}.

At the microscopic level take the low plane $PV_t$ and transport it.
Equation~\eqref{ch18:low-cone} shows that its projection onto the low
coordinates at every later level is injective, and hence onto that
fixed-dimensional coordinate space. At level $m$ subtract its vector
with the same low coordinates as $x_m(v_{N,p})$. The remaining
coefficient has zero low part and norm at most $(1+C_m)\epsilon$.
Repeat the subtraction after each step towards the fixed reference
level, applying \eqref{ch18:high-step} between subtractions. The sum of
all subtracted vectors still lies in the transported microscopic low
plane. Calling its reference label $y_{N,p}$, we obtain
\[
 |v_{N,p}-y_{N,p}|_1\le
 G_m(1+C_m)(a_{r(t)}/a_m)^{t-i}\epsilon,
 \qquad G_m=\prod_{l<m}(1+C_l)C_{l+1}.
\]
Finite initial transfer and norm-equivalence constants are absorbed in
$G_m$. Its logarithm is $o(\tau_m)$ by the prescribed column bounds.

It remains to identify this microscopic low plane with physical lower
degrees; it is not the naive inclusion of unnormalized polynomials.
The labels $y_{N,p}$ are bounded by the displayed inequality.
The expansion in \eqref{ch18:low-cone}, iterated from the microscopic
level, bounds their microscopic low coefficients by
$a_N^{t-i+1}e^{o(\tau_N)}$. Therefore their degree-$t$ physical
representatives can be written with the normalization $a_N^{-(i-1)}$
and coefficient/derivative growth $e^{o(\tau_N)}$, on the same fixed
stencil. Their global value bounds also have this growth, since the
threshold inverses have subpower growth.

The local-trial comparison of Chapter~\ref{ch11:insertions} matches
these representatives to degree-$(i-1)$ reference labels with vanishing
separated residuals. The matched coefficients need not be bounded.
Apply the normalized-null representative argument of
Section~\ref{ch13:quotient}: division by one plus the norm of their
corrected projection makes the equation remainder uniformly controlled,
and its separated responses vanish. Since $y_{N,p}$ has a bounded
reference limit, that argument gives a physical degree-$(i-1)$ field
with exactly its limiting responses. The all-spectator definition and
corrected inclusion thus put every such limit $[y]$ in
$\iota_{i-1,t}F_{i-1}$. Passing to a convergent subsequence of the bounded
$y_{N,p}$ proves \eqref{ch18:graded-distance}. Extending a tail-defined
bounded family by zero outside its tail changes none of these limits.

For $i=0$ no low plane is removed. The reverse coefficient estimate
directly bounds $|v|_1$ by
$e^{o(\tau_m)}(a_{r(t)}/a_m)^t\epsilon$. Passing to the quotient
proves the same conclusion. All fixed reference spacings are absorbed
in $C_{t,i}(m)$.
\end{proof}

The power $t-i$ in \eqref{ch18:graded-distance} is what recovers a
coefficient of degree $i$ from an auxiliary extraction at a much higher
degree $t$. An ungraded bound on $[v]$ would lose precisely that power.
The proof also explains why identifying the lower field requires all
spectator responses, even though its eventual coefficient is measured
by a fixed finite family of probes.

% !TEX root = ../monograph.tex
% May be placed in Chapter13 after the graded quotient estimate.
% This fragment uses neither phase decomposition nor stress-tensor Ward identities.
\section{Local replacement of a cluster on the lattice}
\label{ch18:lattice-products}

We prove a local replacement estimate before using any phase decomposition.
It will also apply to the undecomposed correlations, before any decomposition into phases, with a fixed exterior word.
Observable source variables and local polynomial extraction are standard
renormalization-group tools; see, for example,
\cite{BauerschmidtBrydgesSlade2015}. The quantitative estimates below concern
the block kernels, activity norms and quotient fields constructed
here; that reference does not give these particular estimates.

For $T\to\infty$, write $f\preccurlyeq e^{cT}$ if, for every
$\varepsilon>0$, eventually $|f|\le e^{(c+\varepsilon)T}$ on an ordered
cutoff and volume tail, uniformly on the specified bounded label sets
and geometries. A factor $e^{o(T)}$ is called subpower in this paragraph.
Choose scales $l=l(T)$ with $\tau_l/T\to\theta$, where
$\tau_l=\log(1/a_l)$. The sequence of scales has $\log L_l=o(\tau_l)$, so
first-crossing scales give any positive $\theta$. For $\theta=0$ one may
first cross $\sqrt T$.

\subsection{Signed local transfer}

A local degree-$t$ density with insertions consists of a column coefficient
$x_l$ and a remainder of norm $r_l$ whose extracted coefficient vanishes. Its dimensionless
normalization obeys
$H_l=L_l^t\mathcal B_l H_{l+1}$; the dimensionful density with insertions has
no additional factor. These conventions are those of
Section~\ref{ch11:one-graph}. Keep the full prescribed large-field set and
its normalization by the square root of the large-field bound in this remainder norm.

The backward comparison of Section~\ref{ch11:trial-descent} applies
on any sufficiently late interval $[k_*,m]$, not only from the original
cutoff scale. In detail, a remainder of size $q$ with vanishing extracted coefficient has
output coefficient and remainder at most $L_l^{-2/5}q$.
The regular term uses its zero radial Taylor polynomial through $t$;
the large-field term retains half the suppression of its original
prescribed large-field set. Subtracting its output column restores the vanishing extracted coefficient and
adds at most $(1+2C_l)L_l^{-2/5}q$. The lower bound on the block ratio makes this at
most $L_l^{-3/8}q$. This recentering is an equality of
densities with insertions, including the terms with prescribed large-field sets, and uses linearity rather
than division by an activity.

Add a free column parameter at the finest scale and choose it so that
the transferred residual coefficient vanishes at $k_*$. Invertibility
of the column map permits this on the same history. The cone calculation
in Section~\ref{ch11:trial-descent} then gives
\begin{equation}
 r_l+|x_l|_1\le4(1+C_m)q_{\rm fin}(a_m/a_l)^{1/8},
 \qquad r_m\le2(1+C_m)q_{\rm fin}.
 \label{ch18:signed-transport}
\end{equation}
Here $q_{\rm fin}$ is the combined original coefficient and remainder
bound, before the affine column subtraction. To recall the numerical
point, with $Y=|\mathsf d_lx_{l+1}|_1$, $U=|x_l|_1$ and
$\epsilon_l=L_l^{-3/8}$ the two inequalities are
$Y\le K_l[U+\epsilon_l(Y+r_{l+1})]$ and
$r_l\le K_lY+\epsilon_lr_{l+1}$.
The cone $32(1+C_l)^2U\le r_l$ gives
$Y\le4K_l\epsilon_lr_{l+1}$ and $L_l^{1/8}r_l\le r_{l+1}$.
It starts with $x_{k_*}=0$. At the finest scale absorb the subtracted
column bound to obtain the second inequality, and iterate for the first.
The sign of the local difference has been retained throughout.

\subsection{Grouping insertion points and retaining their supports}

Use the square-zero source algebra of Chapter~\ref{ch11:insertions}.
The total number of copies may be arbitrary, but a reflected cell
contains a fixed number $n$ of insertion points. A group is a set of
these insertion points with a common representative. Grouping is bookkeeping for
supports and source weights: original variables and the coefficient of every source monomial
remain in the expansion. In particular partial modes of a group are
not discarded, and grouping is not asserted to be an algebra
homomorphism replacing their product by a new variable.

Recall that $c_B(X)$ counts the radius-$B$ balls in a smallest connected covering of $X$, and $\sigma$ is its exponential support weight.
If the support $X$ of a term with insertions has to contain new representatives joined
by at most $h$ covering balls per insertion point, then
\[
 c_B(X\cup\{y_i:i\in I\})\le c_B(X)+h|I|.
\]
Shrinking each affected weight by $\eta e^{-\sigma h}$ absorbs this
support enlargement. The incidence norm contributes a second factor: a new
incidence can be attributed to an old insertion point on its joining path.
If at most $\mathfrak N$ representatives occupy a covering ball, there
are at most $(1+\mathfrak N)n$ assignments per representative.
Take $\eta\le[(1+\mathfrak N)n]^{-1}$. Supports without insertions are
unchanged. This is a termwise estimate on the full
series with insertions, hence is uniform in its number of replica cells.

At each scale retain a maximal $B_l^2$-separated set of representatives,
keeping the representative of the center, and join each discarded
representative to a retained neighbor. A merger decreases the number
of groups; there are at most $n-1$ mergers within a cell.
A deliberately generous joining bound is
\[
 h_l=16[1+L_l(B_l+B_{l+1}+1)^2]^4.
\]
Impose $(\sigma_{l+1}+1)h_l=o(\tau_l)$ by choosing the lower bounds on the reference coupling at each stage.
At fixed stage $h_l$ is a polynomial in $\log\widehat\beta_l$
with fixed ratio, so this is a finite admissible request.
The occupancy factors have the same property. The minimum
$\alpha_j\le1$ of all weight-reduction factors for support enlargement on $m\le l<j$
therefore has $-\log\alpha_j=o(\tau_j)$: split the maximum
cost into a fixed initial segment and a tail bounded by
$\varepsilon\tau_j$, and then let $\varepsilon\downarrow0$.

Between mergers a single group loses at most the factor $6L_l$ in
its weight. Here the incoming regular contribution is carried
\emph{without} Taylor subtraction. Its full factor for the change of cover radius
is $T_l\le L_l$, as in \eqref{ch11:cover-transfer}; it
must not be assigned the contraction factor for a vanishing extracted coefficient.
If its incoming mass is $\mathfrak m$, the other local terms with vanishing extracted coefficient have total mass $p_l$ with $2p_l\le\mathfrak m$
after increasing the lower bound on the reference coupling, and the large-field terms have normalized mass
$f_l\le\mathfrak m$. These are the termwise local bounds used
above, retaining the original events. Resolving real and imaginary
parts contributes at most a factor two, so
\[
 2(p_l+T_l\mathfrak m)+f_l
             \le(2L_l+2)\mathfrak m\le6L_l\mathfrak m
             \le L_l^4\mathfrak m.
\]
The last inequality holds for $L_l\ge3$. It applies in the
submultiplicative norm in the source variables to a whole incoming joint factor;
it does not multiply separate bounds for correlated insertions.
With $g$ groups left, the product of their remaining source weights is consequently
at least
\begin{equation}
 (a_j/a_l)^{4n}(\alpha_j/8)^{n(n-g)}\prod_iw_i.
 \label{ch18:group-weight}
\end{equation}
Each step loses at most the factor $L_l^{-4n}$, and a merger uses at most
$n$ of the weight-reduction factors for support enlargement. This proves the formula by
induction, including all data on supports and prescribed large-field sets.

\subsection{Extraction with a freely chosen residual weight}

Put
\[
 \delta_l^{\rm mark}=\widehat\beta_{l+1}^{-2},\quad
 \phi_l=\widehat\beta_{l+1}^{-1},\quad
 \Xi_l=e^{\sqrt{\widehat\beta_{l+1}}},\quad
 y_{t,l}=\widehat\beta_l^{-\nu_t}.
\]
Choose the fixed integer $\nu_t$ to absorb the finite polynomial and
attachment losses of the insertions at degree $t$. The stage lower bounds on the reference coupling ensure
$y_{t,l}C_l\le1$ and
$\phi_l L_l^{t+1}y_{t,l}\le y_{t,l+1}$.
These are finite requests on the lower bound on the reference coupling after the ratio is fixed.
All their logarithms, as well as $\log\Xi_l$, are $o(\tau_l)$.

For a fixed exterior reflection cell at scale $k$, let
\[
 k_* =\min\{l\ge k:B_l^2\le a_k/a_l\},\qquad
 N_k=\exp\{(a_k/a_{k_*})^4q_{k_*}^{\rm site}+5\}.
\]
The stopping-scale argument of
\eqref{ch13:sharp-schedule} gives $\tau_{k_*}/\tau_k\to1$
and $\log N_k=o(\tau_{k_*})$. Indeed if $k_*=l+1>k$,
minimality gives $a_k/a_{k_*}\le L_lB_l^2$; the last condition
of that equation bounds both the normalization exponent and
$\tau_{k_*}-\tau_k$ by $o(\tau_{k_*})$.
When $\tau_k=o(T)$ the same proof gives $\log N_k=o(T)$.

Suppose the cluster has become one group at scale $m+1$, with
product $\omega$ of its remaining source weights. Set
\[
 q_m^{\rm ext}=4\Xi_ma_{m+1}^t/\omega,\qquad
                    E_m=8(1+C_m)q_m^{\rm ext}.
\]
Taylor extraction through radial degree $t$ gives coefficient and
complement each bounded by $q_m^{\rm ext}$. To verify this, divide
by this product of weights, apply the integral Taylor formula in the
local derivative norm, and sum the finite coefficient conversions
and the polynomial factors in the support size with their positive margin in the exponent. Their
bounds are fixed stage constants and powers of logarithmic radii
and inverse thresholds, hence eventually at most $\Xi_m$.
The factor $a_{m+1}^t$ converts the dimensionful density with insertions to the
normalized degree-$t$ density. Large-field terms keep their
original normalized majorants. Extraction of the coefficient of
all insertions in a group is a linear operation on the exact density;
products of partial modes are retained before taking it. Each
term is assigned the union of its supports and original prescribed large-field set,
so this operation never divides by the exact-cover gas of the density without insertions.

Choose a residual weight $\varpi>0$ at $k_*$ and thereafter
$\varpi_{l+1}=\phi_l\varpi_l$. The two sufficient inequalities are
\begin{align}
 \varpi(a_m/a_{k_*})^{t+1}y_{t,m}/y_{t,k_*}
        &\le\delta_m^{\rm mark}\omega/(16\Xi_m),\notag\\
 32(1+C_m)\varpi a_{k_*}^{-t}E_m(a_m/a_{k_*})^{1/8}
        &\le y_{t,k_*}/2.
 \label{ch18:residual-weights}
\end{align}
The first absorbs the local polynomial coefficient bound and the second the bound for its signed transfer.
Indeed the recursion for these bounds gives
$\varpi_l\le\varpi(a_l/a_{k_*})^{t+1}y_{t,l}/y_{t,k_*}$.
Combining this with \eqref{ch18:signed-transport} and the dimensionful
factor $a_l^{-t}$, the ratio of required to allowed weight at
scale $l$ is at most its value at $k_*$ times
$(a_l/a_{k_*})^{7/8}\le1$. Thus no intermediate scale has been
left unchecked.

The extracted matched field has coefficient norm at scale $m$
at most $2E_m$. The interior residual has chessboard seminorm
at most $2N_k/\varpi$. To see the latter on arbitrarily many
copies, put one residual in every reflected cell before integrating.
Taking the coefficient of the required source monomial in the sum with insertions contributes the factor $\varpi^{-1}$ per
cell; the normalized integral bound \eqref{ch11:root-bound}
contributes the factor $N_k$ per cell. Apply this separately to the original local
difference and its compensating column, and use the chessboard
triangle inequality. The rooted component estimates and bounds
used in this calculation depend on the number of insertions per cell, not on the
number of cells. All the terms involving at least two insertions generated during transfer
are included in these sums. The estimate is therefore stronger
than a bound on one value of a correlation functional.

\begin{proposition}[Lattice cluster replacement]
\label{ch18:lattice-replacement}
Let a cluster contain $n$ fields of total degree $W$. Choose
$0\le\theta_-<\theta<\theta_+$ and scales $k,m,j$ with
$\tau_k/T\to\theta_-$, $\tau_m/T\to\theta$, and $\tau_j/T\to\theta_+$. The cluster lies inside the centered block of scale $m+1$,
its points lie on the scale-$j$ grid with separation at least
$B_j^2a_j$, and its dependence sets lie inside the scale-$k$
cell. Put $Z=W\theta_++4n(\theta_+-\theta)$.
For each fixed auxiliary degree $t$ there are output labels
$v_{N,p}^{\rm out}\in V_t$, on the actual histories, satisfying
\begin{align}
 |x_m(v_{N,p}^{\rm out})|_1
       &\preccurlyeq e^{(Z-t\theta)T},\notag\\
 \left\|\prod_iO_{N,p}(v_i,x_i)
                  -O_{N,p}(t,v_{N,p}^{\rm out};0)\right\|_{\mathcal C_k}
       &\preccurlyeq e^{[Z-(t+1/16)(\theta-\theta_-)]T}.
 \label{ch18:lattice-estimate}
\end{align}
The difference is interior. Constants and thresholds are uniform in
cutoff and period on the ordered tails and in the number of copies
of this fixed local cluster.
\end{proposition}
\begin{proof}
The weights of normalized local insertions in
\eqref{ch13:weight-size} give $\prod_iw_i=e^{-W\theta_+T+o(T)}$.
Keep the coefficients of all incoming source monomials above $j$; separation permits their
bounds as terms with insertions but does not assert statistical independence. The grouping
estimate gives $\omega=e^{-ZT+o(T)}$ at $m+1$, and therefore
$E_m=e^{(Z-t\theta)T+o(T)}$.
Set
\[
 \varpi=\exp\{[(t+1/16)(\theta-\theta_-)-Z]T\}.
\]
In the first inequality of \eqref{ch18:residual-weights}, the
left exponent is $-Z-15(\theta-\theta_-)/16$, strictly smaller
than the right exponent $-Z$. In the second, the full left
exponent is $-(\theta-\theta_-)/16<0$ while the right side
has exponent zero. Both strict margins absorb every subpower
factor. The extraction estimate and $\log N_k=o(T)$ give the
claim. All support shifts take place within the specified cell;
reflections and conjugations retain the same bounds.
\end{proof}

\subsection{A consequence for undecomposed correlations without phases}
\label{ch18:mixed-replacement}

Fix finite correlation reconstruction maps onto the filtered spaces $F_D$.
More explicitly, choose a basis in each successive quotient and
extend it to $F_D$, with dual coordinates realized by finite
linear combinations of separated probe correlations. Arrange that
a degree-$i$ coordinate vanishes on $F_{i-1}$, and that reconstruction
fixes every field in $F_D$. Keep each old coordinate functional as
its original correlation formula on every larger field space. When
lifting a new quotient basis, subtract its values under all the old
coordinates, so those coordinates annihilate the new basis vectors.
This is possible by finite-dimensional separation by all spectator
correlations and successive subtraction of lower degrees. Consequently
the degree-$D$ reconstruction of a field in $F_t$, $t>D$, is exactly its
low-degree component in this fixed biorthogonal system. Denote the resulting reconstruction of a word by
$\mathcal P_Dw$. It depends on the cluster, not on its later spectator.
No rotation average or phase measure is used here.

\begin{proposition}[Local replacement for undecomposed correlations]
\label{ch18:mixed-bound}
Let $x_i=o+r z_i$, with the $z_i$ in a fixed compact separated
configuration set, and let $W$ be the total cluster degree.
For any fixed separated exterior word $v$, any $D$ and $\varepsilon>0$,
\begin{align}
 |b_\alpha(w;o)|&\le C r^{|\alpha|-W-\varepsilon},\notag\\
 |S(wv)-S((\mathcal P_D(\tau_{-o}w))(o)v)|
                   &\le C r^{D+1-W-\varepsilon}.
 \label{ch18:mixed-product-rate}
\end{align}
Here $b_\alpha$, $|\alpha|\le D$, are the fixed probe coordinates. The same field
$\mathcal P_D(\tau_{-o}w)$ works for every exterior word. The
estimates are uniform on fixed bounded label sets and separated
geometries, and hold with any fixed finite derivatives in the positions,
with their increased scaling dimensions.
\end{proposition}
\begin{proof}
First use generic points of the scale-$j$ grid. Put $T=\log(1/r)$ and
choose $0<\theta_-<\theta<1<\theta_+$. The cluster then fits
the block and separation conditions for large $T$.
For each fixed geometry, the output labels in
\eqref{ch18:lattice-estimate} are bounded at fixed $m$ and hence
at the reference scale by the inverse-column bound. Take a
subsequence along which the coefficients converge, and call its field $A_t$.
The chessboard inequality bounds its residual paired with any
fixed exterior word by exponent
$Z-(t+1/16)(\theta-\theta_-)$. The same subsequence works for all
spectators, since the coefficient limit was taken before pairing.

The graded quotient estimate \eqref{ch18:graded-distance} shows
that a degree-$i$ coordinate of $A_t$ has exponent $Z-i\theta$.
Its coordinate functional vanishes on $F_{i-1}$; this is where
an ungraded coefficient estimate would lose the desired power.
Choose $t$ so large that
$(t+1/16)(\theta-\theta_-)\ge(D+1)\theta$.
This also bounds every raw coordinate being evaluated, whose degree is
at most $D$. We do not require the residual to bound all coordinates
through the auxiliary degree $t$. Pairing the residual
with the fixed probes bounds the coordinate errors. Reconstruct
$A_t$ and omit degrees $i>D$. Each omitted coordinate has exponent
at most $Z-(D+1)\theta$, and its fixed exterior correlations are
bounded. This gives that exponent for the raw replacement error,
and $Z-i\theta$ for each raw coordinate.

For completeness the grid restriction can be removed without
assuming differentiability of the selected coefficients of $A_t$.
Use the fixed correlation reconstruction instead. Its coefficients and
the pairings of its residual are smooth functions of the cluster positions.
Their derivatives replace an input by a field of degree one higher.
The sharp separated bounds, using fixed cells for exterior probes
and shrinking cells only for cluster factors, give a crude exponent
$W+1$: choose total loss one in the exponent. Thus one can take $c_0=1$.
The grid exponent is $\theta_+W+c_1$, where
$c_1=4n(\theta_+-\theta)-(D+1)\theta$ for the residual, and
$c_1=4n(\theta_+-\theta)-i\theta$ for a degree-$i$ coordinate.
These numbers are independent of $W$ with $n,D,i$ fixed.
Multiplicative constants may depend on each of the finitely many
higher derivative orders used in the following induction. For all sufficiently large total
degrees the crude estimate implies the grid target, because
$\theta_+>1$. Use downward induction on the total degree. Moving to a generic grid
configuration contributes the factor $O(a_j)=e^{-\theta_+T+o(T)}$; the derivative
bound at degree $W+1$ gives precisely the target exponent at
$W$. The connecting paths remain separated. Only finitely many
higher derivatives are needed for each starting degree, so their
constants and thresholds can be chosen together. Finally let
$\theta\uparrow1$ and $\theta_+\downarrow1$ and absorb the
remaining loss into the prescribed $\varepsilon$.
This proves the estimate for undecomposed correlations using neither phase symmetry nor
an operator-product theorem.
\end{proof}

}{}
\setcounter{chapter}{13}
% !TEX root = ../monograph.tex
\chapter[Reflection spaces and Ward identities]{Reflection spaces, Ward identities and Euclidean symmetry}
\label{ch14:ward}

We first use reflection positivity in each coordinate direction
separately. Continuous rotations are not assumed. The resulting
Hilbert spaces and bounded insertions provide the positivity estimates
needed to construct rotational Ward identities. Relativistic
reconstruction requires additional arguments and is not a consequence
of the Hilbert-space construction alone.

Throughout, $S$ is the mixed correlation hierarchy of
Chapter~\ref{ch13:continuum}. Its finite-dimensional field spaces are
$F_d$, with the compatible inclusions and derivative maps
constructed there. Statements using this calculus retain the actual
marked estimates required in that chapter. All correlations in the
present chapter are evaluated at distinct points or against test functions
supported away from collisions.

\section{Reflection spaces for the coordinate axes}
\label{ch14:reflection}

Fix an axis $\mu$. A positive word $w=((x_i,A_i))_{i=1}^n$ has
$x_i^\mu>0$ and distinct points; the empty word is allowed.
Let $\theta_\mu w$ reflect the points across $x^\mu=0$, apply the
signed coordinate reflection to the labels, conjugate them, and
reverse their order. On formal finite linear combinations set
\[
 \langle v,w\rangle_\mu=S(\theta_\mu v\,w).
\]
The form is positive semidefinite by reflection positivity. Its null
space is orthogonal to every vector by Cauchy--Schwarz. The quotient
and completion give $\mathcal H_\mu^{\rm OS}$, with dense word
vectors $[w]_\mu$ and unit vacuum $\Omega_\mu=[\varnothing]_\mu$.
This is the standard Osterwalder--Schrader quotient
\cite[Section3]{OsterwalderSchrader1973}. Only the quotient
construction is used here, not the reconstruction theorem.

Let $\tau_t$ mean translation by $te_\mu$. On formal word sums put
$T_tc=[\tau_tc]_\mu$, $t\ge0$. Translation invariance implies
\begin{equation}\label{ch14:heat-pairing}
 \langle T_sc,T_td\rangle_\mu
             =\langle T_0c,T_{s+t}d\rangle_\mu.
\end{equation}
Every fixed orbit $T_tc$ is bounded: the two reflected word halves
have a fixed positive separation and the separated-correlation
bound of Chapter~\ref{ch13:continuum} is uniform as they move apart.
If $M=\sup_{t\ge0}\|T_tc\|$, Cauchy--Schwarz in
\eqref{ch14:heat-pairing} gives $M^2\le\|T_0c\|M$.
Thus $\|T_tc\|\le\|T_0c\|$, including when $c$ is null.
The operators descend and extend to self-adjoint contractions
$E_\mu(t)$. Their semigroup law holds on the dense word span and
hence everywhere. Correlation continuity gives strong continuity.
The spectral theorem yields
\begin{equation}\label{ch14:heat}
 E_\mu(t)=e^{-tH_\mu},\qquad H_\mu\ge0,
 \qquad \lim_{t\to\infty}E_\mu(t)=P_\mu^0
 \quad\hbox{strongly},
\end{equation}
where $P_\mu^0$ projects onto $\ker H_\mu$.

Transverse translations, by $y$ with $y^\mu=0$, preserve the
reflected form and have inverse translations. They therefore give
strongly continuous unitaries $U_\mu(y)$, commuting with $E_\mu(t)$.
Both leave $\Omega_\mu$ fixed. At this stage we do not claim that
$\ker H_\mu$ is one-dimensional, or that its vectors are fixed by
transverse translations.

The commuting spectral measures of $H_\mu$ and the three transverse
momenta give a joint projection-valued measure $P_\mu$ on
$[0,\infty)\times\RR^3$:
\begin{equation}\label{ch14:coordinate-spectrum}
 E_\mu(t)U_\mu(y)
     =\int e^{-tE+ip\cdot y}\,dP_\mu(E,p).
\end{equation}
This is a coordinate half-space spectrum. No cone bound
$E\ge|p|$ has been proved or used.
For $\operatorname{Re}z>0$, the same spectral integral defines the
holomorphic contraction $e^{-zH_\mu}$; its derivatives are bounded
on every smaller right half-plane. This elementary heat analyticity
will be used before any rotation action is constructed.

\section{Bounded insertions in a slab}
\label{ch14:slabs}

A word $w$ lies in a slab of width $h$ if $0<x_i^\mu<h$ for each
of its points. Its reflected slab word is
$w^\sharp=\tau_h\theta_\mu w$. Formally define
$I_wv=w\tau_hv$. Reflection and translation give
\begin{equation}\label{ch14:formal-adjoint}
 \langle[I_wc]_\mu,[d]_\mu\rangle
       =\langle[c]_\mu,[I_{w^\sharp}d]_\mu\rangle.
\end{equation}
This identity is on formal word sums; neither boundedness nor
preservation of null vectors is assumed yet.

\begin{proposition}[Slab norm bound]\label{ch14:slab-bound}
For labels of bounded degree and norm, a word of length $n\ge1$,
with pairwise separation and slab margins at least $r$, $0<r\le1$,
defines a bounded operator $B_h(w)$ on $\mathcal H_\mu^{\rm OS}$.
There are constants depending only on the label bounds such that
\[
 B_h(w)[v]_\mu=[w\tau_hv]_\mu,
 \qquad \|B_h(w)\|\le C^n(n/r)^{\kappa n}.
\]
The constants are independent of the width and the exterior word.
\end{proposition}

\begin{proof}
Proposition~\ref{ch13:word-bound} first gives the estimate at a
width $s\le h$, with a slightly smaller buffer. Put
$T=I_{w^\sharp}I_w$ at that width and
$D=C^n(n/r)^{\kappa n}$. For every formal finite sum $c$ its
repeated-word estimate gives
$\|[T^jc]_\mu\|\le C_cD^{2j}$. By
\eqref{ch14:formal-adjoint}, $T$ is symmetric for the possibly
degenerate reflected form. Set
$b_j=D^{-2j}\|[T^jc]_\mu\|$ and $M_c=\sup_{j\ge0}b_j$.
Then $b_j^2\le\|[c]_\mu\|b_{2j}$, so
$M_c\le\|[c]_\mu\|$. In particular
\[
 \|[I_wc]_\mu\|^2
 =\langle[c]_\mu,[Tc]_\mu\rangle
 \le D^2\|[c]_\mu\|^2.
\]
The operator preserves null vectors and extends boundedly. Increase
the width to $h$ by right multiplication by $E_\mu(h-s)$.
\end{proof}

The empty word gives $B_h(\varnothing)=E_\mu(h)$. Dense word
matrix elements also prove the exact identities
\begin{align}\label{ch14:slab-laws}
 B_h(w)^*&=B_h(w^\sharp),\\
 B_{h+k}(w\tau_hv)&=B_h(w)B_k(v),\notag\\
 B_{h+t}(w)&=B_h(w)E_\mu(t),\qquad
 B_{t+h}(\tau_tw)=E_\mu(t)B_h(w).\notag
\end{align}
Transverse covariance reads
$U_\mu(y)B_h(w)U_\mu(y)^*=B_h(w+y)$.
These assertions concern bounded slab insertions, not unsmeared
point fields acting on an arbitrary Hilbert vector.

\begin{proposition}[Norm-smooth insertion positions]
\label{ch14:slab-smooth}
On the open configuration space inside a fixed slab, $B_h(w)$ is
smooth in operator norm. A position derivative replaces the
corresponding field label by its derivative in $F_{d+1}$.
\end{proposition}

\begin{proof}
For a fixed configuration choose a smaller common buffer. The norm
bound holds locally for the original tuple and each differentiated
tuple. On dense word matrix elements the derivative identity is
exactly the scalar derivative identity of Chapter~\ref{ch13:continuum}.
If a weak derivative $D A(x)[v]=C(x)[v]$ on a dense domain is locally
bounded in operator norm, the scalar mean-value inequality gives
$\|A(y)-A(x)\|\le\sup_{[x,y]}\|C\|\,|y-x|$.
Apply this first to every differentiated tuple, proving continuity
of its operator. The same inequality after subtracting
$C(x)[y-x]$ proves norm differentiability of $A$. Iteration proves
all orders. This argument explains why weak differentiability alone
would have been insufficient.
\end{proof}

\section{Comparing the mixed fields with free fields}
\label{ch14:comparison}

The comparison is a statement about the actual marked construction,
stronger than equality of dimensions of field spaces. We record its
quantitative form and then derive it from the finite-cutoff comparison. For fixed arity and degree bounds,
let $M_k\to\infty$, $r_k\to0$ and $\epsilon_k\to0$. Suppose the
same coefficient lifts in every slot give rescaled kernels $S_k$,
with, on a compact separated window,
\begin{equation}\label{ch14:grid-input}
 |S_k|\le A_k^n,\quad
 |\nabla S_k|\le C_n V_k^n/\delta_k,
 \quad |S_k-S^0|\le\epsilon_k\quad\hbox{on its grid}.
\end{equation}
The gradient bound is required outside coordinate-equality strips
of width $4\delta_k$; the grid has spacing $s_k$, and $S^0$ is the
free correlation. The necessary small errors are
\begin{equation}\label{ch14:grid-cost}
 A_k^n\delta_k\to0,\qquad V_k^n s_k/\delta_k\to0,
 \qquad s_k=o(\delta_k).
\end{equation}
For example the source construction uses
$r_k=a_kM_k^{\alpha_0}$, $s_k=M_k^{-\alpha_0}$,
$A_k=CM_k^{(\alpha_0-\alpha_1)D}$,
$V_k=CM_k^{(\alpha_0-\alpha_2)D}$ and $\delta_k=M_k^c$,
with strictly negative exponents in \eqref{ch14:grid-cost}.
The marked coefficient inverse, Gaussian comparison and displacement
estimates must all concern these same lifts. The following construction verifies this compatibility.

% !TEX root = ../monograph.tex
\subsection{Obtaining the three comparison bounds from one lattice family}
\label{ch14:lattice-comparison}

Here are the scale choices behind \eqref{ch14:grid-input}. Fix the
upper bound $n$ on the number of insertions, maximum scaling dimension $D$ and a compact set of configurations
$K$ with positive point separation. Enlarge the finite label list to
include its signed coordinate images and conjugates. All constants and
all three exponents below are chosen for this one finite request.
Choose positive numbers $\alpha_0,\alpha_1,\alpha_2$ and $c<0$ with
\begin{align}
 \alpha_1+6&<\alpha_0,&\alpha_2+5&<\alpha_0+c,\notag\\
 nD(\alpha_0-\alpha_1)+c&<0,&
 nD(\alpha_0-\alpha_2)-\alpha_0-c&<0.
 \label{ch14:comparison-exponents}
\end{align}
For example first choose $\alpha_0-\alpha_1=7$, next $-c>7nD$,
then $\alpha_0-\alpha_2>5-c$, and finally increase all three
$\alpha_i$ together until the last inequality holds and all are
positive. Increasing their common part does not alter the preceding
inequalities.

At a comparison scale $k$ whose comparison window lies in a single stage put\mgubi{Issue 10: say for which scales this holds ($k=k_j$ suffices for the limit taken below) and cite Chapter~\ref{ch10:trajectory}.}
$H_k=\exp((\log\log\widehat\beta_k)^2)$ and set
\[
 r_k=a_kH_k^{\alpha_0},\qquad s_k=H_k^{-\alpha_0},\qquad
 \delta_k=H_k^c.
\]
Thus $r_ks_k=a_k$: grid rounding occurs on the construction
lattice. Use one common localization map $T$ from
Section~\ref{ch11:common-extraction}. If $v$ is a fixed real polynomial
representative of a free quotient field, its lift $u_{k,N,p}$ solves
\[
                 S_{g/\theta}T u_{k,N,p}=W_{H_k^{\alpha_0}}v.
\]
Here $S_t$ multiplies a monomial with $f$ curvature factors by $t^f$,
whereas $W_\rho$ multiplies one of scaling dimension $j$ by $\rho^j$;
$g=\beta_{N,p,k}^{-1/2}$ and $\theta$ is the local normalization
used for $T$. These two dilations have different purposes.
The inverse bound in \eqref{ch11:simultaneous-detector-lifts} gives
bounded lifts at each fixed $k$, uniformly in the cutoffs and periods.
Their identities for reality, the identity field and signed symmetries are exact. No separate
choice is made for different spectator words.

For clarity, the use for finitely many insertions of the estimate for the corrected coefficients is
as follows. Write $I_i$ for the lattice insertion in argument $i$. For a subset
$J$ of the insertion indices, let $G_J$ be the free correlation of the
corresponding fields at the rescaled positions. Define recursively
\[
 R_\varnothing=1,\qquad
 R_I=\prod_{i\in I}I_i-
          \sum_{\varnothing\ne J\subseteq I}G_JR_{I\setminus J}.
\]
These are the Wick products of the insertions with respect to the free field (the moments $G_J$ are those of the free field, not of the lattice measure).
Equivalently, introduce commuting variables $s_i$ with $s_i^2=0$ and write
$s_J=\prod_{i\in J}s_i$. Then $R_I$ is the coefficient of $s_I$ in
$\prod_i(1+s_i I_i)/(1+\sum_{J\ne\varnothing}G_Js_J)$. The denominator has
constant coefficient one and hence an exact inverse in the finite
squarefree algebra. For the limiting Gaussian variables every
nonempty coefficient has expectation zero. The extraction of coefficients of source monomials,
covariance comparison and absorption of the larger source weight of
Proposition~\ref{ch11:n-mark-comparison} apply to this \emph{entire}
finite algebra, including its joint incoming terms.
With $q=\log H_k/(16n)$ and source weight $e^{2q}$, the output
normalization is at most $e^q$ per reflected cell, and extraction gives
\begin{equation}
 \|R_I\|_{\mathrm{cell}}
       \le e^{q-2q|I|}\le H_k^{-1/(16n)}
       \quad(\varnothing\ne I\subseteq\{1,\ldots,n\}).
 \label{ch14:subset-root}
\end{equation}
The lower bound on the reference coupling at the stage is chosen after this fixed number of insertions and the constants of its
insertions. The estimate is on the original positive partition
function and includes every reflected copy. Uniformity on $K$ uses
positive point separation in the covariance estimate; it does not
require a positive lower bound on each coordinate difference.
Unwinding the finite recursion and applying the chessboard inequality
therefore gives the grid error $C_KH_k^{-1/(16n)}$ at generic grid
points. Distinct coordinates there differ by at least $a_k$ in continuum
units, so the assignments to lattice sites require no rounding.

Apply the same coefficient map and the same finite comparison request
also at exponents $\alpha_1$ and $\alpha_2$. At its own normalization
a single insertion has chessboard seminorm at most $1+|G_{\{i\}}|$. Changing its
scaling dilation to $\alpha_0$ loses at most a factor
$CH_k^{(\alpha_0-\alpha_j)D}$. Set
\[
 A_k=CH_k^{(\alpha_0-\alpha_1)D},\qquad
 V_k=CH_k^{(\alpha_0-\alpha_2)D}.
\]
The comparison cells have sides at most $a_kH_k^{\alpha_j+5}$.
Consequently the first cells are $o(r_k)$, and the second are
$o(r_k\delta_k)$, by \eqref{ch14:comparison-exponents}.
The chessboard inequality gives $|S_k|\le A_k^n$ on $K$.
The modulus of continuity in the insertion points of Section~\ref{ch11:positional}, applied
with margin $r_k\delta_k$ in continuum units, gives
$|\nabla S_k|\le4nV_k^n/\delta_k$ outside the coincident-coordinate strips.
Include a strip for each configuration taking every coordinate from the point or
from its rounding (finitely many); the total volume of the strips is still $O(\delta_k)$.
The compact configuration sets in continuum units lie inside the smallest period for large $k$,
so the periodic and ordinary separations agree.

Take an ordered limit point of the lifts at each fixed $k$. The
uniform bounds pass to the already chosen continuum hierarchy.
Their exponents give exactly \eqref{ch14:grid-cost}, since
\[
 A_k^n\delta_k=CH_k^{nD(\alpha_0-\alpha_1)+c},\qquad
 V_k^ns_k/\delta_k=CH_k^{nD(\alpha_0-\alpha_2)-\alpha_0-c}.
\]
Both tend to zero, and $s_k/\delta_k\to0$ follows from the second
inequality. A diagonal sequence over the finite requests gives a
single comparison sequence. The bounded coefficient limits do not
alter the chosen correlation hierarchy or its order of limits.


The coordinate-equality strips have volume $O(\delta_k)$ on a fixed
window. Their contribution to any test function of supremum norm at most one
is $O((A_k^n+1)\delta_k)$. Outside them, rounding to the grid and
\eqref{ch14:grid-input} cost
$O(V_k^ns_k/\delta_k)$ plus the free modulus of continuity.
It follows that
\begin{equation}\label{ch14:uniform-tests}
 \sup_{\|f\|_\infty\le1,\ \operatorname{supp}f\subset K}
 \left|\int f(x)(S_k(x)-S^0(x))\,dx\right|\longrightarrow0.
\end{equation}
This is local $L^1$ convergence; it includes coordinate ties which
are not collisions. Testing against each fixed $f$ separately would
be weaker than the displayed assertion.

The lifts can be made exactly compatible with inclusions. Write
$i_d:F_d^0\to F_{d+1}^0$ for the free inclusion,
$j_d$ for the physical inclusion and $L_d$ for a real,
hypercubic-equivariant left inverse of $i_d$. From provisional maps
$U_{k,d}$ define recursively
\[
 \widetilde U_{k,d+1}
 =j_d\widetilde U_{k,d}L_d+U_{k,d+1}(I-i_dL_d).
\]
Exact inclusion identities of correlations preserve
\eqref{ch14:uniform-tests} in every slot. The unit is fixed at degree
zero. Free finite detecting families from Chapter~\ref{ch03:maxwell}
make these maps eventually injective; equality of dimensions in
Chapter~\ref{ch13:continuum} makes them isomorphisms. For each finite
request the tail is chosen once, before the test function. A single
sequence for all requests is obtained only by explicitly diagonalizing
arity, degree, compact windows and errors; the numerical comparison
level must be chosen after those finite requests.

\begin{proposition}[Derivative transport and smooth comparison]
\label{ch14:smooth-comparison}
For the comparison maps just described,
\[
 D_{k,d,a}:=r_kU_{k,d+1}^{-1}\partial_aU_{k,d}
                   \longrightarrow\partial_a^0
\]
in finite-dimensional operator norm. Rescaled correlations converge
in $C^j$ on every compact separated window, for every fixed $j$,
also for convergent moving labels.
\end{proposition}

\begin{proof}
Choose free tested correlations giving an injective map
$L:F_{d+1}^0\to\CC^q$. Their corresponding maps $L_k$
converge to $L$ by \eqref{ch14:uniform-tests}; hence
$\|L_kv\|\ge c\|v\|$ eventually. Integration by parts in the
insertion position gives
$L_k D_{k,d,a}v\to L\partial_a^0v$. No bound on
$D_{k,d,a}$ was assumed: subtracting the limiting derivative and
using the lower bound supplies it and its convergence.
Moving labels are handled by expansion in fixed finite bases.
Every fixed derivative of the correlation error consequently tends
to zero locally in $L^1$. Multiply by a smooth cutoff supported
away from collisions. Integration by parts gives a Fourier bound
by $(1+|p|)^{-M}$ times the sum of its $L^1$ derivative norms
through order $M$. Choose $M>4n+j$ and apply Fourier inversion.
\end{proof}

\section{The conserved stress line}
\label{ch14:stress-line}

Let $\mathcal K^0$ and $\mathcal K$ be the real spaces of centered,
real, symmetric, hypercubic-covariant degree-four tensor arrays in
the free and mixed field spaces. Centered means that each component
has zero expectation. Let $\mathcal V^0,\mathcal V$ be the degree-five
covariant vector spaces. The free divergence calculation of Chapter~\ref{ch03:maxwell}
has rank two. To check that its target is the full covariant vector
space, note that degree five permits only pairings of $F$ with $\partial F$ in $\gLie$. Sign parity and permutations leave the two
expressions
\[
 V_b^{(1)}=\sum_c{:}\langle F_{bc},\partial_bF_{bc}\rangle_{\gLie}{:},\qquad
 V_b^{(2)}=\sum_{c<d,\ c,d\ne b}{:}\langle F_{cd},\partial_bF_{cd}\rangle_{\gLie}{:}.
\]
The apparent further contractions
$\sum_{c,d}{:}\langle F_{bd},\partial_cF_{cd}\rangle_{\gLie}{:}$ and
$\sum_{c,d}{:}\langle F_{cd},\partial_cF_{bd}\rangle_{\gLie}{:}$ reduce to these: use the
Maxwell equation in the first and the Bianchi identity in the second.
The two independent rows of correlation pairings in \eqref{ch03:divergence-matrix}
show that these vectors are independent and in the divergence image.
Thus free divergence is onto, with kernel spanned by $T^0$.

Center the componentwise map $U_{k,4}$ after transport. It gives an
isomorphism $C_k:\mathcal K^0\to\mathcal K$: its inverse transports
back and centers there, and preservation of the unit makes the two
scalar corrections cancel. Let $V_k$ be the corresponding vector
isomorphism. The derivative comparison gives
\[
       V_k^{-1}r_k\operatorname{div}C_k
                          \longrightarrow\operatorname{div}^0.
\]
Compose with a fixed right inverse of free divergence. The result
tends to the identity and is invertible for large $k$, so physical
divergence is onto. Rank--nullity therefore gives
\begin{equation}\label{ch14:conserved-line}
                     \dim\ker(\operatorname{div}|_{\mathcal K})=1.
\end{equation}
This step uses the full low-degree rank computation, not just the
existence of a free conserved tensor.

Choose a nonzero $J$ on this line. Normalize its pulled-back arrays
$r_k^4 U_{k,4}^{-1}J$ to have the norm of $T^0$ and take a convergent
subsequence on that finite-dimensional sphere. Smooth comparison
at arity one ensures that its limit is centered. Derivative transport
ensures conservation. It must therefore be $T^0$ or $-T^0$.
Changing the sign of $J$ once gives positive numbers $c_k$ with
\begin{equation}\label{ch14:balanced-stress}
            U_{k,4}^{-1}(c_k^{-1}r_k^4J)\longrightarrow T^0.
\end{equation}
This fixes an orientation of the physical line. It does not yet fix
its Ward normalization, and no positive lower bound for $c_k$ has
been asserted.

\section{Shrinking shells and local charge fields}
\label{ch14:shells}

Choose $\chi\in C_c^\infty(\RR^4)$ equal to one near an insertion
$x_0$ and zero near all spectators. For a Euclidean Killing field
$v(y)$, either constant or $K(y-x_0)$ with $K^{\mathsf T}=-K$, set
\begin{equation}\label{ch14:flux}
 \mathcal Q^J_v(A;\boldsymbol B)
 =-\int\partial_\mu\chi(y)v_\nu(y)
              S(J_{\mu\nu}(y),A(x_0),\boldsymbol B)\,dy.
\end{equation}
Repeated indices are summed. Conservation and symmetry of $J$ show
that the value is independent of the cutoff as long as it separates
the same distinguished insertion from the spectators. Indeed the
difference is the integral of a total divergence in a collision-free
region, and $(\partial_\mu v_\nu)J_{\mu\nu}=0$. In particular one
can shrink a radial cutoff as in Figure~\ref{ch14:fig-ward-shell}.

\begin{figure}[htbp]
\centering
\begin{tikzpicture}[font=\small,
    inner shell/.style={draw=ymInk,semithick},
    outer shell/.style={draw=ymAccent,semithick,dashed}]
  % A central section of two radial cutoffs in four dimensions.
  \fill[ymAccent!13,even odd rule] (0,0) circle (1.80) (0,0) circle (1.20);
  \fill[ymInk!16,even odd rule] (0,0) circle (.99) (0,0) circle (.66);
  \foreach \r in {.66,.99} \draw[inner shell] (0,0) circle (\r);
  \foreach \r in {1.20,1.80} \draw[outer shell] (0,0) circle (\r);
  \fill (0,0) circle (1.6pt) node[below=2pt] {$A(x_0)$};
  \fill (-2.26,.76) circle (1.6pt) node[above] {$B_1(x_1)$};
  \fill (2.23,-.68) circle (1.6pt) node[below] {$B_2(x_2)$};

  % The horizontal scale here is twice that of the central section.
  \begin{scope}[shift={(3.60,-.85)}]
    \fill[ymInk!12] (1.32,0) rectangle (1.98,1.70);
    \fill[ymAccent!10] (2.40,0) rectangle (3.60,1.70);
    \draw[-{Stealth[length=1.7mm]}] (0,0) -- (5.02,0)
        node[above left=1pt] {$|y-x_0|$};
    \draw[-{Stealth[length=1.7mm]}] (0,0) -- (0,2.12);
    \node[left=2pt] at (0,1.70) {$1$};
    \node[left=2pt] at (0,0) {$0$};
    \draw[inner shell] (0,1.70) -- (1.32,1.70)
        .. controls (1.60,1.70) and (1.70,0) .. (1.98,0) -- (4.42,0);
    \draw[outer shell] (0,1.70) -- (2.40,1.70)
        .. controls (2.90,1.70) and (3.10,0) .. (3.60,0) -- (4.42,0);
    \node[anchor=east] at (1.53,.88) {$\chi_\rho$};
    \node[anchor=west] at (3.18,.88) {$\chi_R$};
    \foreach \r/\name in {1.98/\rho,3.60/R}
      \draw (\r,0) -- +(0,-.09) node[below=1pt] {$\name$};
  \end{scope}
\end{tikzpicture}
\caption[Deformation of a local Ward shell]{Shrinking the Ward shell. The two shaded annuli are planar
sections of the supports of $\nabla\chi_\rho$ and $\nabla\chi_R$.
The radial profiles agree near $x_0$ and near both spectators. Thus
$\chi_R-\chi_\rho$ is supported away from all three insertions, and
conservation gives the same flux on the two shells.}
\label{ch14:fig-ward-shell}
\end{figure}


The actual two-mark estimate needed here has the following form.
If $A\in F_d$, let $\varkappa=0$ for constant $v$ and
$\varkappa=1$ for linear $v$. Put $D=d+1-\varkappa$.
For a finite family of separated spectator correlation functionals $\ell_i$ on
$F_D$ and coefficients $g_i$ annihilating that space,
there must be $\rho_m\downarrow0$ such that
\begin{equation}\label{ch14:annular-input}
 \sup_{1/4\le|z|\le2}
 \left|\sum_i g_i\rho_m^{3+\varkappa}
       S(J_{\mu\nu}(x_0+\rho_m z),A(x_0),\boldsymbol B_i)\right|
                  \longrightarrow0.
\end{equation}
Here the same correction coefficient at the distinguished insertion
must work for every spectator in the finite family.
The power $3+\varkappa$ is the four-dimensional volume minus one
cutoff derivative plus the degree of the Killing field.

The mixed local replacement estimate of
Proposition~\ref{ch18:mixed-bound} supplies this assertion directly,
without using phase decomposition or rotational Ward identities.
Apply it to the two-field cluster $(J_{\mu\nu},A)$, of total degree
$d+4$, and truncate at $D=d+1-\varkappa$. Its replacement field
$C_D(\rho,z)\in F_D$ is chosen by fixed probes before the exterior
spectator. For every spectator in the finite family the residual is
bounded uniformly on $1/4\le|z|\le2$ by
$C\rho^{D+1-(d+4)-\epsilon}$. The same coefficients therefore cancel
under $\sum_i g_i\ell_i=0$. After multiplying by the shell factor,
\[
 \rho^{3+\varkappa}
 \left|\sum_i g_iS(J_{\mu\nu}(x_0+\rho z),A(x_0),\boldsymbol B_i)\right|
 \le C\rho^{1-\epsilon}\longrightarrow0,
 \qquad 0<\epsilon<1.
\]
This proves \eqref{ch14:annular-input} for every sufficiently small
radius, and hence along any sequence $\rho_m\downarrow0$.
The full compact annulus is allowed because the local replacement
proof already removed its lattice-grid restriction by derivative
bounds. Only point collisions are excluded; coordinate ties cause
no additional restriction. This use of local replacement precedes
its phase-uniform operator-product formulation.

\begin{proposition}[Representation of the local flux]
\label{ch14:charges}
There are unique linear maps
$p_b:F_d\to F_{d+1}$ and
$Y_d(K):F_d\to F_d$ such that the constant and
linear fluxes \eqref{ch14:flux} equal the correlations with
$p_bA$ and $Y_d(K)A$ at $x_0$, respectively.
\end{proposition}

\begin{proof}
Every finite linear dependence among the spectator functionals on
$F_D$ annihilates the flux: shrink the shell, rescale its
integration variable and use \eqref{ch14:annular-input}. Choose a
finite separating subfamily of spectator functionals. Its injective
map from $F_D$ has finite-dimensional range; the preceding
annihilator property says that the corresponding flux vector lies
in that range. It therefore specifies a unique field, and adjoining
any other spectator correlation functional proves the identity for it as well.
Continuity extends the identity from generic grid configurations to
every separated configuration. Linearity follows from uniqueness.
\end{proof}

Differentiating the shell identity with the distinguished insertion,
using conservation to move the cutoff and symmetry to remove the
Killing derivative, gives
\begin{equation}\label{ch14:charge-commute}
 Y_{d+1}(K)\partial_bA
   =\partial_bY_d(K)A+\sum_cK_{cb}p_cA,
 \qquad \partial_bp_cA=p_c\partial_bA.
\end{equation}
Equivalently one can differentiate after moving the whole local
shell with $x_0$; the derivative of $K(y-x_0)$ supplies the displayed
constant flux term. Faithfulness of separated spectator correlation functionals turns
these response identities into field identities. The same argument
shows compatibility with inclusions, reality and signed coordinate
symmetries; both charges annihilate the unit.

\section{Coordinate spectral measures and the angular boundary}
\label{ch14:spectral}

We recall that $P_\mu$ in \eqref{ch14:coordinate-spectrum} is the
joint measure of nonnegative heat energy and transverse momentum.
For a point field, translate its insertion by $ue_\mu$, $u>0$,
into the positive half-space and define the undamped measure
\[
 d\nu_\mu^A(E,p)=e^{2uE}
       d\langle[A(ue_\mu)],P_\mu(E,p)[A(ue_\mu)]\rangle.
\]
The semigroup law makes this independent of $u$. It is finite on
$E\le M$; it need not have finite total mass. Word measures are
defined after translating the entire word beyond its extreme point.
Polarization defines mixed measures. For integrable nonnegative
$g_1,g_2$ and $|f|^2\le g_1g_2$, spectral Cauchy--Schwarz gives
\begin{equation}\label{ch14:spectral-CS}
 \left|\int f\,d\nu_\mu^{AB}\right|^2
       \le\int g_1\,d\nu_\mu^A\int g_2\,d\nu_\mu^B.
\end{equation}
Positive damping and the separated derivative bounds imply all
polynomial transverse moments. One direct justification takes
centered differences of the reflected two-point function: their
spectral multiplier is $(2\sin(tp_j/2)/t)^{2r}$. Fatou's lemma
and the integral remainder for finite differences give the moment
bound as $t\to0$. Further heat damping gives energy moments.
For point fields, integration over small damping times against a
sufficiently high power of time also gives
\begin{equation}\label{ch14:polynomial-moments}
       \int (1+E)^{-D}|p_j|^{2r}\,d\nu_\mu^A<\infty
\end{equation}
for some $D$ depending on $A,r$.

\subsection{What opposite reflection signs control}

If a field or word combination $A$ has reflection signs
$\varepsilon_i=\pm1$ in axes in a set $I$, its two-point kernel
has, in time direction $i$, the spectral sign $\varepsilon_i$.
For an even frequency profile $f$ use
\[
 \mathcal Lf(E)=\int e^{-E|t|}\left(\int e^{ipt}f(p)\,dp\right)dt
       =2E\int\frac{f(p)}{p^2+E^2}\,dp.
\]
At $E=0$ the limiting value is $2\pi f(0)$.
To compare representations in different axes without an arbitrary
collision extension, subtract the low moments of the profiles.
Choose even $\sigma\ge0$ of integral one supported in
$1\le|v|\le2$, and the Lagrange polynomials $L_j$ for the nodes
$(j+2)^2$, $0\le j<m$. Define the annular correction by
\[
 \int (\mathcal C_{m,\rho}f)(q)g(q)\,dq
 =\int\sigma(v)f(p)\sum_{j<m}L_j((p/(\rho v))^2)
                           g(\rho v(j+2))\,dp\,dv.
\]
Then $f-\mathcal C_{m,\rho}f$ has its first $2m$ moments zero.
Interpolation of $(p^2+E^2)^{-1}$ gives its exact response
\begin{equation}\label{ch14:positive-residual}
 \mathcal R_{m,\rho}f(E)
 =2E\int\frac{f(p)}{p^2+E^2}\int\sigma(v)
       \prod_{j<m}\frac{(\rho v(j+2))^2-p^2}
                          {(\rho v(j+2))^2+E^2}\,dv\,dp.
\end{equation}
For $f\ge0$ supported in $[-\rho,\rho]$ this is nonnegative.
If $f_w(p)=w^{-1}f(p/w)$ with $f=1$ near zero and
$0<w\le1\le\rho$, it is bounded below by
$c_m/(E+w)$ on $0\le E\le1$. For large $E$ its decay can be made
arbitrarily high by increasing $m$. The upper bound, unlike the
low-energy lower bound, holds also for growing $w$: it requires only
$\rho\ge1$ and $\operatorname{supp}f_w\subset[-\rho,\rho]$.
For $E\le1$ use $\|f_w\|_\infty\le C/w$; for $E\ge1$ the
product in \eqref{ch14:positive-residual} is bounded by
$C_m\min(1,(\rho/E)^{2m})$. Thus
$\mathcal R_{m,\rho}f_w(E)\le C(w^{-1}+\rho^D)(1+E)^{-D}$
for any fixed $D\le2m$. Together with
\eqref{ch14:polynomial-moments} this gives, for fixed $r$,
\[
 \int\mathcal R_{m,\rho}f_w(E)|p_j|^{2r}\,d\nu_i^A
                                      \le C(w^{-1}+\rho^D).
\]
Choose $D$ from the damped moment first, and then $m$ sufficiently
large; $D$ does not grow with the later subtraction order.

We now give the precise comparison. For a word combination $A$, assume
its offsets have zero coordinates in $I$ and that reflection in axis $i$
acts by $\varepsilon_i$ for $i\in I$. These geometric restrictions permit
arbitrarily small positive placements in any of the chosen time axes,
so the integrated moment bound \eqref{ch14:polynomial-moments} holds
in each. They will be relaxed for the line axis below.
For even nonnegative fixed profiles $f_j\in C_c^\infty(\RR)$ write
$f_{j,v}(p)=v^{-1}f_j(p/v)$; in particular the normalization preserves
its integral, while its supremum costs $v^{-1}$. For a product of profiles
$g=(g_k)_{k\in I}$ define
\[
 \operatorname{Resp}_i(g)
   =\int\mathcal Lg_i(E)\prod_{k\in I\setminus\{i\}}g_k(p_k)
                                                     \,d\nu_i^A(E,p).
\]
The same even product test function of the unreflected pair kernel has spectral
values $\varepsilon_i\operatorname{Resp}_i(g)$ in every axis $i$.
First require sufficiently many vanishing moments in two axes, so the
test function vanishes to the required order at the collision. This makes the
comparison an equality of convergent integrals of the original kernel.
If $I$ omits axes, insert in each omitted frequency coordinate a profile
$h(p/T)$ with $h=1$ near zero, and let $T\to\infty$; the damped moment
bounds justify this limit. Thus opposite signs give
$\operatorname{Resp}_i(g)=-\operatorname{Resp}_j(g)$.

Let $a_i=\mathcal C_{m,\rho_i}g_i$ and
$a_j=\mathcal C_{m,\rho_j}g_j$, where each radius contains the support
of its profile and $m$ exceeds the necessary collision orders.
Apply the signed equality to the product with residual profiles in
both $i$ and $j$. Expanding it gives
\begin{align}
 \operatorname{Resp}_i(g_i-a_i;g_{\ne i})
 &\le \operatorname{Resp}_i(g_i-a_i;g_j\mapsto a_j)
       +\operatorname{Resp}_j(g_j-a_j;g_i\mapsto a_i).
 \label{ch14:sign-transfer}
\end{align}
Here the semicolon leaves every unmentioned profile unchanged.
The omitted term is
$-\operatorname{Resp}_j(g_j-a_j;g_{\ne j})\le0$, because its temporal
residual response and all its transverse profiles are nonnegative.
The corrections $a_i,a_j$ themselves can have signs. They occur only
as transverse profiles on the right, so replacing them by nonnegative
majorants gives an upper bound. This explains precisely where positivity
is used.

We claim that when both signs occur in $I$, for $\lambda\in e^{\NN}$,
$q<0$, $\beta_j\le0$, and every $\delta>0$,
\begin{equation}\label{ch14:mixed-sign}
 \int_{E\le1}\frac{\prod_{j\in I\setminus\{i\}}
                         f_{j,\lambda^{\beta_j}}(p_j)}
                        {E+\lambda^q}\,d\nu_i^A
                         =O_\delta(\lambda^\delta).
\end{equation}
The following induction proves the claim while keeping the order of
parameter choices explicit. Choose a walk $i_0,\ldots,i_N$ through $I$
with opposite signs at consecutive entries, starting at $i_0=i$, with
every axis visited before its final entry. Put
$V_s=\{i_0,\ldots,i_{s-1}\}$. Fix $K\ge0$ bounding the magnitudes of
all negative width exponents. Induct backwards on $s$ to find $a_s>0$
such that
\begin{equation}
 \int\mathcal R_{m,R_0\lambda^\eta}
            f_{i_s,\lambda^{\omega_{i_s}}}(E)
          \prod_{j\ne i_s}f_{j,\lambda^{\omega_j}}(p_j)\,d\nu_{i_s}^A
                       \le C\lambda^\delta.
 \label{ch14:sign-induction}
\end{equation}
Here $0\le\omega_j\le a_s$ on $V_s$,
$-K\le\omega_j\le a_s$ elsewhere,
$0\le\eta\le a_s$ and $\omega_{i_s}\le\eta$.
The fixed $R_0\ge1$ contains the support at the current radius;
$m$ may be any sufficiently large fixed subtraction order.
Products here and below range over $I$.

At $s=N$ all widths are at least one, hence the transverse suprema
are bounded. The residual moment estimate with moment order zero
bounds the left side by $C(\lambda^{-\omega_{i_N}}+\lambda^{\eta D})$.
Choose $a_ND<\delta$. For the step from $s+1$ to $s$ write
$i=i_s$, $j=i_{s+1}$ and $p_*=K|I\setminus\{i,j\}|$.
Choose, in this order,
\[
 \alpha=a_{s+1},\qquad
 2r\alpha>p_*+\max(K,\delta)+1,\qquad
 D=D(A,r),\qquad
 0<a_s<\min\{\alpha,\delta/(D+1)\}.
\]
Only after these choices enlarge $m$ to handle the moment exponent $D$
and the collision orders. Correct $g_i$ at radius $R_0\lambda^\eta$
and $g_j$ at a fixed multiple of $\lambda^\alpha$.
The correction in $i$ has a fixed nonnegative annular majorant of
width $\lambda^\eta$, since its coefficients involve only factors
$\lambda^{2r'(\omega_i-\eta)}\le1$. The second term of
\eqref{ch14:sign-transfer} is therefore covered by the next induction
step, with the previously current axis now having nonnegative width
exponent. The correction in $j$ is bounded by
$C_r\lambda^{-2r\alpha}|p_j|^{2r}$ on its annulus; its coefficients
are bounded because $\omega_j<\alpha$. The remaining narrow suprema
cost at most $\lambda^{p_*}$. The first term is consequently at most
\[
 C\lambda^{p_*-2r\alpha+
                    \max(-\omega_i,\eta D)}
 \le C\lambda^{p_*-2r\alpha+\max(K,a_sD)}
                  =O(\lambda^\delta).
\]
This proves \eqref{ch14:sign-induction}. Finally take $s=0$,
$\omega_i=q$, $\omega_j=\beta_j$, and $\eta=0$.
The lower bound $\mathcal R_{m,R_0}f_{i,\lambda^q}(E)
\ge c_m/(E+\lambda^q)$ on $E\le1$, with a bump $f_i=1$ near zero,
gives \eqref{ch14:mixed-sign}. No continuous rotation was used.

\subsection{The separate diagonal-stress estimate}

For $J_{\mu b}$ with $\mu\ne b$, the two odd tensor indices give
opposite signs and \eqref{ch14:mixed-sign} applies. A diagonal
stress component has only positive signs, so it requires conservation.
We describe the contact subtraction needed just for this argument;
it is not a positive extension of the correlation hierarchy.

Let $K_{ab,cd}(x)=S(J_{ab}(x),J_{cd}(0))$ for $x\ne0$, and choose
an integer $D$ above the collision orders of these kernels and their
first derivatives. For one-dimensional Schwartz test functions let $\Pi_D$
retain Taylor orders less than $D$, multiplied by a fixed even
cutoff. On product test functions define
\[
 R_{ab,cd}(h)=\int K_{ab,cd}(x)
       \left(\prod_i h_i(x_i)-\prod_i\Pi_Dh_i(x_i)\right)dx.
\]
Each term of the telescoping product difference has a factor
$O(|x_i|^D)$, ensuring integrability. Using the same cutoff in every
axis preserves signed tensor symmetries and exchange of the index
pairs. Conservation makes its divergence a distribution supported at
zero with finitely many derivatives. Inversion removes its even orders.
Here is an explicit subtraction of the possible linear orders. Put
$\chi_0(x)=\prod_j\chi(x_j)$, using the even Taylor cutoff above, and set
\[
 q_{ab,cd}=\sum_i R_{ib,cd}\bigl(\partial_i(x_a\chi_0)\bigr),
 \qquad R'_{ab,cd}=R_{ab,cd}-q_{ab,cd}\delta_0.
\]
Indeed $\sum_i\partial_iR_{ib,cd}$ evaluated on $x_a\chi_0$ equals
$-q_{ab,cd}$, whereas the divergence of $q_{ib,cd}\delta_0$ has the
same value. Thus their difference has no linear derivative of
$\delta_0$. To check the tensor symmetries, reflection parity makes
$q_{ab,cd}$ vanish unless each coordinate index occurs an even number
of times. The remaining index patterns are $aa,aa$, $aa,bb$ and
$ab,ab$ (including their within-pair reversals); coordinate permutations,
within-pair symmetry of $R$, and exchange of its two pairs give the
same symmetries for $q$. The subtraction is therefore compatible with
all the tensor identities. We henceforth rename $R'$ as $R$.
Its Fourier divergence polynomial starts in degree three.

Write $\widehat R$ for this corrected functional and
$F_{b,cd}=\sum_i p_i\widehat R_{ib,cd}$. For a chosen time axis
$\mu$, set
\begin{align*}
 Q_0&=\widehat R_{\mu\mu,\mu\mu},\\
 Q_1&=\sum_{j,k\ne\mu}p_jp_k\widehat R_{k\mu,j\mu},\\
 Q_2&=\sum_{j,k,r,s\ne\mu}p_jp_kp_rp_s\widehat R_{sk,rj}.
\end{align*}
Conservation on the two index pairs gives
$Q_1=p_\mu^2Q_0+D_1$ and $Q_2=p_\mu^2Q_1+D_2$, where the
polynomials $D_1,D_2$ start in degrees four and six. Sufficiently
many temporal moment subtractions annihilate them exactly.
On a fixed temporal annulus the transverse responses of $Q_0,Q_1,Q_2$
are bounded by $O_\delta(\lambda^\delta)$ times respectively
$1,W^2,W^4$, if the transverse widths are at most $W$. For a diagonal
entry, divide its annular temporal profile by $p_\mu$ and put that
derivative on the label; the resulting vector has mixed signs.
Odd profiles are corrected by
$p^{-1}\mathcal C_{m,\rho}(pf)$, and their residual Laplace absolute
value is bounded by $C(E+w)^{-1}(1+E)^{-2m}$.
Thus \eqref{ch14:mixed-sign} and spectral Cauchy--Schwarz apply.
Removing one transverse moment restriction at a time terminates
after three axes. The annular coefficients remain bounded because
all shrinking width exponents are nonpositive.

For a fixed even annular $\sigma$, evaluate $Q_r$ against
$g_{v,r}(p)=v^{-1}(p/v)^{-2r}\sigma(p/v)$, so
$p^2g_{v,r+1}=v^2g_{v,r}$. Compare the widths $W$ and $2W$ and
subtract their fixed annular corrections. Write $c_r(t)=W^{-2r}\operatorname{Re}Q_r(g_{tW,r})$, with the
transverse profiles fixed, and let $d_r$ be the same normalized real
response to $(I-\mathcal C_{m,8})(g_{W,r}-g_{2W,r})$.
Conservation gives $c_{r+1}(t)=t^2c_r(t)+O_\delta(\lambda^\delta W^2)$
for $t=1,2$. The correction has zero constant moment and higher even
moments $O(W^2)$, so
$d_r=c_r(1)-c_r(2)+O_\delta(\lambda^\delta W^2)$. Eliminating
$c_0(1),c_0(2)$ from these three relations gives
\[
                  4d_0=5d_1-d_2+O_\delta(\lambda^\delta W^2).
\]
The corrected dyadic profile is signed, so its positivity needs
one further observation. For fixed correction nodes
$A_j=(8v(j+2))^2$, the function
\[
 R_E(z)=\frac1{z+E^2}\prod_{j<m}\frac{A_j-z}{A_j+E^2}
\]
is positive and decreasing for $0\le z\le16$, since $|v|\ge1$
and $A_j\ge256$. The dyadic response is the integral of
$2E[R_E(W^2p^2)-R_E(4W^2p^2)]$ against the nonnegative
annular profile $p^{-2r}\sigma(p)$ and the correction profile.
For $W\le1$ it is therefore nonnegative. On $E\le1$ the
product factors have a fixed positive lower bound; retaining the
difference of the two resolvents gives the lower bound
$cEW^2/(E+W)^4$ for $r=0$.
Conservation with the sufficiently many zero moments now gives
$d_0\ge0$, $d_1\le0$, $d_2\ge0$: each temporal multiplier
$p_\mu^2$ becomes $-E^2$ in the Laplace response, with the
contact terms killed exactly. Therefore, for
a nonnegative transverse product profile $\mathcal K_\lambda$,
\begin{equation}\label{ch14:diagonal-mass}
 \int_{E\le1}\frac{EW^4}{(E+W)^4}\mathcal K_\lambda(p)
                         \,d\nu_\mu^{J_{\mu\mu}}
                  =O_\delta(\lambda^\delta W^4).
\end{equation}
The integrand vanishes at $E=0$. The zero-energy part must be kept
separate in the subsequent Cauchy--Schwarz estimate.

\subsection{Vanishing angular flux at infinity}
\label{ch14:collinear}

For $K=E_{ba}-E_{ab}$ and spectator positions on a coordinate line
$\RR e_l$, choose anisotropic exponents
$\omega_a=\omega_b=-100$, $\omega_l=-150$ when $l\notin\{a,b\}$,
and $\omega_j=-300$ on remaining axes. Let
$w_j=\lambda^{1+\omega_j}$,
$\mathcal K_\lambda=\prod_{j\ne\mu}w_j^{-1}\psi(p_j/w_j)$,
with fixed even nonnegative $\psi=1$ near zero, and put
\[
 I_\mu(A;q)=\int_{E\le1}\frac{\mathcal K_\lambda(p)}
                                  {E+\lambda^q}\,d\nu_\mu^A.
\]
Project the spectator word to odd parity in $a,b$ and even parity
in the other axes, calling the result $B$. For a reflected line
axis this is a finite combination of words on the same line.
The mixed-sign estimate gives
\[
 I_\mu(B;-3000)=O(\lambda^{\gamma_\mu}),\quad
 I_\mu(J_{\mu b};\omega_\mu)=O_\delta(\lambda^\delta)\quad(\mu\ne b),
 \quad\gamma_l=1,\quad\gamma_\mu\le150.
\]
For the line time axis this requires a genuine temporal gap.
Choose a nonnegative real even smooth profile $\sigma_0$, let
$c=\sigma_0*\sigma_0$, and normalize
$\Phi=\widehat c=(\widehat\sigma_0)^2$ so $\Phi(0)=1$.
Thus $\Phi\ge0$. If translating the word by $u$ puts it in the
positive half-space, choose an even cutoff $\zeta=1$ on
$[-2u-1,2u+1]$, supported in $[-2u-2,2u+2]$.
Remove from $c_w$ the profile with Fourier transform
$\zeta(t)\Phi(wt)$. Its inverse-Fourier supremum is bounded
uniformly in $w$, whereas the residual has positive response
\[
 \int e^{-E|t|}(1-\zeta(t))\Phi(wt)\,dt
                          \ge c_u/(E+w)\quad(0\le E\le1).
\]
For the lower bound integrate on
$R<t<R+\gamma/(E+w)$, with $R=2u+2$, $wR\le\gamma$ and
$\Phi\ge1/2$ on $[-3\gamma,3\gamma]$.
Apply the sign-transfer inequality using this residual and an
opposite-sign transverse axis. Its correction term is bounded by
the three-axis mixed-sign estimate and the uniform supremum above;
its other term is bounded using a sufficiently high transverse moment,
now with fixed damping $e^{-2uE}$. The parity pattern has both signs
among the transverse axes and a sign opposite to the line axis,
so this proves any positive exponent, in particular $\gamma_l=1$.
For a different time axis use the three transverse-to-line axes
and bound the omitted line profile by its supremum. With width
$\lambda^{1+\omega_l}$ this costs exponent
$-(1+\omega_l)+1=-\omega_l\le150$.
These arguments keep the word measures fixed and include their
zero-energy mass; an arbitrary frequency bump would not have the
positive temporal-gap response used here. The diagonal estimate is
\eqref{ch14:diagonal-mass} with $W=\lambda^{-99}$.
The crude bound $I_\mu(J_{\mu\mu};0)=O(\lambda^{900})$ follows
from finite mass on $E\le1$ and the profile suprema. Outside
$E\le1$ and $|p_j|\le w_j$, positive damping and Schwartz decay
make the relevant test functions $O(\lambda^{-M})$ for every fixed $M$.

Take even $\eta\in C_c^\infty(\RR)$ equal to one near zero,
put $\ell(t)=t\eta(t)$, and define
\[
 f_{ab}(y)=\lambda^{100}\ell(\lambda^{-100}y_a)
                \prod_{j\ne a}\ell'(\lambda^{\omega_j}y_j),
 \qquad V_b=f_{ab},\quad V_a=-f_{ba}.
\]
The other components of $V$ vanish. Near all insertions $V(y)=Ky$.
Conservation and the local shell identities express the total raw
angular response as $-\int\partial_\mu V_c S(J_{\mu c},\boldsymbol A)$.
The terms $\partial_af_{ab}J_{ab}$ and
$\partial_bf_{ba}J_{ba}$ cancel exactly. In the remaining terms,
reflection projects the spectators to $B$. Fourier scaling supplies
$\lambda^3\mathcal K_\lambda$ and the temporal derivative supplies
$e^{-E/\lambda^{\omega_\mu}}$.
Apply \eqref{ch14:spectral-CS}. For the diagonal term split at
$E=\lambda^{-3000}$, using \eqref{ch14:diagonal-mass} above this
value and the crude bound below it. The worst resulting exponents,
including the prefactor $\lambda^{100}$, are
\[
\begin{array}{c|c}
 \mu=l\notin\{a,b\}&100+3-150+(1+1)/2=-46\\
 \mu\notin\{a,b,l\}&100+3-300+(150+1)/2=-121.5\\
 \mu=b,\ E\ge\lambda^{-3000}&100+3+(150+1-396)/2=-19.5\\
 \mu=b,\ E<\lambda^{-3000}&100+3-1500+(150+900)/2=-872.
\end{array}
\]
The exchanged column has the same bounds. Every term tends to zero,
whereas the total shell response is independent of $\lambda$.
Consequently
\begin{equation}\label{ch14:line-ward}
 \sum_i\left\{S(\ldots,Y(K)A_i,\ldots)
       +(Kx_i)_bS(\ldots,p_bA_i,\ldots)\right\}=0
\end{equation}
for every separated configuration on a coordinate line.
This argument used coordinate positivity and conservation, not a
cone spectrum or a pre-existing rotation action.

\section{Normalization and propagation of the Ward identities}
\label{ch14:propagation}

Translate a collinear configuration in its line direction in
\eqref{ch14:line-ward} and subtract the original identity. With a
plane generator mixing that line and a transverse direction one
obtains the translation Ward sum for $p_b$. The commutation law
\eqref{ch14:charge-commute} gives its positional derivatives as well.
The usual positional derivatives $\partial_b$ obey the same identities.

The finite vector correlation functionals of Theorem~\ref{ch03:vector-tests} have a
one-dimensional common kernel in the free vector-map space, spanned
by $\partial^0$. Adjoin a linear functional nonzero on that vector;
the matrix of correlation pairings becomes injective. Smooth comparison transports
this injectivity to the mixed finite-dimensional vector-map space.
Thus its Ward kernel has dimension at most one. It contains both
$p$ and $\partial$, so $p=c_d\partial$ at every sufficiently large
degree. Inclusions and nonzero free differentiation identify these
constants across degrees and extend the identity to the lower ones.
Call the common constant $c_*$. Changing variables in the shell
integrals and using \eqref{ch14:balanced-stress} gives
\begin{equation}\label{ch14:charge-limits}
 \frac{r_k}{c_k}U_{k,d+1}^{-1}p_bU_{k,d}\to\partial_b^0,
 \qquad
 \frac1{c_k}U_{k,d}^{-1}Y_d(K)U_{k,d}\to L_d^0(K).
\end{equation}
A finite separating family of correlation pairings justify these operator limits, even if the
operator norms were not known in advance. Since
$r_kU_{k,d+1}^{-1}\partial_bU_{k,d}\to\partial_b^0\ne0$
in some degree, the first limit gives $c_*\ne0$ and $c_k\to c_*$.
In particular $c_*>0$, since $c_k>0$.
Put $L_d(K)=c_*^{-1}Y_d(K)$. Then
\begin{equation}\label{ch14:derivative-commutator}
 L_{d+1}(K)\partial_b=\partial_bL_d(K)+\sum_cK_{cb}\partial_c.
\end{equation}
Every positional derivative of the normalized angular Ward sum is
therefore another such sum with a differentiated label.

\subsection{Analyticity supplied by a positive Gram kernel}

The missing step from a line to an open set is analyticity, not
smoothness. We establish exactly the analyticity needed.
For $x^a>0$ let $\Psi_A(x)=[A(x)]_a$ in the coordinate reflection
space. The curve $t\mapsto\Psi_A(x+te_b)$ is smooth for $b\ne a$:
its weak derivative against every word vector is
$\Psi_{\partial_b A}(x+te_b)$, which is continuous by translation
continuity. Integrating this continuous vector curve and using the
dense word span proves strong differentiability. Iterate for higher
orders. In the normal direction the curve is analytic away from the
reflection plane, since it is a positive-time heat orbit.

Let $F_{A,B}(x)=S(A(x),B(0))$. Its restriction to $ue_a$, $u>0$,
is analytic by that heat argument. In the plane $(a,b)$ with
$Ke_a=e_b$, $Ke_b=-e_a$, define
\[
 Q_{A,B}(u,z)=F_{e^{-\theta L(K)}A,e^{-\theta L(K)}B}(\rho e_a),
 \quad \rho=(u^2+z^2)^{1/2},\quad\theta=\arctan(z/u).
\]
This is analytic for $u>0$: the label spaces are finite-dimensional
and the radial functions are analytic. It solves the angular
identity because
$(u\partial_z-z\partial_u)\theta=1$ and the same operator kills
$\rho$. It agrees with the actual $F$ at $z=0$.
All transverse jets of the actual angular expression vanish there,
by \eqref{ch14:line-ward} and
\eqref{ch14:derivative-commutator}. For $H=F-Q$, differentiating
the angular expression $n$ times in $z$ gives at $z=0$
\[
 u\partial_z^{n+1}H_{A,B}
 =-\partial_z^n(H_{L(K)A,B}+H_{A,L(K)B})
                        +n\partial_u\partial_z^{n-1}H_{A,B}.
\]
The last term is omitted at $n=0$. Induction proves equality of all
transverse jets of $F$ and $Q$, simultaneously for all labels.

Equality of jets has not yet proved equality of functions. The
positive Gram identity supplies the decisive bound. If a smooth
Hilbert-valued curve $v(t)$ satisfies
$\langle v(s),v(t)\rangle=A(t-s)$ and the derivatives of $A$ at
zero agree with an analytic function, then
\[
 \|v^{(n)}(t)\|^2=|A^{(2n)}(0)|
 \le C R^{-2n}(2n)!
 \le C(2/R)^{2n}(n!)^2.
\]
Hilbert-space Taylor remainders tend to zero for $|h|<R/2$,
uniformly in $t$. Apply this to the single-insertion curve above: its Gram
kernel is $F_{A,A^{\dagger_a}}(2x^ae_a+(t-s)e_b)$, whose jets
are those of $Q$. Thus single-insertion vectors are analytic along every
transverse coordinate as well.

If one insertion is a strict maximum in some coordinate, separate
it from the spectator word by a coordinate hyperplane. Its moment
is a matrix element of this analytic single-insertion curve against a
fixed word vector. Strict minima are handled by the existing signed
reflection. Hence the Ward sum is analytic along any coordinate
segment on which the moving point stays a strict extremum in some
coordinate. This uses no integration of the spin operators into a
rotation group.

\subsection{A finite path through separated configurations}

Configurations with distinct first coordinates are dense. Starting
from one, choose the point with smallest first coordinate among
those not yet moved. Move it in the second coordinate until it
exceeds all other second coordinates; it remains strictly first in
the first coordinate. Next move it in the first coordinate beyond
all other points, remaining extreme in the second. Finally set its
other coordinates to zero, remaining extreme in the first. Repeat
with the remaining points. Every segment is separated and of the
analytic kind just proved. The result lies on the first coordinate
axis. There all Ward sums and all their derivatives vanish.
Reverse the path: the one-variable identity theorem on each segment
propagates the identity and its derivative identities. Continuity
then gives, on every separated configuration,
\begin{equation}\label{ch14:global-ward}
 \sum_i\left\{S(\ldots,L(K)A_i,\ldots;x)
                  +(Kx_i)_b\partial_{x_i^b}S(\boldsymbol A;x)\right\}=0.
\end{equation}
This is why replacing analyticity by smoothness in the propagation
argument would leave a gap.

\section{The rotation action and the normalized tensor}
\label{ch14:rotations-section}

Commuting two identities \eqref{ch14:global-ward} shows that the
internal defect
$[L(K),L(M)]-L([K,M])$ has zero correlation sums. It is real,
alternating, filtration preserving, kills the unit, and has the
signed covariance of a pair of rotation generators. The finite
internal-defect correlation functionals of Theorem~\ref{ch03:internal-tests} are
injective on this defect space. Their injectivity transfers under
the smooth finite-dimensional comparison, exactly as for the vector
correlation functionals above. The defect is therefore zero. Injective inclusions
extend this conclusion from any sufficiently large degree to the
lower ones.

One must still exclude a representation only of the universal cover
of $SO(4)$. Here the free comparison gives an exact intertwiner.
Choose the six plane generators and, at one degree, write their free
matrices as $C_i$ and their transported mixed matrices as $B_{k,i}$.
By \eqref{ch14:charge-limits} these converge to $C_i$. On the matrix
space put
\[
 \mathcal A_{k,i}(X)=B_{k,i}X-XC_i,
 \qquad \mathcal C_k=\sum_i\mathcal A_{k,i}^2.
\]
The bracket identities make $\mathcal C_k$ commute with every
$\mathcal A_{k,i}$. At the limit an invariant inner product makes
$\mathcal C$ self-adjoint and negative semidefinite, with kernel
the commutant of the free action. Its zero eigenvalue is isolated.
Finite-dimensional spectral projections around that eigenvalue give
real projections $P_k\to P$ commuting with all $\mathcal A_{k,i}$.
They also commute with the signed symmetries on matrix space.
The restrictions of $\mathcal A_{k,i}$ to $\operatorname{ran}P_k$
tend to zero. Every plane generator is a commutator of two others,
so their maximal norm satisfies $m_k\le2m_k^2$. Eventually $m_k=0$.
Consequently $f_k=P_kI\to I$ is invertible and
$B_{k,i}f_k=f_kC_i$. The intertwiner $U_{k,d}f_k$ is real and
preserves the signed action. Conjugating the free $O(4)$ action by
it gives the required genuine orthogonal-group representation.
This finite-dimensional argument also fixes the periodicity of
plane rotations, rather than assuming it from Lie brackets.

\begin{theorem}[Euclidean rotations]\label{ch14:rotations}
Under the actual marked comparison and annular estimates used above,
there are continuous complex-linear representations $R_d$ of $O(4)$
on $F_d$, extending the signed coordinate symmetries,
preserving reality and the unit, and compatible with inclusions.
Their infinitesimal generators are $L_d$, their derivatives transform
as vectors, and
\[
 S(R_{d_i}(g)A_i;gx)=S(\boldsymbol A;x),\qquad g\in O(4).
\]
They are isomorphic in each degree to the free label representations.
\end{theorem}

\begin{proof}
The preceding construction supplies the representations. The function
\[
 t\longmapsto S(R(e^{tK})\boldsymbol A;e^{tK}x)
\]
has zero derivative by \eqref{ch14:global-ward}. Plane rotations and
signed reflections
generate $O(4)$. Covariance, the fact that inserting the identity leaves correlations unchanged,
and the exact derivative and inclusion identities show that the corresponding label identities
have zero responses against every spectator. Faithfulness of the
field quotient proves their equality. Thus the independently chosen
finite-degree intertwiners need not themselves commute with inclusions;
the resulting actions do.
\end{proof}

Let $J_0=c_*^{-1}J$. With probability Haar measure on $O(4)$ set
\begin{equation}\label{ch14:averaged-stress}
 T_{ab}=\int_{O(4)}\sum_{i,j}g_{ai}g_{bj}R_4(g)(J_0)_{ij}\,dg.
\end{equation}
\begin{proposition}[Normalized stress tensor]
\label{ch14:normalized-stress}
The tensor $T$ is real, centered, symmetric, conserved and nonzero.
It transforms as a tensor under $R_4$. Its constant and linear local
fluxes are respectively the positional derivative and $L_d(K)$.
Along the same comparison sequence,
$U_{k,4}^{-1}(r_k^4T)\to T^0$.
\end{proposition}

\begin{proof}
Haar invariance gives the tensor law. All other linear properties
pass through the integral, conservation using vector covariance of
derivatives. In a transformed shell change variables by $g$ and use
correlation covariance; its actions are $\partial_b$ and $L_d(K)$.
Averaging preserves these normalized actions. If $T=0$, all tested
positional derivatives would vanish, contradicting smooth comparison
with a free differentiated quadratic field. Finally $T$ belongs to
the same centered conserved line as $J$. Its translation action
forces $J=c_*T$, and \eqref{ch14:balanced-stress} with
$c_k\to c_*$ proves the comparison limit.
\end{proof}

The complexified finite-dimensional rotation representations also
specify the label representations for the later continuation to
Lorentz transformations. No Lorentz action on a physical Hilbert
space, uniqueness of the vacuum, or positive extension of collision
kernels has been constructed in this chapter.

\setcounter{chapter}{14}
% !TEX root = ../monograph.tex
\chapter{Decomposition into phases}
\label{ch15:phases}

A translation-invariant correlation family can be a mixture of states
with different long-distance behavior. We separate this mixture using
operators on the zero-energy subspace of its reflection Hilbert space.
A phase will be a character of a commutative operator algebra, not a
set of gauge configurations. Its correlations are obtained by evaluating
compressed whole words.

The inputs are the common field spaces $F_d$ and separated correlation
family of Chapter~\ref{ch13:continuum}, bounded slab operators with
their adjoint and concatenation laws, and orthogonal covariance of that
same family. The latter is stronger than the original lattice
hypercubic covariance. We also require the bounded real lattice
approximations and corrected moment estimates for separating families specified below.
The assertions retain these analytic hypotheses; a decomposition theorem
for an abstract state alone does not establish them.

\section{The zero-energy subspace}
\label{ch15:ground}

Use time coordinate four. Write $\mathcal H^{\rm OS}$ for the mixed
reflection Hilbert space, $\Omega$ for its unit vacuum,
$E(t)=e^{-tH}$ for the heat semigroup, and $U(y)$ for spatial
translations. Let $P_0$ be the projection onto
$\mathcal H_0=\ker H$. The spectral theorem gives
$E(t)\to P_0$ strongly as $t\to\infty$.

Every spatial translation acts trivially on $\mathcal H_0$.
Here the orthogonal covariance input is essential. For two separated
clusters put $F_{v,w}(x)=S(v\,\tau_xw)$. At sufficiently large time
separation this is a matrix element of $U(\vec x)E(x^4-c)$ between
fixed vectors. Hence a fixed temporal displacement changes it by a
quantity tending to zero, uniformly in arbitrary spatial displacements.
If $b^4>0$, choose an orthogonal map $K$ with $Kb=|b|e_4$.
Then $(Ke_4)^4=b^4/|b|>0$. Rotate the clusters and apply the temporal
statement along $tKe_4$. This proves
$F_{v,w}(te_4+b)-F_{v,w}(te_4)\to0$.
For spatial $y$ use $b=y+e_4$ and remove the extra unit time shift.
Taking strong heat limits and using the density of word vectors gives
\begin{equation}
                              U(y)P_0=P_0.
 \label{ch15:spatial-ground}
\end{equation}
This conclusion concerns the ground space of the mixture; it does not
say that this space is one dimensional.

\section{Compression and separation of clusters}
\label{ch15:compression}

For a word $w$ placed inside a slab of width $h$, let $B_h(w)$ be its
bounded insertion operator and define
\[
 M(w)=P_0B_h(w)|_{\mathcal H_0}.
\]
Increasing the width appends a heat operator, and shifting the word in
time inserts one on the other side. Both disappear after compression.
Spatial placement disappears by \eqref{ch15:spatial-ground}. Thus
$M(w)$ is defined for any distinct-point word independently of its
slab placement, with
\[
 M(\varnothing)=1,\qquad M(w)\Omega=P_0[w],\qquad
                  S(w)=\langle\Omega,M(w)\Omega\rangle.
\]
For positive-time words also put $Y_w=B_h(w)|_{\mathcal H_0}$,
which maps into the full reflection Hilbert space. The slab adjoint
and concatenation laws give
\begin{equation}
 M(\theta v\,w)=Y_v^*Y_w,
 \qquad M(\theta v)=M(v)^*.
 \label{ch15:whole-gram}
\end{equation}

If successive time gaps between translated words tend to infinity,
their correlation tends to the corresponding ordered product of
compressed operators:
\begin{equation}
 \lim_{j\to\infty}S(\tau_{z_1(j)}w_1\cdots\tau_{z_k(j)}w_k)
       =\langle\Omega,M(w_1)\cdots M(w_k)\Omega\rangle,
 \quad (z_{i+1}(j)-z_i(j))^4\to\infty.
 \label{ch15:cluster}
\end{equation}
The spatial parts of the displacements can be arbitrary. Indeed,
if $\psi_j\to\psi$ and $t_j\to\infty$, then
$U(y_j)E(t_j)\psi_j\to P_0\psi$ by contraction, strong heat
convergence and \eqref{ch15:spatial-ground}. Iterating this fact through
the bounded slab operators proves \eqref{ch15:cluster}.

\begin{proposition}[The compression algebra]
\label{ch15:algebra}
The norm-closed unital star algebra $\mathfrak A$ generated by all
$M(w)$ on $\mathcal H_0$ is commutative and separable. The vector
$\Omega$ is cyclic and separating for it.
\end{proposition}
\begin{proof}
We give the particular geometric commutativity argument. Place four
words at
\[
 0,\quad j(e_4+2e_1),\quad j(2e_4+e_1),\quad3j(e_4+e_1).
\]
Let $s$ exchange coordinates one and four. This reverses the time order
of the two middle clusters, as shown in Figure~\ref{ch15:fig-cluster-exchange}.
Covariance and \eqref{ch15:cluster} therefore exchange
the two middle operators, initially with that orthogonal map applied
to all four words. Apply the same cluster limit to four words placed
on the ray $j(e_4+e_1)$, which the exchange fixes. This removes the
common transformation of the labels and internal positions. Thus
\[
 \langle\Omega,M(w_1)[M(w_2),M(w_3)]M(w_4)\Omega\rangle=0.
\]
\begin{figure}[!htbp]
\centering
\begin{tikzpicture}[x=0.87cm,y=0.87cm,font=\small,
  axis/.style={->,>=Stealth,black!65,thin},
  centre/.style={circle,fill=ymInk,inner sep=0pt,minimum size=4.3pt},
  word/.style={fill=white,inner sep=1pt,outer sep=2pt}]
  \foreach \shift/\heading in {0/{Exchange of the middle clusters},6.5/{Centres on the fixed ray}}{
    \begin{scope}[xshift=\shift cm]
      \foreach \k in {1,2,3}{
        \draw[ymGrid,thin] (0,\k) -- (3.25,\k);
        \draw[ymGrid,thin] (\k,0) -- (\k,3.25);
      }
      \draw[axis] (-0.18,0) -- (3.8,0) node[right,black] {$x^1/j$};
      \draw[axis] (0,-0.18) -- (0,3.72) node[above,black] {$x^4/j$};
      \draw[ymInk,dashed] (0,0) -- (3.35,3.35);
      \node[centre] at (0,0) {}; \node[word,below right] at (0,0) {$w_1$};
      \node[centre] at (3,3) {}; \node[word,below right] at (3,3) {$w_4$};
      \node at (1.75,-0.92) {\heading};
    \end{scope}
  }
  \draw[<->,>=Stealth,ymAccent,thick,shorten <=3.5pt,shorten >=3.5pt]
    (2,1) -- (1,2);
  \node[centre] at (2,1) {}; \node[word,right] at (2,1) {$w_2$};
  \node[centre] at (1,2) {}; \node[word,left] at (1,2) {$w_3$};
  \begin{scope}[xshift=6.5cm]
    \node[centre] at (1,1) {}; \node[word,below right] at (1,1) {$w_3$};
    \node[centre] at (2,2) {}; \node[word,below right] at (2,2) {$w_2$};
  \end{scope}
\end{tikzpicture}
\caption[Exchanging two clusters in the phase algebra]{Each point gives the displacement of a whole word. Read the
  points from bottom to top for the original time order and from left to
  right for the time order after the coordinate exchange $s$. On the left
  the order changes from $w_1,w_2,w_3,w_4$ to $sw_1,sw_3,sw_2,sw_4$.
  On the right $s$ fixes all four centres. Here the cluster limit gives
  the same vacuum expectation with $s$ removed from each word in the
  reordered product.}
\label{ch15:fig-cluster-exchange}
\end{figure}

The vectors $M(v)\Omega=P_0[v]$ span a dense subspace of $\mathcal H_0$;
reflected words give the adjoints in the left slot. The bounded
commutator is therefore zero. Density proves cyclicity. If
$A\Omega=0$ for $A\in\mathfrak A$, commutativity gives $AB\Omega=0$
for all $B\in\mathfrak A$, hence $A=0$. Norm continuity of slab
operators, finite-dimensional label spaces, and countable dense
position sets give a countable generating family, hence separability.
\end{proof}

Compression is not multiplication of finite-separation words. In fact
\begin{equation}
 M(\theta v\,w)-M(v)^*M(w)=Y_v^*(1-P_0)Y_w.
 \label{ch15:compression-defect}
\end{equation}
The right side usually survives at finite separation. It is the
whole-word identity \eqref{ch15:whole-gram}, rather than a product of
scalar phase expectations, that will give phase reflection positivity.

\FloatBarrier
\section{Characters and the representing measure}
\label{ch15:characters}

Let $\Xi$ be the space of unital star characters of $\mathfrak A$.
The commutative Gelfand representation identifies
$\mathfrak A\simeq C(\Xi)$ isometrically; see the original
\cite[Section~2, Lemma~1]{GelfandNaimark1943}. The compact space $\Xi$
is metrizable because $\mathfrak A$ is separable. The positive normalized
vacuum functional has a unique representing probability measure
$\nu$ by the Riesz theorem; for this application see
\cite[Theorems~2.4.1 and~2.8.2]{Landsman1998}.
It has full support. Otherwise a nonzero continuous function supported
in a nonempty open null set would correspond to an $A\ne0$ with
$\|A\Omega\|^2=\int|\xi(A)|^2\,d\nu=0$, contradicting separation.

Define
\begin{equation}
 S^\xi(w)=\xi(M(w)),\qquad
 S(w)=\int_\Xi S^\xi(w)\,d\nu(\xi).
 \label{ch15:mixture}
\end{equation}
Every fixed separated value is continuous in $\xi$. This is a
representation of the mixed correlations by characters of the ground
algebra. Vacuum uniqueness in these phases will be proved separately.

We shall repeatedly use the following consequence of full support.
If $A\in\mathfrak A$ and
$\langle\Omega,(A^*A)^j\Omega\rangle\le C^{2j}$ for all $j$, then
\begin{equation}
                        \|A\|=\sup_{\xi\in\Xi}|\xi(A)|\le C.
 \label{ch15:moment-norm}
\end{equation}
This is the usual limit of $L^{2j}$ norms on a probability space,
combined with full support for a continuous function. In applications,
the moment bound comes by taking widely separated copies of the same
lattice expression, applying the chessboard bound at fixed copy number,
and then taking the ordered limits and the cluster limit.
The order of these operations matters.

\section{Symmetries in each phase}
\label{ch15:phase-symmetry}

Symmetry of a mixture need not imply symmetry of every component.
Here an additional lattice variance estimate gives the required
operator identity. Write $\vartheta$ for linear time reflection on
positions and labels, without conjugation. For a generic real grid word $w$ (distinct projections in each
coordinate),
set $Z=M(w)-M(\vartheta w)$. The adjoint law gives $Z^*=-Z$, so
\[
 a:=\langle\Omega,Z^2\Omega\rangle=-\|Z\Omega\|^2\le0.
\]
Corollary~\ref{ch13:bounded-word}, using the complete residual
estimate of Lemma~\ref{ch13:coarse-residual-supplier}, gives the
following bounded real approximation. For widely translated lattice copies $f_i$ of
$w-\vartheta w$, and every $\varepsilon>0$, there are real $g_i$
with $|g_i|\le R_\varepsilon$, uniformly in cutoff, volume and copy
position, such that
\[
 |\langle g_ig_j\rangle-\operatorname{Re}\langle f_if_j\rangle|
                                      \le\varepsilon\quad(i<j)
\]
on the ordered cutoff tails at fixed translation and fixed pair.
There is no corresponding assumed bound on $f_i^2$.
Variance positivity for $q+1$ copies gives
\[
 0\le\left\langle\left(\sum_{i=0}^qg_i\right)^2\right\rangle
 \le(q+1)R_\varepsilon^2+
              2\sum_{i<j}(\operatorname{Re}\langle f_if_j\rangle+\varepsilon).
\]
First take the ordered cutoff limits at fixed $q$ and translations,
then send their separation to infinity. Equation~\eqref{ch15:cluster}
gives $0\le(q+1)[R_\varepsilon^2+q(a+\varepsilon)]$.
Send $q\to\infty$ and then $\varepsilon\downarrow0$.
Thus $a=0$, and separation of $\Omega$ gives $Z=0$.
The cofinal grids are chosen after the degree and lower-level request.
Generic configurations are dense in the full separated configuration
space. Norm continuity and multilinearity
extend the identity to every word.

Orthogonal covariance of the mixture induces simultaneous invariance
of compressed product moments whenever a cluster-separating ray has
positive time before and after rotation. Such a ray exists unless
time is reversed; compose with $\vartheta$ in that case. On the dense
vectors $M(w)\Omega$ this gives a unitary orthogonal action, under
which time reflection now acts as the identity. Its conjugates, all
hyperplane reflections, therefore act trivially as well. Since these
generate $O(4)$,
\begin{equation}
                     M(Kw)=M(w)\qquad(K\in O(4)).
 \label{ch15:phase-invariant}
\end{equation}
Character evaluation proves symmetry in every phase. The proved lattice
approximation estimate is indispensable to this conclusion.

\section{Positivity, field identities and null spaces}
\label{ch15:fields}

For positive words $w_i$, \eqref{ch15:whole-gram} gives the positive
operator
\[
 \sum_{i,j}\overline{c_i}c_jM(\theta w_i\,w_j)
                   =\left(\sum_i c_iY_{w_i}\right)^*
                    \left(\sum_j c_jY_{w_j}\right)\ge0.
\]
It belongs to $\mathfrak A$. Positivity agrees with ambient operator
positivity in a unital C$^*$-subalgebra, and characters are positive;
see \cite[Section~2.6]{Landsman1998}. Hence each phase is reflection
positive, with a Gram form that need not have rank one.
Norm-smooth slab dependence gives smooth phase correlations and their
uniform separated bounds. Field inclusion, derivative and stress
identities pass to phases when written as operator identities, tested
against all exterior words before character evaluation. For a stress
shell the weights avoid collisions, so the norm-continuous operator
integral is defined. The expectation identity alone would not suffice.

A mixed null field has zero slab operator, by testing against the dense
word vectors. It is therefore null in every phase:
$N_d(S)\subset N_d(S^\xi)$. The reverse inclusion uses the quantitative separating families from the
same lattice family. Apply Proposition~\ref{ch11:n-mark-comparison}
with two marks per comparison cell. At a fixed sufficiently deep
comparison level $k$, its corrected coefficients for one and two insertions
have chessboard seminorm at most $B_k=H_k^{-1/32}$, uniformly in
cutoff and period. Choose a simultaneous coefficient limit for the
finitely many bounded lifts of the separating family at this fixed $k$, on a
refinement of the already selected correlation limits.
For every fixed integer $j$, put $2j$ copies of a corrected word
and its adjoint at widely separated translations. The chessboard
bound is $B_k^{2j}$: the two-mark estimate is used with arbitrary
replica census, not rerun with arity $2j$. Its constant is chosen
before $j$; only the cutoff and translation thresholds depend on
that fixed number. The repeated-word identity and the cluster
limit \eqref{ch15:cluster} identify the limiting product with
$\langle\Omega,(\Delta^*\Delta)^j\Omega\rangle$.
Taking successive deeper levels $k=k_\ell$ therefore supplies
compressed errors for one and two insertions $\Delta$ with
\begin{equation}
 \langle\Omega,(\Delta^*\Delta)^j\Omega\rangle
                  \le\gamma_\ell^{2j}\quad(j\ge1),
                  \qquad \gamma_\ell\longrightarrow0.
 \label{ch15:detector-moments}
\end{equation}
Here one may take $\gamma_\ell=B_{k_\ell}$, enlarging it by
a fixed finite row constant if necessary. The free row moments are
bounded, so $\gamma_\ell\to0$. These are estimates uniform in the number of reflected copies,
stronger than a bound on the mixed mean.
The finite reference rows use fixed free labels and convergent
noncollision configurations; their one-point means are uniformly bounded.
To make the pair correction explicit, if the reference one-point means are $s,s'$ and the reference connected pair is $c$, use
\[
 \Delta_1=M(u)-s1,\quad \Delta'_1=M(v)-s'1,\quad
 \Delta_2=M(uv)-sM(v)-s'M(u)+(ss'-c)1.
\]
Equation~\eqref{ch15:moment-norm} bounds each error by
$\gamma_\ell$ in every character. Therefore the phase pair differs
from $ss'+c$ by at most $\gamma_\ell(1+|s|+|s'|)$.
The phase matrices of correlation pairings converge uniformly in $\xi$ to the
injective free matrix of correlation pairings. Since they factor through
$V_d/N_d(S^\xi)$, its dimension is at least $\dim\mathcal M_d^{\rm Max}$.
The inclusion of null spaces gives the opposite bound through
Proposition~\ref{ch13:rank}. Consequently
\begin{equation}
                         N_d(S^\xi)=N_d(S)
                      \quad\hbox{for every }\xi.
 \label{ch15:same-fields}
\end{equation}
All phases have the same field filtration. Neither abstract character
decomposition nor positivity alone proves this equality.

\section{Selecting a phase with a unique ground vector}
\label{ch15:selection}

Construct the reflection Hilbert space and heat semigroup of each
phase using its separated hierarchy. Let $P_{0,\xi}$ denote its ground
projection. For a positive word $w$, the nonnegative heat matrix
elements decrease to $\|P_{0,\xi}[w]_\xi\|^2$ and are dominated by
the continuous zero-time Gram value. Dominated convergence and
\eqref{ch15:mixture} yield
\[
 \int\|P_{0,\xi}[w]_\xi\|^2\,d\nu=\|P_0[w]\|^2
                         =\int|S^\xi(w)|^2\,d\nu.
\]
Their difference is
$\|P_{0,\xi}[w]_\xi-S^\xi(w)\Omega_\xi\|^2$. It vanishes almost
everywhere. Rational positive configurations with basis labels form
one countable family dense in every phase reflection space, by Gram
continuity. Intersect its full-measure sets. Outside the resulting
null set,
\begin{equation}
                       \ker H_\xi=\CC\Omega_\xi.
 \label{ch15:ae-ground}
\end{equation}
Choose and retain one $\xi_*$ in this full-measure set. At this stage
we have not eliminated exceptional phases or shown that $\nu$ is a
point mass.

The slab bounds are uniform in all phases, including exceptional
ones. For the formal iterates $v_j=(A_{w^\dagger}A_w)^jv$, the mixed
maps satisfy
$Y_{v_j}=(B_h(w)^*B_h(w))^jY_v$. Thus
\[
 \|[v_j]_\xi\|^2=\xi(Y_{v_j}^*Y_{v_j})
                   \le\|Y_v\|^2\|B_h(w)\|^{4j}.
\]
The repeated-power argument gives bounded phase slab operators with
$\|B_h^\xi(w)\|\le\|B_h(w)\|$, and their identities and derivatives
follow on the dense word span. These uniform bounds are the input for
relativistic reconstruction and for the later removal of exceptional
phases. No spectral gap has been obtained here.

% !TEX root = ../monograph.tex
\section{Sharp separated bounds in every phase}
\label{ch15:sharp-phase-section}

Recall that \eqref{ch13:sharp-separated} bounds the undecomposed fixed-arity
correlation by the sum of degrees plus any positive loss in the exponent.
The estimate is also uniform in the phase $\xi$. To see this without
bounding phase evaluations by the undecomposed expectation of the same word,
apply the preceding chessboard argument to arbitrarily many separated
copies of that \emph{concatenated} word and their adjoints. The cell scales are
chosen from the classes modulo the tiling and the separations of one fixed word, not
from the total number of copies. Copies placed far apart occupy distinct
cells and have the same one-word factor on the right of the chessboard
inequality. Passing to the separated-cluster limit yields the same
bound to the corresponding power for the vacuum expectations of the powers of its compressed
operator $M(w)^*M(w)$. The norm formula of
\eqref{ch15:moment-norm} and the character bound then give
\begin{equation}
 \sup_{\xi\in\Xi}|S_n^\xi(\boldsymbol A;x)|
      \le C_{\boldsymbol A,n,\epsilon}r^{-\sum_i d_i-\epsilon}.
 \label{ch13:sharp-phase}
\end{equation}
Thus the fixed-arity scaling exponent is the sum of degrees plus an
arbitrarily small loss. No uniform constant as arity tends to infinity
is asserted; the repeated-word estimate remains the appropriate statement
for that different purpose.


\setcounter{chapter}{15}
% !TEX root = ../monograph.tex
\chapter{Relativistic reconstruction in every phase}
\label{ch16:reconstruction}

We now pass from separated Euclidean correlations to relativistic fields.
The order of the argument matters. Bounded slab insertions give analytic
functions of complex time and their joint tempered boundary values.
Lorentz covariance and spectral support then supply a full spacetime
tube, in which Euclidean permutation symmetry becomes locality. Only
after these steps do we form the Hilbert space of field words.

Throughout this chapter we use the conclusions of
Chapters~\ref{ch14:reflection} and~\ref{ch15:phases}: the coordinate
reflection spaces, exact Euclidean symmetries and common phase field
calculus. Write $\Xi$ for the compact phase space, $\nu$ for its
full-support probability measure, and $S^\xi$ for its correlations.
The finite-dimensional spaces $F_d$ have compatible inclusions; their
algebraic direct limit is denoted by $F$. All estimates below concern a
fixed finite list of labels. Uniformity means uniformity in $\xi$, not in
the number or degree of those labels.

The underlying reconstruction principles are those of Wightman and
Osterwalder--Schrader \cite{Wightman1956,OsterwalderSchrader1973,OsterwalderSchrader1975}.
Here the available Euclidean data are smooth functions off the diagonals,
with bounded slab operators. We therefore verify the boundary construction
directly; an unproved extension of these functions across Euclidean
collisions is not an input.

\section{Complex time and a uniform boundary distribution}
\label{ch16:boundary-section}

Use coordinate four as time. In the phase reflection space let
$E_\xi(s)=e^{-sH_\xi^{\rm OS}}$ and $U_\xi(y)$ be the heat and spatial
translation operators. The spectral theorem extends $E_\xi$ to
$\operatorname{Re}s>0$ as a holomorphic contraction. Let $B_\varepsilon^\xi(A)$
be the slab operator of width $2\varepsilon$ with the single label $A$
placed at its center. The phase version of the slab bound gives
\begin{equation}
 \|B_\varepsilon^\xi(A)\|\le C_A\varepsilon^{-q_A},\qquad
 B_{\varepsilon+u}^\xi(A)
       =E_\xi(u)B_\varepsilon^\xi(A)E_\xi(u),
 \quad 0<\varepsilon\le1,
 \label{ch16:slabs}
\end{equation}
where $u\ge0$ in the padding identity. Linearity and compatibility under
$F_d\hookrightarrow F_{d+1}$ follow by testing against all exterior words.

For $z_j=t_j+iy_j$ with $y_1<\cdots<y_n$, set
\[
 G_\varepsilon^\xi(z,x)
       =E_\xi(-iz-2\varepsilon)U_\xi(x).
\]
For $4\varepsilon<\min_j(y_{j+1}-y_j)$ define
\begin{equation}
 \mathsf F_n^\xi(\boldsymbol A;z,\boldsymbol x)
  =\left\langle\Omega_\xi^{\rm OS},
    B_\varepsilon^\xi(A_1)
    \prod_{j=2}^n\left[
       G_\varepsilon^\xi(z_j-z_{j-1},x_j-x_{j-1})
       B_\varepsilon^\xi(A_j)\right]\Omega_\xi^{\rm OS}\right\rangle.
 \label{ch16:heat-function}
\end{equation}
The product is ordered from left to right with increasing $j$.
Padding and the semigroup law show that this expression is independent
of $\varepsilon$: replacing it by $\varepsilon+u$ moves two factors
$E_\xi(u)$ onto each intervening gap, while the exterior factors fix the
vacuum. At purely imaginary times it is $S^\xi$ at the corresponding
Euclidean points. It is holomorphic in the time variables and continuous
in the real spatial variables. No norm continuity of the spatial
unitaries is required.

Contraction and \eqref{ch16:slabs}, with $\varepsilon=\delta/8$, give
\begin{equation}
 |\mathsf F_n^\xi(t+iy,x)|\le C\delta^{-m},\qquad
 0<\delta\le1,\quad y_{j+1}-y_j\ge\delta,
 \label{ch16:gap-bound}
\end{equation}
uniformly in all real positions and in $\xi$.
This is the quantitative regularity needed for reconstruction.

\begin{proposition}[Joint tempered boundary]
\label{ch16:boundary}
For every label tuple there is a tempered distribution
$W_n^\xi(\boldsymbol A)$ in all $4n$ real coordinates, obtained as the
boundary of \eqref{ch16:heat-function} along arbitrary strictly ordered
heights tending to zero. The distributions obey one common Schwartz
seminorm bound, depend continuously on $\xi$ when tested, and have the
mixed correlation distribution as their $\nu$-average. The same assertions
hold at fixed weakly ordered imaginary heights, including tied heights.
\end{proposition}
\begin{proof}
For a joint Schwartz test function $f$ put
$\psi_y^\xi(f)=\int\mathsf F_n^\xi(t+iy,x)f(t,x)\,dt\,dx$.
Let $r=(0,1,\ldots,n-1)$ and $D=-i\sum_jr_j\partial_{t_j}$ acting on test functions.
Cauchy's formula, \eqref{ch16:gap-bound}, and integration by parts give
$\frac d{ds}\psi_{y+sr}^\xi(f)=\psi_{y+sr}^\xi(Df)$.
Taylor's formula, evaluated backwards from $s=1$, is the exact identity
\begin{equation}
 \psi_y^\xi(f)=
 \sum_{k=0}^m\frac{(-1)^k}{k!}\psi_{y+r}^\xi(D^kf)
 -\int_0^1\frac{(-s)^m}{m!}
                   \psi_{y+sr}^\xi(D^{m+1}f)\,ds.
 \label{ch16:taylor-boundary}
\end{equation}
The integrand is bounded by $C\|D^{m+1}f\|_1$ for every weakly ordered
$y$: the factor $s^m$ cancels the worst gap singularity. The right-hand
side therefore defines the extension to all such $y$, continuously by
dominated convergence. Its bound is a fixed finite sum of $L^1$ norms
of derivatives of $f$, hence a Schwartz seminorm. In particular there
is no restriction on ratios of the heights approaching zero.

For strictly ordered heights, phase continuity follows from the normal
family theorem in the time variables: every convergent phase sequence
has locally convergent analytic subsequences, and agreement at imaginary
times identifies each limit with the limiting phase by the identity
theorem. Local boundedness is uniform by \eqref{ch16:gap-bound}.
Dominated convergence after testing gives continuity of $\psi_y^\xi$.
The same uniqueness argument continues the barycenter identity from
imaginary times. Formula~\eqref{ch16:taylor-boundary} transfers both
claims to weakly ordered heights and to the boundary. Arity zero and
one require no gap argument and are covered by their constant bounds.
\end{proof}

\section{Lorentz covariance and the spectral cone}
\label{ch16:covariance-section}

Let $s_4$ be reflection of Euclidean coordinate four. The field adjoint
on $F$ is the antilinear involution
\[
                       A^*=R(s_4)\overline A.
\]
If $L_{ab}$ is the infinitesimal Euclidean rotation action on labels,
the spatial rotation generators are $L_{ij}$ and the boost generators
are $-iL_{4i}$. Their commutation relations give the proper orthochronous
Lorentz action $\rho$. It descends from the covering group because a full
spatial turn already acts trivially under $R$. The above involution
commutes with $\rho$: the sign from conjugating $i$ cancels the reflection
sign on $L_{4i}$. We use spacetime coordinates $(x^1,x^2,x^3,t)$,
metric $\eta=\operatorname{diag}(-1,-1,-1,1)$, and
\[
 \overline V_+=\{(p,E):E\ge|p|\}.
\]
Fourier duality uses the ordinary coordinate pairing $p\cdot x+Et$.
The dual Lorentz action is the inverse transpose, which also preserves
$\overline V_+$.

Differentiate Euclidean rotational covariance at imaginary times and
continue the identity using the Cauchy--Riemann equations. The result
is the infinitesimal Lorentz covariance identity for $\mathsf F_n^\xi$.
Testing and taking its joint boundary proves the distributional identity:
terms containing a height tend to zero because the tested boundary
families have the common seminorm bound of Proposition~\ref{ch16:boundary}.
Integrating the finite-dimensional label action and the corresponding
flow on Schwartz test functions proves Lorentz covariance of $W_n^\xi$.
Translations and spatial rotations continue directly.

Use the Fourier convention
\[
 \widehat f(X)=\int e^{-iX\cdot p}f(p)\,dp,
 \qquad \widehat W(f)=W(\widehat f).
\]
For a compactly supported test function $b$ in momentum space, differentiating simultaneously
the imaginary height and the factor multiplying $b$ gives
\begin{equation}
 \widehat W_n^\xi(b)
   =\psi_y^\xi\!\left(\widehat{e^{\sum_jy_jp_j^4}b}\right).
 \label{ch16:laplace}
\end{equation}
Indeed the two differentiated terms cancel, and their common value at
height zero is the left side. Compact frequency support makes every
intermediate test function a Schwartz function.

Translation invariance restricts support to $\sum_jp_j=0$. Put
$q_j=\sum_{k=j}^np_k$. If the support of $b$ lies in $q_j^4<0$,
let the corresponding imaginary gap tend to infinity in
\eqref{ch16:laplace}, keeping the other gaps positive. Exponential decay
of the weighted test function and the uniform bound at gaps at least one show that its
pairing vanishes. A partition of unity and Lorentz covariance now give
\begin{equation}
 \operatorname{supp}\widehat W_n^\xi\subset K_n,
 \qquad K_n=\left\{p:\sum_jp_j=0,\quad q_j\in\overline V_+ (2\le j\le n)\right\}.
 \label{ch16:cone-support}
\end{equation}
To justify the last step, a vector outside $\overline V_+$ can be made
to have negative time component by a proper orthochronous Lorentz
transformation. Compact subsets of the complement admit finitely many
such neighborhoods.

\section{The spacetime tube and locality}
\label{ch16:locality-section}

The boundary distribution just constructed is denoted by $W_n^\xi$;
$K_n$ in \eqref{ch16:cone-support} is its momentum support cone. Define
\[
 C_n=\{Y:Y_{j+1}-Y_j\in\overline V_+\},\qquad
 \mathcal T_n=\{X+iY:Y_{j+1}-Y_j\in V_+\}.
\]
For $p\in K_n$, telescoping gives
$p\cdot Y=\sum_{j=2}^nq_j\cdot(Y_j-Y_{j-1})\ge0$.
For strictly interior $Y$ this is bounded below by $c_Y|p|$, locally
uniformly in $Y$. One way to see the strict margin is to choose a ball
of radius $2c_Y$ around $Y$ inside the strict-gap cone and apply positivity
to $Y-c_Yp/|p|$.

Choose a smooth cutoff $\chi$ equal to one on the distance-one neighborhood
of $K_n$, supported in its distance-three neighborhood, with all derivatives
bounded. Convolution of the distance-two neighborhood's indicator with
a bump of radius $1/2$ is one such choice. Then
\begin{equation}
 \mathcal W_n^\xi(\boldsymbol A;Z)
   =(2\pi)^{-4n}\widehat W_n^\xi
                    \left(\chi(p)e^{ip\cdot Z}\right),\qquad Z\in\mathcal T_n,
 \label{ch16:full-tube}
\end{equation}
is well defined: the tested function and all its derivatives decay
exponentially on the support of $\chi$. It is holomorphic into Schwartz
space before the distribution is applied. It agrees with
\eqref{ch16:heat-function} at strictly ordered scalar heights, by
\eqref{ch16:laplace} and compact-frequency approximation.

For $Y\in C_n$, the multiplier $\chi e^{-p\cdot Y}$ and its derivatives
are uniformly bounded on compact sets of $Y$. In fact distance at most
three from $K_n$ gives $-p\cdot Y\le3|Y|$. Multiplication by this function
is therefore continuous on Schwartz space, including as $Y$ approaches
a tied gap. This proves the partial boundary formula
\begin{equation}
 W_{n,Y}^\xi(f)=\widehat W_n^\xi
 \left(\chi e^{-p\cdot Y}(2\pi)^{-4n}
                         \int e^{ip\cdot X}f(X)\,dX\right).
 \label{ch16:partial-boundary}
\end{equation}
It has boundary $W_n^\xi$ at $Y=0$. The cutoff is essential here: the
exponential alone need not be a Schwartz multiplier off the spectral
cone. Tube Lorentz covariance follows from boundary covariance; changed
cutoffs agree near $K_n$ and hence give the same pairing.

\begin{proposition}[Adjacent spacelike exchange]
\label{ch16:locality}
Let $\sigma$ exchange slots $j$ and $j+1$. On the open set where
$X_{j+1}-X_j$ is spacelike,
\[
 W_n^\xi(\boldsymbol A;X)
       =W_n^\xi(\sigma\boldsymbol A;\sigma X)
\]
as distributions. This statement permits arbitrary exterior insertions.
\end{proposition}
\begin{proof}
Keep positive scalar imaginary gaps except for one tied pair
$y_j=y_{j+1}$. Write its spacelike real difference as $(v,t)$, and set
$e=v/|v|$. The imaginary part of its image under a boost of parameter
$i\alpha$ in direction $e$ is
$\sin\alpha\,(te,|v|)$, a strictly forward vector for small positive
$\alpha$. All other gaps stay strictly forward by continuity. Thus a
small positive imaginary boost enters the original tube, and a negative
one followed by exchange enters the exchanged tube.

Fix one sufficiently small nonzero boost parameter first. The transformed
tied-gap point is strictly inside the relevant tube, so a whole complex
neighborhood of the original point is mapped into that tube. Pulling
back the holomorphic function, with its entire finite-dimensional label
action, therefore defines a genuine germ on that neighborhood. Use the
opposite sign and the exchanged tube for the second germ. Real Lorentz covariance and the identity
theorem show that each germ extends its original tube function; the
whole short imaginary-boost path stays in that tube at neighboring
interior points.

At vanishing real times these germs coincide. Their restrictions to the
Euclidean slice $(it_k,x_k)$ equal the same smooth, permutation-symmetric
function on its two ordered sides. At the tied pair the spatial points
are distinct, so this Euclidean function is smooth there. Consequently
all derivatives of the two restrictions agree at the tie. The slice
is totally real and its tangent space spans the ambient complex vector
space. The holomorphic Taylor coefficients therefore agree, proving
identity of the germs.

For general real times multiply every real time by $s\in[0,1]$.
The distinguished difference remains spacelike along this path and the
other imaginary gaps remain strict. A finite chain of overlapping germ
neighborhoods propagates the equality from $s=0$ to $s=1$ by uniqueness.
Take the continuous partial boundaries \eqref{ch16:partial-boundary},
and then remove the other imaginary gaps. This proves the identity
against test functions compactly supported in the spacelike open set. Smooth
compact cutoffs converging in Schwartz space give the stated distributional
version. No permutation of coincident Lorentzian boundary values has
been assumed.
\end{proof}

\section{Field words and the reflection space}
\label{ch16:fields-section}

A letter is a pair $(A,f)$ with $A\in F$ and $f\in\mathscr S(\RR^4)$.
The star of a word reverses its order and replaces every letter by
$(A^*,\overline f)$. Let $\mathcal P$ be their algebraic span, including
the empty word $1$. The distributions $W_n^\xi$ define a functional
$w_\xi$ on $\mathcal P$ by evaluation on product test functions.

The essential positivity statement follows inside the phase reflection
space. Give a word strictly increasing positive imaginary heights,
smear its real variables, and use the slab product to obtain a vector
$h_y^\xi(b)$. Its Gram pairing with another such vector is a tested
continued correlation, with the reflected first word on the left.
Proposition~\ref{ch16:boundary} permits the two height lists to tend to
zero independently, so
\[
 \langle h_y^\xi(b),h_z^\xi(c)\rangle
                      \longrightarrow w_\xi(b^*c).
\]
Expanding $\|h_y^\xi(b)-h_z^\xi(b)\|^2$ proves that these vectors are
Cauchy. Their limits $v_\xi(b)$ have Gram form $w_\xi(b^*c)$, which is
therefore positive. The same common seminorm bounds extend positivity
to finite sequences of joint Schwartz test functions by tensor-product density.

The standard word reconstruction \cite{Wightman1956} now gives
\[
 \mathcal N_\xi=\{b:w_\xi(b^*b)=0\},\quad
 \mathcal D_\xi=\mathcal P/\mathcal N_\xi,\quad
 \mathcal H_\xi=\overline{\mathcal D_\xi},\quad\Omega_\xi=[1],
 \qquad \Phi_\xi(A,f)[b]=[(A,f)b].
\]
For clarity, its null-space check is worth recording. If $b$ is null,
Cauchy--Schwarz says it pairs to zero with every word. The concatenation
identity moves a letter to the other side of the pairing, so its product
with $b$ also pairs to zero with every word, including itself. Thus all
fields are well defined on the same invariant dense domain. The same
identity gives
\[
 \Phi_\xi(A^*,\overline f)\subset\Phi_\xi(A,f)^*.
\]
They are closable, and the joint distribution bound gives continuity of
$f\mapsto\Phi_\xi(A,f)v$ in Hilbert norm for every $v\in\mathcal D_\xi$.
We do not assert that these generally unbounded fields are bounded or
that their self-adjoint closures exist without further input.

\begin{theorem}[Fields and unitary equivalence with the reflection-positive space]
\label{ch16:fields}
In every phase the fields above are local, covariant under a strongly
continuous unitary representation of the proper orthochronous Poincar\'e
group, and have joint translation spectrum in $\overline V_+$.
The map
\[
 J_\xi^{\rm OS}[b]=v_\xi(b)
\]
extends to a unitary map from $\mathcal H_\xi$ onto the phase reflection
space. It sends vacuum to vacuum, intertwines spatial translations and
heat semigroups, and the field labeling is faithful.
\end{theorem}
\begin{proof}
Covariance preserves the word Gram form, hence defines unitaries on its
quotient. Schwartz continuity proves strong continuity first on words
and then on their completion. Locality follows by applying
Proposition~\ref{ch16:locality} inside arbitrary word matrix elements.
The joint spectral theorem for the commuting translation unitaries,
together with \eqref{ch16:cone-support} applied to such matrix elements,
places their spectrum in $\overline V_+$.

The Gram identity makes $J_\xi^{\rm OS}$ an isometry. Its range is closed,
but density must still be checked. For a smooth compactly supported
joint frequency test function $b$, multiplication by
$e^{-\sum_j y_jp_j^4}$ is harmless. Equation~\eqref{ch16:laplace} identifies
the norms and mixed pairing of the image of the resulting boundary word
with the smeared positive-height vector at heights $y$. They are the
same vector. Test functions with compact Fourier support are dense in Schwartz space;
therefore the closed range contains every smeared positive-height vector.
Approximate identities, followed by a perturbation of tied times and
permutation of separated Euclidean points, give all Euclidean word
vectors. These span the reflection space. The same positive-height
identities prove heat and spatial intertwining.

Finally a label acting as zero has every boundary correlation zero.
The Laplace identity, Euclidean symmetry and continuity imply that all
its separated Euclidean correlations vanish. The common null-space
identity of Chapter~\ref{ch15:phases}, and the injectivity of the matrix of correlation pairings used
there, force this label to be zero in $F$.
\end{proof}

\section[Vacuum uniqueness in each phase]{The low-energy estimate and uniqueness of the vacuum in each phase}
\label{ch16:vacuum-section}

The phase measure has full support, and the zero-energy space is a line
for almost every phase by \ref{ch15:ae-ground}. To remove the exceptional
set requires a bound uniform in $\xi$. We prove it from locality and the
spectral cone. This argument controls energies near zero, not a positive
mass gap.

For a finite sum $A$ of field words with compactly supported smooth test functions,
let $\mu_\xi$ be the joint spectral measure of $A\Omega_\xi$, and let
$\mu_\xi^*$ be that of $A^*\Omega_\xi$. Their total masses have a common
bound $M$ by Proposition~\ref{ch16:boundary}. Enclose all supports in a
fixed spacetime box. Locality gives some $R<\infty$, independent of phase,
such that
\begin{equation}
 a_\xi(y,t)=\langle\Omega_\xi,[A^*,\alpha_{(y,t)}(A)]\Omega_\xi\rangle=0
                  \quad (|y|>R,\ |t|\le1),
 \label{ch16:local-commutator}
\end{equation}
where $\alpha_z(A)=U_\xi(z)AU_\xi(z)^*$ on the word domain.
The spectral expression is the difference of the Fourier transforms of
$\mu_\xi$ and the reflected measure $\mu_\xi^*$.

Take the odd time average
\[
 \bar a_\xi(y)=\int_0^1
       \frac{a_\xi(y,t)-a_\xi(y,-t)}{2i}\,dt.
\]
It is bounded by $M$ and supported in $\overline B_R$. Its spatial
Fourier measure is the \emph{sum} of the pushforwards of
$\varpi(E)\,d\mu_\xi(p,E)$ and $\varpi(E)\,d\mu_\xi^*(p,E)$ under
$p$ and $-p$, respectively, where
\[
 \varpi(E)=\frac{1-\cos E}{E},\qquad\varpi(0)=0.
\]
Both measures are positive; on $0\le E\le1$,
$\varpi(E)\ge E/\pi$. Choose a smooth momentum bump $b_u$, equal to one
on $B_u$, supported in $B_{2u}$, and between zero and one. Its inverse
Fourier transform is bounded by $Cu^3$. On the shell $u/2<E\le u$,
cone support implies $|p|\le u$, so positivity and Fourier inversion yield
\[
 c u\,\mu_\xi\{u/2<E\le u\}
      \le \left|\int\bar a_\xi(y)\widehat b_u(y)\,dy\right|
      \le CMR^3u^3.
\]
Dyadic summation proves the first assertion below.

\begin{theorem}[Uniform low-energy bound]
\label{ch16:vacuum}
For every such $A$ there is $C_A$, independent of phase, such that
\begin{align}
 \mu_\xi\{0<E\le u\}&\le C_Au^2 &&(0<u\le1),\label{ch16:low-energy}\\
 0\le\langle A\Omega_\xi,(e^{-sH_\xi}-P_{\Omega,\xi})A\Omega_\xi\rangle
       &\le 2C_A' s^{-2} &&(s\ge1),\label{ch16:heat-decay}
\end{align}
where $P_{\Omega,\xi}$ is the projection onto $\CC\Omega_\xi$.
In every phase, $\ker H_\xi=\CC\Omega_\xi$; the same holds in its
reflection space.
\end{theorem}
\begin{proof}
Initially use the full zero-energy projection $Q_\xi$, without assuming
that its range is a line. Increase the constant in
\eqref{ch16:low-energy} so the estimate holds for all $u>0$ by the total
mass bound. Tonelli's theorem gives
\[
 \langle A\Omega_\xi,(e^{-sH_\xi}-Q_\xi)A\Omega_\xi\rangle
 =\int_{E>0}e^{-sE}\,d\mu_\xi
 \le C_A's\int_0^\infty e^{-su}u^2\,du=2C_A's^{-2}.
\]
Theorem~\ref{ch16:fields} transfers almost-everywhere vacuum uniqueness
from the reflection spaces, so $Q_\xi=P_{\Omega,\xi}$ almost everywhere.
For fixed $s>0$ the vacuum-subtracted heat pairing is continuous in $\xi$:
it is the joint distribution at a fixed weakly ordered height with gap
$s$ between its two word halves. This follows either by the independent
height limits or by \eqref{ch16:partial-boundary}. Thus its bound by
$2C_A's^{-2}$ holds everywhere, since any nonempty open violation set
would have positive $\nu$-measure.

Let $s\to\infty$. Strong heat convergence gives
$\|Q_\xi A\Omega_\xi\|^2=|\langle\Omega_\xi,A\Omega_\xi\rangle|^2$.
Since $Q_\xi$ fixes the unit vacuum, its value on each such word vector
lies in the vacuum line. Compactly smeared words are dense: multiply
Schwartz test functions by expanding smooth cutoffs and use joint Gram continuity.
It follows that $Q_\xi=P_{\Omega,\xi}$ on the full space. The onto heat
intertwiner gives the last assertion.
\end{proof}

\section{The weak commutant and local stress charges}
\label{ch16:charges-section}

Irreducibility should be stated in a form appropriate to unbounded fields.
Their common domain is $\mathcal D_\xi$, the span of field words applied
to the vacuum. A bounded operator $B$ belongs to their weak commutant if
\begin{equation}
 \langle\Phi^*u,Bv\rangle=\langle u,B\Phi v\rangle,
 \qquad u,v\in\mathcal D_\xi,
 \label{ch16:weak-commutant}
\end{equation}
for every smeared field $\Phi$. The adjoint inclusion proved above
ensures that these $u$ are in its adjoint domain; no preservation of
$\mathcal D_\xi$ by $B$ is assumed.

\begin{proposition}
\label{ch16:irreducible}
The weak commutant in every phase consists of scalar operators.
\end{proposition}
\begin{proof}
Iterate \eqref{ch16:weak-commutant} to move words through matrix elements.
For a word $A$ put $u=A\Omega_\xi$, $v=A^*\Omega_\xi$. The functions
\[
 p(z)=\langle u,e^{izH_\xi}B\Omega_\xi\rangle,\qquad
 q(z)=\langle B^*\Omega_\xi,e^{izH_\xi}v\rangle
\]
are bounded and holomorphic in the upper half-plane, continuous to its
boundary, and satisfy $p(t)=q(-t)$. Pasting across the real line and
using Morera's theorem gives a bounded entire function, hence a constant.
The heat limit and Theorem~\ref{ch16:vacuum} imply
$B\Omega_\xi=b\Omega_\xi$. Moving words again gives
$\langle A'\Omega_\xi,BA\Omega_\xi\rangle
=b\langle A'\Omega_\xi,A\Omega_\xi\rangle$.
Density proves $B=bI$. Scalars plainly satisfy the condition.
\end{proof}

Let $T_{ab}$ be the normalized Euclidean stress tensor of
Proposition~\ref{ch14:normalized-stress}. Define
\[
 \zeta_i=1\ (i\le3),\qquad\zeta_4=-i,
 \qquad T^{\rm M}_{ab}=-\zeta_a\zeta_bT_{ab}.
\]
We use the same notation for its smeared relativistic field. Thus
$T^{\rm M}_{44}=T_{44}$, $T^{\rm M}_{4i}=iT_{4i}$, and
$T^{\rm M}_{ij}=-T_{ij}$. The adjoint involution, covariance and the
continued derivative identity give a symmetric Hermitian Lorentz tensor
with
\[
 \sum_a\eta_{aa}T^{\rm M}_{ab}(\partial_a f)=0
                      \quad\hbox{on }\mathcal D_\xi.
\]
Indeed continuation of a Euclidean derivative label is
$\Phi_\xi(\partial_a A,f)=-\zeta_a^{-1}\Phi_\xi(A,\partial_a f)$.

For an affine Killing vector $K(x)=k+Vx$, where
$V^{\mathsf T}\eta+\eta V=0$, let $\Sigma_K$ be its infinitesimal
Lorentz label action. Choose real $f\in C_c^\infty(\RR)$ of integral
one and $c\in C_c^\infty(\RR^3)$, and define the local charge
\[
 Q_K(f,c)=\sum_b T^{\rm M}_{4b}\big(K^b(c\otimes f)\big).
\]
Suppose $g\in C_c^\infty(\RR^4)$ and $c=1$ on a ball such that every
point $(y,t)$ outside that ball with $t\in\operatorname{supp}f$ is
spacelike to $\operatorname{supp}g$. Then
\begin{equation}
 i[Q_K(f,c),\Phi_\xi(A,g)]
   =\Phi_\xi(\Sigma_K A,g)
          -\sum_a\Phi_\xi(A,\partial_a(K^ag))
                       \quad\hbox{on }\mathcal D_\xi.
 \label{ch16:charge-law}
\end{equation}
This is a local charge identity, not a claim that the unsmeared spatial
integral of stress exists as an operator.

To justify its continuation, first integrate the Euclidean shell identity
with a spatial cutoff $c$ on the two sides of the distinguished insertion.
Conservation moves the derivative to $c$; that term stays separated from
the insertion. The two current integrals have one-sided limits whose
jump is $\Sigma_K A+K\cdot\partial A$. Retain strictly ordered positive
heights for every spectator and put their complex times on a path
$z_j(\lambda)=iy_j+\lambda t_j$. Place the current immediately above
or below the distinguished time, at distance $\varepsilon$.

On $\operatorname{supp}\nabla c$ the spacelike holomorphic germs from
Proposition~\ref{ch16:locality} bound the current integrals uniformly
as $\varepsilon\downarrow0$, locally along this path. Conservation bounds
the difference between two such integrals by $C|\varepsilon-\varepsilon'|$.
Their limits are therefore holomorphic. At small imaginary $\lambda$
the jump is the Euclidean one; analytic continuation gives the same
jump at $\lambda=1$. Smearing with $f$ and $g$ and taking the joint
partial boundaries proves \eqref{ch16:charge-law}; integration by parts
produces $\partial_a(K^ag)$. General word spectators follow from the
common Schwartz bounds. Locality removes any change in $c$ outside the
stated ball. This retains precisely the normalization established by the
Euclidean shell flux.

\setcounter{chapter}{16}
\chapter{The trace law and its consequences}
\label{ch17:trace}

The leading short-distance limit of the stress tensor is the traceless
Maxwell tensor. It therefore says nothing by itself about the leading
nonzero term in the trace. That term is smaller by one inverse running
coupling, and its two-point function by two. The perturbative origin of
this phenomenon is the trace anomaly of Collins, Duncan and Joglekar
\cite{CollinsDuncanJoglekar1977}. The all-radius assertion below requires
additional estimates for the selected lattice construction; it does not
follow by substituting a perturbative beta function into a formal identity.

Retain the Ward-normalized tensor $T$ of
Proposition~\ref{ch14:normalized-stress}, the phase $\xi_*$ chosen in
Section~\ref{ch15:selection}, and its relativistic reconstruction.
The latter supplies a positive-cone translation spectrum and tempered
boundary correlations. We prove the stronger quantitative current
transfer needed for the trace asymptotic, including its low-moment
cancellations. The construction and reconstruction in the preceding
chapters do not depend on this refinement.

\section{Centering and the positive spectral measure}
\label{ch17:spectrum}

Put $\tau=\sum_{a=1}^4T_{aa}$, $e_{\rm tr}=S^{\xi_*}(\tau)$ and
\[
             \Theta=\tau-e_{\rm tr}1,\qquad
             G_\Theta(r)=S_2^{\xi_*}(\Theta,\Theta;re_4,0).
\]
The real number $e_{\rm tr}$ is the expectation in the selected phase,
which need not equal the mixed expectation. The unit field is retained
through the common filtration. The fact that inserting the identity leaves correlations unchanged extends to
the boundary correlations, so, on the common invariant word domain,
\[
 \Phi_\Theta(f)=\Phi_\tau(f)
                  -e_{\rm tr}\left(\int f\right)I,
             \qquad \langle\Omega,\Phi_\Theta(f)\Omega\rangle=0.
\]
For completeness, the Euclidean correlation with a unit insertion is
constant in that insertion position. Holomorphic uniqueness on ordered
time regions identifies its continuation with the correlation after
deletion. Boundary values and smearing give the displayed operator
identity on dense word matrix elements. This establishes centering on
the domain, as well as for separated expectations.

Write $p=(\boldsymbol p,E)$, $p^2=E^2-|\boldsymbol p|^2$, and
$\overline V_+=\{E\ge|\boldsymbol p|\}$. Use
$\widehat f_+(p)=\int e^{ip\cdot x}f(x)\,dx$ in the translation
spectral convention, and
$\mathcal W\psi(p)=\int e^{ip\cdot(y-x)}\psi(x,y)\,dx\,dy$.
Then $\mathcal W(\overline f\otimes h)=
\overline{\widehat f_+}\widehat h_+$.
Positivity, translation covariance and temperedness give a unique
positive measure $\rho_\Theta$ of polynomial growth, supported in
$\overline V_+$, with
\begin{align}
 W_2(\Theta,\Theta;\psi)&=\int\mathcal W\psi\,d\rho_\Theta,
 \label{ch17:spectral-pairing}\\
 G_\Theta(r)&=\int e^{-rE}\,d\rho_\Theta(p).
 \label{ch17:laplace}
\end{align}
Here is the direct spectral identification needed for these formulas.
For $v_f=\Phi_\Theta(f)\Omega$, covariance gives
$\widehat h_+(P)v_f=\widehat f_+(P)v_h$ by convolution.
On a compact momentum set choose $\widehat f_+$ nonzero and define
$d\rho_\Theta=|\widehat f_+|^{-2}d\mu_{v_f}$ there. The convolution
identity makes the definitions agree on overlaps. Polarization gives
\eqref{ch17:spectral-pairing} first on product test functions and then on joint
Schwartz test functions. Tempered continuity bounds compact masses polynomially.
The positive cone permits exponential damping, and uniqueness of the
holomorphic boundary value identifies its Laplace transform with the
Euclidean function. The scalar transformation law of $\Theta$ makes
$\rho_\Theta$ Lorentz invariant. In particular $G_\Theta$ is real,
nonnegative and nonincreasing. The expression
$r^8\log^2(1/r)G_\Theta(r)$ need not be monotone.

\section{Comparison at prescribed radii}
\label{ch17:radii}

Let $U_{i,d}:\mathcal M_d^{\rm Max}\to F_d$ be real comparison
isomorphisms, preserving the unit, inclusions and signed coordinate
symmetries. They are evaluated at a radius $r_i\downarrow0$:
$U_{i,d}A$ is inserted at $r_ix$. The actual inverse coupling at a
finer insertion level $k_i$ is denoted by $\beta_i$, and
$b_i=\widehat\beta_{k_i}$ is its reference value. The radius level
$n_i$ precedes $k_i$.

The strengthened comparison, proved below, is as follows. Every prescribed
positive null sequence has a subsequence with these data such that
\begin{equation}
 \beta_i/b_i\to1,\qquad
 \frac{\beta_i}{\log(1/r_i)}\to\kappa_G,
 \label{ch17:running-radius}
\end{equation}
and all fixed separated correlations and their fixed derivatives,
with convergent moving free labels, converge to their Maxwell values.
For pairs without derivatives this comparison is uniform over the
phases, including at nearest lattice-rounded points. Shell integrals are
included. It must use the same field quotient and lattice family as
before. A comparison along one freely chosen scale sequence is weaker.

We explain the radius coverage and the relevant quantitative demands.
Write $M_n=M_n^{\rm cut}=\exp((\log\log\widehat\beta_n)^2)$ and
allow
\[
 r(n,u)=a_nM_n^u,\quad \alpha_0\le u\le\alpha_0+1,\qquad
 k(n)=\min\{m>n:\log(a_n/a_m)\ge(\log\widehat\beta_n)^{48}\}.
\]
At fixed degree bound $D$ and arity bound $N_*$, choose comparison
exponents with
\begin{gather*}
 c<0,\quad\alpha_0-\alpha_1>6,\quad
 \alpha_0-\alpha_2>5-c,\\
 (\alpha_0-\alpha_1)DN_*+c<0,\qquad
 (\alpha_0-\alpha_2)DN_*-\alpha_0-c<0.
\end{gather*}
The common size of the $\alpha_j$ can be increased after their
differences and $c$ have been fixed. The marked comparison required
at every sufficiently late level has, uniformly in $u$, single-insertion
seminorm bounds $CM_n^{(\alpha_0-\alpha_j)D}$ and grid errors tending
to zero. With strip width $M_n^c$ and scaled mesh at most
$M_n^{-\alpha_0}$, the two integration errors are bounded by
\[
 C M_n^{(\alpha_0-\alpha_1)DN_*+c},\qquad
 C M_n^{(\alpha_0-\alpha_2)DN_*-\alpha_0-c}.
\]
They tend to zero. Thus the strip removal and derivative-transport
argument of Section~\ref{ch14:comparison} applies to these windows.
Bounds uniform in the number of copies give phase uniformity by
\eqref{ch15:moment-norm}. The next subsection verifies these window estimates, including their
uniformity across changes of stage, from marked propagation.

The coupling thresholds can be chosen to ensure $L_n\le\log\widehat\beta_n$.
Since $M_n$ increases and $L_n\le M_{n+1}$ eventually,
$r(n,\alpha_0)\le r(n+1,\alpha_0+1)$. The intervals therefore
cover every sufficiently small radius. The first crossing places a
prescribed radius in one of them. Diagonalize finite degree, arity,
window and derivative requests, requiring
$(\alpha_0+1)\log M_{n_i}\le(\log\widehat\beta_{n_i})^{48}$.
The overshoot in $k(n)$ is negligible compared with
$\widehat\beta_n$, as is the logarithm of the radius multiplier.
The running law in Chapter~\ref{ch10:trajectory} then gives
\eqref{ch17:running-radius}. This proves radius coverage. It remains to verify the uniform
comparison estimates used in these windows.

% !TEX root = ../monograph.tex
\subsection{Intervals of scales crossing changes of stage}
\label{ch17:variable-windows}

The finite comparison in Proposition~\ref{ch11:n-mark-comparison}
extends to the intervals of radii just described. We give the estimates
because requiring every interval of scales to lie inside a plateau would leave
gaps in radius coverage.

Fix the maximum scaling dimension of the insertions, their number, and
the orders of the derivatives to be estimated. In this paragraph
$C_j$ includes the linear column constants with insertions, the curvature-response
row constants and the finite polynomial-source propagation constants
for this request, all fixed \emph{before} the block ratio of stage $j$.
Write $\mathscr L_j$ for the ratio chosen at stage $j$, so
$L_i=\mathscr L_{J_i}$ with $J_i=J(i+1)$ in the scale convention
of Chapter~\ref{ch10:trajectory}. The ratio choices may require
\[
 \varepsilon_j:=\frac{\log(2+C_j)}{\log\mathscr L_j}\longrightarrow0.
\]
For example one may impose $\varepsilon_j\le j^{-2}$.
Constants depending on the already selected ratio, such as nonlinear
jet bounds of the critical configuration, are not included here. They enter the local
small errors and endpoint constants, and are handled by sufficiently large lower bounds on the reference coupling.
These local errors are summed after transfer; their ratio-dependent
constants are never multiplied once per step. Put $H_n=\exp((\log\log\widehat\beta_n)^2)$. Only a bounded number
of endpoint factors is multiplied, giving $H_n^{o(1)}$.
This distinction avoids a circular ratio choice.

For $b$ in a fixed compact interval of sufficiently large positive
numbers, move backwards from scale $n$, adding $\log L_i$ at each step.
Stop when this sum first reaches $(b+2)\log H_n$. Starting from that
scale, continue until the additional sum first reaches $10\log B_n$. Here $B_n$ is the radius of the restarted representation;
$\log B_n=O(\log\log\widehat\beta_n)$.
The schedule bound $L_i\le\log\widehat\beta_i$ makes each
excess over the target $O(\log\log\widehat\beta_n)=o(\log H_n)$.
The entire interval of scales has logarithmic length $O(\log H_n)$, and the
number of steps is $O(\log H_n)$ since every $L_i\ge3$.
It is therefore shorter than $(\log\log\widehat\beta_n)^4$,
as required by Proposition~\ref{ch11:restart-family}.
All its scale indices tend to infinity. The increment estimate is a
constant times the logarithmic block length plus the
local drift error controlled by the lower bound on the reference coupling. Summing it over this interval gives
\[
 \max_i|\beta_i/\beta_n-1|
 \le \frac{H_n^{o(1)}(1+\log H_n)}{\widehat\beta_n}.
\]
In particular coupling changes are smaller than every fixed inverse
power of $H_n$, not just qualitatively $o(1)$. An interval of scales may cross any number of stages; no bound by two
stages is needed.

For the linear propagation constants the exact bound is
\begin{equation}\label{ch17:window-product}
 \log\prod_{i\text{ in window}}(2+C_{J_i})
 \le\sup_{i\text{ in window}}\varepsilon_{J_i}
                \sum_{i\text{ in window}}\log L_i
 =o(\log H_n).
\end{equation}
Thus they contribute a total factor $H_n^{o(1)}$. In particular the differentiated
curvature-response estimate in Chapter~\ref{ch11:insertions} becomes
\[
 S^{r+2}\sum_y|\Delta^r(\Psi_{n-1}\cdots\Psi_\ell)(x,y)|
 \le 2^d\prod_{i=\ell}^{n-1}R_{J_i}=H_n^{o(1)},
 \qquad S=a_\ell/a_n.
\]
The covariance decomposition telescopes for these unequal
curvature responses just as it does for equal ratios: each conditional covariance
is conjugated by the product of the subsequent curvature responses.
It uses no replacement by a common stage's quadratic form.
The finite-cutoff covariance estimate \eqref{ch11:finite-covariance}
is uniform over the blocking sequence, in particular over the unequal ratios $L_i$ of the interval of scales (Section~\ref{ch11:window-covariance}). Its stage-dependent constants,
protected-path errors, ratio-dependent jet errors and finite incidence
costs are made smaller than the requested powers of $H_n$ by choosing
sufficiently large lower bounds on the reference coupling. Sums over $O(\log H_n)$ steps lose only a factor $H_n^{o(1)}$.
At each step the linear map has its norm bound fixed before the block
ratio is chosen. The remaining nonlinear term is small once the lower bound
on the reference coupling is sufficiently large. Constants depending on the chosen ratio
therefore enter the sum of local errors. The regular and large-field source recurrences therefore retain their
strict negative powers after multiplication by the growth factor
\eqref{ch17:window-product}.

One can now repeat the proof of
Proposition~\ref{ch11:n-mark-comparison} with unequal ratios.
The initial Taylor order and support margins are chosen for the
largest $b$ and the fixed finite request. The entire incoming
polynomial in the source variables is retained. The quadratic covariance
and the polynomial $\sum_IG_Iz_I$ of free moments are common to all subsets;
the density without insertions, and hence the partition function, is unchanged.
Absorbing the larger source weight still uses
$\omega=H_n^{1/(8N_*)}$ and $I=\log H_n/(16N_*)$.
Every additional factor lost in transfer is $H_n^{o(1)}$; the margins in the
$H_n^{-3/2}$ moment error, $H_n^{-31/32}$ connected error and
$H_n^{-1}$ bounds of Lemma~\ref{ch11:restore-regular} absorb it. The chessboard bound \eqref{ch11:n-mark-root} is
therefore at most $H_n^{-c_*}$ for a fixed $c_*>0$ depending only on
the finite arity, for example $c_*=1/(32N_*)$ after increasing the lower bound on the reference coupling.
Its cell side is at most $a_nH_n^{b+5}$.
The lower bound on the partition function is imposed at the stage of the coarsest scale of the interval of scales,
with its constant at most $I$; it is not multiplied once per step.

The common localization map $T$ depends only on the insertion
and its original self covariance. It is independent of the middle
and output stopping scales. Restart on the longest interval of scales and restrict
that restarted history to the shorter ones. This gives the same $T$ for all
$b$ in the compact interval. The finite number of coefficient bounds
above is uniform in $b$: changing $b$ only changes stopping integers
and the scaling dilation $W_{H_n^b}$, whose exponents lie in a
fixed compact interval. All upper bounds used the longest interval of scales and
the largest dilation factor. No compactness over arbitrary incoming
activity arrays is invoked.

Apply this construction to
$b=\alpha_j+u-\alpha_0$, $j=0,1,2$, and
$u\in[\alpha_0,\alpha_0+1]$.
At its own normalization a single insertion has bounded chessboard seminorm.
Changing from $\alpha_j$ to $\alpha_0$ contributes a factor at most
$CH_n^{(\alpha_0-\alpha_j)D}$ by the exact scaling dilation on
the common localization map. These are precisely the size bounds used
in the two estimates above for the coincident-coordinate strips and the grid. Their exponents are
strictly negative, so the harmless factors $H_n^{o(1)}$ do not change
the conclusion. The estimates with constants independent of the number of reflected copies give the same pair bounds
in every phase. The existing argument using separated correlation pairings then gives
fixed derivatives and shell fluxes with the same diagonal subsequence and the same comparison isomorphisms $U_{i,d}$.
Thus the prescribed-radius comparison used in this chapter follows
from the construction with insertions, including intervals of scales that cross
changes of stage.


\section{Tensor coordinates and the regulated current}
\label{ch17:current}

Use the free scalar and tensor
\[
 \mathcal O=\tfrac14\sum_{a,b,\alpha}{:}(F_{ab}^\alpha)^2{:},
 \quad L=T^{\rm Max},\quad D_{ab}=\delta_{ab}\mathcal O,
 \quad X_{ab}=(1-\delta_{ab})\sum_{c,\alpha}{:}F_{ac}^\alpha F_{bc}^\alpha{:}.
\]
The coordinate sum in $\mathcal O$ is over ordered pairs. The symbol
$X$ here denotes the off-diagonal tensor, not a polymer support.
These three tensors span the centered degree-four symmetric
hypercubic tensors. Before imposing the Maxwell equation, their
vector divergences have coordinates in $E=\RR^3$, represented by
\[
 \sum_c F_{bc}\partial_bF_{bc},\qquad
 \sum_{a,c}F_{ac}\partial_bF_{ac},\qquad
 \sum_{a,c}F_{bc}\partial_aF_{ac},
\]
with colour contraction and Wick ordering understood. The free
quotient $s:E\to V$ has kernel spanned by $v=(0,0,1)$, and $V$
is the two-dimensional vector space of Section~\ref{ch14:stress-line}.
For a tensor with diagonal entries
$z_0\sum_cF_{ac}^2+z_1\sum_{b,c}F_{bc}^2$ and off-diagonal entries
$z_2\sum_cF_{ac}F_{bc}$, Bianchi gives divergence coordinates
$(2(z_0-z_2),2z_1+z_2/2,z_2)$. Hence
\begin{equation}
 \rho(x,y,z)=-2x-8y,\qquad
 \rho(\operatorname{div}_{\rm raw}B)\mathcal O=-\operatorname{tr}B.
 \label{ch17:rho}
\end{equation}
In particular $\rho(v)=0$ and the values for $L,X,D$ are $0,0,-4$.

At a lattice site set
$f_{ab}=\operatorname{proj}_{\gLie}(U_{\square(a,b)}-
U_{\square(a,b)}^{-1})/2$ and start with
\[
 J_{ab}=\tfrac14\delta_{ab}\sum_{c,d}\langle f_{cd},f_{cd}\rangle
                  -\sum_c\langle f_{ac},f_{bc}\rangle,
 \qquad (Y_b^x)_e=\mathbf1_{e^-=x}f_{b,e^{\rm dir}}(x).
\]
Here $Y$ is a left link displacement. If $\mathsf S$ is the leading
path action without its coupling, define its current defect by
$\mathcal E_b=\sum_a\nabla_aJ_{ab}-D\mathsf S[Y_b]$, using
\[
 \nabla_a=\sum_{j=1}^4c_j(\mathsf s_a^j-\mathsf s_a^{-j}),
 \qquad(c_1,c_2,c_3,c_4)=(4/5,-1/5,4/105,-1/280).
\]
The operator $\mathsf s_a$ shifts by one site in direction $a$.
Initially the trace and antisymmetric part of $J$ vanish. During
transfer they must remain exact forward differences; smallness of their
values alone would not suffice for the trace rate below.

We fix the coefficient functionals in that calculation. For a local scalar
$F$, its Hessian Fourier kernel is
\[
 \widehat F_{c\alpha,c\alpha}(p,q)=
 \sum_{x,y}e^{i(p\cdot x+q\cdot y)}
 D^2F(0)[e_{x,c,\alpha},e_{y,c,\alpha}],
\]
where $e_{x,c,\alpha}$ is the link-coordinate basis vector. Fix one
colour $\alpha$. For a vector $F_b$ and a scalar action $F$ put
\begin{align*}
 \mathsf l(F)&=\operatorname{Re}\sum_{a\ne c,b}
      i[z^2t]\widehat{F_b}_{c\alpha,c\alpha}(ze_a,-ze_a+te_b),\\
 \mathsf q(F)&=|\Lambda|^{-1}\operatorname{Re}\sum_{a\ne c}
            [z^2]\widehat F_{c\alpha,c\alpha}(ze_a,-ze_a),\\
 \mathsf h(F)&=\mathsf l(F)-\tfrac12\operatorname{Re}\sum_{a\ne c}
            [z^2]\widehat{F_c}_{c\alpha,c\alpha}(ze_a,-ze_a).
\end{align*}
Brackets denote Taylor coefficients including reciprocal factorials.
The volume normalization in $\mathsf q$ is for the translation-invariant
action; $\mathsf l$ evaluates the anchored vector response. Define
$\mathfrak h(Y)=\mathsf h(\operatorname{div}_{\rm Haar}Y)$.
The second term in $\mathsf h$ corrects the lattice derivative.
The regulated current calculation uses
\[
 \Gamma_\mu=\tfrac12\log\det(H+dd^*+\mu^2I)
                 -\log\det(d^*d+\mu^2I),\quad
 R(\mu)=\mathsf q(\Gamma_\mu),\quad
 \mathcal D(\mu)=-\mu R'(\mu).
\]
Here $H$ is the left Hessian of $\mathsf S$, $d$ its infinitesimal
gauge insertion, and $\mathsf q$ is the coefficient functional just defined. Its normalization
counts twelve ordered transverse planes. The actual Ward calculation requires its Haar
coefficient $\mathfrak h(Y)$ to obey
\begin{equation}
 \mathfrak h(Y)=\mathcal D(\mu)+O(\mu)+o(1),\qquad
 R(\mu)=12\kappa_G\log\mu+O(1).
 \label{ch17:regulator-input}
\end{equation}
The $o(1)$ is uniform on each fixed interval $[\delta,1]$ as the
current defect tends to zero and the period tends to infinity;
the constant in $O(\mu)$ is independent of $\mu\le1$.
This identity concerns the actual gauge-fixed Hessian and transferred
current. The perturbative anomaly citation does not supply it.

The next section derives this regulated identity from explicit weighted
current bounds and convergence of the matched current jets. It keeps
the gauge-coordinate corrections and the lattice coefficient correction.
The following current-contraction section then verifies these hypotheses
for repeated stationary transfer.

% !TEX root = ../monograph.tex
% Fragment for Chapter17. Root owns integration and section numbering.
\section{Exact current transport and the scalar regulator}
\label{ch17:current-suppliers}

Several ingredients of the trace transfer can be established before its
quantitative current estimates. We give them separately, since they use
the actual difference operator and the actual path action. They do not
identify a lattice coefficient by comparison with a separately chosen
continuum action.

\subsection{Difference-preserving block averages}

Let $L$ be a block ratio and $s_a$ a unit fine translation. Write
$\delta_a=s_a-1$, $\delta_a^{\rm c}=s_a^L-1$, and let
$\nabla_a=\sum_{j=1}^4c_j(s_a^j-s_a^{-j})$ be the eighth-order
centered difference of Section~\ref{ch17:current}. Write
$\nabla_a^{\rm c}=\sum_{j=1}^4c_j(s_a^{Lj}-s_a^{-Lj})$ for the
corresponding coarse difference, acting on the fine lattice.
The coefficients satisfy
$2\sum_jjc_j=1$. Define finite Laurent polynomials
\begin{align*}
 M_a&=L^{-1}\sum_{v=0}^{L-1}s_a^v,&
 G_a&=\sum_{j=1}^4c_j\sum_{v=-j}^{j-1}s_a^v,\\
 N_a&=G_a(s_a^L)M_a,&
 K_0&=\prod_aM_aN_a,\\
 K_a&=G_aM_a\prod_{c\ne a}M_cN_c,&
 K_a^{\rm e}&=N_a\prod_{c\ne a}M_cN_c.
\end{align*}
All these operators act on the fine anchor of a local function before
stationary substitution. If $\mathcal I F=F(U_{\rm stat}(V))$, coarse
translation of $\mathcal IF$ is $\mathcal I(s_a^LF)$; fine unit shifts
are not silently identified with coarse shifts.

\begin{lemma}[Exact block identities]
\label{ch17:block-identities}
There is a finite Laurent polynomial $K_a^{\rm t}$ such that
\begin{align}
 K_0\nabla_a&=L^{-1}\nabla_a^{\rm c}K_a,&
 K_0\delta_a&=L^{-1}\delta_a^{\rm c}K_a^{\rm e},&
 K_a-K_0&=\delta_a^{\rm c}K_a^{\rm t}.
 \label{ch17:block-algebra}
\end{align}
The kernels have support of radius $CL$ and absolute coefficient sums
bounded by a numerical constant independent of $L$; their masses are one
except that no mass is prescribed for $K_a^{\rm t}$.
\end{lemma}
\begin{proof}
The identities $G_a\delta_a=\nabla_a$ and
$M_a\delta_a=L^{-1}\delta_a^{\rm c}$ prove the first two assertions
by commutation of Laurent polynomials. For the third, work in the
quotient ring with $s_a^L=1$. There $M_a^2=M_a$,
$M_aG_a=G_a(1)M_a=M_a$ and $G_a(s_a^L)=G_a(1)=1$.
Thus $M_a(G_a-G_a(s_a^L)M_a)$ is zero in that quotient, which is
precisely divisibility of $K_a-K_0$ by $s_a^L-1$.
The original kernels have finite products of a fixed number of bounded
coefficient-sum factors, each of range $CL$. For the quotient, group
coefficients in each residue class modulo $L$. Its numerator has zero
sum in each class, and only a bounded number of nonzero entries in that
class. Division by $s_a^L-1$ is a partial sum along this finite sequence,
so increases its absolute sum by only a bounded factor. This also proves
the support bound for $K_a^{\rm t}$.
\end{proof}

For a fine current with $\operatorname{tr}J=\sum_a\delta_aq_a$, set
\[
 J'_{ab}=L^4\mathcal I(K_aJ_{ab}),\qquad
 q'_a=L^3\mathcal I(K_a^{\rm e}q_a)
                       +L^4\mathcal I(K_a^{\rm t}J_{aa}).
\]
Equation~\eqref{ch17:block-algebra} proves exactly
$\operatorname{tr}J'=\sum_a\delta_a^{\rm c}q'_a$.
Indeed split $K_a=K_0+\delta_a^{\rm c}K_a^{\rm t}$ and then apply
the second identity to $K_0\delta_aq_a$.
Similarly, if $J_{ab}-J_{ba}=\sum_c\delta_ct_{abc}$, the transferred
antisymmetric part is the sum of the transferred old differences and
\[
 \delta_a^{\rm c} L^4\mathcal I(K_a^{\rm t}J_{ab})
 -\delta_b^{\rm c} L^4\mathcal I(K_b^{\rm t}J_{ba}).
\]
In particular symmetry may fail pointwise after blocking, but its failure
and the trace remain exact differences. This conclusion is independent
of estimates on the current's Taylor coefficients.

\subsection{What scalar mass regularization changes}

Consider the same admissible path action used in Chapter~\ref{ch07:matching},
with its left Hessian $H$ and infinitesimal gauge insertion $d$.
At the zero field the completed vector operator is
$H+dd^*=|k|^2I+O(|k|^3)$, and the ghost operator is
$d^*d=|k|^2I+O(|k|^4)$, with the global lattice ellipticity stated there.
Define $\Gamma_\mu$ with the scalar mass $\mu^2I$ as above.
The normalized quadratic response of $\Gamma_\mu$ has the same logarithmic
coefficient as the regulated response in Chapter~\ref{ch07:matching}.
This fact concerns its coefficient integral, before any relation to the
Haar divergence of a current.

To see this directly, retain the actual vector and ghost first and second
background insertions $b,c$ and replace only their denominators by the
scalar-mass denominators. In each channel the logarithmic Hessian is
\[
 \operatorname{tr}\bigl(a_\mu(k)^{-1}c(-p,p;k)
 -a_\mu(k)^{-1}b(-p;k+p)a_\mu(k+p)^{-1}b(p;k)\bigr).
\]
The scalar mass has no background insertion. The actual local Taylor
expansion and vertex estimates of Chapter~\ref{ch07:matching} therefore
apply without the regulator-vertex terms. With $T=|k|+\mu$, the
inverse and its $r$ momentum derivatives cost $CT^{-2-r}$, for every
fixed required $r$. Differentiate $Ra=I$ repeatedly to obtain this bound;
the denominator floor is now simply $a_\mu\ge cT^2I$.
Two external derivatives of the displayed expression cost at most
$CT^{-4}$. For $\mu\le|k|\le1$, comparison with the massless continuum
vertices by the resolvent identity costs
\[
 C\bigl(|k|^{-3}+\mu^2|k|^{-6}\bigr).
\]
The first term is the actual lattice Taylor remainder, and the second
is the scalar-mass denominator difference. Both have bounded integrals
in four dimensions, uniformly in $\mu$. The region $|k|<\mu$ also
has bounded integral, since its volume is $O(\mu^4)$ and its integrand
is $O(\mu^{-4})$. Away from zero the integral is uniformly bounded.
Thus the continuum homogeneous term alone determines the logarithm.
With the twelve-plane coefficient convention of this chapter this gives
\begin{equation}
 R(\mu)=12\kappa_G\log\mu+O(1),\qquad0<\mu\le1.
 \label{ch17:scalar-log-supplier}
\end{equation}
This is the same one-loop normalization already calculated using the
actual action in Chapter~\ref{ch07:matching}. At fixed $\mu>0$ the
integrand is periodic analytic on a common smaller strip, so its torus
sums converge to this integral; no uniform rate as $\mu\downarrow0$
is required. This proof does not differentiate the bounded remainder.

\subsection{Exact integration by parts and the remaining current estimate}

For a smooth vector field $Y$ on a finite product of compact link groups,
Haar integration gives the exact identity
\begin{equation}
 \int\bigl[\operatorname{div}_{\rm Haar}Y
                  -D\mathcal A[Y]\bigr]e^{-\mathcal A}\,dU=0.
 \label{ch17:finite-haar}
\end{equation}
It follows by integrating the divergence of $e^{-\mathcal A}Y$.
With a spectator $F$ the right side becomes
$-\int DF[Y]e^{-\mathcal A}\,dU$. In particular it is zero when
$Y$ acts only on links disjoint from those of $F$. Smooth truncations
can be inserted before this identity is used; their differentiated terms
must be retained as marked remainders.

There is also an exact cancellation available for the inherited quadratic
loop term. Chapter~\ref{ch07:matching} represents its Hessian as a sum
with two lattice differences on each leg. Consequently it annihilates
constant and affine polynomial link fields. The degree-three external
momentum coefficient in its variation by the leading current pairs a
constant leg with a quadratic leg, or two affine legs; both vanish.
This proves the zero trace-functional coefficient of that quadratic
loop contribution. It does not bound the nonlinear current remainder.

The missing quantitative step is now quite specific. One needs the
actual stationary pullback of the current defect, with the gauge-covariant block average
frame correction, in the exponentially weighted order-two, -three and
-four coefficient norms. In particular the order-two pullback must
annihilate its affine--affine and affine--quadratic polynomial moment conditions
\emph{before} estimating the remainder. The usual smooth transport bound
$\nabla_f^rW_a=O(L^{-1-r})$ alone does not assert these cancellations.
After gauge removal the relevant affine field costs $L^{-2}$ and its
quadratic term $L^{-3}$; multiplication by the current rescaling $L^5$
would leave the unacceptable factors $L$ and $1$ from their two products.
Only removal of these exact terms leaves the $L^{-1}$ remainder used
in the proposed contraction. Fourth-order action matching must therefore
be combined with an explicit identity for the current's matched jets,
not merely with a bound on the action's matching error.

The group-coordinate identity needed for this comparison is proved next.
Its proof is independent of the contraction of the actual current
sequence, established in Section~\ref{ch17:quadratic-current}.

\subsection{The group-coordinate contact identity}
\label{ch17:contact-section}

Fix a finite lattice and a background $U$ in the product of link groups.
Use the invariant inner product on its link Lie algebra and the chart
$z\mapsto e^zU$ (edge by edge). Let $S$ be a smooth gauge-invariant action,
$H=D_z^2S(e^zU)|_{z=0}$ its symmetric left-coordinate Hessian, and
$d$ the infinitesimal gauge map at $U$. The displacement $Y$ is assumed
gauge covariant, so its local flow
$\Psi_s(U)=e^{sY(U)}U$ commutes with gauge transformations. It need not
be the exact time-$s$ flow; all following identities use only its
first derivative at zero.

Define the chart velocity and its linear part by
\[
 \mathcal Y_U(z)=\left.\partial_s\right|_0
    \log\bigl(\Psi_s(e^zU)\Psi_s(U)^{-1}\bigr),
 \qquad V=D_z\mathcal Y_U(0).
\]
In particular $\mathcal Y_U(0)=0$. This definition includes the
change of the left frame as the base point moves. Replacing $V$ by an
ordinary derivative of the coefficient function $Y$ would omit that
change. Let
\[
 J_U(v,z)=\left.\partial_s\partial_t\right|_{0}
       \log\bigl(\operatorname{gauge}_{e^{sv}}(e^{tz}U)U^{-1}\bigr),
 \qquad N_{i\alpha}=DS(U)[J_U(e_\alpha,e_i)].
\]
Indices $i$ run through link coordinates and $\alpha$ through gauge
coordinates, using the basis of Section~\ref{ch01:gauge} at each link and site. Differentiating gauge invariance first in $v$, then in $z$,
gives $Hd=-N$.

For $\mu>0$ at a background where the following inverses exist, write
\[
 R=(H+dd^*+\mu^2I)^{-1},\qquad
 G=(d^*d+\mu^2I)^{-1},
\]
\[
 \Gamma=\tfrac12\log\det(H+dd^*+\mu^2I)
                            -\log\det(d^*d+\mu^2I).
\]
Near the zero background the positive denominators of the preceding
subsection make these ordinary real analytic functions. Define a link
vector $Z$ by its action against an arbitrary covector $\ell$:
\begin{align}
 \ell(Z)={}&\tfrac12\sum_{i,j}R_{ij}\,
                   \ell\bigl(D_z^2\mathcal Y_U(0)[e_i,e_j]\bigr)\notag\\
 &-\sum_{i,\alpha}
       [Gd^*(V+V^*)R]_{\alpha i}\,\ell(J_U(e_\alpha,e_i)).
 \label{ch17:contact-Z}
\end{align}
All matrices and derivatives here belong to the same background and
coordinate chart.

\begin{proposition}[Exact contact identity]
\label{ch17:contact-identity}
With these definitions,
\begin{equation}
 \tfrac12\operatorname{tr}\bigl(RD^2(DS[Y])\bigr)
       -\operatorname{div}_{\rm Haar}Y
       =D\Gamma[Y]-\mu^2\operatorname{tr}(RV)+DS[Z].
 \label{ch17:contact}
\end{equation}
No stationarity of $U$ is assumed. In particular the $DS[Z]$ term
cannot be discarded before applying the relevant coefficient functional.
\end{proposition}
\begin{proof}
Let a dot denote differentiation along $\Psi_s(U)$, using the left
chart anew at that base point. The chain rule for
$S(\Psi_s(e^zU))$ gives
\[
 D^2(DS[Y])=\dot H+V^*H+HV+C,
 \qquad C_{ij}=DS[D_z^2\mathcal Y_U(0)[e_i,e_j]].
\]
This is an exact chart calculation, not a replacement of the group by
its tangent space. The differential of $\Psi_s$ in these moving charts
is $I+sV+o(s)$. Exponential charts have Haar density one at their origins,
so its Jacobian at the base point has derivative
$\operatorname{div}_{\rm Haar}Y=\operatorname{tr}V$.

Gauge equivariance of $\Psi_s$ gives $\dot d=Vd$: differentiate
$\Psi_s(\operatorname{gauge}_{e^v}U)
=\operatorname{gauge}_{e^v}\Psi_s(U)$ first in $v$ and then in $s$.
Consequently
\[
 D\Gamma[Y]=\tfrac12\operatorname{tr}(R\dot H)
       +\operatorname{tr}(dd^*RV)
       -\operatorname{tr}\bigl(dGd^*(V+V^*)\bigr).
\]
On the other hand the Hessian chain rule and
$RH=I-Rdd^*-\mu^2R$ yield
\[
 \tfrac12\operatorname{tr}(RD^2(DS[Y]))-\operatorname{tr}V
 =\tfrac12\operatorname{tr}(R\dot H)
       -\operatorname{tr}(Rdd^*V)-\mu^2\operatorname{tr}(RV)
       +\tfrac12\operatorname{tr}(RC).
\]
Finally $Hd=-N$ implies the exact resolvent identity
\[
                  Rd=dG+RNG.
\]
Subtract the two preceding displays. The leftover gauge terms are
\[
 -\operatorname{tr}\bigl((Rdd^*+dd^*R-2dGd^*)V\bigr)
       =-\operatorname{tr}\bigl(Gd^*(V+V^*)RN\bigr).
\]
Inserting the definition of $N$ and the term $\frac12\operatorname{tr}(RC)$
identifies their sum with $DS[Z]$ from \eqref{ch17:contact-Z}.
This proves \eqref{ch17:contact}.
\end{proof}

\subsection{Cancellation of the action-gradient terms}

The functional $\mathsf l$ above extracts the coefficient of $z^2t$
from the diagonal polarization pair at momenta
$p=ze_a$, $q=-ze_a+te_b$, and sums over $a\ne c$ and $b$.
We check why it kills the action variations in
\eqref{ch17:contact}. Work at fixed positive $\mu$ so all kernels
are analytic in momentum near zero. A gauge-invariant, translation- and
hypercubic-invariant scalar action $F$ has quadratic kernel
$c_F(|p|^2I-p\otimes p)+O(|p|^3)$ by the Ward calculation in
Chapter~\ref{ch07:matching}. This applies to both $S$ and $\Gamma$.
Let $Z_b$ be any anchored gauge-covariant displacement, including $Y_b$
or the displacement \eqref{ch17:contact-Z}. Simplicity of the gauge
algebra implies $Z_b(0)=0$. The linearized covariance identity says
$DZ_b(0)d_0=0$.

Use $i$ for its output link direction and $c$ for its input polarization.
Its joint momentum symbol, with output dual variable $p$ and input
variable $q$, therefore has expansion
\[
 (DZ_b)_{ic}(p,q)=\sum_j q_j A_{b,ic}^{\,j}+O((|p|+|q|)^2),
 \qquad A_{b,ic}^{\,j}=-A_{b,ij}^{\,c}.
\]
The factors $i$ from Fourier differentiation may be absorbed into $A$.
There is no term depending only on $p$: the identity obtained by
multiplying by the gradient symbol in $q$ forces it to vanish.
For $a\ne c$, the coefficient of $z^2t$ in the quadratic variation
$DF[Z_b]$ is, apart from the common factor $c_F$,
\[
       A_{b,cc}^{\,b}-2\delta_{ab}A_{b,cc}^{\,a}
                              +\delta_{bc}A_{b,ac}^{\,a}.
\]
This follows by inserting the two kinetic kernels on its two legs.
The first sum over $a\ne c$ gives
$3\sum_{b,c}A_{b,cc}^{\,b}$. The second gives minus twice that sum,
since $A_{c,cc}^{\,c}=0$. The last gives minus the same sum after
using $A_{c,ac}^{\,a}=-A_{c,aa}^{\,c}$ and exchanging $a,c$.
Thus
\begin{equation}
                  \mathsf l(DF[Z])=0.
 \label{ch17:gradient-cancellation}
\end{equation}
Higher momentum orders cannot contribute to the coefficient of degree
three. Gauge covariance of \eqref{ch17:contact-Z} follows by transforming
its chart tensors, inverse operators and contracted indices together.
The mass $\mu^2I$ is unchanged by these orthogonal gauge frame changes.
The calculation therefore applies to both action-gradient terms of the
exact contact identity. It is a finite-momentum cancellation, not a
statement that the corresponding functions vanish pointwise.

\subsection{Power counting the regulator remainder}

The remaining estimate has a transparent sufficient numerator condition.
Let $W^{(1)},W^{(0)}$ be the vector and ghost remainder operators obtained
by subtracting the leading curvature displacement from $V$; in the
notation of the current calculation these are
\[
 W_b^{(1)}=\operatorname{sym}(V_b-dM_b-D_b)+\tfrac12T_b^{(1)},
 \qquad W_b^{(0)}=\operatorname{sym}(M_bd)-\tfrac12T_b^{(0)}.
\]
Here $\chi$ is the indicator of the anchored lattice site,
$(M_bh)(x)=\chi(x)h_b(x)$,
$(D_bh)_a(x)=\delta_a\chi(x)h_b(x)$, and
$(T_b^{(r)}h)(x)=\delta_b\chi(x)h(x)$ on $r$-forms.
The full lattice shifts and chart tensors must be retained when
extracting their coefficient; their continuum symbols only identify
the leading terms that cancel.

For that leading calculation the gauge derivative is $d=-\partial$
and the curvature displacement is
$(V_bh)_a=\chi(\partial_bh_a-\partial_ah_b)$.
Then $dM_b+D_b=-\chi\partial_ah_b$ and
$\operatorname{sym}(\chi\partial_b)=-\frac12\partial_b\chi$.
Thus both $W^{(r)}$ have zero first-order principal symbol. Covariantizing
these expressions gives the same cancellation at constant background
to first field order: the skew-adjoint multiplication by
$\operatorname{ad}A_b$ has zero symmetric part, and the remaining
$\operatorname{ad}A_a h_b$ term cancels the corresponding term of $dM_b$.
Accordingly their zero-, one- and two-field symbols have orders
$2,1,0$, respectively. A lattice/current sequence having the stated
matched free jets inherits precisely these vanishings in its limit.
Weighted coefficient bounds control the Taylor remainders.

Here is the resulting estimate without suppressed differentiation orders.
Set $z=|q|+\mu$. For $n\le2$ background derivatives and $r\le3$ external
momentum derivatives the inverse estimates are
$Cz^{-2-n-r}$. The remainder symbols have bounds
$Cz^{2-n-r}$. The quadratic-field coefficient of
$\operatorname{tr}(RW^{(1)})+\operatorname{tr}(GW^{(0)})$, after the
three momentum derivatives defining $\mathsf l$, is therefore bounded
by a sum of terms
\[
 Cz^{-2-n-r}\,z^{2-(2-n)-(3-r)}=Cz^{-5}.
\]
In four dimensions
$\int_{|q|\le1}(|q|+\mu)^{-5}\,dq\le C/\mu$.
The region away from zero is bounded. Hence the regulator factor
$\mu^2$ gives a remainder at most $C\mu$.
For approximate matched jets, first fix $\mu\ge\delta>0$; their
errors then tend to zero uniformly on that regulator interval.
Only afterwards send $\delta$ to zero. This is the order required
by the logarithmic mean-value argument.

\subsection{The corrected coefficient and order of limits}
\label{ch17:regulator-completion}

Retain the lattice correction in $\mathsf h$. For the same unit vector
$u\in\gLie$, set
\[
 \mathsf c(F)=\tfrac12\operatorname{Re}\sum_{a\ne c}[z^2]
 \widehat{F_c}_{cu,cu}(ze_a,-ze_a),
 \qquad \mathsf h=\mathsf l-\mathsf c.
\]
The quadratic variation $DF[Z]$ in the preceding calculation starts at
momentum degree three. Thus $\mathsf c(DF[Z])=0$, in addition to
\eqref{ch17:gradient-cancellation}. Both functionals therefore annihilate
the action-gradient terms of \eqref{ch17:contact}.

Suppose now that the actual current sequence has the following
properties at the fixed output level: uniform exponentially weighted
order-two through order-four coefficient bounds; convergence of the
needed low polynomial jets to the curvature current; the exact trace
and antisymmetry difference identities; and defects tending to zero
through order four in those weighted norms in
\[
                 DS[Y_b]\equiv_4\sum_a\nabla_aJ_{ab}-\mathcal E_b.
\]
The notation $\equiv_4$ refers to equality of field derivatives through
order four at the zero background. It suffices because the contact
expression uses two internal and two external field derivatives.
At fixed $\mu\ge\delta>0$, contracted defect errors tend to zero.
Translation covariance makes contraction by the inverse $R$ commute
with translating the anchor. The $\mathsf l$ coefficient of the
contracted current divergence selects $a=b$ at its sole linear
 total-momentum derivative. The sum is the contracted trace of $J$,
which has a further total-momentum zero because it is a forward
difference. Its coefficient is zero. For $\mathsf c$, the total
momentum is zero from the outset and the divergence vanishes directly.
The exact contact identity consequently yields
\[
 \mathfrak h(Y)=\mu^2\mathsf h\bigl(\operatorname{tr}(RV_b)\bigr)+o(1),
\]
uniformly on every fixed positive regulator interval. All operators here
are functions of the background before the two field derivatives are
extracted.

The definitions of $W^{(1)},W^{(0)}$ give an exact trace decomposition.
Using the same link and site bases induced by Section~\ref{ch01:gauge}, apply $Rd=dG+RNG$ to obtain
\begin{align*}
 \operatorname{tr}(RV_b)
  ={}&\operatorname{tr}\bigl(R(D_b-\tfrac12T_b^{(1)})\bigr)
      +\tfrac12\operatorname{tr}(GT_b^{(0)})+DS[Z_b']
      +\mathrm{Rem}_b,\\
 \mathrm{Rem}_b={}&\operatorname{tr}(RW_b^{(1)})
                                   +\operatorname{tr}(GW_b^{(0)}),\\
 \ell(Z_b')={}&\sum_{i,\alpha}(GM_bR)_{\alpha i}
                                       \ell(J_U(e_\alpha,e_i)).
\end{align*}
For example, the part containing $dM_b$ becomes
$\operatorname{tr}(GM_bd)+\operatorname{tr}(RNGM_b)$;
its symmetric ghost part is precisely
$\frac12\operatorname{tr}(GT_b^{(0)})+\operatorname{tr}(GW_b^{(0)})$,
and its remaining term is $DS[Z_b']$. This also fixes the sign of
the gauge correction.

For the three displayed principal terms the Fourier transform of
$\delta_b\chi$ supplies $e^{-it}-1=-it+O(t^2)$ at total momentum
$te_b$. With the factor $i$ in $\mathsf l$, the $D_b$ terms sum to
$\mathsf q(\operatorname{tr}R)$: the selected derivative has $a=b$
and the component trace then sums its diagonal entries. Each of the
four $T_b^{(1)}$ terms gives $\mathsf q(\operatorname{tr}R)$, and each
$T_b^{(0)}$ term gives $\mathsf q(\operatorname{tr}G)$.
The volume normalization is exactly the translation average of the
single anchor. Their total is therefore
\[
           -\mathsf q(\operatorname{tr}R)
                              +2\mathsf q(\operatorname{tr}G).
\]
Their $\mathsf c$ coefficients are zero, since $\delta_b\chi$ has
zero symbol at total momentum zero. Multiplication by $\mu^2$ gives
$-\mu\partial_\mu\mathsf q(\Gamma)$ by differentiation of the two
determinants. Thus, with the lattice correction fully retained,
\begin{equation}
 \mathfrak h(Y)=\mathcal D_\Lambda(\mu)
                  +\mu^2\mathsf h(\mathrm{Rem})+o(1).
 \label{ch17:corrected-regulated-ward}
\end{equation}
The action-gradient term $DS[Z_b']$ is annihilated by both parts of
$\mathsf h$ as proved above.

The $\mathsf l$ remainder was bounded by $C\mu$ after multiplication
by $\mu^2$. The correction $\mathsf c$ uses only two external momentum
derivatives, so the same product rule costs $Cz^{-4}$. Its integral
is at most $C(1+|\log\mu|)$ in four dimensions, and
$\mu^2(1+|\log\mu|)\le C\mu$ for $0<\mu\le1$.
Consequently \eqref{ch17:corrected-regulated-ward} proves the regulator
estimate, including the anchored correction, from the stated current
hypotheses.

Finally combine it with \eqref{ch17:scalar-log-supplier}. On a fixed
interval $[u-A,u]$ in the variable $\log\mu$, the bounded error in
$R(\mu)-12\kappa_G\log\mu$ implies that some regulator has
$|\mathcal D_\Lambda(\mu)+12\kappa_G|\le(2C+o(1))/A$ after the
period is large. The remainder is at most $Ce^u+o(1)$ on that interval.
First take the current and period limits there, next $u\to-\infty$,
and finally $A\to\infty$. Since $\mathfrak h(Y)$ is independent of
$\mu$, this proves
\[
                         \mathfrak h(Y)\longrightarrow-12\kappa_G.
\]
The group-coordinate contact identity, coefficient cancellation and
regulator estimate are now established. Their application still requires
the actual transferred current sequence to satisfy the weighted bounds
and matched-jet convergence stated at the start of this subsection;
those conditions are not consequences of the contact algebra.

% !TEX root = ../monograph.tex
% Current defect and uniform transfer bounds.
\section{Current conservation and uniform transfer bounds}
\label{ch17:quadratic-current}

The potentially dangerous quadratic defect has an exact cancellation
that should be imposed before estimating its norm. We prove preservation
of that cancellation for the linearization of the critical configuration. This removes
an additive interpretation of the order-two matching error. We then control the nonlinear
jets of the critical configuration and the current itself.

Fix a stage and an admissible action. Let $\mathcal E_b$ be a
gauge-invariant current defect with a base point, with zero constant and linear
field coefficients. Its Hessian at zero has an exponentially local
joint Fourier symbol $E_b(p,q)$, with two independent field momenta.
Gauge invariance gives a Ward identity on each leg. We say that its
quadratic defect has order at least four when
\begin{equation}
                    E_b(p,q)=O((|p|+|q|)^4)
                              \quad\hbox{near }(0,0).
 \label{ch17:defect-four}
\end{equation}
Equivalently, every Taylor coefficient of $E_b$ of total degree less
than four in $(p,q)$ vanishes. This order of vanishing in momentum
implies the cancellations for pairs of affine fields and for an affine
field paired with a quadratic field that are used in the current estimate.

For the initial Wilson current, \eqref{ch17:defect-four} follows from
the leading continuum identity. Its linearized curvature is $dA$;
the leading quadratic action is the curvature square; and the centered
difference $\nabla$ has linear symbol $ip$. The quadratic Taylor
coefficient of
$\sum_a\nabla_aJ_{ab}-DS[Y_b]$ therefore has leading spatial
momentum degree three. At that degree it is exactly
\[
 \sum_a\partial_a\left(\tfrac14\delta_{ab}\sum_{c,d}\|F_{cd}\|_{\gLie}^2
                         -\sum_c \langle F_{ac},F_{bc}\rangle_{\gLie}\right)
       +\sum_{a,c} \langle F_{bc},\partial_aF_{ac}\rangle_{\gLie}=0.
\]
The identity uses the linear Bianchi equation. Lower spatial degrees vanish because each curvature
contains a difference and the action variation contains one more.
All remaining lattice Taylor terms have degree at least four.
The same reasoning applies to any initial path action with the
normalized Taylor polynomial of scaling dimension four already established in
Chapter~\ref{ch07:matching}.

\subsection{Preservation of the order of vanishing}

The first derivative of the nonlinear critical-configuration map, from coarse
to fine link variables with block ratio $L$, is the forest representative
$W_F$ of \eqref{ch05:stationary-tangent}; its higher derivatives
cannot enter the quadratic pullback because the defect has zero first
field derivative. Since $W_F-W_a$ is a fine gradient, the Ward identities
of $E_b$ allow us to compute with the smooth representative $W=W_a$.

For the exact Taylor-order statement we pass once more to a gauge-equivalent representative, the linearization $W_s$ built with the linear averaging operator $M$ (Section~\ref{ch05:actual-tangent}), and test Hessians on polynomial fields.\mgubi{Issue 8: polynomial fields live on $\ZZ^4$, the critical configuration of Chapter~\ref{ch05:gaussian} on a torus. One sentence is missing: the Hessian on the torus is the periodization of the one on $\ZZ^4$. Vanishing on a constant leg (used below) is false on a torus. Counterparts: existence/\allowbreak quadratic-pole-transport.\allowbreak tex; v2 I-section.\allowbreak tex.} By \eqref{ch05:corrected-transport}, $W_a-W_s=-d_{0,f}ERW_s$; since $E$ is bounded and $RW_s$ has a bounded coarse kernel, each column of this difference is the gradient of a bounded fine potential. Let $F$ be a gauge-invariant local function of the fine field with base point $o$ and zero first field derivative, and write $B(a,a')=D^2F(0)[a,a']$, whose kernel decays exponentially in the distances of both legs from $o$.
The Ward identity on each leg, $B(d_{0,f}\phi,a')=B(a',d_{0,f}\phi)=0$, holds for finitely supported $\phi$ and, by this decay, also for bounded $\phi$ and fields $a'$ of polynomial growth: replace $\phi$ by its restriction to a large ball about $o$. Call a one-form with values in $\gLie$ a polynomial field of degree $r$ if its components in a basis of $\gLie$ are polynomials of degree at most $r$ in the coordinates. $B$ converges absolutely on pairs of polynomial fields, and the Taylor coefficient of multidegree $(\alpha,\alpha')$ at zero of the joint symbol of $B$, expanded about $o$, is a nonzero multiple of the value of $B$ on the fields $(x-o)^\alpha e$ and $(x'-o)^{\alpha'}e'$ with $e,e'$ constant. The symbol being analytic, \eqref{ch17:defect-four} for $F$ is therefore equivalent to
\begin{equation}
 B(P,P')=0\qquad\hbox{for polynomial fields $P,P'$ of degrees $r,r'$ with }r+r'\le3.
 \label{ch17:polynomial-four}
\end{equation}
This condition does not depend on the center of the coordinates, since a translate of a polynomial has the same degree. Hence it holds for every fine translate of $E_b$, and so for $F=K_0E_b$, where $K_0$ is the finite combination of fine translates of the base point from Lemma~\ref{ch17:block-identities}; this $F$ is again gauge invariant with zero first field derivative, and its kernel decays exponentially because $K_0$ has range $CL$.
Let $G=\mathcal IF$ be its composition with the critical configuration. For finitely supported coarse fields $a,a'$ the preceding paragraph gives $D^2G(0)[a,a']=B(W_Fa,W_Fa')=B(W_aa,W_aa')$, where the potential of the fine gradient $W_Fa-W_aa$ is bounded because it is a sum of values of $W_aa$ along the spanning tree of one block; and the Ward identity applied to the bounded potentials above gives $D^2G(0)[a,a']=B(W_sa,W_sa')$. The kernel of $W_s$ is $O(L^{-1})$ with exponential coarse decay (Section~\ref{ch05:transports}).
Hence the coarse Hessian kernel of $G$ decays exponentially in the distances of both legs from the coarse block of $o$, its joint symbol is analytic, and for coarse polynomial fields $P,P'$ absolute convergence gives $D^2G(0)[P,P']=B(W_sP,W_sP')$. As for $F$, \eqref{ch17:defect-four} for $G$ is equivalent to the vanishing of these values when $r+r'\le3$.
By \eqref{ch17:polynomial-response}, whose derivation in the proof of Proposition~\ref{ch17:spatial-moments} uses only the formulas of Chapter~\ref{ch05:gaussian}, $W_s$ maps a coarse polynomial field $P$ of degree $r<j$ to a fine polynomial field of degree at most $r$; its homogeneous part of degree $r$ is $L^{-r-1}$ times that of $P$, read in fine coordinates. Here $j$ is the smoothing order of the averaging weight, and $j\ge5$, since \eqref{ch05:actual-bridge} is used below with three fine differences.\mgubi{Issue 8: $j\ge5$ is stronger than the standing $j\ge4$ of Chapter~\ref{ch05:gaussian} and is not entered in the schedule of Chapter~\ref{ch10:trajectory}; nor are the requirements on the smoothing order in Section~\ref{ch11:marked-step} and in the proof of Proposition~\ref{ch17:spatial-moments}.}
If $r+r'\le3$, the two images therefore have degrees with sum at most three, and \eqref{ch17:polynomial-four} for $F$ gives $D^2G(0)[P,P']=0$. Thus $G$, and the pulled-back defect $L^5G$, satisfy \eqref{ch17:defect-four} on the coarse lattice.
The same computation gives the exact rescaling of the first coefficients that need not vanish. Suppose that $B$ vanishes on all pairs of polynomial fields with $r+r'<n$, and let $P,P'$ be homogeneous of degrees $r,r'$ with $r+r'=n<j$. Then the parts of $W_sP$ and $W_sP'$ of lower degree contribute nothing, and $D^2G(0)[P,P']=L^{-n-2}B(P,P')$, with $P,P'$ read in fine coordinates on the right.
No alias decomposition of the linearization and no projection off the fine gradient direction is used: fine gradients are removed by the Ward identity before any polynomial is inserted, and only the polynomial moments of $W_s$ enter.
No bound uniform in $L$, and no inverse of a kinetic symbol at a nonzero alias, is
needed for this exact Taylor-order statement.

The base-point average $K_0$, the average over translates of the base point, from
Lemma~\ref{ch17:block-identities}, is performed before composition with the critical
configuration. Its symbol has a double zero at every nonzero
alias: at least one coordinate of $l$ is nonzero, and the factor
$M_a^2G_a(s_a^L)$ has two factors of the average $M_a$ there.
The argument above does not use this double zero. Of $K_0$ it uses only
that it is a finite combination of fine translates of the base point; the
behaviour of the linearization enters only through \eqref{ch17:polynomial-response}.

Thus the order of vanishing is preserved by an exact computation. The
pulled-back Hessian, evaluated on pairs of coarse polynomial fields of total degree
at most three, equals the Hessian of $K_0\mathcal E_b$ evaluated on their images
under $W_s$. These images are fine polynomial fields of no larger total degree,
on which that Hessian vanishes by \eqref{ch17:polynomial-four}. No estimate,
no limit and no bound uniform in $L$ enters. Hence $L^5\mathcal I(K_0\mathcal E_b)$
satisfies \eqref{ch17:defect-four} on the coarse lattice. This is an
exact identity and holds on the same history.

In particular the quadratic jet is not merely close to having the
required affine cancellations: it has them exactly, provided the
transfer of the current uses the exact pulled-back defect
$\mathcal E'_b=L^5\mathcal I(K_0\mathcal E_b)$ through the relevant
field order. If an implementation redefines that defect by action
matching, equality of these quadratic jets must be retained; a bounded
additive matching error would not preserve this argument.

\subsection{A homogeneous bound for the quadratic channel}

The analytic division argument of Section~\ref{ch07:division} can be
applied to each of the two independent momentum legs of $E_b$.
The Ward identities produce one curvature difference on each leg.
Their homotopy divisions lower Taylor order by one each, so
\eqref{ch17:defect-four} leaves a coefficient of total order at least
two. Coordinatewise Taylor division in the remaining eight variables
factors two additional differences, distributed between the legs.
On a fixed smaller coefficient width these operations are bounded by
constants depending only on the stage and the decrease of the width. Thus $E_b$
is a sum of terms with either one and three differences, or two and
two differences, on its two field legs. The coefficient sum is bounded
by $C_j\|E_b\|$ in the original weighted norm.

Recall that $W=W_a$ is the smooth representative; the quadratic Ward
identities annihilate its gradient difference from $W_F$.
Its fine-difference bounds are
\[
                  \|\nabla_f^rW_a\|\le C_{j,r}L^{-1-r},
                                  \qquad r\le3.
\]
Consequently either allocation of four differences contributes at most
$C_jL^{-6}$ for the two legs. The base-point average $K_0$ has bounded
absolute coefficient sum, not a factor $L^4$; it is an average rather
than an unnormalized sum. Translated coefficient weights are absorbed
by a fixed narrower coarse width, exactly as in the inherited quadratic
action estimate of Chapter~\ref{ch07:matching}. Multiplication by the
current normalization $L^5$ proves
\begin{equation}
             \|D^2\mathcal E'_b(0)\|_{\rm out}
                  \le \frac{C_j}{L}\,
                            \|D^2\mathcal E_b(0)\|_{\rm in}.
 \label{ch17:quadratic-current-contraction}
\end{equation}
This is homogeneous in the defect. No constant inhomogeneous term has
been introduced. Dividing the coefficients by $(\log\widehat\beta)^2$
only adds the adjacent logarithmic ratio, which is bounded and tends
to one on a fixed-stage tail. The selected block ratio makes this
quadratic channel a strict contraction.

\subsection{Ratio-dependent bounds for the nonlinear terms}

The nonlinear jets of the critical configuration enter only below the diagonal in field
order. Their constants may depend on $L$, because the lower bound on the reference coupling is
chosen after the ratio. What must remain uniform before $L$ is the
coarse coefficient width on which these constants are measured.
We obtain that width from the gauge-fixed inverse of
Chapter~\ref{ch07:matching}.

Write $z(v)$ for the fluctuation coordinates of the critical configuration, and
$\mathscr S(v,z)$ for the same joint classical action as there.
Let $T u=(u,Dz(0)u)$, and let $\mathscr H$ be its Hessian in the fluctuation coordinates.
On the slice of the block axial gauge implicit differentiation gives
\begin{align}
 \mathscr H z_2[u,v]
    &=-D^3\mathscr S[T u,T v,\,\cdot\,],\notag\\
 \mathscr H z_3[u,v,w]
    &=-D^4\mathscr S[T u,T v,T w,\,\cdot\,]
      -\sum_{\{ij\}\{k\}}
             D^3\mathscr S[(0,z_2[u_i,u_j]),T u_k,\,\cdot\,].
 \label{ch17:stationary-higher-jets}
\end{align}
There are three terms in the sum. The right sides are fluctuation
covectors, with all mixed derivatives in the averaging operation and in the local gauge transformation included.

These formulas have period-independent bounds $C_{j,L}$ on one
fixed smaller coarse Fourier strip. To check the width statement in this zero-background computation, take the forest of
Section~\ref{ch05:actual-coordinates} to consist of the coarse translates of one reference tree. This fixes a free choice for the present computation only; the average over forests of Section~\ref{ch09:forest-average} is not used here, and independence of the forest is not claimed for the forest-dependent jets $z_2,z_3$ themselves. The bounds are used below only for pullbacks of gauge-invariant functions along the critical-configuration map; that these do not depend on the forest (the slices of the block axial gauge for two forests differ by a gauge transformation equal to one at the block roots) is used here and is not proved in this chapter.\mgubi{Issue 7: v3 has no such lemma. Lean has it near the zero configuration (YM.\allowbreak Work.\allowbreak F05.\allowbreak answer, from graph\_switch, fine\_switch, Forest.\allowbreak fold\_root; the proofs of the first two were not read). For the independent forests of Chapter~\ref{ch09:renormalization}, $\Pi_F$ does not commute with coarse translations.} Let $\Pi_F$ be projection onto the slice of the block axial gauge along the joint gauge
directions; it then commutes with coarse translations. At zero it is computed by finite sums along the block tree
and the local gauge transformation at the root; its coefficients and those of $\Pi_F^*$ have norm
bounded by a finite $C_{j,L}$, but bounded range in coarse blocks.
The dual $\Pi_F^*$ extends a slice covector by zero on gauge directions.
For such a covector the gauge-fixed inverse $\mathscr C^{-1}$ of
Chapter~\ref{ch07:matching} solves the Hessian equation modulo gauge.
Thus the slice inverse is
\[
                       \Pi_F\mathscr C^{-1}\Pi_F^*.
\]
The gauge-fixed inverse has a common positive coarse width chosen before
$L$; multiplication by the finite-range projections does not destroy it.
The path vertices and derivatives of the averaging operation also have a common coarse
width, with finite $C_{j,L}$ norms. Restrict the external strips by a
fixed factor so all intermediate momentum sums lie in this inverse
strip. Formula~\eqref{ch17:stationary-higher-jets} then proves the
asserted bounds for $z_2,z_3$. It does not assert bounds uniform in $L$
for their norms.

For every gauge-invariant function with a base point, evaluate this critical
configuration in a radial gauge at the base point. At first order, the gauge-corrected
link at fine displacement $x$ is a sum of the curvature of the linearization over
coordinate rectangles of total area at most $C(1+|x|)$. The smooth
curvature bound for $W_a$ gives
\[
 |W^{\rm rad}(x,y)|
 \le C_jL^{-2}(1+|x|)
                   e^{C|x|/L}e^{-c_jd(0,y)}.
\]
The polynomial factor and fine exponential are absorbed by the weighted
local coefficient norm. A common strip reduction handles the coarse
weight. Higher derivatives of this gauge change are finite products
along the same contours; the fixed-order factors are polynomial in contour
length. Their period-independent norms are therefore still finite
$C_{j,L}$ on the same smaller coarse width.

The raw higher coefficients $D^3\mathcal E,D^4\mathcal E$ need not
separately satisfy linear Ward identities. We use gauge invariance of
the complete function to make this radial substitution first; its
nonlinear derivatives multiply only lower field derivatives of the
defect.

Now differentiate the exact pullback $\mathcal E'=L^5\mathcal I(K_0\mathcal E)$.
In field order $r$ the term involving $D^r\mathcal E$ has only $r$
factors of the linearization of the critical configuration. For $r=3,4$, radial gauge bounds it by
$C_jL^{5-2r}\|D^r\mathcal E\|$. The quadratic diagonal has the
stronger estimate \eqref{ch17:quadratic-current-contraction}.
Every other chain-rule term involves $D^s\mathcal E$ with $s<r$
and higher jets of the critical configuration or of the gauge change. Hence, for $r=2,3,4$,
\begin{equation}
 e'_r\le a_{j,r}(L)e_r+\sum_{2\le s<r}C_{j,L,r,s}e_s,
 \qquad
 a_{j,2}(L),a_{j,3}(L)\le C_j/L,
 \quad a_{j,4}(L)\le C_j/L^3.
 \label{ch17:unnormalized-defect}
\end{equation}
Here $e_r$ is the weighted norm of the field derivative of order $r$.
The estimates are homogeneous in the defect, because every term in the
composition contains one of its derivatives. There is no action-independent
inhomogeneous term on the right.

Put $\bar e_{m,r}=e_{m,r}/(\log\widehat\beta_m)^r$. For a fine-to-coarse
step from $m+1$ to $m$, the coefficient below the diagonal becomes
\[
 C_{j,L,r,s}
       \frac{(\log\widehat\beta_{m+1})^s}
            {(\log\widehat\beta_m)^r},\qquad s<r.
\]
At fixed $j,L$ this tends to zero as the lower bound on the reference coupling increases.
The diagonal logarithmic ratio tends to one. Choose $L$ first so the
three diagonal coefficients are at most $1/8$, then choose the lower bound on the reference coupling
so each row sum, including the below-diagonal terms, is at most $1/2$.
The finitely many adjacent-stage cases are included before each choice.
This proves a homogeneous contraction of the normalized defect vector.
Starting with the Wilson defect, whose normalized norm is bounded, its
value after $n$ transfers is at most $C2^{-n}$. At a fixed scale
the logarithmic normalization is fixed, so the unnormalized
defect derivatives tend to zero as well.

\subsection{Why action matching introduces no inhomogeneous term}

We specify the map on vector fields, including its sign. At the critical
configuration and coarse local gauge transformation, let $\widetilde Q(U,h)$ be the
averaging operation with its chosen local gauge transformation. Its left-trivialized differential is
\[
 (R_f\eta)_e=\left.\frac d{ds}\right|_0
 \log\bigl(\widetilde Q(e^{s\eta}U,h)_e
                              \widetilde Q(U,h)_e^{-1}\bigr).
\]
Let $b_{y,f}$ be the fixed contour coefficients of the local gauge transformation, defined by
$\sum_f b_{y,f}\eta_f=\tau_y(0;\eta)-\sum_rw_r^{\rm av}\tau_y(r;\eta)$,
where $\tau$ is the signed path average in
\eqref{ch05:frame-linearization}. With the fixed contour $\gamma_{y,f}$
from coarse root $Ly$ to the initial vertex of the fine link $f$, set
\[
 (R_0\eta)(y)=\sum_f b_{y,f}
             \operatorname{Ad}(h(y)U(\gamma_{y,f}))\eta_f,
 \qquad (d_V\phi)_e=\phi(e^-)-\operatorname{Ad}(V_e)\phi(e^+).
\]
At coarse base point zero the transferred coarse vector field is
\begin{equation}
 Y'_b=L^5\sum_xK_0(x)(R_f+d_VR_0)
                              Y_b^x(U_{\rm stat}(V)).
 \label{ch17:actual-current-map}
\end{equation}
Translations give the other base points. The sum is over fine base points;
the output is a vector on coarse links. At the identity $R_0=R$ of
\eqref{ch05:actual-Q}, $R_f=Q=M+d_{0,c}R$, and
$d_V=-d_{0,c}$. In this convention the plus sign in
\eqref{ch17:actual-current-map} cancels the root correction exactly:
\begin{equation}
                         (R_f+d_VR_0)|_{0}=M.
 \label{ch17:flat-current-map}
\end{equation}
The smooth linear averaging operator, rather than a rough row along the root contours,
therefore controls its equal-order Taylor terms. All derivatives of
$R_f+d_VR_0$ have finite $C_{j,L}$ bounds on the common coarse strip.
Their fixed contours have length $O_j(L)$ and remain within bounded
coarse distance; differentiation adds only finite contour-length factors.

For the exact action at the critical configuration the current identity is preserved before
any estimates. The envelope theorem differentiates the stationary value
only through the explicit coarse field. In the penalty term, simultaneous
left multiplication of the averaging operation and the coarse link leaves the
penalty term unchanged. Thus a fine vector field $Y$ induces the coarse
vector field $R_fY$; its root correction $d_VR_0Y$ has zero action
variation by coarse gauge invariance. Together with
\eqref{ch17:block-algebra} this gives
\[
 DS_{\rm stat}[Y'_b]
  =L^5\mathcal I K_0DS[Y_b],\qquad
 \sum_a\nabla_a^{\rm c}J'_{ab}
  =L^5\mathcal I K_0\sum_a\nabla_aJ_{ab}.
\]
Joint stationarity in the fine field and the local gauge transformation is used here, as in the gauge-fixed
quotient calculation of Chapter~\ref{ch07:matching}. Matching replaces
$S_{\rm stat}$ by an action with identical Taylor coefficients through
field order four. Since $Y'(0)=0$, its action variation still agrees
through order four. Therefore the displayed defect pullback is exact
at every field order used in \eqref{ch17:unnormalized-defect}.
A separate additive $C/(\log\widehat\beta)^2$ matching error is neither
needed nor permitted in this argument.

\subsection{Uniform current and difference-primitive bounds}

Recall that $J_{ab}$ is the lattice stress current and $Y_b$ its
associated gauge-covariant vector field. The functions $q_a$ and
$t_{abc}$ are the forward-difference primitives of its trace and
antisymmetric part, respectively, as defined in
Section~\ref{ch17:current-suppliers}. They are distinct from the
small-field threshold $t$ used in the norms with insertions.

For completeness we describe the accompanying bounds needed in the
regulator criterion. Subtract fixed local reference stress and curvature
vector fields, whose leading Taylor coefficients by scaling dimension are those used initially.
The reference functions have finite stencils and fixed coefficient bounds.\mgubi{Issue 8: gauge invariance of the reference stress is not stated; the division by two curls below uses the Ward identity on each leg.}
The lowest jets of the transferred functions agree with them exactly.
For the stress this follows from the exact rescaling at the end of the proof of preservation of the order of vanishing,
applied with $n=2$ to the Hessian $B$ of $K_aJ_{ab}$. Its values on pairs of polynomial fields with $r+r'<2$ vanish: one leg is then constant,
a constant field is the gradient of a linear potential, and the truncation argument given there, using the exponential decay of the kernel, extends the Ward identity to such potentials. If
$r+r'=2$, the rescaling gives $D^2\mathcal I(K_aJ_{ab})(0)[P,P']=L^{-4}B(P,P')$, which the factor $L^4$ in $J'_{ab}$ compensates.
Moreover $K_a$ is a finite combination of fine translates with coefficient sum one, and a translate of a homogeneous polynomial
differs from it by terms of lower degree, on which the Hessian of $J_{ab}$ vanishes; so $B(P,P')$ is the corresponding value for $J_{ab}$.
Consequently the quadratic coefficients of scaling dimension below four remain zero, and the quadratic
coefficient of scaling dimension four is unchanged.

For the vector field one additionally uses the column sum of the linearized averaging operation.
By \eqref{ch17:flat-current-map}, the corrected differential is exactly
$M$ at zero. The alias formula
for $M$ annihilates every nonzero alias at zero coarse momentum, so its
fine column sum in a fixed direction is constant. Its row sum is $L$;
there are $L^4$ positions within a block. The column sum is therefore $L^{-3}$.
This is an identity for the corrected differential, with no residual
root-gradient term. Since the leading fine curvature of $W_a$ is $L^{-2}$ times
the coarse curvature, the factors in the transfer of the current multiply to
$L^5L^{-3}L^{-2}=1$. This preserves the linear curvature jet of $Y$.
Gauge covariance then fixes its accompanying quadratic commutator term:
the covariant polynomial of scaling dimension two in the transferred
vector field is again the curvature.

The stress remainder's quadratic symbol now has order at least three.
Two-leg curl division leaves one extra difference, so its diagonal
transfer is $C_jL^4L^{-5}\le C_j/L$. Higher stress orders have
radial-gauge bounds $C_jL^{4-2r}$.
For $Y$, its zero-order corrected differential $M$ has row sum at most
$C_jL$, while $\|K_0\|_\infty\le C_jL^{-4}$.
To see the resulting local gain, write a translated current coefficient
as $y(f-x;z_1-x,\ldots,z_r-x)$, with output link $f$ and base point $x$.
For each fixed coarse output edge $e$ its absolute sum is bounded by
\[
 \|K_0\|_\infty\sum_f|M_{ef}|
        \sum_x|y(f-x;\,\cdot\,)|
                   \le C_jL^{-3}\|y\|.
\]
Only a bounded number of coarse output edges are within the patches of the averaging
operation and of the base-point average; exponential coefficient tails control the rest.
This proves the same gain for the full local output-link norm.
The linear vector-field remainder has a curvature difference and one
additional difference. If the extra difference acts on its field leg,
it contributes a factor $L^{-1}$ through $W_a$. If it acts on the fine output index,
summation by parts moves it onto $M_{ef}$; the identity for the linear averaging operator gives
$\sum_f|\nabla_fM_{ef}|\le C_j$, again gaining $L^{-1}$ relative to
its row sum $C_jL$. Thus its diagonal factor is
$C_jL^5L^{-3}L^{-3}=C_j/L$.
For field order $r\ge2$ its diagonal factor is $C_jL^{2-2r}$.
Derivatives of the averaging operation or of the correction from its local gauge transformation multiply lower
field orders of $Y$, so their fixed-$L$ constants are absorbed by the
same logarithmic normalization as above. Evaluating in a local radial
gauge is legitimate for this covariant vector field: transform the local gauge
transformation at its output and the differential of the averaging operation together. The additional
derivatives of the local gauge transformation again occur only below the field-order diagonal.

The trace and antisymmetry primitives start at field order two. Their
old terms carry $L^3$ and hence gain $C_jL^{3-2r}\le C_j/L$.
Their new terms are the explicit $L^4K_a^{\rm t}J$ terms proved in
Lemma~\ref{ch17:block-identities}; the equal-order inhomogeneous term from $J$
is bounded independently of $L$. Choose a fixed smaller weight for
these primitive norms so this one-way inhomogeneous term has norm at most $1/8$.
The other diagonal terms are then made small by $L$, and all lower-order
terms by the lower bound on the reference coupling. Reference functions contribute finite
$C_{j,L}$ inhomogeneous terms, divided by positive powers of $\log\widehat\beta$;
a sufficiently large lower bound on the reference coupling can make these terms at most $1/4$.

Thus a weighted sum of the normalized current and primitive norms obeys
$X_m\le\frac12X_{m+1}+\frac14$. The initial Wilson current and zero
primitives have bounded normalized norms. Uniform bounds follow throughout
the history, and in particular at each fixed scale. At that scale
these are the unnormalized finite constants required by the regulator
criterion, since $\widehat\beta_m$ is fixed. The needed free scaling-dimension
jets are unchanged, and the defect tends to zero by the homogeneous
contraction above. All estimates use the same averaging operation, critical
configuration, action matching and periodized kernels as the trajectory with insertions.


\section{Quantitative stability of the trace functional}
\label{ch17:stability}

The finite-dimensional argument is elementary, but its rate is essential.
Suppose $s:E\to V$ has right inverse $t$ and one-dimensional kernel
$\RR v$. Let $s_i:E\to V$, $t_i:V\to E$ satisfy
$s_it_i=I$, $t_i\to t$, and let $v_i\to v$. If $e_i\ge1$ and
\[
 e_i\|s_iv_i\|\to0,\qquad e_i|\rho(v_i)|\to0,
\]
then, for $\lambda_i=\rho t_i$,
\begin{equation}
 s_i\to s,\quad \sup_i\|\lambda_i\|<\infty,
              \quad e_i\|\lambda_is_i-\rho\|\to0.
 \label{ch17:stable-functional}
\end{equation}
Indeed $k_i=v_i-t_is_iv_i\to v$ lies in $\ker s_i$. The maps
$(h,c)\mapsto t_ih+ck_i$ converge to the invertible map
$(h,c)\mapsto th+cv$, so their inverses are uniformly bounded.
Writing $y=t_is_iy+c_i(y)k_i$ gives $|c_i(y)|\le C\|y\|$ and
$(\lambda_is_i-\rho)y=-c_i(y)\rho(k_i)$, which proves the claim.

Here are the complete rate estimates needed for the trace conclusion;
the next section proves them for the common family.
Write
\[
 e_i=\exp((\log\log b_i)^2/5),\qquad
 q_i=\exp((\log\log b_i)^2/6),
\]
so $q_i/e_i\to0$ and both are smaller than every positive power of
$b_i$. The letter $q_i$ in this section is a scalar bound, unrelated
to the inserted densities of Chapter~\ref{ch11:insertions}.
Let $B_i^L,B_i^X,B_i^D$ be the centered transferred tensors and set
\[
 A_i=\beta_i r_i^4 U_{i,4}^{-1},\qquad
 K_i=\beta_i r_i^5 U_{i,5}^{-1}\operatorname{div}.
\]
The required same-family current transfer estimates are
\begin{align}
 &\|\beta_iK_iB_i^L\|\le Cq_i,\quad
             \lambda_i(\beta_iK_iB_i^L)\to-\kappa_G,
 \label{ch17:transfer}\\
 &K_iB_i^Z\to\operatorname{div}^{\rm Max}Z,\quad
 e_i[\lambda_i(K_iB_i^Z)-\rho(\operatorname{div}_{\rm raw}Z)]\to0
                       \quad(Z=X,D),\notag\\
 &A_iB_i^Z\to Z\quad(Z=L,X,D),\quad
 b_i^p\operatorname{tr}(A_iB_i^Z)\to0
                       \quad(Z=L,X,\ p\in\NN).\notag
\end{align}
These bounds are stronger than smooth free comparison. The next section
proves them using the kernel identities and a separate low-moment
estimate for spatial remainders.

% !TEX root = ../monograph.tex
\section{From the current identity to coefficient estimates}
\label{ch17:marked-rates}

The integration by parts needed here is performed before the
large-field expansion. This order matters: a large-field activity is
only measurable, so its derivative cannot be inserted into the
expansion as an ordinary bounded function.

\subsection{A differentiated block kernel}

In the exact kernel $k^{(L)}_\beta$ of \eqref{ch04:kernel}, condition
on the fine field $U$ and all auxiliary choices $\alpha$. The conditional density of $V_e$ is
$q_{b\beta}(V_e Q_e^{-1})$, where
$q_\gamma(u)=Z_\gamma^{-1}\exp[-\gamma w(u)]$.
For every fixed order $r$, compactness and smoothness of $w$ give
\begin{equation}\label{ch17:kernel-score}
 |D^r q_\gamma(u)|\le C_r(1+\gamma)^r q_\gamma(u).
\end{equation}
Here $D^r$ denotes $r$ successive left derivatives in unit Lie algebra
directions. Dividing by $q_\gamma$, the result is a polynomial in
derivatives of $w$ with at most $r$ factors of $\gamma$. The normalization is
independent of $u$. For a product of link kernels the corresponding
bound has an additional polynomial in the number of differentiated
links. No lower bound on $q_\gamma$ is used.

First truncate the exponentially local current and the kernels of the critical-configuration map (the coefficients of its Taylor expansion in the coarse field) on a finite neighborhood. For the vector field use a finite
gauge-covariant extension of its controlled jets through field order four,
not the uncontrolled higher derivatives of a deep transferred current.
Use the coefficientwise gauge-covariant extension: replace each
fixed-order curvature or covariant-difference monomial by its plaquette
and contour-variable representative, and correct its lower jets
triangularly through order four. Each fixed number of link derivatives
adds only a fixed power of the contour length. The extension therefore
has a bounded weighted norm at a strictly narrower exponential width;
it is not an arbitrary jet section on a space of growing dimension.
Complete in that weighted norm first and then truncate. All additional
fixed-order derivative bounds depend only on the controlled jets and
the fixed stage. Extend this local small-field formula smoothly to
the compact product group. One can use a smooth cutoff
supported strictly inside the logarithmic chart, equal to one on the
smaller chart used for Taylor extraction, and then average over the
compact local gauge group. Equivalently extend its finite covariant
holonomy jets. This preserves the required jets and covariance.
The finite derivative factors are inverse powers of the chart width and
polynomials in the support size. They are allowed in the bounds below.
Only after this extension do we integrate by parts. For an exponentially
local current do this first for finite support truncations. Fixed powers
of the support length are absorbed by a strict margin in the
exponential coefficient width. The derivative bounds and tail errors
are consequently summable, so the finite-period identity passes to
the weighted tail limit.

Let $k(U,V)$ be that exact kernel, let $f(U)$ be the entire incoming
density, and put $F(V)=\int k(U,V)f(U)\,dU$. All original fine large-field indicators
and compatibility conditions are part of $f$ or of the integral over the local gauge transformations;
none is differentiated. For a smooth coarse vector field $Y$, define
\[
 \mathcal D_Y k=(\operatorname{div}_{\rm Haar}Y)k+D_Vk[Y].
\]
Dominated differentiation, using \eqref{ch17:kernel-score}, gives the
\emph{density identity}
\begin{equation}\label{ch17:kernel-divergence}
 \operatorname{div}_{\rm Haar}(FY)
       =\int\mathcal D_Y k(U,V)f(U)\,dU.
\end{equation}
Consequently its integral against every spectator independent of the
varied coarse links is zero. This is an identity for the full differentiated kernel,
including its nonpolynomial remainder, not an assertion that a finite
Taylor polynomial alone has zero exterior moments, that is, zero integral against every spectator independent of the varied links.

The right hand side is an admissible conditional integral with the differentiated kernel inserted.
Its extra factor is a finite sum of smooth local kernel scores times
$Y$ or its derivative. On a protected dependence set the incoming
measurable product remains fixed under the change of variables in
Chapter~\ref{ch06:large-fields}. Additional coarse derivatives act
on the score and the common density ratio only. The bounded pointwise-jet estimate
\eqref{ch06:actual-jet-integral} therefore applies with a fixed
polynomial factor in $\beta$, support size and current derivative bounds.
Rooting the connected expansion at this extra factor gives the same
expansion with insertions as in Chapter~\ref{ch11:insertions}. In particular
large-field terms retain their original prescribed large-field sets.
This argument does not differentiate the event indicators introduced when that
conditional integral is subsequently decomposed.

The small-field contribution needs a sharper identity than the absolute
score bound. Fix the full dependence set of a rooted connected component and use one
additional coarse derivative beyond those prescribed for the extracted
coefficient. On its smaller small-field chart, make exactly the conditional
Gaussian change of variables used to construct the small-field expansion.
Write its normalized conditional integrand as $p(V,z)$ and its integral
as $Z_{\rm ord}(V)$. For finite small-field selections, differentiation
under this integral gives
\[
 D_VZ_{\rm ord}[Y]=\int D_Vp[Y]\,dz,
 \qquad
 \int D_Vp[Y]\,dz
 =Z_{\rm ord}D_V\log Z_{\rm ord}[Y].
\]
Every Jacobian, relative determinant and conditional normalization is
included in $p$; thus this is not the derivative of an unnormalized
Gaussian moment. The absolutely convergent small-field logarithm, with
one additional derivative, passes this identity to the countable expansion.
By its definition in Chapter~\ref{ch09:renormalization},\mgubi{Issue 9: $Z_{\rm ord}$ is defined only above, for one component; $E_{\rm reg}$ is defined nowhere. Chapter~\ref{ch09:renormalization} gives this form with a constant only for related objects (the pure small-field block integral; the constant extracted once, globally; part (a) of Assumption~\ref{ch09:forest-assumption}, unproved). A constant for one component needs an argument.}
\[
 Z_{\rm ord}=\mathcal C
       \exp[-\beta S-\mathcal L-E_{\rm reg}],
 \qquad D_V\log\mathcal C=0.
\]
If an intermediate local normalization depends on $V$, it is included
in $\mathcal L+E_{\rm reg}$, not in $\mathcal C$. Consequently the
small-field score is \emph{exactly}
$-\beta DS[Y]-D\mathcal L[Y]-DE_{\rm reg}[Y]$.
The regular term is estimated using its already proved differentiated
regular norm. We do not estimate it by the polynomial-in-$\beta$
absolute score bound.

This argument uses the full rooted component, not just the score
stencil. In the rooted connected expansion, small-field factors in
components disjoint from the component of the differentiated kernel cancel
between the differentiated numerator and its normalization. Every term with a large fluctuation-field event anywhere in the connected
component of the differentiated kernel, including a distant such event joined to the score, is kept in the complementary
sum. The smooth cutoffs used above are constant on the smaller small-field
chart; their derivatives are supported where at least one enlarged
chart condition fails. Attach this condition to its dependence set. The same
Gaussian or energy estimate for the nested cutoffs bounds this new
large fluctuation-field event, while all incoming prescribed large-field sets are retained.
The absolute score bound is used only for these complementary terms,
where the exponentially small probability of the large fluctuation-field event absorbs its polynomial
$\beta$ factor. No derivative of a hard indicator is taken. This proves
the small-field cancellation and the complementary large-field estimate
separately, which is essential to preserve the regular power gain.

\subsection{The finite degree-five coefficient}

At the comparison scale $k$ write
\[
 h=(\log\log\widehat\beta_k)^2,\qquad H=e^h,
 \qquad t=\widehat\beta_k^{-1/2}H.
\]
All constants in this paragraph are for a fixed stage and the fixed
finite derivative orders and degrees; they are fixed before its lower bound on the
reference coupling. Let $\mathcal P_5$ denote radial Taylor extraction through
degree five and let $f$ be the number of curvature factors of a homogeneous
coefficient. Set $\mathcal P_{5,t}=t^{f-2}\mathcal P_5$, followed by
real signed-permutation averaging. There are no invariant linear forms on $\gLie$. Thus the nonzero vector coefficients under consideration
have $f\ge2$, and this normalization does not introduce a negative
power of $t$ beyond the displayed $t^{-2}$.

Apply the exact step with insertions to \eqref{ch17:kernel-divergence}, retaining
its entire remainder, and use the current identity through order four.
Let $j_k$ be the radial coefficient of the current divergence, in the
same threshold coordinates. The classical term $-\beta DS[Y]$ is
$-\beta j_k$ up to the contracting current defect. Thus the coefficient
of the full differentiated kernel, with zero exterior moments, is $\Sigma_k-\beta j_k$;
$\Sigma_k$ alone is not asserted to have zero exterior moments. Its
remaining degree-five scalar-vector coefficient can be written
\begin{equation}\label{ch17:finite-sigma}
 \Sigma_k=\mathcal P_{5,t}
       (\operatorname{div}_{\rm Haar}Y-D\mathcal L[Y])
       +E_{\rm stat}+E_{\rm tail}+E_{\rm reg}+E_{\rm mark}.
\end{equation}
Here the kernels of the critical-configuration map and the loop kernels are truncated at
$C_j(\log\widehat\beta_k)^2$. Their common exponential decay makes
$E_{\rm stat}$ and $E_{\rm tail}$ at most $H^{-1}$ after finite
projection factors. The regular error satisfies
\[
 t^{-2}\epsilon_k^{\rm reg}
       =\widehat\beta_k^{-1/10}H^{-2}.
\]
A directional derivative along $Y$ has the same bound up to the fixed
current norm, since the derivative bound contributes a factor $t^{-1}$
and $Y=O(t)$ on the small-field chart. The previously proved logarithmic current bounds
and fixed localization factors are absorbed by increasing the lower bound on the reference coupling.

For the remainder with insertions one must first cancel the whole classical
current divergence, not merely use the absolute kernel-score bound.
The current identity agrees through field order four. The finite
gauge-covariant extension just chosen has controlled higher derivatives,
so the remaining classical error starts at order five and contributes
\[
 t^{-2}\beta t^5=O(\widehat\beta_k^{-1/2}H^3).
\]
Fixed current and localization factors preserve a bound of the form
$C_j\widehat\beta_k^{-1/10}H^{C_j}$ after increasing the lower bound on the reference coupling. The small-field loop
and Jacobian Taylor remainders have no leading factor $\beta$ and
gain an additional small-field factor. This controls the high-field-order remainder. It does not control a
low-field-order term with many spatial differences: vanishing Taylor coefficients through the specified scaling dimension alone give no additional small-field factor. The remaining low-field-order spatial terms are controlled by exact
polynomial preservation, rather than by an extra power of the field.

\begin{proposition}[Low-moment estimate for spatial remainders]
\label{ch17:spatial-moments}
For the linearization of the critical configuration and the regular local terms of the differentiated kernel just constructed,
a remainder whose radial Taylor coefficients through scaling dimension five vanish
has extracted degree-five coefficient $O(H^{-1/2})$. The estimate is
uniform in period and in the finite current truncations, after the
stage orders and thresholds have been chosen with one extra derivative
order. It applies to both the nonclassical coefficient and the
normalized current divergence.
\end{proposition}
\begin{proof}
First identify the linearization of the critical configuration on polynomials. Recall the notation of
Chapter~\ref{ch05:gaussian}: $D_l,D$ are the rescaled fine and coarse
covariance symbols, $x_l,x$ their gradient columns, and
$\alpha=\sum_l|m_l|^2\tr D_l$.
Write $s=\lambda_{M_c}(p)$ and
$D=s^{-1}\operatorname{diag}(B_a)$, where the $B_a$ are the
regularized diagonal covariance factors. The scalar correction
\[
 \chi=\frac{x^*\operatorname{diag}(B_a^{-1})x}{s}
\]
is analytic and nonzero at zero, as in the pole cancellation for the
formulas for the linearization. Put
$\sigma=\chi^{-1}x^*\operatorname{diag}(B_a^{-1})$, so that
$\pi=s^{-1}\sigma$ in those formulas. For a nonzero alias define
\[
 g_{l,a}=\overline{m_l}
 \left(\frac{\tr D_l}{s\alpha}
       -\frac{(D_l)_{aa}|r_L(\theta_{l,a})|^2}{B_a}\right).
\]
Substitution into the linearization built with the linear averaging operator gives the regular formula
\[
 \mathcal A_l=\mathcal B_l+
          \operatorname{diag}(g_{l,a})x_l\sigma,
 \qquad \mathcal B_l=D_lF_{1,l}^*D^{-1}.
\]
Here $D_l$ is regular at a nonzero alias, and $s\alpha$ and
$B_a$ are analytic and nonzero at zero. Thus no uncancelled factor
$s^{-1}$ remains. Both terms contain $\overline{m_l}$, with
$m_l=w_L(\theta_l)^j$ and $w_L(\theta_l(0))=0$.
Every derivative of $\mathcal A_l$ of order less than $j$ therefore
vanishes at zero. The column at the zero alias has value $I$. Through \eqref{ch17:polynomial-response}, this
order-$j$ vanishing is also what the preservation of the order of vanishing of the
current defect in Section~\ref{ch17:quadratic-current} uses. Stars and conjugates in these symbols mean analytic
transpose at the opposite momentum, as in Chapter~\ref{ch05:gaussian}.

Let $A_x^{ab}(z)$ be the physical row from coarse direction $b$ at
$\lfloor x/L\rfloor+z$ to fine direction $a$ at $x$.
Differentiating the Fourier inversion formula through order less than $j$
therefore gives, for every vector polynomial $p$ of degree less than $j$,
\begin{equation}\label{ch17:polynomial-response}
 (W_sp)(x)_a=\sum_{b,z}A_0^{ab}(z)p_b(x/L+z),
 \qquad \sum_z A_0^{ab}(z)=L^{-1}\delta_{ab}.
\end{equation}
The exponential row bound justifies the derivatives and the sum.
Polarization gives mixed polynomial moments. In particular the output
is a global polynomial of no larger degree, not a polynomial separately
on each block. Its homogeneous degree-$r$ leading term is
$L^{-r-1}p(x)$; lower coefficients obey the same ratio-independent row bounds.

Choose $j$ above the order of the coefficient evaluations. In a polarized coefficient
evaluation take only tuples of potential polynomials whose summed scaling
dimension is at most five. Equation~\eqref{ch17:polynomial-response}
preserves or lowers each polynomial degree, so every resulting tuple
remains within that same finite list. The polynomial-moment
cancellation of Section~\ref{ch11:one-graph} makes its deterministic
coefficient evaluation exactly zero. This is not a claim of vanishing on arbitrary
polynomial inputs of degree less than $j$. Truncating the coarse field to the ball of radius
$C_j(\log\widehat\beta_k)^2$ changes this coefficient evaluation only by the
exponential row tail, times fixed polynomial directional and support
factors. That tail is smaller than every prescribed power of $H^{-1}$.

The smooth representative of the linearization of the critical configuration is
$W_a=W_s-d_fERW_s$. Its difference from $W_s$ is a fine gradient.
The coefficient of these regular terms is a spacetime vector with trivial adjoint action on its values,
invariant under local fine gauge transformations. Indeed gauge
equivariance of $Y$ and invariance of Haar volume imply invariance of
$\operatorname{div}_{\rm Haar}Y$; the action and loop variations
are invariant for the same reason. Apply the corresponding fine gauge
transformation before evaluating its local jet. The Baker--Campbell--Hausdorff formula with the same local gauge choice compares the two exponential
representatives with an error $C_jP(B)t^2$ in the link logarithms,
where $P$ is a fixed polynomial. Dividing by the fine threshold gives
$C_jP(B)t$. This argument uses the exact correction to the averaging operation, not an
unmatched linear averaging operator.

It remains to compare the conditional insertion itself to this
polynomial coefficient evaluation. Normalize the input link logarithms by $t$.
The nonlinear critical configuration, the same local gauge coordinate change,
and the deterministic truncation contribute $C_jP(B)t$ and the exponential
tail just described. Gaussian fluctuations contribute
$C_jP(B)g/t$, where $g=\beta^{-1/2}$, by the fundamental theorem of
calculus and the one additional derivative. Explicitly, if
$p(z)\ge1$ is a polynomial majorant for the finite Gaussian coordinates,
the difference between the conditional and the deterministic evaluation in
each prescribed derivative is bounded by
\[
 C_j N_k P(B)\bigl(t+g/t+e^{-c_jB}\bigr)\,p(z)^{C_j},
\]
with an appropriate enlargement of the truncation radius in the last term.
Here $N_k$ is the norm of the already $t^{-2}$-normalized regular
term of the differentiated kernel and its derivatives in normalized link logarithms, including
its current and loop coefficients. The extraction normalized by the number of curvature factors has
already been performed; no further $t^{-2}$ is applied to this bound. For this bound subtract the deterministic coefficient evaluation
inside the same normalized conditional integral. Its derivative
through the prescribed scaling dimension is zero. Use the joint
polynomial moments and differentiated density scores of that one
conditional law to estimate the evaluation difference; no small
total-variation distance between the nonlinear law and a Gaussian
is asserted. The pointwise jet and moment estimates are those of \eqref{ch06:actual-jet-integral};
no differentiation of incoming large-field factors is involved.
The complement is estimated by the already retained large-field bound. Finite coefficient recovery adds a factor depending only on the stage.

There are no invariant terms that are constant spacetime vectors or linear in the $\gLie$-valued field. Hence these regular terms of the differentiated kernel start at field order two. Their
values and normalized-link derivatives on the threshold ball contain
a factor $t^2$; the explicit $t^{-2}$ removes exactly this factor.
The classical part has already been cancelled and its order-five
remainder treated separately above. Thus for the present regular terms
$N_k\le H^{\theta_j+o(1)}$, with
$\theta_j\to0$: current norms are logarithmic and the retained loop
norm has this bound. Since $B$ and all truncation radii are fixed powers
of $\log\widehat\beta_k$, their polynomial factors are $H^{o(1)}$.
The coupling band gives $g/t\le C H^{-1}$, whereas $t$ and
the truncation tail decay faster than every fixed inverse power of $H$.
The resulting coefficient error is $H^{-1+\theta_j+o(1)}$, hence
at most $CH^{-1/2}$ after increasing the lower bound on the reference coupling. Large supports are first truncated
using a margin in the support exponent; its tail obeys the same
bound. This proves the proposition for the full local sum.
\end{proof}

The differentiated-kernel
bound \eqref{ch17:kernel-score} applies to the complementary part.
The complementary part retains the square root
$\exp(-H^{1/5}/2)$ of the bound \eqref{ch04:activity-bounds} on the large-field activities. Combining Proposition~\ref{ch17:spatial-moments} with the high-field
remainder, the nonwrapping contributions are bounded by
\begin{equation}\label{ch17:sigma-error}
 C H^{-1/2}+C_j\widehat\beta_k^{-1/10}H^{C_j}
       +C_j\widehat\beta_k^{C_j}H^{C_j}e^{-H^{1/5}/2}.
\end{equation}
The fixed-order score has changed the polynomial constants, not the
eligibility event \eqref{ch04:eligibility} of a large-field activity or its exponential gain. If $N_{\rm per}$ is the
period in scale-$k$ units, a wrapping term has covering number at least
$N_{\rm per}/[10(2B_k+1)]-1$. Retaining two units of the support exponent gives the
additional error
\[
 C_j\widehat\beta_k^{C_j}H^{C_j}
       \exp\bigl(-N_{\rm per}/[5(2B_k+1)]+2\bigr).
\]
It vanishes in the ordered period limit at fixed $k$.

The two terms other than $CH^{-1/2}$ in
\eqref{ch17:sigma-error} are $o(H^{-1})$. To check the
order of limits explicitly, put $x=\log\log\widehat\beta_k$.
The logarithm of the first bound multiplied by $H$ is at most
$C_j+(C_j+1)x^2-e^x/10$; that of the second is at most
$C_j+C_je^x+(C_j+1)x^2-e^{x^2/5}/2$.
Both tend to minus infinity at each fixed stage. Increasing its lower bound on the
reference coupling makes the bounds uniform on that stage.

The loop coefficient has zero $\rho$ by its two moment cancellations:
its quadratic kernel annihilates both constant and affine links, whereas
the cubic momentum coefficient functional pairs a constant with a quadratic link or two
affine links. This uses the loop kernel of Chapter~\ref{ch07:matching},
not only gauge invariance.
The calculation with mass regularization identifies the remaining coefficient as
$\mathfrak h(Y)/12$. One can check the factor directly in the three
vector coordinates of \eqref{ch17:rho}, before the quotient by the equation of motion: the cubic momentum
functional takes values $-24,-96,0$, while $\rho$ takes values
$-2,-8,0$. The lattice correction is the matching lower-order
Taylor correction in $\mathfrak h$.

Choose the number of transfers and then the period at each fixed $k$ so
that the error of the mass regularization and the finite momentum approximation are
smaller than $\widehat\beta_k^{-30}$. The retained loop bound is
$2H^{\theta_j}$, with $\theta_j\to0$, and current derivative norms
are logarithmic. Sufficiently large lower bounds on the reference coupling therefore give joint ordered limit
coefficients satisfying
\begin{equation}\label{ch17:sigma-bound}
 \|\Sigma_k\|\le C H^{1/6},\qquad
 \rho(\Sigma_k)\longrightarrow-\kappa_G.
\end{equation}
All coefficient maps and the underlying correlations are retained
in the same joint limit. Compactness is used only after the uniform
finite-cutoff estimates have been established.

\subsection{The kernel vector and the necessary rates}

Let $j_k$ be the degree-five coefficient of the transferred current's
divergence, in the same threshold coordinates. Its local Taylor
coefficient has zero $\rho$: the trace is a forward difference of a
degree-three primitive, whose invariant projection is a multiple of the identity field.
The same localization and error bounds with insertions compare the transferred
coefficient with this Taylor coefficient within $CH^{-1/4}$.
Consequently, in the notation of the stability argument, with $v=(0,0,1)$ the kernel vector of the free quotient $s$,
\[
 j_i\to-v,\quad |\rho(j_i)|\le CH_i^{-1/4},\qquad
 v_i^{\rm eq}=\beta_i^{-1}\Sigma_{k_i}-j_i\to v.
\]
With $e_i=H_i^{1/5}$ this gives $e_i\rho(v_i^{\rm eq})\to0$.

There is also a stronger, independent bound on the image of this
vector. The differentiated kernel is carried as a density-valued factor, without
dividing it by $F$ or claiming that the resulting quotient is local.
More directly, for every incoming coefficient $f_A(U)$ with insertions,
including all the original fine spectator factors,
\[
 \int \mathcal D_Y k(U,V)\,dV=0
 \quad\hbox{for each }U,\qquad
 \int \mathcal D_Y k(U,V)f_A(U)\,dU\,dV=0.
\]
The same differentiated kernel is used for every $A$, so the joint rooted
expansion retains all pushed spectator cross-modes and this exact
zero-moment statement. There is no smoothness requirement on $f_A$.
For a smooth local coarse spectator $A(V)$ one also has
\[
 \int A(V)\!\int\mathcal D_Y k(U,V)f(U)\,dU\,dV
 =-\int D_VA[Y]\!\int k(U,V)f(U)\,dU\,dV.
\]
It is zero when the spectator does not read the varied links. For an
exponentially local spectator that need not have finite range, first truncate its dependence set
and use this identity on disjoint terms. Retain the remaining terms
in the same joint expansion with insertions. A connected term reaching from
the differentiated kernel to a spectator at distance $R$ has covering number at least
$R/(2B_k+1)-1$. Retaining two units of the support exponent bounds these terms by
$C\exp[-2R/(2B_k+1)]$, with their normalized large-field majorants as well.
For spectator configurations separated by a fixed positive multiple of $a_n$, their
distance in scale-$k(n)$ units is at least $c\exp(T_n)$. This tail
is therefore smaller than $e^{-T_n/100}$ after all fixed factors.
Every spectator subset uses the same density without insertions.
Write $\tau_j=\log(1/a_j)$ and
$T_n=(\log\widehat\beta_n)^{48}$. Choose $k(n)$ to be the smallest
integer $k\ge n$ for which $\tau_k-\tau_n\ge T_n$. Match the full
coefficient and remainder backwards from $k(n)$ to $n$. Equation~\eqref{ch11:descent-error} gives
$e^{-T_n/8}$ before multiplication by its constants. Fixed projection and
threshold factors have logarithm $O(\log\widehat\beta_n)$; accumulated
ratio-independent coefficient-map factors are $o(T_n)$. Here it is only the latter
constants that are multiplied along the interval of scales. Ratio-dependent nonlinear
jet constants have already been absorbed into the errors controlled by the lower bounds on the reference coupling of the stages.

Place a copy in each reflected cell at the first backward crossing of $B_n^2$, so
$B_n^2\le a_m/a_n\le L_*B_n^2$. The lower bound on the partition function
in \eqref{ch11:replicated-bound} contributes at most
\[
 L_*^4B_n^8q_n+6=O((\log\widehat\beta_n)^{22})=o(T_n),
\]
using $L_*\le\log\widehat\beta_n$, $B_n=O((\log\widehat\beta_n)^2)$
and the per-site normalization exponent $q_n=O((\log\widehat\beta_n)^2)$.
The inverse source weights in \eqref{ch11:replicated-bound} have logarithm $O(\log\widehat\beta_n)$.
The dimensionful normalization of the insertion cancels the factor $a_n^{-5}$.
Thus the normalized residual of \eqref{ch11:descent-error} is at most $e^{-T_n/100}$.
This argument uses the full differentiated kernel with zero exterior moments; its Taylor remainder
is carried by the graph of the backward comparison, not discarded.

In the choice of radii along the diagonal subsequence require
$r_i/a_{n_i}\le\exp((\log\log\widehat\beta_{n_i})^3)$.
Finite separating matrices of free two-point correlation pairings have uniformly bounded left
inverses. For every fixed $p$ the image therefore satisfies
\[
 b_i^p\|s_i(v_i^{\rm eq})\|
 \le C b_i^p\exp\bigl(5(\log\log\widehat\beta_{n_i})^3-T_{n_i}/100\bigr)
 \longrightarrow0.
\]
This proves the quantitative kernel hypothesis in the stability lemma.
Its bounded right inverses give bounded $\lambda_i$, with
$e_i\|\lambda_i s_i-\rho\|\to0$.

The same identity for the differentiated kernel gives
$\beta_iK_iB_i^L-s_i\Sigma_{k_i}\to0$.
For $Z=X,D$, the ordinary backward comparison gives
$K_iB_i^Z-s_i\operatorname{div}_{\rm raw}Z=o(b_i^{-p})$ for every
fixed $p$. Apply the bounded functional and use
$H_i^{1/6}/e_i\to0$ in \eqref{ch17:sigma-bound}.
This proves the first two lines of \eqref{ch17:transfer}.

Finally the trace of $X$ is identically zero, while the transferred
trace of $L$ is an exact forward difference. Match its degree-three
primitive at a fixed late pair of cell scale and stopping scale. The backward comparison gives endpoint
error $C e^{-(\tau_k-\tau_{k_0})/8}$, uniformly near the configurations used in the finite separating
correlation pairings. Subexponential coefficient conversions and support truncation
make this $O(a_k^\eta)$ for any fixed $0<\eta<1/8$.
Average over $m=\lfloor a_k^{-1+\eta/2}\rfloor$ consecutive differences.
The identity-field parts of the two endpoint values are equal and cancel, leaving at most
$2Ca_k^\eta/(ma_k)=O(a_k^{\eta/2})$.
The averaging length $ma_k$ in continuum units tends to zero, so the finite matrix of
correlation pairings converges to its original invertible matrix. This proves the
last rate in \eqref{ch17:transfer}, including multiplication by every
fixed power of $b_i$. Ordinary Taylor matching gives its remaining
convergence $A_iB_i^Z\to Z$.


\section{The coefficient and the all-radius limit}
\label{ch17:limit}

\begin{proposition}[The trace coefficient]
\label{ch17:coefficient}
For the prescribed-radius comparison maps constructed above,
\[
          \beta_i r_i^4 U_{i,4}^{-1}\Theta
                       \longrightarrow-\kappa_G\mathcal O.
\]
\end{proposition}
\begin{proof}
The free divergences of $X,D$ form a basis of $V$. Thus the maps
$f_i(c)=c_0K_iB_i^X+c_1K_iB_i^D$ have bounded inverses eventually.
Put $c_i=-f_i^{-1}(\beta_iK_iB_i^L)$, so $|c_i|\le Cq_i$ and
$Z_i=\beta_iB_i^L+c_{i,0}B_i^X+c_{i,1}B_i^D$ is conserved.
Applying $\lambda_i$ to its zero divergence gives
\[
 4c_{i,1}=\lambda_i(\beta_iK_iB_i^L)
 +c_{i,0}\lambda_i(K_iB_i^X)
 +c_{i,1}[\lambda_i(K_iB_i^D)+4].
\]
The last two terms are $O(q_i)o(e_i^{-1})=o(1)$, whence
$c_{i,1}\to-\kappa_G/4$. Since $q_i/\beta_i\to0$,
$r_i^4U_{i,4}^{-1}Z_i\to L$. The faster trace bounds give
$\beta_i\operatorname{tr}(r_i^4U_{i,4}^{-1}Z_i)
\to-\kappa_G\mathcal O$.

The conserved line gives $Z_i=\nu_iT$. Choose one free translation
flux pairing with nonzero value. Its flux on the rescaled $Z_i$ tends
to that value. The Ward identity for $T$ and derivative comparison
make its corresponding flux tend to the same nonzero value.
Their ratio is $\nu_i$, so $\nu_i\to1$. Finally centering subtracts
$e_{\rm tr}\beta_i r_i^4 1\to0$ by \eqref{ch17:running-radius}.
\end{proof}

\begin{theorem}[The all-radius trace law]
\label{ch17:all-radii}
The centered trace in the selected phase satisfies
\begin{equation}
 \lim_{r\downarrow0}r^8\log^2(1/r)G_\Theta(r)
              =\mathcal A_G:=\frac{3\dim_\RR\gLie}{\pi^4}>0.
 \label{ch17:trace-law}
\end{equation}
\end{theorem}
\begin{proof}
The free normalization is
$S_2^{\rm Max}(\mathcal O,\mathcal O;x,0)=\mathcal A_G|x|^{-8}$.
One can check its factor without a tensor convention ambiguity:
the curvature covariance is
$(H_{ai}H_{bj}-H_{aj}H_{bi})/(\pi^2|x|^4)$ with
$H=I-2xx^{\mathsf T}/|x|^2$ orthogonal. The sum of its squared
numerators over ordered pairs is $24$. The two Wick pairings and
the two factors $1/4$ give $2\cdot24/16=3$ per colour.

Start with any prescribed null sequence of radii, and take its
comparison subsequence. Proposition~\ref{ch17:coefficient} gives
convergent moving labels, while \eqref{ch17:running-radius} cancels
the factor $\kappa_G^2$ in their pair correlation. For a generic
pair at distance $t$, lattice rounding therefore yields
\[
 r_i^8\log^2(1/r_i)G_\Theta(r_it_i)\to\mathcal A_Gt^{-8},
                         \qquad t_i\to t.
\]
Choose the pair $0,(t/2,t/2,t/2,t/2)$ for $1/2<t<2$.
For $t<1$ the monotonicity of \eqref{ch17:laplace} bounds the upper
limit at $r_i$ by $\mathcal A_Gt^{-8}$; $t>1$ bounds the lower
limit. Send $t$ to one from each side. Every prescribed null
sequence has a subsequence with limit $\mathcal A_G$, which proves
\eqref{ch17:trace-law} by the subsequence criterion.
\end{proof}

\section{What the trace law excludes}
\label{ch17:exclusion}

Specify the comparison class before using this asymptotic. Consider finite
linear combinations of finite products of the mass-shell Laplace kernels
\[
 K_{m,k,h}(r)=\int_{\RR^3}
 \frac{\omega_m(\boldsymbol p)^k h(\boldsymbol p)}{2\omega_m(\boldsymbol p)}
 e^{-r\omega_m(\boldsymbol p)}\,d^3\boldsymbol p,
 \qquad\omega_m=(m^2+|\boldsymbol p|^2)^{1/2},
\]
where $m\ge0$, $k$ is a nonnegative integer and $h$ is a polynomial;
empty products are constants. The measure has no implicit Fourier
normalization. All these integrals are locally integrable at zero
and converge for $r>0$.

Each such profile has, near zero, a finite sum of terms
$r^a(\log r)^b$, with integer $a$ and nonnegative integer $b$,
plus a bounded remainder. This is the standard Laplace asymptotic expansion
\cite{Wong2001}, applied as follows: decompose $h$ into homogeneous terms,
use polar coordinates, and expand
$(\omega^2-m^2)^{(j+1)/2}$ for $\omega\ge2m$.
Choose enough terms that the remainder has as many energy moments
as needed for the other factors in a product. The retained integer
powers integrate to powers and logarithms; the bounded energy
interval contributes an analytic function. Thus only nonnegative
powers of the logarithm occur.

After multiplication by $r^8\log^2(1/r)$, a nonzero leading term
with $a\le-8$ diverges in magnitude, while every term with $a>-8$
and the bounded remainder tend to zero. Therefore a finite limit,
if it exists, is zero. The positive limit in
\eqref{ch17:trace-law} proves that $G_\Theta$ cannot agree with any
such profile on any interval $(0,\delta)$.

This excludes that specified finite mass-shell polynomial comparison.
It does not exclude a generalized free field chosen to have the
same positive measure $\rho_\Theta$, nor an infinite Wick expansion.
No mass gap follows from a short-distance trace law.

\setcounter{chapter}{17}
% !TEX root = ../monograph.tex
\chapter[Operator products and contact extensions]{Operator products and extensions across diagonals}
\label{ch18:products}

There are two distinct questions about colliding insertions. An operator
product expansion replaces a small cluster by finitely many local fields
when tested against separated spectators. A distributional extension
assigns values to test functions whose supports meet collision diagonals. The first requires
quantitative information about the marked renormalization group. The
second already follows, with a conservative finite order, from the
polynomial separated bounds. Neither construction defines a product of
point fields by simply setting their positions equal.

We use the common filtration $F_d$, its inclusions $\iota_{d,e}$,
and the mixed correlations $S$ from Chapter~\ref{ch13:continuum}.
The orthogonal actions on these spaces and the phase correlations
$S^\xi$ are those of Chapters~\ref{ch14:ward} and
\ref{ch15:phases}. Write $\mathfrak A$ for the
commutative compression algebra of Chapter~\ref{ch15:phases} and
$M(w)$ for the compression of a whole separated word $w$.
We first convert the lattice replacement estimate into the norm
bound needed for the operator product expansion. The contact extension constructed later does not use it.

The OPE originates with Wilson~\cite{Wilson1969}.
Fr\"ob and Holland establish the Yang--Mills OPE to every fixed order
of perturbation theory~\cite{FrobHolland2016}. That result explains the
relevant structure, but does not supply the nonperturbative estimate
required below. The standard distribution-extension argument is the
normal Taylor subtraction of Brunetti and Fredenhagen
\cite[Section~5.2, Theorems~5.2--5.3]{BrunettiFredenhagen2000}.
We use that procedure with an explicit choice adapted to the union of
all pair diagonals; no reconstruction theorem requires this particular
extension.

\section{A fixed prescription for all finite degrees}
\label{ch18:prescription}

A probe is a fixed word $v$ at distinct nonzero positions. It defines
$\lambda(A)=S(A(0)v)$ on each field space on which $A$ is defined.
These functionals separate $F_d$, by the definition of correlation-null
fields. Since $F_d$ is finite dimensional, finitely many probes span
its dual. A probe chosen at one degree continues to act at every higher
degree by the same correlation, using the physical inclusions.

One can choose fields $u_\alpha$ and finite linear combinations
$b_\alpha$ of probes, with a degree $|\alpha|$ attached to each, such that
\begin{equation}
 b_\alpha(u_\beta)=\delta_{\alpha\beta},\qquad
 A=\sum_{|\alpha|\le d}b_\alpha(A)u_\alpha\quad(A\in F_d),
 \label{ch18:dual}
\end{equation}
where fields in lower degrees are included in $F_d$ when necessary.
All these identities hold simultaneously at every finite degree.
Here is the useful linear-algebra construction. Suppose the choices
through $d$ have been made and let $R:F_d\to F_{d+1}$ be inclusion.
The old probes give a map $B:F_{d+1}\to F_d$ with $BR=1$.
For any section $S_0$ of the quotient map
$q:F_{d+1}\to F_{d+1}/RF_d$, replace it by
$S_0-RBS_0$. The corrected section is killed by every old probe.
Lift a quotient basis through it, and represent its dual quotient
coordinates by finite probe combinations. This extends
\eqref{ch18:dual} without changing an earlier field or probe.
In particular $b_\alpha$ vanishes on all degrees below $|\alpha|$.
This is only finite-dimensional linear algebra; no infinite expansion
of a field is being asserted.

For a cluster word $w$, define the unsymmetrized field reconstruction
\[
 H_d(w)=\sum_{|\alpha|\le d}b_\alpha(w)u_\alpha,
\]
where $b_\alpha(w)$ means replacing the distinguished field in its
probe correlations by the entire cluster. The mixed table is always
used for this measurement. Thus the coefficients have no dependence
on a phase or on spectators later placed outside the cluster.
At a center $o$, apply this definition to $\tau_{-o}w$.
For each finite collection of probes there is a positive clearance:
a cluster sufficiently close to $o$ cannot meet a probe point, even
after rotating the cluster. All subsequent formulas are used on this
clear region.

Write $R_d(K)$ for the orthogonal action on $F_d$, with normalized
Haar measure on $O(4)$. Symmetrize the reconstruction by
\[
 A_dw=\int_{O(4)}R_d(K)H_d(K^{-1}w)\,dK,
 \qquad B_dw=\tfrac12\{A_dw+\overline{A_d\bar w}\}.
\]
Both maps fix every single field of degree at most $d$ inserted at zero.
They are smooth on the clear configuration region, since finitely many
smooth correlation functions are integrated over a compact group.
They preserve reality and intertwine rotations. Translation covariance
is implemented by translating the center as above.

Averaging must be followed by a compatible removal of lower degrees.
Let $I_qA=A(0)$, put $L_{i,q}=B_iI_q$, and define
\begin{equation}
 Q_0=1,\qquad
 Q_j=1-\sum_{i<j}\iota_{i,j}Q_iL_{i,j},\qquad
 E_j=\bigcap_{i<j}\ker(Q_iL_{i,j}).
 \label{ch18:graded}
\end{equation}
The inclusion identities give, by induction,
\[
 Q_jL_{j,q}=0\quad(q<j),\qquad
 \sum_{i\le j}\iota_{i,j}Q_iL_{i,q}=\iota_{q,j}\quad(q\le j).
\]
Consequently $Q_j$ is the projection onto $E_j$ along
$\iota_{j-1,j}F_{j-1}$. These spaces are invariant under rotations
and conjugation. Fix real bases $(O_{j,a})_a$ of $E_j$, taking the
unit itself in degree zero, and denote their coordinate functionals
by $\omega_{j,a}$. The coefficients used in this chapter are
\begin{equation}
 c_{j,a}(w;o)=\omega_{j,a}Q_jB_j(\tau_{-o}w).
 \label{ch18:coefficients}
\end{equation}
The fields, probes and projections are fixed once, before any
truncation, phase, spectator list or small-distance limit.


\section{The marked replacement estimate}
\label{ch18:input}

The expansion of Chapter~\ref{ch11:insertions} keeps every joint
nilpotent marker coefficient. To replace a cluster by a single field,
one needs more than the one-insertion graph contraction. Markers may
merge as their separation becomes smaller than a block. Their source
weights multiply; an incoming term involving at least two insertions is kept as one factor,
with its support, prescribed set of large-field sites and eligibility condition. Its stronger support
norm absorbs the loss when several marks are assigned to the same coarse
region. After local Taylor subtraction, the signed residual must be
transported on the same history and compared in the physical quotient.
Taking absolute values before that subtraction loses the cancellation.

Here is the quantitative replacement estimate. It is formulated
so that its two uses, determination of local polynomial coefficients
and phase-uniform error control, are visible. Put $w_0=\sum_i d_i$ for a fixed cluster degree
list, and let
\[
 x_i=o+r z_i,\qquad |z_i|\le M,\qquad |z_i-z_j|\ge\eta>0.
\]
Cluster labels range over fixed bounded subsets of their finite-dimensional
spaces. Spectator labels have bounded degrees and coefficients; their
positions stay at distance at least $c>0$ from $o$ and from one another.
The constants below may depend on these finite lists, $M,\eta,c$ and
a requested error exponent, but not on positions, phases or $r$.

\begin{proposition}[Same-history marked replacement]
\label{ch18:replacement-input}
For every finite degree $D$ and every $\epsilon>0$, on the geometry
just specified and for all sufficiently small $r$, the unsymmetrized reconstruction
coefficients obey
\[
 |b_\alpha(w;o)|\le C r^{|\alpha|-w_0-\epsilon}.
\]
For each spectator word $v$, set
\[
 Z_D=M(wv)-M((H_D(\tau_{-o}w))(o)\,v).
\]
The actual lattice marked replacement, with its signed residual and
all original terms involving at least two insertions retained, gives the bound uniform in the number of reflected copies
\begin{equation}
 \langle\Omega,(Z_D^*Z_D)^k\Omega\rangle
       \le (C r^{D+1-w_0-\epsilon})^{2k},\qquad k\ge1.
 \label{ch18:replacement-moments}
\end{equation}
The constant and the small-$r$ threshold precede $k$; cutoff thresholds
may depend on each fixed $k$. The same estimates hold for every fixed
finite list of differentiated insertions, with the corresponding
scaling dimensions, and for the finitely many probe families above,
uniformly under their rotations and conjugations.
\end{proposition}

This is an analytic estimate on the selected family, not an assertion
that any bounded marked expansion has an OPE. Its construction requires
the signed local transport, the stronger anchored marked support bound,
and the graded quotient comparison. In particular, if a moving field
has coefficient size $\varepsilon_m$ at scale $m$, the needed quotient
estimate is of the form
\[
 \operatorname{dist}_{F_t}(A,\iota_{i-1,t}F_{i-1})
       \le a_m^{-(t-i)}e^{o(\log(1/a_m))}\varepsilon_m .
\]
The lower field must be identified by all spectator responses, using the
normalized-null lemma of Chapter~\ref{ch13:continuum}. A bound on a chosen
coordinate representative alone does not supply this statement.
The estimate uniform in the number of reflected copies also cannot be replaced by a bound on the mixed
mean of $Z_D$. These are supplied by the lattice estimate of Section~\ref{ch18:lattice-products} and its passage to the compression algebra below.

% !TEX root = ../monograph.tex
\begin{proof}[Proof of Proposition~\ref{ch18:replacement-input}]
\label{ch18:raw-replacement}
The lattice estimate in Proposition~\ref{ch18:lattice-replacement}
was proved before phase decomposition. We use its uniformity in
replica number to obtain the algebra-valued statement for the
reconstructions $H_D$ fixed above.

Let $F$ be an interior lattice difference with chessboard seminorm
at most $B$, and let its exterior factors have interior seminorms
at most $B_1,\ldots,B_q$. Put copies of the complete difference
and exterior factors in separated cells before taking limits.
The repeated-word identity and chessboard inequality give, for its
compressed limit $Z$ and each integer $n\ge1$,
\[
 \int_\Xi|\xi(Z)|^{2n}\,d\nu(\xi)
                   \le(B\prod_{l=1}^qB_l)^{2n}.
\]
The seminorm bounds $B,B_1,\ldots,B_q$ are independent of $n$.
For each fixed $n$, the cutoff and the separation of the copies may
need to be increased.
This uses the complete signed density, with every term involving at least two insertions,
from the proof of \eqref{ch18:lattice-estimate}.
Equation~\eqref{ch15:moment-norm} then gives
$\|Z\|\le B\prod_lB_l$. Positivity here belongs to the limiting
commutative algebra and the original Wilson law, not to an
individual signed effective activity.

Apply this observation to the auxiliary field $A_t$ in the proof
of Proposition~\ref{ch18:mixed-bound}. On grid configurations its
replacement error has norm with exponent
$Z_0-(t+1/16)(\theta-\theta_-)$, where
$Z_0=W\theta_++4n_{\rm cl}(\theta_+-\theta)$.
Fixed exterior factors cost a bounded factor. The graded quotient
bound gives exponent $Z_0-i\theta$ for a degree-$i$ component of
$A_t$; all the finitely many fields used in its reconstruction
have uniformly bounded compressed exterior words. Testing the
same error against the fixed probes gives the probe coordinate
errors. With $t$ chosen as in that proof, the omitted degrees
and the low-coordinate errors therefore give
\[
 \|M(wv)-M((H_D(\tau_{-o}w))(o)v)\|
                \preccurlyeq e^{[Z_0-(D+1)\theta]T}.
\]
There is no spectator-dependent choice of $H_D$: the only
geometry-dependent auxiliary field has disappeared from this
inequality.

The derivative descent of Proposition~\ref{ch18:mixed-bound}
works in this norm. Equivalently, apply it to every character,
using the same sharp chessboard seminorm bounds, and take their supremum.
The crude and grid constants are uniform in the character.
A positional derivative changes its input degree by one; a finite
descending induction removes the grid restriction, also for any
fixed list of differentiated fields. Letting the two resolution
heights tend to one yields
\[
 \|M(wv)-M((H_D(\tau_{-o}w))(o)v)\|
                      \le C r^{D+1-W-\varepsilon}.
\]
The coefficient assertion is exactly the corresponding fixed-probe
estimate in Proposition~\ref{ch18:mixed-bound}. Compact label
sets, conjugation, and the finite probe families do not change
its constants. For their rotation orbits the same uniform estimates
apply on the compact set of rotated labels and positions; choose
$r$ so all these probes remain separated from the cluster.

Finally apply the vacuum state to $(Z_D^*Z_D)^n$. Its norm is
at most $\|Z_D\|^{2n}$, giving
\eqref{ch18:replacement-moments} with the same $C$ and the same
small-$r$ threshold for every $n$. Thus the order of quantifiers
is obtained from the norm estimate, rather than from an
uncontrolled sequence of separate moment constants.
\end{proof}


\section{Finite-order operator product expansion}
\label{ch18:ope}

\begin{theorem}[Phase-uniform finite-order expansion]
\label{ch18:ope-theorem}
The fixed coefficients
\eqref{ch18:coefficients} satisfy, for every finite $D$ and
$\epsilon>0$,
\begin{align}
 |c_{j,a}(w;o)|&\le C r^{j-w_0-\epsilon},\qquad j\le D,\notag\\
 \left\|M(wv)-\sum_{j\le D,a}c_{j,a}(w;o)
                        M(O_{j,a}(o)v)\right\|
      &\le C r^{D+1-w_0-\epsilon}.
 \label{ch18:ope-bound}
\end{align}
The constants and threshold are uniform on the bounded label sets and
separated geometries of Section~\ref{ch18:input}. Evaluating a character,
or the mixed vacuum functional, gives the same bound for its correlations.
\end{theorem}
\begin{proof}
Full support of the phase measure converts
\eqref{ch18:replacement-moments} into the operator-norm estimate by
\eqref{ch15:moment-norm}. To compare the unsymmetrized and symmetric reconstructions,
write $h=H_H(\tau_{-o}w)$ at an auxiliary finite degree $H$.
Every component of $B_jw-L_{j,H}h$ is an average of finitely many
replacement errors tested against rotated probe words. Their clearance
is uniform by compactness of $O(4)$. Hence
\[
 |B_j(\tau_{-o}w)-L_{j,H}h|
                         \le C r^{H+1-w_0-\epsilon}.
\]
This identity uses covariance of the mixed correlations to rotate the
probes, as well as the cluster; rotating labels alone would not suffice.

Apply $Q_j$. In the expansion of $h$, every field of degree below $j$
is killed. The remaining finitely many terms have powers at least
$r^{j-w_0-\epsilon}$, by the raw coefficient estimate. Taking $H\ge j$
proves the first bound. For the second, take $H=D$ and sum the reconstruction
comparison through $j=D$. The reconstruction identity following
\eqref{ch18:graded} returns $h$ exactly. Multiplying the coefficient
errors by $M(O_{j,a}(o)v)$ costs a uniform constant on the separated
spectator geometry. Add the raw operator error. Finitely many sums and
averages are controlled by first using a smaller error exponent.
\end{proof}

For two fields of degrees $d_A,d_B$, the degree-$j$ coefficient thus
has size at most $C r^{j-d_A-d_B-\epsilon}$. The remainder may still
grow when $D+1<d_A+d_B$; its order nevertheless improves with $D$.
No assertion concerns the convergence of the infinite sum in $D$.

The derivative version uses the differentiated marked replacement estimate. One first fixes the finite derivative order and the degree
margin, then carries out the same finite probe functionals. The derivative
identities of the field filtration identify the resulting derivatives.
Uniform differentiated bounds permit passage from a dense grid to all
separated positions by the mean-value theorem. Mere pointwise grid
estimates without those bounds would not justify that passage.

For two nested clusters at radii $\rho<R$, replacements can be composed
by telescoping the signed difference. The first error is tested against
the unchanged second cluster; the next is multiplied by the already
produced coefficients. Accordingly a uniform nested estimate requires
the replacement bound when exterior separation also tends to zero,
with its dependence on $R$ retained. Fixed-positive-clearance
Proposition~\ref{ch18:replacement-input} by itself does not give that
stronger result. For any prescribed finite tree, the required additional
statement is a signed bound at every vertex, with degree-$q$ coefficient
cost $\rho^{q-d}$ and residual gain $(\rho/R)^{D+1}$, up to arbitrary
small power losses, including the coefficients and the operators that its descendants apply to the test function. With those bounds, finite telescoping proves the
corresponding finite-tree expansion. This stronger multi-scale assertion
is not needed for the distributional extension below.

\section{An explicit finite-order extension across every diagonal}
\label{ch18:extension}

We now use only a uniform polynomial separated bound. Fix an arity
$n\ge2$ and a field list $\boldsymbol A$. Write
\[
 \Delta=\bigcup_{i<j}\{x_i=x_j\},\qquad
 \delta(x)=\min_{i<j}|x_i-x_j|.
\]
Suppose, as supplied by the separated bounds and the uniform phase slab
bounds, that for some finite $P\ge0$,
\begin{equation}
 |S_n^\xi(\boldsymbol A;x)|\le C_K\min(1,\delta(x))^{-P}
 \quad(x\in K\setminus\Delta),
 \label{ch18:polynomial}
\end{equation}
uniformly in $\xi$ on every compact $K$. The mixed family has the same
bound. No optimal value of $P$ is required. Choose once an integer
$k\ge0$ with $k+1>P$, and a length $s>0$.

We specify which normal jet is subtracted near simultaneous collisions.
Let $m=n(n-1)/2$ be the number of pairs and put, on $\Conf_n$,
\[
 \rho(x)=\left(\sum_{i<j}|x_i-x_j|^{-2}\right)^{-1/2}.
\]
Then $\delta/\sqrt m\le\rho\le\delta$.
Choose a fixed nonnegative smooth function $\varphi$ on $[0,\infty)$,
equal to one on $[0,2\sqrt m]$ and zero on $[4\sqrt m,\infty)$.
For a pair $h=\{i,j\}$ define
\begin{equation}
 a_h(x)=\frac{\varphi(|x_i-x_j|/\rho(x))}
                   {\sum_{u<v}\varphi(|x_u-x_v|/\rho(x))}.
 \label{ch18:pair-weights}
\end{equation}
The denominator is at least one, because a closest pair has ratio at
most $\sqrt m$. Thus $a_h\ge0$ and $\sum_h a_h=1$.
Only pairs whose separation is at most $4\sqrt m\rho$ are selected.
These weights are smooth off $\Delta$, invariant under simultaneous
translations and rotations, and permuted with the slots.

Let $\pi_hx$ replace $x_i,x_j$ by their centroid, leaving the other
positions unchanged, and put $z_h=x-\pi_hx$. This is the orthogonal
normal displacement from the pair diagonal; its norm is
$|x_i-x_j|/\sqrt2$. Define the normal Taylor polynomial of a test function by
\begin{equation}
 T_{h,k}f(x)=\sum_{\ell=0}^k\frac1{\ell!}
             D^\ell f(\pi_hx)[z_h,\ldots,z_h].
 \label{ch18:normal-Taylor}
\end{equation}
Only normal directions occur; the centroid and all external positions
remain parameters. Finally fix $\chi\in C^\infty([0,\infty))$ with
$0\le\chi\le1$, $\chi=1$ on $[0,1]$ and $\chi=0$ on $[2,\infty)$,
and set
\begin{equation}
 (\mathcal W f)(x)=f(x)-\chi(\rho(x)/s)
                                  \sum_h a_h(x)T_{h,k}f(x).
 \label{ch18:subtraction}
\end{equation}
This formula, including $k,s,\varphi,\chi$, is the subtraction convention.
It is the same in every phase and is fixed before the test function.
It deliberately subtracts more jets than an optimal scaling-degree
prescription may require.

\begin{theorem}[A common contact extension]
\label{ch18:extension-theorem}
Under \eqref{ch18:polynomial}, the formula
\begin{equation}
 \mathfrak D_{\boldsymbol A}^\xi(f)
   =\int_{\Conf_n} S_n^\xi(\boldsymbol A;x)
                                      (\mathcal Wf)(x)\,dx
 \label{ch18:extension-formula}
\end{equation}
defines a distribution of local order at most $k+1$.
It extends the separated correlation, is uniform in the phase on
compact supports of test functions, and is continuous in $\xi$ on each test function.
Averaging over phases gives the same formula for the mixed correlations.
Translations, orthogonal covariance and simultaneous permutation
symmetry are preserved by this prescription.
\end{theorem}
\begin{proof}
This is the normal Taylor extension argument of
\cite[Theorem~5.3]{BrunettiFredenhagen2000}; the following verification
explains why one subtraction treats every partial diagonal here.
When $\rho\le s$, the weights sum to one, so
\[
 \mathcal Wf=\sum_h a_h(f-T_{h,k}f).
\]
For each selected pair, $|z_h|\le 2\sqrt{2m}\rho$.
Taylor's integral remainder therefore bounds this expression by
$C\rho^{k+1}\|f\|_{C^{k+1}}$. Since $\delta\ge\rho$,
\eqref{ch18:polynomial} bounds its product with the kernel by
$C\rho^{k+1-P}\|f\|_{C^{k+1}}$, which is bounded and hence locally
integrable. This remains true when arbitrarily many other separations
vanish at independent rates: the selected pair always has displacement
comparable to the smallest separation. No fixed collision pattern or
single radial approach is assumed.

On the transition region $s\le\rho\le2s$, all separations are bounded
below by $s$, and the finite Taylor polynomials are bounded by the same
seminorm of the test function. Outside this region the integrand uses $f$ itself. If $f$ is supported
in a compact $K$, a nonzero subtracted jet requires $\pi_hx\in K$ and
$|z_h|\le C_ns$. Thus the integrand is supported in one fixed compact
enlargement $K'$ of $K$. We obtain
\[
 |\mathfrak D_{\boldsymbol A}^\xi(f)|
        \le C_{K'}\max_{|\alpha|\le k+1}
                              \|\partial^\alpha f\|_\infty.
\]
This proves continuity and the local order assertion. If $f$ is supported
away from $\Delta$, all its jets at every $\pi_hx\in\Delta$ vanish,
so \eqref{ch18:extension-formula} is the original separated integral.
Continuity in the phase and the mixture formula follow by dominated
convergence with the same bound. The normal projections, pair weights
and radial cutoff commute with the stated geometric symmetries.
\end{proof}

For $n=1$ use the original smooth one-point function, and for $n=0$
use the normalized unit. A consistent multilinear family on the whole
filtration is obtained using the fixed invariant summands $E_j$ in
\eqref{ch18:graded}. For each finite tuple $(j_1,\ldots,j_n)$ choose
$P$ uniformly on the finite basis lists of $E_{j_1}\times\cdots\times
E_{j_n}$, and choose the same $P,k$ for every permutation of that tuple.
Use this one subtraction on the tuple and extend multilinearly there.
Every field has a unique finite sum of components in the $E_j$, so
extension by finite multilinearity gives a prescription on all of $F$.
It is compatible with the inclusions and preserves rotations because
each $E_j$ is invariant. Choosing an order independently for each vector,
or redefining it on each larger degree space, would not have these
properties.

\section{Sharper subtraction and dependence on the prescription}
\label{ch18:scope}

The explicit extension above is conservative. A degree-sensitive
alternative assigns to every collision subset $h$ the normal dimension
$q_h=4(|h|-1)$, the total input degree $d_h=\sum_{i\in h}d_i$, and
\[
 L_h=\max(0,d_h+1-q_h).
\]
If a contact projection extracts a degree-$q$ field from that cluster,
one subtracts the normal Taylor jet of the test function through order $L_h-1-q$ when
$q<L_h$. A single shell then has the combined power
\[
 \rho^{q-d_h}\,\rho^{L_h-q}\,\rho^{q_h}
                  =\rho^{q_h+L_h-d_h},
\]
with exponent at least one. This explains the source prescription and
its proposed sharper local-order bound. For several collisions the
shell estimate must be established for the complete signed expression,
including nested subtractions and all child operators acting on test functions. The parent
Taylor operator acts first on the test function; the child operators act after
it. They cannot be commuted for nested subsets. A partition of unity
cannot be moved through these normal Taylor operators without its
resulting diagonal jets. A one-shell estimate alone does not verify
that stronger forest prescription. We have therefore used the fully
specified pair-weighted subtraction \eqref{ch18:subtraction}, for which
absolute convergence at all collision patterns was proved directly.

Two extensions agreeing on $\Conf_n$ differ by a distribution supported
on $\Delta$. Along a smooth stratum, a finite-order difference is a sum
of normal derivatives of delta distributions with distributional
coefficients. This is the usual local renormalization freedom
\cite[Section~5.2]{BrunettiFredenhagen2000}. Changing the cutoff or
subtraction order can change contact terms while preserving separated
correlations and the OPE estimates. Contact reflection positivity,
contact Ward identities, uniqueness, unit removal and identification
with unsubtracted coincident-point lattice limits require additional
normalization arguments.

\setcounter{chapter}{18}
% !TEX root = ../monograph.tex
\chapter{Assembly and the scope of the construction}
\label{ch19:assembly}

The preceding chapters distinguish finite algebraic identities, analytic
estimates for the lattice transformation, and consequences for the
continuum theory. This chapter states how they fit together. The construction and finite-order operator products use the estimates
proved in the preceding chapters. The sharper trace asymptotic is
separated below because it requires the additional current estimates
of Chapter~\ref{ch17:trace}.

\section{One family and the order of its limits}
\label{ch19:family}

Choose the compact matrix group $G$ with simple nonabelian real Lie
algebra, its trace inner product as in Section~\ref{ch01:gauge}, and a
positive soft-block penalty $b$. Fix a sufficiently large coarse inverse
coupling $\bar\beta$. The common schedule of
Chapter~\ref{ch10:trajectory} selects all block lengths, derivative
orders, support weights, coupling thresholds and admissible physical
periods $R_p$. At cutoff $N$ and period index $p$, shooting from the
Wilson measure gives its bare coupling $\beta_{N,p,N}$ with prescribed
coarse endpoint $\bar\beta$.

The same finite histories must define the coefficient normalizations
of every local insertion. One cannot separately choose convenient
histories for convergence, for free separating families and for the trace law.
The microscopic sections are fixed before $N,p$, and their comparison
maps act simultaneously on all labels and spectator correlation pairings. Any additional
finite stage requirements, such as the equation calibration floor,
belong to the selection of this one schedule.

For each fixed physical volume the first limit is the full limit
$N\to\infty$. Only afterwards is a subsequence of physical volumes
extracted. Given any increasing sequence of period indices, a single
refinement supplies all arities, degree bounds, label basis elements and
derivatives of the separated kernels. Auxiliary coefficient limit points
may be selected on ordered tails; they do not change this full inner
limit of the correlations. Scaling comparison sequences and the choice
of a phase are later operations on the same continuum family.

\section{Assembly of the construction}
\label{ch19:theorem-section}

The hypotheses used along the construction have now been supplied for
one common family. Chapters~\ref{ch09:renormalization} and
\ref{ch10:trajectory} give the actual block estimates and fixed-input
convergence. Chapter~\ref{ch11:insertions} supplies the marked comparisons,
including restarted histories and arbitrary fixed numbers of insertions.
Chapter~\ref{ch13:continuum} constructs the bounded approximations and
local product estimates. The free comparison and annular estimates in
Chapter~\ref{ch14:ward}, and the repeated estimate for a separating family in
Chapter~\ref{ch15:phases}, use these same objects. Thus they are not
additional assumptions in the following assembly statement.

\begin{theorem}[Assembly of the common Wilson family]
\label{ch19:conditional}
The chosen Wilson family has the following
properties. For every prescribed increasing volume sequence there is
one subsequence along which the cutoff limits converge to a hierarchy $S$ of smooth separated Euclidean
correlations and compatible finite-dimensional field spaces $F_d$.
Their dimensions and orthogonal representations agree with those of
the invariant free Maxwell comparison spaces. The filtration $(F_d)$ contains the unit field. Its inclusion maps are
injective, and its four coordinate derivatives commute.

There is a compact metrizable phase space $\Xi$, a full-support
probability measure $\nu$, and phase hierarchies $S^\xi$ with
$S=\int S^\xi\,d\nu$. All phases have the same field-null spaces,
Euclidean covariance and reflection positivity. Each phase reconstructs
a local relativistic quantum field theory with tempered Wightman
correlations, a faithful common field filtration, a positive-cone
translation spectrum and a unique invariant vacuum. Its weak field
commutant is scalar. These statements are uniform over the phases in
the finite-list senses specified in the preceding chapters.

The common real stress tensor is symmetric and conserved, with its
normalization fixed by local translation and rotation Ward identities.
Its relativistic continuation gives the local charge law
\eqref{ch16:charge-law}. The fixed OPE prescription of
Section~\ref{ch18:prescription} gives the phase-uniform finite-order
coefficient and remainder bounds \eqref{ch18:ope-bound}. Independently,
the fixed subtraction prescription
of Section~\ref{ch18:extension} extends all separated correlations to
finite-order distributions across collision diagonals, compatibly
with the field inclusions, translations, rotations and permutations.
\end{theorem}
\begin{proof}
Chapter~\ref{ch12:ultraviolet} takes the full cutoff
limit at each fixed period. The diagonal subsequence in the volume index works for every fixed
number of insertions and every scaling dimension. The two rank bounds
in Chapter~\ref{ch13:continuum} identify the dimensions of the resulting
field spaces, and its derivative maps give their differential calculus. That chapter also proves the repeated-word bound,
so bounded slab insertions in Chapter~\ref{ch14:ward} do not depend on
later reconstruction or operator products.

Coordinate reflection positivity, the free comparison and the annular
estimates give the conserved tensor, its local charges, Euclidean
rotations and Ward normalization in Chapter~\ref{ch14:ward}.
Chapter~\ref{ch15:phases} then compresses whole words to the mixed
ground space. The real approximation and moments for separating families prove
symmetries and the equality of field spaces in every character, beyond
what abstract commutative algebra alone gives.

The slab bounds supply joint boundary distributions in
Chapter~\ref{ch16:reconstruction}. Covariance gives the spectral cone,
analytic continuation gives locality, and word reconstruction gives
the relativistic fields. The unitary equivalence of the reconstructed and reflection-positive Hilbert spaces and the
phase-uniform low-energy estimate remove exceptional vacua. The same
chapter proves weak irreducibility and continues the stress charge law.

The finite OPE and contact extension follow from their separate arguments in
Chapter~\ref{ch18:products}. Each step preserves the originally chosen
family and the already selected correlation limits.
\end{proof}

\begin{corollary}[The trace law]
\label{ch19:conditional-trace}
In the phase $\xi_*$ of Section~\ref{ch15:selection}, let
$\Theta=\sum_aT_{aa}-S^{\xi_*}(\sum_aT_{aa})1$ be the centered trace and
$G_\Theta(r)=S_2^{\xi_*}(\Theta,\Theta;re_4,0)$ its two-point function.
Then
\[
 \lim_{r\downarrow0}r^8\log^2(1/r)G_\Theta(r)
                    =\frac{3\dim_\RR\gLie}{\pi^4}.
\]
Consequently the centered trace cannot agree near zero with any finite linear combination of products of the
mass-shell Laplace kernels specified in Section~\ref{ch17:exclusion}.
\end{corollary}
\begin{proof}
Apply the rate-sensitive tensor cancellation and prescribed-radius
squeeze of Chapter~\ref{ch17:trace}. The positive spectral measure
required there already follows from reconstruction.
\end{proof}

\section{Conclusions that are not supplied}
\label{ch19:scope}

The logarithmic ultraviolet law and the low-energy estimate describe
different regimes. The latter bounds spectral mass near zero by a
power; it does not show that the spectrum avoids an interval above the
vacuum. A positive mass gap is not proved.

Vacuum uniqueness within each phase does not imply uniqueness of the
phase, a point-mass phase measure, or uniqueness of the infinite-volume
limit. The construction uses a refinement of a prescribed volume
sequence. It does not prove a simultaneous continuum and infinite-volume
limit, independence of the block penalty or all normalization choices,
or a separate lattice realization of each character phase.

The trace comparison excludes the finite class explicitly defined in
Chapter~\ref{ch17:trace}. It does not rule out every generalized free
field or every infinite Wick expansion on the basis of one two-point
function. The OPE is an expansion to each fixed finite degree; no
infinite series convergence has been asserted. The contact extension
has a specified subtraction convention and may have larger local order
than an optimized prescription. Positivity at contacts, contact Ward
identities, the fact that inserting the identity leaves correlations unchanged, and uniqueness of contact
terms require additional normalization conditions.

The trace argument uses more than convergence to the free tensor. Its
low-moment estimate distinguishes field degree from scaling dimension:
a quadratic expression with extra spatial differences is still quadratic
in the field. Exact polynomial preservation by the linear response,
together with the normalized Gaussian fluctuation estimate, supplies the
required decay in Proposition~\ref{ch17:spatial-moments}. This is why
that additional calculation is needed after reconstruction.

\backmatter
\chapter{Notation in the current draft}

This index lists the main symbols used in the written chapters.
Auxiliary symbols are defined where used. The chapter and equation
references give the precise conventions; in particular, lattice distances
and physical distances must not be interchanged.

{\small
\begin{longtable}{@{}p{.22\textwidth}p{.73\textwidth}@{}}
\toprule
Symbol & Meaning and definition \\
\midrule
\endfirsthead
\toprule
Symbol & Meaning and definition \\
\midrule
\endhead
$G,\gLie,n_G$ & Compact matrix gauge group, its real Lie algebra and matrix size; Section~\ref{ch01:gauge}.\\[2pt]
$\|\cdot\|_G$ & Normalized matrix norm, $\|A\|_G^2=n_G^{-1}\tr(A^*A)$; \eqref{ch01:matrix-norm}.\\[2pt]
$C_{\gLie},\kappa_G$ & Adjoint Casimir and inverse-coupling logarithmic coefficient; \eqref{ch01:casimir}--\eqref{ch01:slope}.\\[2pt]
$\Lambda_s,E_s$ & Periodic lattice with $s$ sites per direction and its positive links; Section~\ref{ch01:wilson-section}.\\[2pt]
$U_e,U_{y;\mu\nu}$ & Link variable and oriented plaquette holonomy; \eqref{ch01:plaquette}.\\[2pt]
$E_s^{\mathrm W}$ & Wilson energy, the sum of $\tfrac12\|U_{y;\mu\nu}-1\|_G^2$ over plaquettes; \eqref{ch01:wilson-energy}.\\[2pt]
$Z_{s,\beta},\mu_{s,\beta}$ & Partition function and Wilson probability measure; Definition~\ref{ch01:wilson}.\\[2pt]
$R_p,L_m$ & Physical period and odd block ratio, respectively; Definition~\ref{ch01:scales}.\\[2pt]
$D_m,a_m,\tau_m$ & Cumulative lattice denominator, spacing $D_m^{-1}$ and logarithmic scale $\log D_m$; \eqref{ch01:scale-sequence}.\\[2pt]
$N,m,p$ & Ultraviolet cutoff, renormalization level and physical-volume index; Section~\ref{ch01:scales-section}.\\[2pt]
$\Conf_n$ & Configurations of $n$ distinct points in $\RR^4$; Definition~\ref{ch01:separated-test}.\\[2pt]
$V_d$ & Invariant radial-jet polynomials of scaling dimension at most $d$; \eqref{ch02:label-space}.\\[2pt]
$N_{\rm fac}$ & Polynomial factor-count operator, distinct from scaling dimension; \eqref{ch02:factor-scaling}.\\[2pt]
$P_d^{\rm loc}$ & Finite-stencil representative of a local polynomial; Theorem~\ref{ch02:local-sections}.\\[2pt]
$S_n^{\rm Max}$ & Wick correlations in the free Maxwell comparison theory; Section~\ref{ch03:wick}.\\[2pt]
$\mathcal S_j,\langle\cdot,\cdot\rangle_j$ & Homogeneous degree-$j$ Maxwell solution jets (Section~\ref{ch03:shell}) and their coefficient inner product (Section~\ref{ch03:strict}).\\[2pt]
$\omega_d$ & Restriction from radial-potential labels to free solution jets; Section~\ref{ch03:shell}.\\[2pt]
$b$ & Positive penalty in the soft block constraint, distinct from inverse coupling $\beta$; Section~\ref{ch04:kernel-section}.\\[2pt]
$\mathcal B^{(L)}_\beta$ & Exact block transformation from $\Lambda_{Ls}$ to $\Lambda_s$ at ratio $L$; \eqref{ch04:block-operator}.\\[2pt]
$\widehat\beta_m,\mathfrak B_m$ & Reference inverse coupling and the permitted interval around it; Section~\ref{ch04:density-coordinates}.\\[2pt]
$M_m^{\rm cut},t_m,B_m$ & Cutoff amplitude, plaquette threshold and radius used in support weights; \eqref{ch04:thresholds}.\\[2pt]
$E_i,X_i,o_i$ & Regular activity, its support and anchor; Section~\ref{ch04:activities}.\\[2pt]
$K_i,Y_i,I_i$ & Large-field activity, exclusion support and prescribed set of large-field sites; Section~\ref{ch04:activities}. The full link-dependence region can be larger than $Y_i$.\\[2pt]
$\mathcal E_m,\mathcal K_m$ & Weighted regular and large-field activity sums; \eqref{ch04:weighted-sums}.\\[2pt]
$\Omega^r,d_r$ & Lattice $r$-forms and forward exterior derivative; Section~\ref{ch05:forms}.\\[2pt]
$H,K$ & Curvature metric and kinetic operator $d_1^*Hd_1$; \eqref{ch05:kinetic}. These are operators, not polymer activities.\\[2pt]
$Q,A^{-1},W$ & Linear map from fine to coarse fields, covariance of the normalized fine fluctuations, and map from a coarse field to its minimizing fine field; \eqref{ch05:schur-data}.\\[2pt]
$T$ & Two-form transport associated with $W$; \eqref{ch05:physical-transports}.\\
$\mathcal Q,W_a,W_F$ & Linearized gauge-covariant block average and two representatives of the minimizing fine field, respectively in the smooth gauge and in the forest coordinates; Section~\ref{ch05:actual-tangent}.\\[2pt]
$\mathscr Z_{\rm blk}$ & Exact blocked density before its local reorganization; Section~\ref{ch09:integral}.\\[2pt]
$J(m),\mathscr L_j$ & Stage at level $m$ and odd block ratio selected for stage $j$; Section~\ref{ch10:schedule}.\\[2pt]
$\bar\beta,\beta_{N,p,m}$ & Prescribed coarse coupling and actual coupling along the finite history; Theorem~\ref{ch10:finite-history}.\\[2pt]
$z_i,\mathcal A_{\mathcal I}$ & Commuting nilpotent source variables and their finite coefficient algebra; Chapter~\ref{ch11:insertions}.\\[2pt]
$q_m$ & Density with coefficients indexed by subsets of insertions in Chapter~\ref{ch11:insertions}; \eqref{ch11:source-density}. In Chapter~\ref{ch10:trajectory}, $q_m$ instead denotes the local scalar coupling correction.\\[2pt]
$M^{\rm UV}_{p,v}$ & Smooth fixed-volume correlation after the full ultraviolet limit; Theorem~\ref{ch12:uv-theorem}.\\[2pt]
$S_n(v;x)$ & Infinite-volume separated correlation along the selected volume subsequence; \eqref{ch13:volume-limit}.\\[2pt]
$N_d(S),F_d$ & Polynomials that vanish in every separated correlation, and the quotient by this null space; \eqref{ch13:null}.\\[2pt]
$\Pi_d$ & Limiting corrected coefficient projection; Proposition~\ref{ch13:rank}.\\[2pt]
$\iota_{d,e},\partial_\mu$ & Inclusions of filtered field spaces and derivatives with respect to an insertion point; Section~\ref{ch13:calculus}.\\[2pt]
$\theta w,w^\dagger$ & Antilinearly reflected word and its translated reflected word; Sections~\ref{ch13:positivity} and~\ref{ch13:repeated}.\\
$\mathcal H_\mu^{\rm OS},E_\mu(t),P_\mu^0$ & Coordinate reflection space, heat semigroup and zero-energy projection; Section~\ref{ch14:reflection}.\\[2pt]
$B_h(w)$ & Bounded insertion of a word in a slab of width $h$; Proposition~\ref{ch14:slab-bound}.\\[2pt]
$J,T$ & Conserved tensor and its normalization by the Ward identities; \eqref{ch14:conserved-line} and Proposition~\ref{ch14:normalized-stress}.\\[2pt]
$M(w),\mathfrak A$ & Insertion operator for a word, compressed to the ground space, and the commutative algebra generated as the word varies; Section~\ref{ch15:compression}.\\[2pt]
$\Xi,\nu,S^\xi$ & Character space, full-support representing measure and phase correlations; Section~\ref{ch15:characters}.\\[2pt]
$\xi_*$ & Retained phase with a one-dimensional ground space; Section~\ref{ch15:selection}.\\
$\mathcal H_\xi,\mathcal D_\xi,\Phi_\xi$ & Reconstructed Hilbert space, common word domain and smeared fields; Section~\ref{ch16:fields-section}.\\[2pt]
$\mathsf F_n^\xi,\mathcal W_n^\xi,W_n^\xi$ & Ordered-time holomorphic function, full tube function and joint boundary distribution; \eqref{ch16:heat-function}, \eqref{ch16:full-tube} and Proposition~\ref{ch16:boundary}.\\[2pt]
$\Theta,G_\Theta,\rho_\Theta$ & Centered trace, its Euclidean radial two-point function and positive spectral measure in the retained phase; Section~\ref{ch17:spectrum}.\\[2pt]
$E_j,Q_j,O_{j,a}$ & Complement of the fields of lower scaling dimension in $F_j$, projection onto this complement, and its chosen basis for the operator product expansion; \eqref{ch18:graded}.\\[2pt]
$c_{j,a}(w;o)$ & Coefficient for a cluster word $w$ about center $o$ in the fixed prescription; \eqref{ch18:coefficients}.\\[2pt]
$\mathfrak D_{\boldsymbol A}^\xi$ & Distributional contact extension with specified normal Taylor subtraction; \eqref{ch18:extension-formula}.\\
\bottomrule
\end{longtable}
}

\bibliographystyle{alpha}
\bibliography{references}
\end{document}
